% VOLUME I - CHAPTER 2 FINAL MASTERPIECE PASS (2026-06-14)
% Rules 1-15: classical permanence; math, operators, spaces frozen.

\chapter{Maxwell Operators and Radiation Operators}
\label{ch:02}
\setcounter{section}{-1}

\section{Introduction}
\label{sec:ch20}

Antenna theory, inverse scattering, near-field metrology, and radiation synthesis all reduce,
at fixed frequency, to two engineering questions that recur on every range and in every
synthesis code \cite{harrington2001}: \emph{What field can this source produce?}
and \emph{What source produced this field?} The first question governs forward design---a horn
mouth, patch feed, or array excitation is specified and the far-zone pattern on a large sphere
$\Gamma_R$ is sought \cite{stratton1941}. The second governs diagnosis and inversion---probes
on a scanner record tangential fields and one asks which impressed currents could have generated
the observed export \cite{devaney2004}. Maxwell's equations answer both questions in
principle; the purpose of Chapter~2 is to organize those answers so that laboratory
measurements, reactive storage, and non-radiating degrees of freedom are not conflated
\cite{maxwell1873}.
% PHYSICAL-HISTORY:sec:ch20

Maxwell's equations fixed the local balances from which every antenna measurement still begins \cite{maxwell1873}. Poynting's identification of directed energy flux made the wattmeter on an enclosing sphere a primary instrument rather than a post-processing afterthought \cite{poynting1884}. Stratton's treatment of the radiation condition and Green-function propagation turned open-region horn and scatterer problems into the boundary-value laboratory practice that compact ranges still emulate \cite{stratton1941}. Chu later sharpened the same finite-support picture into a counting question: electrical size fixes how much angular structure a bounded radiator can export before bandwidth and stored energy dominate \cite{chu1948}. Harrington's time-harmonic field theory and method-of-moments habit carried those limits into the discrete source-to-pattern maps laboratories use today \cite{harrington2001}. The present chapter stands in that measurement line: it asks which Maxwell states bounded impressed currents export after outgoing selection, which remain in reactive storage, and which finite instruments can resolve on a declared far shell \cite{stratton1941}. The governing thesis, carried forward from Volume~I, is that \emph{finite geometry limits realizable, accessible, and observable angular electromagnetic complexity} \cite{pal2026}. Forward design, inverse scattering, and near-field metrology differ in hardware, yet they share the same historical obligation---declare support, enforce outgoing export, then report watts and patterns only within the rank the outline and probe count permit \cite{harrington2001}.

% PHYSICAL-MEMORY:sec:ch20

\begin{physicalvignette}[Pyramidal horn on a compact range]
Consider Pyramidal horn on a compact range. Before any coefficient table is printed, technicians record frequency, cable loss, and the mechanical outline that fixes electrical size. The first walk-through of the range is always tactile: flange bolts, probe standoff, turntable boresight, and the wattmeter reference on the receive chain \cite{hansen1988}.

On this hardware the near zone is visually impressive---standing reactive structure, rapid phase rotation, sometimes a visible VSWR ripple at the feed---yet the export question is measured only after outgoing selection on a large sphere \cite{harrington2001}. What the camera sees in the chamber is not automatically what the far-zone pattern claims \cite{stratton1941}.

Engineers therefore learn to separate three readings: what the feed can synthesize on declared support, what propagating content survives on the radiative range, and what a finite instrument can recover at declared tolerance \cite{devaney2004}. Confusing those readings produces the rank gaps and watt failures that later sections quantify \cite{pal2026}.

Rayleigh and Stratton treated such ranges as boundary-value laboratories long before digital near-field transforms became routine \cite{stratton1941}. Harrington's method-of-moments culture reinforced the same discipline: metal outline first, harmonic headline second \cite{harrington2001}.

Subsection~\ref{sec:ch20} states the formal export condition that this range scene presupposes \cite{stratton1941}.
\end{physicalvignette}

The governing idea of \ref{ch:02} is \emph{radiation as accessibility}: which Maxwell states
generated by bounded sources remain observable after outgoing selection on a declared far-zone
shell $\Gamma_R$ \cite{maxwell1873}. The radiation operator reports what finite
sources export once Sommerfeld-outgoing qualification is enforced; it does not create angular
structure \cite{maxwell1873}.

At fixed frequency, Maxwell's equations define a source-to-field operator
\(\mathcal{S}_\omega\)---the harmonic \emph{solution map} that sends impressed
current \(\Jimp\) to the field pair \((\E,\H)\) consistent with sources and
boundary termination---and, after outgoing selection, a radiation map
\(\mathcal{R}_\omega=\Pi_{\mathrm{rad}}\circ\mathcal{S}_\omega\), where
\(\Pi_{\mathrm{rad}}\) is the outgoing selector that retains only far-zone
transverse content \cite{maxwell1873}. Green's retarded potentials and
Sommerfeld's radiation condition made the exterior problem precise long before
modern mathematical theory was applied to antennas and scattering
\cite{stratton1941}. Hertz \cite{hertz1888} demonstrated that Maxwell's equations predict propagating electromagnetic waves in the laboratory. The scientific question of this chapter is
classical: which field content generated by finite impressed currents
\(\Jimp\in\mathcal{J}_\Omega\) is exported to the far zone, which remains in
reactive storage, and which lies in the null space of \(\mathcal{R}_\omega\)?
\ref{ch:01} fixed Maxwell consistency, conservation laws, gauge limits on
local readouts, and the measurement sequence; here those results become the radiation-selection maps \cite{stratton1941}.
No new electromagnetic law is introduced in this organization;
the maps rename consequences of Maxwell theory together with outgoing selection and finite
instrumentation \cite{maxwell1873}.

\phantomsection\label{prin:ch2.radiative-accessibility}
\emph{Radiative probe rank.} Far-zone observability is a composed map: radiation selection, geometric regime, and finite measurement rank act in sequence. Evanescent and reactive content may dominate near scans without enlarging the angular rank-$N$ recovery of export on $\mathcal{F}_{\mathrm{rad}}$ \cite{stratton1941}.

\begin{figure}[t]
\centering
\ChOneFigBlock{%
\begin{tikzpicture}[ch1fig, node distance=5mm and 7mm]
 \node[ch1box, minimum width=28mm] (src) {Finite\\sources};
 \node[ch1box, minimum width=28mm, below=of src] (max) {Maxwell\\operator};
 \node[ch1box, minimum width=28mm, below=of max] (grn) {Green\\propagation};
 \node[ch1box, minimum width=28mm, below=of grn] (bnd) {Boundary\\closure};
 \node[ch1box, minimum width=28mm, below=of bnd] (out) {Outgoing\\selection};
 \node[ch1box, minimum width=28mm, below=of out] (rad) {Radiation\\operator};
 \node[ch1box, minimum width=28mm, below=of rad] (meas) {Measurement\\operator};
 \node[ch1box, minimum width=28mm, below=of meas] (ang) {Observable\\angular structure};
 \draw[ch1flow] (src) -- (max) -- (grn) -- (bnd) -- (out) -- (rad) -- (meas) -- (ang);
\end{tikzpicture}%
}
\caption{Finite impressed currents excite Maxwell evolution; Green propagation and boundary completion select a physical field branch; outgoing selection and the radiation map $\mathcal{R}_\omega$---the chapter's central link---export far-zone structure; finite measurement rank compresses what laboratories may report \cite{maxwell1873}.}
\label{fig:ch2.master-spine}
\end{figure}

\ref{ch:01} established admissibility; \ref{ch:02} establishes $\mathcal{E}_{\mathrm{acc}}^{(N)}$---which Maxwell-generated content survives outgoing radiation condition and reaches the far zone \cite{maxwell1873}. Maxwell field specification fixes what can exist on a declared domain; radiation fixes what leaves that domain under outgoing selection \cite{maxwell1873}. A cavity mode and a compact-range pattern may share the same frequency yet differ completely in export; the distinction is not visible in a single near-field snapshot \cite{harrington2001}.

Consider a shorted coaxial stub and an open-ended dipole at $1\,\mathrm{GHz}$ in the same
anechoic cell. Only the radiating structure supports watt completion on $\Gamma_R$ once
outgoing qualification is enforced; the stub stores reactive energy without net export
\cite{stratton1941}. That contrast is the operational content of the operator split developed
below: reactive storage is diagnosed on $\mathcal{S}_\omega[\Jimp]$, while far-zone reporting
attaches to $\mathcal{R}_\omega[\Jimp]$. In laboratory measurements the distinction is
immediately visible. A network analyzer may report substantial stored energy and strong local
fields in both structures, yet only the dipole produces sustained outward power flow and a
measurable far-field pattern \cite{harrington2001}. Radiation operators express
this experimentally familiar distinction \cite{maxwell1873}.

Source specification, Maxwell completion, radiation selection, observation, and inference form
one chain, not a sequence of unrelated topics \cite{maxwell1873}. Not every Maxwell field
radiates; not every radiating field is observable; not every observable field identifies its
source \cite{stratton1941}. The three subsections of Section~\ref{sec:ch20} develop that
chain in order: Subsection~\ref{sec:ch201} names the maps and their domains;
Subsection~\ref{sec:ch202} packages Maxwell's equations as the operator whose inversion
produces \(\mathcal{S}_\omega\); Subsection~\ref{sec:ch203} records continuity with the
$\mathcal{D}_J=1$ and measurement criteria of \ref{ch:01}
\cite{stratton1941}. Section~\ref{sec:ch21} then supplies the function spaces on which
those maps act with mathematical precision.

\subsection{Radiation as a source-to-field operator}
\label{sec:ch201}

% CH02-NARR-EXP:02.201
A network analyzer $|S_{11}|$, a spherical near-field scan on $\Gamma_R$, and a far-zone pattern cut are three readings on one ordered map: impressed data, Maxwell closure, outgoing export, finite observation \cite{maxwell1873,stratton1941}. Misplacing a measurement on that chain---reporting reactive storage as radiated power---remains the common compact-range failure mode \cite{poynting1884,hansen1988}.


Section~\ref{sec:ch20} posed radiation as a question of accessibility rather than of
field existence alone. The laboratory scene below makes that distinction operational by naming
the maps that connect impressed currents to Maxwell field pairs and, through outgoing
selection, to far-zone observables \cite{stratton1941}. The distinction between field
existence and field $\mathcal{E}_{\mathrm{acc}}^{(N)}$ is central to radiation theory \cite{stratton1941}.
Maxwell's equations determine which fields may arise from a specified source distribution,
whereas the radiation map $\mathcal{R}_\omega$ determines which portion of those fields
survives as observable outgoing content \cite{maxwell1873}. From an antenna
engineer's perspective, $\mathcal{S}_\omega$ answers what the structure \emph{stores and
circulates} locally, while $\mathcal{R}_\omega$ answers what the range may \emph{meter and
pattern} after the reactive shell is left behind \cite{harrington2001}.
% PHYSICAL-MEMORY:sec:ch201

\begin{physicalvignette}[Parabolic dish on a three-axis turntable]
Consider Parabolic dish on a three-axis turntable. Before any coefficient table is printed, technicians record frequency, cable loss, and the mechanical outline that fixes electrical size. The first walk-through of the range is always tactile: flange bolts, probe standoff, turntable boresight, and the wattmeter reference on the receive chain \cite{hansen1988}.

On this hardware the near zone is visually impressive---standing reactive structure, rapid phase rotation, sometimes a visible VSWR ripple at the feed---yet the export question is measured only after outgoing selection on a large sphere \cite{harrington2001}. What the camera sees in the chamber is not automatically what the far-zone pattern claims \cite{stratton1941}.

Engineers therefore learn to separate three readings: what the feed can synthesize on declared support, what propagating content survives on the radiative range, and what a finite instrument can recover at declared tolerance \cite{devaney2004}. Confusing those readings produces the rank gaps and watt failures that later sections quantify \cite{pal2026}.

the prior chapter asks which source-generated structure survives Maxwell propagation, outgoing selection, and finite measurement . 

Rayleigh and Stratton treated such ranges as boundary-value laboratories long before digital near-field transforms became routine \cite{stratton1941}. Harrington's method-of-moments culture reinforced the same discipline: metal outline first, harmonic headline second \cite{harrington2001}.

Subsection~\ref{sec:ch201} states the formal export condition that this range scene presupposes \cite{stratton1941}.
\end{physicalvignette}

A network analyzer $|S_{11}|$, a spherical near-field scan on $\Gamma_R$, and a far-zone pattern cut are three readings on one ordered map: impressed data, Maxwell field specification, outgoing export, finite observation \cite{maxwell1873}.
Each instrument samples a different stage of that map, and confusion between stages---for
example, reporting near-field reactive content as radiated power---remains the common
compact-range laboratory error \cite{poynting1884}.

Laboratory measurements presuppose an excitation model even when it is left implicit: a port
drive, a probe sample, or a termination fixes how energy is partitioned between storage and
escape \cite{stratton1941}. Impressed data from Section~\ref{sec:ch01.1.2} map through Maxwell
field specification to \((\E,\H)\), and only then is the outward-propagating fraction
identified as the carrier of far-zone power and angular flux
(Eq.~\eqref{eq:ch1.maxwell-local}; Sections~\ref{sec:ch01.1.4}--\ref{sec:ch01.1.5}). At fixed
$\omega$ the maps are the harmonic solution operator $\mathcal{S}_\omega$ and radiation
operator $\mathcal{R}_\omega$ under the \(\exp(-\jj\omega t)\) convention of
Chapter~\ref{ch:01}.

Let \(\Omega\subset\mathbb{R}^{3}\) be a piecewise smooth domain with boundary
\(\Gamma=\partial\Omega\). Phasor fields \(\E(\mathbf{r})\) and \(\H(\mathbf{r})\)
are suppressed in time unless causality is restored below. Impressed electric
current \(\Jimp\), charge density \(\rho\), and the data tuple \(\mathfrak{D}\) of
Eq.~\eqref{eq:ch1.data-specification-tuple} are not redefined; they are the
inputs on which the maps act. Finite-source admissibility is assumed in the sense of the
\ref{ch:01} $\mathcal{D}_J=1$ conditions
(Eq.~\eqref{eq:ch1.continuity-admissibility}): only configurations with
\(\mathcal{D}_{J}=1\) enter the domain of the radiation chain developed here
\cite{stratton1941}.

\begin{figure}[t]
\centering
\VolIFigInclude{ch2-source-radiation-sequence}
\caption{Excitation-to-observable chain: \(\mathcal{S}_\omega\) enforces Maxwell
closure with inherited boundary data; \(\Pi_{\mathrm{rad}}\) isolates the
radiative component on which far-zone flux integrals are meaningful.}
\label{fig:ch2.source-radiation-sequence}
\end{figure}

Function spaces are declared before the maps. Let $\mathcal{F}_\Omega$ be the
space of phasor pairs \((\E,\H)\) on \(\Omega\) whose components have
\(L^{2}\)-integrable curls and divergences in the sense required by
Eq.~\eqref{eq:ch1.maxwell-operator-block}, and whose boundary traces lie in the
domain of \(\mathcal{B}_\Gamma\). Let \(\mathcal{J}_\Omega\subset L^{2}(\Omega)^{3}\)
be the impressed-current space compatible with finite-energy class
(Section~\ref{sec:ch01.0.3}), charge continuity
(Eq.~\eqref{eq:ch1.charge-continuity-harmonic}), and the same boundary class.
Refined Sobolev and distributional extensions of these spaces are supplied in
Section~\ref{sec:ch211}; here the declaration is minimal but precise enough to
state the maps \cite{stratton1941}. The harmonic Maxwell \emph{solution map} is the
first quantitative object in the radiation chain: it records, for each admissible impressed
current \(\Jimp\in\mathcal{J}_\Omega\), the Maxwell field pair that the governing
equations and boundary data jointly permit. Physically, this map is what a full-wave solver
computes when it inverts Maxwell's equations on a declared domain; the output is a \emph{complete} field state, not yet a radiative report
\cite{stratton1941}. Experimentally one would observe, on a closed measurement surface
enclosing the source, both reactive circulation and outward-propagating flux in the
solution unless outgoing selection is applied afterward \cite{harrington2001}. The solution
map sends impressed data to a field state:
\begin{equation}
\label{eq:ch2.maxwell-solution-map}
\mathcal{S}_\omega:\mathcal{J}_\Omega\longrightarrow\mathcal{F}_\Omega,
\qquad
\Jimp\longmapsto(\E,\H),
\end{equation}
defined by \(\mathcal{S}_\omega[\Jimp]=(\E,\H)\) satisfying the Maxwell operator
residuals of Eq.~\eqref{eq:ch1.maxwell-operator-block} with sources \(\mathfrak{D}\)
and boundary data \(\mathbf{g}_\Gamma\) \cite{stratton1941}.
Equation~\eqref{eq:ch2.maxwell-solution-map} is the \ref{ch:02} statement
of the initial-boundary value problem proven in Section~\ref{sec:ch01.1.1}; it
does not alter the law, only the input--output viewpoint.
Injectivity fails in general: distinct \(\Jimp\) can produce identical \((\E,\H)\)
on observational subdomains, and equivalent source distributions related by
divergence-free reparametrizations inside \(\Omega\) can leave exterior patterns
unchanged when \(\mathcal{B}_\Gamma\) permits reactive storage
(Section~\ref{sec:ch01.2.4}). Those equivalences are not nuisances to be removed
by definition; they foreshadow null-space and realizability analyses in
Sections~\ref{sec:ch27} and \ref{ch:04}. Radiation accounting therefore
never begins from a guessed modal field; it begins from \(\mathcal{S}_\omega[\Jimp]\)
with explicit completion.
The full state typically contains both stored and exported electromagnetic energy.
Reactive fields exchange power with sources and terminations at distances small
compared with \(k_0^{-1}\); radiative fields carry net outward flux at infinity
in open configurations. Let \(\mathcal{F}_{\mathrm{rad}}\subset\mathcal{F}_\Omega\)
denote pairs whose far-zone asymptotics satisfy outgoing selection (made analytic in
Section~\ref{sec:ch24}). Radiation accounting therefore requires a second map that
discards non-outgoing storage before pattern or transport claims are made. The outgoing
selector \(\Pi_{\mathrm{rad}}:\mathcal{F}_\Omega\to\mathcal{F}_{\mathrm{rad}}\) performs
that discard: it accepts a completed Maxwell pair and returns only the component whose
leading far-zone asymptotics satisfy Sommerfeld-outgoing qualification
\cite{stratton1941}. From a measurement standpoint, \(\Pi_{\mathrm{rad}}\) is the
mathematical counterpart of moving a probe beyond the reactive shell and reporting only
the transverse outgoing field that persists with \(1/r\) decay \cite{harrington2001}.
Composing source excitation with that selector defines the \emph{radiation operator}
\begin{equation}
\label{eq:ch2.radiation-operator-def}
\mathcal{R}_\omega:\mathcal{J}_\Omega\longrightarrow\mathcal{F}_{\mathrm{rad}},
\qquad
\mathcal{R}_\omega[\Jimp]=\Pi_{\mathrm{rad}}\,\mathcal{S}_\omega[\Jimp].
\end{equation}
Equation~\eqref{eq:ch2.radiation-operator-def} composes source excitation with
outgoing selection: \(\mathcal{R}_\omega\) accepts the same impressed currents as
\(\mathcal{S}_\omega\) but returns only the radiative image
\(\Pi_{\mathrm{rad}}(\E,\H)\in\mathcal{F}_{\mathrm{rad}}\), where
\(\mathcal{F}_{\mathrm{rad}}\) is the subspace of field pairs whose leading
far-zone behavior is outgoing and transverse.
The composition law records that radiation is never computed independently of
Maxwell completion:
\begin{equation}
\label{eq:ch2.operator-composition}
\mathcal{R}_\omega
=
\Pi_{\mathrm{rad}}\circ\mathcal{S}_\omega,
\qquad
\mathcal{D}_{J}=1
\Longrightarrow
\mathcal{S}_\omega:\mathcal{J}_\Omega\to\mathcal{F}_\Omega,
\quad
\mathcal{R}_\omega:\mathcal{J}_\Omega\to\mathcal{F}_{\mathrm{rad}}.
\end{equation}
where $\mathcal{S}_\omega$ is the harmonic Maxwell solution map, $\Pi_{\mathrm{rad}}$ the
outgoing selector, $\mathcal{R}_\omega$ the radiation map, $\mathcal{J}_\Omega$ the impressed-current
space, $\mathcal{F}_\Omega$ the admissible field space, $\mathcal{F}_{\mathrm{rad}}$ the radiative
subspace, and $\mathcal{D}_J=1$ the transport qualification inherited from
\ref{ch:01} \cite{stratton1941}. Equation~\eqref{eq:ch2.operator-composition} composes $\mathcal{S}_\omega$ with outgoing
selection $\Pi_{\mathrm{rad}}$ on $\mathcal{J}_\Omega$ \cite{stratton1941}. On a ball
$B_a$ with $k_0 a=0.5$, reactive-to-radiative field-energy ratios on $S_R$ with $R=2a$
typically exceed unity by an order of magnitude; applying pattern or transport metrics to
$\mathcal{S}_\omega[\Jimp]$ without $\Pi_{\mathrm{rad}}$ therefore overstates export by
roughly $10\times$ in magnitude bookkeeping before phase errors accumulate
\cite{jackson1999}. Superposition of radiative images is legitimate only after each
constituent satisfies Eq.~\eqref{eq:ch2.source-charge-continuous} and outgoing radiation
condition \cite{harrington2001}. Until Subsection~\ref{sec:ch224} fixes
$\Pi_{\mathrm{rad}}=\mathcal{O}_{\mathrm{out}}\circ\Pi_{\mathrm{prop}}$, treat
$\Pi_{\mathrm{rad}}$ as an abstract outgoing selector \cite{stratton1941}. Angular-flux
transport from Sections~\ref{sec:ch01.2.3}--\ref{sec:ch01.2.4} and structured-beam
interpretations in Section~\ref{sec:ch01.6.2} attach to $\mathcal{R}_\omega[\Jimp]$;
coefficients read from reactive components of $\mathcal{S}_\omega[\Jimp]$ describe
near-zone storage, not exported flux \cite{allen1992}.
When Green operators are available (Section~\ref{sec:ch23}), \(\mathcal{R}_\omega\)
admits a retarded integral representation that makes the propagation step explicit.
The integral representation below is the frequency-domain statement of what a
method-of-moments code evaluates when it integrates the dyadic Green function against
impressed currents: each source element at \(\mathbf{r}'\) contributes to the observation
point \(\mathbf{r}\) with phase and amplitude fixed by outgoing propagation
\cite{stratton1941}. Physically, the integral superposes the radiative
images of all current elements; only the portion of \(\mathcal{S}_\omega[\Jimp]\) that
\(\Pi_{\mathrm{rad}}\) retains enters the integrand in the completed theory developed in
Section~\ref{sec:ch23}. For \(\mathbf{r}\notin\Omega_s\) with
\(\operatorname{supp}\Jimp=\Omega_s\Subset\Omega\),
\begin{equation}
\label{eq:ch2.radiation-integral}
\E_{\mathrm{rad}}(\mathbf{r})
=
-\jj\omega\mu_0
\int_{\Omega_s}
\boldsymbol{\Gamma}(\mathbf{r},\mathbf{r}')\cdot\Jimp(\mathbf{r}')\,d^{3}r',
\end{equation}
In Eq.~\eqref{eq:ch2.radiation-integral},
\(\E_{\mathrm{rad}}(\mathbf{r})\) is the radiative electric phasor at
\(\mathbf{r}\) (V\,m\(^{-1}\));
\(\boldsymbol{\Gamma}(\mathbf{r},\mathbf{r}')\) is the outgoing dyadic Green tensor
(m\(^{-1}\)) that propagates unit current dipoles from \(\mathbf{r}'\) to
\(\mathbf{r}\);
\(\mu_0=4\pi\times 10^{-7}\,\mathrm{H\,m^{-1}}\) is the vacuum permeability;
\(\mathbf{r}'\) ranges over the source volume \(\Omega_s\); and the volume element
\(d^{3}r'\) carries SI units m\(^{3}\).
The factor \(-\jj\omega\mu_0\) arises from the harmonic Amp\`ere law under the
locked phasor convention.
Here \(\boldsymbol{\Gamma}\) is the outgoing dyadic Green function
\cite{stratton1941}. Equation~\eqref{eq:ch2.radiation-integral} is a
frequency-domain representation of the radiative field generated by a compact source
distribution through the outgoing dyadic Green function \cite{stratton1941}.
Every application of this representation assumes the $\mathcal{D}_J=1$ and boundary conditions
stated above. Unless noted otherwise, \(\omega>0\) is fixed,
materials are linear isotropic passive regions, and \(\Omega\) is either bounded
with declared \(\mathcal{B}_\Gamma\) or an exterior domain with compact
\(\operatorname{supp}\Jimp\) and outgoing radiation condition at infinity. Maps are defined
only on \(\mathcal{J}_\Omega\) satisfying Eq.~\eqref{eq:ch1.charge-continuity-harmonic};
generalized sources appear in Section~\ref{sec:ch232}. Maxwell linearity on \(\Omega\) induces linearity of both maps on fixed boundary
data: for \(\Jimp^{(1)},\Jimp^{(2)}\in\mathcal{J}_\Omega\) and scalars \(a_{1},a_{2}\),
\begin{equation}
\label{eq:ch2.radiation-linearity}
\mathcal{R}_\omega\!\left[a_{1}\Jimp^{(1)}+a_{2}\Jimp^{(2)}\right]
=
a_{1}\mathcal{R}_\omega[\Jimp^{(1)}]
+a_{2}\mathcal{R}_\omega[\Jimp^{(2)}],
\end{equation}
In Eq.~\eqref{eq:ch2.radiation-linearity},
\(a_1,a_2\in\mathbb{C}\) are complex amplitudes and
\(\Jimp^{(1)},\Jimp^{(2)}\in\mathcal{J}_\Omega\) are admissible impressed currents;
\(\mathcal{R}_\omega\) is the radiation map; and linearity holds when boundary data
\(\mathbf{g}_\Gamma\) and the material layout are held fixed \cite{maxwell1873}.
Composite-source analysis decomposes \(\Jimp\) into constituents,
maps each through \(\mathcal{R}_\omega\), and superposes radiative outputs---a
discipline that prevents reactive interference from masquerading as transport
(Section~\ref{sec:ch01.6.4}).
Write \(\operatorname{ran}\mathcal{S}_\omega\subset\mathcal{F}_\Omega\) for all
field pairs reachable from finite sources and
\(\operatorname{ran}\mathcal{R}_\omega\subset\mathcal{F}_{\mathrm{rad}}\) for
far-zone radiative outputs. Strict inclusion
\(\operatorname{ran}\mathcal{R}_\omega\subsetneq\operatorname{ran}\mathcal{S}_\omega\)
occurs whenever reactive storage is nontrivial: distinct near-zone states can share
identical radiative images (Section~\ref{sec:ch01.2.4}). Pattern measurements
probe \(\operatorname{ran}\mathcal{R}_\omega\) and cannot invert
\(\mathcal{S}_\omega\) without additional near-zone data---a limitation revisited in
Section~\ref{sec:ch244}.
Radiative power is metered with inherited Poynting machinery.
The practical question is how much time-averaged power crosses a closed surface that
encloses the source and lies beyond the reactive shell; that flux is the scalar
quantity laboratories report as radiated power.
On a sphere
\(S_R=\{|\mathbf{r}|=R\}\) enclosing \(\operatorname{supp}\Jimp\), with outward
unit vector \(\hat{\mathbf{r}}\) and complex Poynting vector
\(\mathbf{S}=\tfrac{1}{2}\E\times\H^*\) from Section~\ref{sec:ch01.1.4},
\begin{equation}
\label{eq:ch2.radiated-power-screen}
P_{\mathrm{rad}}(R)
=
\Re\oint_{S_R}\mathbf{S}\cdot\hat{\mathbf{r}}\,dS.
\end{equation}
In Eq.~\eqref{eq:ch2.radiated-power-screen},
\(P_{\mathrm{rad}}(R)\) is the time-averaged outward power (W) through the sphere
of radius \(R\) (m);
\(\mathbf{S}\) is the complex Poynting vector (W\,m\(^{-2}\)) built from the
radiative field pair;
\(\hat{\mathbf{r}}=\mathbf{r}/|\mathbf{r}|\) is the outward unit normal on
\(S_R\); and \(dS\) is the area element (m\(^{2}\)).
The real part \(\Re\{\cdot\}\) extracts the time-averaged power flux carried by the
phasor fields under the \(\exp(-\jj\omega t)\) convention.
Equation~\eqref{eq:ch2.radiated-power-screen} inherits the Poynting law of
Section~\ref{sec:ch01.1.4}. For \(\mathcal{R}_\omega[\Jimp]\), \(P_{\mathrm{rad}}(R)\) becomes
$R$-independent beyond the reactive shell \cite{poynting1884}. Angular-momentum budgets from
Section~\ref{sec:ch01.1.5} attach to the same radiative component.
\begin{figure}[t]
\centering
\ChOneFigBlock{%
\begin{tikzpicture}[ch1fig, node distance=5mm and 7mm]
 \node[ch1box, minimum width=22mm] (core) at (0,0) {Source\\$\Omega_s$};
 \node[ch1box, minimum width=28mm] (react) at (0,-12mm) {Reactive shell\\$1/r^2$, $1/r^3$};
 \node[ch1box, minimum width=32mm] (sel) at (0,-24mm) {Selector\\$\Pi_{\mathrm{rad}}$};
 \node[ch1box, minimum width=28mm] (far) at (0,-36mm) {Far zone\\$1/r$ outgoing};
 \draw[ch1flow] (core) -- (react) -- (sel) -- (far);
\end{tikzpicture}%
}
\caption{Spatial nesting seen by the maps: \(\mathcal{S}_\omega\) populates all
regions; \(\Pi_{\mathrm{rad}}\) passes only the outgoing transverse content
compatible with Section~\ref{sec:ch01.2.4}.}
\label{fig:ch2.reactive-radiative-split}
\end{figure}
The solution map energizes reactive and radiative regions alike; the
radiation map discards non-outgoing storage before pattern or transport claims are
made. Misidentifying reactive circulation as exported angular structure is a
recurring error in high-gradient near-zone beams \cite{allen1992}. Finite measurement channels (Section~\ref{sec:ch01.5}) implement linear functionals
on field space. Composing a declared far-zone observable with \(\mathcal{R}_\omega\)
yields a linear functional of \(\Jimp\); composing the same observable with
\(\mathcal{S}_\omega\) without selection generally includes reactive contamination.
No new measurement axioms are introduced---only the correct map for the declared
observable class \cite{stratton1941}. Non-radiating currents illustrate kernel visibility. If
\(\Jimp^{\mathrm{(NR)}}\in\mathcal{J}_\Omega\) has vanishing far-zone transverse
fields while maintaining near-zone storage (Section~\ref{sec:ch01.2.4}), then
\begin{equation}
\label{eq:ch2.null-radiation}
\mathcal{R}_\omega[\Jimp^{\mathrm{(NR)}}]=\mathbf{0}
\quad\text{in}\ \mathcal{F}_{\mathrm{rad}},
\end{equation}
In Eq.~\eqref{eq:ch2.null-radiation},
\(\Jimp^{\mathrm{(NR)}}\) is an admissible impressed current (A\,m\(^{-2}\)) that satisfies
\(\mathcal{D}_{J}=1\) yet carries no outgoing transverse export;
\(\mathcal{F}_{\mathrm{rad}}\) is the far-zone radiative field space selected by
\(\Pi_{\mathrm{rad}}\); and \(\ker\mathcal{R}_\omega\) denotes the set of currents for
which the radiation map vanishes on that space,
while \(\mathcal{S}_\omega[\Jimp^{\mathrm{(NR)}}]\neq\mathbf{0}\) in general.
Equation~\eqref{eq:ch2.null-radiation} targets the null space analyzed in
Section~\ref{sec:ch272}; ignoring \(\ker\mathcal{R}_\omega\) invites spurious modal
assignments to invisible sources.
Membership in \(\mathcal{J}_\Omega\) requires more than square integrability:
the Maxwell-closed screen below states, in compact form, which impressed currents
may enter the radiation chain without contradicting charge continuity or boundary
compatibility. \(\mathcal{D}_{J}\) is not a mathematical
regularity condition alone: it encodes the laboratory requirement that impressed currents
be charge-consistent and boundary-compatible before Maxwell field specification is attempted
\cite{maxwell1873}. Violating the screen produces formal solutions that no
passive feed can realize \cite{harrington2001}:
\begin{equation}
\label{eq:ch2.source-charge-continuous}
\mathcal{D}_{J}[\Jimp]=1
\iff
\left(
\Jimp\in L^{2}(\Omega)^{3},\
\divg\Jimp=-\jj\omega\rho,\
\Jimp\ \text{satisfies BC compatibility on }\Gamma
\right),
\end{equation}
In Eq.~\eqref{eq:ch2.source-admissibility},
\(\mathcal{D}_{J}\) is the source Maxwell-closed screen (dimensionless) inherited from
Eq.~\eqref{eq:ch1.continuity-charge-continuous};
\(\Jimp\) is the impressed electric current density (A\,m\(^{-2}\)) on compact
\(\Omega\subset\mathbb{R}^{3}\);
\(\rho\) is the impressed charge density (C\,m\(^{-3}\)) linked by
\(\divg\Jimp=-\jj\omega\rho\); and \(\Gamma=\partial\Omega\) is the boundary on which
tangential traces must match the declared boundary class.
with \(\mathcal{D}_{J}\) as in Eq.~\eqref{eq:ch1.continuity-admissibility}.
Violating \(\divg\Jimp=-\jj\omega\rho\) contradicts
Eq.~\eqref{eq:ch1.charge-continuity-harmonic} before Maxwell field specification is attempted.
Boundary compatibility ensures induced surface equivalents do not contradict
Eq.~\eqref{eq:ch1.boundary-data-operator} \cite{stratton1941}.
Causality constrains harmonic representations.
A frequency-domain phasor is physically meaningful only when it can be traced to a
retarded time-domain evolution that respects finite propagation speed.
With time-domain
current \(j(\mathbf{r},t)\) and retarded dyadic impulse
\(\mathcal{G}(\mathbf{r},\mathbf{r}',t-t')\),
\begin{equation}
\label{eq:ch2.causality-retarded}
\E(\mathbf{r},t)
=
-\mu_0\,\frac{\partial}{\partial t}
\int
\mathcal{G}(\mathbf{r},\mathbf{r}',t-t')\cdot j(\mathbf{r}',t')\,d^{3}r'\,dt',
\end{equation}
integrated only over \(t'\le t\) \cite{stratton1941}.
In Eq.~\eqref{eq:ch2.causality-retarded},
\(\E(\mathbf{r},t)\) is the real time-domain electric field (V\,m\(^{-1}\));
\(j(\mathbf{r}',t')\) is the impressed current density (A\,m\(^{-2}\)) at source
point \(\mathbf{r}'\) and retarded time \(t'\);
\(\mathcal{G}(\mathbf{r},\mathbf{r}',t-t')\) is the causal dyadic Green impulse
response (m\(^{-1}\)\,s\(^{-1}\)) that vanishes for \(t-t'<0\); and the temporal
integral enforces that the field at \((\mathbf{r},t)\) depends only on sources in
the past light cone.
The steady-state operator
$\mathcal{R}_\omega$ may therefore be viewed as the frequency-domain specialization of
the retarded causal evolution \cite{stratton1941}. The equivalence between
the time-domain and harmonic descriptions follows under the standard assumptions of
linearity, stationarity, and finite-energy excitation \cite{maxwell1873}.
Vacuum scales from Eq.~\eqref{eq:ch1.vacuum-impedance-wavenumber} fix the carrier
on which selection acts: \(k_0=\omega\sqrt{\mu_0\varepsilon_0}\),
\(\eta_0=\sqrt{\mu_0/\varepsilon_0}\). Outgoing components retained by
\(\Pi_{\mathrm{rad}}\) carry \(\exp(-\jj k_0 r)\) phase fronts at leading order;
components whose phase does not lock to outgoing \(k_0 r\) remain in
\(\mathcal{S}_\omega[\Jimp]\) but outside \(\mathcal{F}_{\mathrm{rad}}\). Power
metered by Eq.~\eqref{eq:ch2.radiated-power-screen} stabilizes only on the latter
\cite{stratton1941}. Distance labels become meaningful through operator action. For compact
\(\operatorname{supp}\Jimp\), \(\mathcal{S}_\omega[\Jimp]\) contains
\(1/r^{2}\) and \(1/r^{3}\) reactive terms; \(\mathcal{R}_\omega[\Jimp]\) retains
leading \(1/r\) transverse fields with \(kr\) phase in the far zone
(Section~\ref{sec:ch01.2.4}). Asymptotic partitions are sharpened in
Section~\ref{sec:ch252}.
Observability enters as functional composition rather than as a new postulate.
Let \(\mathcal{O}_{\mathrm{far}}\) denote a linear functional that extracts a
declared far-zone scalar or vector observable---pattern voltage, axial ratio
proxy, or a single spherical-harmonic coefficient---from a field pair in
\(\mathcal{F}_{\mathrm{rad}}\). The measured quantity associated with impressed current \(\Jimp\) is the scalar
\(\mathcal{O}_{\mathrm{far}}\big[\mathcal{R}_\omega[\Jimp]\big]\), not
\(\mathcal{O}_{\mathrm{far}}\big[\mathcal{S}_\omega[\Jimp]\big]\) unless the
observable has been shown to vanish on reactive content. Because both
\(\mathcal{R}_\omega\) and \(\mathcal{O}_{\mathrm{far}}\) are linear on their
domains, the laboratory reading is a linear functional
\(\Lambda_{\mathrm{rad}}=\mathcal{O}_{\mathrm{far}}\circ\mathcal{R}_\omega\)
on \(\mathcal{J}_\Omega\). That composition expresses in map language
the measurement limitations in Section~\ref{sec:ch01.5}: finite apertures and
finite modal sets appear as rank-deficient \(\mathcal{O}_{\mathrm{far}}\), not
as modifications of Maxwell law \cite{stratton1941}. Angular-transport claims
require further that \(\Lambda_{\mathrm{rad}}\) be compatible with the flux
budgets of Section~\ref{sec:ch01.1.5} when the observable is interpreted as
momentum or power export.
The outgoing selector \(\Pi_{\mathrm{rad}}\) is idempotent on
\(\mathcal{F}_{\mathrm{rad}}\): applying selection twice leaves a radiative state
unchanged. On the larger space \(\mathcal{F}_\Omega\), \(\Pi_{\mathrm{rad}}\) is
a projection only after the Sommerfeld-class uniqueness of Section~\ref{sec:ch243}
is imposed; without that completion, multiple outgoing decompositions can coexist
for the same truncated near-zone data. The idempotence condition is therefore a
\emph{joint} property of field space and boundary class, not a algebraic trick
independent of physics \cite{harrington2001}. Numerical methods that
window fields in space and call the result ``radiative'' without outgoing radiation condition are approximating \(\Pi_{\mathrm{rad}}\) by a non-idempotent mask; such outputs
must not be fed into transport criteria without an explicit error budget.
The composition law used downstream is
\begin{equation}
\label{eq:ch2.measurement chain-composition}
\mathcal{R}_\omega
=
\Pi_{\mathrm{rad}}\circ\mathcal{S}_\omega,
\qquad
\mathcal{D}_{J}=1
\Longrightarrow
\mathcal{S}_\omega:\mathcal{J}_\Omega\to\mathcal{F}_\Omega,
\quad
\mathcal{R}_\omega:\mathcal{J}_\Omega\to\mathcal{F}_{\mathrm{rad}}.
\end{equation}
In Eq.~\eqref{eq:ch2.measurement chain-composition},
\(\mathcal{S}_\omega\) is the Maxwell solution map from admissible sources to field pairs,
\(\Pi_{\mathrm{rad}}\) is the outgoing selector that retains far-zone transverse export,
\(\mathcal{R}_\omega=\Pi_{\mathrm{rad}}\circ\mathcal{S}_\omega\) is the radiation map,
\(\mathcal{D}_{J}=1\) denotes source Maxwell-closed from
Eq.~\eqref{eq:ch2.source-charge-continuous}, and \(\mathcal{J}_\Omega\),
\(\mathcal{F}_\Omega\), \(\mathcal{F}_{\mathrm{rad}}\) are the source and field spaces
declared in Subsection~\ref{sec:ch201}.
Green kernels realize \(\mathcal{S}_\omega\); Sommerfeld conditions define
\(\Pi_{\mathrm{rad}}\); spectral theory in Section~\ref{sec:ch26} studies
eigenstructure on spaces refined in Section~\ref{sec:ch21}
\cite{harrington2001}. Nothing in
Eqs.~\eqref{eq:ch2.maxwell-solution-map}--\eqref{eq:ch2.measurement chain-composition}
re-proves results listed in Section~\ref{sec:ch01.7.1}; continuity
Eq.~\eqref{eq:ch1.continuity-map} is honored by citation-only continuity
of curl--curl algebra, stress flux theorems, and gauge lemmas.
Subsection~\ref{sec:ch201} has named \(\mathcal{S}_\omega\), \(\Pi_{\mathrm{rad}}\), and
\(\mathcal{R}_\omega\), introduced their Green and Poynting realizations, and connected
finite measurement rank to composition with far-zone observables. Subsection~\ref{sec:ch202}
identifies the Maxwell operator whose inversion, under declared boundary completion,
produces \(\mathcal{S}_\omega\).
\subsection{Maxwell equations in operator form} \label{sec:ch202}
% HARRINGTON-OPEN:sec:ch202
A standard-gain horn on a compact-range turntable fixes the laboratory scene for `Maxwell equations in operator form'': frequency, probe standoff, and wattmeter closure on $\Gamma_R$ precede the formal displays below \cite{hansen1988}. Maxwell's retarded export to $\Gamma_R$ remains the governing constraint \cite{maxwell1873}.
Subsection~\ref{sec:ch201} named the solution map $\mathcal{S}_\omega$ and the outgoing
selector $\Pi_{\mathrm{rad}}$; the display below identifies the Maxwell operator whose
inversion, under declared boundary completion, produces that map \cite{stratton1941}. The governing question is which field pair $(\E,\H)$ is
consistent with impressed currents $\Jimp$ and with the termination class $\mathcal{B}_\Gamma$
on $\Gamma=\partial\Omega$---the same question every full-wave solver answers before any
pattern is extracted \cite{harrington2001}. Harmonic Maxwell theory on $\Omega$ is packaged
as a first-order residual system on a six-component state,
augmented by boundary traces through the operator
Eq.~\eqref{eq:ch1.boundary-data-operator} imported from
Section~\ref{sec:ch01.1.2}. Local curl and divergence balances remain those of
Eq.~\eqref{eq:ch1.maxwell-operator-block}; the Maxwell export sequence developed here does not
alter those local balances \cite{stratton1941}.
First-order Maxwell theory couples electric and magnetic fields through circulation;
packaging both components in a single state vector is the standard device that permits
resolvent notation and variational formulations without abandoning the physical coupling
\cite{stratton1941}. \(\mathbf{u}_\omega\) is
not an arbitrary six-vector: each admissible state satisfies Gauss constraints and carries
tangential traces compatible with the boundary class declared with the impressed data tuple
\(\mathfrak{D}\). From an antenna engineer's perspective, \(\mathbf{u}_\omega\) is what a
full-wave solver stores at each mesh node---the coupled \((\E,\H)\) phasor pair at fixed
\(\omega\)---before any post-processing extracts patterns or port quantities
\cite{harrington2001}. Collect phasor fields into the Maxwell state
% PHYSICAL-MEMORY:sec:ch202
\begin{physicalvignette}[Spherical near-field scanner with dual-polarization probes]
Consider Spherical near-field scanner with dual-polarization probes. Before any coefficient table is printed, technicians record frequency, cable loss, and the mechanical outline that fixes electrical size. The first walk-through of the range is always tactile: flange bolts, probe standoff, turntable boresight, and the wattmeter reference on the receive chain \cite{hansen1988}.
On this hardware the near zone is visually impressive---standing reactive structure, rapid phase rotation, sometimes a visible VSWR ripple at the feed---yet the export question is measured only after outgoing selection on a large sphere \cite{harrington2001}. What the camera sees in the chamber is not automatically what the far-zone pattern claims \cite{stratton1941}.
Engineers therefore learn to separate three readings: what the feed can synthesize on declared support, what propagating content survives on the radiative range, and what a finite instrument can recover at declared tolerance \cite{devaney2004}. Confusing those readings produces the rank gaps and watt failures that later sections quantify \cite{pal2026}.
Rayleigh and Stratton treated such ranges as boundary-value laboratories long before digital near-field transforms became routine \cite{stratton1941}. Harrington's method-of-moments culture reinforced the same discipline: metal outline first, harmonic headline second \cite{harrington2001}.
Subsection~\ref{sec:ch202} states the formal export condition that this range scene presupposes \cite{stratton1941}.
\end{physicalvignette}
\begin{equation}
\label{eq:ch2.maxwell-state-vector}
\mathbf{u}_\omega=
\begin{bmatrix}\E\\ \H\end{bmatrix}
\in\mathcal{F}_\Omega,
\end{equation}
where \(\mathbf{u}_\omega\) denotes the six-component Maxwell state at fixed
\(\omega\), with \(\E\) and \(\H\) carrying the SI units noted above and lying in
the admissible field space \(\mathcal{F}_\Omega\) declared in Subsection~\ref{sec:ch201}.
The pair \((\E,\H)\) is constrained jointly: divergence
rows of Eq.~\eqref{eq:ch1.maxwell-operator-block} must hold, and tangential traces on
\(\Gamma=\partial\Omega\) must lie in the domain of the boundary class declared with
\(\mathfrak{D}\). Treating \(\mathbf{u}_\omega\) as a single unknown is what permits
resolvent notation, variational forms, and the spectral compression developed in
Section~\ref{sec:ch26} \cite{harrington2001}.
Every boundary-value problem begins by distinguishing prescribed quantities from unknown
quantities. In Maxwell theory the impressed currents, charges, and boundary data constitute
the prescribed excitation, whereas the electromagnetic field remains the unknown state to
be determined \cite{maxwell1873}. The operator formulation merely records this
distinction in a compact mathematical form suitable for analysis \cite{harrington2001}.
Impressed data enter through a forcing vector aligned with the Maxwell residual rows.
The forcing packages charge and current drive on the rows where impressed sources
appear, leaving divergence constraints to be satisfied by the field unknown.
For
the tuple Eq.~\eqref{eq:ch1.data-specification-tuple},
\begin{equation}
\label{eq:ch2.maxwell-forcing-vector}
\mathbf{s}_\omega(\mathfrak{D})=
\begin{bmatrix}\mathbf{0}\\ -\Jimp\\ -\rho\\ \mathbf{0}\end{bmatrix},
\end{equation}
In Eq.~\eqref{eq:ch2.maxwell-forcing-vector},
\(\mathbf{s}_\omega(\mathfrak{D})\) is the Maxwell forcing six-vector (SI units per row)
assembled from the data tuple \(\mathfrak{D}\);
\(\Jimp\) enters the Amp\`ere row (A\,m\(^{-2}\)) and \(\rho\) the Gauss row
(C\,m\(^{-3}\)); the remaining rows are zero for the impressed-current class treated here;
row order matches Eq.~\eqref{eq:ch1.maxwell-operator-block}. Equation
\eqref{eq:ch2.maxwell-forcing-vector} records drive on the Amp\`ere and Gauss rows
only for the impressed-current class treated here; sheet and magnetic excitations
extend the vector through Eq.~\eqref{eq:ch1.interface-conditions} when trace spaces
are enriched in Section~\ref{sec:ch211} \cite{stratton1941}. The forcing is \emph{data}, not a
field variable: it specifies how \(\mathcal{S}_\omega\) is sourced while
\(\mathbf{u}_\omega\) remains the unknown. Operator language separates \emph{states}
\((\E,\H)\) from \emph{data} \((\Jimp,\rho,\mathfrak{D})\): the feed specifies what is
driven, the solver returns what Maxwell allows, and radiation operators report what is
exported \cite{maxwell1873}. Treating a tangential-field template as independent forcing
without a matching current repeats the realizability failure already flagged in
Chapter~1 \cite{stratton1941}.
Let \(\mathcal{M}_\omega\) denote the harmonic Maxwell operator of
Eq.~\eqref{eq:ch1.maxwell-operator-block}. The volume subsystem reads
\begin{equation}
\label{eq:ch2.maxwell-driven-operator}
\mathcal{M}_\omega[\mathbf{u}_\omega]=\mathbf{s}_\omega(\mathfrak{D}),
\end{equation}
and is satisfied when each residual row vanishes. Linearity in \(\mathbf{u}_\omega\)
follows from linearity of Eq.~\eqref{eq:ch1.maxwell-local}; spatially varying
\(\varepsilon(\mathbf{r})\) and \(\mu(\mathbf{r})\) modify coefficients inside
\(\mathcal{M}_\omega\) without changing the first-order block shape
(Section~\ref{sec:ch222}) \cite{stratton1941}. In uniform vacuum regions,
\(\mathcal{M}_\omega\) has translation-invariant symbol, the property exploited by
dyadic Green kernels in Section~\ref{sec:ch23}.
\begin{figure}[t]
\centering
\ChOneFigBlock{%
\begin{tikzpicture}[ch1fig, node distance=5mm and 7mm]
 \node[ch1box] (vol) {Volume\\$\mathcal{M}_\omega\mathbf{u}=\mathbf{s}$};
 \node[ch1box, below=5mm of vol] (bnd) {Boundary\\$\mathcal{B}_\Gamma\gamma_\Gamma\mathbf{u}=\mathbf{g}$};
 \node[ch1box, left=18mm of vol] (u) {State $\mathbf{u}_\omega$};
 \draw[ch1flow] (u.west) -- (u.west |- vol.east);
 \draw[ch1flow] (u.west |- vol.east) -- (vol.east);
 \draw[ch1flow] (u.east) -- (u.east |- bnd.west);
 \draw[ch1flow] (u.east |- bnd.west) -- (bnd.west);
\end{tikzpicture}%
}
\caption{Parallel volume and boundary channels for the closed harmonic problem:
\(\mathbf{u}_\omega\) simultaneously satisfies Maxwell residuals and prescribed
traces.}
\label{fig:ch2.maxwell-boundary-stack}
\end{figure}
~\eqref{eq:ch2.maxwell-driven-operator} and
\eqref{eq:ch1.boundary-data-operator}. A solver that enforces only volume rows
produces a formal solution compatible with Maxwell \emph{sources} but not necessarily
with termination physics; conversely, boundary data without volume balance violate
conservation. Subsection~\ref{sec:ch201} applies \(\mathcal{R}_\omega\) only after
both channels are closed \cite{harrington2001}. Tangential traces are extracted by the restriction
\begin{equation}
\label{eq:ch2.trace-restriction}
\gamma_\Gamma:\mathcal{F}_\Omega\to\mathcal{T}_\Gamma,
\qquad
\gamma_\Gamma[\mathbf{u}_\omega]=
\begin{bmatrix}
\left.\hat{\mathbf{n}}\times\E\right|_{\Gamma}\\
\left.\hat{\mathbf{n}}\times\H\right|_{\Gamma}
\end{bmatrix},
\end{equation}
with outward unit normal \(\hat{\mathbf{n}}\) on \(\Gamma\) and trace space
\(\mathcal{T}_\Gamma\) induced by interior states. The map
\(\gamma_\Gamma\) is not defined on arbitrary \(L^{2}\) six-tuples; Sobolev and
distribution refinements in Section~\ref{sec:ch211} specify its domain
\cite{stratton1941}. The closed operator stacks volume and boundary residuals:
\begin{equation}
\label{eq:ch2.maxwell-boundary-composite}
\mathcal{L}_\omega[\mathbf{u}_\omega]=
\begin{bmatrix}
\mathcal{M}_\omega[\mathbf{u}_\omega]\\
\mathcal{B}_\Gamma[\gamma_\Gamma[\mathbf{u}_\omega]]
\end{bmatrix}.
\end{equation}
A physically closed harmonic problem requires both volume Maxwell balance and boundary
termination to hold simultaneously; neither channel alone suffices
\cite{stratton1941}. The composite operator \(\mathcal{L}_\omega\) encodes that joint
requirement: its first block tests whether \((\E,\H)\) satisfies the curl--curl structure
with impressed sources removed to the forcing vector, while its second block tests whether
tangential traces on \(\Gamma\) match the declared termination (PEC, impedance, or outgoing)
\cite{harrington2001}. Experimentally, a solver that minimizes volume residuals while
ignoring port or chamber walls produces a formal Maxwell solution that may disagree with
what a network analyzer reports on the same structure \cite{stratton1941}. The harmonic
boundary-value problem is to find \(\mathbf{u}_\omega\in\mathcal{F}_\Omega\)
such that
\begin{equation}
\label{eq:ch2.closed-maxwell-bvp}
\mathcal{L}_\omega[\mathbf{u}_\omega]=
\begin{bmatrix}\mathbf{s}_\omega(\mathfrak{D})\\ \mathbf{g}_\Gamma\end{bmatrix}.
\end{equation}
Perfect conductors impose \(\hat{\mathbf{n}}\times\E=\mathbf{0}\); radiation and
Leontovich completions encode outgoing or dissipative relations on
\(\gamma_\Gamma[\mathbf{u}_\omega]\). The operator \(\mathcal{L}_\omega\) is the
Maxwell balance operator for this coupled system \cite{harrington2001}.
Hence $\mathcal{L}_\omega$ enforces Maxwell balance with impressed drive and boundary termination; without outgoing selection it does not distinguish export from storage, and every near-field probe reading $\mathcal{S}_\omega[\Jimp]$ presupposes its inversion on the declared domain \cite{stratton1941}.
Equation~\eqref{eq:ch2.maxwell-solution-map} is the resolvent relation
\begin{equation}
\label{eq:ch2.solution-map-resolvent}
\mathcal{S}_\omega[\Jimp]=\mathbf{u}_\omega
\quad\text{with}\quad
\mathcal{L}_\omega[\mathbf{u}_\omega]=
\begin{bmatrix}\mathbf{s}_\omega(\Jimp,\rho)\\ \mathbf{g}_\Gamma\end{bmatrix},
\end{equation}
where \(\rho\) is tied to \(\Jimp\) through
Eq.~\eqref{eq:ch1.charge-continuity-harmonic} when only electric impressed currents
drive the system. Invertibility of \(\mathcal{L}_\omega\) is not universal: existence
and uniqueness depend on \(\mathcal{B}_\Gamma\) and on whether outgoing selection is
imposed (Section~\ref{sec:ch243} for exterior domains; Section~\ref{sec:ch261} for
cavities) \cite{stratton1941}. \(\mathcal{S}_\omega\) is therefore a \emph{physical
branch} of \(\mathcal{L}_\omega^{-1}\), not an abstract algebraic inverse on all of
\(\mathcal{F}_\Omega\). \begin{figure}[t]
\centering
\ChOneFigBlock{%
\begin{tikzpicture}[ch1fig, node distance=10mm]
 \node[ch1box] (E) at (0,6mm) {$\curl\E$};
 \node[ch1box] (H) at (0,-6mm) {$\curl\H$};
 \node[ch1box, right=20mm of E] (r1) {$-\jj\omega\mu_0\H$};
 \node[ch1box, right=20mm of H] (r2) {$\Jimp+\jj\omega\varepsilon_0\E$};
 \draw[ch1flow] (E.east) -- (E.east |- r1.west);
 \draw[ch1flow] (E.east |- r1.west) -- (r1.west);
 \draw[ch1flow] (H.east) -- (H.east |- r2.west);
 \draw[ch1flow] (H.east |- r2.west) -- (r2.west);
 \draw[ch1flow, dashed] (r1.south) to[bend left=25] (H.east);
 \draw[ch1flow, dashed] (r2.north) to[bend left=25] (E.east);
\end{tikzpicture}%
}
\caption{Circulation coupling in \(\mathcal{M}_\omega\): each curl row feeds the
partner field; \(\Jimp\) appears as external drive on the magnetic row.}
\label{fig:ch2.first-order-curl-coupling}
\end{figure}
Single-field Helmholtz reductions equivalent to
Eqs.~\eqref{eq:ch1.vacuum-wave-source-driven}--\eqref{eq:ch1.vacuum-curl-second-order}
are used later when a scalar symbol simplifies kernels; the elimination belongs to
\ref{ch:01} and is not repeated. The electric-wave symbol used in
\ref{ch:02} is recorded as
\begin{equation}
\label{eq:ch2.electric-wave-operator}
\mathcal{W}_{E,\omega}\E\equiv
\curl\!\left(\mu_0^{-1}\curl\E\right)-\omega^2\varepsilon_0\E
= -\jj\omega\mu_0\Jimp+\mu_0^{-1}\curl\mathbf{J}^{\mathrm{(mag)}}_{\mathrm{eff}},
\end{equation}
with \(\mathbf{J}^{\mathrm{(mag)}}_{\mathrm{eff}}\) collecting equivalent magnetic
drive when present \cite{stratton1941}. Subsection~\ref{sec:ch222} supplies the
magnetic counterpart and material-weighted generalizations. Variational structure is inherited from Section~\ref{sec:ch01.1.4}. For test state
\(\mathbf{v}_\omega=(\E',\H')\),
\begin{equation}
\label{eq:ch2.maxwell-energy-form}
\mathfrak{a}_\omega(\mathbf{u}_\omega,\mathbf{v}_\omega)=
\int_\Omega\Big[\mu_0\,\H\!\cdot\!\curl\E'-\varepsilon_0\,\E\!\cdot\!\curl\H'\Big]d^{3}r.
\end{equation}
Integration by parts connects Eq.~\eqref{eq:ch2.maxwell-energy-form} to surface
fluxes and to the weak formulation
\begin{equation}
\label{eq:ch2.maxwell-weak-form}
\mathfrak{a}_\omega(\mathbf{u}_\omega,\mathbf{v}_\omega)=\mathfrak{f}_\omega(\mathbf{v}_\omega)
\quad\forall\,\mathbf{v}_\omega\in\mathcal{F}_\Omega^{0},
\end{equation}
where \(\mathcal{F}_\Omega^{0}\) carries homogeneous boundary traces compatible with
\(\mathcal{B}_\Gamma\) and \(\mathfrak{f}_\omega\) is induced by
\(\mathbf{s}_\omega(\mathfrak{D})\) \cite{poynting1884}. Equation
\eqref{eq:ch2.maxwell-weak-form} is the variational face of
Eq.~\eqref{eq:ch2.closed-maxwell-bvp}; finite-element discretizations favor this
route when material coefficients are piecewise constant
\cite{harrington2001}.
Well-posedness of Eq.~\eqref{eq:ch2.maxwell-weak-form} is not automatic on
unbounded domains or on multiply connected cavities. On a bounded region with
passive \(\mathcal{B}_\Gamma\) and homogeneous traces in
\(\mathcal{F}_\Omega^{0}\), coercivity fails only along divergence-free
cavity modes that satisfy both curl-curl balance and homogeneous tangential
data---precisely the null space that Section~\ref{sec:ch214} isolates before
imposing gauge or polarization fixes \cite{harrington2001}. For
exterior problems, coercivity is recovered only after restricting to the
outgoing subspace selected by Sommerfeld-type completion; without that restriction,
the bilinear form is indefinite and the weak problem admits non-physical
growing solutions. The practical rule for antenna analysis is therefore
threefold: declare \(\mathcal{D}(\mathcal{L}_\omega)\), specify
\(\mathcal{B}_\Gamma\) on every connected piece of \(\Gamma\), and confirm that
the impressed tuple lies in the compatible range of
Eq.~\eqref{eq:ch2.source-compatibility}. Numerical codes that enforce volume
Maxwell rows on truncated domains but leave exterior faces open implement an
\emph{implicit} \(\mathcal{B}_\Gamma\); unless that implicit completion matches
the laboratory termination, \(\mathcal{S}_\omega\) in the code is a different
operator from the one measured on the bench \cite{harrington2001}. Domain specification cannot remain at \(C^\infty\) when sources are distributional.
The minimal energy class declared here is
\begin{equation}
\label{eq:ch2.operator-domain-declaration}
\mathcal{D}(\mathcal{L}_\omega)\subset
\big\{\mathbf{u}_\omega=(\E,\H):
\E,\H\in L^{2}(\Omega)^{3},\ \curl\E,\curl\H\in L^{2}(\Omega)^{3}\big\},
\end{equation}
where $\mathcal{L}_\omega$ is the composite Maxwell operator on the six-component state
$\mathbf{u}_\omega=(\E,\H)$, $\mathcal{D}(\mathcal{L}_\omega)$ its domain of square-integrable
curl fields compatible with $\mathcal{B}_\Gamma$, and $\Omega$ the analysis region
\cite{stratton1941}. Section~\ref{sec:ch211} refines embeddings and compactness; evaluating
Eq.~\eqref{eq:ch2.maxwell-driven-operator} outside
\(\mathcal{D}(\mathcal{L}_\omega)\) generates spurious boundary terms that corrupt
flux identities \cite{stratton1941}.
Piecewise smooth \(\varepsilon(\mathbf{r})\) and \(\mu(\mathbf{r})\) modify
coefficients inside \(\mathcal{M}_\omega\) while preserving the block template of
Eq.~\eqref{eq:ch1.maxwell-operator-block}. Interface jumps satisfy
Eq.~\eqref{eq:ch1.interface-conditions}; multi-region Green operators in
Section~\ref{sec:ch234} realize \(\mathcal{L}_\omega^{-1}\) on composite domains
\cite{stratton1941}. The coupling structure described above extends naturally to practical
terminal conditions. Finite feed structures and external circuits frequently impose
impedance relations between tangential field components on the boundary
\cite{harrington2001}:
\begin{equation}
\label{eq:ch2.boundary-impedance-operator}
\mathcal{Z}_\Gamma\!
\begin{bmatrix}
\left.\hat{\mathbf{n}}\times\E\right|_{\Gamma}\\
\left.\hat{\mathbf{n}}\times\H\right|_{\Gamma}
\end{bmatrix}=\mathbf{0},
\end{equation}
In Eq.~\eqref{eq:ch2.boundary-impedance-operator},
\(\mathcal{Z}_\Gamma\) is a linear boundary impedance operator (SI units absorbed in its
definition) acting on the tangential Maxwell traces
\(\hat{\mathbf{n}}\times\E|_{\Gamma}\) and \(\hat{\mathbf{n}}\times\H|_{\Gamma}\) on
\(\Gamma=\partial\Omega\); admissible states satisfy the terminal or absorbing relation
\(\mathcal{Z}_\Gamma[\text{traces}]=\mathbf{0}\).
encoding Leontovich, absorbing, or mode-matched completions \cite{harrington2001}.
When Eq.~\eqref{eq:ch2.boundary-impedance-operator} replaces explicit
\(\mathbf{g}_\Gamma\) in Eq.~\eqref{eq:ch2.closed-maxwell-bvp}, admissible states
satisfy the impedance manifold on \(\Gamma\). Exterior outgoing problems correspond
to a limit in which \(\mathcal{Z}_\Gamma\) selects the Sommerfeld branch of
Section~\ref{sec:ch243}; conflating a finite-chamber termination with that limit
without normalization corrupts far-zone patterns \cite{stratton1941}.
Forcing must lie in the range implied by charge conservation. Divergence of the
Amp\`ere row together with Eq.~\eqref{eq:ch1.charge-continuity-harmonic} yields
\begin{equation}
\label{eq:ch2.source-compatibility}
\divg\mathbf{s}_{\omega,2}+\jj\omega\,\mathbf{s}_{\omega,3}=0
\end{equation}
on the second and third blocks of Eq.~\eqref{eq:ch2.maxwell-forcing-vector}.
Inconsistent \((\Jimp,\rho)\) pairs lie outside physically realizable drive;
numerical solvers may still return vectors that violate
Eq.~\eqref{eq:ch2.source-compatibility} \cite{stratton1941}. A compact
current filament with vanishing \(\rho\) automatically satisfies
Eq.~\eqref{eq:ch2.source-compatibility} because \(\mathbf{s}_{\omega,3}=0\).
When a volumetric \(\rho\) is prescribed independently of \(\Jimp\), the pair
must be adjusted so that their divergence balance matches the Amp\`ere--Gauss
rows; otherwise the closed problem is mathematically posed but physically
inconsistent. Feed models that inject \(\Jimp\) without updating \(\rho\) are
common in method-of-moments implementations; they remain valid only when the
discretization enforces continuity on the impressed grid or when the problem
is magnetostatic in the relevant band \cite{harrington2001}.
Linking to Subsection~\ref{sec:ch201}, radiation acts on closed solutions:
\begin{equation}
\label{eq:ch2.maxwell-radiation-composition}
\mathcal{R}_\omega=\Pi_{\mathrm{rad}}\circ\mathcal{S}_\omega
=\Pi_{\mathrm{rad}}\circ\mathcal{L}_\omega^{-1}\circ
\begin{bmatrix}\mathbf{s}_\omega(\cdot)\\ \mathbf{g}_\Gamma\end{bmatrix},
\end{equation}
with \(\mathcal{L}_\omega^{-1}\) understood on the outgoing-admissible branch fixed by
\(\mathcal{B}_\Gamma\) and with \(\Pi_{\mathrm{rad}}\) from
Eq.~\eqref{eq:ch2.radiation-operator-def} \cite{stratton1941}. Completion must precede
radiative selection: applying \(\Pi_{\mathrm{rad}}\) to non-closed states mixes
reactive content into far-zone observables \cite{harrington2001}.
The energy form supports an integration-by-parts identity
\begin{equation}
\label{eq:ch2.maxwell-adjoint-pairing}
\mathfrak{a}_\omega(\mathbf{u}_\omega,\mathbf{v}_\omega)=
\langle\mathcal{M}_\omega[\mathbf{u}_\omega],\mathbf{v}_\omega\rangle_{\Omega}
+\int_\Gamma\big(\hat{\mathbf{n}}\times\E\cdot\H'-\hat{\mathbf{n}}\times\E'\cdot\H\big)dS,
\end{equation}
where \(\langle\cdot,\cdot\rangle_{\Omega}\) is the componentwise \(L^{2}\) pairing
\cite{stratton1941}. Self-adjoint truncations on bounded domains require that
\(\mathcal{B}_\Gamma\) cancel the boundary integrand on admissible traces; criteria
are supplied in Section~\ref{sec:ch214}. Equation
\eqref{eq:ch2.maxwell-adjoint-pairing} is the local spine for those arguments.
Fredholm structure note: for bounded \(\Omega\) with passive \(\mathcal{B}_\Gamma\), \(\mathcal{L}_\omega\) is a Fredholm operator on
\(\mathcal{D}(\mathcal{L}_\omega)\) modulo cavity modes, with index tied to
topology and boundary class. Exterior domains replace discrete spectra with
continuous radiation modes, and \(\mathcal{L}_\omega^{-1}\) becomes a resolvent on
the outgoing subspace rather than an ordinary inverse. Section~\ref{sec:ch263}
develops the spectral/Fredholm foundation; here the point is that invertibility questions
for \(\mathcal{S}_\omega\) are questions about \(\mathcal{L}_\omega\) on the
\emph{declared} domain, not about formal elimination of \(\H\) alone
\cite{harrington2001}.
Linearity of \(\mathcal{L}_\omega\) in \(\mathbf{u}_\omega\) implies that
superposition holds for fixed boundary class: if
\(\mathcal{L}_\omega[\mathbf{u}_\omega^{(a)}]=[\mathbf{s}^{(a)},\mathbf{g}^{(a)}]\)
and \(\mathcal{L}_\omega[\mathbf{u}_\omega^{(b)}]=[\mathbf{s}^{(b)},\mathbf{g}^{(b)}]\)
with identical \(\mathcal{B}_\Gamma\), then
\(\mathcal{L}_\omega[\mathbf{u}_\omega^{(a)}+\mathbf{u}_\omega^{(b)}]=
[\mathbf{s}^{(a)}+\mathbf{s}^{(b)},\mathbf{g}^{(a)}+\mathbf{g}^{(b)}]\).
Changing \(\mathcal{B}_\Gamma\) while holding \(\Jimp\) fixed breaks this
principle because boundary rows are part of the operator, not of the drive.
Laboratory comparisons that swap terminations without re-solving the coupled
system therefore compare different \(\mathcal{L}_\omega\), not different
``modes'' of the same map \cite{stratton1941}. The radiation
operator \(\mathcal{R}_\omega\) inherits the same sensitivity: composition
Eq.~\eqref{eq:ch2.maxwell-radiation-composition} is linear in impressed data
only after \(\mathcal{B}_\Gamma\) and the outgoing branch are held fixed
\cite{stratton1941}.
Subsection~\ref{sec:ch202} has packaged harmonic Maxwell theory as the composite operator
\(\mathcal{L}_\omega\), recorded its resolvent relation to \(\mathcal{S}_\omega\), and
linked radiation composition Eq.~\eqref{eq:ch2.maxwell-radiation-composition} to operator
inversion on the outgoing-admissible branch. Subsection~\ref{sec:ch203} records how those
maps inherit the admissibility, regime, and measurement criteria established in
\ref{ch:01}.
\subsection{Admissibility and Maxwell export sequence}
\label{sec:ch203}
% HARRINGTON-OPEN:sec:ch203
A coax-fed pyramidal horn over a ground plane defines the impressed-current and boundary-data tuple technicians declare before Maxwell export is computed \cite{stratton1941}. Maxwell's retarded export to $\Gamma_R$ remains the governing constraint \cite{maxwell1873}.
Subsections~\ref{sec:ch201}--\ref{sec:ch202} introduced the radiation maps and Maxwell
export sequence; the display below invokes \ref{ch:01} $\mathcal{D}_J=1$
\cite{stratton1941}. The present chapter builds on
established in \ref{ch:01}. Rather, it applies those foundations to radiation problems
\cite{maxwell1873}. The Maxwell-closed criteria, conservation laws, and
measurement principles developed previously remain in force; the present chapter introduces
the operator calculus through which admissible sources are connected to observable
radiative fields \cite{harrington2001}. The operators $\mathcal{S}_\omega$,
$\mathcal{R}_\omega$, and $\mathcal{L}_\omega$ are applied on admissible inputs without
reopening the \ref{ch:01} proofs \cite{stratton1941}. Impressed data enter the
solution domain only when the \ref{ch:01} source conditions hold; radiative
reports attach to \(\mathcal{R}_\omega[\Jimp]=\Pi_{\mathrm{rad}}\mathcal{S}_\omega[\Jimp]\),
not to the full solution pair by default \cite{stratton1941}.
The continuity relations below restate, in operator language, the \ref{ch:01}
regime and measurement criteria that qualify when angular-flux and transport reports may
attach to \(\mathcal{R}_\omega[\Jimp]\) rather than to the unreconstructed solution pair
\cite{stratton1941}. Physically, this means that near-field reactive content diagnosed on
\(\mathcal{S}_\omega[\Jimp]\) must not be reported as exported radiative flux unless the
regime screen \(\mathcal{D}_{J,\mathrm{reg}}\) from \ref{ch:01} is satisfied on
the measurement shell \cite{harrington2001}.
% PHYSICAL-MEMORY:sec:ch203
\begin{physicalvignette}[Beth torsion balance in a microwave anechoic chamber]
Consider Beth torsion balance in a microwave anechoic chamber. Before any coefficient table is printed, technicians record frequency, cable loss, and the mechanical outline that fixes electrical size. The first walk-through of the range is always tactile: flange bolts, probe standoff, turntable boresight, and the wattmeter reference on the receive chain \cite{hansen1988}.
On this hardware the near zone is visually impressive---standing reactive structure, rapid phase rotation, sometimes a visible VSWR ripple at the feed---yet the export question is measured only after outgoing selection on a large sphere \cite{harrington2001}. What the camera sees in the chamber is not automatically what the far-zone pattern claims \cite{stratton1941}.
Engineers therefore learn to separate three readings: what the feed can synthesize on declared support, what propagating content survives on the radiative range, and what a finite instrument can recover at declared tolerance \cite{devaney2004}. Confusing those readings produces the rank gaps and watt failures that later sections quantify \cite{pal2026}.
Rayleigh and Stratton treated such ranges as boundary-value laboratories long before digital near-field transforms became routine \cite{stratton1941}. Harrington's method-of-moments culture reinforced the same discipline: metal outline first, harmonic headline second \cite{harrington2001}.
Subsection~\ref{sec:ch203} states the formal export condition that this range scene presupposes \cite{stratton1941}.
\end{physicalvignette}
\begin{equation}
\label{eq:ch2.continuity-operator-domain}
\mathcal{D}_{\mathrm{ch1}}=1
\ \Longrightarrow\
\Jimp\in\mathcal{J}_\Omega
\ \wedge\
\mathcal{S}_\omega[\Jimp]\in\mathcal{F}_\Omega,
\end{equation}
\begin{equation}
\label{eq:ch2.transport-qualified-radiation}
\mathcal{D}_{J,\mathrm{reg}}=1
\ \Longrightarrow\
\text{angular-flux reports attach to}\ \mathcal{R}_\omega[\Jimp],
\ \text{not to}\ \mathcal{S}_\omega[\Jimp]\ \text{by default},
\end{equation}
% Forward registry: \ref{ch:01} labels imported at Ch.~2 continuity (proofs remain in Ch.~1).
\phantomsection\label{sec:ch01.3.3}
\phantomsection\label{fig:ch1.wp-active-passive-rotation}
% PHYSICAL-MEMORY:sec:ch01.3.3
\begin{physicalvignette}[Outdoor far-field range with mast and ground reflection]
Consider Outdoor far-field range with mast and ground reflection. Before any coefficient table is printed, technicians record frequency, cable loss, and the mechanical outline that fixes electrical size. The first walk-through of the range is always tactile: flange bolts, probe standoff, turntable boresight, and the wattmeter reference on the receive chain \cite{hansen1988}.
On this hardware the near zone is visually impressive---standing reactive structure, rapid phase rotation, sometimes a visible VSWR ripple at the feed---yet the export question is measured only after outgoing selection on a large sphere \cite{harrington2001}. What the camera sees in the chamber is not automatically what the far-zone pattern claims \cite{stratton1941}.
Engineers therefore learn to separate three readings: what the feed can synthesize on declared support, what propagating content survives on the radiative range, and what a finite instrument can recover at declared tolerance \cite{devaney2004}. Confusing those readings produces the rank gaps and watt failures that later sections quantify \cite{pal2026}.
Rayleigh and Stratton treated such ranges as boundary-value laboratories long before digital near-field transforms became routine \cite{stratton1941}. Harrington's method-of-moments culture reinforced the same discipline: metal outline first, harmonic headline second \cite{harrington2001}.
Subsection~\ref{sec:ch01.3.3} states the formal export condition that this range scene presupposes \cite{stratton1941}.
\end{physicalvignette}
\begin{equation}
\label{eq:ch1.local-regime-decomposition}
\Pi(\theta,\phi;r)
=
\Pi_{\mathrm{a}}(\theta,\phi;r)+\Pi_{\mathrm{re}}(\theta,\phi;r),
\qquad
\Pi_{\mathrm{a}}=\Re\{\mathbf{S}\cdot\hat{\mathbf{n}}\},
\end{equation}
with reactive sector \(\Pi_{\mathrm{re}}\) the remainder on the declared control surface.
\begin{equation}
\label{eq:ch1.reactive-flux-sphere}
Q_{\mathrm{re}}(r)
=
\Im\oint_{S_r}\mathbf{S}\cdot\hat{\mathbf{r}}\,dS,
\qquad
P_{\mathrm{rad}}(r)
=
\Re\oint_{S_r}\mathbf{S}\cdot\hat{\mathbf{r}}\,dS,
\end{equation}
\begin{equation}
\label{eq:ch1.regime-index}
\chi(r)
=
\frac{|Q_{\mathrm{re}}(r)|}{P_{\mathrm{rad}}(r)+\epsilon_P},
\qquad \epsilon_P>0.
\end{equation}
\begin{equation}
\label{eq:ch1.radiative-regime-condition}
P_{\mathrm{rad}}(r)>0
\ \wedge\
\chi(r)\le\chi_{\mathrm{th}},
\end{equation}
\begin{equation}
\label{eq:ch1.crossover-radius}
r_c:\ \chi(r_c)=\chi_{\mathrm{th}},
\end{equation}
\begin{equation}
\label{eq:ch1.regime-qualified-decision}
\mathcal{D}_{J,\mathrm{reg}}(r)=1
\iff
\Big(\chi(r)\le\chi_{\mathrm{th}}\ \wedge\ \mathcal{D}_J(r)=1\Big),
\end{equation}
\begin{equation}
\label{eq:ch1.rotation-power-invariance}
\oint_{S_r}\Re\{\mathbf{S}_R\cdot\hat{\mathbf{r}}\}\,dS
=
\oint_{S_r}\Re\{\mathbf{S}\cdot\hat{\mathbf{r}}\}\,dS,
\end{equation}
\begin{equation}
\label{eq:ch1.scalar-observable-invariance}
\mathcal{I}\big[\mathcal{P}_R\mathcal{F}\big]=\mathcal{I}[\mathcal{F}],
\end{equation}
for scalar observables \(\mathcal{I}\) under passive frame rotation \(\mathcal{P}_R\).
\begin{equation}
\label{eq:ch1.covariance-residual}
\mathcal{R}_{\mathcal{F}}
=
\left\|
\mathcal{R}_\omega\big[\mathcal{T}_{R,\mathbf{a}}\Jimp\big]
-
\mathcal{T}_{R,\mathbf{a}}\,\mathcal{R}_\omega[\Jimp]
\right\|,
\end{equation}
\begin{equation}
\label{eq:ch1.intrinsic-angular-observable}
\mathbf{J}_{\mathrm{int}}
=
\mathbf{J}_{\mathrm{em}}
-
(\mathbf{r}-\mathbf{r}_c)\times\mathbf{P}_{\mathrm{tot}},
\end{equation}
\begin{equation}
\label{eq:ch1.transverse-components}
\E_\perp
=
\E-(\hat{\mathbf{r}}\cdot\E)\hat{\mathbf{r}},
\qquad
\H_\perp
=
\H-(\hat{\mathbf{r}}\cdot\H)\hat{\mathbf{r}},
\end{equation}
\begin{equation}
\label{eq:ch1.stokes-local}
S_0=|E_x|^2+|E_y|^2,\quad
S_1=|E_x|^2-|E_y|^2,\quad
S_2=2\Re(E_x E_y^*),\quad
S_3=2\Im(E_x E_y^*),
\end{equation}
\begin{equation}
\label{eq:ch1.poynting-vorticity}
\Omega_P
=
\hat{\mathbf{n}}\cdot(\nabla\times\mathbf{P}),
\qquad
\mathbf{P}=\Re\{\mathbf{S}\},
\end{equation}
where \(\mathcal{D}_{J,\mathrm{reg}}\) is the regime flag from
Eq.~\eqref{eq:ch1.regime-qualified-decision}.
The laboratory chain composes inherited \ref{ch:01} constraints with radiation
selection and finite observation:
\begin{equation}
\label{eq:ch2.continuity-composition-law}
\mathfrak{A}_{\mathrm{ch2}}^{\mathrm{(report)}}
=
\mathfrak{A}_{\mathrm{meas}}
\circ
\mathfrak{A}_{\mathrm{g}}
\circ
\mathfrak{A}_{\mathrm{cov}}
\circ
\mathfrak{A}_{\mathrm{tr}}
\circ
\Pi_{\mathrm{rad}}
\circ
\mathcal{S}_\omega,
\end{equation}
\begin{equation}
\label{eq:ch2.operational-measurement-chain}
\Lambda_{\mathrm{lab}}
=
\mathcal{O}_{\mathrm{meas}}\circ\mathcal{O}_{\mathrm{far}}\circ\mathcal{R}_\omega,
\qquad
\mathcal{Q}\in\{\text{direct},\ \text{surrogate},\ \text{inference}\},
\end{equation}
% Missing labeled equations restored at natural narrative anchors (math frozen).
\begin{equation}
\label{eq:ch2.covariance-operator-test}
\mathcal{R}_\omega\big[\mathcal{T}_{R,\mathbf{a}}\Jimp\big]
=
\mathcal{T}_{R,\mathbf{a}}\,\mathcal{R}_\omega[\Jimp]
\quad\text{on admissible data,}
\end{equation}
with residual Eq.~\eqref{eq:ch1.covariance-residual} serving as the numerical diagnostic when
discretization breaks Eq.~\eqref{eq:ch2.covariance-operator-test} \cite{stratton1941}.
Function spaces, Green functions, and radiation operators exist because Maxwell field
specification requires them---not as parallel formulations \cite{maxwell1873}. Section~\ref{sec:ch20}
has established the radiation maps, Maxwell export sequence, and \ref{ch:01}
continuity; Section~\ref{sec:ch21} now supplies the Hilbert and Sobolev apparatus through
which those maps acquire precise domains, traces, and adjoints \cite{stratton1941}.
\section{Function Spaces and Maxwell Operators}
\label{sec:ch21}
Every practical electromagnetic system occupies finite space, is driven by finite sources,
and is observed through finite apertures \cite{stratton1941}. These physical
limitations impose restrictions on the fields that can be generated, measured, and transported.
Before the Sobolev symbols \(\HCurl(\Omega)\) and \(L^{2}(\Omega)\) appear, ask what kind of
electromagnetic object lives here: not an abstract Hilbert vector, but a horn aperture, a slot
on a ground plane, a finite feed current, or a near-field scan surface whose tangential
\(\E\) must be the trace of a volume field \cite{stratton1941}. A horn feed launches aperture currents whose rim
singularity is visible in near-zone scans; a patch or slot on a ground plane carries
sheet-like \(\Jimp\) in method-of-moments assembly; a finite-element port model enforces
tangential \(\hat{\mathbf{n}}\times\E\) on conductors while impressing currents drawn from
\(\mathcal{J}_\Omega\). Each construction asks whether a field pair \((\E,\H)\) can satisfy
Maxwell's equations and match what probes record on \(\Gamma=\partial\Omega\). The solution
map \(\mathcal{S}_\omega:\mathcal{J}_\Omega\to\mathcal{F}_\Omega\), the composite Maxwell
operator \(\mathcal{L}_\omega\) on \(\mathcal{D}(\mathcal{L}_\omega)\), and the radiation
map \(\mathcal{R}_\omega=\Pi_{\mathrm{rad}}\circ\mathcal{S}_\omega\) were named in
Section~\ref{sec:ch201}; here \(\mathcal{F}_\Omega\) is the admissible field space and
\(\Pi_{\mathrm{rad}}\) the outgoing projection that qualifies far-zone content
\cite{harrington2001}.
Function spaces provide the mathematical language through which those restrictions are
expressed. They do not replace Maxwell's equations; rather, they specify the class of fields
and sources to which Maxwell's equations may be consistently applied
\cite{maxwell1873}.
% PHYSICAL-MEMORY:sec:ch21
\begin{physicalvignette}[Phased array panel on a positioner]
Consider Phased array panel on a positioner. Before any coefficient table is printed, technicians record frequency, cable loss, and the mechanical outline that fixes electrical size. The first walk-through of the range is always tactile: flange bolts, probe standoff, turntable boresight, and the wattmeter reference on the receive chain \cite{hansen1988}.
On this hardware the near zone is visually impressive---standing reactive structure, rapid phase rotation, sometimes a visible VSWR ripple at the feed---yet the export question is measured only after outgoing selection on a large sphere \cite{harrington2001}. What the camera sees in the chamber is not automatically what the far-zone pattern claims \cite{stratton1941}.
Engineers therefore learn to separate three readings: what the feed can synthesize on declared support, what propagating content survives on the radiative range, and what a finite instrument can recover at declared tolerance \cite{devaney2004}. Confusing those readings produces the rank gaps and watt failures that later sections quantify \cite{pal2026}.
Rayleigh and Stratton treated such ranges as boundary-value laboratories long before digital near-field transforms became routine \cite{stratton1941}. Harrington's method-of-moments culture reinforced the same discipline: metal outline first, harmonic headline second \cite{harrington2001}.
Subsection~\ref{sec:ch21} states the formal export condition that this range scene presupposes \cite{stratton1941}.
\end{physicalvignette}
Section~\ref{sec:ch21} now supplies the Hilbert and Sobolev apparatus---domain, range, adjoint,
boundedness, and completion---through which finite source extent and finite observation apertures
appear as accessibility constraints on Maxwell maps rather than as after-the-fact modeling
choices \cite{stratton1941}. Spectral decomposition enters only after outgoing selection; until then, avoid spectral
decay or modal rank; those consequences of compact geometry are developed once in
Section~\ref{sec:ch26} and invoked thereafter by reference \cite{chu1948}. The finite-energy
criterion of Eq.~\eqref{eq:ch1.energy-class} and admissible state class
Eq.~\eqref{eq:ch1.admissible-space} from Subsection~\ref{sec:ch01.0.3} remain the physical
floor. Subsection~\ref{sec:ch01.1.2} fixed source tuples, boundary operators, and charge
continuity; every space declaration below must remain compatible with
Eq.~\eqref{eq:ch2.continuity-operator-domain} and must not reopen \ref{ch:01}
proofs \cite{stratton1941}. The guiding constraint from Subsection~\ref{sec:ch203} persists
without restatement here.
Subsection~\ref{sec:ch211} defines source, field, and dual spaces for phasor
Maxwell theory, including distributional sources and trace spaces on
\(\Gamma=\partial\Omega\), importing finite-energy charge-continuous from
Section~\ref{sec:ch01.0.3} and source--boundary coupling from
Section~\ref{sec:ch01.1.2}. Subsection~\ref{sec:ch212} treats linear operators on
those spaces: domain, range, graph norm, and adjoint pairing compatible with
Eq.~\eqref{eq:ch2.maxwell-adjoint-pairing}. Subsection~\ref{sec:ch213} develops
boundedness and compactness criteria, rank deficiency of observation maps, and the
spectral compression viewpoint developed for radiation operators in
Section~\ref{sec:ch26}. Subsection~\ref{sec:ch214} addresses self-adjoint
realizations, Friedrichs and other completions, and the $\mathcal{D}_J=1$ constraints
under which truncated Maxwell operators inherit coercivity or passivity on bounded
domains \cite{harrington2001}. Several constructions deliberately remain outside Section~\ref{sec:ch21}. Dyadic
Green kernels and retarded inverses belong to Section~\ref{sec:ch23}; Sommerfeld
outgoing selection and exterior uniqueness to Section~\ref{sec:ch24}; radiation
integrals and flux operators at the radiative boundary to Section~\ref{sec:ch25}.
Section~\ref{sec:ch21} prepares the domains on which those maps act. Attempting to
evaluate \(\mathcal{L}_\omega^{-1}\) through a Green function before trace spaces
and dual pairings are fixed produces kernel formulas that ignore boundary rows and
corrupt flux identities \cite{stratton1941}. Conversely, staying at the
\(L^{2}\)-level of Eq.~\eqref{eq:ch2.operator-domain-declaration} without
Sobolev refinement is adequate for uniqueness statements on smooth domains but
insufficient for sheet sources, impedance boundaries, and finite-element
discretizations where \(\curl\) and \(\hat{\mathbf{n}}\times\) must be closed
simultaneously \cite{harrington2001}.
The energy bilinear form Eq.~\eqref{eq:ch2.maxwell-energy-form} and weak
formulation Eq.~\eqref{eq:ch2.maxwell-weak-form} from Subsection~\ref{sec:ch202}
anchor the Hilbert-space choices. Section~\ref{sec:ch21} identifies the minimal
completion in which \(\mathfrak{a}_\omega\) is continuous and coercive modulo
cavity modes, and in which the forcing functional \(\mathfrak{f}_\omega\) induced
by \(\mathbf{s}_\omega(\mathfrak{D})\) is bounded on test states with homogeneous
traces. Well-posedness of the weak problem on bounded domains then follows from
Lax--Milgram arguments once null spaces aligned with divergence-free resonances are quotiented; the quotient construction is owned by Subsection~\ref{sec:ch214}
and feeds the cavity-mode analysis of Section~\ref{sec:ch263}
\cite{stratton1941}. Exterior domains require outgoing subspaces rather
than full Hilbert self-adjointness; that distinction is preserved throughout so
that interior cavity spectra are not conflated with continuous radiation spectra
\cite{harrington2001}.
Rank deficiency enters at two levels that Section~\ref{sec:ch21} must keep separate.
Instrument maps---finite apertures, modal truncations, probe coupling---are
finite-rank operators on field space even when \(\mathcal{L}_\omega\) is
injective on the physical branch. Equivalent source distributions in
\(\ker\mathcal{R}_\omega\) produce identical far-zone patterns while differing in
reactive near fields; that null-space structure is analyzed in
Section~\ref{sec:ch27}, but the compact embedding lemmas used there are developed
in Subsection~\ref{sec:ch213}. Measurement surrogates from
Section~\ref{sec:ch01.5} therefore appear as rank-deficient compositions on the
right of Eq.~\eqref{eq:ch2.continuity-composition-law}, not as modifications of
Maxwell law \cite{stratton1941}. Both are
valid only when the discrete spaces inherit the continuity and trace structure
defined in Subsections~\ref{sec:ch211}--\ref{sec:ch212}. Section~\ref{sec:ch21}
states the continuous targets those discretizations approximate; it does not
substitute for mesh convergence analysis \cite{harrington2001}.
Material inhomogeneity and interface conditions from
Eq.~\eqref{eq:ch1.interface-conditions} enter through piecewise constant
coefficients in \(\mathcal{M}_\omega\) without changing the space template:
fields remain in product Sobolev spaces with jump conditions enforced on traces.
Multi-region formulations in Section~\ref{sec:ch234} rely on the interface trace
classes introduced in Subsection~\ref{sec:ch211}. Perfect conductors impose
homogeneous tangential \(\E\); impedance and Leontovich boundaries modify the
boundary pairing in Eq.~\eqref{eq:ch2.maxwell-adjoint-pairing} and are treated as
closure operators in Subsection~\ref{sec:ch214} \cite{stratton1941}. Subsection~\ref{sec:ch214} completes the function-space foundation begun here with
self-adjoint truncations and Maxwell-closed completions that connect directly to the
spectral/Fredholm program of Section~\ref{sec:ch263}. The subsection sequence
begins with the space definitions that make \(\mathcal{D}(\mathcal{L}_\omega)\) in
Eq.~\eqref{eq:ch2.operator-domain-declaration} precise enough for the Green,
radiation, and spectral sections that follow.
\subsection{Spaces for sources, fields, and distributions}
\label{sec:ch211} 
% HARRINGTON-OPEN:sec:ch211
Function spaces describe admissible currents and fields before any antenna synthesis \cite{stratton1941}. Compact support and finite power are laboratory facts, not optional regularity \cite{maxwell1873}. Maxwell's retarded export to $\Gamma_R$ remains the governing constraint \cite{maxwell1873}.
The physical problem is source realizability on finite hardware: every antenna campaign begins with currents and boundary data that occupy bounded regions, carry finite power, and meet conductors where tangential \(\E\) or \(\H\) is measured \cite{stratton1941}. A horn feed, patch edge, or waveguide flange excites currents that are not smooth functions on all of $\mathbb{R}^{3}$: sheet sources, rim currents, and port terminals appear as distributions whose support is bounded yet whose derivatives fail at edges \cite{stratton1941}. Without the Sobolev and trace spaces declared here, one obtains Green formulas that ignore boundary rows, mis-state flux through $\Gamma=\partial\Omega$, and corrupt near-field transforms on spherical scanners \cite{harrington2001}. 
Function spaces describe the admissible currents and fields from which any antenna
synthesis must begin \cite{stratton1941}. Sobolev and trace classes encode the physical
constraints imposed by finite power, bounded supports, conducting boundaries, and
realizable feed structures \cite{harrington2001}. A synthesis algorithm that ignores these
classes may produce mathematically sparse currents that no laboratory feed can realize
\cite{devaney2004}.
% PHYSICAL-MEMORY:sec:ch211
\begin{physicalvignette}[Yagi array aimed at a fixed receive horn]
Consider Yagi array aimed at a fixed receive horn. Before any coefficient table is printed, technicians record frequency, cable loss, and the mechanical outline that fixes electrical size. The first walk-through of the range is always tactile: flange bolts, probe standoff, turntable boresight, and the wattmeter reference on the receive chain \cite{hansen1988}.
On this hardware the near zone is visually impressive---standing reactive structure, rapid phase rotation, sometimes a visible VSWR ripple at the feed---yet the export question is measured only after outgoing selection on a large sphere \cite{harrington2001}. What the camera sees in the chamber is not automatically what the far-zone pattern claims \cite{stratton1941}.
Engineers therefore learn to separate three readings: what the feed can synthesize on declared support, what propagating content survives on the radiative range, and what a finite instrument can recover at declared tolerance \cite{devaney2004}. Confusing those readings produces the rank gaps and watt failures that later sections quantify \cite{pal2026}.
Function spaces describe admissible currents and fields before any antenna synthesis . Compact support and finite power are laboratory facts, not optional regularity .
Rayleigh and Stratton treated such ranges as boundary-value laboratories long before digital near-field transforms became routine \cite{stratton1941}. Harrington's method-of-moments culture reinforced the same discipline: metal outline first, harmonic headline second \cite{harrington2001}.
Subsection~\ref{sec:ch211} states the formal export condition that this range scene presupposes \cite{stratton1941}.
\end{physicalvignette}
A PEC-backed slot fed by $\Jimp\in L^2(\Omega_s)$ and an idealized $1/r$ current template
illustrate the admissibility distinction at fixed frequency. Only the compact $L^2$ source
class supports finite radiated power; divergent templates violate the finite-energy
$\mathcal{D}_J=1$ conditions inherited from Chapter~1 \cite{stratton1941}. Compact support and
finite power are laboratory facts, not optional regularity \cite{maxwell1873}.
Method-of-moments and finite-element antenna models represent impressed currents on
surface patches and volume filaments, then test Maxwell's equations with localized basis
functions \cite{harrington2001}. That procedure requires square-integrable
\(\curl\) fields, well-defined tangential traces on conductors and apertures, and
distributional extensions when sources collapse to sheets---precisely the analytic gaps that
naive \(L^{2}\) current models leave open. Subsections~\ref{sec:ch201}--\ref{sec:ch202}
declared minimal spaces \(\mathcal{F}_\Omega\) and \(\mathcal{J}_\Omega\) adequate to
\emph{name} the maps \(\mathcal{S}_\omega\), \(\mathcal{R}_\omega\), and
\(\mathcal{L}_\omega\); Subsection~\ref{sec:ch211} supplies the Sobolev, trace, and
distributional structure those declarations presuppose. Finite-energy Maxwell-closed from
Subsection~\ref{sec:ch01.0.3} and the source--boundary tuple of
Subsection~\ref{sec:ch01.1.2} remain authoritative; embeds
them in spaces where \(\curl\), \(\divg\), and tangential restriction are
continuous maps with explicit domains \cite{stratton1941}. Nothing
here re-proves charge continuity, interface jumps, or outer-boundary classes already
fixed in \ref{ch:01}. Let \(\Omega\subset\mathbb{R}^{3}\) be a bounded Lipschitz domain or a finite
union of Lipschitz patches \(\{\Omega_a\}\) with outward boundary
\(\Gamma=\partial\Omega\). Phasor fields use the \(\exp(-\jj\omega t)\) convention
throughout.
In the display below, \(L^{2}(\Omega)\) denotes square-integrable field energy on \(\Omega\);
\(\HCurl(\Omega)\) and \(\Hdiv(\Omega)\) are the minimal vector Sobolev classes in which
curl-coupled Maxwell balances close weakly; \(\mathcal{J}_\Omega\) collects impressed
currents of compact support and finite power; \(\mathcal{F}_\Omega\) is the admissible
Maxwell field class; and \(\mathcal{S}_\omega\), \(\mathcal{R}_\omega\), and
\(\mathcal{L}_\omega\) are the solution, radiation, and composite operators declared in
Section~\ref{sec:ch201} \cite{stratton1941}.
Square-integrability alone is too weak for curl-coupled Maxwell theory: one needs
control of both the field and its curl so that circulation balances are meaningful
in the \(L^{2}\) sense.
An engineer measuring tangential \(\E\) on a patch boundary is, in effect, asking whether
the recorded trace is the restriction of a field with square-integrable curl; a
method-of-moments code that tests \(\curl\E\) against rooftop functions presupposes the
same control. The space \(\HCurl(\Omega)\) names precisely that charge-continuous class:
elements of \(L^{2}(\Omega)^{3}\) whose curls also lie in \(L^{2}(\Omega)^{3}\), so
Faraday and Amp\`ere balances can be posed weakly without assuming pointwise smoothness
at horn rims or slot edges \cite{stratton1941}. Readers should notice that
\(\HCurl(\Omega)\) is stronger than \(L^{2}\) alone yet weaker than requiring classical
\(C^{1}\) fields everywhere; it is the minimal completion in which curl-coupled
Maxwell operators close on bounded \(\Omega\) \cite{harrington2001}. The companion
\(\Hdiv(\Omega)\) plays the parallel role for Gauss constraints when divergence must be
controlled in the \(L^{2}\) sense. With that physical reading fixed, the formal definitions
are:
Single-vector Sobolev classes encode the minimal regularity that curl-coupled Maxwell
theory requires on bounded domains. Membership in
\(\HCurl(\Omega)\) controls both field amplitude and circulation in the \(L^{2}\) sense,
whereas plain \(L^{2}\) control would not guarantee that Faraday and Amp\`ere balances
can be posed weakly at horn rims or slot edges \cite{stratton1941}. From a measurement
standpoint, these spaces establish that tangential fields recorded on \(\Gamma\) are traces
of physically admissible volume states, not arbitrary boundary templates
\cite{harrington2001}. Single-vector Sobolev classes are
\begin{equation}
\label{eq:ch2.hcurl-space}
\HCurl(\Omega)
=
\left\{
\mathbf{u}\in L^{2}(\Omega)^{3}:
\curl\mathbf{u}\in L^{2}(\Omega)^{3}
\right\},
\end{equation}
\begin{equation}
\label{eq:ch2.hdiv-space}
\Hdiv(\Omega)
=
\left\{
\mathbf{u}\in L^{2}(\Omega)^{3}:
\divg\mathbf{u}\in L^{2}(\Omega)
\right\},
\end{equation}
where \(\divg\mathbf{u}\) is the divergence of \(\mathbf{u}\) (m\(^{-2}\) for current-like
fields, C\,m\(^{-3}\) when \(\mathbf{u}\) represents a displacement-related quantity),
required to be square integrable so that Gauss-type constraints can be enforced
weakly.
each equipped with the graph norm
\(\|\mathbf{u}\|_{\HCurl}^{2}=\|\mathbf{u}\|_{L^{2}}^{2}+\|\curl\mathbf{u}\|_{L^{2}}^{2}\)
and the analogous \(\Hdiv\) norm \cite{stratton1941}. Equations
\eqref{eq:ch2.hcurl-space}--\eqref{eq:ch2.hdiv-space} are the vector analogues of
the finite-energy floor Eq.~\eqref{eq:ch1.energy-class}: they permit Maxwell
curl-coupling without assuming pointwise smoothness of \(\E\) or \(\H\). The consequence
for antenna modeling is direct: horn rims, slot edges, and patch boundaries may host
singularities, yet the impressed current and recovered field remain admissible as long as
they lie in the declared Sobolev classes \cite{harrington2001}. For
piecewise constant \(\varepsilon,\mu\), fields are understood patchwise in
Eqs.~\eqref{eq:ch2.hcurl-space}--\eqref{eq:ch2.hdiv-space} with interface traces
constrained by Eq.~\eqref{eq:ch1.interface-conditions} \cite{harrington2001}.
The Maxwell state space is the product
\begin{equation}
\label{eq:ch2.maxwell-product-space}
\mathcal{U}_\omega
=
\HCurl(\Omega)\times\HCurl(\Omega),
\qquad
\mathbf{u}_\omega=(\E,\H)\in\mathcal{U}_\omega.
\end{equation}
Equation~\eqref{eq:ch2.maxwell-product-space} refines the six-component packaging
Eq.~\eqref{eq:ch2.maxwell-state-vector}: both components carry square-integrable
curls, matching the first-order structure of
Eq.~\eqref{eq:ch1.maxwell-operator-block}. A weighted energy norm compatible with
Eq.~\eqref{eq:ch2.maxwell-energy-form} is
\begin{equation}
\label{eq:ch2.field-energy-norm}
\|\mathbf{u}_\omega\|_{\mathcal{U}}^{2}
=
\|\E\|_{L^{2}}^{2}+\|\H\|_{L^{2}}^{2}
+\alpha\,\|\curl\E\|_{L^{2}}^{2}+\beta\,\|\curl\H\|_{L^{2}}^{2},
\end{equation}
with positive weights \(\alpha,\beta\) fixed by constitutive scaling in each patch.
Norm~\eqref{eq:ch2.field-energy-norm} controls the bilinear form
\(\mathfrak{a}_\omega\) of Eq.~\eqref{eq:ch2.maxwell-weak-form} on bounded domains
and makes \(\mathcal{U}_\omega\) a Hilbert space \cite{stratton1941}.
\begin{figure}[t]
\centering
\ChOneFigBlock{%
\begin{tikzpicture}[ch1fig, node distance=5mm and 7mm]
 \node[ch1box] (l2) {$L^{2}(\Omega)^{3}$};
 \node[ch1box, above=5mm of l2] (hc) {$\HCurl(\Omega)$};
 \node[ch1box, below=5mm of l2] (hd) {$\Hdiv(\Omega)$};
 \node[ch1box, right=18mm of hc] (tr) {$\mathcal{T}_\Gamma$};
 \draw[ch1flow] (l2.north) -- (l2.north |- hc.south);
 \draw[ch1flow] (l2.north |- hc.south) -- (hc.south);
 \draw[ch1flow] (l2.south) -- (l2.south |- hd.north);
 \draw[ch1flow] (l2.south |- hd.north) -- (hd.north);
 \draw[ch1flow] (hc.east) -- node[above,font=\scriptsize]{$\gamma_\tau$} (tr.west);
\end{tikzpicture}%
}
\caption{Embedding ladder for Maxwell fields: \(\HCurl\) and \(\Hdiv\) sit above
\(L^{2}\); tangential traces map \(\HCurl\) fields to boundary data compatible with
Eq.~\eqref{eq:ch1.boundary-data-operator}.}
\label{fig:ch2.sobolev-embedding-ladder}
Sobolev traces, adjoints, and Fredholm structure establish Maxwell boundary data and export rank---they are not optional ornaments on the source-to-field pathway \cite{stratton1941}.
\end{figure}
Divergence control enters through source spaces and through Gauss constraints rather
than through requiring both components to lie in \(\Hdiv\); forcing \(\E\in\Hdiv\)
and \(\H\in\Hdiv\) simultaneously is stronger than Maxwell's first-order system
requires and is reserved for scalar-potential reductions deferred to later chapters
\cite{stratton1941}. Impressed sources refine the declaration of Subsection~\ref{sec:ch201}.
The admissible impressed-current space collects currents of finite support and
finite power that satisfy the \ref{ch:01} $\mathcal{D}_J=1$ screen.
Let
\(\Omega_s\subseteq\Omega\) denote compact source support. The admissible impressed
current space is
\begin{equation}
\label{eq:ch2.source-space-refined}
\mathcal{J}_\Omega
=
\left\{
\Jimp\in \HCurl(\Omega):
\operatorname{supp}\Jimp\subseteq\Omega_s,\
\mathcal{D}_{J}[\Jimp]=1
\right\},
\end{equation}
where \(\mathcal{D}_{J}\) is the criterion of
Eq.~\eqref{eq:ch2.source-admissibility} inherited from
Eq.~\eqref{eq:ch1.continuity-charge-continuous} and
Eq.~\eqref{eq:ch1.source-admissibility} \cite{stratton1941}. Equation
\eqref{eq:ch2.source-space-refined} strengthens square integrability to
\(\HCurl\) regularity so that \(\curl\Jimp\) is defined in \(L^{2}\) whenever
numerical or analytical measurement sequences differentiate sources. Charge-neutral subclasses
used in radiation null-space analysis satisfy
\begin{equation}
\label{eq:ch2.charge-neutral-source}
\mathcal{J}_\Omega^{(0)}
=
\left\{
\Jimp\in\mathcal{J}_\Omega:
\divg\Jimp=0
\right\},
\end{equation}
with \(\rho=0\) implied on \(\Omega_s\). Equation~\eqref{eq:ch2.charge-neutral-source}
is the subspace on which divergence-free reparametrizations act without altering
\(\mathcal{S}_\omega[\Jimp]\) in simply connected regions; null-space structure is
developed in Section~\ref{sec:ch272} \cite{harrington2001}.
Boundary data from the tuple Eq.~\eqref{eq:ch1.data-specification-tuple} live in
trace spaces rather than in volume Sobolev classes. Define
\begin{equation}
\label{eq:ch2.trace-space-definition}
\mathcal{T}_\Gamma
=
H^{-1/2}(\Gamma)^{3}\times H^{-1/2}(\Gamma)^{3},
\end{equation}
as the product space for tangential electric and magnetic traces induced by
\(\HCurl(\Omega)\) fields on Lipschitz \(\Gamma\) \cite{stratton1941}. The
tangential trace operators extend Eq.~\eqref{eq:ch2.trace-restriction}:
\begin{equation}
\label{eq:ch2.tangential-trace-maps}
\gamma_{\tau,E}:\HCurl(\Omega)\to H^{-1/2}(\Gamma)^{3},
\quad
\gamma_{\tau,E}[\E]=\hat{\mathbf{n}}\times\E|_{\Gamma},
\end{equation}
\begin{equation}
\label{eq:ch2.tangential-trace-magnetic}
\gamma_{\tau,H}:\HCurl(\Omega)\to H^{-1/2}(\Gamma)^{3},
\quad
\gamma_{\tau,H}[\H]=\hat{\mathbf{n}}\times\H|_{\Gamma},
\end{equation}
with outward unit normal \(\hat{\mathbf{n}}\). Equations
\eqref{eq:ch2.tangential-trace-maps}--\eqref{eq:ch2.tangential-trace-magnetic}
make Eq.~\eqref{eq:ch1.boundary-data-operator} an operator between declared
Banach spaces: \(\mathcal{B}_\Gamma:\mathcal{T}_\Gamma\to\mathcal{Y}_\Gamma\) for
data space \(\mathcal{Y}_\Gamma\) fixed by the boundary class of
Eq.~\eqref{eq:ch1.outer-boundary-classes} or outgoing traces of
Eq.~\eqref{eq:ch1.radiation-trace-condition} \cite{stratton1941}.
\begin{figure}[t]
\centering
\ChOneFigBlock{%
\begin{tikzpicture}[ch1fig, node distance=5mm and 7mm]
 \node[ch1box] (vol) {Volume\\$\HCurl(\Omega)$}; \node[ch1box, right=16mm of vol] (tr) {Traces\\$\mathcal{T}_\Gamma$};
 \node[ch1box, right=16mm of tr] (dat) {Data\\$\mathcal{Y}_\Gamma$};
 \draw[ch1flow] (vol.east) -- node[above,font=\scriptsize]{$\gamma_\tau$} (tr.west);
 \draw[ch1flow] (tr.east) -- node[above,font=\scriptsize]{$\mathcal{B}_\Gamma$} (dat.west);
\end{tikzpicture}%
}
\caption{Trace--data chain for boundary completion: volume fields map to
\(\mathcal{T}_\Gamma\); the boundary operator from
Eq.~\eqref{eq:ch1.boundary-data-operator} selects admissible data.}
\label{fig:ch2.trace-boundary-chain}
\end{figure}
Figure~\ref{fig:ch2.maxwell-boundary-stack}.
Surjectivity of \(\gamma_\tau\) onto
compatible trace subclasses fails globally---only fields in specified
\(\HCurl\) subspaces with prescribed divergence pairings map to physically realizable
\(\mathbf{g}_\Gamma\) \cite{stratton1941}. Subsection~\ref{sec:ch214} identifies
self-adjoint boundary realizations; here the point is that \(\mathcal{L}_\omega\) of
Eq.~\eqref{eq:ch2.maxwell-boundary-composite} acts on \(\mathcal{U}_\omega\) only when
\(\gamma_\tau[\mathbf{u}_\omega]\) lies in the domain of \(\mathcal{B}_\Gamma\). Sheet currents and interface channels from Eq.~\eqref{eq:ch1.interface-conditions}
require distributional extensions. Let \(\mathcal{D}'(\Omega)^{3}\) denote
temporal-harmonic distributional three-vector fields. Impressed volume data remain
\(\Jimp\in\mathcal{J}_\Omega\), while surface equivalents \(\mathbf{J}_{s,ab}\) on
\(\Sigma_{ab}\) embed as
\begin{equation}
\label{eq:ch2.distributional-source-extension}
\langle\mathbf{J}_{s},\boldsymbol{\phi}\rangle
=
\sum_{ab}\int_{\Sigma_{ab}}\mathbf{J}_{s,ab}\cdot\boldsymbol{\phi}\,dS,
\qquad
\boldsymbol{\phi}\in C_c^\infty(\Omega)^{3},
\end{equation}
linking the data tuple Eq.~\eqref{eq:ch1.data-specification-tuple} to a continuous
linear functional on smooth test fields \cite{stratton1941}. Equation
\eqref{eq:ch2.distributional-source-extension} does not alter the jump laws of
Eq.~\eqref{eq:ch1.interface-conditions}; it specifies the functional-analytic class
in which those jumps appear in weak formulations. Interface trace spaces on
\(\Sigma_{ab}\) are denoted
\begin{equation}
\label{eq:ch2.interface-trace-space}
\mathcal{T}_{\Sigma_{ab}}
=
H^{-1/2}(\Sigma_{ab})^{3},
\end{equation}
with tangential continuity of \(\E\) and jump of \(\hat{\mathbf{n}}\times\H\) encoded
as constraints on \(\gamma_\tau\) taken from each patch \cite{harrington2001}.
Dual pairings for source and test fields use the \(L^{2}\) inner product
\begin{equation}
\label{eq:ch2.dual-source-pairing}
\langle\Jimp,\boldsymbol{\psi}\rangle_{\Omega}
=
\int_\Omega \Jimp\cdot\boldsymbol{\psi}^{*}\,d^{3}r,
\qquad
\boldsymbol{\psi}\in \HCurl(\Omega),
\end{equation}
which induces a bounded functional on test states when \(\Jimp\in L^{2}(\Omega)^{3}\)
and extends to Eq.~\eqref{eq:ch2.distributional-source-extension} for sheet data.
Forcing in Eq.~\eqref{eq:ch2.maxwell-forcing-vector} enters the weak form
Eq.~\eqref{eq:ch2.maxwell-weak-form} through
\(\mathfrak{f}_\omega(\mathbf{v}_\omega)=\langle\mathbf{s}_\omega,\mathbf{v}_\omega\rangle\)
once \(\mathbf{v}_\omega\in\mathcal{U}_\omega\) carries homogeneous boundary traces
\cite{stratton1941}. The symbol \(\mathcal{F}_\Omega\) names field pairs
\(\mathbf{u}_\omega=(\E,\H)\) that are not merely square integrable but carry controlled
curls and admissible tangential traces on \(\Gamma\); every output of
\(\mathcal{S}_\omega[\Jimp]\) on the physical branch must lie in this sharpened class
\cite{harrington2001}. The admissible Maxwell field space compatible with \ref{ch:01} is
\begin{equation}
\label{eq:ch2.admissible-field-space}
\mathcal{F}_\Omega
=
\left\{
\mathbf{u}_\omega=(\E,\H)\in\mathcal{U}_\omega:
\gamma_\tau[\mathbf{u}_\omega]\in\mathcal{D}\big(\mathcal{B}_\Gamma\big),\
(\E,\H)\ \text{satisfy patchwise Gauss constraints}
\right\},
\end{equation}
where \(\mathcal{D}(\mathcal{B}_\Gamma)\) is the domain of boundary data compatible
with Eq.~\eqref{eq:ch1.admissible-data-class} imported from
Section~\ref{sec:ch01.1.2}. Equation~\eqref{eq:ch2.admissible-field-space} refines
the minimal declaration of Subsection~\ref{sec:ch201}: same symbol
\(\mathcal{F}_\Omega\), sharpened topology. Membership implies
Eq.~\eqref{eq:ch1.energy-class} but adds curl and trace control needed for
integration by parts in Eq.~\eqref{eq:ch2.maxwell-adjoint-pairing}
\cite{stratton1941}. The composite operator $\mathcal{L}_\omega$ couples volume
residuals with boundary rows; its domain $\mathcal{D}(\mathcal{L}_\omega)$ is the
closed Maxwell solution class on which both $\mathcal{M}_\omega$ and
$\mathcal{B}_\Gamma$ are simultaneously evaluable in the declared data spaces
\cite{harrington2001}. The closed Maxwell solution domain becomes
\begin{equation}
\label{eq:ch2.domain-refinement-L}
\mathcal{D}(\mathcal{L}_\omega)
=
\left\{
\mathbf{u}_\omega\in\mathcal{F}_\Omega:
\mathcal{M}_\omega[\mathbf{u}_\omega]\in L^{2}(\Omega)^{4}
\right\},
\end{equation}
which specializes Eq.~\eqref{eq:ch2.operator-domain-declaration} to Lipschitz
domains with patchwise smooth coefficients \cite{stratton1941}. On
\(\mathcal{D}(\mathcal{L}_\omega)\), the volume residuals of
Eq.~\eqref{eq:ch2.maxwell-driven-operator} are square integrable, and the boundary
row of Eq.~\eqref{eq:ch2.maxwell-boundary-composite} is evaluated in
\(\mathcal{T}_\Gamma\). Exterior domains append outgoing subspaces rather than
changing Eqs.~\eqref{eq:ch2.hcurl-space}--\eqref{eq:ch2.admissible-field-space};
that selection is function-space compatible with weighted norms that decay at
infinity and is detailed in Section~\ref{sec:ch243} \cite{harrington2001}.
Free-charge density paired with \(\Jimp\) in
Eq.~\eqref{eq:ch1.data-specification-tuple} is taken in the scalar class
\(\rho\in L^{2}(\Omega)\) whenever volumetric charge is prescribed, with
distributional interpretation \(\divg\Jimp+\jj\omega\rho=0\) understood in
\(\mathcal{D}'(\Omega)\) as in Eq.~\eqref{eq:ch1.source-admissibility}. Scalar
\(L^{2}\) control is sufficient because Gauss laws in
Eq.~\eqref{eq:ch1.maxwell-operator-block} involve \(\divg\) on \(\varepsilon_0\E\)
and \(\divg\) on \(\mu_0\H\), not gradients of \(\rho\). When \(\rho\) is absent and
\(\Jimp\in\mathcal{J}_\Omega^{(0)}\), the electric Gauss row reduces to a
divergence constraint on \(\E\) enforced weakly through test functions in
\(\HCurl(\Omega)\) \cite{stratton1941}. Approximation properties matter for numerical projection. Smooth compactly supported
fields are dense in \(\HCurl(\Omega)\) and in \(\mathcal{U}_\omega\) equipped with
norm~\eqref{eq:ch2.field-energy-norm}; every admissible state is a limit of
\(C_c^\infty\) test pairs modified to satisfy boundary class on \(\Gamma\)
\cite{stratton1941}. Density underpins Galerkin convergence: edge-element spaces
approximate \(\mathcal{F}_\Omega\) when their degrees of freedom match tangential
trace continuity on \(\Gamma_{\mathrm{PEC}}\), where
\begin{equation}
\label{eq:ch2.pec-subspace}
\mathcal{U}_\omega^{\mathrm{PEC}}
=
\left\{
(\E,\H)\in\mathcal{U}_\omega:
\hat{\mathbf{n}}\times\E|_{\Gamma_{\mathrm{PEC}}}=\mathbf{0}
\right\},
\end{equation}
encodes perfect-conductor tangential data from
Eq.~\eqref{eq:ch1.outer-boundary-classes} without assuming \(\E=\mathbf{0}\) in
\(\Omega\) \cite{harrington2001}. Integration by parts on
\(\mathcal{U}_\omega^{\mathrm{PEC}}\) yields the boundary-term cancellation used in
Eq.~\eqref{eq:ch2.maxwell-adjoint-pairing} when \(\mathcal{B}_\Gamma\) selects
homogeneous tangential electric trace.
Compact embeddings used later for Fredholm arguments are recorded as
\begin{equation}
\label{eq:ch2.embedding-chain}
\HCurl(\Omega)\hookrightarrow L^{2}(\Omega)^{3}
\hookleftarrow H^{-1/2}(\Gamma)^{3},
\end{equation}
with the first inclusion compact on bounded \(\Omega\) and the trace map continuous
but not compact \cite{stratton1941}. Equation~\eqref{eq:ch2.embedding-chain}
explains why interior cavity spectra are discrete while boundary observation maps are
rank deficient: compact embeddings produce accumulating eigenvalues in volume,
whereas non-compact traces map infinite-dimensional field spaces to finite
observation vectors \cite{harrington2001}. The quantitative singular-value decay of
radiation maps on compact support is developed in Section~\ref{sec:ch26} and cited
thereafter \cite{stratton1941}. Subsection~\ref{sec:ch213} develops the
consequences for radiation accessibility and finite observation.
Subsection~\ref{sec:ch211} has refined \(\mathcal{F}_\Omega\), \(\mathcal{J}_\Omega\), and
\(\mathcal{D}(\mathcal{L}_\omega)\) from minimal declarations to Sobolev, trace, and
distributional classes compatible with \ref{ch:01} charge-continuous. The consequence
for measurement is immediate: near-field scanners and method-of-moments codes inherit
these spaces as the continuous targets their discretizations must approximate.
Subsection~\ref{sec:ch212}
treats the Maxwell maps as linear operators with explicit domains, ranges, and adjoints on
those spaces \cite{stratton1941}.
As a worked example, suppose \(\mathbf{g}_\Gamma\) prescribes tangential \(\E\) on \(\Gamma\) with
\(H^{1/2}\) regularity but is not the trace of any \(\HCurl\) field in \(\Omega\).
Then no \(\mathbf{u}_\omega\in\mathcal{F}_\Omega\) satisfies
\(\gamma_{\tau,E}[\E]=\mathbf{g}_\Gamma\), regardless of how accurately a volume
solver minimizes Maxwell residuals. The failure is a trace-space mismatch, not a
Green-function error; correcting it requires either lifting \(\mathbf{g}_\Gamma\) into
\(\mathcal{T}_\Gamma\) as an image under \(\gamma_\tau\) or relaxing the boundary
class through impedance completion in Subsection~\ref{sec:ch214}
\cite{stratton1941}.
\subsection{Linear operators: domain, range, and adjoint}
\label{sec:ch212}
The physical problem is operator accountability: a network analyzer, finite-element solver, or
near-field scanner implicitly assumes that each map in the source-to-field pathway accepts
only admissible excitations and returns only measurable field pairs \cite{stratton1941}.
Every linear map---$\mathcal{S}_\omega$, $\mathcal{L}_\omega$, $\mathcal{R}_\omega$, and
finite-rank measurement surrogates---requires explicit domain, range, and adjoint pairings
\cite{stratton1941}. Recall that $\mathcal{S}_\omega:\mathcal{J}_\Omega\to\mathcal{F}_\Omega$
is the harmonic solution map, $\mathcal{R}_\omega=\Pi_{\mathrm{rad}}\circ\mathcal{S}_\omega$
selects outgoing radiative content, and $\mathcal{L}_\omega$ packages volume and boundary
residuals jointly \cite{harrington2001}. Reciprocity is a symmetry principle of accessibility:
network analyzers swap transmit and receive ports, and the adjoint is the mathematical shadow
of that hardware symmetry when sources remain passive \cite{harrington2001}. Treating ``adjoint''
as jargon detached from range calibration misses this link \cite{devaney2004}. When a synthesis
residual persists under port reversal, the fault often lies in domain mismatch between trial
currents and measured tangential data \cite{hansen1988}.
Continuity and divergence constraints tie impressed currents to charge conservation on real feeds \cite{maxwell1873}.
When those constraints are violated, numerical solvers may still return finite vectors,
but the corresponding fields exhibit unphysical reactive blow-up that no anechoic
termination can absorb.
Domain, range, and adjoint language records which excitations the operator accepts and
which field pairs are returned \cite{stratton1941}.
Subsection~\ref{sec:ch211} fixed the Hilbert and trace spaces on which Maxwell maps
act. Subsection~\ref{sec:ch212} treats those maps as linear operators with explicit
domains, ranges, and adjoints \cite{stratton1941}. The volume
operator \(\mathcal{M}_\omega\), boundary operator \(\mathcal{B}_\Gamma\), composite
\(\mathcal{L}_\omega\), and solution map \(\mathcal{S}_\omega\) introduced in
Subsections~\ref{sec:ch201}--\ref{sec:ch202} are not re-defined; their
Maxwell field structure is recorded here as the definitive construction of
domain--range--adjoint theory. Weak and energy forms of
Eqs.~\eqref{eq:ch2.maxwell-weak-form} and
\eqref{eq:ch2.maxwell-adjoint-pairing} supply the duality pairings below and are
cited, \cite{stratton1941}.
Let \(\mathcal{X}\) and \(\mathcal{Y}\) be Hilbert spaces. The linear map representing
\(T:\mathcal{D}(T)\to\mathcal{Y}\) is specified by its domain
\(\mathcal{D}(T)\subseteq\mathcal{X}\) and range
\(\mathcal{R}(T)=T\mathcal{D}(T)\subseteq\mathcal{Y}\). For Maxwell operators,
\(\mathcal{X}=\mathcal{U}_\omega\) from Eq.~\eqref{eq:ch2.maxwell-product-space}
and \(\mathcal{Y}\) is a product of volume and boundary data spaces. When a finite-element
solver reports residual norms, it is effectively measuring a graph norm: how large the
field is in \(\mathcal{U}_\omega\) together with how large Maxwell's residuals become when
that field is inserted into \(\mathcal{L}_\omega\). Engineers should read the graph norm
as a joint flux measurement record of field energy and equation balance, not as a single \(L^{2}\)
amplitude \cite{harrington2001}. A trial vector that is small in \(L^{2}\) but drives
large curl residuals cannot lie in \(\mathcal{D}(\mathcal{L}_\omega)\) on the physical
branch; conversely, a state with controlled graph norm is the continuous target that
Galerkin discretizations on edge elements approximate \cite{stratton1941}. The graph
norm
% PHYSICAL-MEMORY:sec:ch212
\begin{physicalvignette}[Rectangular waveguide aperture radiating into free space]
Consider Rectangular waveguide aperture radiating into free space. Before any coefficient table is printed, technicians record frequency, cable loss, and the mechanical outline that fixes electrical size. The first walk-through of the range is always tactile: flange bolts, probe standoff, turntable boresight, and the wattmeter reference on the receive chain \cite{hansen1988}.
On this hardware the near zone is visually impressive---standing reactive structure, rapid phase rotation, sometimes a visible VSWR ripple at the feed---yet the export question is measured only after outgoing selection on a large sphere \cite{harrington2001}. What the camera sees in the chamber is not automatically what the far-zone pattern claims \cite{stratton1941}.
Engineers therefore learn to separate three readings: what the feed can synthesize on declared support, what propagating content survives on the radiative range, and what a finite instrument can recover at declared tolerance \cite{devaney2004}. Confusing those readings produces the rank gaps and watt failures that later sections quantify \cite{pal2026}.
Continuity and divergence constraints tie impressed currents to charge conservation on real feeds. Violations appear as unphysical reactive blow-up in a mismatched numerical mesh model.
Rayleigh and Stratton treated such ranges as boundary-value laboratories long before digital near-field transforms became routine \cite{stratton1941}. Harrington's method-of-moments culture reinforced the same discipline: metal outline first, harmonic headline second \cite{harrington2001}.
Subsection~\ref{sec:ch212} states the formal export condition that this range scene presupposes \cite{stratton1941}.
\end{physicalvignette}
\begin{equation}
\label{eq:ch2.graph-norm}
\|\mathbf{u}\|_{\mathrm{gr}(T)}^{2}
=
\|\mathbf{u}\|_{\mathcal{X}}^{2}+\|T\mathbf{u}\|_{\mathcal{Y}}^{2},
\qquad
\mathbf{u}\in\mathcal{D}(T),
\end{equation}
turns \(\mathcal{D}(T)\) into a Hilbert space when \(T\) is closed
\cite{stratton1941}. Equation~\eqref{eq:ch2.graph-norm} is the diagnostic used in
Subsection~\ref{sec:ch214} to construct self-adjoint completions; here it organizes
domain control for \(\mathcal{L}_\omega\). 
Each Maxwell map acts from a declared domain of admissible states into a data space of residuals or traces; that domain--range discipline applies to each Maxwell map once
\(\mathcal{X}\) and \(\mathcal{Y}\) are chosen. Confusing \(\mathcal{D}(T)\) with
all of \(\mathcal{X}\) is a common source of spurious boundary terms in flux
identities \cite{stratton1941}.
The volume Maxwell operator $\mathcal{M}_\omega$ packages the four local Maxwell balance
rows that a harmonic solver tests row by row; it returns residuals in
$\mathcal{Y}_{\mathrm{vol}}$, not the field itself, and therefore acts only on states whose
curls are controlled enough that those residuals lie in $L^{2}$ \cite{stratton1941}. When a
finite-element log reports ``Maxwell residual,'' the engineer is observing an approximate
norm of $\mathcal{M}_\omega[\mathbf{u}_\omega]$; convergence to zero on the physical
branch is the numerical signature of membership in $\mathcal{D}(\mathcal{L}_\omega)$
\cite{harrington2001}. The volume Maxwell operator of Eq.~\eqref{eq:ch2.maxwell-driven-operator} acts as
\begin{equation}
\label{eq:ch2.volume-operator-domain-range}
\mathcal{M}_\omega:\mathcal{D}(\mathcal{M}_\omega)\longrightarrow
\mathcal{Y}_{\mathrm{vol}},
\qquad
\mathcal{Y}_{\mathrm{vol}}=L^{2}(\Omega)^{3}\times L^{2}(\Omega)^{3}
\times L^{2}(\Omega)\times L^{2}(\Omega)^{3},
\end{equation}
with \(\mathcal{D}(\mathcal{M}_\omega)\subset\mathcal{F}_\Omega\) the subspace on
which all four residual rows of Eq.~\eqref{eq:ch1.maxwell-operator-block} lie in
\(\mathcal{Y}_{\mathrm{vol}}\) \cite{stratton1941}. Operand interpretation:
\(\mathcal{M}_\omega\) evaluates curl-coupled Maxwell residuals pointwise in the
\(L^{2}\) sense; physically, it returns whether a candidate \((\E,\H)\) satisfies
local Maxwell balance with impressed data removed to the forcing vector
Eq.~\eqref{eq:ch2.maxwell-forcing-vector}. The map is linear in
\(\mathbf{u}_\omega\) and bounded on bounded subsets of
\(\mathcal{D}(\mathcal{M}_\omega)\) but is not bounded on all of
\(\mathcal{U}_\omega\) because curl norms are not controlled by \(L^{2}\) norms alone
\cite{harrington2001}.
Tangential traces form bounded boundary operators on \(\HCurl(\Omega)\):
\begin{equation}
\label{eq:ch2.trace-operator-bounded}
\gamma_\tau:\HCurl(\Omega)\longrightarrow H^{-1/2}(\Gamma)^{3},
\qquad
\gamma_\tau[\mathbf{u}]=\hat{\mathbf{n}}\times\mathbf{u}|_{\Gamma},
\end{equation}
with \(\|\gamma_\tau\|_{\mathrm{op}}<\infty\) on Lipschitz \(\Gamma\)
\cite{stratton1941}. Equation~\eqref{eq:ch2.trace-operator-bounded} specializes
Eqs.~\eqref{eq:ch2.tangential-trace-maps}--\eqref{eq:ch2.tangential-trace-magnetic}
to a single template. The boundary operator of
Eq.~\eqref{eq:ch1.boundary-data-operator} is
\begin{equation}
\label{eq:ch2.boundary-operator-domain-range}
\mathcal{B}_\Gamma:\mathcal{D}(\mathcal{B}_\Gamma)\longrightarrow\mathcal{Y}_\Gamma,
\qquad
\mathcal{D}(\mathcal{B}_\Gamma)\subseteq\mathcal{T}_\Gamma,
\end{equation}
where \(\mathcal{T}_\Gamma\) and \(\mathcal{Y}_\Gamma\) are from
Subsection~\ref{sec:ch211}. Impedance and outgoing classes from
Eqs.~\eqref{eq:ch1.outer-boundary-classes}--\eqref{eq:ch1.radiation-trace-condition}
restrict \(\mathcal{D}(\mathcal{B}_\Gamma)\) to admissible trace pairs compatible
with Eq.~\eqref{eq:ch1.admissible-data-class} \cite{stratton1941}.
The composite operator of Eq.~\eqref{eq:ch2.maxwell-boundary-composite} is
\begin{equation}
\label{eq:ch2.composite-operator-domain}
\mathcal{L}_\omega:\mathcal{D}(\mathcal{L}_\omega)\longrightarrow
\mathcal{Y}_{\mathrm{vol}}\times\mathcal{Y}_\Gamma,
\end{equation}
with \(\mathcal{D}(\mathcal{L}_\omega)\) given by
Eq.~\eqref{eq:ch2.domain-refinement-L}. Range membership is \emph{coupled}: a state
cannot satisfy volume rows while violating boundary rows in the product sense of
Eq.~\eqref{eq:ch2.closed-maxwell-bvp}. The graph norm
\eqref{eq:ch2.graph-norm} on \(\mathcal{L}_\omega\) controls both residuals
simultaneously and is the natural norm for convergence tests in numerical Maxwell
solvers \cite{harrington2001}.
The harmonic solution map is an operator on source space:
\begin{equation}
\label{eq:ch2.solution-operator-domain}
\mathcal{S}_\omega:\mathcal{D}(\mathcal{S}_\omega)\longrightarrow\mathcal{F}_\Omega,
\qquad
\mathcal{D}(\mathcal{S}_\omega)=\mathcal{J}_\Omega,
\end{equation}
defined by Eq.~\eqref{eq:ch2.solution-map-resolvent} as
\(\mathcal{L}_\omega[\mathcal{S}_\omega[\Jimp]]=[\mathbf{s}_\omega(\Jimp),\mathbf{g}_\Gamma]\).
Well-posedness requires invertibility of \(\mathcal{L}_\omega\) on the physical
branch selected by \(\mathcal{B}_\Gamma\); then
\begin{equation}
\label{eq:ch2.resolvent-operator}
\mathcal{S}_\omega=\mathcal{L}_\omega^{-1}\circ
\big(\mathbf{s}_\omega(\cdot),\mathbf{g}_\Gamma\big),
\end{equation}
with \(\mathcal{L}_\omega^{-1}\) understood as the outgoing or cavity-truncated
resolvent of Sections~\ref{sec:ch243} and~\ref{sec:ch263}, not as an algebraic
inverse on all of \(\mathcal{U}_\omega\) \cite{stratton1941}. Radiation and
observation operators compose accordingly:
\begin{equation}
\label{eq:ch2.radiation-operator-composition}
\mathcal{R}_\omega=\Pi_{\mathrm{rad}}\circ\mathcal{S}_\omega,
\qquad
\Lambda_{\mathrm{rad}}=\mathcal{O}_{\mathrm{far}}\circ\mathcal{R}_\omega,
\end{equation}
with \(\Pi_{\mathrm{rad}}\) from Eq.~\eqref{eq:ch2.radiation-operator-def} and
\(\mathcal{O}_{\mathrm{far}}\) from Subsection~\ref{sec:ch201}
\cite{harrington2001}.
\begin{figure}[t]
\centering
\ChOneFigBlock{%
\begin{tikzpicture}[ch1fig, node distance=5mm and 7mm]
 \node[ch1box] (j) {$\mathcal{J}_\Omega$};
 \node[ch1box, right=7mm of j] (s) {$\mathcal{S}_\omega$};
 \node[ch1box, right=7mm of s] (f) {$\mathcal{F}_\Omega$};
 \node[ch1box, right=7mm of f] (r) {$\mathcal{R}_\omega$};
 \draw[ch1flow] (j) -- (s) -- (f) -- node[above,font=\scriptsize] {$\Pi_{\mathrm{rad}}$} (r);
\end{tikzpicture}%
}
\caption{Excitation-to-flux sequence on declared domains:
\(\mathcal{S}_\omega:\mathcal{J}_\Omega\to\mathcal{F}_\Omega\),
\(\mathcal{R}_\omega:\mathcal{J}_\Omega\to\mathcal{F}_{\mathrm{rad}}\).}
\label{fig:ch2.solution-radiation-operator-composition}
\end{figure}
~\eqref{eq:ch2.measurement chain-composition} in domain--range language. Each arrow is meaningful only on the declared domain; extending \(\mathcal{S}_\omega\) to
non-admissible sources produces elements of \(\mathcal{U}_\omega\) that need not
lie in \(\mathcal{F}_\Omega\) \cite{stratton1941}. Adjoints are defined through duality pairings compatible with
Eq.~\eqref{eq:ch2.dual-source-pairing}. For the volume operator, the formal adjoint
\( \mathcal{M}_\omega^{\dagger}\) satisfies
\begin{equation}
\label{eq:ch2.volume-adjoint-definition}
\langle\mathcal{M}_\omega[\mathbf{u}_\omega],\mathbf{v}_\omega\rangle_{\mathrm{vol}}
=
\langle\mathbf{u}_\omega,\mathcal{M}_\omega^{\dagger}[\mathbf{v}_\omega]\rangle_{\mathcal{U}}
+\mathcal{B}_{\mathrm{surf}}(\mathbf{u}_\omega,\mathbf{v}_\omega),
\end{equation}
where \(\langle\cdot,\cdot\rangle_{\mathrm{vol}}\) is the componentwise
\(L^{2}\) pairing on \(\mathcal{Y}_{\mathrm{vol}}\),
\(\langle\cdot,\cdot\rangle_{\mathcal{U}}\) is induced by
Eq.~\eqref{eq:ch2.field-energy-norm}, and
\(\mathcal{B}_{\mathrm{surf}}\) collects boundary terms from integration by parts
\cite{stratton1941}. On \(\mathcal{D}(\mathcal{M}_\omega)\cap\mathcal{U}_\omega^{\mathrm{PEC}}\)
from Eq.~\eqref{eq:ch2.pec-subspace}, Eq.~\eqref{eq:ch2.maxwell-adjoint-pairing}
identifies \(\mathcal{B}_{\mathrm{surf}}\) with the tangential trace pairing of
Eq.~\eqref{eq:ch2.maxwell-adjoint-pairing}; that identity is imported, not
re-proved \cite{stratton1941}. The composite adjoint satisfies the block structure
\begin{equation}
\label{eq:ch2.composite-adjoint-block}
\mathcal{L}_\omega^{\dagger}
=
\begin{bmatrix}
\mathcal{M}_\omega^{\dagger} & 0\\
0 & \mathcal{B}_\Gamma^{\dagger}
\end{bmatrix}
\quad\text{on compatible pairs,}
\end{equation}
with cross-coupling entering only through shared trace data when boundary and
volume domains are not synchronized. Self-adjoint \emph{realizations} of truncated
Maxwell problems require choosing \(\mathcal{B}_\Gamma\) so that
\(\mathcal{B}_{\mathrm{surf}}\) cancels on \(\mathcal{D}(\mathcal{L}_\omega)\);
Forcing functionals induced by impressed data are bounded linear operators on test
spaces. The map
\begin{equation}
\label{eq:ch2.forcing-functional-operator}
\mathfrak{f}_\omega:\mathcal{U}_\omega^{0}\longrightarrow\mathbb{C},
\qquad
\mathfrak{f}_\omega[\mathbf{v}_\omega]=
\langle\mathbf{s}_\omega(\mathfrak{D}),\mathbf{v}_\omega\rangle_{\mathrm{vol}},
\end{equation}
is continuous when \(\mathbf{s}_\omega\in\mathcal{Y}_{\mathrm{vol}}\) and
\(\mathbf{v}_\omega\in\mathcal{U}_\omega^{0}\) carries homogeneous boundary traces
compatible with \(\mathcal{B}_\Gamma\) \cite{stratton1941}. Equation
\eqref{eq:ch2.forcing-functional-operator} is the operator reading of the right-hand
side of Eq.~\eqref{eq:ch2.maxwell-weak-form}. Source-to-field operators inherit
continuity on bounded subsets: if \(\|\Jimp\|_{\mathcal{J}}\le C\) and
\(\mathcal{L}_\omega\) is boundedly invertible on a fixed branch, then
\begin{equation}
\label{eq:ch2.solution-map-stability}
\|\mathcal{S}_\omega[\Jimp]\|_{\mathcal{U}}
\le
C_{\mathrm{res}}\,\|\mathbf{s}_\omega(\Jimp)\|_{\mathrm{vol}},
\end{equation}
with \(C_{\mathrm{res}}=\|\mathcal{L}_\omega^{-1}\|\) on the selected subspace
\cite{harrington2001}. Inequality~\eqref{eq:ch2.solution-map-stability}
is the operator form of solver stability under source perturbations.
Kernel and range spaces are recorded for later null-space analysis:
\begin{equation}
\label{eq:ch2.kernel-range-L}
\ker\mathcal{L}_\omega
=
\left\{
\mathbf{u}_\omega\in\mathcal{D}(\mathcal{L}_\omega):
\mathcal{L}_\omega[\mathbf{u}_\omega]=\mathbf{0}
\right\},
\end{equation}
\begin{equation}
\label{eq:ch2.range-L}
\mathcal{R}(\mathcal{L}_\omega)
=
\left\{
[\mathbf{s},\mathbf{g}]^{\mathsf{T}}\in
\mathcal{Y}_{\mathrm{vol}}\times\mathcal{Y}_\Gamma:
\exists\,\mathbf{u}_\omega\in\mathcal{D}(\mathcal{L}_\omega):\
\mathcal{L}_\omega[\mathbf{u}_\omega]=[\mathbf{s},\mathbf{g}]^{\mathsf{T}}
\right\}.
\end{equation}
On bounded domains with passive \(\mathcal{B}_\Gamma\), \(\ker\mathcal{L}_\omega\)
contains cavity modes; on exterior domains with outgoing radiation condition,
\(\ker\mathcal{L}_\omega\) is trivial on the outgoing branch
\cite{stratton1941}. Section~\ref{sec:ch272} analyzes
\(\ker\mathcal{R}_\omega\); Subsection~\ref{sec:ch213} treats compactness and
Fredholm index on \(\mathcal{R}(\mathcal{L}_\omega)\).
Extension operators lift boundary data into volume fields. A right inverse
\(E_\Gamma:\mathcal{Y}_\Gamma\to\HCurl(\Omega)\) satisfying
\(\gamma_\tau\circ E_\Gamma=\mathrm{id}\) on compatible subclasses exists only for
\(\mathbf{g}_\Gamma\) in the range of \(\gamma_\tau\); otherwise no
\(\mathbf{u}_\omega\in\mathcal{F}_\Omega\) meets the boundary row
\cite{stratton1941}. Extension operators enter Calder\'on preconditioners and
boundary-element formulations developed in later sections; here they clarify that
\(\mathcal{B}_\Gamma\) acts on traces, not on arbitrary \(L^{2}\) boundary vectors
\cite{harrington2001}.
As a worked example, fix \(\Jimp\in\mathcal{J}_\Omega\) and two distinct outgoing radiation conditions
\(\mathcal{B}_\Gamma^{(a)}\) and \(\mathcal{B}_\Gamma^{(b)}\) on the same exterior
domain. The operators \(\mathcal{L}_\omega^{(a)}\) and
\(\mathcal{L}_\omega^{(b)}\) have the same volume domain but different boundary
domains; Eq.~\eqref{eq:ch2.resolvent-operator} yields two resolvents
\(\mathcal{S}_\omega^{(a)}\neq\mathcal{S}_\omega^{(b)}\) in general. Comparing
analysis to measurement without matching \(\mathcal{D}(\mathcal{B}_\Gamma)\) is
therefore a comparison of different operators, not of different discretizations of
one operator \cite{stratton1941}.
Subsection~\ref{sec:ch212} has packaged domain, range, graph norm, adjoint pairings, and
the composition $\mathcal{R}_\omega=\Pi_{\mathrm{rad}}\circ\mathcal{S}_\omega$ on declared
spaces. Subsection~\ref{sec:ch213} records boundedness, compact embeddings, and the rank
deficiency that finite instruments impose on $\mathcal{F}_{\mathrm{rad}}$ \cite{stratton1941}.
\subsection{Boundedness, compactness, rank deficiency, spectral compression}
\label{sec:ch213}
% HARRINGTON-OPEN:sec:ch213
Spectral compression on compact support limits how sharply a finite aperture can beam \cite{jackson1999}. Sidelobe floors on a range trace the same law \cite{stratton1941}. Maxwell's retarded export to $\Gamma_R$ remains the governing constraint \cite{maxwell1873}.
The physical problem is observability under finite geometry: a compact source, a bounded
domain, and a finite aperture cannot couple to infinitely many independent field degrees of
freedom \cite{hansen1988}. Subsections~\ref{sec:ch211}--\ref{sec:ch212} fixed spaces and operator domains.
Subsection~\ref{sec:ch213} records boundedness and compactness properties, rank
deficiency of observation maps, and the viewpoint from which Section~\ref{sec:ch26}
develops spectral compression for radiation operators \cite{stratton1941}. The physical
content is familiar on a compact range: finite sources and finite probes cannot couple to
infinitely many independent far-zone degrees of freedom, so accessible rank is limited by
geometry and observation extent \cite{hansen1988}. Fredholm index computation
and complete spectral decomposition remain in Sections~\ref{sec:ch263}
and~\ref{sec:ch261}; here the constructive development establishes the lemmas on
boundedness, compactness, and rank deficiency those sections invoke by reference
\cite{stratton1941}.
In the display below, $\mathcal{F}_{\mathrm{rad}}\subset\mathcal{F}_\Omega$ denotes the
outgoing radiative subspace selected by $\Pi_{\mathrm{rad}}$; $\mathcal{S}_\omega$ and
$\mathcal{R}_\omega$ are the solution and radiation maps on $\mathcal{J}_\Omega$
\cite{harrington2001}.
The linear map representing \(T:\mathcal{X}\to\mathcal{Y}\) between Hilbert spaces is
\emph{bounded} if
% PHYSICAL-MEMORY:sec:ch213
\begin{physicalvignette}[Circular patch on a handset ground plane]
Consider Circular patch on a handset ground plane. Before any coefficient table is printed, technicians record frequency, cable loss, and the mechanical outline that fixes electrical size. The first walk-through of the range is always tactile: flange bolts, probe standoff, turntable boresight, and the wattmeter reference on the receive chain \cite{hansen1988}.
On this hardware the near zone is visually impressive---standing reactive structure, rapid phase rotation, sometimes a visible VSWR ripple at the feed---yet the export question is measured only after outgoing selection on a large sphere \cite{harrington2001}. What the camera sees in the chamber is not automatically what the far-zone pattern claims \cite{stratton1941}.
Engineers therefore learn to separate three readings: what the feed can synthesize on declared support, what propagating content survives on the radiative range, and what a finite instrument can recover at declared tolerance \cite{devaney2004}. Confusing those readings produces the rank gaps and watt failures that later sections quantify \cite{pal2026}.
Spectral compression on compact support limits how sharply a finite aperture can beam . Sidelobe floors on a range trace the same law .
Rayleigh and Stratton treated such ranges as boundary-value laboratories long before digital near-field transforms became routine \cite{stratton1941}. Harrington's method-of-moments culture reinforced the same discipline: metal outline first, harmonic headline second \cite{harrington2001}.
Subsection~\ref{sec:ch213} states the formal export condition that this range scene presupposes \cite{stratton1941}.
\end{physicalvignette}
\begin{equation}
\label{eq:ch2.operator-boundedness}
\|T\|_{\mathrm{op}}
=
\sup_{\|\mathbf{x}\|_{\mathcal{X}}=1}\|T\mathbf{x}\|_{\mathcal{Y}}
<\infty.
\end{equation}
Boundedness is equivalent to continuity at the origin and is the default
hypothesis for composition with measurement readouts \cite{stratton1941}. On the
Maxwell chain, operators split into bounded and unbounded classes. Trace and observation maps are bounded on declared spaces. The tangential trace
Eq.~\eqref{eq:ch2.trace-operator-bounded} satisfies
\(\|\gamma_\tau\|_{\mathrm{op}}<\infty\) from \(\HCurl(\Omega)\) to
\(H^{-1/2}(\Gamma)^{3}\) \cite{stratton1941}. Far-zone readouts from
Subsection~\ref{sec:ch201} define bounded functionals when restricted to
\(\mathcal{F}_{\mathrm{rad}}\):
\begin{equation}
\label{eq:ch2.far-observation-bounded}
\|\mathcal{O}_{\mathrm{far}}\|_{\mathrm{op}}
=
\sup_{\|\mathbf{u}\|_{\mathcal{U}}=1}
\big|\mathcal{O}_{\mathrm{far}}[\mathbf{u}]\big|
<\infty
\quad\text{on}\ \mathcal{F}_{\mathrm{rad}}.
\end{equation}
Laboratory maps \(\mathcal{O}_{\mathrm{meas}}\circ\mathcal{O}_{\mathrm{far}}\) in
Eq.~\eqref{eq:ch2.operational-measurement-chain} are bounded on the same subspace
because composition of bounded operators is bounded \cite{harrington2001}.
Volume and composite Maxwell operators are unbounded on full
\(\mathcal{U}_\omega\). The map \(\mathcal{M}_\omega\) of
Eq.~\eqref{eq:ch2.volume-operator-domain-range} has
\(\|\mathcal{M}_\omega\|_{\mathrm{op}}=\infty\) when
\(\mathcal{D}(\mathcal{M}_\omega)=\mathcal{U}_\omega\) because controlling
\(\|\mathcal{M}_\omega[\mathbf{u}]\|_{\mathrm{vol}}\) would require uniform curl
bounds beyond \(L^{2}\) control \cite{stratton1941}. Restriction to
\(\mathcal{D}(\mathcal{L}_\omega)\) with graph norm
Eq.~\eqref{eq:ch2.graph-norm} restores boundedness of the \emph{resolvent}
\(\mathcal{L}_\omega^{-1}\) on the physical branch where invertibility holds:
\begin{equation}
\label{eq:ch2.resolvent-bounded}
\|\mathcal{L}_\omega^{-1}\|_{\mathrm{op}}
<\infty
\quad\Longleftrightarrow\quad
\mathcal{L}_\omega\ \text{bijective on}\ \mathcal{D}(\mathcal{L}_\omega)/\ker\mathcal{L}_\omega.
\end{equation}
Equation~\eqref{eq:ch2.resolvent-bounded} specializes
Eq.~\eqref{eq:ch2.solution-map-stability}; exterior outgoing problems use
resolvent bounds on weighted subspaces rather than on all of
\(\mathcal{U}_\omega\) \cite{harrington2001}.
\begin{figure}[t]
\centering
\ChOneFigBlock{%
\begin{tikzpicture}[ch1fig, node distance=5mm and 7mm]
 \node[ch1box] (hc) {$\HCurl(\Omega)$};
 \node[ch1box, right=16mm of hc] (l2) {$L^{2}(\Omega)^{3}$};
 \node[ch1box, below=5mm of l2] (tr) {$H^{-1/2}(\Gamma)^{3}$};
 \draw[ch1flow] (hc.east) -- node[above,font=\scriptsize]{compact} (l2.west);
 \draw[ch1flow] (hc.south) -- node[right,font=\scriptsize]{$\gamma_\tau$} (tr.north);
\end{tikzpicture}%
}
\caption{Compact embedding \(\HCurl(\Omega)\hookrightarrow L^{2}(\Omega)^{3}\) on
bounded \(\Omega\); trace map is bounded but not compact
(Eq.~\eqref{eq:ch2.embedding-chain}).}
\label{fig:ch2.compact-embedding-schematic}
\end{figure}
~\eqref{eq:ch2.embedding-chain}. Rellich-type compactness for
\(\HCurl(\Omega)\hookrightarrow L^{2}(\Omega)^{3}\) on bounded domains implies
that bounded sequences in \(\HCurl\) have \(L^{2}\)-convergent subsequences
\cite{stratton1941}. Finite geometry cannot export arbitrarily fine angular complexity;
propagation therefore acts as a compression mechanism on electromagnetic structure.
An operator \(K:\mathcal{X}\to\mathcal{Y}\) is
\emph{compact} if bounded sets in \(\mathcal{X}\) map to relatively compact sets
in \(\mathcal{Y}\):
\begin{equation}
\label{eq:ch2.compact-operator-def}
K\ \text{compact}
\iff
K(B_{\mathcal{X}}(0,1))\ \text{has compact completion in}\ \mathcal{Y}.
\end{equation} 
Compact perturbations of identity produce Fredholm structure used in
Section~\ref{sec:ch263}; Subsection~\ref{sec:ch213} records the embedding lemmas
without computing index \cite{harrington2001}.
Radiation and Green maps on fixed bounded domains often decompose as
\begin{equation}
\label{eq:ch2.green-compact-split}
\mathcal{S}_\omega
=
G_{\mathrm{sing}}+K_{\mathrm{cpt}},
\end{equation}
with \(G_{\mathrm{sing}}\) capturing singular near-field kernels and
\(K_{\mathrm{cpt}}\) compact on \(\mathcal{J}_\Omega\) when sources have compact
support and \(\Omega\) is bounded \cite{stratton1941}. Equation
\eqref{eq:ch2.green-compact-split} is the operator template for
Section~\ref{sec:ch23}; kernel construction is not repeated here. The split
separates non-compact singular behavior at coincident points from compact
smoothing at finite separation \cite{harrington2001}. Rank deficiency enters through finite observation. An aperture-limited measurement
map \(\mathcal{O}_{\mathrm{meas}}:\mathcal{F}_{\mathrm{rad}}\to\mathbb{C}^{N}\)
has finite-dimensional range with \(N<\infty\):
\begin{equation}
\label{eq:ch2.finite-rank-observation}
\mathcal{O}_{\mathrm{meas}}
=
\sum_{n=1}^{N}\langle\cdot,\mathbf{f}_n\rangle_{\mathcal{U}}\,\mathbf{e}_n,
\end{equation}
for dual vectors \(\mathbf{f}_n\) and coordinate directions \(\mathbf{e}_n\).
Therefore \(\mathrm{rank}(\mathcal{O}_{\mathrm{meas}})\le N\) and
\begin{equation}
\label{eq:ch2.observation-null-space}
\ker\mathcal{O}_{\mathrm{meas}}
=
\left\{
\mathbf{u}\in\mathcal{F}_{\mathrm{rad}}:
\mathcal{O}_{\mathrm{meas}}[\mathbf{u}]=\mathbf{0}
\right\}
\end{equation}
is infinite-dimensional whenever \(\mathcal{F}_{\mathrm{rad}}\) is
infinite-dimensional \cite{stratton1941}. Equation
\eqref{eq:ch2.observation-null-space} is the operator form of the measurement
instrument rank limits in Section~\ref{sec:ch01.5}: finite modal sets cannot resolve all
angular content. Rank deficiency is structural, not a discretization artifact alone; the
spectral compression and singular-value decay law for radiation maps on compact support are
owned by Section~\ref{sec:ch26} \cite{stratton1941}.
\begin{figure}[t]
\centering
\ChOneFigBlock{%
\begin{tikzpicture}[ch1fig, node distance=5mm and 7mm]
 \node[ch1box] (f) {$\mathcal{F}_{\mathrm{rad}}$};
 \node[ch1box, right=7mm of f] (o) {$\mathcal{O}_{\mathrm{meas}}$};
 \node[ch1box, right=7mm of o] (cn) {$\mathbb{C}^{N}$};
 \node[ch1box, below=5mm of f] (ker) {$\ker\mathcal{O}_{\mathrm{meas}}$};
 \draw[ch1flow] (f) -- (o) -- (cn);
 \draw[ch1flow, dashed] (ker.south) -- (f.north);
\end{tikzpicture}%
}
\caption{Finite-rank infinite-dimensional field space maps to
\(\mathbb{C}^{N}\); \(\ker\mathcal{O}_{\mathrm{meas}}\) is infinite-dimensional for
\(N<\infty\).}
\label{fig:ch2.rank-deficient-observation}
\end{figure}
~\eqref{eq:ch2.radiation-operator-composition}.
Angular-spectrum reports that omit unresolved harmonics lie in the range of
\(\mathcal{O}_{\mathrm{meas}}\) but not in its injective sector
\cite{stratton1941}. Subsection~\ref{sec:ch213} does not classify report types;
Section~\ref{sec:ch01.5} remains authoritative. Singular values quantify rank deficiency. For compact
\(K:\mathcal{X}\to\mathcal{Y}\), the singular value sequence
\(\{\sigma_n(K)\}_{n\ge 1}\) decreases to zero, and
\begin{equation}
\label{eq:ch2.singular-value-decay}
\|K-R_N K\|_{\mathrm{op}}
=
\inf_{\mathrm{rank}\,R_N\le N}\|K-R_N\|_{\mathrm{op}}
\sim \sigma_{N+1}(K),
\end{equation}
where \(R_N\) ranges over rank-\(N\) operators \cite{stratton1941}. For radiation
maps truncated to \(N\) spherical or plane-wave modes, Eq.~\eqref{eq:ch2.singular-value-decay}
controls how much field content lies outside the retained modal subspace; explicit
VSH eigenbases are deferred to Section~\ref{sec:ch261} and Chapter~3 (\emph{Angular Momentum Operators and SO(3)})
\cite{harrington2001}. \emph{Spectral compression} names the reduction of infinite-dimensional radiation
dynamics to a finite or countable spectral model on the physical branch. Let
\(\{\phi_n\}\) be a spectral family in \(\mathcal{F}_\Omega\) and
\(\{P_n\}\) the associated one-dimensional projectors. A truncated radiation operator
\begin{equation}
\label{eq:ch2.spectral-compression-operator}
\mathcal{R}_\omega^{(N)}
=
\sum_{n=1}^{N}\lambda_n\,P_n\circ\mathcal{S}_\omega
\end{equation}
retains only modes with eigenvalues \(\lambda_n\) above a declared threshold
\cite{stratton1941}. Equation
\eqref{eq:ch2.spectral-compression-operator} is the abstract template for
Section~\ref{sec:ch26}; concrete eigenfunctions are not constructed here. On
bounded cavities, compact embeddings yield accumulating eigenvalues
\(\lambda_n\to\infty\) for interior resonances; on exterior domains, continuous
spectra require outgoing projectors rather than discrete sums
\cite{stratton1941}.
Compression error is bounded in operator norm when \(K\) is compact:
\begin{equation}
\label{eq:ch2.compression-error-bound}
\|\mathcal{R}_\omega-\mathcal{R}_\omega^{(N)}\|_{\mathrm{op}}
\le
C\,\sigma_{N+1}(K_{\mathrm{cpt}}),
\end{equation}
with \(K_{\mathrm{cpt}}\) from Eq.~\eqref{eq:ch2.green-compact-split} and
\(C\) independent of impressed data on bounded \(\mathcal{J}_\Omega\)
\cite{harrington2001}. Physically, Eq.~\eqref{eq:ch2.compression-error-bound}
states that truncating modal content beyond the \((N+1)\)st singular value cannot
be recovered by any rank-\(N\) surrogate operator; the quantitative decay law is
as in Section~\ref{sec:ch26} \cite{stratton1941}.
Rank deficiency of \(\mathcal{R}_\omega\) itself is distinct from observation rank
deficiency. Non-radiating sources satisfy
\(\mathcal{R}_\omega[\Jimp]=\mathbf{0}\) while
\(\mathcal{S}_\omega[\Jimp]\neq\mathbf{0}\) (Eq.~\eqref{eq:ch2.null-radiation});
Section~\ref{sec:ch272} develops that null space. Here the point is compositional:
\begin{equation}
\label{eq:ch2.rank-composition-law}
\mathrm{rank}(\mathcal{O}_{\mathrm{meas}}\mathcal{R}_\omega)
\le
\min\big\{\mathrm{rank}(\mathcal{O}_{\mathrm{meas}}),
\dim\mathcal{R}(\mathcal{R}_\omega)\big\},
\end{equation}
In Eq.~\eqref{eq:ch2.rank-composition-law},
\(\mathcal{O}_{\mathrm{meas}}\) is a finite-rank linear measurement functional
(pattern probe bank, modal fitter, or spherical-harmonic extractor) composed with
\(\mathcal{R}_\omega\); \(\mathrm{rank}(\cdot)\) is the dimension of the accessible
report subspace; and \(\mathcal{R}(\mathcal{R}_\omega)\) is the range of the radiation map
on \(\mathcal{J}_\Omega\).
So measurement rank caps reported angular content even when
\(\mathcal{R}_\omega\) is injective on admissible sources
\cite{stratton1941}. Both deficits must appear in uncertainty
budgets for angular-transport claims imported from Section~\ref{sec:ch01.2.3}. Fredholm structure note: on bounded \(\Omega\) with passive
\(\mathcal{B}_\Gamma\), \(\mathcal{L}_\omega\) is Fredholm on
\(\mathcal{D}(\mathcal{L}_\omega)\) modulo cavity kernels when written as
\begin{equation}
\label{eq:ch2.fredholm-split}
\mathcal{L}_\omega
=
\mathcal{L}_0+K,
\qquad
K\ \text{compact on graph domain},
\end{equation}
with \(\mathcal{L}_0\) an invertible reference operator on the quotient by
\(\ker\mathcal{L}_\omega\) \cite{stratton1941}. Index and solvability criteria are
developed in Section~\ref{sec:ch263}; Eq.~\eqref{eq:ch2.fredholm-split} supplies the
split Subsection~\ref{sec:ch212} stated the introductory relation in
Subsection~\ref{sec:ch202}.
As a worked example, a spherical near-field scanner records \(N=50\) tangential field samples on a
closed surface. The map \(\mathcal{O}_{\mathrm{meas}}\) has rank at most \(50\), so
angular harmonics beyond the resolved set lie in
\(\ker\mathcal{O}_{\mathrm{meas}}\). Applying Eq.~\eqref{eq:ch2.singular-value-decay}
to the compact smoothing operator from sample interpolation shows that recovered
far-zone coefficients error is controlled by \(\sigma_{51}\), not by Maxwell
residual alone \cite{stratton1941}. Increasing frequency without
increasing \(N\) lowers \(\sigma_{N+1}\) slowly for fixed aperture, worsening
compression error even when \(\mathcal{L}_\omega[\mathbf{u}_\omega]=\mathbf{0}\)
in the volume.
As a worked example, in a closed cavity, compact embedding yields discrete resonances and
Eq.~\eqref{eq:ch2.spectral-compression-operator} sums over cavity modes. In an
exterior problem with outgoing radiation condition, the same formula requires a continuous
spectral integral rather than a finite sum; using cavity compression on an
open antenna model conflates discrete storage modes with continuous radiation
modes \cite{stratton1941}. Section~\ref{sec:ch243} selects the outgoing branch;
Subsection~\ref{sec:ch213} only records that compression templates differ by domain class
\cite{stratton1941}.
Finite-geometry interpretation (compactness and rank). Compact embeddings and finite-rank observation maps express one enduring electromagnetic
truth: finite geometry destroys information. A bounded source support $\Omega_s$ cannot
excite arbitrarily fine angular structure at every observation point; a finite aperture
integrates over unresolved directions; a finite modal set truncates accessible rank.
Inaccessible angular structure is therefore not a numerical inconvenience but a consequence
of Maxwell propagation from compact sources through outgoing selection to finite
instruments. flux report is fundamentally limited, and the mathematics of compactness
makes that limit quantifiable through decaying singular values and infinite-dimensional
observation null spaces \cite{stratton1941}.
Subsection~\ref{sec:ch213} has established boundedness classes, compact embeddings, and
finite-rank observation structure. The quantitative singular-value decay for radiation maps
on compact support is developed in Section~\ref{sec:ch26}; here only the structural lemmas
are recorded. Subsection~\ref{sec:ch214}
completes the function-space foundation with self-adjoint realizations and admissibility
completions on bounded domains \cite{stratton1941}.
\subsection{Self-adjointness, completions, and admissibility}
\label{sec:ch214}
The physical problem is passive realizability on bounded hardware: cavity walls, matched
terminations, and dissipative linings must not permit unbounded storage of reactive energy
under admissible excitation \cite{stratton1941}. Friedrichs and other
completions extend domains just enough that $\curl$ and $\hat{\mathbf{n}}\times$ are
simultaneously closed on the same space---precisely what finite-element antenna models
require when sheet sources sit on interior interfaces \cite{harrington2001}. Improper
domain extension predicts unphysical growth of reactive storage---a numerical symptom that
must be corrected before radiation maps are trusted \cite{maxwell1873}. Completion
choices are radiation-physics choices about which reactive storage remains inside $\Omega$
and which energy is exported \cite{sommerfeld1949}.
% PHYSICAL-MEMORY:sec:ch214
\begin{physicalvignette}[Helical feed on a satellite reflector]
Consider Helical feed on a satellite reflector. Before any coefficient table is printed, technicians record frequency, cable loss, and the mechanical outline that fixes electrical size. The first walk-through of the range is always tactile: flange bolts, probe standoff, turntable boresight, and the wattmeter reference on the receive chain \cite{hansen1988}.
On this hardware the near zone is visually impressive---standing reactive structure, rapid phase rotation, sometimes a visible VSWR ripple at the feed---yet the export question is measured only after outgoing selection on a large sphere \cite{harrington2001}. What the camera sees in the chamber is not automatically what the far-zone pattern claims \cite{stratton1941}.
Engineers therefore learn to separate three readings: what the feed can synthesize on declared support, what propagating content survives on the radiative range, and what a finite instrument can recover at declared tolerance \cite{devaney2004}. Confusing those readings produces the rank gaps and watt failures that later sections quantify \cite{pal2026}.
Rayleigh and Stratton treated such ranges as boundary-value laboratories long before digital near-field transforms became routine \cite{stratton1941}. Harrington's method-of-moments culture reinforced the same discipline: metal outline first, harmonic headline second \cite{harrington2001}.
Subsection~\ref{sec:ch214} states the formal export condition that this range scene presupposes \cite{stratton1941}.
\end{physicalvignette}
Subsections~\ref{sec:ch211}--\ref{sec:ch213} supplied spaces, solution domains,
compactness, and spectral compression templates. Subsection~\ref{sec:ch214}
identifies self-adjoint and coercive realizations of truncated Maxwell operators
on bounded domains, records completion and quotient constructions that remove cavity
kernels, and states the $\mathcal{D}_J=1$ embedding that links Chapter~\ref{ch:01}
finite-energy criteria to the solution domains declared here
\cite{stratton1941}. Weak and energy forms
Eqs.~\eqref{eq:ch2.maxwell-energy-form}--\eqref{eq:ch2.maxwell-weak-form} and the
adjoint pairing Eq.~\eqref{eq:ch2.maxwell-adjoint-pairing} are presupposed; they
are \cite{stratton1941}.
A densely defined operator \(T\) on Hilbert space \(\mathcal{H}\) is
\emph{symmetric} if
\(\langle T\mathbf{u},\mathbf{v}\rangle=\langle\mathbf{u},T\mathbf{v}\rangle\)
for all \(\mathbf{u},\mathbf{v}\in\mathcal{D}(T)\), and \emph{self-adjoint} if
\(\mathcal{D}(T^{\dagger})=\mathcal{D}(T)\) with equality of adjoints
\cite{stratton1941}. Maxwell operators are symmetric on suitable domains because
Eq.~\eqref{eq:ch2.maxwell-adjoint-pairing} cancels boundary pairings when
\(\mathcal{B}_\Gamma\) is passive and homogeneous on test states. Self-adjointness
requires matching domain and range of adjoints; that match fails on exterior
domains without restricting to outgoing subspaces \cite{harrington2001}. On bounded \(\Omega\) with passive \(\mathcal{B}_\Gamma\), the sesquilinear form
\(\mathfrak{a}_\omega\) of Eq.~\eqref{eq:ch2.maxwell-energy-form} is coercive on
the quotient by cavity modes:
\begin{equation}
\label{eq:ch2.coercivity-quotient}
\mathfrak{a}_\omega(\mathbf{u},\mathbf{u})
\ge
\alpha\,\|\mathbf{u}\|_{\mathcal{U}}^{2},
\qquad
\mathbf{u}\in\mathcal{U}_\omega^{0}/\ker\mathcal{L}_\omega,
\end{equation}
for some \(\alpha>0\) when \(\ker\mathcal{L}_\omega\) is finite-dimensional and
consists of divergence-free resonances compatible with boundary class
\cite{stratton1941}. Equation~\eqref{eq:ch2.coercivity-quotient}
is the Lax--Milgram hypothesis that makes
Eq.~\eqref{eq:ch2.maxwell-weak-form} well posed on the quotient. Coercivity fails
on the unquotiented space because cavity modes satisfy
\(\mathfrak{a}_\omega(\mathbf{u},\mathbf{u})=0\) while
\(\|\mathbf{u}\|_{\mathcal{U}}>0\).
\begin{figure}[t]
\centering
\ChOneFigBlock{%
\begin{tikzpicture}[ch1fig, node distance=5mm and 7mm]
 \node[ch1box] (u) {$\mathcal{U}_\omega^{0}$}; \node[ch1box, right=15mm of u] (q) {$\mathcal{U}_\omega^{0}/\ker\mathcal{L}_\omega$};
 \node[ch1box, below=5mm of q] (sol) {Unique weak solution};
 \draw[ch1flow] (u.east) -- node[above,font=\scriptsize]{quotient} (q.west);
 \draw[ch1flow] (q.south) -- (q.south |- sol.north);
 \draw[ch1flow] (q.south |- sol.north) -- (sol.north);
\end{tikzpicture}%
}
\caption{Cavity-mode quotient: coercivity and uniqueness hold on the quotient space
after removing \(\ker\mathcal{L}_\omega\).}
\label{fig:ch2.cavity-quotient-schematic}
\end{figure}
Physical excitation at a cavity frequency
requires impedance or source coupling that breaks pure standing-wave balance;
Subsection~\ref{sec:ch214} records the linear-algebraic quotient; Section~\ref{sec:ch263}
develops spectral consequences \cite{stratton1941}. Impedance and Leontovich completions from
Eq.~\eqref{eq:ch2.boundary-impedance-operator} yield self-adjoint boundary
realizations when surface impedance is passive:
\begin{equation}
\label{eq:ch2.passive-impedance-class}
\Re\{Z_s\}\ge 0\ \text{on}\ \Gamma_Z
\quad\Longrightarrow\quad
\mathcal{B}_\Gamma\ \text{symmetric on}\ \mathcal{T}_\Gamma.
\end{equation}
Passive \(\Re\{Z_s\}\ge 0\) dissipates power through the boundary rather than
injecting net energy; active or negative-real parts require revisiting Maxwell-closed
as noted in Section~\ref{sec:ch01.1.2} \cite{stratton1941}. Perfect conductors
select Eq.~\eqref{eq:ch2.pec-subspace}; impedance boundaries enlarge
\(\mathcal{D}(\mathcal{B}_\Gamma)\) while preserving symmetry when
Eq.~\eqref{eq:ch2.passive-impedance-class} holds \cite{harrington2001}.
When \(\mathcal{B}_\Gamma\) is symmetric but the Maxwell operator is only densely
defined, the Friedrichs extension supplies a self-adjoint realization:
\begin{equation}
\label{eq:ch2.friedrichs-extension}
\mathcal{L}_\omega^{\mathrm{F}}
=
\overline{\mathcal{L}_\omega}\big|_{\mathcal{D}(\mathfrak{a}_\omega)},
\end{equation}
the completion of \(\mathcal{L}_\omega\) on the energy domain of
\(\mathfrak{a}_\omega\) \cite{stratton1941}. Equation
\eqref{eq:ch2.friedrichs-extension} is the locked definition for for
self-adjoint Maxwell truncations on bounded domains. Exterior outgoing problems
use a different completion template: restrict to the outgoing subspace selected in
Section~\ref{sec:ch243} rather than applying Eq.~\eqref{eq:ch2.friedrichs-extension}
to all of \(\mathcal{U}_\omega\) \cite{harrington2001}. Well-posedness on the quotient combines coercivity and bounded forcing:
\begin{equation}
\label{eq:ch2.lax-milgram-wellposed}
\exists!\,\mathbf{u}_\omega\in\mathcal{U}_\omega^{0}/\ker\mathcal{L}_\omega:
\mathfrak{a}_\omega(\mathbf{u}_\omega,\mathbf{v})
=
\mathfrak{f}_\omega(\mathbf{v})
\quad\forall\,\mathbf{v}\in\mathcal{U}_\omega^{0}/\ker\mathcal{L}_\omega,
\end{equation}
with \(\mathfrak{f}_\omega\) from Eq.~\eqref{eq:ch2.forcing-functional-operator}.
Stability follows from Eqs.~\eqref{eq:ch2.coercivity-quotient} and
\eqref{eq:ch2.solution-map-stability} \cite{stratton1941}. Resolvent bound
Eq.~\eqref{eq:ch2.resolvent-bounded} holds on the quotient when
\(\mathcal{L}_\omega^{\mathrm{F}}\) is bijective there.
\begin{figure}[t]
\centering
\ChOneFigBlock{%
\begin{tikzpicture}[ch1fig, node distance=5mm and 7mm]
 \node[ch1box] (a) {$\mathfrak{a}_\omega$};
 \node[ch1box, right=7mm of a] (lm) {Lax--Milgram};
 \node[ch1box, right=7mm of lm] (sw) {$\mathcal{S}_\omega$};
 \node[ch1box, below=5mm of a] (b) {$\mathcal{B}_\Gamma$ passive};
 \draw[ch1flow] (a) -- (lm) -- (sw);
 \draw[ch1flow] (b.south) -- (b.south |- a.north);
 \draw[ch1flow] (b.south |- a.north) -- (a.north);
\end{tikzpicture}%
}
\caption{Self-adjoint completion pathway on bounded domains: passive
\(\mathcal{B}_\Gamma\) supports coercive \(\mathfrak{a}_\omega\) and stable
\(\mathcal{S}_\omega\) on the cavity quotient.}
\label{fig:ch2.self-adjoint-closure-pathway}
\end{figure}
Changing \(\mathcal{B}_\Gamma\) while leaving
\(\mathfrak{a}_\omega\) formally unchanged can destroy coercivity or symmetry;
that sensitivity was already visible in the resolvent-branch illustration of
Subsection~\ref{sec:ch212} \cite{stratton1941}. charge-continuous from Subsection~\ref{sec:ch01.0.3} embeds into solution domains.
The \ref{ch:01} finite-energy class
Eq.~\eqref{eq:ch1.admissible-space} requires
\((\E,\H)\in L^{2}(\Omega)^{3}\times L^{2}(\Omega)^{3}\) with compatible boundary
and source data \cite{stratton1941}. The \ref{ch:02} refinement is
\begin{equation}
\label{eq:ch2.admissibility-embedding}
\mathcal{X}_{\mathrm{ch1}}(\Omega)\cap\mathcal{U}_\omega
\subseteq
\mathcal{F}_\Omega
\subseteq
\mathcal{D}(\mathcal{L}_\omega),
\end{equation}
with strict inclusions in general because \(\HCurl\) control and trace membership
strengthen \(L^{2}\) Maxwell-closed without replacing
Eq.~\eqref{eq:ch1.continuity-charge-continuous} \cite{stratton1941}. Equation
\eqref{eq:ch2.admissibility-embedding} specializes
Eq.~\eqref{eq:ch2.continuity-operator-domain}: when
\(\mathcal{D}_{\mathrm{ch1}}=1\), impressed data lie in \(\mathcal{J}_\Omega\) and
closed solutions lie in \(\mathcal{F}_\Omega\). Operator Maxwell-closed adds
graph-norm control:
\begin{equation}
\label{eq:ch2.graph-charge-continuous-criterion}
\mathbf{u}_\omega\in\mathcal{F}_\Omega
\ \Longrightarrow\
\|\mathbf{u}_\omega\|_{\mathrm{gr}(\mathcal{L}_\omega)}<\infty
\ \text{when}\ \mathcal{L}_\omega[\mathbf{u}_\omega]\ \text{is defined}.
\end{equation}
Inequality~\eqref{eq:ch2.graph-admissibility-criterion} is the quantitative form of
``physically realizable'' field pairs for variational solvers
\cite{harrington2001}.
Source-side Maxwell-closed remains coupled to
Eq.~\eqref{eq:ch2.source-space-refined} and
Eq.~\eqref{eq:ch1.source-charge-continuous} by reference. Completion operators must not
introduce incompatible traces:
\begin{equation}
\label{eq:ch2.closure-admissibility-lock}
\mathcal{B}_\Gamma\ \text{admissible}
\iff
\mathcal{R}(\gamma_\tau|_{\mathcal{F}_\Omega})\supseteq
\mathcal{D}(\mathcal{B}_\Gamma).
\end{equation}
Failure of Eq.~\eqref{eq:ch2.closure-admissibility-lock} means prescribed
\(\mathbf{g}_\Gamma\) is not the trace of any \(\HCurl\) field, reproducing the
trace-mismatch illustration of Subsection~\ref{sec:ch211} \cite{stratton1941}. Exterior domains: outgoing radiation condition from Section~\ref{sec:ch243} defines a
\emph{physically self-adjoint} branch rather than a Hilbert self-adjoint operator
on all of \(\mathcal{U}_\omega\):
\begin{equation}
\label{eq:ch2.outgoing-nonselfadjoint-split}
\mathcal{U}_\omega
=
\mathcal{U}_{\mathrm{out}}\oplus\mathcal{U}_{\mathrm{re}},
\qquad
\mathcal{L}_\omega\ \text{coercive on}\ \mathcal{U}_{\mathrm{out}},
\end{equation}
with reactive storage content in \(\mathcal{U}_{\mathrm{re}}\) controlled by
regime criteria of Section~\ref{sec:ch01.2.4} rather than by full coercivity on
\(\mathcal{U}_\omega\) \cite{stratton1941}. Conflating
Eq.~\eqref{eq:ch2.friedrichs-extension} with exterior radiation problems produces
incorrect cavity-like spectra in open regions.
With Subsections~\ref{sec:ch211}--\ref{sec:ch214} complete, Section~\ref{sec:ch21}
has moved from minimal declarations in Section~\ref{sec:ch20} to a closed operator
calculus on Sobolev spaces: traces, adjoints, compactness, and self-adjoint
truncations compatible with \ref{ch:01} charge-continuous \cite{stratton1941}.
Section~\ref{sec:ch22} applies those spaces to time--frequency equivalences;
Sections~\ref{sec:ch23}--\ref{sec:ch24} inherit $\mathcal{D}(\mathcal{L}_\omega)$ and
$\mathcal{F}_\Omega$ as fixed input classes. \ref{ch:04} states which operator outputs
survive finite source support and manufacturable geometry \cite{devaney2004}.
\section{Time- and Frequency-Domain Maxwell Operators}
\label{sec:ch22}
Function spaces and Maxwell operators in Section~\ref{sec:ch21} fix admissible sources and field pairs before time-harmonic reduction. Section~\ref{sec:ch22} carries the same physics into frequency-domain maps used on every fixed-frequency range \cite{stratton1941}. The laboratory never measures abstract Hilbert norms; it measures phasor fields at declared $\omega$ \cite{maxwell1873}. An oscilloscope trace records $\E(\mathbf{r},t)$ and $\H(\mathbf{r},t)$ as functions of time; a vector network analyzer records complex amplitudes at each stepped tone. Both instruments interrogate the same structure, but only the time-domain picture makes retarded support visible before any phasor reduction is declared \cite{harrington2001}.
% PHYSICAL-MEMORY:sec:ch22
\begin{physicalvignette}[Wire dipole over a finite ground plane]
Consider Wire dipole over a finite ground plane. Before any coefficient table is printed, technicians record frequency, cable loss, and the mechanical outline that fixes electrical size. The first walk-through of the range is always tactile: flange bolts, probe standoff, turntable boresight, and the wattmeter reference on the receive chain \cite{hansen1988}.
On this hardware the near zone is visually impressive---standing reactive structure, rapid phase rotation, sometimes a visible VSWR ripple at the feed---yet the export question is measured only after outgoing selection on a large sphere \cite{harrington2001}. What the camera sees in the chamber is not automatically what the far-zone pattern claims \cite{stratton1941}.
Engineers therefore learn to separate three readings: what the feed can synthesize on declared support, what propagating content survives on the radiative range, and what a finite instrument can recover at declared tolerance \cite{devaney2004}. Confusing those readings produces the rank gaps and watt failures that later sections quantify \cite{pal2026}.
Rayleigh and Stratton treated such ranges as boundary-value laboratories long before digital near-field transforms became routine \cite{stratton1941}. Harrington's method-of-moments culture reinforced the same discipline: metal outline first, harmonic headline second \cite{harrington2001}.
Subsection~\ref{sec:ch22} states the formal export condition that this range scene presupposes \cite{stratton1941}.
\end{physicalvignette}
Fixed-$\omega$ phasors are frequency slices of retarded evolution, valid when excitation
and media are narrowband on the declared interval \cite{stratton1941}.
Each harmonic operator $\mathcal{L}_\omega$ is a labeled
cross-section of a single causal evolution law, not as an independent postulate.
A vector network
analyzer sweep therefore samples one harmonic operator family $\mathcal{L}_\omega$ at each
tone---not a separate physics from time-domain causality
\cite{maxwell1873}. Every compact-range calibration and harmonic
finite-difference cell presupposes that reduction.
The harmonic operators $\mathcal{L}_\omega$ and $\mathcal{S}_\omega$ introduced earlier
describe individual frequency slices of Maxwell evolution. The purpose of the present
section is to show how those slices emerge from a single causal time-domain theory
\cite{stratton1941}. Section~\ref{sec:ch22} develops the time--frequency equivalences that
connect harmonic operators to causal evolution, Helmholtz symbols, and
propagating-versus-evanescent splits under the $\exp(-\jj\omega t)$ convention locked in
\ref{ch:01}. It records when that reduction is valid, how retarded time-domain maps
relate to $\mathcal{R}_\omega$, and which spectral components carry far-zone transport
content importable from Section~\ref{sec:ch01.2.4}.
Harmonic Maxwell theory in Section~\ref{sec:ch20} treated $\omega>0$ as a fixed parameter.
Time-domain formulations restore the full operator pencil $\partial/\partial t$ and clarify
why frequency-domain solvers implement causal limits when inversion is performed correctly.
Subsection~\ref{sec:ch221} gives the \ref{ch:02} account of time-domain evolution;
Subsection~\ref{sec:ch222} packages frequency-domain Maxwell and curl--curl Helmholtz
operators without re-displaying Eq.~\eqref{eq:ch1.maxwell-local}; Subsection~\ref{sec:ch223}
proves equivalence between retarded Green propagation and $\mathcal{S}_\omega$ stated in
Eq.~\eqref{eq:ch2.causality-retarded}; Subsection~\ref{sec:ch224} applies the
reactive-radiative partition of Section~\ref{sec:ch01.2.4} to propagating, evanescent, and
non-propagating subspaces in operator form \cite{stratton1941}.
Several related constructions are deferred deliberately. In particular, dyadic
Green-function development and inversion theory are postponed until Section~\ref{sec:ch23},
where the harmonic operators introduced here acquire explicit integral representations
\cite{stratton1941}. Section~\ref{sec:ch22} does not re-open function-space declarations
from Section~\ref{sec:ch21} or local Maxwell derivations from \ref{ch:01}.
Helmholtz and curl--curl reductions developed in Subsection~\ref{sec:ch222} are
single-field tools: they simplify kernel construction but do not replace the first-order
$\mathcal{L}_\omega$ packaging of Eq.~\eqref{eq:ch2.maxwell-boundary-composite} when
boundary rows and impedance closures must be enforced jointly \cite{harrington2001}.
Within the locked harmonic band, $\mathcal{S}_\omega$ is the $\omega$-slice of a causal
time-domain map whenever impressed data and $\mathcal{B}_\Gamma$ are held fixed across the
transform \cite{stratton1941}. Finite-difference time-domain schemes approximate
Subsection~\ref{sec:ch221}; frequency-domain finite-element solvers approximate
Subsection~\ref{sec:ch222} on the domains of Eq.~\eqref{eq:ch2.domain-refinement-L}.
Modal and plane-wave decompositions in Subsection~\ref{sec:ch224} filter which components
of $\mathcal{S}_\omega[\Jimp]$ may enter $\mathcal{R}_\omega$ under
Eq.~\eqref{eq:ch2.transport-qualified-radiation} \cite{harrington2001}.
\subsection{Time-domain evolution}
\label{sec:ch221}
% HARRINGTON-OPEN:sec:ch221
Time-harmonic phasors match fixed-frequency range measurements; $\exp(-\j\omega t)$ fixes the sign of reactive terms \cite{stratton1941}. Maxwell's retarded export to $\Gamma_R$ remains the governing constraint \cite{maxwell1873}.
The physical problem is causal propagation: a switched feed, a pulsed radar, or a transient
immunity test excites fields whose support in $(\mathbf{r},t)$ must respect finite signal
speed \cite{maxwell1873}. Time-harmonic phasors are how network analyzers, compact
ranges, and chamber receivers report steady-state radiation at each frequency step
\cite{stratton1941}. The $\exp(-\jj\omega t)$ convention fixed in Chapter~\ref{ch:01} aligns
Poynting flux with wattmeter readings and VNA phase references \cite{poynting1884}. Phasor
analysis is best viewed as a steady-state reduction of a fundamentally causal theory rather
than as the primary formulation of electromagnetic dynamics \cite{harrington2001}. Even small
phase errors can rotate synthesized patterns when near-field data are back-propagated
\cite{hansen1988}. The time-domain picture below states the causal support that every VNA tone
must inherit when $\mathcal{L}_\omega$ is interpreted as a frequency slice rather than an
independent postulate \cite{stratton1941}.
% PHYSICAL-MEMORY:sec:ch221
\begin{physicalvignette}[Slot array in a flat plate]
Consider Slot array in a flat plate. Before any coefficient table is printed, technicians record frequency, cable loss, and the mechanical outline that fixes electrical size. The first walk-through of the range is always tactile: flange bolts, probe standoff, turntable boresight, and the wattmeter reference on the receive chain \cite{hansen1988}.
On this hardware the near zone is visually impressive---standing reactive structure, rapid phase rotation, sometimes a visible VSWR ripple at the feed---yet the export question is measured only after outgoing selection on a large sphere \cite{harrington2001}. What the camera sees in the chamber is not automatically what the far-zone pattern claims \cite{stratton1941}.
Engineers therefore learn to separate three readings: what the feed can synthesize on declared support, what propagating content survives on the radiative range, and what a finite instrument can recover at declared tolerance \cite{devaney2004}. Confusing those readings produces the rank gaps and watt failures that later sections quantify \cite{pal2026}.
Time-harmonic phasors match fixed-frequency range measurements; $\exp(-\jj\omega t)$ fixes the sign of reactive terms \cite{stratton1941}. 
Rayleigh and Stratton treated such ranges as boundary-value laboratories long before digital near-field transforms became routine \cite{stratton1941}. Harrington's method-of-moments culture reinforced the same discipline: metal outline first, harmonic headline second \cite{harrington2001}.
Subsection~\ref{sec:ch221} states the formal export condition that this range scene presupposes \cite{stratton1941}.
\end{physicalvignette}
Subsection~\ref{sec:ch221} supplies the time-domain foundation underlying every harmonic
formulation used later in the chapter. It introduces the first-order Maxwell evolution
system, the associated evolution operator, and the causal support restrictions that all
physically admissible frequency-domain solutions must inherit \cite{stratton1941}.
Time-domain evolution specifies retarded trajectories: the field observed at
$(\mathbf{r},t)$ depends only on source events contained within the corresponding past
light cone \cite{stratton1941}. Any numerical or analytical approximation that violates
this support structure introduces nonphysical precursors and loses its connection to
Maxwell causality \cite{harrington2001}. Local Maxwell balances remain those of
Eq.~\eqref{eq:ch1.ibvp-time-domain}; they are cited. Function spaces from
Section~\ref{sec:ch21} enter through time-dependent embeddings developed below; Sobolev and
trace declarations are not repeated \cite{stratton1941}.
Let \(\Omega\subset\mathbb{R}^{3}\) be a Lipschitz domain with boundary
\(\Gamma=\partial\Omega\). Real time-domain fields
\(\E(\mathbf{r},t)\), \(\H(\mathbf{r},t)\) use SI units and the same constitutive
class as \ref{ch:01}.
The evolution problem is posed on trajectories of the electromagnetic state rather
than on isolated snapshots at a single frequency.
A pulsed or switched excitation on a feed line makes the distinction visible: the field at
$(\mathbf{r},t)$ cannot respond until source events inside the backward light cone have
occurred. Time-domain evolution is therefore the laboratory picture behind every harmonic
slice $\mathcal{S}_\omega$ used later in the chapter \cite{stratton1941}. The state vector
$\mathbf{U}(t)$ packages the coupled $(\E,\H)$ fields at frozen time $t$ into a single
object on which the Maxwell generator acts; this packaging mirrors the phasor state $\mathbf{u}_\omega$ of Eq.~\eqref{eq:ch2.maxwell-state-vector},
with time retained rather than suppressed \cite{harrington2001}. Collect the
electromagnetic state at time \(t\) as
\begin{equation}
\label{eq:ch2.time-state-vector}
\mathbf{U}(t)
=
\begin{bmatrix}\E(\cdot,t)\\ \H(\cdot,t)\end{bmatrix}
\in\mathcal{V}(\Omega),
\end{equation}
where \(\mathcal{V}(\Omega)\) is the time-dependent product space of square-integrable
field pairs compatible with Eq.~\eqref{eq:ch1.admissible-space} at each frozen time
slice \cite{stratton1941}. Equation~\eqref{eq:ch2.time-state-vector} is the
time-domain analogue of Eq.~\eqref{eq:ch2.maxwell-state-vector}; phasor symbols are
reserved for the \(\exp(-\jj\omega t)\) reduction in Subsection~\ref{sec:ch222}. The time-domain Maxwell
system is the causal law from which every harmonic operator in Section~\ref{sec:ch22} is
obtained by phasor substitution. Physically, Eqs.~\eqref{eq:ch2.time-maxwell-evolution}
state that changing $\E$ in space drives $\H$ in time through Faraday's law, and changing
$\H$ in space drives $\E$ in time through Amp\`ere's law with impressed current as external
drive \cite{maxwell1873}. From an antenna engineer's perspective, this is the
system that finite-difference time-domain codes march forward in time; the harmonic solvers
of Subsection~\ref{sec:ch222} are steady-state limits of the same law when all data are
single-frequency \cite{harrington2001}. The time-domain Maxwell system imported from
Eq.~\eqref{eq:ch1.ibvp-time-domain} reads in vacuum constitutive form as
\begin{equation}
\label{eq:ch2.time-maxwell-evolution}
\begin{aligned}
\curl\E(\mathbf{r},t)&=-\mu_0\,\partial_t\H(\mathbf{r},t),\\
\curl\H(\mathbf{r},t)&=\mathbf{j}_{\mathrm{imp}}(\mathbf{r},t)+\varepsilon_0\,\partial_t\E(\mathbf{r},t),\\
\divg\E(\mathbf{r},t)&=\frac{\rho(\mathbf{r},t)}{\varepsilon_0},\qquad
\divg\H(\mathbf{r},t)=0,
\end{aligned}
\end{equation}
with impressed current density \(\mathbf{j}_{\mathrm{imp}}(\mathbf{r},t)\) and charge
density \(\rho(\mathbf{r},t)\). Equation~\eqref{eq:ch2.time-maxwell-evolution}
packages local dynamics as a first-order system in \((\E,\H)\) and \(\partial_t\);
the hyperbolic coupling is the time-domain face of the curl structure in
Eq.~\eqref{eq:ch1.maxwell-operator-block} \cite{stratton1941}. Charge continuity
requires
\begin{equation}
\label{eq:ch2.charge-continuity-time}
\divg\mathbf{j}_{\mathrm{imp}}(\mathbf{r},t)+\partial_t\rho(\mathbf{r},t)=0
\end{equation}
in the sense of distributions when \(\rho\) is present \cite{stratton1941}. Pairs
\((\mathbf{j}_{\mathrm{imp}},\rho)\) violating
Eq.~\eqref{eq:ch2.charge-continuity-time} lie outside the time-domain admissible
class parallel to Eq.~\eqref{eq:ch2.source-admissibility}.
\begin{figure}[t]
\centering
\ChOneFigBlock{%
\begin{tikzpicture}[ch1fig, node distance=5mm and 7mm]
 \node[ch1box] (ic) {Initial data\\$(\E_0,\H_0)$};
 \node[ch1box, right=7mm of ic] (src) {Impressed\\$\mathbf{j}_{\mathrm{imp}}(\cdot,t)$};
 \node[ch1box, right=7mm of src] (ev) {Evolution\\$\mathcal{T}(t)$};
 \node[ch1box, right=7mm of ev] (st) {State $\mathbf{U}(t)$};
 \draw[ch1flow] (ic.east) -- (ic.east |- ev.west);
 \draw[ch1flow] (ic.east |- ev.west) -- (ev.west);
 \draw[ch1flow] (src.east) -- (src.east |- ev.west);
 \draw[ch1flow] (src.east |- ev.west) -- (ev.west);
 \draw[ch1flow] (ev.east) -- (ev.east |- st.west);
 \draw[ch1flow] (ev.east |- st.west) -- (st.west);
\end{tikzpicture}%
}
\caption{Time-domain evolution: initial data and impressed sources drive the
semigroup \(\mathcal{T}(t)\) to field state \(\mathbf{U}(t)\) under boundary class
\(\mathcal{B}_\Gamma\).}
\label{fig:ch2.time-evolution-measurement chain}
\end{figure}
Boundary data enter through the same 
\(\mathcal{B}_\Gamma\) of Eq.~\eqref{eq:ch1.boundary-data-operator}, now evaluated
on time-dependent traces \(\gamma_\tau[\mathbf{U}(t)]\in\mathcal{T}_\Gamma\) for
each \(t\ge 0\) \cite{stratton1941}. The closed time-domain problem
is to find \(\mathbf{U}(t)\) such that
\begin{equation}
\label{eq:ch2.time-ibvp-package}
\begin{cases}
\partial_t\mathbf{U}=\mathcal{A}\,\mathbf{U}+\mathbf{F}(t), & \mathbf{r}\in\Omega,\ t>0,\\
\mathbf{U}(0)=\mathbf{U}_0, & \mathbf{r}\in\Omega,\\
\mathcal{B}_\Gamma[\mathbf{U}(t)]=\mathbf{g}_\Gamma(t), & \mathbf{r}\in\Gamma,\ t\ge 0,
\end{cases}
\end{equation}
where \(\mathcal{A}\) is the Maxwell generator induced by
Eq.~\eqref{eq:ch2.time-maxwell-evolution} on
\(\mathcal{D}(\mathcal{A})\subset\mathcal{V}(\Omega)\) and
\(\mathbf{F}(t)\) collects impressed-source rows \cite{stratton1941}. Equation
\eqref{eq:ch2.time-ibvp-package} is the time-domain counterpart of
Eq.~\eqref{eq:ch2.closed-maxwell-bvp}; uniqueness on a fixed branch is inherited
from Section~\ref{sec:ch01.1.3} when data classes coincide with the harmonic limit
\cite{stratton1941}. Define the evolution operator (semigroup) \(\mathcal{T}(t)\) by
\begin{equation}
\label{eq:ch2.time-semigroup-definition}
\mathcal{T}(t):\mathbf{U}_0\longmapsto \mathbf{U}(t),
\qquad
\mathcal{T}(0)=\mathrm{id},\qquad
\mathcal{T}(t+s)=\mathcal{T}(t)\mathcal{T}(s),
\end{equation}
when \(\mathbf{F}\equiv\mathbf{0}\) and \(\mathcal{B}_\Gamma\) is homogeneous. With
forcing,
\begin{equation}
\label{eq:ch2.time-evolution-forced}
\mathbf{U}(t)
=
\mathcal{T}(t)\mathbf{U}_0
\int_0^t \mathcal{T}(t-s)\,\mathbf{F}(s)\,ds,
\end{equation}
the Duhamel representation on \(\mathcal{D}(\mathcal{A})\)
\cite{stratton1941}. Equation~\eqref{eq:ch2.time-evolution-forced}
makes causality explicit: \(\mathbf{U}(t)\) depends only on
\(\mathbf{F}(s)\) for \(s\le t\). The time-domain \emph{solution map} sends impressed
data and initial state to trajectories:
\begin{equation}
\label{eq:ch2.time-domain-solution-map}
\mathcal{S}_{t}:\big(\mathbf{U}_0,\mathbf{j}_{\mathrm{imp}}\big)\longmapsto
\{\mathbf{U}(t)\}_{t\ge 0}.
\end{equation}
Harmonic \(\mathcal{S}_\omega\) of Eq.~\eqref{eq:ch2.maxwell-solution-map} is a
single-frequency slice of Eq.~\eqref{eq:ch2.time-domain-solution-map} when all data
are proportional to \(\exp(-\jj\omega t)\) and \(\mathcal{B}_\Gamma\) is fixed
\cite{stratton1941}; the equivalence is proved in Subsection~\ref{sec:ch223}.
Retarded representation specializes Eq.~\eqref{eq:ch2.time-evolution-forced} when
Green kernels exist:
\begin{equation}
\label{eq:ch2.retarded-convolution-form}
\E(\mathbf{r},t)
=
\E_{\mathrm{hom}}(\mathbf{r},t)
-\mu_0\,\partial_t
\int_0^t\!\!\int_\Omega
\mathcal{G}(\mathbf{r},\mathbf{r}',t-t')\cdot
\mathbf{j}_{\mathrm{imp}}(\mathbf{r}',t')\,d^{3}r'\,dt',
\end{equation}
with \(\mathcal{G}(\mathbf{r},\mathbf{r}',t-t')=\mathbf{0}\) for \(t-t'<0\)
\cite{stratton1941}. Equation~\eqref{eq:ch2.retarded-convolution-form} refines the
Eq.~\eqref{eq:ch2.causality-retarded} without constructing \(\mathcal{G}\)
here; kernel development is in Section~\ref{sec:ch23}. The support condition
\begin{equation}
\label{eq:ch2.causal-support-constraint}
\operatorname{supp}_{(\mathbf{r},t)}\E
\subseteq
\big\{(\mathbf{r},t): t\ge t_{\mathrm{ret}}(\mathbf{r})\big\}
\end{equation}
expresses finite propagation speed \(c_0=1/\sqrt{\mu_0\varepsilon_0}\) on vacuum
regions \cite{stratton1941}. Figure~\ref{fig:ch2.causal-light-cone}
schematizes Eq.~\eqref{eq:ch2.causal-support-constraint}.
\begin{figure}[t]
\centering
\ChOneFigBlock{%
\begin{tikzpicture}[ch1fig, node distance=5mm and 7mm]
 \draw[ch1flow] (0,0) -- (2.2,2.2);
 \draw[ch1flow] (0,0) -- (2.2,-2.2);
 \node[ch1oncanvas, right=2mm] at (2.2,2.2) {$t=t_{\mathrm{ret}}$};
 \node[ch1box] at (0,0) {Source event};
\end{tikzpicture}%
}
\caption{Causal support: field at \((\mathbf{r},t)\) depends on sources inside the
backward light cone; retarded integration uses \(t'\le t\) only.}
\label{fig:ch2.causal-light-cone}
\end{figure}
Energy balance in time domain imports the Poynting and storage laws of
Section~\ref{sec:ch01.1.4}. With volume energy density
\begin{equation}
\label{eq:ch2.time-energy-density}
u(\mathbf{r},t)
=
\frac{\varepsilon_0}{2}\|\E(\mathbf{r},t)\|^2
+\frac{\mu_0}{2}\|\H(\mathbf{r},t)\|^2,
\end{equation}
integration over a control volume \(V\subset\Omega\) yields
\begin{equation}
\label{eq:ch2.time-poynting-balance}
\frac{d}{dt}\int_V u\,dV
+\int_{\partial V}\mathbf{S}\cdot\hat{\mathbf{n}}\,dS
=
-\int_V \E\cdot\mathbf{j}_{\mathrm{imp}}\,dV,
\end{equation}
with \(\mathbf{S}=\E\times\H\) \cite{stratton1941}. Equation
\eqref{eq:ch2.time-poynting-balance} is cited; it supplies the physical
reading of Eq.~\eqref{eq:ch2.time-evolution-forced}: impressed work on the right
drives stored and exported energy on the left. Reactive exchange with boundaries
appears in the surface term and motivates the regime split of
Subsection~\ref{sec:ch224} \cite{harrington2001}. Harmonic operators follow from the substitution
\(\E(\mathbf{r},t)=\Re\{\underline{\E}(\mathbf{r})\eim\}\) with
\(\underline{\H}\) analogous, yielding
\begin{equation}
\label{eq:ch2.time-to-frequency-map}
\partial_t \longrightarrow -\jj\omega,
\qquad
\mathcal{A}\longrightarrow \mathcal{M}_\omega,
\qquad
\mathcal{S}_t\big|_{\omega}\longrightarrow \mathcal{S}_\omega,
\end{equation}
on admissible data of Eq.~\eqref{eq:ch2.admissibility-embedding}
\cite{stratton1941}. Equation~\eqref{eq:ch2.time-to-frequency-map} is symbolic;
Subsection~\ref{sec:ch223} states equivalence conditions under which frequency-domain
solvers implement causal time-domain maps. Exterior domains require outgoing
boundary conditions in time domain compatible with
Eq.~\eqref{eq:ch1.radiation-trace-condition}; those completions are sharpened in
Section~\ref{sec:ch243} \cite{stratton1941}.
Time-domain Maxwell-closed couples Eq.~\eqref{eq:ch1.admissible-space} at each time
slice to Eq.~\eqref{eq:ch2.charge-continuity-time}:
\begin{equation}
\label{eq:ch2.time-charge-continuous-coupling}
\mathbf{U}(t)\in\mathcal{F}_\Omega\ \text{a.e.\ }t
\ \wedge\
(\mathbf{j}_{\mathrm{imp}},\rho)\ \text{satisfy charge continuity}
\ \Longrightarrow\
\mathbf{U}\in\mathcal{D}(\mathcal{S}_{t}).
\end{equation}
Finite-difference time-domain schemes discretize
Eq.~\eqref{eq:ch2.time-maxwell-evolution} on staggered grids; stability requires
discrete analogues of Eq.~\eqref{eq:ch2.causal-support-constraint} and consistent
boundary enforcement of \(\mathcal{B}_\Gamma\) \cite{harrington2001}. Subsection
~\ref{sec:ch221} states the continuous targets those schemes approximate.
Subsection~\ref{sec:ch221} has posed the time-domain Maxwell system, the evolution
semigroup \(\mathcal{T}(t)\), and the retarded convolution form that enforces causal
support. Subsection~\ref{sec:ch222} carries the same physics into fixed-frequency
operators and Helmholtz symbols on the spaces of Section~\ref{sec:ch21}
\cite{stratton1941}.
\subsection{Frequency-domain Maxwell operators}
\label{sec:ch222}
% HARRINGTON-OPEN:sec:ch222
A standard-gain horn on a compact-range turntable fixes the laboratory scene for `Frequency-domain Maxwell operators'': frequency, probe standoff, and wattmeter closure on $\Gamma_R$ precede the formal displays below \cite{hansen1988}. Maxwell's retarded export to $\Gamma_R$ remains the governing constraint \cite{maxwell1873}.
The physical problem is steady-state interrogation at fixed frequency: every calibrated
$S$-parameter, resonator $Q$ measurement, and harmonic finite-element solve assumes that
excitation, media, and boundary class are frozen while $\omega$ is stepped
\cite{harrington2001}. A vector network analyzer sweep is the canonical laboratory observation that precedes
every harmonic operator declaration. The instrument steps $\omega$ through a calibrated
band, records complex $S$-parameters at each tone, and treats each frequency as a
steady-state slice of the same passive structure \cite{harrington2001}. The composite
operator $\mathcal{L}_\omega$ introduced in Section~\ref{sec:ch21} is precisely that
slice: Maxwell's equations with $\partial_t\to-\jj\omega$ on impressed data and boundary
class held fixed. Readers should therefore read $\mathcal{L}_\omega$ not as an abstract
Hilbert map alone but as the operator whose inversion reproduces the phasor fields a VNA
implicitly assumes when it reports $S_{21}(\omega)$ after de-embedding \cite{stratton1941}.
Frequency-domain formulations dominate practical electromagnetics because measurements,
simulations, and design specifications are usually organized around frequency. Network
analyzers, resonator measurements, and antenna-pattern ranges all interrogate electromagnetic
systems one frequency at a time \cite{harrington2001}. The harmonic operators
introduced below state that experimentally familiar viewpoint while remaining consistent
with the causal evolution described in Subsection~\ref{sec:ch221}.
% PHYSICAL-MEMORY:sec:ch222
\begin{physicalvignette}[Compact SNF chamber with rolled AUT mount]
Consider Compact SNF chamber with rolled AUT mount. Before any coefficient table is printed, technicians record frequency, cable loss, and the mechanical outline that fixes electrical size. The first walk-through of the range is always tactile: flange bolts, probe standoff, turntable boresight, and the wattmeter reference on the receive chain \cite{hansen1988}.
On this hardware the near zone is visually impressive---standing reactive structure, rapid phase rotation, sometimes a visible VSWR ripple at the feed---yet the export question is measured only after outgoing selection on a large sphere \cite{harrington2001}. What the camera sees in the chamber is not automatically what the far-zone pattern claims \cite{stratton1941}.
Engineers therefore learn to separate three readings: what the feed can synthesize on declared support, what propagating content survives on the radiative range, and what a finite instrument can recover at declared tolerance \cite{devaney2004}. Confusing those readings produces the rank gaps and watt failures that later sections quantify \cite{pal2026}.
Rayleigh and Stratton treated such ranges as boundary-value laboratories long before digital near-field transforms became routine \cite{stratton1941}. Harrington's method-of-moments culture reinforced the same discipline: metal outline first, harmonic headline second \cite{harrington2001}.
Subsection~\ref{sec:ch222} states the formal export condition that this range scene presupposes \cite{stratton1941}.
\end{physicalvignette}
Subsection~\ref{sec:ch221} fixed time-domain evolution and
Eq.~\eqref{eq:ch2.time-to-frequency-map}. Subsection~\ref{sec:ch222} gives the
Chapter~\ref{ch:02} account of harmonic Maxwell operators and curl--curl Helmholtz
symbols on the spaces of Section~\ref{sec:ch21} \cite{stratton1941}. Recall that
$\mathcal{S}_\omega:\mathcal{J}_\Omega\to\mathcal{F}_\Omega$ is the harmonic solution map,
$\mathcal{R}_\omega=\Pi_{\mathrm{rad}}\circ\mathcal{S}_\omega$ selects outgoing radiative
content, and $\mathcal{L}_\omega$ packages volume and boundary residuals jointly
\cite{harrington2001}.
The first-order system Eq.~\eqref{eq:ch1.maxwell-operator-block} and composite
\(\mathcal{L}_\omega\) of Eq.~\eqref{eq:ch2.maxwell-boundary-composite} are
presupposed; local Maxwell derivations from \ref{ch:01} are not repeated.
The electric-wave operator Eq.~\eqref{eq:ch2.electric-wave-operator} was recorded in
Subsection~\ref{sec:ch202}; here it is embedded in domain declarations and extended
by magnetic and material-weighted counterparts. Full vector-field separation is
deferred to Section~3.3 \cite{stratton1941}.
Under the locked \(\exp(-\jj\omega t)\) convention, differentiation with respect to
time converts to multiplication by \(-\jj\omega\) on phasor quantities, and the
second time derivative maps to \(-\omega^{2}\).
Every phasor displayed on a network analyzer trace has already undergone this
substitution: a time derivative in Maxwell's equations becomes a multiplication by
$-\jj\omega$ on the locked $\exp(-\jj\omega t)$ convention. Engineers should notice
that the substitution is not a notational convenience alone; it is the algebraic
condition under which a time-domain IBVP and a frequency-domain BVP describe the same
passive structure at fixed $\omega$ \cite{stratton1941}. The substitution fails when
excitation is broadband without superposition, when $\mathcal{B}_\Gamma$ varies in time,
or when material response is not single-frequency on the declared band
\cite{harrington2001}. With those limits stated, the formal rule is:
Time derivatives map to
\begin{equation}
\label{eq:ch2.frequency-substitution-rule}
\partial_t \longrightarrow -\jj\omega,
\qquad
\frac{\partial^2}{\partial t^2}\longrightarrow -\omega^2,
\end{equation}
on phasor fields \(\E(\mathbf{r})\), \(\H(\mathbf{r})\) and impressed data
\(\Jimp(\mathbf{r})\), \(\rho(\mathbf{r})\) \cite{stratton1941}.
In Eq.~\eqref{eq:ch2.frequency-substitution-rule},
\(\partial_t\) acting on a time-harmonic quantity is replaced by \(-\jj\omega\),
where \(\omega\) is in rad\,s\(^{-1}\); and \(\partial_t^{2}\) is replaced by
\(-\omega^{2}\).
This substitution is valid only when all time dependence is single-frequency or when
a superposition of such terms is taken with consistent phasor definitions.
Equation
\eqref{eq:ch2.frequency-substitution-rule} specializes
Eq.~\eqref{eq:ch2.time-to-frequency-map} and converts
Eq.~\eqref{eq:ch1.ibvp-time-domain} to the harmonic BVP of
Eq.~\eqref{eq:ch1.ibvp-harmonic-bvp} by reference.
The harmonic residual collects the four Maxwell rows in frequency form.
A finite-element or method-of-moments solve at fixed $\omega$ is, operationally, the search
for $\mathbf{u}_\omega=(\E,\H)$ that drives $\mathcal{M}_\omega[\mathbf{u}_\omega]$ to zero
in the volume while satisfying the boundary row enforced by $\mathcal{B}_\Gamma$.
The engineer observes residuals indirectly through solver convergence and port matching;
the mathematician packages the same requirement as membership in
$\mathcal{D}(\mathcal{L}_\omega)$ with $\mathcal{L}_\omega[\mathbf{u}_\omega]$ balanced
against impressed data \cite{harrington2001}. Readers should notice that
Eq.~\eqref{eq:ch2.harmonic-maxwell-residual} is not a new physics postulate but the
frequency image of Eq.~\eqref{eq:ch2.time-maxwell-evolution} once $\partial_t\to-\jj\omega$
on admissible data \cite{stratton1941}. The harmonic Maxwell residual
operator \(\mathcal{M}_\omega\) of Eq.~\eqref{eq:ch2.maxwell-driven-operator}
acts on \(\mathbf{u}_\omega=(\E,\H)\in\mathcal{U}_\omega\):
\begin{equation}
\label{eq:ch2.harmonic-maxwell-residual}
\mathcal{M}_\omega[\mathbf{u}_\omega]
=
\begin{bmatrix}
\curl\E+\jj\omega\mu_0\H\\
\curl\H-\jj\omega\varepsilon_0\E-\Jimp\\
\divg(\varepsilon_0\E)-\rho\\
\divg(\mu_0\H)
\end{bmatrix},
\end{equation}
row-aligned with Eq.~\eqref{eq:ch1.maxwell-operator-block} \cite{stratton1941}.
Each row of Eq.~\eqref{eq:ch2.harmonic-maxwell-residual} is a residual that must
vanish for a physically consistent harmonic state:
the first row enforces Faraday's law with phasor substitution;
the second enforces Amp\`ere's law with impressed current \(\Jimp\);
the third enforces Gauss's law for electric flux with charge density \(\rho\);
and the fourth enforces the magnetic Gauss law.
Equation~\eqref{eq:ch2.harmonic-maxwell-residual} is the frequency-domain packaging
of Eq.~\eqref{eq:ch2.time-maxwell-evolution}; completion and traces remain in
\(\mathcal{L}_\omega[\mathbf{u}_\omega]=[\mathbf{s}_\omega,\mathbf{g}_\Gamma]^{\mathsf{T}}\)
from Eq.~\eqref{eq:ch2.closed-maxwell-bvp}.
\begin{figure}[t]
\centering
\ChOneFigBlock{%
\begin{tikzpicture}[ch1fig, node distance=5mm and 7mm]
 \node[ch1box] (fo) {First-order\\$\mathcal{M}_\omega$};
 \node[ch1box, right=7mm of fo] (el) {$\mathcal{W}_{E,\omega}$};
 \node[ch1box, below=5mm of fo] (mag) {$\mathcal{W}_{H,\omega}$};
 \node[ch1box, right=7mm of el] (sym) {Symbol $k(\mathbf{r})$};
 \draw[ch1flow] (fo.east) -- node[above,font=\scriptsize]{eliminate $\H$} (el.west);
 \draw[ch1flow] (fo.south) -- node[left,font=\scriptsize]{eliminate $\E$} (mag.north);
 \draw[ch1flow] (el.east) -- (el.east |- sym.west);
 \draw[ch1flow] (el.east |- sym.west) -- (sym.west);
\end{tikzpicture}%
}
\caption{Frequency-domain reduction ladder: first-order \(\mathcal{M}_\omega\)
yields electric and magnetic Helmholtz operators when partner fields are eliminated
under stated domain and gauge constraints.}
\label{fig:ch2.frequency-operator-ladder}
\end{figure}
Single-field Helmholtz forms equivalent to
Eqs.~\eqref{eq:ch1.vacuum-wave-source-driven}--\eqref{eq:ch1.vacuum-curl-second-order}
were established in \ref{ch:01} and are; elimination
steps are not repeated \cite{stratton1941}. The \ref{ch:02} electric-wave
operator Eq.~\eqref{eq:ch2.electric-wave-operator} is the vacuum electric Helmholtz
map on \(\mathcal{D}(\mathcal{W}_{E,\omega})\subset\HCurl(\Omega)\):
\begin{equation}
\label{eq:ch2.electric-helmholtz-domain}
\mathcal{D}(\mathcal{W}_{E,\omega})
=
\big\{\E\in\HCurl(\Omega):
\mathcal{W}_{E,\omega}\E\in L^{2}(\Omega)^{3}\big\}.
\end{equation}
The magnetic counterpart is
\begin{equation}
\label{eq:ch2.magnetic-wave-operator}
\mathcal{W}_{H,\omega}\H\equiv
\curl\!\left(\varepsilon_0^{-1}\curl\H\right)-\omega^2\mu_0\H
= \jj\omega\mu_0\Jimp-\varepsilon_0^{-1}\curl\Jimp,
\end{equation}
on \(\mathcal{D}(\mathcal{W}_{H,\omega})\subset\HCurl(\Omega)\) analogously
\cite{stratton1941}. Equations
\eqref{eq:ch2.electric-wave-operator} and~\eqref{eq:ch2.magnetic-wave-operator}
are dual single-field symbols; using one without the other is valid only when the
eliminated partner is recoverable from curl relations without ambiguity. In piecewise smooth media with \(\varepsilon(\mathbf{r})\), \(\mu(\mathbf{r})\),
\begin{equation}
\label{eq:ch2.material-helmholtz-electric}
\mathcal{W}_{E,\omega}^{(\varepsilon,\mu)}\E\equiv
\curl\!\left(\mu^{-1}\curl\E\right)-\omega^2\varepsilon\E
= -\jj\omega\mu\Jimp+\mu^{-1}\curl\mathbf{J}^{\mathrm{(mag)}}_{\mathrm{eff}},
\end{equation}
\begin{equation}
\label{eq:ch2.material-helmholtz-magnetic}
\mathcal{W}_{H,\omega}^{(\varepsilon,\mu)}\H\equiv
\curl\!\left(\varepsilon^{-1}\curl\H\right)-\omega^2\mu\H
= \jj\omega\mu\Jimp-\varepsilon^{-1}\curl\Jimp,
\end{equation}
with interface jumps on \(\mathcal{D}(\mathcal{W}_{\cdot,\omega}^{(\varepsilon,\mu})\)
fixed by Eq.~\eqref{eq:ch1.interface-conditions} \cite{stratton1941}. Material
weights modify symbols but not the first-order coupling pattern of
Figure~\ref{fig:ch2.first-order-curl-coupling}; multi-region Green operators in
Section~\ref{sec:ch234} invert Eqs.~\eqref{eq:ch2.material-helmholtz-electric}--\eqref{eq:ch2.material-helmholtz-magnetic}
on composite domains \cite{harrington2001}.
The local Helmholtz symbol is
\begin{equation}
\label{eq:ch2.helmholtz-symbol}
k(\mathbf{r})=\omega\sqrt{\mu(\mathbf{r})\varepsilon(\mathbf{r})},
\qquad
k_0=\omega\sqrt{\mu_0\varepsilon_0},
\end{equation}
from Eq.~\eqref{eq:ch1.vacuum-impedance-wavenumber} in uniform vacuum
\cite{stratton1941}. In source-free regions,
Eqs.~\eqref{eq:ch1.vacuum-helmholtz-pair} apply with
\(k(\mathbf{r})\) replacing \(k_0\). The symbol \(k(\mathbf{r})\) labels propagating
versus evanescent wavenumber locally; it does not replace Sommerfeld outgoing selection on
open ranges \cite{stratton1941}. Evanescent sheets with large $|\E|$ can
satisfy the symbol while exporting no watts through $\Gamma_R$ \cite{chu1948}.
Single-field reduction is admissible only under explicit constraints:
\begin{equation}
\label{eq:ch2.single-field-reduction-validity}
\text{eliminate }\H
\ \Longrightarrow\
\Jimp\in\HCurl(\Omega),\ \rho\ \text{compatible},\ 
\mathcal{B}_\Gamma\ \text{expressible in }\E\text{-traces alone},
\end{equation}
with the magnetic elimination requiring the parallel conditions on \(\H\)-traces
\cite{stratton1941}. Failure of Eq.~\eqref{eq:ch2.single-field-reduction-validity}
---for example impedance boundaries coupling both tangential components---forces
retention of the first-order \(\mathcal{L}_\omega\) system
Eq.~\eqref{eq:ch2.maxwell-boundary-composite} \cite{harrington2001}.
Forcing maps between impressed data and Helmholtz right-hand sides are linear:
\begin{equation}
\label{eq:ch2.helmholtz-forcing-map}
\mathbf{f}_{E,\omega}
=
-\jj\omega\mu\Jimp+\mu^{-1}\curl\mathbf{J}^{\mathrm{(mag)}}_{\mathrm{eff}},
\qquad
\mathcal{W}_{E,\omega}^{(\varepsilon,\mu)}\E=\mathbf{f}_{E,\omega},
\end{equation}
with \(\mathbf{f}_{E,\omega}\in L^{2}(\Omega)^{3}\) when
\(\Jimp\in\mathcal{J}_\Omega\) from Eq.~\eqref{eq:ch2.source-space-refined}
\cite{stratton1941}. The harmonic solution operator on the electric Helmholtz branch
is
\begin{equation}
\label{eq:ch2.helmholtz-solution-operator}
\mathcal{G}_{E,\omega}:\mathbf{f}_{E,\omega}\longmapsto\E,
\qquad
\mathcal{W}_{E,\omega}^{(\varepsilon,\mu)}\mathcal{G}_{E,\omega}=\mathrm{id}
\end{equation}
on the physical resolvent branch selected by \(\mathcal{B}_\Gamma\); explicit
\(\mathcal{G}_{E,\omega}\) construction is developed in Section~\ref{sec:ch23}
\cite{harrington2001}. Composition with the full map yields
\begin{equation}
\label{eq:ch2.frequency-solution-composition}
\mathcal{S}_\omega[\Jimp]
=
\big(\mathcal{G}_{E,\omega}\mathbf{f}_{E,\omega},\,
-\frac{1}{\jj\omega\mu}\curl\E\big)
\quad\text{when }\H\ \text{recovered from Amp\`ere row},
\end{equation}
linking Eq.~\eqref{eq:ch2.maxwell-solution-map} to Helmholtz inversion without
re-displaying Eq.~\eqref{eq:ch2.solution-map-resolvent} \cite{stratton1941}.
Spectral viewpoint note: on bounded domains with discrete cavity spectra,
\(\mathcal{W}_{E,\omega}^{(\varepsilon,\mu)}\) admits eigenpairs
\((\lambda_n,\mathbf{e}_n)\) accumulating at infinity because
\(\HCurl(\Omega)\hookrightarrow L^{2}(\Omega)^{3}\) compactly
(Eq.~\eqref{eq:ch2.embedding-chain}) \cite{stratton1941}. On exterior domains,
continuous spectrum replaces discrete eigenvalues; Section~\ref{sec:ch263} develops
Fredholm and spectral consequences. Vector spherical harmonic separation of
eigenmodes is deferred to Section~3.3 \cite{harrington2001}.
\begin{figure}[t]
\centering
\ChOneFigBlock{%
\begin{tikzpicture}[ch1fig, node distance=5mm and 7mm]
 \node[ch1box] (l) {$\mathcal{L}_\omega$\\first-order};
 \node[ch1box, right=16mm of l] (w) {$\mathcal{W}_{E,\omega}$\\curl--curl};
 \node[ch1box, below=5mm of w] (b) {Boundary row\\$\mathcal{B}_\Gamma$};
 \draw[ch1flow] (l.east) -- node[above,font=\scriptsize]{eliminate} (w.west);
 \draw[ch1flow, dashed] (b.east) -- (l.west);
 \draw[ch1flow, dashed] (b.south) -- (w.north);
\end{tikzpicture}%
}
\caption{First-order \(\mathcal{L}_\omega\) retains coupled boundary rows; single-field
\(\mathcal{W}_{E,\omega}\) is valid only when \(\mathcal{B}_\Gamma\) can be
projected without losing impedance coupling.}
\label{fig:ch2.first-order-vs-helmholtz}
\end{figure}
Impedance and outgoing radiation conditions modify both volume and trace operators
simultaneously \cite{stratton1941}. Weak form of the electric Helmholtz operator uses the same sesquilinear structure as
Eq.~\eqref{eq:ch2.maxwell-weak-form} restricted to \(\E\)-only test spaces when
reduction is valid \cite{stratton1941}. Coercivity on the quotient by cavity modes
from Eq.~\eqref{eq:ch2.coercivity-quotient} applies to
\(\mathcal{W}_{E,\omega}^{(\varepsilon,\mu)}\) on bounded domains with passive
\(\mathcal{B}_\Gamma\); Subsection~\ref{sec:ch214} develops the self-adjoint realizations
\cite{harrington2001}.
Subsection~\ref{sec:ch222} has packaged harmonic Maxwell residuals, electric and magnetic
Helmholtz operators, and the conditions under which single-field reduction is admissible.
Subsection~\ref{sec:ch223} proves that \(\mathcal{S}_\omega\) is the frequency-domain
image of a causal time-domain map when impressed data, media, and boundary class are held
fixed \cite{stratton1941}.
\subsection{Causality and time-frequency equivalence}
\label{sec:ch223}
% HARRINGTON-OPEN:sec:ch223
A standard-gain horn on a compact-range turntable fixes the laboratory scene for `Causality and time-frequency equivalence'': frequency, probe standoff, and wattmeter closure on $\Gamma_R$ precede the formal displays below \cite{hansen1988}. Maxwell's retarded export to $\Gamma_R$ remains the governing constraint \cite{maxwell1873}.
The physical problem is retarded response: no field at $(\mathbf{r},t)$ may depend on source
events that have not yet had time to propagate from $\mathbf{r}'$ to $\mathbf{r}$
\cite{maxwell1873}. On the range, this is the difference between a
measured impulse response that vanishes before the light-travel delay and an unphysical
advanced template that leads the feed. A pulsed compact-range measurement makes the
distinction visible: energy arrives only after transit time $|\mathbf{r}-\mathbf{r}'|/c_0$;
a single-frequency phasor can hide violations that broadband synthesis would expose
\cite{harrington2001}. Causality is therefore the principle that links time-domain evolution
to harmonic analysis before any Laplace or resolvent notation is introduced
\cite{stratton1941}. A time-domain operator is causal when no response can occur before the
corresponding excitation \cite{stratton1941}. In electromagnetics this requirement appears
through retarded propagation: disturbances generated at one location influence another only
after the finite transit time imposed by the speed of light \cite{sommerfeld1949}.
Subsection~\ref{sec:ch223} supplies the definitive proof that $\mathcal{S}_\omega$ and
$\mathcal{R}_\omega$ are the frequency-domain images of causal time-domain maps when
impressed data, media, and boundary class are held fixed \cite{stratton1941}. The retarded
Eq.~\eqref{eq:ch2.causality-retarded} and symbolic relation
Eq.~\eqref{eq:ch2.time-to-frequency-map} are upgraded to an equivalence theorem;
Green-function construction remains in Section~\ref{sec:ch23} \cite{stratton1941}.
Causal support. A time-domain operator \(\mathcal{K}\) is \emph{causal} if
% PHYSICAL-MEMORY:sec:ch223
\begin{physicalvignette}[Anechoic turntable with azimuth-over-elevation stage]
Consider Anechoic turntable with azimuth-over-elevation stage. Before any coefficient table is printed, technicians record frequency, cable loss, and the mechanical outline that fixes electrical size. The first walk-through of the range is always tactile: flange bolts, probe standoff, turntable boresight, and the wattmeter reference on the receive chain \cite{hansen1988}.
On this hardware the near zone is visually impressive---standing reactive structure, rapid phase rotation, sometimes a visible VSWR ripple at the feed---yet the export question is measured only after outgoing selection on a large sphere \cite{harrington2001}. What the camera sees in the chamber is not automatically what the far-zone pattern claims \cite{stratton1941}.
Engineers therefore learn to separate three readings: what the feed can synthesize on declared support, what propagating content survives on the radiative range, and what a finite instrument can recover at declared tolerance \cite{devaney2004}. Confusing those readings produces the rank gaps and watt failures that later sections quantify \cite{pal2026}.
Rayleigh and Stratton treated such ranges as boundary-value laboratories long before digital near-field transforms became routine \cite{stratton1941}. Harrington's method-of-moments culture reinforced the same discipline: metal outline first, harmonic headline second \cite{harrington2001}.
Subsection~\ref{sec:ch223} states the formal export condition that this range scene presupposes \cite{stratton1941}.
\end{physicalvignette}
\begin{equation}
\label{eq:ch2.causal-operator-support}
\mathcal{K}[f](t)=0\ \text{for}\ t<t_0
\quad\text{whenever}\ 
\operatorname{supp}_t f\subseteq[t_0,\infty),
\end{equation}
for admissible test functions \(f\) \cite{stratton1941}. Retarded propagation forbids
response before transit time $|\mathbf{r}-\mathbf{r}'|/c$; frequency-domain Green
functions used on ranges are analytic continuations of that support, while broadband
pulses expose violations that single-frequency phasors can hide \cite{sommerfeld1949}.
The retarded dyadic kernel
in Eq.~\eqref{eq:ch2.retarded-convolution-form} satisfies
\begin{equation}
\label{eq:ch2.retarded-kernel-support}
\mathcal{G}(\mathbf{r},\mathbf{r}',t-t')=\mathbf{0}\quad\text{for}\ t-t'<0,
\end{equation}
which enforces Eq.~\eqref{eq:ch2.causal-support-constraint} on \(\E\)
\cite{stratton1941}. Advanced kernels with
\(\mathcal{G}(\mathbf{r},\mathbf{r}',t-t')\neq\mathbf{0}\) for \(t-t'<0\) violate
Eq.~\eqref{eq:ch2.causal-operator-support} and are excluded from the equivalence
class treated here.
\begin{figure}[t]
\centering
\ChOneFigBlock{%
\begin{tikzpicture}[ch1fig, node distance=5mm and 7mm]
 \node[ch1box] (td) {$\mathcal{S}_t$\\time domain};
 \node[ch1box, right=7mm of td] (lt) {Laplace\\$s=-\jj\omega$};
 \node[ch1box, right=7mm of lt] (fd) {$\mathcal{S}_\omega$\\harmonic};
 \draw[ch1flow] (td) -- (lt) -- (fd);
\end{tikzpicture}%
}
\caption{Causality link: causal \(\mathcal{S}_t\) maps to harmonic
\(\mathcal{S}_\omega\) under Laplace/Fourier evaluation on the locked phasor axis
\(s=-\jj\omega\), \(\omega>0\).}
\label{fig:ch2.laplace-resolvent-linkage}
\end{figure}
Let
\(\mathbf{j}_{\mathrm{imp}}(\mathbf{r},t)\) be an impressed current with compact
support in \(\Omega_s\times[0,\infty)\) and let \(\mathbf{U}_0=\mathbf{0}\) for
simplicity. The Laplace transform in \(s\) with \(\Re\{s\}\) sufficiently large is
\begin{equation}
\label{eq:ch2.laplace-source-transform}
\widetilde{\mathbf{j}}(\mathbf{r},s)
=
\int_0^\infty \mathbf{j}_{\mathrm{imp}}(\mathbf{r},t)\,e^{-st}\,dt.
\end{equation}
On the phasor axis \(s=-\jj\omega\) with \(\omega>0\),
\begin{equation}
\label{eq:ch2.fourier-phasor-identification}
\widetilde{\mathbf{j}}(\mathbf{r},-\jj\omega)
=
\int_0^\infty \mathbf{j}_{\mathrm{imp}}(\mathbf{r},t)\,e^{\jj\omega t}\,dt,
\end{equation}
which is the Fourier transform restricted to causal support
\cite{stratton1941}. For \(\mathbf{j}_{\mathrm{imp}}(\mathbf{r},t)=\Re\{\underline{\mathbf{j}}(\mathbf{r})e^{-\jj\omega_0 t}\}\)
with fixed \(\omega_0>0\), Eq.~\eqref{eq:ch2.fourier-phasor-identification}
selects the phasor \(\underline{\mathbf{j}}\) under the locked convention
\(\exp(-\jj\omega t)\) \cite{stratton1941}. Applying the Laplace transform to Eq.~\eqref{eq:ch2.time-ibvp-package} with fixed
\(\mathcal{B}_\Gamma\) and homogeneous initial data yields, formally,
\begin{equation}
\label{eq:ch2.laplace-maxwell-resolvent}
\widetilde{\mathbf{U}}(\mathbf{r},s)
=
\mathcal{R}(s)\,\widetilde{\mathbf{F}}(s),
\qquad
\mathcal{R}(s)=\mathcal{L}_s^{-1},
\end{equation}
where \(\mathcal{L}_s\) is the Laplace-domain Maxwell operator obtained from
\(\mathcal{A}\) by Eq.~\eqref{eq:ch2.frequency-substitution-rule} with
\(\partial_t\to s\) \cite{stratton1941}. On the axis \(s=-\jj\omega\),
\(\mathcal{L}_s=\mathcal{L}_\omega\) from Eq.~\eqref{eq:ch2.composite-operator-domain},
so
\begin{equation}
\label{eq:ch2.resolvent-green-equivalence}
\mathcal{R}(-\jj\omega)=\mathcal{L}_\omega^{-1}
\quad\text{on the outgoing/resolvent branch,}
\end{equation}
and the retarded convolution Eq.~\eqref{eq:ch2.retarded-convolution-form} is the
inverse Laplace transform of Eq.~\eqref{eq:ch2.laplace-maxwell-resolvent} when
\(\mathcal{G}\) is the causal Green operator constructed in Section~\ref{sec:ch23}
\cite{stratton1941}. Equation~\eqref{eq:ch2.resolvent-green-equivalence}
is the operator form of Eq.~\eqref{eq:ch2.causality-retarded}: harmonic
inversion is not an independent postulate but the \(s=-\jj\omega\) evaluation of a
causal resolvent.
The equivalence theorem requires explicit assumptions:
\begin{equation}
\label{eq:ch2.equivalence-assumptions}
\begin{aligned}
&\text{(i)}\ \mathcal{B}_\Gamma\ \text{and}\ \varepsilon,\mu\ \text{fixed in time},\\
&\text{(ii)}\ \mathbf{j}_{\mathrm{imp}}\ \text{causal with}\ \operatorname{supp}_t\subseteq[0,\infty),\\
&\text{(iii)}\ \text{outgoing/resolvent branch selected for exterior domains},\\
&\text{(iv)}\ \mathbf{U}_0=\mathbf{0}\ \text{or compatible with harmonic steady limit}.
\end{aligned}
\end{equation}
Under Eqs.~\eqref{eq:ch2.equivalence-assumptions},
\begin{equation}
\label{eq:ch2.harmonic-equivalence-theorem}
\mathcal{S}_\omega[\underline{\mathbf{j}}]
=
\eval{\widetilde{\mathbf{U}}(\mathbf{r},s)}_{s=-\jj\omega}
=
\eval{\mathcal{L}_\omega^{-1}\widetilde{\mathbf{F}}(s)}_{s=-\jj\omega},
\end{equation}
where \(\underline{\mathbf{j}}\) is the phasor of
\(\mathbf{j}_{\mathrm{imp}}(\mathbf{r},t)=\Re\{\underline{\mathbf{j}}(\mathbf{r})e^{-\jj\omega t}\}\)
and \(\widetilde{\mathbf{F}}\) is the Laplace forcing induced by
\(\widetilde{\mathbf{j}}\) \cite{stratton1941}. Equation
\eqref{eq:ch2.harmonic-equivalence-theorem} is the formal statement that
frequency-domain solvers implement causal time-domain maps when
Eqs.~\eqref{eq:ch2.equivalence-assumptions} hold \cite{harrington2001}. For
single-frequency excitation with damping or outgoing radiation,
\begin{equation}
\label{eq:ch2.steady-state-harmonic-limit}
\lim_{t\to\infty}\mathbf{U}(t)
=
\Re\big\{\mathcal{S}_\omega[\underline{\mathbf{j}}]\,e^{-\jj\omega t}\big\},
\end{equation}
which justifies steady-state phasor measurements after switch-on transients decay
\cite{stratton1941}. Because $\Pi_{\mathrm{rad}}$ is linear and commutes with harmonic
extraction on the outgoing subspace,
\begin{equation}
\label{eq:ch2.radiation-time-frequency-equiv}
\mathcal{R}_\omega[\underline{\mathbf{j}}]
=
\Pi_{\mathrm{rad}}\,\mathcal{S}_\omega[\underline{\mathbf{j}}],
\end{equation}
with transport and regime qualifications from
Eq.~\eqref{eq:ch2.transport-qualified-radiation} unchanged
\cite{stratton1941}. Frequency-domain codes that invert $\mathcal{L}_\omega$ without
outgoing or passive completion implement a different resolvent branch than the causal
retarded $\mathcal{G}$ of Eq.~\eqref{eq:ch2.retarded-convolution-form}
\cite{harrington2001}. Advanced Green functions satisfy
\begin{equation}
\label{eq:ch2.advanced-kernel-violation}
\mathcal{G}_{\mathrm{adv}}(\mathbf{r},\mathbf{r}',t-t')\neq\mathbf{0}
\ \text{for some}\ t-t'<0,
\end{equation}
and cannot satisfy Eq.~\eqref{eq:ch2.harmonic-equivalence-theorem} for any causal
$\mathbf{j}_{\mathrm{imp}}$ \cite{stratton1941}. Broadband synthesis requires
superposition over $\omega$, and time-varying $\mathcal{B}_\Gamma(t)$ breaks assumption
(i) in Eq.~\eqref{eq:ch2.equivalence-assumptions} \cite{stratton1941}.
Subsections~\ref{sec:ch221}--\ref{sec:ch223} fixed evolution, harmonic operators, and
causality. Subsection~\ref{sec:ch224} supplies spectral partitions into propagating,
evanescent, and non-propagating subspaces in operator form \cite{stratton1941}.
Reactive versus radiative regime criteria from Section~\ref{sec:ch01.2.4} are invoked by
reference; the outgoing selector $\Pi_{\mathrm{rad}}$ of
Eq.~\eqref{eq:ch2.radiation-operator-def} is refined spectrally, not replaced
\cite{poynting1884}.
The physical problem is spectral transport: a near-field scanner records both traveling and
decaying longitudinal content, yet only outgoing propagating modes may enter far-zone power
reports \cite{hansen1988}. Subsection~\ref{sec:ch224} defines propagating, evanescent, and non-propagating spectral
subspaces and refines \(\Pi_{\mathrm{rad}}=\mathcal{O}_{\mathrm{out}}\circ\Pi_{\mathrm{prop}}\)
so that only outgoing propagating content qualifies for far-zone transport reports
\cite{stratton1941}. Recall that $\mathcal{S}_\omega:\mathcal{J}_\Omega\to\mathcal{F}_\Omega$ produces the
full harmonic field pair, $\mathcal{R}_\omega=\Pi_{\mathrm{rad}}\circ\mathcal{S}_\omega$
retains only radiative content, and $\Pi_{\mathrm{rad}}$ acts after propagating selection
\cite{harrington2001}. A near-field scanner routinely measures both propagating and evanescent spectral content.
The distinction becomes visible when reconstructed fields are propagated away from the
measurement surface: propagating components persist, whereas evanescent components decay
exponentially with distance \cite{hansen1988}. The spectral decomposition
developed below states that experimentally familiar observation in spectral language.
In a homogeneous region with symbol \(k_0\) from Eq.~\eqref{eq:ch2.helmholtz-symbol},
plane-wave trials decompose the field into lateral and longitudinal wavenumber components.
The dispersion relation below is the local criterion that
separates traveling from decaying longitudinal behavior: when lateral wavenumber
\(k_t=\sqrt{k_x^2+k_y^2}\) exceeds \(k_0\), no real propagating longitudinal root exists
and the mode becomes evanescent \cite{stratton1941}. Experimentally one would observe that
near-field scan content with \(k_t>k_0\) dies away within a penetration depth
\(\delta_{\mathrm{ev}}=1/\alpha(k_t)\) as the field is propagated outward
\cite{hansen1988}. Plane-wave trials \(\mathbf{u}_{\mathbf{k}}(\mathbf{r})=\mathbf{u}_0\,e^{-\jj\mathbf{k}\cdot\mathbf{r}}\)
satisfy the Helmholtz condition when
\begin{equation}
\label{eq:ch2.dispersion-relation}
|\mathbf{k}|=k_0,\qquad
k_z^2=k_0^2-k_x^2-k_y^2,
\end{equation}
with lateral wavenumber \(k_t=\sqrt{k_x^2+k_y^2}\) \cite{stratton1941}. Equation
\eqref{eq:ch2.dispersion-relation} splits spectral content by longitudinal behavior:
\begin{equation}
\label{eq:ch2.longitudinal-wavenumber-split}
k_z=
\begin{cases}
\pm\sqrt{k_0^2-k_t^2}, & k_t\le k_0\ \text{(propagating)},\\[4pt]
-\jj\,\alpha(k_t), & k_t>k_0\ \text{(evanescent)},\ \alpha>0,
\end{cases}
\end{equation}
where the evanescent root is chosen with decay in the \(+z\) half-space
\cite{stratton1941}. For \(k_t>k_0\) the decay rate is
\begin{equation}
\label{eq:ch2.evanescent-decay-rate}
\alpha(k_t)=\sqrt{k_t^2-k_0^2},
\end{equation}
so a pure evanescent trial on \(z>0\) scales as
\(e^{-\alpha(k_t)z}\) with penetration depth \(\delta_{\mathrm{ev}}=1/\alpha(k_t)\)
\cite{stratton1941}. Propagating modes carry real \(k_z=\pm\sqrt{k_0^2-k_t^2}\) and
oscillatory phase along the longitudinal axis; evanescent modes carry imaginary
\(k_z\) and store reactive energy in the near zone without far-zone power export
\cite{harrington2001}. At fixed \((k_x,k_y)\), the Helmholtz symbol
\(\mathcal{L}_\omega\) from Subsection~\ref{sec:ch222} acts on each longitudinal
branch as a scalar multiplier in the plane-wave basis, which is why spectral
partitions commute with harmonic extraction once outgoing radiation condition is fixed on the
propagating sheet \cite{stratton1941}.
Transverse polarization completes the mode description. For a trial electric field
\(\widetilde{\E}(k_x,k_y)\,e^{-\jj(k_x x+k_y y+k_z z)}\), source-free admissibility
requires \(\mathbf{k}\cdot\widetilde{\E}=0\) in the homogeneous region
\cite{stratton1941}. Decompose \(\mathbf{k}=\mathbf{k}_t+k_z\hat{\mathbf{z}}\)
with \(\mathbf{k}_t=k_x\hat{\mathbf{x}}+k_y\hat{\mathbf{y}}\). Then
\begin{equation}
\label{eq:ch2.transverse-polarization-constraint}
k_x\widetilde{E}_x+k_y\widetilde{E}_y+k_z\widetilde{E}_z=0,
\end{equation}
which leaves two independent transverse amplitudes on each branch of
Eq.~\eqref{eq:ch2.longitudinal-wavenumber-split}. On the propagating sheet,
Eq.~\eqref{eq:ch2.transverse-polarization-constraint} supports traveling power
along \(\hat{\mathbf{z}}\) when outgoing \(k_z\) is selected; on the evanescent sheet
the same constraint forces predominantly tangential storage with longitudinal
\(E_z\) tied to \(\alpha(k_t)\) \cite{stratton1941}. Subsection~\ref{sec:ch224}
does not classify TE/TM guide nomenclature---that labeling is deferred to
waveguide sections---but Eq.~\eqref{eq:ch2.transverse-polarization-constraint}
is the local polarization criterion paired with the dispersion criterion of
Eq.~\eqref{eq:ch2.dispersion-relation} \cite{harrington2001}.
\begin{figure}[t]
\centering
\ChOneFigBlock{%
\begin{tikzpicture}[ch1fig, node distance=5mm and 7mm]
 \node[ch1box] (kt) {$k_t$};
 \node[ch1box, right=7mm of kt] (prop) {$k_t\le k_0$\\propagating};
 \node[ch1box, below=5mm of kt] (ev) {$k_t>k_0$\\evanescent};
 \draw[ch1flow] (kt.east) -- (kt.east |- prop.west);
 \draw[ch1flow] (kt.east |- prop.west) -- (prop.west);
 \draw[ch1flow] (kt.south) -- (kt.south |- ev.north);
 \draw[ch1flow] (kt.south |- ev.north) -- (ev.north);
\end{tikzpicture}%
}
\caption{Dispersion split Eq.~\eqref{eq:ch2.longitudinal-wavenumber-split}: lateral
wavenumber \(k_t\) selects propagating versus evanescent longitudinal behavior at
fixed \(k_0\).}
\label{fig:ch2.propagating-evanescent-split}
\end{figure}
Angular-spectrum representations in
Section~\ref{sec:ch29} expand fields in these modes; here the subspaces are defined
as solution domains \cite{stratton1941}. Let \(\mathcal{F}_{\mathrm{prop}}\subset\mathcal{F}_\Omega\) be the completion of
finite superpositions of propagating plane waves with outgoing \(k_z\) selected by
exterior completion, let \(\mathcal{F}_{\mathrm{ev}}\subset\mathcal{F}_\Omega\) be the
closure of evanescent superpositions, and let
\(\mathcal{F}_{\mathrm{non}}\subset\mathcal{F}_\Omega\) collect modes that are
neither outgoing-propagating nor evanescent-decaying in the declared half-space---for
example cutoff guide modes and purely reactive standing content tied to terminals
\cite{stratton1941}. 
\phantomsection\label{sec:ch224}
Orthogonal spectral projectors (formal) satisfy
% PHYSICAL-MEMORY:sec:ch224
\begin{physicalvignette}[Radar feed horn illuminating a corner reflector]
Consider Radar feed horn illuminating a corner reflector. Before any coefficient table is printed, technicians record frequency, cable loss, and the mechanical outline that fixes electrical size. The first walk-through of the range is always tactile: flange bolts, probe standoff, turntable boresight, and the wattmeter reference on the receive chain \cite{hansen1988}.
On this hardware the near zone is visually impressive---standing reactive structure, rapid phase rotation, sometimes a visible VSWR ripple at the feed---yet the export question is measured only after outgoing selection on a large sphere \cite{harrington2001}. What the camera sees in the chamber is not automatically what the far-zone pattern claims \cite{stratton1941}.
Engineers therefore learn to separate three readings: what the feed can synthesize on declared support, what propagating content survives on the radiative range, and what a finite instrument can recover at declared tolerance \cite{devaney2004}. Confusing those readings produces the rank gaps and watt failures that later sections quantify \cite{pal2026}.
Rayleigh and Stratton treated such ranges as boundary-value laboratories long before digital near-field transforms became routine \cite{stratton1941}. Harrington's method-of-moments culture reinforced the same discipline: metal outline first, harmonic headline second \cite{harrington2001}.
Subsection~\ref{sec:ch224} states the formal export condition that this range scene presupposes \cite{stratton1941}.
\end{physicalvignette}
\begin{equation}
\label{eq:ch2.spectral-projectors}
\Pi_{\mathrm{prop}}+\Pi_{\mathrm{ev}}+\Pi_{\mathrm{non}}=\mathrm{id}
\quad\text{on}\ \widehat{\mathcal{F}}_\Omega,
\end{equation}
with \(\widehat{\mathcal{F}}_\Omega\) the spectral completion of admissible states
\cite{stratton1941}. Idempotency and mutual annihilation read
\begin{equation}
\label{eq:ch2.spectral-projector-algebra}
\Pi_{\mathrm{prop}}^2=\Pi_{\mathrm{prop}},\quad
\Pi_{\mathrm{ev}}^2=\Pi_{\mathrm{ev}},\quad
\Pi_{\mathrm{non}}^2=\Pi_{\mathrm{non}},\quad
\Pi_{\mathrm{prop}}\Pi_{\mathrm{ev}}=\Pi_{\mathrm{ev}}\Pi_{\mathrm{prop}}=0,
\end{equation}
with analogous zero products for \(\Pi_{\mathrm{non}}\) \cite{stratton1941}. On
\(\mathcal{D}(\mathcal{L}_\omega)\) from Subsection~\ref{sec:ch222}, each projector
block is an invariant subspace of the scalar Helmholtz symbol in the plane-wave
basis,
\medskip
Equations~\eqref{eq:ch2.spectral-projectors}--\eqref{eq:ch2.spectral-projector-algebra} are
\emph{formal} on the spectral completion \(\widehat{\mathcal{F}}_\Omega\): idempotency and
orthogonality are justified in the rigged Hilbert space treatment of Chapter~3 (\emph{Angular Momentum Operators and SO(3)})
Subsections~3.1.3--3.5.2 \cite{stratton1941}. Until that completion,
projector algebra is used only on admissible finite superpositions and on
\(\mathcal{R}_\omega[\Jimp]\), where outgoing selection is already enforced
\cite{harrington2001}. \begin{equation}
\label{eq:ch2.helmholtz-subspace-commutation}
\mathcal{L}_\omega\big(\Pi_{\mathrm{prop}}\mathbf{u}\big)
=
\Pi_{\mathrm{prop}}\big(\mathcal{L}_\omega\mathbf{u}\big),
\end{equation}
and similarly for \(\Pi_{\mathrm{ev}}\) and \(\Pi_{\mathrm{non}}\) when boundary
closure does not mix sheets \cite{harrington2001}. Equation
\eqref{eq:ch2.helmholtz-subspace-commutation} is the operator reason harmonic
solvers can treat propagating and evanescent content separately before radiation
selection; Sommerfeld outgoing radiation condition in Section~\ref{sec:ch24} acts only on
\(\Pi_{\mathrm{prop}}\) \cite{stratton1941}.
The decomposition
\begin{equation}
\label{eq:ch2.propagating-evanescent-decomposition}
\mathcal{S}_\omega[\Jimp]
=
\Pi_{\mathrm{prop}}\mathcal{S}_\omega[\Jimp]
+\Pi_{\mathrm{ev}}\mathcal{S}_\omega[\Jimp]
+\Pi_{\mathrm{non}}\mathcal{S}_\omega[\Jimp]
\end{equation}
holds when the impressed spectrum excites all three classes. Only
\(\Pi_{\mathrm{prop}}\mathcal{S}_\omega[\Jimp]\) is eligible for far-zone radiation
content after outgoing radiation condition; evanescent and non-propagating components remain in
reactive storage exchange \cite{stratton1941}. Energy norms from
Eq.~\eqref{eq:ch2.field-energy-norm} decompose orthogonally in the spectral basis
when the impressed angular spectrum has sufficient decay in \((k_x,k_y)\)
\cite{stratton1941}.
\begin{equation}
\label{eq:ch2.rad-projection-composition}
\Pi_{\mathrm{rad}}
=
\mathcal{O}_{\mathrm{out}}\circ\Pi_{\mathrm{prop}},
\end{equation}
where \(\mathcal{O}_{\mathrm{out}}\) enforces outgoing selection on propagating
content compatible with Eq.~\eqref{eq:ch1.radiation-trace-condition}
\cite{stratton1941}. Equation~\eqref{eq:ch2.rad-projection-composition} refines
\(\Pi_{\mathrm{rad}}\): radiative subspace \(\mathcal{F}_{\mathrm{rad}}\) is a
subset of \(\Pi_{\mathrm{prop}}\mathcal{F}_\Omega\), not of all of
\(\mathcal{F}_\Omega\) \cite{harrington2001}. Evanescent and non-propagating
components lie in the null space of far-zone radiation:
\begin{equation}
\label{eq:ch2.rad-null-on-evanescent}
\Pi_{\mathrm{rad}}\Pi_{\mathrm{ev}}=0,\qquad
\Pi_{\mathrm{rad}}\Pi_{\mathrm{non}}=0,
\end{equation}
while \(\Pi_{\mathrm{rad}}\Pi_{\mathrm{prop}}=\Pi_{\mathrm{rad}}\) after outgoing
closure \cite{stratton1941}. Equations~\eqref{eq:ch2.rad-null-on-evanescent} and
\eqref{eq:ch2.rad-projection-composition} restate, in projector form, the regime
logic of Section~\ref{sec:ch01.2.4}: reactive near fields may be large even when
\(\Pi_{\mathrm{rad}}\mathcal{S}_\omega[\Jimp]=\mathbf{0}\) \cite{harrington2001}.
Causality from Subsection~\ref{sec:ch223} ensures retarded support but does not,
by itself, promote evanescent content into \(\mathcal{F}_{\mathrm{rad}}\).
\medskip
For all \(r>r_{\min}(a,\lambda_0)\) and every transport, flux, near-to-far, and angular-spectrum
statement in Sections~\ref{sec:ch25}--\ref{sec:ch210},
\begin{equation*}
\Pi_{\mathrm{rad}}=\mathcal{O}_{\mathrm{out}}\circ\Pi_{\mathrm{prop}}
\quad\text{as in Eq.~\eqref{eq:ch2.rad-projection-composition}.}
\end{equation*}
Earlier uses of \(\Pi_{\mathrm{rad}}\) in Section~\ref{sec:ch20} are to be read through this
refinement \cite{stratton1941}. Transport qualification from Section~\ref{sec:ch01.2.4} enters as an operator criterion:
\begin{equation}
\label{eq:ch2.regime-operator-screen}
\mathcal{D}_{J,\mathrm{reg}}=1
\ \Longrightarrow\
\text{transport reports use}\ 
\Pi_{\mathrm{prop}}\mathcal{R}_\omega[\Jimp],
\ \text{not}\ \Pi_{\mathrm{ev}}\mathcal{S}_\omega[\Jimp],
\end{equation}
with \(\mathcal{D}_{J,\mathrm{reg}}\) from
Eq.~\eqref{eq:ch1.regime-qualified-decision} \cite{stratton1941}. Equation
\eqref{eq:ch2.regime-operator-screen} is the operator restatement of
Eq.~\eqref{eq:ch2.transport-qualified-radiation}; regime indices \(\chi(r)\) and
\(\chi_\Omega\) diagnose when Eq.~\eqref{eq:ch2.regime-operator-screen} fails on a
measurement shell \cite{stratton1941}.
Waveguide and cutoff structures illustrate non-propagating subspaces. For a uniform
guide with lateral eigenvalue \(\kappa_n\), mode \(n\) propagates only when
\begin{equation}
\label{eq:ch2.cutoff-condition}
k_0^2>\kappa_n^2
\quad\Longleftrightarrow\quad
\omega>\omega_{c,n}=\frac{c_0\kappa_n}{\sqrt{\varepsilon_r\mu_r}},
\end{equation}
with \(\omega_{c,n}\) the cutoff angular frequency \cite{stratton1941}. Below
cutoff, the mode lies in \(\mathcal{F}_{\mathrm{ev}}\) or
\(\mathcal{F}_{\mathrm{non}}\) depending on boundary termination; applying far-zone
transport formulas to such fields violates Eq.~\eqref{eq:ch2.regime-operator-screen}
as in the TE\(_{10}\) illustration of Section~\ref{sec:ch01.1.2}
\cite{harrington2001}. In a rectangular guide of width \(a\), lateral eigenvalues
\(\kappa_{m,n}=\sqrt{(m\pi/a)^2+(n\pi/b)^2}\) replace \(k_t\) in
Eq.~\eqref{eq:ch2.cutoff-condition}; each mode index \((m,n)\) switches from
\(\mathcal{F}_{\mathrm{non}}\) to \(\mathcal{F}_{\mathrm{prop}}\) only when
\(\omega\) exceeds its cutoff \cite{stratton1941}. Near cutoff the propagating
group velocity vanishes and stored reactive energy dominates even though the mode
has crossed into \(\Pi_{\mathrm{prop}}\); Eq.~\eqref{eq:ch1.radiative-regime-condition}
remains the far-zone transport criterion, not Eq.~\eqref{eq:ch2.cutoff-condition} alone
\cite{harrington2001}.
Angular-spectrum expansion: expand \(\E(\mathbf{r})\) on \(z=\mathrm{const}\) as
\begin{equation}
\label{eq:ch2.angular-spectrum}
\E(\mathbf{r})
=
\int\!\!\int \widetilde{\E}(k_x,k_y)\,
e^{-\jj(k_x x+k_y y+k_z(k_t) z)}\,dk_x\,dk_y,
\end{equation}
with \(k_z(k_t)\) from Eq.~\eqref{eq:ch2.longitudinal-wavenumber-split}
\cite{stratton1941}. Equation~\eqref{eq:ch2.angular-spectrum} is the
\ref{ch:02} spectral template for Section~\ref{sec:ch29}; full radiation
operator action on angular spectra is developed there. Subsection~\ref{sec:ch224} only
defines the propagating/evanescent partition induced by
Eq.~\eqref{eq:ch2.dispersion-relation} \cite{harrington2001}.
\begin{figure}[t]
\centering
\ChOneFigBlock{%
\begin{tikzpicture}[ch1fig, node distance=5mm and 7mm]
 \node[ch1box] (s) {$\mathcal{S}_\omega[\Jimp]$};
 \node[ch1box, right=7mm of s] (pp) {$\Pi_{\mathrm{prop}}$};
 \node[ch1box, right=7mm of pp] (pr) {$\Pi_{\mathrm{rad}}$}; \node[ch1box, below=5mm of s] (pe) {$\Pi_{\mathrm{ev}}$};
 \node[ch1box, below=5mm of pe] (pn) {$\Pi_{\mathrm{non}}$};
 \draw[ch1flow] (s) -- (pp) -- (pr);
 \draw[ch1flow] (s.south) -- (s.south |- pe.north);
 \draw[ch1flow] (s.south |- pe.north) -- (pe.north);
 \draw[ch1flow] (s.east) -- (s.east |- pn.west);
 \draw[ch1flow] (s.east |- pn.west) -- (pn.west);
\end{tikzpicture}%
}
\caption{Spectral criterion: only \(\Pi_{\mathrm{prop}}\) content passes to
\(\Pi_{\mathrm{rad}}\); evanescent and non-propagating components remain reactive.}
\label{fig:ch2.spectral-radiation-criterion}
\end{figure}
Figure~\ref{fig:ch2.source-radiation-sequence}. Causality from
Subsection~\ref{sec:ch223} does not automatically place evanescent energy in
\(\mathcal{F}_{\mathrm{rad}}\); spectral and regime qualifications are additional
\cite{stratton1941}. Near-zone reactive fields decompose into large evanescent and non-propagating
fractions even when \(\mathcal{S}_\omega[\Jimp]\) is exact. As \(r\to\infty\),
\(\Pi_{\mathrm{prop}}\) content dominates \(P_{\mathrm{rad}}\) when
Eq.~\eqref{eq:ch1.radiative-regime-condition} holds \cite{stratton1941}. Operator
norms on subspaces satisfy
\begin{equation}
\label{eq:ch2.subspace-energy-inequality}
\|\Pi_{\mathrm{ev}}\mathcal{S}_\omega[\Jimp]\|_{\mathcal{U}}
\gg
\|\Pi_{\mathrm{prop}}\mathcal{S}_\omega[\Jimp]\|_{\mathcal{U}}
\quad\text{on near shells,}
\end{equation}
while the reverse inequality governs far-zone radiative observables
\cite{harrington2001}. Subsection~\ref{sec:ch224} does not compute
shell integrals; it supplies the spectral vocabulary for
Section~\ref{sec:ch252} and flux operators in Section~\ref{sec:ch25}.
With Subsections~\ref{sec:ch221}--\ref{sec:ch224} complete, Section~\ref{sec:ch22}
has connected time-domain evolution, harmonic Maxwell and Helmholtz operators,
causal equivalence, and propagating--evanescent spectral split on the function spaces
of Section~\ref{sec:ch21}. Green-function inversion in Section~\ref{sec:ch23}
realizes \(\mathcal{L}_\omega^{-1}\) and \(\mathcal{G}_{E,\omega}\) with retarded
causality; Sommerfeld selection in Section~\ref{sec:ch24} fixes outgoing branches for
\(\Pi_{\mathrm{prop}}\) and \(\Pi_{\mathrm{rad}}\); radiation integrals in
Section~\ref{sec:ch25} meter flux on the propagating subspace qualified by
Eq.~\eqref{eq:ch2.regime-operator-screen} \cite{stratton1941}.
\section{Green Functions and Electromagnetic Inversion}
\label{sec:ch23}
Time- and frequency-domain operators in Section~\ref{sec:ch22} specify how impressed currents enter Maxwell's equations. Section~\ref{sec:ch23} supplies Green-function propagation---the computational heart of method-of-moments and near-field scanners \cite{stratton1941}. Propagation must be understood before radiation is selected \cite{sommerfeld1949}.
% PHYSICAL-MEMORY:sec:ch23
\begin{physicalvignette}[Radio telescope dish with Cassegrain feed]
Consider Radio telescope dish with Cassegrain feed. Before any coefficient table is printed, technicians record frequency, cable loss, and the mechanical outline that fixes electrical size. The first walk-through of the range is always tactile: flange bolts, probe standoff, turntable boresight, and the wattmeter reference on the receive chain \cite{hansen1988}.
On this hardware the near zone is visually impressive---standing reactive structure, rapid phase rotation, sometimes a visible VSWR ripple at the feed---yet the export question is measured only after outgoing selection on a large sphere \cite{harrington2001}. What the camera sees in the chamber is not automatically what the far-zone pattern claims \cite{stratton1941}.
Engineers therefore learn to separate three readings: what the feed can synthesize on declared support, what propagating content survives on the radiative range, and what a finite instrument can recover at declared tolerance \cite{devaney2004}. Confusing those readings produces the rank gaps and watt failures that later sections quantify \cite{pal2026}.
Rayleigh and Stratton treated such ranges as boundary-value laboratories long before digital near-field transforms became routine \cite{stratton1941}. Harrington's method-of-moments culture reinforced the same discipline: metal outline first, harmonic headline second \cite{harrington2001}.
Subsection~\ref{sec:ch23} states the formal export condition that this range scene presupposes \cite{stratton1941}.
\end{physicalvignette}
Maxwell's equations couple impressed currents to fields throughout a domain; in practice one
needs the field at a specified observation point produced by sources elsewhere.
The Green function is the impulse response that reduces that coupling to a single
propagator integral rather than a fresh boundary-value solve at every new observation
coordinate \cite{stratton1941}. Before dyadic kernels are tabulated, fix the
lowest-dimensional laboratory calibration every range engineer uses: a short dipole or probe
feed at $k_0 a\ll 1$ on a turntable, with watt closure on a sphere $\Gamma_R$ and a
transverse $\sin\theta$ pattern that is invariant under rotation \cite{jackson1999}.
That Hertzian test is the physical image of $\mathcal{S}_\omega[\Jimp]$ when
$\Jimp$ collapses to a point dipole moment; Subsection~\ref{sec:ch23.signature-dipole}
closes the calculation, but the range observation precedes the dyad formulas
\cite{stratton1941}. Two laboratory questions follow. Given currents on a feed
or equivalent surface, how is the field computed at an observation point? Given field
samples on a scan surface or range contour, how may the equivalent source be recovered?
Method-of-moments assembly, near-field back-propagation, compact-range transformation, and
microwave imaging differ in discretization yet implement the same propagation law
\cite{harrington2001}. Scalar, vector, and dyadic kernels, singular
regularization, retarded causality, and inversion on composite domains are developed below;
outgoing selection at infinity belongs to Section~\ref{sec:ch24}. The export map below
$\mathcal{S}_\omega[\Jimp]$ realized through $G$ and $\overline{\mathbf{G}}_{E}$ is the
physical response of the medium to declared impressed currents; radiation export through
$\mathcal{R}_\omega$ and $\Pi_{\mathrm{rad}}$ is deferred until Section~\ref{sec:ch24} has fixed
the outgoing branch on which those operators act \cite{stratton1941}.
\subsection{Scalar, vector, and dyadic Green functions}
\label{sec:ch231}
Physically, $G$ and $\overline{\mathbf{G}}_{E}$ are the medium's response functions: each
records how a unit impulse at $\mathbf{r}'$ launches scalar or vector fields to $\mathbf{r}$ under
the same boundary completion as $\mathcal{L}_\omega$ \cite{stratton1941}. Every compact-range campaign begins with a dipole elementary check: a electrically small feed
must reproduce the transverse far-zone law before multipole catalogs are trusted on larger
apertures \cite{jackson1999}. That observation anchors the Green-function program
developed here---first the scalar monopole response, then vector potentials, then the
electric dyad that maps $\Jimp$ to $\E$ directly \cite{stratton1941}. The dyadic Green
function is the electromagnetic impulse response of the governing Maxwell
operator: the field at an observation point produced by an elemental current located
elsewhere. In uniform media it is also the continuous limit of one radiating cell in a
phased array and the local kernel that probe correction applies when near-field data are
translated to a reference surface \cite{stratton1941}. Admissible kernels are
retarded---field response follows source excitation only after finite propagation delay---so
localized impressed currents remain consistent with Maxwell causality
\cite{sommerfeld1949}. Scalar Helmholtz functions describe monopole-like propagation;
polarization coupling on vector feeds requires the full dyadic \cite{collin2001}.
% PHYSICAL-MEMORY:sec:ch231
\begin{physicalvignette}[Half-space ground plane with image currents]
Consider Half-space ground plane with image currents. Before any coefficient table is printed, technicians record frequency, cable loss, and the mechanical outline that fixes electrical size. The first walk-through of the range is always tactile: flange bolts, probe standoff, turntable boresight, and the wattmeter reference on the receive chain \cite{hansen1988}.
On this hardware the near zone is visually impressive---standing reactive structure, rapid phase rotation, sometimes a visible VSWR ripple at the feed---yet the export question is measured only after outgoing selection on a large sphere \cite{harrington2001}. What the camera sees in the chamber is not automatically what the far-zone pattern claims \cite{stratton1941}.
Engineers therefore learn to separate three readings: what the feed can synthesize on declared support, what propagating content survives on the radiative range, and what a finite instrument can recover at declared tolerance \cite{devaney2004}. Confusing those readings produces the rank gaps and watt failures that later sections quantify \cite{pal2026}.
Green functions propagate impressed currents to fields before radiation is selected . Near-field probes see this propagation; far-field patterns see its outgoing part .
Rayleigh and Stratton treated such ranges as boundary-value laboratories long before digital near-field transforms became routine \cite{stratton1941}. Harrington's method-of-moments culture reinforced the same discipline: metal outline first, harmonic headline second \cite{harrington2001}.
Subsection~\ref{sec:ch231} states the formal export condition that this range scene presupposes \cite{stratton1941}.
\end{physicalvignette}
The \emph{kernel formulas} developed below must be distinguished from
and the \emph{domain declarations} that make those formulas meaningful on a declared
boundary-value problem \cite{stratton1941}. Unless noted otherwise, \(\omega>0\) is fixed,
\(\Jimp\in\mathcal{J}_\Omega\) from Eq.~\eqref{eq:ch2.source-space-refined}, charge continuity
Eq.~\eqref{eq:ch1.charge-continuity-harmonic} holds, and \(\mathcal{B}_\Gamma\) is declared on
\(\Gamma=\partial\Omega\). In uniform vacuum regions the wavenumber is \(k_0\) from
Eq.~\eqref{eq:ch2.helmholtz-symbol}. Kernels are stated first in operator form on
\(\mathcal{D}(\mathcal{W}_{E,\omega})\subset\HCurl(\Omega)\) from
Eq.~\eqref{eq:ch2.electric-helmholtz-domain}; piecewise \(k(\mathbf{r})\) extensions use the
same template with patchwise coefficients \cite{stratton1941}. The harmonic solution operator
$\mathcal{G}_{E,\omega}$ introduced in Subsection~\ref{sec:ch222} is realized below by
convolution with $G$ and $\overline{\mathbf{G}}_{E}$; observation coordinate $\mathbf{r}$ and
source coordinate $\mathbf{r}'$ must remain distinct throughout every range transformation
\cite{harrington2001}.
Imagine exciting a uniform medium with the smallest source one can conceive: a unit
impulse concentrated at a single point $\mathbf{r}'$. The field that appears at every
other location $\mathbf{r}$ is the scalar Green function $G(\mathbf{r},\mathbf{r}')$---the
impulse response of the Helmholtz operator at fixed frequency
\cite{stratton1941}. Laboratory engineers never place a mathematical
$\delta$-function on a bench, yet every extended feed, patch, or volume current is assembled
from such elemental responses by superposition: each infinitesimal slice of $\Jimp(\mathbf{r}')$
radiates as if it were a point source scaled by the local source strength, and the total
field is the sum of those contributions weighted by propagation from $\mathbf{r}'$ to
$\mathbf{r}$ \cite{harrington2001}. Green functions are therefore the natural language of
propagation: they replace repeated boundary-value solves with a single propagator integral
that method-of-moments codes, near-field scanners, and imaging algorithms all discretize
\cite{stratton1941}.
The physical significance of Eq.~\eqref{eq:ch2.scalar-helmholtz-green-pde} is that it isolates
the response of the medium to the smallest source that can be imagined---a unit point excitation.
Once this elementary response is known, more complicated sources may be synthesized by
superposition \cite{jackson1999}. The scalar Green function $G(\mathbf{r},\mathbf{r}')$
inverts the Helmholtz operator $(\nabla^2+k_0^2)$ on scalar potentials; $\mathbf{r}$ denotes
the observation point and $\mathbf{r}'$ the source point, with $R=|\mathbf{r}-\mathbf{r}'|$
the geometric separation in metres. On a compact range, the scalar kernel is the monopole
reference against which probe calibration and layered-substrate corrections are measured
\cite{hansen1988}. It satisfies
\begin{equation}
\label{eq:ch2.scalar-helmholtz-green-pde}
\big(\nabla^{2}+k_0^{2}\big)\,G(\mathbf{r},\mathbf{r}')
=
-\delta(\mathbf{r}-\mathbf{r}'),
\end{equation}
with \(\delta\) the three-dimensional Dirac distribution and boundary conditions on
\(\Gamma\) inherited from \(\mathcal{B}_\Gamma\) \cite{stratton1941}. Equation
\eqref{eq:ch2.scalar-helmholtz-green-pde} is the local symbol of spectral inversion:
for fixed \(\mathbf{r}'\), \(G(\cdot,\mathbf{r}')\) is the response to a unit point
source at \(\mathbf{r}'\) under the same completion as \(\mathcal{L}_\omega\)
\cite{jackson1999}. Convolution with \(G\) therefore propagates an arbitrary scalar source
to the field \(\phi\) that satisfies the inhomogeneous Helmholtz equation on the same domain.
The associated solution operator is
\begin{equation}
\label{eq:ch2.scalar-green-solution-operator}
\phi(\mathbf{r})
=
\int_{\Omega} G(\mathbf{r},\mathbf{r}')\,q(\mathbf{r}')\,d^{3}r',
\qquad
\big(\nabla^{2}+k_0^{2}\big)\phi=q,
\end{equation}
for scalar sources \(q\) in the dual class paired with \(\phi\)
\cite{stratton1941}.
Readers should interpret Eq.~\eqref{eq:ch2.scalar-green-solution-operator} as a weighted
accumulation of responses arriving from every source point inside the domain $\Omega$.
The Green function determines how strongly a source at $\mathbf{r}'$ influences the
observation point $\mathbf{r}$, while $q(\mathbf{r}')$ specifies how much source strength is
available at that location \cite{stratton1941}. In plain language: $\phi(\mathbf{r})$ is the
scalar field one would measure or compute at the observation point; $q(\mathbf{r}')$ is the
source density distributed throughout the impressed region $\Omega$; and $d^{3}r'$ sums
contributions from every infinitesimal source element \cite{jackson1999}. Mathematically,
$\phi(\mathbf{r})$ carries units that depend on the physical quantity represented (potential,
pressure, or a component precursor to $\E$); $q(\mathbf{r}')$ is the corresponding source
density per cubic metre; integration extends only over $\Omega$, the region where sources are
nonzero \cite{stratton1941}.
On exterior domains, $G$ must satisfy outgoing decay compatible
with Section~\ref{sec:ch24}; on bounded domains, Dirichlet, Neumann, or impedance
data on \(\Gamma\) replace the free-space radiation template
\cite{harrington2001}.
Before examining Eq.~\eqref{eq:ch2.free-space-scalar-green} itself, it is useful to recall
the familiar physical picture of a disturbance spreading outward from a localized source.
An expanding spherical wavefront carries phase $e^{-\jj k_0 R}$ and amplitude dilution
proportional to $1/R$; the same geometric spreading responsible for the inverse-square law
in power appears here as the inverse-distance amplitude factor \cite{jackson1999}. Phase
accumulates linearly with propagation distance $R$, and outgoing energy transport requires
the retarded phase sign fixed by the locked phasor convention $\exp(-\jj\omega t)$
\cite{stratton1941}. In unbounded vacuum with outgoing selection at infinity, Sommerfeld
radiation at large $R$ singles out the spherical wave whose amplitude falls as $1/R$ and whose
phase advances with propagation distance.
\begin{equation}
\label{eq:ch2.free-space-scalar-green}
G_0(\mathbf{r},\mathbf{r}')
=
\frac{e^{-\jj k_0 R}}{4\pi R},
\qquad
R=|\mathbf{r}-\mathbf{r}'|,
\end{equation}
with \(R>0\) and the branch of \(k_0\) fixed by the locked phasor axis
\cite{stratton1941}. Equation~\eqref{eq:ch2.free-space-scalar-green}
satisfies Eq.~\eqref{eq:ch2.scalar-helmholtz-green-pde} for
\(\mathbf{r}\neq\mathbf{r}'\); the distributional definition at coincidence is
completed in Subsection~\ref{sec:ch232}. Physically,
Eq.~\eqref{eq:ch2.free-space-scalar-green} is the spherical outgoing wave produced
by a unit monopole source at \(\mathbf{r}'\); its phase gradient defines local
propagation direction on the propagating sheet of Subsection~\ref{sec:ch224}
\cite{stratton1941}. Green kernels realize the resolvent
Eq.~\eqref{eq:ch2.resolvent-operator} developed in Section~\ref{sec:ch22}: the scalar and
dyadic definitions below implement $\mathcal{G}_{E,\omega}$ of
Eq.~\eqref{eq:ch2.helmholtz-solution-operator} and feed the radiation integral
Eq.~\eqref{eq:ch2.radiation-integral} \cite{stratton1941}. Electric-wave operators
and material-weighted symbols from Subsection~\ref{sec:ch222} are presupposed; gauge enters
only through the Lorenz reduction below. Coincidence limits belong to
Subsection~\ref{sec:ch232}; outgoing completion at infinity to Section~\ref{sec:ch24}.
Vector potential and vector kernels. Direct field inversion from measured $\E$ alone is difficult because Maxwell's equations
couple $\E$ and $\H$ through curls, material response, and boundary data; the map from
sources to fields is nonlocal and orientation-dependent \cite{stratton1941}. Vector
potentials are introduced not as additional fields one would place on a probe calibration
sheet but as mathematical intermediaries that simplify reconstruction. In Lorenz gauge,
each Cartesian component of $\mathbf{A}$ satisfies the scalar Helmholtz template of
Eq.~\eqref{eq:ch2.scalar-helmholtz-green-pde} with source $-\mu_0\Jimp$
\cite{jackson1999}. The vector potential is not an
additional physical field but as a mathematical intermediary that simplifies the
reconstruction of electromagnetic quantities; gauge choices remove redundant degrees of
freedom so that propagation integrals take a standard convolution form while what an
engineer measures remains $\E$ and $\H$ \cite{stratton1941}.
Lorenz-gauge vector potentials \(\mathbf{A}\) reduce vector Maxwell inversion to a
vector Helmholtz system driven by \(\Jimp\). Under the gauge and continuity constraints
fixed in Chapter~\ref{ch:01}, each Cartesian component of \(\mathbf{A}\) satisfies
Eq.~\eqref{eq:ch2.scalar-helmholtz-green-pde} with \(k_0\) and source \(-\mu_0\Jimp\)
\cite{stratton1941}.
In Lorenz gauge the vector potential \(\mathbf{A}(\mathbf{r})\) carries units of
T\,m (equivalently Wb\,m\(^{-1}\) in SI); each Cartesian component satisfies the scalar
Helmholtz template with source \(-\mu_0\Jimp\). The electric field phasor follows from
\(\mathbf{A}\) and its divergence through
\begin{equation}
\label{eq:ch2.lorenz-field-from-potential}
\E(\mathbf{r})
=
\jj\omega\mu_0\,\mathbf{A}(\mathbf{r})
+\frac{1}{\jj\omega\varepsilon_0}\,\grad\,\divg\mathbf{A}(\mathbf{r}),
\end{equation}
with \(\E\) in V\,m\(^{-1}\), \(\varepsilon_0\) the vacuum permittivity (F\,m\(^{-1}\)),
and \(\H\) obtained from \(\curl\E\) when the first-order system is required
\cite{harrington2001}. Equation~\eqref{eq:ch2.lorenz-field-from-potential}
is a vector reduction tool, not a substitute for the dyadic definition below when
\(\mathcal{B}_\Gamma\) couples \(\E\) and \(\H\) at boundaries: impedance and PEC
closures are enforced on \(\mathcal{L}_\omega\), not on \(\mathbf{A}\) alone
\cite{stratton1941}. A vector Green function \(\mathbf{G}^{A}(\mathbf{r},\mathbf{r}')\)
may be assembled from three scalar kernels; the electric dyadic is preferred when
\(\Jimp\) is the primary excitation because it maps \(\Jimp\) to \(\E\) directly
without intermediate gauge potentials \cite{harrington2001}.
The scalar Green function describes how a point disturbance propagates. Real electromagnetic
sources, however, possess orientation and polarization. A current element directed along
$\hat{\mathbf{x}}$ generally produces field components in multiple directions. The scalar
kernel is therefore insufficient, and the propagation law must be upgraded to a tensorial
form \cite{stratton1941}. Let \(\overline{\mathbf{G}}_{E}(\mathbf{r},\mathbf{r}')\) denote a rank-two tensor
(dyad) acting on vector sources by dot product from the left; $\mathbf{r}$ is the
observation coordinate and $\mathbf{r}'$ the source coordinate.
The electric dyadic Green function answers, in dyadic form, the impulse-response
question: what electric field \(\E(\mathbf{r})\) is produced at \(\mathbf{r}\) when a
vector current impulse is placed at \(\mathbf{r}'\)?
Method-of-moments assembly implements Eq.~\eqref{eq:ch2.electric-dyadic-convolution}
by testing this kernel against patch basis functions; near-field scanners apply the same
propagation law when translating probe-corrected samples to a reference surface
\cite{harrington2001}. Readers should notice that the dyad encodes polarization
coupling absent from the scalar $G$: a transverse feed excites both radiative and reactive
dyadic components, which is why the Hertzian dipole test at $k_0 a\ll 1$ remains the
lowest-dimensional calibration before aperture catalogs are trusted \cite{jackson1999}.
It inverts the electric-wave operator
Eq.~\eqref{eq:ch2.electric-wave-operator} in the sense
\begin{equation}
\label{eq:ch2.dyadic-green-defining-operator}
\mathcal{W}_{E,\omega}\,
\overline{\mathbf{G}}_{E}(\cdot,\mathbf{r}')
=
\overline{\mathbf{I}}\,\delta(\mathbf{r}-\mathbf{r}'),
\end{equation}
where \(\overline{\mathbf{I}}\) is the identity dyad and the operator
\(\mathcal{W}_{E,\omega}\) acts on the \(\mathbf{r}\)-dependence of
\(\overline{\mathbf{G}}_{E}\) \cite{stratton1941}. Equation
\eqref{eq:ch2.dyadic-green-defining-operator} is the dyadic analogue of
Eq.~\eqref{eq:ch2.scalar-helmholtz-green-pde}: each column of
\(\overline{\mathbf{G}}_{E}\) is the electric field produced by an impulsive vector
source aligned with a Cartesian unit vector at \(\mathbf{r}'\)
\cite{jackson1999}.
Every method-of-moments code and every near-field scanner relies on the convolution
structure below: given an impressed current distribution $\Jimp$, the electric field at
observation points follows by superposition through the dyadic propagator
$\overline{\mathbf{G}}_{E}$. The physical measurement problem is forward prediction---what
field does a known or assumed current produce?---which is the prerequisite for calibration,
coupling analysis, and array synthesis \cite{harrington2001}. Inversion asks the inverse
question from sparse $\E$ samples and is substantially harder even though forward
propagation is linear and explicit \cite{hansen1988}.
Superposition over an extended impressed current \(\Jimp\) then yields the electric field
through the dyadic convolution below; the prefactor \(-\jj\omega\mu_0\) ties the integral
to the locked phasor convention \(\exp(-\jj\omega t)\).
The electric field due to impressed current is
\begin{equation}
\label{eq:ch2.electric-dyadic-convolution}
\E(\mathbf{r})
=
-\jj\omega\mu_0
\int_{\Omega}
\overline{\mathbf{G}}_{E}(\mathbf{r},\mathbf{r}')\cdot\Jimp(\mathbf{r}')\,d^{3}r'
\;+\;\E_{\mathrm{hom}}(\mathbf{r}),
\end{equation}
The importance of Eq.~\eqref{eq:ch2.electric-dyadic-convolution} lies not in the integral
itself but in the asymmetry between prediction and reconstruction. Given $\Jimp$, the field
follows directly. Given only sparse measurements of $\E$, however, recovering $\Jimp$
becomes substantially more difficult \cite{hansen1988}. Here
\(\E_{\mathrm{hom}}\) is the homogeneous solution fixed by \(\mathcal{B}_\Gamma\)
\cite{stratton1941}.
In Eq.~\eqref{eq:ch2.electric-dyadic-convolution}, \(\E(\mathbf{r})\) is the electric
field phasor in V\,m\(^{-1}\), \(\Jimp(\mathbf{r}')\) is the impressed current density in
A\,m\(^{-2}\), \(\mu_0\) is the vacuum permeability (H\,m\(^{-1}\)), and \(\omega\) is the
angular frequency in rad\,s\(^{-1}\); the volume integral extends over the source region
\(\Omega\), and \(\E_{\mathrm{hom}}\) collects excitations not attributable to \(\Jimp\),
including incident fields and cavity modes fixed by \(\mathcal{B}_\Gamma\).
Equation~\eqref{eq:ch2.electric-dyadic-convolution} is the constructive export map every
method-of-moments impedance matrix discretizes: patch currents $\Jimp$ tested against
$\overline{\mathbf{G}}_{E}$ produce the $\E$ that near-field scanners back-propagate to
equivalent surfaces \cite{harrington2001}. The consequence for inversion is
immediate: specifying $\Jimp$ determines $\E$ through a single propagator integral, while
recovering $\Jimp$ from sparse $\E$ samples is ill-posed and requires the regularization
program of Subsection~\ref{sec:ch232} \cite{stratton1941}. Equation~\eqref{eq:ch2.electric-dyadic-convolution} is the
formal realization of Eq.~\eqref{eq:ch2.radiation-integral} with outgoing dyad
\(\boldsymbol{\Gamma}=\overline{\mathbf{G}}_{E}\) on the physical branch
\cite{harrington2001}. The prefactor \(-\jj\omega\mu_0\) matches the locked phasor
convention and the forcing map Eq.~\eqref{eq:ch2.helmholtz-forcing-map}.
The operator $\mathcal{G}_{E,\omega}$ in Eq.~\eqref{eq:ch2.dyadic-green-integral-operator}
should be read as the machine that converts admissible forcing functions into electromagnetic
responses: it accepts a current distribution in the forcing class and returns the
corresponding field throughout the observation domain \cite{harrington2001}. The integral operator \(\mathcal{G}_{E,\omega}\) of
Eq.~\eqref{eq:ch2.helmholtz-solution-operator} is realized by
\begin{equation}
\label{eq:ch2.dyadic-green-integral-operator}
\big(\mathcal{G}_{E,\omega}\mathbf{f}\big)(\mathbf{r})
=
\int_{\Omega}
\overline{\mathbf{G}}_{E}(\mathbf{r},\mathbf{r}')\cdot\mathbf{f}(\mathbf{r}')\,d^{3}r',
\end{equation}
for \(\mathbf{f}\) in the forcing class of Eq.~\eqref{eq:ch2.helmholtz-forcing-map}
\cite{stratton1941}. When \(\mathbf{f}_{E,\omega}\) is the electric Helmholtz drive
induced by \(\Jimp\), Eqs.~\eqref{eq:ch2.electric-dyadic-convolution}
and~\eqref{eq:ch2.dyadic-green-integral-operator} coincide up to the linear map from
\(\Jimp\) to \(\mathbf{f}_{E,\omega}\) \cite{harrington2001}. The solution map
$\mathcal{S}_\omega$ is not merely an abstract mapping notation---it represents the complete
electromagnetic experiment beginning with impressed current and ending with observable
fields, including homogeneous background $\E^{\mathrm{hom}}$ and boundary data
$\mathcal{B}_\Gamma$ when present \cite{stratton1941}. The solution map on
impressed currents is therefore
\begin{equation}
\label{eq:ch2.solution-map-dyadic-realization}
\mathcal{S}_\omega[\Jimp]
=
\mathbf{u}_\omega
\quad\text{with}\quad
\E\ \text{from Eq.~\eqref{eq:ch2.electric-dyadic-convolution}},\
\H\ \text{from}\ \curl\E\ \text{when single-field reduction is valid},
\end{equation}
on domains satisfying Eq.~\eqref{eq:ch2.single-field-reduction-validity}; otherwise
the first-order dyadic system of Section~\ref{sec:ch234} applies
\cite{stratton1941}.
\begin{figure}[t]
\centering
\ChOneFigBlock{%
\begin{tikzpicture}[ch1fig, node distance=5mm and 7mm]
 \node[ch1box] (j) {$\Jimp(\mathbf{r}')$};
 \node[ch1box, right=7mm of j] (g) {$\overline{\mathbf{G}}_{E}$};
 \node[ch1box, right=7mm of g] (e) {$\E(\mathbf{r})$};
 \node[ch1box, below=5mm of g] (bc) {$\mathcal{B}_\Gamma$};
 \node[ch1box, below=5mm of e] (h) {$\H=\mu_0^{-1}\jj\omega^{-1}\curl\E$};
 \draw[ch1flow] (j) -- node[above,font=\scriptsize] {conv.} (g) -- (e);
 \draw[ch1flow] (bc.south) -- (bc.south |- g.north);
 \draw[ch1flow] (bc.south |- g.north) -- (g.north);
 \draw[ch1flow] (e.south) -- (e.south |- h.north);
 \draw[ch1flow] (e.south |- h.north) -- (h.north);
\end{tikzpicture}%
}
\caption{Dyadic inversion chain Eq.~\eqref{eq:ch2.electric-dyadic-convolution}:
impressed current convolves with \(\overline{\mathbf{G}}_{E}\) fixed by
\(\mathcal{B}_\Gamma\); \(\H\) follows from curl relations when reduction is admissible.}
\label{fig:ch2.dyadic-convolution-measurement chain}
\end{figure}
Equation~\eqref{eq:ch2.electric-dyadic-convolution} is a Fredholm integral of the first kind
when \(\mathbf{r}\) is fixed and \(\Jimp\) is unknown, and a constructive export map when
\(\Jimp\) is specified. The prefactor \(-\jj\omega\mu_0\) ties the kernel to the locked
phasor axis \(\exp(-\jj\omega t)\); reversing the sign convention would flip incoming versus
outgoing content without changing the local curl--curl symbol. In vacuum, for \(k_0 R=3\)
with \(R=|\mathbf{r}-\mathbf{r}'|\), the transverse dyadic component of
Eq.~\eqref{eq:ch2.free-space-dyadic-green} contributes \(\approx 70\%\) of
\(|\overline{\mathbf{G}}_{E,0}|\) in operator norm, while the \(R^{-2}\) and \(R^{-3}\)
reactive terms still matter for \(k_0 a\gtrsim 1\) near-zone scans \cite{stratton1941}.
Boundary data \(\mathcal{B}_\Gamma\) enter through the kernel that satisfies the same traces,
not through a separate post-multiplication on \(\E\) alone
\cite{harrington2001}.
In uniform vacuum with outgoing condition at infinity, let
\(\mathbf{R}=\mathbf{r}-\mathbf{r}'\), \(R=|\mathbf{R}|\), and
\(\widehat{\mathbf{R}}=\mathbf{R}/R\).
When observation and source points are separated (\(R>0\)), the outgoing vacuum dyad may
be written explicitly in terms of these geometric quantities. The projector
\(\overline{\mathbf{I}}-\widehat{\mathbf{R}}\widehat{\mathbf{R}}\) represents transverse
radiation: only field components perpendicular to the separation vector carry energy to
infinity in the far zone \cite{stratton1941}. The inverse-power corrections in $k_0 R$
represent stored reactive energy---magnetic and electric energy oscillating locally without
net outward transport at short range. An engineer measuring fields close to an antenna would
observe significant contributions from these reactive terms; far from the source they become
progressively less important, leaving only the transverse radiative component
\cite{harrington2001}.
The standard outgoing electric dyad is
\begin{equation}
\label{eq:ch2.free-space-dyadic-green}
\overline{\mathbf{G}}_{E,0}(\mathbf{r},\mathbf{r}')
=
\frac{e^{-\jj k_0 R}}{4\pi R}
\left[
\overline{\mathbf{I}}
-\widehat{\mathbf{R}}\widehat{\mathbf{R}}
+\left(\frac{1}{\jj k_0 R}-\frac{1}{k_0^{2}R^{2}}\right)
\Big(\overline{\mathbf{I}}-3\,\widehat{\mathbf{R}}\widehat{\mathbf{R}}\Big)
\right],
\end{equation}
valid for \(R>0\) \cite{stratton1941}. The
\(\overline{\mathbf{I}}-\widehat{\mathbf{R}}\widehat{\mathbf{R}}\) term is transverse
to \(\mathbf{R}\) and dominates in the radiative far zone; the $R^{-2}$ and $R^{-3}$ terms
in the parentheses supply near-field reactive structure through inverse powers of \(k_0 R\)
and dominate on electrically small probes placed within a wavelength of the feed
\cite{stratton1941}. Equation~\eqref{eq:ch2.free-space-dyadic-green} satisfies
Eq.~\eqref{eq:ch2.dyadic-green-defining-operator} for
\(\mathbf{r}\neq\mathbf{r}'\); coincidence evaluation is deferred to
Subsection~\ref{sec:ch232}. The same kernel appears in
Eq.~\eqref{eq:ch2.radiation-integral} once \(\mathbf{r}\) lies outside
\(\operatorname{supp}\Jimp\) and outgoing radiation condition is active
\cite{harrington2001}. The magnetic field may be reconstructed from the electric field through Maxwell curl relations
when the observation region is source-free and simply connected: $\H$ is not an independent
measurement channel in the Lorenz reduction but follows from the spatial derivatives of
$\E$ \cite{stratton1941}. This reconstruction is valid in homogeneous regions away from
impressed currents and fails across material interfaces or inside the source support where
curl relations must be supplemented by impressed-current terms \cite{harrington2001}.
The magnetic field may be generated by the electric dyad through Maxwell curl relations
when \(\Jimp\) is divergence-free in the source region:
\begin{equation}
\label{eq:ch2.magnetic-from-electric-dyad}
\H(\mathbf{r})
=
\frac{1}{\jj\omega\mu_0}\,\curl\E(\mathbf{r}),
\qquad
\E\ \text{from Eq.~\eqref{eq:ch2.electric-dyadic-convolution}},
\end{equation}
on simply connected source-free observation regions \cite{stratton1941}. A magnetic
dyadic \(\overline{\mathbf{G}}_{H}(\mathbf{r},\mathbf{r}')\) (A\,m\(^{-1}\)) inverts
Eq.~\eqref{eq:ch2.magnetic-wave-operator} analogously to
Eq.~\eqref{eq:ch2.dyadic-green-defining-operator}, with operator
$\mathcal{W}_{H,\omega}$ in place of $\mathcal{W}_{E,\omega}$: it is the machine that
maps magnetic impressed sources to $\H$, while $\overline{\mathbf{G}}_{H}$ denotes the
corresponding tensorial impulse response \cite{harrington2001}. Duality
\(\varepsilon_0\E\leftrightarrow\mu_0\H\), \(\Jimp\leftrightarrow\mathbf{J}^{\mathrm{(mag)}}_{\mathrm{eff}}\)
maps \(\overline{\mathbf{G}}_{E,0}\leftrightarrow\overline{\mathbf{G}}_{H,0}\) in
vacuum without re-deriving the magnetic-wave symbol. Subsection~\ref{sec:ch231} states
\(\overline{\mathbf{G}}_{H}\) only where magnetic impressed currents are primary;
otherwise Eq.~\eqref{eq:ch2.magnetic-from-electric-dyad} avoids duplicate kernel
tabulations \cite{stratton1941}.
The preceding formulas assume a uniform medium. Real structures seldom satisfy this
assumption. Substrates, radomes, biological tissues, and layered materials alter
propagation and therefore modify the Green response \cite{collin2001}.
In piecewise media, local symbols \(k(\mathbf{r})\), \(\varepsilon(\mathbf{r})\),
\(\mu(\mathbf{r})\) modify Eq.~\eqref{eq:ch2.dyadic-green-defining-operator} to
\begin{equation}
\label{eq:ch2.material-dyadic-defining-operator}
\mathcal{W}_{E,\omega}^{(\varepsilon,\mu)}\,
\overline{\mathbf{G}}_{E}^{(\varepsilon,\mu)}(\cdot,\mathbf{r}')
=
\overline{\mathbf{I}}\,\delta(\mathbf{r}-\mathbf{r}'),
\end{equation}
with \(\mathcal{W}_{E,\omega}^{(\varepsilon,\mu)}\) from
Eq.~\eqref{eq:ch2.material-helmholtz-electric} \cite{stratton1941}. Free-space
Eq.~\eqref{eq:ch2.free-space-dyadic-green} applies patchwise only when
\(k(\mathbf{r})\) is constant in each patch and interface jumps are enforced by
continuity of tangential \(\E,\H\) across shared boundaries
(Eq.~\eqref{eq:ch1.interface-conditions}); multi-patch assembly is
Subsection~\ref{sec:ch234} \cite{harrington2001}. Reciprocity expresses a fundamental
symmetry of passive electromagnetic environments: interchanging transmitter and receiver
does not alter the propagation law itself \cite{stratton1941}. On passive media without
magneto-electric coupling, reciprocity requires that interchanging observation and
source points leaves the dyadic invariant up to transposition:
\begin{equation}
\label{eq:ch2.dyadic-reciprocity-swap}
\overline{\mathbf{G}}_{E}(\mathbf{r},\mathbf{r}')
=
\overline{\mathbf{G}}_{E}^{\mathsf{T}}(\mathbf{r}',\mathbf{r}),
\end{equation}
at fixed $\omega$; the full proof is in Section~\ref{sec:ch28} \cite{stratton1941}.
Equations~\eqref{eq:ch2.electric-dyadic-convolution}--\eqref{eq:ch2.solution-map-dyadic-realization}
hold when \(\Jimp\) lies in the domain of the forcing map, when
\(\mathcal{B}_\Gamma\) is time-independent and identical under phasor reduction, and
when the selected kernel branch matches the resolvent of
Eq.~\eqref{eq:ch2.resolvent-green-equivalence}. They fail for resonant cavities
without quotienting, for incoming rather than outgoing exterior conditions, and for
sources with \(\divg\Jimp\neq 0\) unless charge density \(\rho\) is included in the
complete source tuple of Section~\ref{sec:ch01.1.2} \cite{stratton1941}.
Near \(\mathbf{r}=\mathbf{r}'\), Eqs.~\eqref{eq:ch2.free-space-scalar-green}
and~\eqref{eq:ch2.free-space-dyadic-green} are pointwise undefined; convolution
meaning is distributional (Subsection~\ref{sec:ch232}).
The scalar monopole kernel $G$ supplies the lowest-dimensional propagation law;
the electric dyad $\overline{\mathbf{G}}_{E}$ is the vector generalization demanded
by polarization on real feeds. Subsection~\ref{sec:ch23.signature-dipole} closes the
loop in the opposite pedagogical order: a measured dipole pattern on $\Gamma_R$
precedes the dyadic tabulation, so the laboratory observation remains the anchor
when the full tensor formulas are applied to extended apertures
\cite{jackson1999}. Scalar monopole law $G$, vector dyad
$\overline{\mathbf{G}}_{E}$, and operator realization $\mathcal{G}_{E,\omega}$ are now fixed with
$\mathbf{r}$ as observation and $\mathbf{r}'$ as source throughout; coincidence limits and
retarded causality refine how those kernels may be convolved with extended $\Jimp$
\cite{stratton1941}.
\subsection{Point electric dipole and the dyadic Green operator}
\label{sec:ch23.signature-dipole}
Engineering reading (Hertzian chain). \emph{Physical setup:} a short wire or probe feed whose electrical size satisfies
\(k_0a\ll 1\) on a turntable range.
\emph{Objective:} verify rank-one far-zone export before multipole catalogs are applied
to larger apertures.
\emph{Expected intuition:} a transverse \(\sin\theta\) pattern with a single dominant
radiative channel.
\emph{Measurement:} watt closure on a sphere \(\Gamma_R\) and turntable invariance of the
total radiated power \cite{jackson1999}.
\emph{Misconception:} equating mesh unknowns in a numerical model with the radiative rank
of the physical source \cite{chu1948}.
% PHYSICAL-MEMORY:sec:ch23.signature-dipole
\begin{physicalvignette}[Rectangular horn mouth in a ground-plane flange]
Consider Rectangular horn mouth in a ground-plane flange. Before any coefficient table is printed, technicians record frequency, cable loss, and the mechanical outline that fixes electrical size. The first walk-through of the range is always tactile: flange bolts, probe standoff, turntable boresight, and the wattmeter reference on the receive chain \cite{hansen1988}.
On this hardware the near zone is visually impressive---standing reactive structure, rapid phase rotation, sometimes a visible VSWR ripple at the feed---yet the export question is measured only after outgoing selection on a large sphere \cite{harrington2001}. What the camera sees in the chamber is not automatically what the far-zone pattern claims \cite{stratton1941}.
Engineers therefore learn to separate three readings: what the feed can synthesize on declared support, what propagating content survives on the radiative range, and what a finite instrument can recover at declared tolerance \cite{devaney2004}. Confusing those readings produces the rank gaps and watt failures that later sections quantify \cite{pal2026}.
The Hertzian dipole is the lowest-dimensional test of outgoing selection and far-zone flux \cite{jackson1999}. Every compact feed should be checked against it at low $ka$ \cite{chu1948}.
Rayleigh and Stratton treated such ranges as boundary-value laboratories long before digital near-field transforms became routine \cite{stratton1941}. Harrington's method-of-moments culture reinforced the same discipline: metal outline first, harmonic headline second \cite{harrington2001}.
Subsection~\ref{sec:ch23.signature-dipole} states the formal export condition that this range scene presupposes \cite{stratton1941}.
\end{physicalvignette}
On a turntable compact range, the dipole test is performed before any aperture catalog:
rotate the electrically small feed, integrate $|\E|^2$ on $\Gamma_R$, and confirm that the
total radiated power is invariant and that the far-zone pattern carries a single transverse
lobe \cite{hansen1988}. That watt closure is the laboratory observation that
precedes the dyadic formulas of Subsection~\ref{sec:ch231}; when it fails, the fault lies
in calibration, outgoing branch selection, or source admissibility---not in the multipole
catalog applied later \cite{stratton1941}.
The Hertzian dipole is the canonical image of the dyadic Green operator: a localized
source, a retarded kernel, a radiation map, and a far-zone angular spectrum linked in
one closed calculation \cite{stratton1941}. Short-wire or probe feeds at $k_0a\ll1$ on
a turntable should reproduce transverse $\sin\theta$ export and turntable-invariant watt
totals before multipole catalogs are applied to larger apertures
\cite{jackson1999}. Every compact feed should be checked against
this lowest-dimensional outgoing test at low $ka$ \cite{chu1948}.
The range observation fixed above---watt closure on $\Gamma_R$, turntable invariance,
and a single transverse lobe---is what a correct Green propagation must reproduce before
any multipole catalog is trusted on larger structures \cite{hansen1988}.
Only after that calibration does it pay to write the dyadic convolution in closed form:
replace the extended source $\Jimp(\mathbf{r}')$ by a rank-one impulse and read off the
field pair that every method-of-moments code must match at electrically small $k_0 a$
\cite{harrington2001}.
The Hertzian dipole is chosen as the canonical calibration source because it is the
smallest oriented radiator that still exercises the full dyadic structure: one moment
vector, one propagation direction, and a transverse far-zone pattern that every radiation
theory must reproduce \cite{jackson1999}. Compact feeds on turntable ranges,
probe references, and method-of-moments self-tests all reduce to this limit at
$k_0 a\ll 1$; when a code or measurement sequence cannot match the dipole field pair, no
larger aperture catalog built on the same Green operator should be trusted
\cite{hansen1988}.
The dipole moment \(\mathbf{p}=\int \Jimp\,d^{3}r'\) (units: A\,m, equivalently C\,m for
equivalent polarization sources) localizes the impressed current at the origin; $\mathbf{r}$
denotes the observation point and $\mathbf{r}'$ the source coordinate in the display below.
Substituting \(\Jimp(\mathbf{r}')=\mathbf{I}\,\delta(\mathbf{r}')\) into
Eq.~\eqref{eq:ch2.electric-dyadic-convolution} collapses the volume integral to a single
dyadic--vector product and yields the closed-form field pair
\begin{equation}
\label{eq:ch2.signature-dipole-field}
\E(\mathbf{r})
=
-\jj\omega\mu_0\,\overline{\mathbf{G}}_{E,0}(\mathbf{r},\mathbf{0})\cdot\mathbf{p},
\qquad
\H(\mathbf{r})=\frac{1}{\jj\omega\mu_0}\,\curl\E(\mathbf{r}),
\end{equation}
with outgoing dyad Eq.~\eqref{eq:ch2.free-space-dyadic-green}. This is
\(\mathcal{S}_\omega[\Jimp]\) in closed form; \(\mathcal{R}_\omega[\Jimp]\) is the
same pair after outgoing selection on the exterior domain
\cite{stratton1941}. Physically, Eq.~\eqref{eq:ch2.signature-dipole-field} states that
every Cartesian component of $\mathbf{p}$ excites a column of $\overline{\mathbf{G}}_{E,0}$,
so polarization coupling is already present in the rank-one source that calibrates compact
feeds \cite{stratton1941}.
\begin{figure}[t]
\centering
\ChOneFigBlock{%
\begin{tikzpicture}[ch1fig, node distance=5mm and 7mm]
 \node[ch1box] (p) {Dipole $\mathbf{p}$};
 \node[ch1box, right=7mm of p] (g) {$\overline{\mathbf{G}}_{E,0}$};
 \node[ch1box, right=7mm of g] (e) {$\E,\H$};
 \node[ch1box, right=7mm of e] (r) {$\mathcal{R}_\omega$};
 \node[ch1box, right=7mm of r] (a) {$a(\widehat{\mathbf{k}})$};
 \draw[ch1flow] (p) -- (g) -- (e) -- (r) -- (a);
\end{tikzpicture}%
}
\caption{Signature chain for a Hertzian dipole: impressed moment convolves with the
outgoing dyadic Green function, producing \((\E,\H)\); the radiation map extracts the
far-zone amplitude and angular spectrum.}
\label{fig:ch2.signature-hertzian-dipole}
\end{figure}
In the radiative far zone \(k_0 r\gg 1\), with \(r=|\mathbf{r}|\) the observation radius
in metres, the reactive \(R^{-2}\) and \(R^{-3}\) terms in
Eq.~\eqref{eq:ch2.free-space-dyadic-green} become negligible and
Eq.~\eqref{eq:ch2.signature-dipole-field} reduces to
Eq.~\eqref{eq:ch2.far-field-amplitude-integral} with
\begin{equation}
\label{eq:ch2.signature-dipole-far-field}
\E(\mathbf{r})\approx
-\jj\omega\mu_0\,\frac{e^{-\jj k_0 r}}{4\pi r}\,
\overline{\mathbf{I}}-\widehat{\mathbf{r}}\widehat{\mathbf{r}}\,\cdot\mathbf{p},
\end{equation}
so the transverse radiative image is fixed by \(\widehat{\mathbf{r}}\times(\widehat{\mathbf{r}}\times\mathbf{p})\).
This expression reveals that only transverse components survive at large distance. The far
field therefore forgets most local source details and retains only radiatively accessible
information---precisely the content a compact-range pattern measurement can resolve on
$\Gamma_R$ \cite{stratton1941}.
Power through sphere \(S_r\) follows Eq.~\eqref{eq:ch2.radiated-power-screen} with
\(P_{\mathrm{rad}}\propto k_0^{4}|\mathbf{p}|^{2}\) in the electrically small limit
\cite{stratton1941}. Experimentally, the angular spectrum is the decomposition of the
far-zone field into plane waves propagating in different directions $\widehat{\mathbf{k}}$:
each outgoing direction receives a distinct amplitude fixed by how the dipole moment
projects transversely onto that direction \cite{hansen1988}. Pattern measurements on a
spherical scanner therefore sample this angular distribution; synthesis codes invert the
same lateral Fourier image when reconstructing aperture fields from far-zone data
\cite{stratton1941}. The plane-wave angular spectrum of
Subsection~\ref{sec:ch291} reads
\begin{equation}
\label{eq:ch2.signature-dipole-angular-spectrum}
a(\widehat{\mathbf{k}})
=
-\jj k_0^{2}\,\sqrt{\mu_0/\varepsilon_0}\,
\big(\overline{\mathbf{I}}-\widehat{\mathbf{k}}\widehat{\mathbf{k}}\big)\cdot\mathbf{p},
\end{equation}
where \(a(\widehat{\mathbf{k}})\) is the far-zone plane-wave amplitude (V\,m\(^{-1}\))
associated with propagation direction \(\widehat{\mathbf{k}}\) (dimensionless unit vector)
and \(\widehat{\mathbf{k}}\) parametrizes the angular-spectrum parametrization of
Subsection~\ref{sec:ch291}. Equation~\eqref{eq:ch2.signature-dipole-angular-spectrum}
is the lateral Fourier image of the same export channel: each outgoing direction receives
transverse amplitude fixed by $(\overline{\mathbf{I}}-\widehat{\mathbf{k}}\widehat{\mathbf{k}})\cdot\mathbf{p}$.
One source, one kernel, one radiation map, one far field, one angular spectrum---every
abstract operator statement in Chapter~\ref{ch:02} reduces to this geometry when
\(\Jimp=\mathbf{p}\,\delta(\mathbf{r}')\) \cite{stratton1941}. The dipole test on
$\Gamma_R$ is therefore not a numerical afterthought; it is the lowest-dimensional
witness that $G$, $\overline{\mathbf{G}}_{E}$, $\mathcal{S}_\omega$, and
$\mathcal{R}_\omega$ describe the same hardware \cite{hansen1988}.
\subsection{Singularities, distributions, and regularization}
\label{sec:ch232}
Layered substrates in patch and handset designs modify Green propagation through
reflection coefficients at material interfaces \cite{stratton1941}. Printed antennas and
handset boards propagate through layered media; the free-space dyadic is the reference
impulse, not the laboratory default when a probe hovers above a dielectric stack
\cite{collin2001}. Probe correction files must include substrate Green
response, not free-space alone, because the governing question---how a point source on a
feed line launches fields to a probe above the substrate---already presupposes outgoing
propagation in each half-space \cite{sommerfeld1949}. Kernels transport impressed
currents; they do not specify which currents the feed can realize \cite{devaney2004}.
Subsection~\ref{sec:ch231} defined scalar and dyadic Green kernels that satisfy
Eqs.~\eqref{eq:ch2.scalar-helmholtz-green-pde} and~\eqref{eq:ch2.dyadic-green-defining-operator}
for \(\mathbf{r}\neq\mathbf{r}'\) and deferred coincidence evaluation to the present
subsection \cite{stratton1941}. Subsection~\ref{sec:ch232}
supplies the definitive account of distributional convolutions, principal-value
and finite-part prescriptions, static--dynamic kernel splits, and self-term
regularization that make Eq.~\eqref{eq:ch2.electric-dyadic-convolution} well-defined
on \(\Jimp\in\mathcal{J}_\Omega\) including sheet and filamentary support. The
Sobolev and distributional source classes of Subsection~\ref{sec:ch211} are presupposed
through Eqs.~\eqref{eq:ch2.distributional-source-extension} and~\eqref{eq:ch2.dual-source-pairing};
their functional-analytic construction is not repeated. Retarded time-domain support
belongs to Subsection~\ref{sec:ch233}; outgoing branch selection at infinity belongs
to Section~\ref{sec:ch24}. Let \(\Omega\Subset\mathbb{R}^{3}\) be Lipschitz with piecewise smooth internal
interfaces, \(\omega>0\) fixed, and \(\Jimp\) in the refined source class
Eq.~\eqref{eq:ch2.source-space-refined} or its distributional extension
Eq.~\eqref{eq:ch2.distributional-source-extension}. Kernels are the vacuum templates
Eqs.~\eqref{eq:ch2.free-space-scalar-green} and~\eqref{eq:ch2.free-space-dyadic-green}
unless \(\mathcal{B}_\Gamma\) or piecewise \(k(\mathbf{r})\) modify them patchwise.
Observation points \(\mathbf{r}\) may approach \(\operatorname{supp}\Jimp\) but
remain in the domain where \(\E\in\HCurl(\Omega)\) is sought.
For \(\mathbf{r}\notin\operatorname{supp}\Jimp\), the convolution
Eq.~\eqref{eq:ch2.electric-dyadic-convolution} is a classical integral. What happens when
observation and source points become arbitrarily close? The free-space dyad develops a
non-integrable leading singularity: field amplitudes grow as inverse powers of separation
$R$, reflecting reactive storage and self-interaction between neighboring source elements
\cite{stratton1941}. These singularities do not imply infinite radiated power---they signal
that pointwise source models and naive quadrature break down, while extended patches with
finite extent carry finite reactive energy \cite{jackson1999}. As
\(\mathbf{r}\to\mathbf{r}'\in\operatorname{supp}\Jimp\), the free-space dyad
Eq.~\eqref{eq:ch2.free-space-dyadic-green} develops a non-integrable leading singularity;
the local expansion in powers of \(k_0 R\) separates static storage from dynamic radiation
and is the analytic starting point for every regularization scheme in the display below.
With separation vector \(\mathbf{R}=\mathbf{r}-\mathbf{r}'\), distance \(R=|\mathbf{R}|\) in
metres, observation point $\mathbf{r}$ and source point $\mathbf{r}'$, and unit direction
\(\widehat{\mathbf{R}}=\mathbf{R}/R\),
% PHYSICAL-MEMORY:sec:ch232
\begin{physicalvignette}[Open-ended coax probe above a calibration table]
Consider Open-ended coax probe above a calibration table. Before any coefficient table is printed, technicians record frequency, cable loss, and the mechanical outline that fixes electrical size. The first walk-through of the range is always tactile: flange bolts, probe standoff, turntable boresight, and the wattmeter reference on the receive chain \cite{hansen1988}.
On this hardware the near zone is visually impressive---standing reactive structure, rapid phase rotation, sometimes a visible VSWR ripple at the feed---yet the export question is measured only after outgoing selection on a large sphere \cite{harrington2001}. What the camera sees in the chamber is not automatically what the far-zone pattern claims \cite{stratton1941}.
Engineers therefore learn to separate three readings: what the feed can synthesize on declared support, what propagating content survives on the radiative range, and what a finite instrument can recover at declared tolerance \cite{devaney2004}. Confusing those readings produces the rank gaps and watt failures that later sections quantify \cite{pal2026}.
Rayleigh and Stratton treated such ranges as boundary-value laboratories long before digital near-field transforms became routine \cite{stratton1941}. Harrington's method-of-moments culture reinforced the same discipline: metal outline first, harmonic headline second \cite{harrington2001}.
Subsection~\ref{sec:ch232} states the formal export condition that this range scene presupposes \cite{stratton1941}.
\end{physicalvignette}
\begin{equation}
\label{eq:ch2.dyadic-singular-expansion}
\overline{\mathbf{G}}_{E,0}(\mathbf{r},\mathbf{r}')
=
\frac{1}{4\pi R^{3}}\,\overline{\mathbf{P}}_{t}(\widehat{\mathbf{R}})
\;+\;
\frac{\jj k_0}{4\pi R^{2}}\,\overline{\mathbf{Q}}(\widehat{\mathbf{R}})
\;+\;
\frac{k_0^{2}}{4\pi R}\,\overline{\mathbf{S}}(\widehat{\mathbf{R}})
\;+\;\mathcal{O}(k_0^{3}),
\end{equation}
where \(\overline{\mathbf{P}}_{t}(\widehat{\mathbf{R}})=\overline{\mathbf{I}}-3\widehat{\mathbf{R}}\widehat{\mathbf{R}}\)
is transverse to \(\mathbf{R}\) and \(\overline{\mathbf{Q}}\), \(\overline{\mathbf{S}}\)
collect the next two anisotropic terms in the locked phasor convention
\cite{harrington2001}. Equation~\eqref{eq:ch2.dyadic-singular-expansion} matches the
\(R^{-3}\), \(R^{-2}\), and \(R^{-1}\) pieces of
Eq.~\eqref{eq:ch2.free-space-dyadic-green} and separates reactive near storage from
radiative \(R^{-1}\) decay. The \(R^{-3}\) term is non-integrable in a neighborhood
of coincidence; convolution with \(\Jimp\) therefore requires a distributional or
regularized definition, not pointwise Riemann integration
\cite{stratton1941}. The scalar kernel Eq.~\eqref{eq:ch2.free-space-scalar-green}
carries the parallel ordering \(R^{-1}\), \(R^{0}\), \ldots\ with
\(R^{-1}\) non-integrable on neighborhoods of the source.
In the electrostatic limit \(k_0\to 0\) at fixed separation \(R>0\) (metres), the dynamic
terms in Eq.~\eqref{eq:ch2.dyadic-singular-expansion} vanish and only the \(R^{-3}\)
dipolar projector survives; that static piece is extracted analytically in
method-of-moments codes before the smooth remainder is integrated numerically.
Eq.~\eqref{eq:ch2.dyadic-singular-expansion} reduces to
\begin{equation}
\label{eq:ch2.static-dyadic-limit}
\overline{\mathbf{G}}_{E,0}^{\mathrm{(st)}}(\mathbf{r},\mathbf{r}')
=
\frac{1}{4\pi R^{3}}\Big(\overline{\mathbf{I}}-3\,\widehat{\mathbf{R}}\widehat{\mathbf{R}}\Big),
\end{equation}
the field of a unit dipole aligned with each Cartesian source direction
\cite{jackson1999}. Physically, Eq.~\eqref{eq:ch2.static-dyadic-limit} captures static-field
dominance at small $k_0 R$: reactive storage and self-interaction between neighboring
current elements overwhelm the radiative $R^{-1}$ term until separation or frequency
increases \cite{stratton1941}. Equation~\eqref{eq:ch2.static-dyadic-limit} is the
leading singular part extracted in method-of-moments static--dynamic splits: dynamic
corrections are \(\mathcal{O}(k_0 R)\) relative to the static term when
\(k_0 R\ll 1\) on source patches \cite{harrington2001}. For divergence-free
\(\Jimp\), the \(\grad\,\divg\) structure of Lorenz reduction
Eq.~\eqref{eq:ch2.lorenz-field-from-potential} removes monopole artifacts; the
dominant self-coupling on compact patches is dipolar and transverse to the separation
vector between observation and source centroids \cite{stratton1941}.
\begin{figure}[t]
\centering
\ChOneFigBlock{%
\begin{tikzpicture}[ch1fig, node distance=5mm and 7mm]
 \node[ch1box] (src) {Source patch\\$\operatorname{supp}\Jimp$};
 \node[ch1box, right=16mm of src] (near) {Near zone\\$R\to 0$};
 \node[ch1box, right=16mm of near] (reg) {Regularized\\convolution};
 \node[ch1box, below=5mm of near] (sing) {$G_{\mathrm{sing}}$};
 \node[ch1box, below=5mm of reg] (cpt) {$K_{\mathrm{cpt}}$};
 \draw[ch1flow] (src) -- (near) -- (reg);
 \draw[ch1flow] (near.south) -- (near.south |- sing.north);
 \draw[ch1flow] (near.south |- sing.north) -- (sing.north);
 \draw[ch1flow] (reg.south) -- (reg.south |- cpt.north);
 \draw[ch1flow] (reg.south |- cpt.north) -- (cpt.north);
\end{tikzpicture}%
}
\caption{Singular--regular split at coincidence: non-integrable \(G_{\mathrm{sing}}\)
is extracted analytically; remainder \(K_{\mathrm{cpt}}\) is smooth for Galerkin or
point evaluation away from the extracted neighborhood.}
\label{fig:ch2.singular-regularization-split}
\end{figure}
~\eqref{eq:ch2.green-compact-split} at kernel level. The decomposition separates analytically known local physics from the smoother long-range interaction that can be handled numerically: $\overline{\mathbf{G}}_{E}^{\mathrm{(sing)}}$ collects the non-integrable near-zone coupling, while $\overline{\mathbf{G}}_{E}^{\mathrm{(reg)}}$ is the remainder suitable for standard quadrature \cite{harrington2001}. Write
\begin{equation}
\label{eq:ch2.kernel-singular-regular-split}
\overline{\mathbf{G}}_{E}(\mathbf{r},\mathbf{r}')
=
\overline{\mathbf{G}}_{E}^{\mathrm{(sing)}}(\mathbf{r},\mathbf{r}')
+\overline{\mathbf{G}}_{E}^{\mathrm{(reg)}}(\mathbf{r},\mathbf{r}'),
\end{equation} with \(\overline{\mathbf{G}}_{E}^{\mathrm{(sing)}}\) collecting the non-integrable
static and leading dynamic singularities of Eq.~\eqref{eq:ch2.dyadic-singular-expansion}
and \(\overline{\mathbf{G}}_{E}^{\mathrm{(reg)}}\) bounded on
\(\Omega\times\Omega\) \cite{stratton1941}. The corresponding
solution operator split is
\begin{equation}
\label{eq:ch2.solution-singular-compact-split}
\mathcal{S}_\omega
=
\mathcal{G}_{\mathrm{sing}}+\mathcal{K}_{\mathrm{cpt}},
\qquad
\mathcal{G}_{\mathrm{sing}}[\Jimp]
=
-\jj\omega\mu_0\int_{\Omega}
\overline{\mathbf{G}}_{E}^{\mathrm{(sing)}}(\mathbf{r},\mathbf{r}')\cdot\Jimp(\mathbf{r}')\,d^{3}r',
\end{equation}
with \(\mathcal{K}_{\mathrm{cpt}}\) compact on bounded \(\Omega\) when
\(\overline{\mathbf{G}}_{E}^{\mathrm{(reg)}}\) is square-integrable in both variables
\cite{stratton1941}. Here $\mathcal{G}_{\mathrm{sing}}$ is the singular propagation
operator carrying non-integrable near-zone coupling, $\mathcal{K}_{\mathrm{cpt}}$ is the
compact remainder suitable for stable numerical quadrature, and
$\overline{\mathbf{G}}_{E}^{\mathrm{(sing)}}$ and $\overline{\mathbf{G}}_{E}^{\mathrm{(reg)}}$
are the corresponding kernel split defined in Eq.~\eqref{eq:ch2.kernel-singular-regular-split}
\cite{harrington2001}. Equations~\eqref{eq:ch2.kernel-singular-regular-split}
and~\eqref{eq:ch2.solution-singular-compact-split} realize
Eq.~\eqref{eq:ch2.green-compact-split} for Green-based inversion; Fredholm consequences
are developed in Section~\ref{sec:ch263} \cite{harrington2001}.
When \(\Jimp\) is a bounded function, Eq.~\eqref{eq:ch2.electric-dyadic-convolution}
defines \(\E\in L^{2}_{\mathrm{loc}}(\Omega)^{3}\) away from the source. For sheet
currents Eq.~\eqref{eq:ch2.distributional-source-extension}, the convolution is defined
by duality: for all smooth compactly supported test fields \(\boldsymbol{\phi}\) (dimensionless
weighting fields that probe how the electric field pairs against admissible test states),
\begin{equation}
\label{eq:ch2.distributional-dyadic-pairing}
\langle\E,\boldsymbol{\phi}\rangle_{\Omega}
=
-\jj\omega\mu_0
\sum_{ab}\int_{\Sigma_{ab}}\int_{\Omega}
\mathbf{J}_{s,ab}(\mathbf{r}')
\cdot\overline{\mathbf{G}}_{E}^{\mathsf{T}}(\mathbf{r}',\mathbf{r})\cdot\boldsymbol{\phi}(\mathbf{r})\,d^{3}r\,dS',
\end{equation}
Here $\boldsymbol{\phi}$ is a smooth test field used to probe the electric field when
pointwise evaluation fails; $\Sigma_{ab}$ is a surface patch carrying sheet current
$\mathbf{J}_{s,ab}$ (A\,m\(^{-1}\)), the surface density of impressed current on patch
$(a,b)$; and $\mathbf{J}_{s,ab}(\mathbf{r}')$ is the physical current flowing on that
patch at source point $\mathbf{r}'$ \cite{stratton1941}, with appropriate Fubini interpretation and outgoing radiation condition on exterior problems
\cite{stratton1941}. Equation~\eqref{eq:ch2.distributional-dyadic-pairing} is the
weak form of Eq.~\eqref{eq:ch2.electric-dyadic-convolution} when pointwise evaluation
fails; it pairs with Eq.~\eqref{eq:ch2.maxwell-weak-form} once \(\boldsymbol{\phi}\)
is replaced by admissible test states in \(\mathcal{F}_\Omega^{0}\)
\cite{harrington2001}. Physically, sheet sources produce jump conditions on
\(\hat{\mathbf{n}}\times\E\) and \(\hat{\mathbf{n}}\times\H\); the distributional
definition ensures those jumps arise from the same kernel as volume currents, not from
ad hoc post-factors \cite{stratton1941}.
For volume \(\Jimp\) and observation points \(\mathbf{r}\notin\operatorname{supp}\Jimp\)
approaching the support, one may attempt to define the field as the symmetric limit of
integrals that exclude a small ball about the observation point. The principal value is that
symmetric exclusion limit: ordinary integration fails because the $R^{-3}$ pole is not
integrable over a neighborhood of coincidence, but symmetric cancellation about the
observation point can still yield a finite field when the singularity is weak enough
\cite{jackson1999}. Define the principal-value convolution
\begin{equation}
\label{eq:ch2.principal-value-convolution}
\E(\mathbf{r})
=
-\jj\omega\mu_0\,\mathrm{p.v.}\int_{\Omega}
\overline{\mathbf{G}}_{E}(\mathbf{r},\mathbf{r}')\cdot\Jimp(\mathbf{r}')\,d^{3}r',
\end{equation}
as the limit of integrals over \(\Omega\setminus B_{\varepsilon}(\mathbf{r})\) as
\(\varepsilon\downarrow 0\) when that limit exists in \(\HCurl(\Omega)\)
\cite{jackson1999}. For observation on the source support---for example
\(\mathbf{r}\in\Sigma_{ab}\) for a surface patch---the principal value generally
diverges because of the \(R^{-3}\) term; the \emph{Hadamard finite part} subtracts
the singular expansion explicitly and retains the finite reactive coupling between
extended patches. Finite-part subtraction works because the non-integrable piece is known
analytically from Eq.~\eqref{eq:ch2.dyadic-singular-expansion} and can be removed before
the remainder is integrated numerically---the engineering prescription MoM codes implement
as static--dynamic splitting \cite{harrington2001}:
\begin{equation}
\label{eq:ch2.finite-part-convolution}
\E(\mathbf{r})
=
-\jj\omega\mu_0\lim_{\varepsilon\to 0}
\left[
\int_{\Omega\setminus B_{\varepsilon}(\mathbf{r})}
\overline{\mathbf{G}}_{E}(\mathbf{r},\mathbf{r}')\cdot\Jimp(\mathbf{r}')\,d^{3}r'
\;+\;
\int_{B_{\varepsilon}(\mathbf{r})\cap\Omega}
\overline{\mathbf{G}}_{E}^{\mathrm{(sing)}}(\mathbf{r},\mathbf{r}')\cdot\Jimp(\mathbf{r}')\,d^{3}r'
\right],
\end{equation}
provided \(\Jimp\) is smooth enough on the patch that the limit exists
\cite{harrington2001}. Equation~\eqref{eq:ch2.finite-part-convolution} is the
prescription for on-surface field evaluation; it differs
from Eq.~\eqref{eq:ch2.principal-value-convolution} by the explicit analytic
subtraction of \(\overline{\mathbf{G}}_{E}^{\mathrm{(sing)}}\), not by numerical
mesh refinement alone \cite{stratton1941}. Mathematically, the finite part is the
constant term in the Laurent expansion of the integral as \(\varepsilon\to 0\);
physically, it removes infinite self-energy of pointwise source models while retaining
finite reactive coupling between extended patches \cite{jackson1999}.
Method-of-moments and boundary-element codes replace the continuum convolution
Eq.~\eqref{eq:ch2.electric-dyadic-convolution} by discrete impedance entries obtained
by testing the dyadic kernel against pairs of basis functions on source patches. In plain
language: $\mathbf{f}_m$ and $\mathbf{f}_n$ are the $m$th and $n$th current basis functions
on patches $\Omega_m$ and $\Omega_n$; $\Omega_m$ and $\Omega_n$ are the surface or volume
patches where those basis functions live; and $Z_{mn}$ is the mutual impedance---the
complex voltage response at patch $m$ per unit current on patch $n$---that a numerical code
assembles from the double integral below \cite{harrington2001}. What a MoM code actually
computes is this bilinear pairing of basis functions through the dyadic propagator, not a
pointwise field evaluation at patch centers unless a separate approximation is introduced.
The Galerkin entry for patches \(\Omega_m\) and \(\Omega_n\) reads
\begin{equation}
\label{eq:ch2.mom-dyadic-entry}
Z_{mn}
=
-\jj\omega\mu_0\int_{\Omega_m}\int_{\Omega_n}
\mathbf{f}_m(\mathbf{r})\cdot
\overline{\mathbf{G}}_{E}(\mathbf{r},\mathbf{r}')\cdot
\mathbf{f}_n(\mathbf{r}')\,d^{3}r\,d^{3}r',
\end{equation}
for basis functions \(\mathbf{f}_m,\mathbf{f}_n\) on patches \(\Omega_m,\Omega_n\)
\cite{harrington2001}. When \(m=n\), the integrand is non-integrable at
\(\mathbf{r}=\mathbf{r}'\); the Chapter~\ref{ch:02} regularized self-term is
\begin{equation}
\label{eq:ch2.mom-self-term-regularized}
Z_{mm}^{\mathrm{(reg)}}
=
-\jj\omega\mu_0\int_{\Omega_m}\int_{\Omega_m}
\mathbf{f}_m(\mathbf{r})\cdot
\overline{\mathbf{G}}_{E}^{\mathrm{(reg)}}(\mathbf{r},\mathbf{r}')\cdot
\mathbf{f}_m(\mathbf{r}')\,d^{3}r\,d^{3}r'
\;+\;
Z_{mm}^{\mathrm{(sing)}},
\end{equation}
where $Z_{mm}^{\mathrm{(reg)}}$ is the numerically integrated smooth remainder,
$Z_{mm}^{\mathrm{(sing)}}$ is the analytic integral of
\(\overline{\mathbf{G}}_{E}^{\mathrm{(sing)}}\) over the coincident patch---often
reduced to closed forms for flat triangles and rectangles---and the sum is the quantity a
code must assemble so that neighboring mesh elements do not produce near-singular divergences
when observation and source points approach one another \cite{stratton1941}. Point
matching---evaluating Eq.~\eqref{eq:ch2.electric-dyadic-convolution} at patch centers
with \(\overline{\mathbf{G}}_{E}\) replaced by
\(\overline{\mathbf{G}}_{E}^{\mathrm{(reg)}}\) only---is \emph{not} equivalent to
Eq.~\eqref{eq:ch2.mom-self-term-regularized} unless basis overlap and singular
subtraction match the same finite-part convention \cite{harrington2001}. The
distinction is one of source realizability: Galerkin testing defines a bounded bilinear form on
\(\mathcal{J}_\Omega\times\mathcal{J}_\Omega\); raw point evaluation does not
\cite{stratton1941}.
When \(\mathbf{r}\) and \(\mathbf{r}'\) lie on distinct patches separated by small
distance \(d\ll \lambda_0=2\pi/k_0\), direct quadrature of
Eq.~\eqref{eq:ch2.free-space-dyadic-green} is ill-conditioned. Physically, neighboring mesh
elements interact through nearly coincident kernel samples: the separation is small enough
that reactive $R^{-3}$ and $R^{-2}$ terms vary rapidly across the integration domain, and
standard quadrature underestimates or overestimates the coupling unless the singular part is
subtracted and integrated analytically \cite{harrington2001}. The standard
subtraction template is
\begin{equation}
\label{eq:ch2.near-singular-subtraction}
\int_{\Omega_n}
\overline{\mathbf{G}}_{E}(\mathbf{r},\mathbf{r}')\cdot\Jimp(\mathbf{r}')\,d^{3}r'
=
\int_{\Omega_n}
\Big[\overline{\mathbf{G}}_{E}-\overline{\mathbf{G}}_{E}^{\mathrm{(near)}}(\mathbf{r},\mathbf{r}')\Big]
\cdot\Jimp(\mathbf{r}')\,d^{3}r'
\;+\;
\int_{\Omega_n}
\overline{\mathbf{G}}_{E}^{\mathrm{(near)}}(\mathbf{r},\mathbf{r}')\cdot\Jimp(\mathbf{r}')\,d^{3}r',
\end{equation}
where \(\overline{\mathbf{G}}_{E}^{\mathrm{(near)}}\) matches the first few terms of
Eq.~\eqref{eq:ch2.dyadic-singular-expansion} about \(\mathbf{r}\) and the integral of
\(\overline{\mathbf{G}}_{E}^{\mathrm{(near)}}\) is evaluated analytically
\cite{harrington2001}. Equation~\eqref{eq:ch2.near-singular-subtraction}
does not alter the continuum operator; it stabilizes numerical approximation of
\(\mathcal{K}_{\mathrm{cpt}}\) in Eq.~\eqref{eq:ch2.solution-singular-compact-split}.
Scalar singularities are weaker than dyadic singularities: the monopole kernel diverges only
as $R^{-1}$, which is barely non-integrable in three dimensions, whereas the dyadic kernel
carries an $R^{-3}$ pole that breaks naive integration more severely \cite{jackson1999}. This
distinction matters computationally because scalar potential formulations can sometimes use
simpler self-term corrections, while vector MoM codes must implement full static--dynamic
dyadic subtraction \cite{harrington2001}. The scalar Green function satisfies, in distributional sense,
\begin{equation}
\label{eq:ch2.scalar-coincidence-limit}
\int_{B_{\varepsilon}(\mathbf{r})}
G_0(\mathbf{r},\mathbf{r}')\,d^{3}r'
=
\frac{\varepsilon}{2}
\;+\;\mathcal{O}(\varepsilon^{3}),
\end{equation}
for small balls \(B_{\varepsilon}(\mathbf{r})\) in \(\mathbb{R}^{3}\), reflecting the
\(R^{-1}\) integrability borderline in three dimensions \cite{jackson1999}. Equation
\eqref{eq:ch2.scalar-coincidence-limit} explains why monopole self-terms in scalar
potential formulations require either finite source extent or explicit regularization;
vector dyads inherit stronger \(R^{-3}\) divergence captured by
Eq.~\eqref{eq:ch2.static-dyadic-limit} \cite{stratton1941}.
Equations~\eqref{eq:ch2.finite-part-convolution}--\eqref{eq:ch2.mom-self-term-regularized}
require patchwise smooth \(\Jimp\), Lipschitz \(\Omega\), and fixed \(\omega\) in the
passive linear band. They fail for point dipoles modeled as \(\delta\)-supported
\(\Jimp\) without accompanying finite-part renormalization, and for non-Lipschitz
edges where \(\HCurl\) membership of \(\E\) requires additional edge singularities not captured by Eq.~\eqref{eq:ch2.dyadic-singular-expansion}
\cite{stratton1941}. Lossy media with complex \(k_0\) shift singular
expansion coefficients but not the order of the \(R^{-3}\) pole; Subsection~\ref{sec:ch232}
does not re-tabulate complex-\(k\) subtractors \cite{stratton1941}. Regularized convolutions
now make $\mathcal{S}_\omega[\Jimp]$ well-defined when observation $\mathbf{r}$ approaches source
support; retarded time-domain support fixes \emph{when} those spatial kernels may act, completing
the causal picture begun in the harmonic dyads with $\mathbf{r}$ as observation and $\mathbf{r}'$
as source \cite{stratton1941}.
\subsection{Retardation and causal propagation}
\label{sec:ch233}
Subsection~\ref{sec:ch233} supplies the definitive account of \emph{retarded} time-domain
Green operators, their causal support, and the kernel-level upgrade of the resolvent
equivalence Eq.~\eqref{eq:ch2.resolvent-green-equivalence} and harmonic theorem
Eq.~\eqref{eq:ch2.harmonic-equivalence-theorem} from Subsection~\ref{sec:ch223}
\cite{stratton1941}. Outgoing causality and Laplace--Fourier transforms are presupposed
through Eqs.~\eqref{eq:ch2.causal-operator-support}--\eqref{eq:ch2.fourier-phasor-identification};
the display below constructs the Green kernels that realize those statements.
Harmonic formulas Eqs.~\eqref{eq:ch2.free-space-scalar-green} and~\eqref{eq:ch2.free-space-dyadic-green}
are cited; they appear as $s=-\jj\omega$ evaluations of causal Laplace
transforms under Sommerfeld outgoing selection at infinity \cite{stratton1941}.
% PHYSICAL-MEMORY:sec:ch233
\begin{physicalvignette}[Rotating scatterer on a lighthouse-style beam chopper]
Consider Rotating scatterer on a lighthouse-style beam chopper. Before any coefficient table is printed, technicians record frequency, cable loss, and the mechanical outline that fixes electrical size. The first walk-through of the range is always tactile: flange bolts, probe standoff, turntable boresight, and the wattmeter reference on the receive chain \cite{hansen1988}.
On this hardware the near zone is visually impressive---standing reactive structure, rapid phase rotation, sometimes a visible VSWR ripple at the feed---yet the export question is measured only after outgoing selection on a large sphere \cite{harrington2001}. What the camera sees in the chamber is not automatically what the far-zone pattern claims \cite{stratton1941}.
Engineers therefore learn to separate three readings: what the feed can synthesize on declared support, what propagating content survives on the radiative range, and what a finite instrument can recover at declared tolerance \cite{devaney2004}. Confusing those readings produces the rank gaps and watt failures that later sections quantify \cite{pal2026}.
Rayleigh and Stratton treated such ranges as boundary-value laboratories long before digital near-field transforms became routine \cite{stratton1941}. Harrington's method-of-moments culture reinforced the same discipline: metal outline first, harmonic headline second \cite{harrington2001}.
Subsection~\ref{sec:ch233} states the formal export condition that this range scene presupposes \cite{stratton1941}.
\end{physicalvignette}
Let $\mathbf{j}_{\mathrm{imp}}(\mathbf{r},t)$ be causal with
\(\operatorname{supp}_{(\mathbf{r},t)}\subseteq\Omega_s\times[0,\infty)\),
\(\Omega_s\Subset\Omega\), and let \(\mathcal{B}_\Gamma\), \(\varepsilon\), and
\(\mu\) be time-independent. Initial data are homogeneous,
\(\mathbf{U}_0=\mathbf{0}\), unless noted for switch-on illustrations. Vacuum
regions use \(c_0=1/\sqrt{\mu_0\varepsilon_0}\) and \(k_0=\omega/c_0\) from
Eq.~\eqref{eq:ch2.helmholtz-symbol} \cite{stratton1941}. The equivalence class
excludes advanced kernels satisfying Eq.~\eqref{eq:ch2.advanced-kernel-violation}.
The causal scalar Green function for the wave operator
\((\nabla^{2}-\varepsilon_0\mu_0\,\partial_t^{2})\) in vacuum satisfies the
inhomogeneous wave equation below, with the time argument \(t-t'\) measuring elapsed
interval between source event and field observation (seconds, SI).
\begin{equation}
\label{eq:ch2.retarded-scalar-wave-operator}
\Big(\nabla^{2}-\varepsilon_0\mu_0\,\partial_t^{2}\Big)\,
G(\mathbf{r},\mathbf{r}',t-t')
=
-\delta(\mathbf{r}-\mathbf{r}')\,\delta(t-t'),
\end{equation}
with retarded causality $G=\mathbf{0}$ for $t-t'<0$: no signal precedes the light cone
\cite{stratton1941}. Time-harmonic kernels inherit this support; acausal branches appear
as standing-wave contamination in open domains \cite{harrington2001}.
The free-space retarded kernel concentrates its support on the light cone
\(t-t'=R/c_0\), where \(c_0=1/\sqrt{\mu_0\varepsilon_0}\) is the vacuum speed of light
(m\,s\(^{-1}\)) and \(R=|\mathbf{r}-\mathbf{r}'|\) is the geometric separation in metres.
The free-space outgoing solution is
\begin{equation}
\label{eq:ch2.retarded-scalar-green-time}
G_0(\mathbf{r},\mathbf{r}',t-t')
=
\frac{\delta\!\big(t-t'-R/c_0\big)}{4\pi R},
\qquad
R=|\mathbf{r}-\mathbf{r}'|,
\end{equation}
which transports impulse support on the light cone \(t-t'=R/c_0\)
\cite{stratton1941}. Equation~\eqref{eq:ch2.retarded-scalar-green-time} is the
time-domain counterpart of Eq.~\eqref{eq:ch2.free-space-scalar-green}: the locked
phasor field \(\exp(-\jj\omega t)\) selects the \(s=-\jj\omega\) Laplace transform
of Eq.~\eqref{eq:ch2.retarded-scalar-green-time} on the outgoing branch
\cite{harrington2001}. Physically, Eq.~\eqref{eq:ch2.retarded-scalar-green-time}
states that a unit impulse at \((\mathbf{r}',t')\) influences \(\mathbf{r}\) only
after travel time \(R/c_0\), enforcing finite propagation speed and
Eq.~\eqref{eq:ch2.causal-support-constraint} on potentials generated by scalar
convolution \cite{stratton1941}.
The retarded electric dyadic \(\overline{\mathbf{G}}_{E}(\mathbf{r},\mathbf{r}',t-t')\)
upgrades Eq.~\eqref{eq:ch2.electric-dyadic-convolution} to the time domain.
Time differentiation of the causal kernel enforces that the electric field responds only
after the impressed current has acted for a retarded interval; the convolution below is
the dyadic analogue of the scalar retarded integral.
The causal Maxwell field is
\begin{equation}
\label{eq:ch2.retarded-dyadic-convolution}
\E(\mathbf{r},t)
=
\E_{\mathrm{hom}}(\mathbf{r},t)
-\mu_0\,\partial_t\int_{0}^{t}\!\!\int_{\Omega}
\overline{\mathbf{G}}_{E}(\mathbf{r},\mathbf{r}',t-t')\cdot
\mathbf{j}_{\mathrm{imp}}(\mathbf{r}',t')\,d^{3}r'\,dt',
\end{equation}
which specializes Eq.~\eqref{eq:ch2.retarded-convolution-form} once the dyadic kernel
is identified \cite{stratton1941}. Retarded causality requires
\begin{equation}
\label{eq:ch2.retarded-dyadic-support}
\overline{\mathbf{G}}_{E}(\mathbf{r},\mathbf{r}',t-t')
=
\mathbf{0}\quad\text{for}\ t-t'<0,
\end{equation}
analogous to Eq.~\eqref{eq:ch2.retarded-kernel-support} \cite{stratton1941}. For
\(\mathbf{r}\notin\Omega_s\), Eq.~\eqref{eq:ch2.retarded-dyadic-convolution} produces
\(\E(\mathbf{r},t)=\mathbf{0}\) until the earliest retarded arrival time
\(t_{\mathrm{ret}}(\mathbf{r})\), defined as the infimum over all source events
\((\mathbf{r}',t')\) in the support of \(\mathbf{j}_{\mathrm{imp}}\) of the sum
\(t'+R/c_0\):
\begin{equation}
\label{eq:ch2.light-cone-arrival-time}
t_{\mathrm{ret}}(\mathbf{r})
=
\inf_{(\mathbf{r}',t')\in\operatorname{supp}\mathbf{j}_{\mathrm{imp}}}
\big\{t'\,+\,R/c_0\big\},
\end{equation}
which is the geometric light-cone support condition underlying
Eq.~\eqref{eq:ch2.causal-support-constraint} \cite{jackson1999}. The
harmonic dyad Eq.~\eqref{eq:ch2.free-space-dyadic-green} is not repeated here; it is
the \(s=-\jj\omega\) symbol of \(\overline{\mathbf{G}}_{E,0}\) under
Eqs.~\eqref{eq:ch2.equivalence-assumptions} and outgoing radiation condition
\cite{harrington2001}.
\begin{figure}[t]
\centering
\ChOneFigBlock{%
\begin{tikzpicture}[ch1fig, node distance=5mm and 7mm]
 \node[ch1box] (src) {$(\mathbf{r}',t')$\\source event};
 \node[ch1box, right=16mm of src] (cone) {Light cone\\$t-t'=R/c_0$};
 \node[ch1box, right=16mm of cone] (obs) {$(\mathbf{r},t)$\\observation};
 \node[ch1box, below=5mm of cone] (adv) {Advanced\\excluded};
 \draw[ch1flow] (src) -- (cone) -- (obs);
 \draw[ch1flow,dashed] (src) to[bend right=20] (adv);
\end{tikzpicture}%
}
\caption{Retarded support Eq.~\eqref{eq:ch2.retarded-dyadic-support}: field at
\((\mathbf{r},t)\) depends only on sources with \(t'\le t-R/c_0\); advanced
coupling is excluded.}
\label{fig:ch2.retarded-light-cone}
\end{figure}
~\eqref{eq:ch2.retarded-dyadic-convolution}.
Subsection~\ref{sec:ch232} treated coincidence in \(\mathbf{r}\); the display below
completes the causal picture by fixing the temporal ordering enforced by retarded support
\cite{stratton1941}. Regularized self-terms from
Eq.~\eqref{eq:ch2.mom-self-term-regularized} apply to harmonic convolutions; time
stepping discretizes Eq.~\eqref{eq:ch2.retarded-dyadic-convolution} with the same
kernel class \cite{harrington2001}. Let \(\widetilde{\overline{\mathbf{G}}}_{E}(\mathbf{r},\mathbf{r}',s)\) denote the
Laplace transform of \(\overline{\mathbf{G}}_{E}(\mathbf{r},\mathbf{r}',t)\) in
\(t-t'\) with \(\Re\{s\}\) large enough for convergence.
The complex variable \(s\) (s\(^{-1}\)) encodes temporal frequency; evaluation at
\(s=-\jj\omega\) recovers the harmonic dyad used throughout the phasor convention.
Causal support
Eq.~\eqref{eq:ch2.retarded-dyadic-support} implies
\(\widetilde{\overline{\mathbf{G}}}_{E}\) is analytic in the right half-plane and,
under outgoing radiation condition,
\begin{equation}
\label{eq:ch2.retarded-laplace-harmonic-link}
\overline{\mathbf{G}}_{E,\omega}(\mathbf{r},\mathbf{r}')
=
\eval{\widetilde{\overline{\mathbf{G}}}_{E}(\mathbf{r},\mathbf{r}',s)}_{s=-\jj\omega},
\end{equation}
which is the kernel form of Eq.~\eqref{eq:ch2.resolvent-green-equivalence}
\cite{stratton1941}. Applying the Laplace transform to
Eq.~\eqref{eq:ch2.retarded-dyadic-convolution} with causal
\(\mathbf{j}_{\mathrm{imp}}\) yields
\begin{equation}
\label{eq:ch2.laplace-dyadic-convolution}
\widetilde{\E}(\mathbf{r},s)
=
\widetilde{\E}_{\mathrm{hom}}(\mathbf{r},s)
-\mu_0 s\int_{\Omega}
\widetilde{\overline{\mathbf{G}}}_{E}(\mathbf{r},\mathbf{r}',s)\cdot
\widetilde{\mathbf{j}}(\mathbf{r}',s)\,d^{3}r',
\end{equation}
with \(\widetilde{\mathbf{j}}\) from Eq.~\eqref{eq:ch2.laplace-source-transform}
\cite{stratton1941}. On \(s=-\jj\omega\) under Eqs.~\eqref{eq:ch2.equivalence-assumptions},
Eq.~\eqref{eq:ch2.laplace-dyadic-convolution} reduces to
Eq.~\eqref{eq:ch2.electric-dyadic-convolution} with
\(\overline{\mathbf{G}}_{E,\omega}\) from Eq.~\eqref{eq:ch2.retarded-laplace-harmonic-link}
and \(\Jimp=\eval{\widetilde{\mathbf{j}}(\mathbf{r},s)}_{s=-\jj\omega}\)
\cite{harrington2001}. The Chapter~\ref{ch:02} causal propagation theorem at kernel
level is
\begin{equation}
\label{eq:ch2.retarded-harmonic-kernel-theorem}
\mathcal{S}_\omega[\Jimp]
=
\eval{\mathcal{S}_t[\mathbf{j}_{\mathrm{imp}}](\mathbf{r},t)}_{\text{harmonic at }\omega}
\quad\Longleftrightarrow\quad
\overline{\mathbf{G}}_{E,\omega}
=
\eval{\widetilde{\overline{\mathbf{G}}}_{E}}_{s=-\jj\omega}
\end{equation}
for the same \(\mathcal{B}_\Gamma\) and outgoing branch, with \(\mathcal{S}_t\) the
time-domain map Eq.~\eqref{eq:ch2.time-domain-solution-map} \cite{stratton1941}.
Equation~\eqref{eq:ch2.retarded-harmonic-kernel-theorem} does not replace
Eq.~\eqref{eq:ch2.harmonic-equivalence-theorem}; it identifies the Green operators
whose resolvent relation Eq.~\eqref{eq:ch2.resolvent-green-equivalence} implements
\cite{harrington2001}.
Advanced kernels and branch failure. Advanced dyads satisfying Eq.~\eqref{eq:ch2.advanced-kernel-violation} produce
Laplace symbols that cannot equal
Eq.~\eqref{eq:ch2.retarded-laplace-harmonic-link} for any causal
\(\mathbf{j}_{\mathrm{imp}}\) \cite{stratton1941}. Time-harmonic solvers that
implicitly use incoming Green functions therefore fail
Eq.~\eqref{eq:ch2.retarded-harmonic-kernel-theorem} even when interior Maxwell rows
vanish \cite{harrington2001}. Exterior problems require outgoing radiation condition consistent
with Section~\ref{sec:ch24}; interior cavity problems may require loss or quotienting
as in Eq.~\eqref{eq:ch2.coercivity-quotient} before a causal Laplace contour can be
closed \cite{stratton1941}. For finite-bandwidth excitation, the harmonic slice at fixed \(\omega\) is insufficient.
Causal fields superpose as
\begin{equation}
\label{eq:ch2.broadband-retarded-synthesis}
\E(\mathbf{r},t)
=
\int_{-\infty}^{\infty}
\widetilde{\E}(\mathbf{r},-\jj\omega)\,e^{-\jj\omega t}\,\frac{d\omega}{2\pi},
\end{equation}
with \(\widetilde{\E}(\mathbf{r},-\jj\omega)\) the Fourier transform induced by
Eq.~\eqref{eq:ch2.laplace-dyadic-convolution} on causal data
\cite{stratton1941}. Equation~\eqref{eq:ch2.broadband-retarded-synthesis}
is the inverse of the single-frequency reduction used in the phasor convention; dispersive
\(\varepsilon(\omega)\), \(\mu(\omega)\) modify
\(\widetilde{\overline{\mathbf{G}}}_{E}(\mathbf{r},\mathbf{r}',s)\) and are deferred
beyond the locked single-\(\omega\) harmonic band except where loss enters as
\(\Im\{k_0\}\neq 0\) \cite{harrington2001}. Within the harmonic band,
Eq.~\eqref{eq:ch2.steady-state-harmonic-limit} describes the long-time limit after
switch-on when outgoing radiation or damping removes transients
\cite{stratton1941}.
Because \(\Pi_{\mathrm{rad}}\) commutes with harmonic extraction under the ordering
assumed in Subsection~\ref{sec:ch223},
\begin{equation}
\label{eq:ch2.retarded-radiation-kernel-composition}
\mathcal{R}_t[\mathbf{j}_{\mathrm{imp}}]
=
\Pi_{\mathrm{rad}}\circ\mathcal{S}_t[\mathbf{j}_{\mathrm{imp}}],
\qquad
\mathcal{R}_\omega[\Jimp]
=
\Pi_{\mathrm{rad}}\circ\mathcal{S}_\omega[\Jimp],
\end{equation}
with \(\mathcal{R}_\omega\) from Eq.~\eqref{eq:ch2.radiation-time-frequency-equiv}
\cite{stratton1941}. Causality does not promote evanescent content to
radiative export: Eq.~\eqref{eq:ch2.regime-operator-screen} remains the transport criterion
after time-domain evolution \cite{stratton1941}. Retarded propagation ensures support
constraints only; spectral and regime qualifications from Subsection~\ref{sec:ch224} are
additional \cite{harrington2001}.
Equations~\eqref{eq:ch2.retarded-dyadic-convolution}--\eqref{eq:ch2.retarded-harmonic-kernel-theorem}
require Eqs.~\eqref{eq:ch2.equivalence-assumptions}, outgoing or passive completion on
exterior domains, and \(\mathbf{j}_{\mathrm{imp}}\) in the causal class of
Subsection~\ref{sec:ch221}. They fail for time-varying \(\mathcal{B}_\Gamma(t)\),
moving media, and nonlinear sources. Coincidence and self-term evaluation of
\(\overline{\mathbf{G}}_{E,\omega}\) remains governed by Subsection~\ref{sec:ch232}
even when Eq.~\eqref{eq:ch2.retarded-dyadic-convolution} is discretized in time
\cite{harrington2001}. Retarded support is therefore the temporal counterpart of the spatial
singular split in Subsection~\ref{sec:ch232}: causality fixes when propagation may act, while
regularization fixes how strongly a source couples to itself at coincidence
\cite{stratton1941}. Harmonic kernels $G$ and $\overline{\mathbf{G}}_{E}$ are therefore
retarded response functions in space and time; constructive inversion on composite domains asks
which resolvent branch those responses realize when interfaces and exterior completions are
assembled \cite{stratton1941}.
\subsection{Inversion and solvability}
\label{sec:ch234}
% HARRINGTON-OPEN:sec:ch234
A standard-gain horn on a compact-range turntable fixes the laboratory scene for `Inversion and solvability'': frequency, probe standoff, and wattmeter closure on $\Gamma_R$ precede the formal displays below \cite{hansen1988}. Maxwell's retarded export to $\Gamma_R$ remains the governing constraint \cite{maxwell1873}.
Subsection~\ref{sec:ch234} supplies the definitive account of constructive inversion and
solvability on single- and multi-region domains with interface data from
Eq.~\eqref{eq:ch1.interface-conditions} \cite{stratton1941}. Weak well-posedness on
bounded quotients through Eqs.~\eqref{eq:ch2.lax-milgram-wellposed} and~\eqref{eq:ch2.coercivity-quotient}
from Subsection~\ref{sec:ch214} is presupposed; Lax--Milgram and Friedrichs arguments
are not repeated. Outgoing uniqueness on exterior domains is finalized in
Section~\ref{sec:ch243}; here the resolvent branch and its Green realization are
stated as hypotheses \cite{stratton1941}. Material Helmholtz symbols
Eqs.~\eqref{eq:ch2.material-helmholtz-electric}--\eqref{eq:ch2.material-helmholtz-magnetic}
and trace spaces Eqs.~\eqref{eq:ch2.interface-trace-space}--\eqref{eq:ch2.domain-refinement-L}
from Sections~\ref{sec:ch211}--\ref{sec:ch222} are fixed input classes.
% PHYSICAL-MEMORY:sec:ch234
\begin{physicalvignette}[Dual-polarization horn pair on a compact range]
Consider Dual-polarization horn pair on a compact range. Before any coefficient table is printed, technicians record frequency, cable loss, and the mechanical outline that fixes electrical size. The first walk-through of the range is always tactile: flange bolts, probe standoff, turntable boresight, and the wattmeter reference on the receive chain \cite{hansen1988}.
On this hardware the near zone is visually impressive---standing reactive structure, rapid phase rotation, sometimes a visible VSWR ripple at the feed---yet the export question is measured only after outgoing selection on a large sphere \cite{harrington2001}. What the camera sees in the chamber is not automatically what the far-zone pattern claims \cite{stratton1941}.
Engineers therefore learn to separate three readings: what the feed can synthesize on declared support, what propagating content survives on the radiative range, and what a finite instrument can recover at declared tolerance \cite{devaney2004}. Confusing those readings produces the rank gaps and watt failures that later sections quantify \cite{pal2026}.
Rayleigh and Stratton treated such ranges as boundary-value laboratories long before digital near-field transforms became routine \cite{stratton1941}. Harrington's method-of-moments culture reinforced the same discipline: metal outline first, harmonic headline second \cite{harrington2001}.
Subsection~\ref{sec:ch234} states the formal export condition that this range scene presupposes \cite{stratton1941}.
\end{physicalvignette}
Let \(\Omega=\cup_{j=1}^{N_p}\Omega_j\) be a finite union of Lipschitz patches with
pairwise interfaces \(\Sigma_{ab}=\overline{\Omega_a}\cap\overline{\Omega_b}\),
piecewise constants \(\varepsilon_j,\mu_j>0\) on \(\Omega_j\), and outer boundary
\(\Gamma=\partial\Omega\). Impressed data satisfy
\(\Jimp\in\mathcal{J}_\Omega\), charge continuity
Eq.~\eqref{eq:ch1.charge-continuity-harmonic}, and boundary class
\(\mathcal{B}_\Gamma\) compatible with Eq.~\eqref{eq:ch2.admissible-field-space}.
Regularized convolutions follow Subsection~\ref{sec:ch232}; retarded causality follows
Subsection~\ref{sec:ch233} \cite{stratton1941}.
Composite domain and patchwise inversion. Inversion separates invertible propagation from compact perturbations and null directions
\cite{stratton1941}. Ill-posed spherical near-field campaigns often fail here: the
export map is well defined, but the inverse amplifies noise along small singular values
\cite{hansen1988}.
On a composite domain the impressed-source region is partitioned into material patches
\(\Omega_j\) separated by interfaces \(\Sigma_{ab}\); each patch carries its own wavenumber
\(k_j=\omega\sqrt{\mu_j\varepsilon_j}\) (rad\,m\(^{-1}\)).
Decompose the solution map as
\begin{equation}
\label{eq:ch2.composite-domain-decomposition}
\Omega
=
\bigsqcup_{j=1}^{N_p}\Omega_j
\ \cup\
\bigsqcup_{a<b}\Sigma_{ab},
\qquad
k_j=\omega\sqrt{\mu_j\varepsilon_j},
\end{equation}
with wavenumbers \(k_j\) from Eq.~\eqref{eq:ch2.helmholtz-symbol} on each patch
\cite{stratton1941}. On \(\Omega_j\), the electric field satisfies the patchwise
Helmholtz drive of Eq.~\eqref{eq:ch2.material-helmholtz-electric} and is represented
by the regularized dyadic convolution stated below; the patch kernel
\(\overline{\mathbf{G}}_{E}^{(j)}\) incorporates local \(\varepsilon_j,\mu_j\) and the
finite-part prescription of Subsection~\ref{sec:ch232} when observation approaches the
source support:
\begin{equation}
\label{eq:ch2.patchwise-dyadic-inversion}
\E(\mathbf{r})
=
\E_{\mathrm{hom}}(\mathbf{r})
+\jj\omega\mu_j
\int_{\Omega_j}
\overline{\mathbf{G}}_{E}^{(j)}(\mathbf{r},\mathbf{r}')\cdot\Jimp(\mathbf{r}')\,d^{3}r',
\qquad
\mathbf{r}\in\Omega_j,
\end{equation}
where \(\overline{\mathbf{G}}_{E}^{(j)}\) satisfies
Eq.~\eqref{eq:ch2.material-dyadic-defining-operator} with patch coefficients and
finite-part evaluation when \(\mathbf{r}\in\Omega_j\) approaches
\(\operatorname{supp}\Jimp\) \cite{harrington2001}. In Eq.~\eqref{eq:ch2.patchwise-dyadic-inversion},
$\mathbf{r}$ is the observation point in patch $\Omega_j$ and $\mathbf{r}'$ the source point
integrated over the impressed support; $\overline{\mathbf{G}}_{E}^{(j)}(\mathbf{r},\mathbf{r}')$
is the patchwise outgoing response function \cite{stratton1941}. Equation
\eqref{eq:ch2.patchwise-dyadic-inversion} is the constructive content of
Eq.~\eqref{eq:ch2.solution-map-dyadic-realization} on uniform subdomains; it does
not yet enforce interface jumps across \(\Sigma_{ab}\) \cite{stratton1941}.
Fields in adjacent patches are linked by tangential continuity and normal-flux jumps
imported from Eq.~\eqref{eq:ch1.interface-conditions} \cite{stratton1941}.
The unit normal \(\hat{\mathbf{n}}_{ab}\) on interface \(\Sigma_{ab}\) points from patch
\(a\) into patch \(b\); tangential \(\E\) and \(\H\) must match across the interface so
that the composite field remains a single physical solution.
In
operator form,
\begin{equation}
\label{eq:ch2.interface-dyadic-continuity}
\hat{\mathbf{n}}_{ab}\times\E^{(a)}\big|_{\Sigma_{ab}}
=
\hat{\mathbf{n}}_{ab}\times\E^{(b)}\big|_{\Sigma_{ab}},
\qquad
\hat{\mathbf{n}}_{ab}\times\H^{(a)}\big|_{\Sigma_{ab}}
=
\hat{\mathbf{n}}_{ab}\times\H^{(b)}\big|_{\Sigma_{ab}},
\end{equation}
with \(\E^{(j)}\) from Eq.~\eqref{eq:ch2.patchwise-dyadic-inversion} and
\(\H^{(j)}\) from Eq.~\eqref{eq:ch2.magnetic-from-electric-dyad} when valid
\cite{harrington2001}. Equations~\eqref{eq:ch2.interface-dyadic-continuity} replace
free-space Eq.~\eqref{eq:ch2.free-space-dyadic-green} by patch Green functions
\(\overline{\mathbf{G}}_{E}^{(j)}\) plus interface unknowns on
\(\mathcal{T}_{\Sigma_{ab}}\) from Eq.~\eqref{eq:ch2.interface-trace-space}
\cite{stratton1941}. The assembled multi-region system is
\begin{equation}
\label{eq:ch2.multi-region-inversion-system}
\begin{bmatrix}
\mathcal{I}+\mathcal{K}_{11} & \cdots & \mathcal{K}_{1N_p}\\
\vdots & \ddots & \vdots\\
\mathcal{K}_{N_p1} & \cdots & \mathcal{I}+\mathcal{K}_{N_pN_p}
\end{bmatrix}
\begin{bmatrix}\mathbf{x}_1\\ \vdots\\ \mathbf{x}_{N_p}\end{bmatrix}
=
\begin{bmatrix}\mathbf{b}_1\\ \vdots\\ \mathbf{b}_{N_p}\end{bmatrix},
\end{equation}
where \(\mathbf{x}_j\) collects tangential field or equivalent surface currents on
patch boundaries, \(\mathcal{K}_{jk}\) are dyadic layer potentials generated by
\(\overline{\mathbf{G}}_{E}^{(k)}\), and \(\mathbf{b}_j\) encodes impressed sources
and outer \(\mathcal{B}_\Gamma\) data \cite{harrington2001}.
Equation~\eqref{eq:ch2.multi-region-inversion-system} is the constructive statement of
multi-region solvability: invertibility of the block operator is the
Green-function analogue of Eq.~\eqref{eq:ch2.lax-milgram-wellposed} on composite
domains \cite{stratton1941}.
\begin{figure}[t]
\centering
\ChOneFigBlock{%
\begin{tikzpicture}[ch1fig, node distance=5mm and 7mm]
 \node[ch1box] (p1) {$\Omega_1$\\$k_1$};
 \node[ch1box, right=7mm of p1] (sig) {$\Sigma_{12}$};
 \node[ch1box, right=7mm of sig] (p2) {$\Omega_2$\\$k_2$};
 \node[ch1box, below=5mm of sig] (g) {$\overline{\mathbf{G}}_{E}^{(j)}$};
 \node[ch1box, below=5mm of p2] (bc) {$\mathcal{B}_\Gamma$};
 \draw[ch1flow] (p1) -- (sig) -- (p2);
 \draw[ch1flow] (g.east) -- (g.east |- p1.west);
 \draw[ch1flow] (g.east |- p1.west) -- (p1.west);
 \draw[ch1flow] (g.east) -- (g.east |- p2.west);
 \draw[ch1flow] (g.east |- p2.west) -- (p2.west);
 \draw[ch1flow] (bc.south) -- (bc.south |- p2.north);
 \draw[ch1flow] (bc.south |- p2.north) -- (p2.north);
\end{tikzpicture}%
}
\caption{Multi-region inversion Eq.~\eqref{eq:ch2.multi-region-inversion-system}:
patchwise dyadic kernels coupled by interface continuity
Eq.~\eqref{eq:ch2.interface-dyadic-continuity} and exterior
\(\mathcal{B}_\Gamma\).}
\label{fig:ch2.multi-region-green-assembly}
\end{figure}
Subsection~\ref{sec:ch232} supplies
self-term regularization on each patch; Subsection~\ref{sec:ch233} supplies causal
time marching when Eq.~\eqref{eq:ch2.multi-region-inversion-system} is embedded in
\(\mathcal{S}_t\) \cite{harrington2001}. On a bounded Lipschitz domain with passive \(\mathcal{B}_\Gamma\) satisfying
Eq.~\eqref{eq:ch2.passive-impedance-class}, coercivity on the quotient
\(\mathcal{U}_\omega^{0}/\ker\mathcal{L}_\omega\) from
Eq.~\eqref{eq:ch2.coercivity-quotient} implies existence of a unique weak solution
Eq.~\eqref{eq:ch2.lax-milgram-wellposed} \cite{stratton1941}. Green-function
inversion is equivalent when the regularized operator
\begin{equation}
\label{eq:ch2.green-bvp-equivalence}
\mathcal{S}_\omega[\Jimp]
=
\mathbf{u}_\omega
\quad\Longleftrightarrow\quad
\mathcal{L}_\omega[\mathbf{u}_\omega]
=
\begin{bmatrix}\mathbf{s}_\omega(\Jimp,\rho)\\ \mathbf{g}_\Gamma\end{bmatrix},
\end{equation}
with \(\mathbf{u}_\omega=(\E,\H)\) obtained from
Eqs.~\eqref{eq:ch2.patchwise-dyadic-inversion}--\eqref{eq:ch2.mom-self-term-regularized}
on simply connected patches and from
Eq.~\eqref{eq:ch2.multi-region-inversion-system} otherwise
\cite{harrington2001}. Equation~\eqref{eq:ch2.green-bvp-equivalence}
expresses Eq.~\eqref{eq:ch2.solution-map-resolvent} with constructive kernels; stability
is controlled by Eq.~\eqref{eq:ch2.solution-map-stability} and resolvent bound
Eq.~\eqref{eq:ch2.resolvent-bounded} on the physical branch
\cite{stratton1941}.
On exterior \(\Omega\) with compact \(\operatorname{supp}\Jimp\), existence of
\(\mathcal{S}_\omega\) requires outgoing Green functions whose far zone matches
Section~\ref{sec:ch24} \cite{stratton1941}. Without that selection, the block system
Eq.~\eqref{eq:ch2.multi-region-inversion-system} admits incoming solutions violating
Eq.~\eqref{eq:ch2.retarded-harmonic-kernel-theorem} \cite{harrington2001}. Uniqueness
on the outgoing branch is
\begin{equation}
\label{eq:ch2.exterior-uniqueness-hypothesis}
\mathcal{L}_\omega[\mathbf{u}_\omega]=\mathbf{0}
\ \text{and outgoing}\ \mathcal{B}_\Gamma
\ \Longrightarrow\
\mathbf{u}_\omega=\mathbf{0},
\end{equation}
with \(\mathbf{u}_\omega=(\E,\H)\) the total field phasor pair. Eq.~\eqref{eq:ch2.exterior-uniqueness-hypothesis} is the solvability criterion for
exterior Green inversion \cite{stratton1941}. The split
Eq.~\eqref{eq:ch2.outgoing-nonselfadjoint-split} reminds that exterior problems are
not Friedrichs self-adjoint on all of \(\mathcal{U}_\omega\)
\cite{harrington2001}.
Resonant cavity spectra. Bounded cavities with PEC walls possess discrete spectra \(\{\omega_n\}\).
At a resonant frequency \(\omega=\omega_n\) the Helmholtz operator acquires a
non-trivial null space and the Green resolvent ceases to be bounded on the full
admissible field space:
\begin{equation}
\label{eq:ch2.cavity-resonance-failure}
\ker\mathcal{L}_\omega\neq\{\mathbf{0}\}
\quad\Longrightarrow\quad
\mathcal{L}_\omega^{-1}\ \text{unbounded on}\ \mathcal{U}_\omega,
\end{equation}
and Eq.~\eqref{eq:ch2.green-bvp-equivalence} fails unless quotienting by standing
modes as in Eq.~\eqref{eq:ch2.coercivity-quotient} or introducing loss
\cite{stratton1941}. Poles of the resolvent appear as
\begin{equation}
\label{eq:ch2.resolvent-pole-cavity}
\det\big(\mathcal{I}+\mathcal{K}\big)=0
\quad\Longleftrightarrow\quad
\omega=\omega_n,
\end{equation}
for the coupled operator \(\mathcal{K}\) in
Eq.~\eqref{eq:ch2.multi-region-inversion-system} on closed resonant structures
\cite{stratton1941}. Impressed \(\Jimp\) at resonance requires compatibility with
\(\ker\mathcal{L}_\omega\); otherwise amplitudes diverge even when Green formulas
are formally applied \cite{harrington2001}. Incoming exterior excitation produces a
distinct failure: Eq.~\eqref{eq:ch2.exterior-uniqueness-hypothesis} is violated
without additional incident-field subtraction in Section~\ref{sec:ch251}
\cite{stratton1941}.
Numerical equivalence. Finite-element stiffness systems and boundary-element impedance matrices are
Galerkin discretizations of Eq.~\eqref{eq:ch2.maxwell-weak-form} and
Eq.~\eqref{eq:ch2.multi-region-inversion-system}, respectively:
\begin{equation}
\label{eq:ch2.discrete-green-equivalence}
\mathbf{Z}\mathbf{I}=\mathbf{V}
\quad\Longleftrightarrow\quad
\mathcal{K}\mathbf{x}=\mathbf{b},
\end{equation}
when basis functions match trace spaces and self-terms implement
Eq.~\eqref{eq:ch2.mom-self-term-regularized} \cite{harrington2001}. Equation
\eqref{eq:ch2.discrete-green-equivalence} is the same inversion problem in matrix form.
Ill-conditioning near
Eq.~\eqref{eq:ch2.resolvent-pole-cavity} mirrors practical difficulty at cavity
modes and near-spurious resonances in open-region formulations
\cite{harrington2001}.
Equations~\eqref{eq:ch2.patchwise-dyadic-inversion}--\eqref{eq:ch2.multi-region-inversion-system}
require Lipschitz interfaces, passive linear media on each patch, and fixed
\(\omega\) in the harmonic band. They fail for time-varying boundaries, nonlinear
constitutive laws, and anisotropic magneto-electric coupling without retensoring
\(\overline{\mathbf{G}}_{E}^{(j)}\) \cite{stratton1941}. Point sources without
Subsection~\ref{sec:ch232} regularization produce formal but non-admissible fields;
\(\Jimp\) must remain in \(\mathcal{J}_\Omega\) \cite{harrington2001}. With Subsections~\ref{sec:ch231}--\ref{sec:ch234} complete, Section~\ref{sec:ch23}
has constructed scalar, vector, and dyadic Green kernels, regularized their
singular convolutions, tied them to retarded causal propagation, and stated
constructive inversion on composite domains compatible with
Eq.~\eqref{eq:ch2.resolvent-operator}. Sommerfeld outgoing selection in
Section~\ref{sec:ch24} fixes the exterior branch assumed in
Eq.~\eqref{eq:ch2.exterior-uniqueness-hypothesis}; radiation integrals in
Section~\ref{sec:ch25} evaluate far-zone observables from the propagating subspace
qualified by Eq.~\eqref{eq:ch2.regime-operator-screen}; spectral and Fredholm
consequences of Eq.~\eqref{eq:ch2.green-compact-split} continue in
Section~\ref{sec:ch263} \cite{stratton1941}. Green-function inversion recovers field
states from sources under declared completion; whether a desired angular pattern is
source-realizable under finite impressed currents is established in Chapter~\ref{ch:04}
(\emph{Source Realizability and Spectral Truncation}) \cite{devaney2004}.
\section{Radiation Conditions and Uniqueness}
\label{sec:ch24}
Green-function inversion in Section~\ref{sec:ch23} reconstructs fields from sources on compact support. Section~\ref{sec:ch24} imposes outgoing radiation conditions that separate measurable far-zone content from reactive storage \cite{sommerfeld1949}. A dense near-field map without outgoing filtering can fit non-radiating currents \cite{devaney2004}.
% PHYSICAL-MEMORY:sec:ch24
\begin{physicalvignette}[Waveguide TE10 mouth in a flange radiating to half-space]
Consider Waveguide TE10 mouth in a flange radiating to half-space. Before any coefficient table is printed, technicians record frequency, cable loss, and the mechanical outline that fixes electrical size. The first walk-through of the range is always tactile: flange bolts, probe standoff, turntable boresight, and the wattmeter reference on the receive chain \cite{hansen1988}.
On this hardware the near zone is visually impressive---standing reactive structure, rapid phase rotation, sometimes a visible VSWR ripple at the feed---yet the export question is measured only after outgoing selection on a large sphere \cite{harrington2001}. What the camera sees in the chamber is not automatically what the far-zone pattern claims \cite{stratton1941}.
Engineers therefore learn to separate three readings: what the feed can synthesize on declared support, what propagating content survives on the radiative range, and what a finite instrument can recover at declared tolerance \cite{devaney2004}. Confusing those readings produces the rank gaps and watt failures that later sections quantify \cite{pal2026}.
Rayleigh and Stratton treated such ranges as boundary-value laboratories long before digital near-field transforms became routine \cite{stratton1941}. Harrington's method-of-moments culture reinforced the same discipline: metal outline first, harmonic headline second \cite{harrington2001}.
Subsection~\ref{sec:ch24} states the formal export condition that this range scene presupposes \cite{stratton1941}.
\end{physicalvignette}
Maxwell's equations alone do not distinguish fields radiating away from a source from
fields arriving from infinity, yet every laboratory measurement assumes the former.
Antennas are fed by finite sources, compact ranges suppress external illumination, and
receivers attribute detected power to identifiable radiators rather than unexplained
incoming waves \cite{sommerfeld1949}. Section~\ref{sec:ch23} supplied Green
kernels that map impressed currents to fields; on exterior domains those kernels remain
ambiguous until an outgoing branch is fixed at infinity \cite{jackson1999}. Retarded
causality narrows the admissible class but does not fully replace Sommerfeld selection
when harmonic solvers invert $\mathcal{L}_\omega$ at fixed $\omega$
\cite{stratton1941}. The subsections below state outgoing radiation conditions, embed them
as resolvent selection, prove uniqueness on the radiative branch, and record
reconstruction limits when that branch is omitted.
\subsection{Sommerfeld radiation condition}
\label{sec:ch241}
Before Sommerfeld asymptotics are written, the physical ambiguity must be named. Harmonic
Helmholtz solutions on open domains superpose $e^{-\jj k_0 r}/r$ outgoing and
$e^{+\jj k_0 r}/r$ incoming spherical waves with equal standing rights in the volume
equations alone \cite{sommerfeld1949}. A near-field scan or pattern cut on
$\Gamma_R$ cannot distinguish these completions without an explicit outgoing selection: wattmeters
and probes then read mixtures of export, storage, and energy arriving from the computational or
chamber boundary \cite{hansen1988}. Sommerfeld radiation condition is the asymptotic
statement that restores laboratory identification---radiation launched locally and leaving the
domain---before any uniqueness argument or far-zone integral is evaluated
\cite{sommerfeld1949}.
% PHYSICAL-MEMORY:sec:ch241
\begin{physicalvignette}[Microstrip patch array on a finite substrate]
Consider Microstrip patch array on a finite substrate. Before any coefficient table is printed, technicians record frequency, cable loss, and the mechanical outline that fixes electrical size. The first walk-through of the range is always tactile: flange bolts, probe standoff, turntable boresight, and the wattmeter reference on the receive chain \cite{hansen1988}.
On this hardware the near zone is visually impressive---standing reactive structure, rapid phase rotation, sometimes a visible VSWR ripple at the feed---yet the export question is measured only after outgoing selection on a large sphere \cite{harrington2001}. What the camera sees in the chamber is not automatically what the far-zone pattern claims \cite{stratton1941}.
Engineers therefore learn to separate three readings: what the feed can synthesize on declared support, what propagating content survives on the radiative range, and what a finite instrument can recover at declared tolerance \cite{devaney2004}. Confusing those readings produces the rank gaps and watt failures that later sections quantify \cite{pal2026}.
Sommerfeld radiation conditions remove incoming waves from open-region models . Compact-range anechoic chambers approximate the same selection .
Rayleigh and Stratton treated such ranges as boundary-value laboratories long before digital near-field transforms became routine \cite{stratton1941}. Harrington's method-of-moments culture reinforced the same discipline: metal outline first, harmonic headline second \cite{harrington2001}.
Subsection~\ref{sec:ch241} states the formal export condition that this range scene presupposes \cite{stratton1941}.
\end{physicalvignette}
Before any asymptotic formula is written, the range engineer already decides whether the
measurement enclosure behaves as an open exterior domain \cite{hansen1988}.
A compact range places the device under test inside an anechoic chamber whose absorbing
linings are designed to let export escape instead of reflecting back as multipath
\cite{collin2001}. Pattern probes on a large sphere $\Gamma_R$ and shell wattmeters report
tangential field and power only when chamber absorption and computational termination both
approximate free-space outgoing behavior \cite{poynting1884}. Incomplete
absorption produces ripple on azimuth cuts; unstable spherical near-field transforms when
spectra are truncated show the same branch failure at the data-processing stage
\cite{hansen1988}. Until outgoing selection is restored, $\Gamma_R$ meters read cavity
storage and standing-wave content rather than watt-bearing export \cite{maxwell1873}.
Sommerfeld radiation condition below is the asymptotic mathematical statement of what the
range engineer already assumes when attributing a measured pattern to the device under test
rather than to energy arriving from infinity \cite{sommerfeld1949}.
An open antenna measurement is meaningful only if the field on the range can be identified
with radiation from the device under test rather than with energy arriving from infinity.
Sommerfeld outgoing selection enforces that identification: disturbances launched locally
propagate outward and are not returned from the far boundary of the model
\cite{sommerfeld1949}. When the condition is omitted, the same near-field scan or pattern
cut admits formal completions with incoming or standing-wave content, and inversion may
assign equivalent sources that would require excitation mechanisms absent from the physical
problem and therefore lie outside the admissible solution class \cite{devaney2004}.
Free-space kernels
Eqs.~\eqref{eq:ch2.free-space-scalar-green} and~\eqref{eq:ch2.free-space-dyadic-green}
supply the asymptotic templates; outer-boundary classes from Chapter~\ref{ch:01} are
presupposed. Uniqueness proofs and reconstruction limits belong to
Subsections~\ref{sec:ch243} and~\ref{sec:ch244}.
Let \(\Omega_{\mathrm{ext}}\) denote an exterior region surrounding compact \(\operatorname{supp}\Jimp\subset\Omega_s\), with \(\omega>0\) fixed and vacuum
constants \(\varepsilon_0,\mu_0\) in the far zone. Spherical coordinates
\((r,\theta,\phi)\) are centered so that \(\Omega_s\subset B_R(0)\) for some
\(R_0>0\). Fields are phasors under \(\exp(-\jj\omega t)\); outgoing means
\(e^{-\jj k_0 r}/r\) decay with \(k_0=\omega\sqrt{\mu_0\varepsilon_0}\) from
Eq.~\eqref{eq:ch2.helmholtz-symbol} \cite{stratton1941}. Interior interfaces and
finite domains use Eq.~\eqref{eq:ch1.outer-boundary-classes}; Sommerfeld completion
applies on artificial or physical spheres at large \(r\) \cite{harrington2001}.
On $\Gamma_R$ with $R\gg\lambda_0$, a range engineer expects the measured phase to advance as
$k_0 r$ increases---outgoing fronts receding from the enclosure---and the amplitude to fall no
faster than $r^{-1}$ in three dimensions \cite{jackson1999}. That observation is what
Sommerfeld selection encodes at infinity: the scalar template states what every outgoing Helmholtz
response must look like at large distance from a localized source. Phase fronts recede with the
observer; amplitude decay compatible with energy export to infinity is $r^{-1}$ in three
dimensions \cite{jackson1999}. For scalar potentials and Helmholtz kernels, outgoing radiation
requires, with $r=|\mathbf{r}|$ the observation radius from an origin enclosing the source,
$k_0=\omega\sqrt{\mu_0\varepsilon_0}$ the vacuum wavenumber (rad\,m$^{-1}$), and
$A(\theta,\phi,\mathbf{r}')$ the angular amplitude that carries directional information while
$\mathbf{r}'$ labels the fixed source point during the limit $r\to\infty$:
\begin{equation}
\label{eq:ch2.sommerfeld-scalar-asymptotic}
G(\mathbf{r},\mathbf{r}')
=
A(\theta,\phi,\mathbf{r}')\,\frac{e^{-\jj k_0 r}}{r}
\;+\;
\mathcal{O}(r^{-2}),
\qquad
r=|\mathbf{r}|\to\infty,
\end{equation}
The remainder \(\mathcal{O}(r^{-2})\) collects subleading reactive and curvature corrections
that need not vanish at finite range \cite{stratton1941}. Equation
\eqref{eq:ch2.sommerfeld-scalar-asymptotic} is the \ref{ch:02} scalar
Sommerfeld condition: it excludes incoming \(e^{+\jj k_0 r}/r\) and standing
\(e^{\pm\jj k_0 r}/r\) superpositions that degenerate exterior uniqueness
\cite{harrington2001}. The free-space template Eq.~\eqref{eq:ch2.free-space-scalar-green}
satisfies Eq.~\eqref{eq:ch2.sommerfeld-scalar-asymptotic} for fixed \(\mathbf{r}'\)
\cite{stratton1941}. Physically, Eq.~\eqref{eq:ch2.sommerfeld-scalar-asymptotic}
states that monopole-like responses propagate outward with phase fronts
\(k_0 r=\mathrm{const}\) and amplitude decay \(r^{-1}\), the only radial
combination compatible with energy flux to infinity in three dimensions
\cite{jackson1999}.
Vector fields must satisfy a transverse impedance relation at infinity: the electric and
magnetic components that carry export must stand in the same ratio as a plane wave
propagating along \(\hat{\mathbf{r}}\), with vacuum impedance \(\eta_0\) as the proportionality
constant \cite{stratton1941}. For harmonic Maxwell fields in vacuum far from sources,
\begin{equation}
\label{eq:ch2.sommerfeld-field-condition}
\mathbf{r}\times\E
\;+\;
\eta_0\,(\mathbf{r}\times\H)\times\hat{\mathbf{r}}
=
\mathbf{o}(1),
\qquad
r\to\infty,
\end{equation}
with \(\eta_0=\sqrt{\mu_0/\varepsilon_0}\) the vacuum wave impedance,
\(\hat{\mathbf{r}}=\mathbf{r}/r\) the unit radial direction, and \(\mathbf{o}(1)\) denoting
any vector whose Cartesian components vanish in the limit \(r\to\infty\)
\cite{stratton1941}. Physically, Eq.~\eqref{eq:ch2.sommerfeld-field-condition} requires that
the transverse electric field and the transverse magnetic field on a distant sphere lock
phase and magnitude as if energy were leaving the domain along \(\hat{\mathbf{r}}\), rather
than arriving from infinity \cite{harrington2001}. Equation~\eqref{eq:ch2.sommerfeld-field-condition}
is the \ref{ch:02} vector Sommerfeld radiation condition. It implies
transverse radiative fields with asymptotic ratio
\begin{equation}
\label{eq:ch2.sommerfeld-impedance-ratio}
\H_{\mathrm{rad}}
\sim
\eta_0^{-1}\,\hat{\mathbf{r}}\times\E_{\mathrm{rad}},
\qquad
r\to\infty,
\end{equation}
and matches the outgoing-trace condition Eq.~\eqref{eq:ch1.radiation-trace-condition}
from Subsection~\ref{sec:ch01.1.2} when tangential traces are taken on a sphere of
large radius \cite{stratton1941}. Mathematically,
Eq.~\eqref{eq:ch2.sommerfeld-field-condition} selects the subspace on which the
exterior Laplacian/Helmholtz resolvent coincides with retarded Green functions
Eq.~\eqref{eq:ch2.retarded-harmonic-kernel-theorem}; physically, it enforces
power flow outward rather than inward at infinity \cite{harrington2001}.
Far-zone transversality and propagation. Sommerfeld-admissible fields decompose at leading order as
\begin{equation}
\label{eq:ch2.far-zone-transverse-radiation}
\E(\mathbf{r})
=
\E_0(\theta,\phi)\,\frac{e^{-\jj k_0 r}}{r}
\;+\;\mathcal{O}(r^{-2}),
\qquad
\E_0\cdot\hat{\mathbf{r}}=0,
\end{equation}
with \(\E_0\) the radiation amplitude pattern and \(\E_0\cdot\hat{\mathbf{r}}=0\)
the transversality constraint \cite{stratton1941}. Equation
\eqref{eq:ch2.far-zone-transverse-radiation} connects Sommerfeld completion to
\(\Pi_{\mathrm{prop}}\) from Subsection~\ref{sec:ch224}: far-zone content lies on
the propagating sheet with real \(k_r=k_0\) and outgoing phase
\cite{harrington2001}. Evanescent and reactive components may still be large at
finite \(r\) while satisfying Eq.~\eqref{eq:ch2.sommerfeld-field-condition}; regime
qualification Eq.~\eqref{eq:ch2.regime-operator-screen} remains separate from
Sommerfeld selection \cite{stratton1941}.
\begin{figure}[t]
\centering
\ChOneFigBlock{%
\begin{tikzpicture}[ch1fig, node distance=5mm and 7mm]
 \node[ch1box] (src) {Compact\\$\Omega_s$};
 \node[ch1box, right=16mm of src] (near) {Near/reactive\\$kr\lesssim 1$};
 \node[ch1box, right=16mm of near] (far) {Sommerfeld\\$e^{-jk_0 r}/r$};
 \node[ch1box, below=5mm of near] (ev) {$\Pi_{\mathrm{ev}}$};
 \draw[ch1flow] (src) -- (near) -- (far);
 \draw[ch1flow] (near.south) -- (near.south |- ev.north);
 \draw[ch1flow] (near.south |- ev.north) -- (ev.north);
\end{tikzpicture}%
}
\caption{Sommerfeld far zone Eq.~\eqref{eq:ch2.sommerfeld-field-condition}: compact
sources drive near fields with evanescent content; at large \(r\) only outgoing
transverse \(e^{-jk_0 r}/r\) waves satisfy exterior completion.}
\label{fig:ch2.sommerfeld-far-zone}
\end{figure}
Equation~\eqref{eq:ch2.sommerfeld-field-condition} is a local impedance statement at
infinity: transverse \(\E\) and \(\H\) must satisfy the plane-wave ratio
Eq.~\eqref{eq:ch2.sommerfeld-impedance-ratio} with \(\eta_0\) from
Eq.~\eqref{eq:ch1.vacuum-impedance-wavenumber}. Mathematically, it selects the
eigensheet of the exterior Helmholtz resolvent with outgoing radial phase; physically,
it enforces net power flow toward infinity rather than inward from infinity. On a
truncation sphere \(\Gamma_R\) with \(R=8\lambda_0\) enclosing a source with \(k_0 a=1.5\),
Silver--M\"uller surrogate Eq.~\eqref{eq:ch2.silver-muller-trace-approx} typically
introduces spurious reflected power below \(1\%\) of \(P_{\mathrm{rad}}\) while reactive
near fields at \(r=1.2a\) may differ by \(\mathcal{O}(1)\) from the exact Sommerfeld solution
\cite{harrington2001}. Sommerfeld completion is therefore necessary but not
sufficient for transport claims---regime qualification Eq.~\eqref{eq:ch2.regime-operator-screen}
remains mandatory \cite{stratton1941}. Figure~\ref{fig:ch2.sommerfeld-far-zone} separates Sommerfeld completion from regime qualification. Subsection~\ref{sec:ch241} fixes the asymptotic \emph{boundary at
infinity}; Subsection~\ref{sec:ch224} and Section~\ref{sec:ch01.2.4} govern when
far-zone samples may be interpreted as transport \cite{stratton1941}. Green-function asymptotics. Dyadic kernels from Subsection~\ref{sec:ch231} inherit scalar Sommerfeld decay.
For \(r=|\mathbf{r}-\mathbf{r}'|\to\infty\) at fixed \(\mathbf{r}'\),
\begin{equation}
\label{eq:ch2.sommerfeld-dyadic-asymptotic}
\overline{\mathbf{G}}_{E,0}(\mathbf{r},\mathbf{r}')
=
\overline{\mathbf{G}}_{\infty}(\theta,\phi,\mathbf{r}')\,
\frac{e^{-\jj k_0 r}}{r}
\;+\;\mathcal{O}(r^{-2}),
\end{equation}
with \(\overline{\mathbf{G}}_{\infty}\) transverse in the \(\mathbf{r}\)-direction
\cite{stratton1941}. Equation~\eqref{eq:ch2.sommerfeld-dyadic-asymptotic}
is the kernel form of Eq.~\eqref{eq:ch2.sommerfeld-field-condition} under
Eq.~\eqref{eq:ch2.electric-dyadic-convolution}: impressed currents that satisfy
Eq.~\eqref{eq:ch1.source-admissibility} generate fields whose far zone is built from
outgoing dyadic radiation \cite{stratton1941}. The leading \(r^{-1}\) term in
Eq.~\eqref{eq:ch2.free-space-dyadic-green} is the explicit vacuum instance of
Eq.~\eqref{eq:ch2.sommerfeld-dyadic-asymptotic} \cite{jackson1999}. Subsection~\ref{sec:ch242}
promotes Eq.~\eqref{eq:ch2.sommerfeld-dyadic-asymptotic} to an operator selector on
\(\mathcal{S}_\omega\); Subsection~\ref{sec:ch243} uses it in uniqueness energy
integrals \cite{harrington2001}.
Collecting the asymptotic and impedance conditions above yields the admissible function class
on which exterior uniqueness arguments and radiation export are defined. A field pair belongs
to this class only if it satisfies the vector Sommerfeld condition and the transverse far-zone
template simultaneously; incoming or standing-wave completions are excluded by construction
\cite{harrington2001}. Define the Sommerfeld outgoing subspace on an exterior shell
\(\Gamma_R=\{r=R\}\):
\begin{equation}
\label{eq:ch2.sommerfeld-outgoing-class}
\mathcal{F}_{\mathrm{SF}}
=
\left\{
(\E,\H)\in\HCurl(\Omega_{\mathrm{ext}}):
\text{Eqs.~\eqref{eq:ch2.sommerfeld-field-condition}--\eqref{eq:ch2.far-zone-transverse-radiation} hold}
\right\}.
\end{equation}
In Eq.~\eqref{eq:ch2.sommerfeld-outgoing-class},
\(\mathcal{F}_{\mathrm{SF}}\subset\HCurl(\Omega_{\mathrm{ext}})\) is the Sommerfeld
outgoing field class on exterior \(\Omega_{\mathrm{ext}}\): pairs \((\E,\H)\) whose
far-zone asymptotics satisfy outgoing phase and transverse radiation templates;
incoming and standing-wave completions are excluded by membership in this set.
Equation~\eqref{eq:ch2.sommerfeld-outgoing-class} embeds Eq.~\eqref{eq:ch1.radiation-trace-condition}
into the \ref{ch:02} function-space setting of
Eq.~\eqref{eq:ch2.admissible-field-space} \cite{stratton1941}. The physical
solution operator \(\mathcal{S}_\omega\) maps \(\mathcal{J}_\Omega\) into
\(\mathcal{F}_{\mathrm{SF}}\) when \(\mathcal{B}_\Gamma\) includes outgoing radiation condition
\cite{harrington2001}. Incoming solutions lie outside
\(\mathcal{F}_{\mathrm{SF}}\) even if they satisfy volume Maxwell rows; that exclusion is the Sommerfeld radiation condition \cite{stratton1941}.
Finite computations cannot place a boundary at infinity; the Silver--M\"uller trace below
approximates the Sommerfeld impedance relation on a truncation sphere large enough that
spurious reflections are small compared with exported power \cite{harrington2001}. On a finite
truncation sphere \(\Gamma_R\), Eq.~\eqref{eq:ch2.sommerfeld-impedance-ratio}
motivates the first-order absorbing condition
\begin{equation}
\label{eq:ch2.silver-muller-trace-approx}
\hat{\mathbf{n}}\times\H
\;-\;
\eta_0^{-1}(\hat{\mathbf{n}}\times\E)\times\hat{\mathbf{n}}
=
\mathbf{0}\quad\text{on }\Gamma_R,
\end{equation}
which approximates Eq.~\eqref{eq:ch1.radiation-trace-condition} at finite \(R\)
\cite{stratton1941}. Equation~\eqref{eq:ch2.silver-muller-trace-approx}
is not identical to Sommerfeld completion at infinity: spurious reflections from
\(\Gamma_R\) contaminate near fields unless \(R\) is large compared with source
extent and wavelength \cite{harrington2001}. Section~\ref{sec:ch24} treats
Eq.~\eqref{eq:ch2.silver-muller-trace-approx} as a practical surrogate; the
continuous target remains Eq.~\eqref{eq:ch2.sommerfeld-field-condition}
\cite{stratton1941}. Perfectly matched layers and radiation boundary integrals are
discrete refinements of the same outgoing-trace logic \cite{harrington2001}.
Once Sommerfeld closure is in place, the radiation export map $\mathcal{R}_\omega$ and projector
$\Pi_{\mathrm{rad}}$ act on the same outgoing class the range attributes to measured patterns:
$\mathcal{R}_\omega[\Jimp]=\Pi_{\mathrm{rad}}\,\mathcal{S}_\omega[\Jimp]$ with full field pair
$\mathcal{S}_\omega[\Jimp]$ carrying reactive content that $\Pi_{\mathrm{rad}}$ removes before
far-zone metering \cite{stratton1941}. The outgoing selector \(\Pi_{\mathrm{rad}}\) of
Eq.~\eqref{eq:ch2.radiation-operator-def} is idempotent on fields satisfying
Eq.~\eqref{eq:ch2.sommerfeld-outgoing-class} once uniqueness is established in
Subsection~\ref{sec:ch243} \cite{stratton1941}. Prior to that proof, Sommerfeld
closure defines the domain on which \(\Pi_{\mathrm{rad}}\) acts as projection rather than heuristic windowing \cite{harrington2001}. Composition with
Eq.~\eqref{eq:ch2.transport-qualified-radiation} yields
\begin{equation}
\label{eq:ch2.sommerfeld-radiation-composition}
\mathcal{R}_\omega[\Jimp]
=
\Pi_{\mathrm{rad}}\,\mathcal{S}_\omega[\Jimp],
\qquad
\mathcal{S}_\omega[\Jimp]\in\mathcal{F}_{\mathrm{SF}},
\end{equation}
with transport qualification from Section~\ref{sec:ch01.2.4} applied after
Eq.~\eqref{eq:ch2.sommerfeld-radiation-composition} \cite{stratton1941}. Equation
\eqref{eq:ch2.sommerfeld-radiation-composition} expresses
Eq.~\eqref{eq:ch2.radiation-time-frequency-equiv} on the Sommerfeld domain without
re-proving causality from Subsection~\ref{sec:ch223} \cite{harrington2001}. Consequence for
range practice: pattern meters and gain reports must compose with $\mathcal{R}_\omega$, not with
raw $\mathcal{S}_\omega[\Jimp]$, whenever reactive near fields are non-negligible on the
measurement shell \cite{stratton1941}. Subsection~\ref{sec:ch241} therefore states
what the compact range already assumes before any far-zone integral is evaluated in
Section~\ref{sec:ch25} \cite{stratton1941}.
Equations~\eqref{eq:ch2.sommerfeld-scalar-asymptotic}--\eqref{eq:ch2.sommerfeld-outgoing-class}
require linear vacuum far zones, fixed \(\omega>0\), and compact \(\operatorname{supp}\Jimp\).
They fail for anisotropic or magneto-electric exterior media without modified
asymptotics, for lossless closed cavities where Eq.~\eqref{eq:ch1.outer-boundary-classes}
uses PEC rather than outgoing radiation condition, and for time-harmonic fields with net incoming
power at infinity \cite{stratton1941}. Loss shifts
Eq.~\eqref{eq:ch2.sommerfeld-scalar-asymptotic} to complex \(k_0\) but preserves
outgoing \emph{selection} as distinct from incoming \cite{stratton1941}.
As a worked example, replace Eq.~\eqref{eq:ch2.sommerfeld-field-condition} by
Eq.~\eqref{eq:ch2.silver-muller-trace-approx} on \(\Gamma_R\) with \(R=5\lambda_0\).
Fields inside \(\Omega_s\) differ from the exact Sommerfeld solution by
\(\mathcal{O}((k_0 R)^{-1})\) in typical radiative components while reactive near
content may differ by order unity \cite{stratton1941}. Declaring a
transport channel from \(\Gamma_R\) probes without
Eq.~\eqref{eq:ch2.regime-operator-screen} conflates truncation error with angular
transport \cite{stratton1941}.
\subsection{Outgoing-wave selection}
\label{sec:ch242}
Finite computational domains terminate open-region problems artificially. On a physical compact
range, anechoic absorbers play the same role as absorbing boundary conditions in equivalent analytical models:
they mimic terminations that let export escape instead of reflecting into the near field
\cite{hansen1988}. Silver--M\"uller and related traces approximate outgoing selection
on meshes used for horns, cavities, and reflector feeds
\cite{stratton1941}. A finite-element model employing inadequate
absorbing boundaries permits artificial reflections that contaminate field amplitudes and
reconstructed radiation characteristics---the same standing-wave signatures that appear when chamber
absorption is imperfect on a physical range \cite{hansen1988}. Absorbing-boundary
parameters are therefore a radiation-physics choice about exported energy, not a mesh-only tuning
parameter \cite{stratton1941}. The operator $\mathcal{O}_{\mathrm{out}}$ records that choice on
$\mathcal{S}_\omega$ before $\Pi_{\mathrm{rad}}$ and $\mathcal{R}_\omega$ attach pattern meters.
% PHYSICAL-MEMORY:sec:ch242
\begin{physicalvignette}[Conical horn feeding a reflector edge]
Consider Conical horn feeding a reflector edge. Before any coefficient table is printed, technicians record frequency, cable loss, and the mechanical outline that fixes electrical size. The first walk-through of the range is always tactile: flange bolts, probe standoff, turntable boresight, and the wattmeter reference on the receive chain \cite{hansen1988}.
On this hardware the near zone is visually impressive---standing reactive structure, rapid phase rotation, sometimes a visible VSWR ripple at the feed---yet the export question is measured only after outgoing selection on a large sphere \cite{harrington2001}. What the camera sees in the chamber is not automatically what the far-zone pattern claims \cite{stratton1941}.
Engineers therefore learn to separate three readings: what the feed can synthesize on declared support, what propagating content survives on the radiative range, and what a finite instrument can recover at declared tolerance \cite{devaney2004}. Confusing those readings produces the rank gaps and watt failures that later sections quantify \cite{pal2026}.
Rayleigh and Stratton treated such ranges as boundary-value laboratories long before digital near-field transforms became routine \cite{stratton1941}. Harrington's method-of-moments culture reinforced the same discipline: metal outline first, harmonic headline second \cite{harrington2001}.
Subsection~\ref{sec:ch242} states the formal export condition that this range scene presupposes \cite{stratton1941}.
\end{physicalvignette}
The outgoing-wave selector $\mathcal{O}_{\mathrm{out}}$ implements Sommerfeld completion
Eq.~\eqref{eq:ch2.sommerfeld-outgoing-class} as an operator on $\mathcal{F}_\Omega$, on
harmonic Green kernels, and on the resolvent branch of $\mathcal{L}_\omega^{-1}$
\cite{stratton1941}. The composition relation
Eq.~\eqref{eq:ch2.rad-projection-composition} from Subsection~\ref{sec:ch224} named
$\mathcal{O}_{\mathrm{out}}$ without constructing it; the definitions below supply that
operator and its links to $\Pi_{\mathrm{prop}}$, $\Pi_{\mathrm{rad}}$, and retarded
Green functions Eq.~\eqref{eq:ch2.retarded-harmonic-kernel-theorem}. Sommerfeld field
conditions
Eqs.~\eqref{eq:ch2.sommerfeld-field-condition}--\eqref{eq:ch2.sommerfeld-outgoing-class}
are presupposed; uniqueness proofs belong to Subsection~\ref{sec:ch243}.
Assumptions. Exterior regions use vacuum far-zone constants with \(k_0>0\) real in the locked
harmonic band, compact \(\operatorname{supp}\Jimp\), and time-independent
\(\mathcal{B}_\Gamma\) including outgoing radiation condition at infinity or on
Eq.~\eqref{eq:ch2.silver-muller-trace-approx} at large finite \(R\)
\cite{stratton1941}. Propagating and evanescent projectors
Eqs.~\eqref{eq:ch2.spectral-projectors}--\eqref{eq:ch2.spectral-projector-algebra}
from Subsection~\ref{sec:ch224} are fixed on \(\widehat{\mathcal{F}}_\Omega\)
\cite{harrington2001}. Incoming and advanced branches satisfying
Eq.~\eqref{eq:ch2.advanced-kernel-violation} are excluded.
Let \(\mathcal{F}_{\mathrm{in}}\) denote the incoming counterpart of
\(\mathcal{F}_{\mathrm{SF}}\) obtained by reversing outgoing phase in the far zone.
The outgoing-wave selector is the idempotent map
\begin{equation}
\label{eq:ch2.outgoing-selector-def}
\mathcal{O}_{\mathrm{out}}:\mathcal{F}_\Omega\to\mathcal{F}_{\mathrm{SF}},
\qquad
\mathcal{O}_{\mathrm{out}}[\mathbf{u}_\omega]
=
\mathbf{u}_{\mathrm{out}},
\end{equation}
where \(\mathbf{u}_{\mathrm{out}}\) is the unique member of
\(\mathcal{F}_{\mathrm{SF}}\) whose propagating spectral content matches the outgoing
\(k_z\) branch of Eq.~\eqref{eq:ch2.longitudinal-wavenumber-split} and whose
incoming content is annihilated \cite{stratton1941}. Equation
\eqref{eq:ch2.outgoing-selector-def} is the locked definition for operator
realization of Sommerfeld completion: it acts \emph{before} far-zone measurement and
\emph{after} volume completion by \(\mathcal{L}_\omega\), selecting the physical
resolvent branch assumed in Eq.~\eqref{eq:ch2.resolvent-green-equivalence}
\cite{stratton1941}. On \(\mathcal{F}_{\mathrm{SF}}\),
\begin{equation}
\label{eq:ch2.outgoing-idempotency}
\mathcal{O}_{\mathrm{out}}^2=\mathcal{O}_{\mathrm{out}},
\qquad
\mathcal{O}_{\mathrm{out}}\big|_{\mathcal{F}_{\mathrm{SF}}}=\mathrm{id},
\end{equation}
while incoming states satisfy \(\mathcal{O}_{\mathrm{out}}[\mathbf{u}_{\mathrm{in}}]=\mathbf{0}\)
when extended by continuity from the far zone \cite{harrington2001}. Physically,
\(\mathcal{O}_{\mathrm{out}}\) removes incoming spherical and plane-wave components
that satisfy Maxwell equations but violate retarded causality
\cite{jackson1999}.
On homogeneous regions, outgoing selection in lateral wavenumber \((k_x,k_y)\) uses
the longitudinal root from Eq.~\eqref{eq:ch2.longitudinal-wavenumber-split}:
\begin{equation}
\label{eq:ch2.outgoing-plane-wave-criterion}
k_z=
\begin{cases}
+\sqrt{k_0^2-k_t^2}, & k_t\le k_0\ \text{(outgoing propagating)},\\[4pt]
-\jj\,\alpha(k_t), & k_t>k_0\ \text{(decaying evanescent)},
\end{cases}
\end{equation}
with the sign of the propagating root fixed by exterior completion and the locked
\(\exp(-\jj\omega t)\) convention \cite{stratton1941}. Equation
\eqref{eq:ch2.outgoing-plane-wave-criterion} is the spectral action of
\(\mathcal{O}_{\mathrm{out}}\) on each plane-wave mode; reversing
\(k_z\to -k_z\) on the propagating sheet produces incoming content annihilated by
\(\mathcal{O}_{\mathrm{out}}\) \cite{stratton1941}. Evanescent modes are unchanged
in decay direction but remain outside \(\Pi_{\mathrm{rad}}\) through
Eq.~\eqref{eq:ch2.rad-null-on-evanescent} \cite{harrington2001}. The propagating
projector \(\Pi_{\mathrm{prop}}\) from Subsection~\ref{sec:ch224} retains all
outgoing and incoming propagating modes; \(\mathcal{O}_{\mathrm{out}}\) acts only
after \(\Pi_{\mathrm{prop}}\) in the radiation chain
Eq.~\eqref{eq:ch2.rad-projection-composition} \cite{stratton1941}.
\begin{figure}[t]
\centering
\ChOneFigBlock{%
\begin{tikzpicture}[ch1fig, node distance=5mm and 7mm]
 \node[ch1box] (f) {$\mathcal{F}_\Omega$};
 \node[ch1box, right=7mm of f] (pp) {$\Pi_{\mathrm{prop}}$};
 \node[ch1box, right=7mm of pp] (oo) {$\mathcal{O}_{\mathrm{out}}$};
 \node[ch1box, right=7mm of oo] (pr) {$\Pi_{\mathrm{rad}}$};
 \node[ch1box, below=5mm of f] (ev) {$\Pi_{\mathrm{ev}}$};
 \draw[ch1flow] (f) -- (pp) -- (oo) -- (pr);
 \draw[ch1flow] (f.south) -- (f.south |- ev.north);
 \draw[ch1flow] (f.south |- ev.north) -- (ev.north);
\end{tikzpicture}%
}
\caption{Outgoing selection Eq.~\eqref{eq:ch2.outgoing-selector-def} sits between
propagating projection and radiation export:
\(\Pi_{\mathrm{rad}}=\mathcal{O}_{\mathrm{out}}\circ\Pi_{\mathrm{prop}}\).}
\label{fig:ch2.outgoing-selector-measurement chain}
\end{figure}
~\eqref{eq:ch2.rad-projection-composition} with explicit operator order. Regime
qualification Eq.~\eqref{eq:ch2.regime-operator-screen} applies after
\(\Pi_{\mathrm{rad}}\), not at the \(\mathcal{O}_{\mathrm{out}}\) stage
Green-function branch selection. Harmonic kernels from Section~\ref{sec:ch23}---with observation coordinate $\mathbf{r}$ and
source coordinate $\mathbf{r}'$ in $G$ and $\overline{\mathbf{G}}_{E}$---admit incoming and
outgoing analytic continuations. The outgoing dyadic branch is
\begin{equation}
\label{eq:ch2.outgoing-green-branch}
\overline{\mathbf{G}}_{E,\omega}^{\mathrm{(out)}}
=
\mathcal{O}_{\mathrm{out}}\,\overline{\mathbf{G}}_{E,\omega},
\qquad
\overline{\mathbf{G}}_{E,\omega}^{\mathrm{(out)}}
\ \text{satisfies}\ 
\text{Eq.~\eqref{eq:ch2.sommerfeld-dyadic-asymptotic}},
\end{equation}
with \(\mathcal{O}_{\mathrm{out}}\) acting on the observation-point dependence
\cite{stratton1941}. Equation~\eqref{eq:ch2.outgoing-green-branch}
identifies the kernel used in Eq.~\eqref{eq:ch2.electric-dyadic-convolution} on
exterior problems with the retarded symbol of
Eq.~\eqref{eq:ch2.retarded-harmonic-kernel-theorem} \cite{stratton1941}. Incoming
kernels \(\overline{\mathbf{G}}_{E,\omega}^{\mathrm{(in)}}\) are related by
\(k_0\to -k_0\) on the propagating Riemann sheet and produce fields outside
\(\mathcal{F}_{\mathrm{SF}}\) \cite{jackson1999}. Subsection~\ref{sec:ch242} does
not re-tabulate Eq.~\eqref{eq:ch2.free-space-dyadic-green}; it fixes the sheet on
which that formula is evaluated \cite{harrington2001}. The physical resolvent branch is
\begin{equation}
\label{eq:ch2.outgoing-resolvent-selection}
\mathcal{L}_\omega^{-1}\big|_{\mathrm{phys}}
=
\mathcal{O}_{\mathrm{out}}\circ\mathcal{L}_\omega^{-1},
\end{equation}
understood on \(\mathcal{D}(\mathcal{L}_\omega)\) with outgoing \(\mathcal{B}_\Gamma\)
\cite{stratton1941}. The solution map on impressed currents becomes
\begin{equation}
\label{eq:ch2.outgoing-solution-operator}
\mathcal{S}_\omega[\Jimp]
=
\mathcal{O}_{\mathrm{out}}\,
\Big(-\jj\omega\mu_0\int_{\Omega}
\overline{\mathbf{G}}_{E,\omega}^{\mathrm{(out)}}(\mathbf{r},\mathbf{r}')\cdot\Jimp(\mathbf{r}')\,d^{3}r'
\;+\;\E_{\mathrm{hom}}\Big),
\end{equation}
with regularization from Subsection~\ref{sec:ch232} when \(\mathbf{r}\) approaches
\(\operatorname{supp}\Jimp\) \cite{harrington2001}. Equations
\eqref{eq:ch2.outgoing-resolvent-selection} and~\eqref{eq:ch2.outgoing-solution-operator}
realize Eq.~\eqref{eq:ch2.exterior-uniqueness-hypothesis} once
Subsection~\ref{sec:ch243} proves uniqueness on the selected branch
\cite{stratton1941}. Causality from Subsection~\ref{sec:ch233} is consistent with
Eq.~\eqref{eq:ch2.outgoing-green-branch} but not sufficient when harmonic solvers
bypass time-domain construction \cite{harrington2001}. The volume integral in
Eq.~\eqref{eq:ch2.outgoing-solution-operator} sums impressed current contributions at each
source point \(\mathbf{r}'\in\Omega\) against the outgoing dyadic evaluated at observation
point \(\mathbf{r}\); \(\E_{\mathrm{hom}}\) carries any exterior illumination not generated
by \(\Jimp\) and must vanish for pure radiation problems \cite{stratton1941}.
On propagating content, outgoing selection precedes radiation export in the measurement sequence
fixed in Section~\ref{sec:ch24}: $\mathcal{R}_\omega=\Pi_{\mathrm{rad}}\circ\mathcal{S}_\omega$ with
$\Pi_{\mathrm{rad}}=\mathcal{O}_{\mathrm{out}}\circ\Pi_{\mathrm{prop}}$ \cite{stratton1941}. On
propagating content,
\begin{equation}
\label{eq:ch2.outgoing-radiation-refinement}
\Pi_{\mathrm{rad}}
=
\mathcal{O}_{\mathrm{out}}\circ\Pi_{\mathrm{prop}},
\qquad
\mathcal{R}_\omega
=
\Pi_{\mathrm{rad}}\circ\mathcal{S}_\omega,
\end{equation}
restating Eqs.~\eqref{eq:ch2.rad-projection-composition} and~\eqref{eq:ch2.radiation-operator-def}
with the locked definition of \(\mathcal{O}_{\mathrm{out}}\)
\cite{stratton1941}. Radiative subspace
\(\mathcal{F}_{\mathrm{rad}}\subset\mathcal{F}_{\mathrm{SF}}\) consists of fields
with nonzero far-zone transverse amplitude Eq.~\eqref{eq:ch2.far-zone-transverse-radiation};
\(\mathcal{O}_{\mathrm{out}}\) maps into \(\mathcal{F}_{\mathrm{SF}}\) but not
automatically into \(\mathcal{F}_{\mathrm{rad}}\) when evanescent content dominates
near sources \cite{stratton1941}. Transport reporting therefore requires
Eq.~\eqref{eq:ch2.regime-operator-screen} after
Eq.~\eqref{eq:ch2.outgoing-radiation-refinement} \cite{harrington2001}.
Outgoing selection on a compact range is implemented twice in practice: absorbers suppress
multipath reflections that would mimic incoming waves, and harmonic solvers fix the propagating
$k_z$ branch of Eq.~\eqref{eq:ch2.outgoing-plane-wave-criterion} on exterior meshes
\cite{hansen1988}. When either layer fails, $\mathcal{O}_{\mathrm{out}}$ acts on
the wrong field orbit even though volume Maxwell rows still vanish \cite{harrington2001}.
Incoming--outgoing split. Define \(\mathcal{O}_{\mathrm{in}}\) analogously on \(\mathcal{F}_{\mathrm{in}}\).
On the propagating subspace,
\begin{equation}
\label{eq:ch2.incoming-outgoing-complement}
\mathrm{id}
=
\mathcal{O}_{\mathrm{out}}+\mathcal{O}_{\mathrm{in}}
\quad\text{on}\ \Pi_{\mathrm{prop}}\mathcal{F}_\Omega,
\qquad
\mathcal{O}_{\mathrm{out}}\mathcal{O}_{\mathrm{in}}=0,
\end{equation}
with evanescent and non-propagating components orthogonal to both when decay
direction is fixed by Eq.~\eqref{eq:ch2.outgoing-plane-wave-criterion}
\cite{stratton1941}. Equation~\eqref{eq:ch2.incoming-outgoing-complement}
explains branch degeneracy in unconstrained Helmholtz solvers: both
\(\mathcal{O}_{\mathrm{out}}[\mathbf{u}]\) and
\(\mathcal{O}_{\mathrm{in}}[\mathbf{u}]\) satisfy volume equations unless one is
selected by boundary or radiation completion \cite{jackson1999}.
Equations~\eqref{eq:ch2.outgoing-selector-def}--\eqref{eq:ch2.outgoing-solution-operator}
require linear media, fixed \(\omega\), and outgoing \(\mathcal{B}_\Gamma\) on
exterior problems. They fail on closed PEC cavities where
Eq.~\eqref{eq:ch1.outer-boundary-classes} replaces Sommerfeld completion, and on
lossless open resonators where discrete spectra require quotienting
Eq.~\eqref{eq:ch2.coercivity-quotient} \cite{stratton1941}.
Absorbing layers approximate \(\mathcal{O}_{\mathrm{out}}\) but introduce spurious
modes unless tuned to Eq.~\eqref{eq:ch2.sommerfeld-field-condition} at large distance
\cite{harrington2001}.
\subsection{Uniqueness of radiative solutions}
\label{sec:ch243}
% CH02-NARR-EXP:02.243
Outgoing selection is not a post hoc spatial mask: Sommerfeld radiation and uniqueness on $\mathcal{B}_\Gamma$ decide which Maxwell solutions carry watt-bearing export before harmonic charts attach \cite{sommerfeld1949,stratton1941}. Reactive near-field energy may be locally large while contributing negligibly to $P_{\mathrm{rad}}$ on $\Gamma_R$ once $\Pi_{\mathrm{rad}}$ is fixed \cite{harrington2001}.
Without exterior uniqueness, two distinct field pairs could satisfy the same feed data while
differing on the measurement sphere---a failure no wattmeter or pattern probe could resolve
\cite{sommerfeld1949}. The proof below is therefore not a mathematical luxury
imposed for elegance; it is the physical condition under which laboratory attribution of
radiation to declared sources becomes meaningful. Outgoing selection is not a post hoc spatial
mask: Sommerfeld radiation and uniqueness on
$\mathcal{B}_\Gamma$ decide which Maxwell solutions carry watt-bearing export before harmonic
charts attach \cite{sommerfeld1949}. On $\Gamma_R$, reactive near-field energy may be locally
large while contributing negligibly to $P_{\mathrm{rad}}$ once $\mathcal{R}_\omega$ and
$\Pi_{\mathrm{rad}}$ are fixed \cite{harrington2001}. Exterior uniqueness ties measured patterns
to declared sources within the radiative class \cite{stratton1941}. Solvers that omit outgoing
filters fit non-radiating currents that compact ranges cannot distinguish from export without
watt completion on $\mathcal{R}_\omega[\Jimp]$ rather than on $\mathcal{S}_\omega[\Jimp]$ alone
\cite{jackson1999}. The proof below imports the difference-field energy method from
Subsection~\ref{sec:ch01.1.3} and specializes it to the Sommerfeld outgoing class fixed in
Subsection~\ref{sec:ch241}.
% PHYSICAL-MEMORY:sec:ch243
\begin{physicalvignette}[Planar near-field scanner above a handset]
Consider Planar near-field scanner above a handset. Before any coefficient table is printed, technicians record frequency, cable loss, and the mechanical outline that fixes electrical size. The first walk-through of the range is always tactile: flange bolts, probe standoff, turntable boresight, and the wattmeter reference on the receive chain \cite{hansen1988}.
On this hardware the near zone is visually impressive---standing reactive structure, rapid phase rotation, sometimes a visible VSWR ripple at the feed---yet the export question is measured only after outgoing selection on a large sphere \cite{harrington2001}. What the camera sees in the chamber is not automatically what the far-zone pattern claims \cite{stratton1941}.
Engineers therefore learn to separate three readings: what the feed can synthesize on declared support, what propagating content survives on the radiative range, and what a finite instrument can recover at declared tolerance \cite{devaney2004}. Confusing those readings produces the rank gaps and watt failures that later sections quantify \cite{pal2026}.
Exterior uniqueness ties measured patterns to declared sources within the radiation class . Inverse solvers that skip outgoing filters invent non-physical sources .
Rayleigh and Stratton treated such ranges as boundary-value laboratories long before digital near-field transforms became routine \cite{stratton1941}. Harrington's method-of-moments culture reinforced the same discipline: metal outline first, harmonic headline second \cite{harrington2001}.
Subsection~\ref{sec:ch243} states the formal export condition that this range scene presupposes \cite{stratton1941}.
\end{physicalvignette}
Subsection~\ref{sec:ch243} proves that exterior radiative solutions are unique on the
outgoing branch, upgrading hypothesis
Eq.~\eqref{eq:ch2.exterior-uniqueness-hypothesis} to exterior radiative uniqueness
\cite{stratton1941}. The difference-field reduction
Eqs.~\eqref{eq:ch1.uniqueness-difference-system}--\eqref{eq:ch1.uniqueness-null} and
energy identity Eq.~\eqref{eq:ch1.uniqueness-energy-identity} from
Subsection~\ref{sec:ch01.1.3} are presupposed; the passive uniqueness argument and
coercive estimate Eq.~\eqref{eq:ch1.uniqueness-estimate} are not re-proved.
Subsection~\ref{sec:ch243} specializes the energy method to exterior domains with
Eqs.~\eqref{eq:ch2.sommerfeld-field-condition}--\eqref{eq:ch2.sommerfeld-outgoing-class},
outgoing resolvent selection
Eqs.~\eqref{eq:ch2.outgoing-resolvent-selection}--\eqref{eq:ch2.outgoing-solution-operator},
and radiation refinement Eq.~\eqref{eq:ch2.outgoing-radiation-refinement}. Source
reconstruction limits belong to Subsection~\ref{sec:ch244}. Let \(\Omega_{\mathrm{ext}}\) denote an exterior domain with compact
\(\operatorname{supp}\Jimp\subset B_R(0)\), vacuum constants \(\varepsilon_0,\mu_0\)
in the far zone, and fixed \(\omega>0\) in the harmonic band of
Subsection~\ref{sec:ch223}. Data are admissible per
Eq.~\eqref{eq:ch1.admissible-data-class} with outgoing radiation condition
Eq.~\eqref{eq:ch2.sommerfeld-outgoing-class} at infinity, equivalently
Eq.~\eqref{eq:ch1.radiation-trace-condition} on spherical shells
\cite{stratton1941}. Constitutive response in the source region is passive:
\(\Im\{\varepsilon\}\ge 0\), \(\Im\{\mu\}\ge 0\) or equivalent loss so homogeneous
difference energy is dissipative \cite{stratton1941}. Incoming and advanced branches
are excluded by Eq.~\eqref{eq:ch2.outgoing-selector-def} and
Eq.~\eqref{eq:ch2.incoming-outgoing-complement} \cite{harrington2001}. Fields lie in
\(\mathcal{F}_{\mathrm{SF}}\) as defined in Subsection~\ref{sec:ch241}; radiative
export uses \(\mathcal{F}_{\mathrm{rad}}\subset\mathcal{F}_{\mathrm{SF}}\) from
Eq.~\eqref{eq:ch2.far-zone-transverse-radiation} \cite{stratton1941}.
Uniqueness is established by showing that two outgoing solutions for the same impressed data
cannot differ nontrivially: the proof therefore studies the homogeneous problem satisfied by
their difference \cite{stratton1941}. Let \((\E_1,\H_1)\) and \((\E_2,\H_2)\) be two harmonic
solutions on the outgoing branch for the same impressed data and \(\mathcal{B}_\Gamma\). Define
\(\Delta\E=\E_1-\E_2\), \(\Delta\H=\H_1-\H_2\), and
\(\mathbf{u}_\Delta=(\Delta\E,\Delta\H)\). Linearity and
Eq.~\eqref{eq:ch1.uniqueness-difference-system} give the homogeneous exterior system
\begin{equation}
\label{eq:ch2.exterior-difference-system}
\curl\Delta\E=-\jj\omega\mu\,\Delta\H,\quad
\curl\Delta\H=\jj\omega\varepsilon\,\Delta\E,\quad
\mathcal{B}_\Gamma[\Delta\E,\Delta\H]=\mathbf{0},
\end{equation}
with \(\varepsilon,\mu\) matching the physical model in \(\Omega_s\) and
\(\varepsilon_0,\mu_0\) outside \cite{stratton1941}. Physically, if two outgoing solutions
for the same feed existed, their difference would satisfy this homogeneous system while still
carrying export on $\Gamma_R$---a situation the flux argument below excludes
\cite{stratton1941}. Outgoing radiation condition requires
\begin{equation}
\label{eq:ch2.exterior-difference-outgoing}
\mathcal{O}_{\mathrm{out}}[\mathbf{u}_\Delta]=\mathbf{u}_\Delta,
\qquad
\mathbf{u}_\Delta\in\mathcal{F}_{\mathrm{SF}},
\end{equation}
which is the exterior specialization of the third row of
Eq.~\eqref{eq:ch1.uniqueness-difference-system} when
\(\mathcal{B}_\Gamma\) includes Sommerfeld radiation \cite{harrington2001}. Equations
\eqref{eq:ch2.exterior-difference-system}--\eqref{eq:ch2.exterior-difference-outgoing}
state that uniqueness testing reduces to showing only the trivial pair satisfies the
homogeneous outgoing problem \cite{stratton1941}. Incoming difference modes would
violate Eq.~\eqref{eq:ch2.exterior-difference-outgoing} even if they satisfied the
curl rows \cite{jackson1999}.
On truncated exterior shells, form the complex Poynting vector for difference fields,
\(\mathbf{S}_\Delta=\tfrac12\Delta\E\times\Delta\H^{*}\), as in
Eq.~\eqref{eq:ch1.uniqueness-energy-identity} \cite{stratton1941}. For a large
sphere \(\Gamma_R=\{r=R\}\) enclosing \(\Omega_s\), let
\(\Omega_R=B_R\setminus\Omega_s\). The divergence theorem on \(\Omega_R\) yields
\begin{equation}
\label{eq:ch2.exterior-flux-identity}
\int_{\Gamma_R}\hat{\mathbf{r}}\cdot\mathbf{S}_\Delta\,dS
=
-\jj\omega\!\int_{\Omega_R}
\Big(u_{\Delta H}-u_{\Delta E}\Big)\,dV
\;+\;\mathcal{P}_{\Delta,s},
\end{equation}
where \(u_{\Delta E}=\tfrac14\varepsilon|\Delta\E|^{2}\),
\(u_{\Delta H}=\tfrac14\mu|\Delta\H|^{2}\), and \(\mathcal{P}_{\Delta,s}\) collects
source-region passive dissipation when \(\Im\{\varepsilon\},\Im\{\mu\}>0\)
\cite{stratton1941}. Equation~\eqref{eq:ch2.exterior-flux-identity}
is Eq.~\eqref{eq:ch1.uniqueness-energy-identity} with \(\Delta\Jimp=\mathbf{0}\) and
\(\Gamma=\Gamma_R\); its real part controls uniqueness through the sign of the
boundary flux \cite{stratton1941}. On a compact range, the real part of the $\Gamma_R$ integral
is what a wattmeter would read from the difference field if such a superposition were
physically present; outgoing selection forces that reading to nonnegative export or zero
\cite{harrington2001}. In lossless vacuum on
\(\Omega_R\setminus\Omega_s\), the volume term is purely imaginary at steady state,
so \(\Re\int_{\Gamma_R}\hat{\mathbf{r}}\cdot\mathbf{S}_\Delta\,dS\) is determined
entirely by far-zone radiation \cite{harrington2001}.
Sommerfeld condition Eq.~\eqref{eq:ch2.sommerfeld-field-condition} implies that on
\(\Gamma_R\) with \(R\gg a=\max|\mathbf{r}|\) over \(\operatorname{supp}\Jimp\), the
leading transverse contribution to \(\Delta\E,\Delta\H\) scales as
\(e^{-\jj k_0 R}/R\) with \(k_0=\omega\sqrt{\varepsilon_0\mu_0}\)
\cite{stratton1941}. The far-zone transverse radiation law
Eq.~\eqref{eq:ch2.far-zone-transverse-radiation} then gives
\begin{equation}
\label{eq:ch2.exterior-flux-asymptotic}
\Re\int_{\Gamma_R}\hat{\mathbf{r}}\cdot\mathbf{S}_\Delta\,dS
=
P_{\Delta,\mathrm{rad}}(R)
=
\frac{1}{2\eta_0}\oint_{\Gamma_R}
\big|\mathbf{E}_{\Delta,t}\big|^{2}\,dS
\;+\;\mathcal{O}(R^{-1}),
\end{equation}
with \(\eta_0=\sqrt{\mu_0/\varepsilon_0}\) and \(\mathbf{E}_{\Delta,t}\) the
transverse spherical component of \(\Delta\E\) \cite{stratton1941}. Taking
\(R\to\infty\) under outgoing radiation condition,
\begin{equation}
\label{eq:ch2.exterior-sommerfeld-flux-limit}
\lim_{R\to\infty}P_{\Delta,\mathrm{rad}}(R)
=
P_{\Delta,\mathrm{rad}}\ge 0,
\end{equation}
where \(P_{\Delta,\mathrm{rad}}\) is the net outgoing power carried by the difference
field \cite{stratton1941}. Equation
\eqref{eq:ch2.exterior-sommerfeld-flux-limit} is the exterior boundary term induced
by Eq.~\eqref{eq:ch1.radiation-trace-condition}: outgoing difference fields radiate
nonnegative power to infinity, while incoming difference fields would produce negative
flux and are excluded by \(\mathcal{O}_{\mathrm{out}}\)
\cite{jackson1999}. Dyadic asymptotics
Eq.~\eqref{eq:ch2.sommerfeld-dyadic-asymptotic} are presupposed; they are not
re-tabulated here \cite{harrington2001}.
\begin{figure}[t]
\centering
\ChOneFigBlock{%
\begin{tikzpicture}[ch1fig, node distance=5mm and 7mm]
 \node[ch1box] (s) {Compact\\$\Omega_s$};
 \node[ch1box, right=16mm of s] (d) {$\Delta\E,\Delta\H$\\outgoing};
 \node[ch1box, right=16mm of d] (gr) {$\Gamma_R$\\flux};
 \node[ch1box, right=16mm of gr] (out) {$P_{\Delta,\mathrm{rad}}\ge 0$};
 \draw[ch1flow] (s) -- (d) -- (gr) -- (out);
\end{tikzpicture}%
}
\caption{Exterior uniqueness flux argument: homogeneous difference fields on the
outgoing branch satisfy Eqs.~\eqref{eq:ch2.exterior-flux-identity}--\eqref{eq:ch2.exterior-sommerfeld-flux-limit}
on spheres \(\Gamma_R\).}
\label{fig:ch2.exterior-uniqueness-flux}
\end{figure}
Section~\ref{sec:ch01.3}: volume terms in
Eq.~\eqref{eq:ch2.exterior-flux-identity} are controlled by passivity when loss is
present, and vanish in lossless vacuum away from sources \cite{stratton1941}. \phantomsection\label{lem:ch2.zero-flux-difference}
Let \(\mathbf{u}_\Delta=(\Delta\E,\Delta\H)\) satisfy
Eqs.~\eqref{eq:ch2.exterior-difference-system}--\eqref{eq:ch2.exterior-difference-outgoing}
on an exterior domain with vacuum far zone. If
\(P_{\Delta,\mathrm{rad}}=0\) in Eq.~\eqref{eq:ch2.exterior-sommerfeld-flux-limit},
then \(\mathbf{E}_{\Delta,t}=\mathbf{0}\) on every sphere \(\Gamma_R\) for
\(R\) large enough, and \(\mathbf{u}_\Delta\) has no propagating outgoing content in
\(\mathcal{F}_{\mathrm{rad}}\) \cite{stratton1941}.
Eq.~\eqref{eq:ch2.exterior-flux-asymptotic} with \(P_{\Delta,\mathrm{rad}}=0\) forces
\(\oint_{\Gamma_R}|\mathbf{E}_{\Delta,t}|^{2}\,dS=0\) for all sufficiently large
\(R\). Outgoing transverse fields with zero spherical flux cannot carry
\(1/R\) radiation; the only remaining homogeneous outgoing transverse modes are
identically zero on \(\mathcal{F}_{\mathrm{rad}}\) \cite{stratton1941}. \phantomsection\label{thm:ch2.exterior-radiative-uniqueness}
Let \(\Omega_{\mathrm{ext}}\) be an exterior domain with compact impressed sources,
vacuum far zone, and outgoing Sommerfeld completion
Eqs.~\eqref{eq:ch2.sommerfeld-outgoing-class}--\eqref{eq:ch2.outgoing-selector-def}.
If \((\E_1,\H_1)\) and \((\E_2,\H_2)\) are harmonic Maxwell solutions on the outgoing
branch for the same admissible data, then \(\E_1=\E_2\) and \(\H_1=\H_2\).
Equivalently,
\begin{equation}
\label{eq:ch2.exterior-uniqueness-theorem}
\mathcal{L}_\omega[\mathbf{u}_\omega]=\mathbf{0}
\ \text{and}\ 
\mathbf{u}_\omega\in\mathcal{F}_{\mathrm{SF}}
\ \Longrightarrow\
\mathbf{u}_\omega=\mathbf{0},
\end{equation}
which upgrades Eq.~\eqref{eq:ch2.exterior-uniqueness-hypothesis} to the Chapter~\ref{ch:02}
definitive uniqueness statement \cite{stratton1941}.
Let \(\mathbf{u}_\Delta\) be the difference of two outgoing solutions. Then
Eqs.~\eqref{eq:ch2.exterior-difference-system}--\eqref{eq:ch2.exterior-difference-outgoing}
hold. Apply Eq.~\eqref{eq:ch2.exterior-flux-identity} and take
\(\Re\{\cdot\}\). In lossless vacuum outside \(\Omega_s\), the volume integral is
purely imaginary, so \(\Re P_{\Delta,\mathrm{rad}}(R)\) equals the real boundary flux.
Outgoing radiation condition forces \(\Re P_{\Delta,\mathrm{rad}}\ge 0\); passivity in
\(\Omega_s\) adds nonnegative \(\Re\mathcal{P}_{\Delta,s}\), so the only way to
satisfy Eq.~\eqref{eq:ch2.exterior-sommerfeld-flux-limit} with equality throughout is
\(P_{\Delta,\mathrm{rad}}=0\) \cite{stratton1941}. The auxiliary relation (\ref{lem:ch2.zero-flux-difference})
eliminates outgoing transverse content. Evanescent components remain in
\(\mathcal{F}_{\mathrm{SF}}\) but are tied to the same source-free outgoing solution
operator; with no radiative export and homogeneous forcing, the unique retarded
exterior state is trivial by Eq.~\eqref{eq:ch2.outgoing-solution-operator} with
\(\Jimp=\mathbf{0}\) \cite{harrington2001}. Hence
\(\mathbf{u}_\Delta=\mathbf{0}\), proving Eq.~\eqref{eq:ch2.exterior-uniqueness-theorem}.
Exterior radiative uniqueness (Subsection~\ref{sec:ch243}) is the exterior counterpart of
the passive uniqueness argument (Subsection~\ref{sec:ch01.1.3}): instead of a bounded passive
\(\Gamma\) with coercivity Eq.~\eqref{eq:ch1.uniqueness-estimate}, the exterior uses
Sommerfeld flux Eq.~\eqref{eq:ch2.exterior-sommerfeld-flux-limit} as the coercive
boundary channel \cite{stratton1941}. The uniqueness index
Eq.~\eqref{eq:ch1.uniqueness-condition-index} is finite on exterior outgoing problems
away from embedded resonances in finite scatterers; it diverges on bounded cavities at
Eq.~\eqref{eq:ch2.cavity-resonance-failure} \cite{harrington2001}. Exterior uniqueness
therefore replaces bounded coercivity with an \emph{open} boundary condition that
selects decaying outgoing phase \cite{stratton1941}.
Exterior radiative uniqueness (Subsection~\ref{sec:ch243}) makes the physical resolvent
injective on the outgoing subspace:
\begin{equation}
\label{eq:ch2.exterior-resolvent-injective}
\mathcal{L}_\omega^{-1}\big|_{\mathrm{phys}}:
\mathcal{Y}_\omega\to\mathcal{F}_{\mathrm{SF}}
\ \text{is single-valued},
\qquad
\mathcal{S}_\omega[\Jimp]
=
\mathcal{L}_\omega^{-1}\big|_{\mathrm{phys}}
\begin{bmatrix}\mathbf{s}_\omega(\Jimp,\rho)\\ \mathbf{g}_\Gamma\end{bmatrix},
\end{equation}
consistent with Eq.~\eqref{eq:ch2.green-bvp-equivalence}, outgoing Green branch
Eq.~\eqref{eq:ch2.outgoing-green-branch}, and solution selection
Eq.~\eqref{eq:ch2.outgoing-solution-operator} \cite{stratton1941}.
\phantomsection\label{cor:ch2.outgoing-green-unique}
For fixed \(\Jimp\) in the admissible class, the exterior field
\(\mathcal{S}_\omega[\Jimp]\) obtained from
Eq.~\eqref{eq:ch2.outgoing-solution-operator} is unique in
\(\mathcal{F}_{\mathrm{SF}}\) \cite{stratton1941}.
Two convolutions differing by a homogeneous outgoing state violate
exterior radiative uniqueness (Subsection~\ref{sec:ch243}) \cite{harrington2001}.
Subsection~\ref{sec:ch201} declared \(\Pi_{\mathrm{rad}}\) idempotent on
\(\mathcal{F}_{\mathrm{rad}}\); the joint property with outgoing radiation condition can now be
stated precisely. On radiative states,
\begin{equation}
\label{eq:ch2.rad-idempotency-radiative}
\Pi_{\mathrm{rad}}^2=\Pi_{\mathrm{rad}}
\quad\text{on}\ \mathcal{F}_{\mathrm{rad}},
\end{equation}
as developed in Subsection~\ref{sec:ch201} \cite{stratton1941}. After
exterior radiative uniqueness (Subsection~\ref{sec:ch243}), the measurement sequence
Eq.~\eqref{eq:ch2.outgoing-radiation-refinement} defines a single-valued map
\(\mathcal{R}_\omega=\Pi_{\mathrm{rad}}\circ\mathcal{S}_\omega\) without branch
ambiguity from duplicate outgoing decompositions \cite{harrington2001}. Evanescent
components remain in \(\ker\Pi_{\mathrm{rad}}\) through
Eq.~\eqref{eq:ch2.rad-null-on-evanescent}; uniqueness of \(\mathcal{S}_\omega\) does
not imply uniqueness of \(\mathcal{R}_\omega[\Jimp]\) when near-zone evanescent
storage dominates export \cite{stratton1941}.
Outgoing selector idempotency Eq.~\eqref{eq:ch2.outgoing-idempotency} and resolvent
selection Eq.~\eqref{eq:ch2.outgoing-resolvent-selection} are consistent with
exterior radiative uniqueness (Subsection~\ref{sec:ch243}):
\begin{equation}
\label{eq:ch2.outgoing-resolvent-unique}
\mathcal{O}_{\mathrm{out}}\circ\mathcal{L}_\omega^{-1}
=
\mathcal{L}_\omega^{-1}\big|_{\mathrm{phys}}
\quad\text{is unique on}\ \mathcal{Y}_\omega,
\end{equation}
because any two selections differing by a homogeneous outgoing mode would violate
Eq.~\eqref{eq:ch2.exterior-uniqueness-theorem} \cite{stratton1941}.
Equation~\eqref{eq:ch2.outgoing-resolvent-unique} closes the excitation-to-flux sequence of
Figure~\ref{fig:ch2.outgoing-selector-measurement chain} on exterior domains: once data are
fixed, \(\mathcal{O}_{\mathrm{out}}\), \(\mathcal{S}_\omega\), and
\(\Pi_{\mathrm{rad}}\) act on a single field orbit in \(\mathcal{F}_{\mathrm{SF}}\)
\cite{stratton1941}.
On bounded domains with passive \(\mathcal{B}_\Gamma\), uniqueness follows from
Eq.~\eqref{eq:ch2.lax-milgram-wellposed} and
the passive uniqueness argument (Subsection~\ref{sec:ch01.1.3}) after quotienting
Eq.~\eqref{eq:ch2.coercivity-quotient} \cite{stratton1941}. At cavity eigenfrequencies,
\begin{equation}
\label{eq:ch2.exterior-vs-cavity-uniqueness}
\text{bounded PEC cavity:}\ \ker\mathcal{L}_\omega\neq\{\mathbf{0}\}
\quad\text{vs.}\quad
\text{exterior outgoing:}\ \ker\mathcal{L}_\omega|_{\mathcal{F}_{\mathrm{SF}}}=\{\mathbf{0}\},
\end{equation}
so Eq.~\eqref{eq:ch2.exterior-uniqueness-theorem} does not apply to closed resonators
where Eq.~\eqref{eq:ch2.cavity-resonance-failure} holds \cite{harrington2001}.
Exterior uniqueness is therefore a property of \(\mathcal{F}_{\mathrm{SF}}\) with
Sommerfeld completion, not of arbitrary \(\mathcal{B}_\Gamma\) from
Eq.~\eqref{eq:ch1.outer-boundary-classes} \cite{stratton1941}. A finite scatterer
embedded in free space may support Mie resonances in the near field without breaking
exterior uniqueness: those modes are driven by \(\Jimp\) and do not appear as free
homogeneous outgoing states \cite{jackson1999}.
Laboratory errors and non-applicability. Exterior radiative uniqueness (Subsection~\ref{sec:ch243}) requires linear passive media,
fixed \(\omega\) away from pathological embedded spectra, and exact outgoing radiation condition
rather than approximate ABC alone \cite{stratton1941}. First-order
absorbing boundaries that do not enforce Eq.~\eqref{eq:ch2.sommerfeld-field-condition}
at large \(R\) can admit spurious incoming content; uniqueness is then a property of
the \emph{augmented} truncated problem, not of the physical exterior
\cite{harrington2001}. Anisotropic, moving, or nonlinear media require modified flux
functionals; time-varying \(\mathcal{B}_\Gamma(t)\) breaks the harmonic reduction of
Subsection~\ref{sec:ch223} \cite{stratton1941}. Uniqueness of \(\mathcal{S}_\omega\)
does not imply uniqueness of \(\Jimp\) in \(\mathcal{R}_\omega\)---that is
Subsection~\ref{sec:ch244} \cite{harrington2001}. Regime qualification for transport
remains Eq.~\eqref{eq:ch2.regime-operator-screen}, independent of field uniqueness
\cite{stratton1941}.
As a worked example, truncating at finite \(R\) with a first-order ABC may return two numerically distinct
fields that agree on \(\Gamma_R\) but differ in spurious incoming content reflected
between \(\Omega_s\) and the ABC \cite{harrington2001}. Those differences are not
counterexamples to exterior radiative uniqueness (Subsection~\ref{sec:ch243}) on
\(\Omega_{\mathrm{ext}}\); they signal that the truncated BVP lacks
Eq.~\eqref{eq:ch2.sommerfeld-outgoing-class}. Silver--Müller surrogates
Eq.~\eqref{eq:ch2.silver-muller-trace-approx} improve the completion but still require
\(R\gg\lambda_0\) for uniqueness in practice \cite{stratton1941}.
\subsection{Implications for source reconstruction}
\label{sec:ch244}
% HARRINGTON-OPEN:sec:ch244
A coax-fed pyramidal horn over a ground plane defines the impressed-current and boundary-data tuple technicians declare before Maxwell export is computed \cite{stratton1941}. Maxwell's retarded export to $\Gamma_R$ remains the governing constraint \cite{maxwell1873}.
Exterior radiative uniqueness (Subsection~\ref{sec:ch243}) fixes uniqueness of
\(\mathcal{S}_\omega[\Jimp]\) on the outgoing branch. Subsection~\ref{sec:ch244}
records the inverse-source consequences: field uniqueness does not imply source
identifiability. Uniqueness of \((\E,\H)\) under outgoing conditions does not make
$\mathcal{R}_\omega$ injective---distinct $\Jimp$ can share the same far pattern
\cite{devaney2004}. Source reconstruction is a conditioned inverse, not a direct readout
of coefficients \cite{hansen1988}.
\medskip
\(\mathcal{R}_\omega\) is generally non-injective, and finite observation channels
compress rank further \cite{stratton1941}. Observation null-space
Eq.~\eqref{eq:ch2.observation-null-space}, rank composition
Eq.~\eqref{eq:ch2.rank-composition-law}, and non-radiating relation
Eq.~\eqref{eq:ch2.null-radiation} from Subsection~\ref{sec:ch213} are
presupposed; complete structure of \(\ker\mathcal{R}_\omega\) belongs to
Section~\ref{sec:ch272} \cite{harrington2001}. Measurement classification from
Section~\ref{sec:ch01.5} applies here; angular-momentum observables from
Chapter~\ref{ch:01} are not re-opened \cite{stratton1941}. Field uniqueness versus source identifiability. Forward exterior completion yields
% PHYSICAL-MEMORY:sec:ch244
\begin{physicalvignette}[Cylindrical near-field arch with probe carriage]
Consider Cylindrical near-field arch with probe carriage. Before any coefficient table is printed, technicians record frequency, cable loss, and the mechanical outline that fixes electrical size. The first walk-through of the range is always tactile: flange bolts, probe standoff, turntable boresight, and the wattmeter reference on the receive chain \cite{hansen1988}.
On this hardware the near zone is visually impressive---standing reactive structure, rapid phase rotation, sometimes a visible VSWR ripple at the feed---yet the export question is measured only after outgoing selection on a large sphere \cite{harrington2001}. What the camera sees in the chamber is not automatically what the far-zone pattern claims \cite{stratton1941}.
Engineers therefore learn to separate three readings: what the feed can synthesize on declared support, what propagating content survives on the radiative range, and what a finite instrument can recover at declared tolerance \cite{devaney2004}. Confusing those readings produces the rank gaps and watt failures that later sections quantify \cite{pal2026}.
Rayleigh and Stratton treated such ranges as boundary-value laboratories long before digital near-field transforms became routine \cite{stratton1941}. Harrington's method-of-moments culture reinforced the same discipline: metal outline first, harmonic headline second \cite{harrington2001}.
Subsection~\ref{sec:ch244} states the formal export condition that this range scene presupposes \cite{stratton1941}.
\end{physicalvignette}
\begin{equation}
\label{eq:ch2.field-source-uniqueness-split}
\underbrace{\mathcal{S}_\omega[\Jimp^{(1)}]=\mathcal{S}_\omega[\Jimp^{(2)}]}_{\text{same impressed data}}
\ \Longrightarrow\
\Jimp^{(1)}=\Jimp^{(2)}
\quad\text{false in general},
\end{equation}
whereas exterior radiative uniqueness (Subsection~\ref{sec:ch243}) gives
\begin{equation}
\label{eq:ch2.field-uniqueness-forward}
\mathcal{S}_\omega[\Jimp^{(1)}]=\mathcal{S}_\omega[\Jimp^{(2)}]
\ \Longrightarrow\
(\E_1,\H_1)=(\E_2,\H_2)
\quad\text{on}\ \mathcal{F}_{\mathrm{SF}}
\end{equation}
for the same boundary class and outgoing branch \cite{stratton1941}. Equations
\eqref{eq:ch2.field-source-uniqueness-split}--\eqref{eq:ch2.field-uniqueness-forward}
separate \emph{state} uniqueness from \emph{excitation} uniqueness: many admissible
currents can produce the same radiative exterior field once reactive degrees of
freedom are quotiented \cite{harrington2001}. Inverse problems therefore target
equivalence classes of \(\Jimp\), not individual source distributions
\cite{stratton1941}.
The radiative export map from Subsection~\ref{sec:ch201} is
\begin{equation}
\label{eq:ch2.radiation-forward-map}
\mathcal{R}_\omega:\mathcal{J}_\Omega\longrightarrow\mathcal{F}_{\mathrm{rad}},
\qquad
\mathcal{R}_\omega[\Jimp]=\Pi_{\mathrm{rad}}\,\mathcal{S}_\omega[\Jimp],
\end{equation}
restating Eq.~\eqref{eq:ch2.radiation-operator-def} after outgoing radiation condition
Eqs.~\eqref{eq:ch2.outgoing-radiation-refinement}--\eqref{eq:ch2.outgoing-resolvent-unique}
\cite{stratton1941}. The source-side null space recorded in
Eq.~\eqref{eq:ch2.null-radiation} is
\begin{equation}
\label{eq:ch2.ker-R-omega}
\ker\mathcal{R}_\omega
=
\left\{
\Jimp\in\mathcal{J}_\Omega:
\Pi_{\mathrm{rad}}\,\mathcal{S}_\omega[\Jimp]=\mathbf{0}
\right\},
\end{equation}
generally infinite-dimensional on open domains \cite{harrington2001}.
Section~\ref{sec:ch272} develops the complete null-space analysis; Subsection~\ref{sec:ch244}
records the reconstruction consequences of Eq.~\eqref{eq:ch2.ker-R-omega}
without constructing non-radiating bases \cite{stratton1941}. Elements of
\(\ker\mathcal{R}_\omega\) are invisible to far-zone radiative observables yet may
carry substantial near-zone reactive energy through \(\mathcal{S}_\omega\)
\cite{harrington2001}.
\begin{figure}[t]
\centering
\ChOneFigBlock{%
\begin{tikzpicture}[ch1fig, node distance=5mm and 7mm]
 \node[ch1box] (j) {$\mathcal{J}_\Omega$};
 \node[ch1box, right=7mm of j] (s) {$\mathcal{S}_\omega$};
 \node[ch1box, right=7mm of s] (f) {$\mathcal{F}_{\mathrm{SF}}$};
 \node[ch1box, right=7mm of f] (r) {$\mathcal{F}_{\mathrm{rad}}$};
 \node[ch1box, below=5mm of j] (ker) {$\ker\mathcal{R}_\omega$};
 \draw[ch1flow] (j) -- (s) -- (f) -- node[above,font=\scriptsize] {$\Pi_{\mathrm{rad}}$} (r);
 \draw[ch1flow, dashed] (ker.south) -- (j.north);
\end{tikzpicture}%
}
\caption{Forward chain for reconstruction: \(\mathcal{S}_\omega\) is single-valued on
\(\mathcal{F}_{\mathrm{SF}}\) after exterior radiative uniqueness (Subsection~\ref{sec:ch243});
\(\mathcal{R}_\omega\) collapses many sources to the same radiative output.}
\label{fig:ch2.reconstruction-forward-chain}
\end{figure}
Observation composition and inverse problem. Laboratory readouts compose with $\mathcal{R}_\omega$ before any source inversion
is attempted: pattern probes and near-field scanners observe a finite-rank image of impressed
currents, not $\Jimp$ itself \cite{hansen1988}. From
Eq.~\eqref{eq:ch2.radiation-operator-composition} and finite-rank observation
Eq.~\eqref{eq:ch2.finite-rank-observation},
\begin{equation}
\label{eq:ch2.reconstruction-observation-map}
\mathcal{M}_\omega
=
\mathcal{O}_{\mathrm{meas}}\circ\mathcal{O}_{\mathrm{far}}\circ\mathcal{R}_\omega
=
\mathcal{O}_{\mathrm{meas}}\circ\Lambda_{\mathrm{rad}}:
\mathcal{J}_\Omega\longrightarrow\mathbb{C}^{N},
\end{equation}
with \(\Lambda_{\mathrm{rad}}=\mathcal{O}_{\mathrm{far}}\circ\mathcal{R}_\omega\)
\cite{stratton1941}. Given measurement vector
\(\mathbf{y}\in\mathbb{C}^{N}\), the inverse source problem seeks
\(\Jimp\in\mathcal{J}_\Omega\) such that
\begin{equation}
\label{eq:ch2.inverse-source-equation}
\mathcal{M}_\omega[\Jimp]=\mathbf{y},
\qquad
\mathcal{D}_{J}[\Jimp]=1,
\end{equation}
subject to admissibility Eq.~\eqref{eq:ch2.source-admissibility}
\cite{harrington2001}. Equation~\eqref{eq:ch2.inverse-source-equation} is ill-posed
whenever \(\ker\mathcal{M}_\omega\neq\{\mathbf{0}\}\), which occurs if either
\(\ker\mathcal{R}_\omega\neq\{\mathbf{0}\}\) or \(\ker\mathcal{O}_{\mathrm{meas}}\)
is nontrivial on \(\operatorname{ran}\mathcal{R}_\omega\)
\cite{stratton1941}. Rank cap Eq.~\eqref{eq:ch2.rank-composition-law} applies directly:
\begin{equation}
\label{eq:ch2.reconstruction-rank-cap}
\mathrm{rank}(\mathcal{M}_\omega)
\le
\min\big\{N,\,\dim\mathcal{R}(\mathcal{R}_\omega)\big\},
\end{equation}
so no measurement protocol with \(N\) channels resolves more than \(N\) independent
radiative degrees of freedom \cite{harrington2001}. The composed null space satisfies
\begin{equation}
\label{eq:ch2.reconstruction-null-space}
\ker\mathcal{M}_\omega
\supseteq
\ker\mathcal{R}_\omega
\ \cup\
\mathcal{R}_\omega^{-1}\!\big(\ker\mathcal{O}_{\mathrm{meas}}\big),
\end{equation}
with set inclusion on admissible sources \cite{stratton1941}. Identifiable sources
modulo null space form the quotient
\begin{equation}
\label{eq:ch2.source-identifiability-quotient}
\mathcal{J}_\Omega\big/\ker\mathcal{M}_\omega
\ \cong\
\operatorname{ran}\mathcal{M}_\omega
\subset\mathbb{C}^{N},
\end{equation}
so reconstruction recovers at best a finite-dimensional image, not a unique
\(\Jimp(\mathbf{r})\) \cite{harrington2001}. Equivalent-current models that expand
\(\Jimp\) on auxiliary surfaces inherit the same non-uniqueness: distinct surface
distributions related by \(\ker\mathcal{R}_\omega\) produce identical far-zone
patterns \cite{stratton1941}. Boundary-element inversion therefore
requires regularization or additional constraints beyond Maxwell consistency alone
\cite{harrington2001}.
\begin{figure}[t]
\centering
\ChOneFigBlock{%
\begin{tikzpicture}[ch1fig, node distance=5mm and 7mm]
 \node[ch1box] (j1) {$\Jimp^{(1)}$};
 \node[ch1box, right=7mm of j1] (j2) {$\Jimp^{(2)}$};
 \node[ch1box, right=7mm of j2] (plus) {$\Jimp^{(1)}+\delta\Jimp$};
 \node[ch1box, below=5mm of plus] (y) {Same $\mathbf{y}$};
 \node[ch1box, below=5mm of j1] (delta) {$\delta\Jimp\in\ker\mathcal{M}_\omega$};
 \draw[ch1flow] (j1.east) -- (j1.east |- plus.west);
 \draw[ch1flow] (j1.east |- plus.west) -- (plus.west);
 \draw[ch1flow] (delta.east) -- (delta.east |- plus.west);
 \draw[ch1flow] (delta.east |- plus.west) -- (plus.west);
 \draw[ch1flow] (plus.south) -- (plus.south |- y.north);
 \draw[ch1flow] (plus.south |- y.north) -- (y.north);
\end{tikzpicture}%
}
\caption{Non-uniqueness: adding \(\delta\Jimp\in\ker\mathcal{M}_\omega\) leaves
measurement vector \(\mathbf{y}\) unchanged.}
\label{fig:ch2.reconstruction-null-class}
\end{figure}
~\eqref{eq:ch2.observation-null-space} pulled back to source space through
Eq.~\eqref{eq:ch2.reconstruction-observation-map} \cite{stratton1941}. When Eq.~\eqref{eq:ch2.inverse-source-equation} is ill-posed, stable inversion uses
a penalized least-squares or Tikhonov template on admissible sources:
\begin{equation}
\label{eq:ch2.tikhonov-source-reconstruction}
{\Jimp}_{\mathrm{reg}}
=
\arg\min_{\Jimp\in\mathcal{J}_\Omega}
\left\{
\|\mathcal{M}_\omega[\Jimp]-\mathbf{y}\|^{2}
+\alpha\,\mathcal{P}[\Jimp]
\right\},
\end{equation}
with regularization parameter \(\alpha>0\) and penalty
\(\mathcal{P}\) enforcing smoothness, minimum norm, or sparsity on an equivalent
surface \cite{harrington2001}. Equation
\eqref{eq:ch2.tikhonov-source-reconstruction} does not remove
\(\ker\mathcal{M}_\omega\); it selects one representative from the affine coset
\(\mathcal{M}_\omega^{-1}(\mathbf{y})\) \cite{stratton1941}. The choice of
\(\mathcal{P}\) is physics-adjacent modeling, not a Maxwell identity; different
penalties yield different \({\Jimp}_{\mathrm{reg}}\) with identical radiative
predictions when \(\delta\Jimp\in\ker\mathcal{R}_\omega\) \cite{harrington2001}. Transport and angular reports attach to \(\mathcal{R}_\omega[\Jimp]\), not to raw
\(\mathcal{S}_\omega[\Jimp]\), per Eq.~\eqref{eq:ch2.transport-qualified-radiation}
\cite{stratton1941}. regime qualification Eq.~\eqref{eq:ch2.regime-operator-screen} criteria
whether \(\mathcal{R}_\omega[\Jimp]\) lies in the propagating sector on which
far-zone observables were designed:
\begin{equation}
\label{eq:ch2.reconstruction-regime-criterion}
\mathcal{D}_{J,\mathrm{reg}}=1
\ \Longrightarrow\
\text{inverse claims target}\ \mathcal{R}_\omega[\Jimp],
\ \text{not reactive near fields},
\end{equation}
with \(\mathcal{D}_{J,\mathrm{reg}}\) from Eq.~\eqref{eq:ch1.regime-qualified-decision}
\cite{stratton1941}. Attempting to reconstruct \(\Jimp\) from
near-zone reactive samples without regime qualification conflates storage with export,
the laboratory error flagged in Section~\ref{sec:ch01.2.4} \cite{stratton1941}. Outgoing
uniqueness from exterior radiative uniqueness (Subsection~\ref{sec:ch243}) stabilizes the
\emph{field} side of Eq.~\eqref{eq:ch2.inverse-source-equation}; it does not promote
evanescent-dominated states to identifiable radiative sources
\cite{harrington2001}.
Formal adjoint back-projection provides a reconstruction \emph{template} without
claiming uniqueness:
\begin{equation}
\label{eq:ch2.adjoint-backprojection}
\widetilde{\Jimp}
=
\mathcal{M}_\omega^{\dagger}\,\mathbf{y},
\end{equation}
where \(\mathcal{M}_\omega^{\dagger}\) is the Hilbert-space adjoint on declared
domains \cite{stratton1941}. Equation
\eqref{eq:ch2.adjoint-backprojection} is useful for sensitivity analysis and
iterative inversion but generally satisfies
\(\mathcal{M}_\omega[\widetilde{\Jimp}]\neq\mathbf{y}\) when
\(\ker\mathcal{M}_\omega\neq\{\mathbf{0}\}\) \cite{harrington2001}. Consistent
reconstruction requires projecting onto \(\operatorname{ran}\mathcal{M}_\omega^{\dagger}\)
or applying Eq.~\eqref{eq:ch2.tikhonov-source-reconstruction}
\cite{stratton1941}. Full integral representations for
Eq.~\eqref{eq:ch2.adjoint-backprojection} belong to Section~\ref{sec:ch251}; they are
not repeated here \cite{harrington2001}.
Equivalent sources and Huygens surfaces. Equivalent-current inversion on a closed surface \(\Gamma_{\mathrm{eq}}\) seeks
\(\mathbf{J}_{\mathrm{eq}}\) on \(\Gamma_{\mathrm{eq}}\) such that exterior fields
match measured data \cite{stratton1941}. Uniqueness of exterior
\(\mathcal{S}_\omega\) from exterior radiative uniqueness (Subsection~\ref{sec:ch243})
fixes the \emph{field} on \(\mathcal{F}_{\mathrm{SF}}\) but not
\(\mathbf{J}_{\mathrm{eq}}\): any non-radiating tangential distribution on
\(\Gamma_{\mathrm{eq}}\) can be added without changing far-zone samples
\cite{harrington2001}. The operator form is
\begin{equation}
\label{eq:ch2.equivalent-source-ambiguity}
\mathbf{J}_{\mathrm{eq},1}-\mathbf{J}_{\mathrm{eq},2}
\in
\ker\mathcal{R}_\omega\big|_{\Gamma_{\mathrm{eq}}},
\end{equation}
with restriction to surface currents in the trace class of
Subsection~\ref{sec:ch211} \cite{stratton1941}. Numerical equivalence in
Eq.~\eqref{eq:ch2.discrete-green-equivalence} inherits the same null-space structure
at finite matrix rank \cite{harrington2001}.
Near-field versus far-field data. Near-field scanning samples \(\mathcal{S}_\omega[\Jimp]\) on a surface enclosing
\(\Omega_s\), not \(\mathcal{R}_\omega[\Jimp]\) alone \cite{stratton1941}. Because
\(\mathcal{S}_\omega\) is injective on outgoing fields but many sources share the
same \(\mathcal{R}_\omega\) image, near-field data can \emph{reduce} null-space
dimension relative to far-zone-only inversion without eliminating
\(\ker\mathcal{R}_\omega\) in general \cite{harrington2001}. Outgoing radiation condition still
required: two incoming/outgoing superpositions with identical near traces on a finite
surface need not agree in the far zone unless exterior radiative uniqueness (Subsection~\ref{sec:ch243})
applies to the full exterior problem \cite{stratton1941}. Near-field
inversion therefore couples Subsection~\ref{sec:ch243} field uniqueness with
Subsection~\ref{sec:ch244} source non-uniqueness \cite{harrington2001}. Equations~\eqref{eq:ch2.inverse-source-equation}--\eqref{eq:ch2.tikhonov-source-reconstruction}
require linear Maxwell theory, fixed \(\omega\), admissible
Eq.~\eqref{eq:ch2.source-admissibility}, and outgoing forward completion
Eq.~\eqref{eq:ch2.outgoing-resolvent-unique} \cite{stratton1941}.
Nonlinear, time-varying, or dispersive inversion breaks the harmonic reduction of
Subsection~\ref{sec:ch223}. Single-frequency inversion cannot resolve broadband
excitation without superposition across \(\omega\) \cite{stratton1941}. Loss improves
conditioning of Eq.~\eqref{eq:ch2.inverse-source-equation} by damping resonances but
does not make \(\mathcal{R}_\omega\) injective on all of \(\mathcal{J}_\Omega\)
\cite{harrington2001}. Complete null-space bases and NR source examples belong to
Section~\ref{sec:ch27} \cite{stratton1941}.
As a worked example, with \(N\) probe channels, Eq.~\eqref{eq:ch2.reconstruction-rank-cap} limits
recoverable modal content to rank \(N\) regardless of mesh density
\cite{stratton1941}. Refining a boundary-element mesh without adding independent
measurements does not resolve components in
\(\ker\mathcal{O}_{\mathrm{meas}}\circ\Lambda_{\mathrm{rad}}\)
\cite{harrington2001}. This is the operator reading of
Eq.~\eqref{eq:ch2.observation-null-space}, distinct from discretization error alone
\cite{stratton1941}. As a worked example, a source with \(\mathcal{D}_{J,\mathrm{reg}}=0\) may produce strong near-zone
\(\mathcal{S}_\omega[\Jimp]\) while \(\mathcal{R}_\omega[\Jimp]\approx\mathbf{0}\)
\cite{stratton1941}. Inverting near samples without
Eq.~\eqref{eq:ch2.reconstruction-regime-criterion} assigns spurious radiative labels to
storage-dominated fields \cite{harrington2001}. regime qualification
Eq.~\eqref{eq:ch2.regime-operator-screen} must precede angular-transport claims on
reconstructed sources \cite{stratton1941}. \medskip
\noindent
Section~\ref{sec:ch24} closes with the excitation-to-flux sequence fixed for open regions:
Sommerfeld asymptotics and outgoing class
Eqs.~\eqref{eq:ch2.sommerfeld-outgoing-class}--\eqref{eq:ch2.sommerfeld-field-condition}
(Subsection~\ref{sec:ch241}), outgoing selector and resolvent branch
Eqs.~\eqref{eq:ch2.outgoing-selector-def}--\eqref{eq:ch2.outgoing-solution-operator}
(Subsection~\ref{sec:ch242}), exterior radiative uniqueness
Eq.~\eqref{eq:ch2.exterior-uniqueness-theorem} and resolvent injectivity
Eq.~\eqref{eq:ch2.exterior-resolvent-injective} (Subsection~\ref{sec:ch243}), and
inverse reconstruction limits
Eqs.~\eqref{eq:ch2.inverse-source-equation}--\eqref{eq:ch2.ker-R-omega}
(Subsection~\ref{sec:ch244}) \cite{stratton1941}. Green
kernels from Section~\ref{sec:ch23} are now tied to a single physical branch;
\(\mathcal{S}_\omega\) is single-valued on \(\mathcal{F}_{\mathrm{SF}}\), while
\(\mathcal{R}_\omega\) and composed measurement readouts remain rank-deficient on
\(\mathcal{J}_\Omega\) \cite{harrington2001}. Section~\ref{sec:ch25} evaluates
radiation integrals, flux meters, and near/intermediate/far partitions on the
qualified propagating subspace; Section~\ref{sec:ch27} completes non-radiating
null-space structure \cite{stratton1941}. Regime qualification
Eq.~\eqref{eq:ch2.regime-operator-screen} and transport criterion
Eq.~\eqref{eq:ch2.transport-qualified-radiation} remain mandatory before
angular-transport claims on operator outputs \cite{stratton1941}.
Outgoing-wave selection fixes the exterior resolvent branch assumed throughout radiation
integrals; realizability limits---which accessible modes finite geometry can
excite---are established in \ref{ch:04} (\emph{Source $\mathcal{S}_{\mathrm{real}}$ and Spectral Truncation}).
\section{Radiation Integrals and Source-to-Field Representations}
\label{sec:ch25}
Uniqueness and radiation conditions in Section~\ref{sec:ch24} fix the exterior problem. Section~\ref{sec:ch25} writes the source integrals that range engineers discretize on wire and surface meshes \cite{stratton1941}. The same integrals appear in compact-range far-field transforms \cite{hansen1988}.
% PHYSICAL-MEMORY:sec:ch25
% PHYSICAL-MEMORY:sec:ch25-b2
\begin{physicalvignette}[Offset-fed reflector on a satellite bus]
A second bench observation matters on Corporate-feed phased array behind a radome. When technicians swap only the receive probe while holding the transmit chain fixed, reciprocity checks whether the reported coupling is property of the structure or artifact of a mismatched flux gauge \cite{stratton1941}. Failed reciprocity usually traces to a near-field tap, a radome standing-wave, or a probe polarization basis that was never aligned with the transmit standard \cite{hansen1988}.
Rayleigh and Stratton treated such ranges as boundary-value laboratories long before digital near-field transforms became routine \cite{stratton1941}. Harrington's method-of-moments culture reinforced the same discipline: metal outline first, harmonic headline second \cite{harrington2001}.
The measurement logic above is made precise in Subsection~\ref{sec:ch25} \cite{stratton1941}.
\end{physicalvignette}
\begin{physicalvignette}[Offset-fed reflector on a satellite bus]
Consider Offset-fed reflector on a satellite bus. Before any coefficient table is printed, technicians record frequency, cable loss, and the mechanical outline that fixes electrical size. The first walk-through of the range is always tactile: flange bolts, probe standoff, turntable boresight, and the wattmeter reference on the receive chain \cite{hansen1988}.
On this hardware the near zone is visually impressive---standing reactive structure, rapid phase rotation, sometimes a visible VSWR ripple at the feed---yet the export question is measured only after outgoing selection on a large sphere \cite{harrington2001}. What the camera sees in the chamber is not automatically what the far-zone pattern claims \cite{stratton1941}.
Engineers therefore learn to separate three readings: what the feed can synthesize on declared support, what propagating content survives on the radiative range, and what a finite instrument can recover at declared tolerance \cite{devaney2004}. Confusing those readings produces the rank gaps and watt failures that later sections quantify \cite{pal2026}.
Rayleigh and Stratton treated such ranges as boundary-value laboratories long before digital near-field transforms became routine \cite{stratton1941}. Harrington's method-of-moments culture reinforced the same discipline: metal outline first, harmonic headline second \cite{harrington2001}.
Subsection~\ref{sec:ch25} states the formal export condition that this range scene presupposes \cite{stratton1941}.
\end{physicalvignette}
Section~\ref{sec:ch24} fixed the outgoing branch under which exterior solutions are unique.
An engineer on a range nonetheless needs evaluated fields and pattern amplitudes at declared
observation points---not proof of existence alone \cite{poynting1884}. Before any
far-zone asymptotic is applied, the field surrounding a compact source occupies three spatial
bands: a reactive near zone where $1/r^{3}$ and $1/r^{2}$ dyadic terms dominate what probes
measure through $\mathcal{S}_\omega[\Jimp]$; an intermediate transition where storage and export
compete; and a radiative far zone where transverse $e^{-\jj k_0 r}/r$ waves support pattern and
gain definitions through $\mathcal{R}_\omega[\Jimp]=\Pi_{\mathrm{rad}}\,\mathcal{S}_\omega[\Jimp]$
\cite{stratton1941}. Distance labels and regime indices are fixed in
Subsection~\ref{sec:ch252} before asymptotic reductions; the integral representations below are
the constructive route to $\mathcal{S}_\omega$ in all three bands \cite{stratton1941}. A gain trace
on $\Gamma_R$ is a pattern meter reading on $\mathcal{R}_\omega[\Jimp]$ after export selection;
a port $|S_{11}|$ meter on a horn feed reads coupling through $\mathcal{S}_\omega[\Jimp]$ and
may include reactive storage unrelated to far-zone export \cite{harrington2001}. A
method-of-moments solution on a horn mouth illustrates the distinction: port coupling integrals
supply $|S_{11}|$, while a radiation integral on $\Gamma_R$ supplies far-zone gain; both descend
from the same outgoing Green convolution but meter different observables
\cite{harrington2001}.
The excitation-to-flux sequence $\Jimp\mapsto\mathcal{S}_\omega[\Jimp]\mapsto\mathcal{R}_\omega[\Jimp]$
is single-valued on $\mathcal{F}_{\mathrm{SF}}$ at the field level yet rank-deficient at the
source level through Eq.~\eqref{eq:ch2.ker-R-omega} \cite{stratton1941}. This
section develops the radiation integral Eq.~\eqref{eq:ch2.radiation-integral} to
evaluable outgoing integrals and ties range meters to
$\mathcal{R}_\omega[\Jimp]=\Pi_{\mathrm{rad}}\,\mathcal{S}_\omega[\Jimp]$ rather than to
the full pair returned by Green convolution
Eq.~\eqref{eq:ch2.electric-dyadic-convolution} \cite{stratton1941}. Poynting flux and
transport indices from \ref{ch:01} apply here; dyadic kernels and
Sommerfeld completion remain in Sections~\ref{sec:ch23}--\ref{sec:ch24}.
Let \(\Omega_s=\operatorname{supp}\Jimp\) be compact with characteristic scale \(a\),
vacuum wavenumber \(k_0=\omega\sqrt{\varepsilon_0\mu_0}\), and wavelength
\(\lambda_0=2\pi/k_0\). Section~\ref{sec:ch25} organizes evaluation in four stages:
\begin{equation}
\label{eq:ch2.section25-evaluation-chain}
\Jimp
\ \xrightarrow{\ \mathcal{S}_\omega\ }
(\E,\H)\in\mathcal{F}_{\mathrm{SF}}
\ \xrightarrow{\ \Pi_{\mathrm{rad}}\ }
\mathcal{R}_\omega[\Jimp]
\ \xrightarrow{\ \mathcal{F}_{\mathrm{flux}}\ }
\text{power, pattern, norm reports},
\end{equation}
where \(\mathcal{F}_{\mathrm{flux}}\) denotes the flux and operator-norm functionals
developed below by reference to \ref{ch:01} flux identities
\cite{stratton1941}. Equation~\eqref{eq:ch2.section25-evaluation-chain} is the section-level
restatement of Eq.~\eqref{eq:ch2.measurement chain-composition}: stage (iii) export must precede
stage (iv) metering \cite{harrington2001}. Applying \(\mathcal{F}_{\mathrm{flux}}\) to
\(\mathcal{S}_\omega[\Jimp]\) without \(\Pi_{\mathrm{rad}}\) imports reactive circulation
into quantities designed for outgoing transverse radiation
(Section~\ref{sec:ch01.2.4}) \cite{stratton1941}.
The outgoing integral realization of \(\mathcal{R}_\omega\) for
\(\mathbf{r}\notin\Omega_s\) is
\begin{equation}
\label{eq:ch2.radiation-integral-outgoing-upgrade}
\E_{\mathrm{rad}}(\mathbf{r})
=
-\jj\omega\mu_0
\int_{\Omega_s}
\overline{\mathbf{G}}_{E,\omega}^{\mathrm{(out)}}(\mathbf{r},\mathbf{r}')\cdot\Jimp(\mathbf{r}')\,d^{3}r',
\qquad
\H_{\mathrm{rad}}
=
-\frac{1}{\jj\omega\mu_0}\,\curl\E_{\mathrm{rad}},
\end{equation}
In Eq.~\eqref{eq:ch2.radiation-integral-outgoing-upgrade}, \(\mathbf{r}\) is the observation
point, \(\mathbf{r}'\) ranges over the source volume \(\Omega_s\), \(\Jimp(\mathbf{r}')\) is
the impressed electric current density at the source point, \(\mu_0\) is the vacuum permeability,
\(\omega\) is the angular frequency, and \(\overline{\mathbf{G}}_{E,\omega}^{\mathrm{(out)}}\)
is the outgoing electric dyadic from Eq.~\eqref{eq:ch2.outgoing-green-branch}; the magnetic
field \(\H_{\mathrm{rad}}\) follows from \(\E_{\mathrm{rad}}\) by the source-free Maxwell curl
relation in the observation variable \(\mathbf{r}\) \cite{stratton1941}. Equation
\eqref{eq:ch2.radiation-integral-outgoing-upgrade} upgrades
Eq.~\eqref{eq:ch2.radiation-integral} after Sections~\ref{sec:ch23}--\ref{sec:ch24};
the integral subsection below supplies magnetic couplings, equivalent surfaces, and routes
that reduce to the same $\mathcal{R}_\omega$ image \cite{harrington2001}. Near coincidence
$\mathbf{r}\to\Omega_s$ requires regularization from Subsection~\ref{sec:ch232}; evaluation
uses Eq.~\eqref{eq:ch2.radiation-integral-outgoing-upgrade} only where the observation
point lies outside $\Omega_s$ or on declared equivalent surfaces with trace matching
\cite{stratton1941}. Regime qualification
Eq.~\eqref{eq:ch2.regime-operator-screen} must be satisfied before flux meters attach to
any shell integral \cite{stratton1941}.
As a worked example, for \(\Jimp\) supported in a ball of radius \(a\ll\lambda_0\), \(\mathcal{S}_\omega[\Jimp]\)
includes dominant reactive \(1/r^{2}\), \(1/r^{3}\) terms on spheres with \(r\sim a\),
while \(\mathcal{R}_\omega[\Jimp]\) carries leading \(1/r\) transverse radiation with
\(k_0 r\) phase for \(r\gg a\) \cite{jackson1999}. Integrating Poynting flux on \(S_R\) with
\(R\sim a\) conflates storage with export; only when \(k_0 R\gg 1\) and \(R\gg a\) does shell
power approach the \(R\)-independent value predicted by far-zone reduction
\cite{harrington2001}.
As a worked example, near-field probes sample \(\mathcal{S}_\omega[\Jimp]\) on a surface enclosing \(\Omega_s\)
\cite{stratton1941}. Transforming those samples to a far-zone pattern implements
Eq.~\eqref{eq:ch2.radiation-integral-outgoing-upgrade} only when outgoing radiation condition
and equivalent-source uniqueness on the exterior side are satisfied---that is, when exterior
radiative uniqueness (Subsection~\ref{sec:ch243}) holds \cite{harrington2001}.
Distinct impressed currents in \(\ker\mathcal{R}_\omega\) produce identical patterns while
differing in reactive near fields; near scans alone do not resolve that null space without
additional constraints (Subsection~\ref{sec:ch244}) \cite{stratton1941}.
\subsection{Source-to-field integral representations}
\label{sec:ch251}
% HARRINGTON-OPEN:sec:ch251
Retarded source integrals are what method-of-moments codes discretize on antennas \cite{harrington2001}. Far-zone extraction is a limit on the same integral \cite{jackson1999}. Maxwell's retarded export to $\Gamma_R$ remains the governing constraint \cite{maxwell1873}.
The representation begins from retarded potentials driven by $\Jimp$ on compact $\Omega_s$, not from a plane-wave expansion \cite{maxwell1873}. The quantity being organized is the outgoing field pair $(\E,\H)$; angular labels attach only after $\Pi_{\mathrm{rad}}$ \cite{sommerfeld1949}. An alternative route through static potentials would omit radiation damping and misstate reactive storage \cite{jackson1999}. On a horn range, this step is what a MoM mesh implements before any pattern plot \cite{harrington2001}.
% PHYSICAL-MEMORY:sec:ch251
\begin{physicalvignette}[Pyramidal horn on a compact range]
Consider a standard-gain pyramidal horn bolted to a compact-range feed mount at ten gigahertz. Range technicians first measure flange return loss, then walk the quiet zone with a field probe on a linear slide. The near zone within one wavelength of the mouth shows rapid phase rotation and a visually ``fat'' reactive collar that never appears on the published far-zone cut. Students often mistake that collar for extra directivity; in practice it is stored energy that the outward flux integral on a large sphere must exclude before any watt report is trusted \cite{hansen1988}.
The same horn on a turntable illustrates the single-integral picture used throughout this section. Rotating the hardware through a few degrees slides the main lobe on the receiver display while total transmitted power, corrected for cable loss, remains flat within calibration uncertainty. That invariance is the laboratory face of a retarded source integral: one impressed-current model, one propagation law, two observation domains depending on whether the probe sits in the reactive zone or in the Fraunhofer region \cite{harrington2001}.
Rayleigh and Stratton treated such feeds as canonical apertures long before digital near-field transforms became routine \cite{stratton1941}. Harrington's culture of method-of-moments practice reinforced the same habit: discretize the metal outline that actually radiates, then extract the far pattern only after the scan distance satisfies the range equation recorded on the calibration sheet \cite{harrington2001}.
Subsection~\ref{sec:ch251} states the source-to-field integral that this compact-range scene presupposes \cite{stratton1941}.
\end{physicalvignette}
Before any volume quadrature is written, the range engineer already knows which observable the
campaign will report \cite{harrington2001}. A calibrated pattern probe on $\Gamma_R$ integrates
tangential $\E$ or $\H$ against a known polarization and yields a voltage proportional to a
directional amplitude---a discrete sample of the far-zone pattern carried by
$\mathcal{R}_\omega[\Jimp]$, not of the full pair $\mathcal{S}_\omega[\Jimp]$
\cite{stratton1941}. A shell wattmeter closes the same budget through
$\mathcal{P}_{\mathrm{rad}}$ on $\Gamma_R$ after $\Pi_{\mathrm{rad}}$ \cite{stratton1941}. Near-field
scanners on a hull enclosing $\Omega_s$ record $\mathcal{S}_\omega[\Jimp]$ and require a transform
before pattern meters apply \cite{hansen1988}. The volume integrals below are the constructive
route to those meter readings: they assemble $(\E,\H)$ from impressed $\Jimp$, after which export
selection maps the result to the pattern quantities the range records \cite{stratton1941}.
Radiation integrals assemble far-zone fields as superpositions of current moments weighted by
outgoing phase \cite{stratton1941}. Four laboratory settings motivate the
representations that follow. On an antenna range, probes on a large sphere $\Gamma_R$ sample
the field produced by a feed in the compact range; the measured voltage is a discretization of
a surface integral over $\Gamma_R$, not a direct reading of impressed $\Jimp$
\cite{harrington2001}. Horn illumination fixes $\Jimp$ on the mouth through a waveguide mode or
aperture integral; the pattern on the range wall appears only after those currents convolve
against the outgoing Green function \cite{collin2001}. Aperture radiation replaces the volume
support by equivalent currents on $\Gamma_{\mathrm{eq}}$ when the interior is inaccessible
\cite{stratton1941}. Near-field scanners reconstruct far-zone quantities from dense samples on
a hull enclosing $\Omega_s$---a transform that succeeds only when outgoing uniqueness and export
selection are respected \cite{hansen1988}.
The representation begins from retarded potentials driven by $\Jimp$ on compact $\Omega_s$,
organizes the outgoing pair $(\E,\H)$ before angular labels attach, and coincides with what
method-of-moments codes discretize before far-zone extraction
\cite{maxwell1873}. Static-potential routes omit radiation damping
and misstate reactive storage; they are not used here \cite{jackson1999}. Outgoing kernels,
regularization, and branch selection are presupposed from Sections~\ref{sec:ch23}--\ref{sec:ch24};
distance-regime asymptotics and flux meters follow in the subsections below.
Assumptions. Fix \(\omega>0\), compact \(\operatorname{supp}\Jimp=\Omega_s\Subset\Omega\), admissible
data Eq.~\eqref{eq:ch1.admissible-data-class}, charge continuity
Eq.~\eqref{eq:ch1.charge-continuity-harmonic}, and outgoing radiation condition
Eqs.~\eqref{eq:ch2.sommerfeld-outgoing-class}--\eqref{eq:ch2.outgoing-selector-def} on
exterior problems \cite{stratton1941}. Observation points \(\mathbf{r}\) lie outside
\(\Omega_s\) or on equivalent surfaces with trace matching declared below; coincidence
limits use Subsection~\ref{sec:ch232} \cite{harrington2001}. Kernels
\(\overline{\mathbf{G}}_{E,\omega}^{\mathrm{(out)}}\),
\(\overline{\mathbf{G}}_{H,\omega}^{\mathrm{(out)}}\), and scalar
\(G_\omega^{\mathrm{(out)}}\) denote outgoing branches fixed by
Eq.~\eqref{eq:ch2.outgoing-green-branch} \cite{stratton1941}.
The electric volume integral is the constructive form of the dyadic convolution route: each
source element \(\Jimp(\mathbf{r}')\,d^{3}r'\) contributes a dipole-like response that
propagates outward with retarded phase, and the total field is their coherent sum
\cite{stratton1941}. On a MoM mesh, patch currents discretize the same sum; on $\Gamma_R$, pattern
probes sample the far-zone limit after $\mathcal{R}_\omega$ has removed reactive storage that the
volume integral still carries at finite $r$ \cite{harrington2001}. In the symbols below,
$\mathbf{r}$ is the observation coordinate, $\mathbf{r}'$ the source coordinate integrated over
$\Omega_s$, $\E_{\mathrm{hom}}$ any exterior illumination not generated by $\Jimp$, and the
prefactor $-\jj\omega\mu_0$ ties the integral to the locked phasor convention
$\exp(-\jj\omega t)$ \cite{stratton1941}. The primary volume representation is
\begin{equation}
\label{eq:ch2.outgoing-electric-volume-integral}
\E(\mathbf{r})
=
-\jj\omega\mu_0
\int_{\Omega_s}
\overline{\mathbf{G}}_{E,\omega}^{\mathrm{(out)}}(\mathbf{r},\mathbf{r}')\cdot\Jimp(\mathbf{r}')\,d^{3}r'
\;+\;\E_{\mathrm{hom}}(\mathbf{r}),
\end{equation}
\begin{equation}
\label{eq:ch2.outgoing-magnetic-from-curl-integral}
\H(\mathbf{r})
=
\frac{1}{\jj\omega\mu_0}\,\curl\E(\mathbf{r}),
\end{equation}
Equations~\eqref{eq:ch2.outgoing-electric-volume-integral}--\eqref{eq:ch2.outgoing-magnetic-from-curl-integral}
package the dyadic route as an evaluable export map: the electric field is a volume
convolution against the \emph{outgoing} dyad, while the magnetic field is a curl derived
from that electric field in source-free observation regions \cite{stratton1941}. For a
square loop of side \(2a\) with \(k_0 a=1.2\) observed at \(\mathbf{r}\) with
\(k_0|\mathbf{r}|=6\), the integral typically requires \(16\)--\(64\) wavelength-scale
subdivisions of \(\Omega_s\) for \(2\%\) amplitude accuracy when raw
Eq.~\eqref{eq:ch2.free-space-dyadic-green} is used patchwise; self-term corrections from
Subsection~\ref{sec:ch232} dominate error when \(|\mathbf{r}|\to a\)
\cite{harrington2001}. The homogeneous term \(\E_{\mathrm{hom}}\) carries exterior
illumination and must vanish for pure radiation problems \cite{stratton1941}. Equations
\eqref{eq:ch2.outgoing-electric-volume-integral}--\eqref{eq:ch2.outgoing-magnetic-from-curl-integral}
specialize Eqs.~\eqref{eq:ch2.electric-dyadic-convolution}--\eqref{eq:ch2.magnetic-from-electric-dyad}
with outgoing dyads for \(\mathbf{r}\notin\Omega_s\) on simply connected observation
regions; they coincide with Eq.~\eqref{eq:ch2.radiation-integral-outgoing-upgrade} when
\(\E_{\mathrm{hom}}=\mathbf{0}\) \cite{stratton1941}. The prefactor
\(-\jj\omega\mu_0\) matches Eq.~\eqref{eq:ch2.helmholtz-forcing-map} and the locked
phasor axis \cite{stratton1941}. What propagates to a distant probe is the coherent sum of
outgoing dyadic responses $\overline{\mathbf{G}}_{E,\omega}^{\mathrm{(out)}}(\mathbf{r},\mathbf{r}')$
from each source element; what remains reactive at finite $r$ is the non-transverse and
inverse-power content retained in $\mathcal{S}_\omega[\Jimp]$ until $\Pi_{\mathrm{rad}}$ and
regime qualification isolate export \cite{stratton1941}. A near-field scanner
records the full pair $\mathcal{S}_\omega[\Jimp]$; a pattern meter on $\Gamma_R$ with
$k_0 R\gg 1$ approximates $\mathcal{R}_\omega[\Jimp]$ \cite{hansen1988}.
When magnetic impressed sources
\(\mathbf{J}^{\mathrm{(mag)}}_{\mathrm{imp}}\) are primary, duality maps
\(\varepsilon_0\E\leftrightarrow\mu_0\H\) and
\(\Jimp\leftrightarrow\mathbf{J}^{\mathrm{(mag)}}_{\mathrm{imp}}\) yields
\begin{equation}
\label{eq:ch2.outgoing-magnetic-volume-integral}
\H(\mathbf{r})
=
\jj\omega\varepsilon_0
\int_{\Omega_s}
\overline{\mathbf{G}}_{H,\omega}^{\mathrm{(out)}}(\mathbf{r},\mathbf{r}')\cdot
\mathbf{J}^{\mathrm{(mag)}}_{\mathrm{imp}}(\mathbf{r}')\,d^{3}r'
\;+\;\H_{\mathrm{hom}}(\mathbf{r}),
\end{equation}
\begin{equation}
\label{eq:ch2.outgoing-electric-from-curl-magnetic}
\E(\mathbf{r})
=
-\frac{1}{\jj\omega\varepsilon_0}\,\curl\H(\mathbf{r}),
\end{equation}
with \(\overline{\mathbf{G}}_{H,\omega}^{\mathrm{(out)}}\) the magnetic dyad defined
duality-wise from Subsection~\ref{sec:ch231} without re-tabulating its closed form
\cite{stratton1941}. Equations
\eqref{eq:ch2.outgoing-magnetic-volume-integral}--\eqref{eq:ch2.outgoing-electric-from-curl-magnetic}
complete the dual route to Eqs.~\eqref{eq:ch2.outgoing-electric-volume-integral}--\eqref{eq:ch2.outgoing-magnetic-from-curl-integral}
for aperture magnetic currents and dual-equivalent models \cite{stratton1941}.
When both electric and magnetic impressed sources are present, superposition assembles the
full Maxwell pair from the dual volume routes without re-opening gauge or static-potential
paths \cite{stratton1941}. Superposition gives the coupled volume representation
\begin{equation}
\label{eq:ch2.combined-impressed-source-integral}
\begin{aligned}
\E(\mathbf{r})
&=
-\jj\omega\mu_0\!\int_{\Omega_s}\!
\overline{\mathbf{G}}_{E,\omega}^{\mathrm{(out)}}(\mathbf{r},\mathbf{r}')\cdot\Jimp(\mathbf{r}')\,d^{3}r'
-\frac{1}{\jj\omega\varepsilon_0}\,\curl\H^{(\mathrm{mag})}(\mathbf{r})
+\E_{\mathrm{hom}}(\mathbf{r}),\\[4pt]
\H(\mathbf{r})
&=
\jj\omega\varepsilon_0\!\int_{\Omega_s}\!
\overline{\mathbf{G}}_{H,\omega}^{\mathrm{(out)}}(\mathbf{r},\mathbf{r}')\cdot
\mathbf{J}^{\mathrm{(mag)}}_{\mathrm{imp}}(\mathbf{r}')\,d^{3}r'
+\frac{1}{\jj\omega\mu_0}\,\curl\E^{(\mathrm{elec})}(\mathbf{r})
+\H_{\mathrm{hom}}(\mathbf{r}),
\end{aligned}
\end{equation}
where \(\E^{(\mathrm{elec})},\H^{(\mathrm{mag})}\) denote the single-source contributions
from Eqs.~\eqref{eq:ch2.outgoing-electric-volume-integral} and
\eqref{eq:ch2.outgoing-magnetic-volume-integral} \cite{stratton1941}.
Equation~\eqref{eq:ch2.combined-impressed-source-integral} is the first-order Maxwell
integral system on outgoing kernels; it reduces to either branch when one impressed
source vanishes \cite{stratton1941}. Linearity
Eq.~\eqref{eq:ch2.radiation-linearity} follows from Eq.~\eqref{eq:ch2.combined-impressed-source-integral}
on fixed \(\mathcal{B}_\Gamma\) \cite{harrington2001}. The combined integral returns the full
admissible Maxwell pair $\mathcal{S}_\omega[\mathfrak{D}]$: electric and magnetic impressed
sources superpose at observation $\mathbf{r}$ with each contribution evaluated at source
$\mathbf{r}'$ through the outgoing dyads $G$ and $\overline{\mathbf{G}}_{E}$,
$\overline{\mathbf{G}}_{H}$. Pattern, gain, and shell wattmeters do not read this sum directly;
they act on $\mathcal{R}_\omega[\mathfrak{D}]=\Pi_{\mathrm{rad}}\,\mathcal{S}_\omega[\mathfrak{D}]$
after export selection removes reactive and non-propagating fractions that the volume integral
still carries at finite range \cite{stratton1941}.
\begin{figure}[t]
\centering
\ChOneFigBlock{%
\begin{tikzpicture}[ch1fig, node distance=5mm and 7mm]
 \node[ch1box] (jev) {$\Jimp$ volume};
 \node[ch1box, right=7mm of jev] (jmv) {$\mathbf{J}^{\mathrm{(mag)}}_{\mathrm{imp}}$};
 \node[ch1box, right=7mm of jmv] (surf) {Huygens\\$\Gamma_{\mathrm{eq}}$};
 \node[ch1box, below=5mm of jmv] (ge) {$\overline{\mathbf{G}}_{E}^{\mathrm{(out)}}$};
 \node[ch1box, below=5mm of surf] (gh) {$\overline{\mathbf{G}}_{H}^{\mathrm{(out)}}$};
 \node[ch1box, right=7mm of jmv] (ef) {$(\E,\H)$};
 \draw[ch1flow] (jev) -- (ge) -- (ef);
 \draw[ch1flow] (jmv) -- (gh) -- (ef);
 \draw[ch1flow] (surf.east) -- (surf.east |- ef.west);
 \draw[ch1flow] (surf.east |- ef.west) -- (ef.west);
\end{tikzpicture}%
}
\caption{Integral routes: electric volume, magnetic volume, and Huygens surface
representations converge to the same \(\mathcal{S}_\omega\) image on the outgoing branch when
trace matching holds.}
\label{fig:ch2.source-to-field-integral-routes}
\end{figure}
When \(\divg\Jimp=-\jj\omega\rho\) with nonzero \(\rho\), the complete electric field
includes a scalar-potential gradient driven by \(\rho\):
\begin{equation}
\label{eq:ch2.charge-augmented-electric-integral}
\E(\mathbf{r})
=
-\jj\omega\mu_0\!\int_{\Omega_s}\!
\overline{\mathbf{G}}_{E,\omega}^{\mathrm{(out)}}(\mathbf{r},\mathbf{r}')\cdot\Jimp(\mathbf{r}')\,d^{3}r'
-\grad\!\int_{\Omega_s} G_\omega^{\mathrm{(out)}}(\mathbf{r},\mathbf{r}')\,
\frac{\rho(\mathbf{r}')}{\varepsilon_0}\,d^{3}r'
\;+\;\E_{\mathrm{hom}}(\mathbf{r}),
\end{equation}
with \(G_\omega^{\mathrm{(out)}}\) the outgoing scalar kernel from
Eq.~\eqref{eq:ch2.scalar-green-solution-operator} \cite{stratton1941}. Equation
\eqref{eq:ch2.charge-augmented-electric-integral} enforces Gauss coupling without leaving
\(\mathcal{J}_\Omega\); when Eq.~\eqref{eq:ch1.charge-continuity-harmonic} holds and
\(\rho\) is determined by \(\Jimp\), the gradient term is fixed by the same impressed data
\cite{stratton1941}. The charge-augmented integral still evaluates
$\mathcal{S}_\omega[\mathfrak{D}]$: the scalar kernel $G_\omega^{\mathrm{(out)}}(\mathbf{r},\mathbf{r}')$
propagates charge density from $\mathbf{r}'$ to scalar potential at $\mathbf{r}$, while the
dyadic term propagates $\Jimp$; shell meters on $\mathcal{R}_\omega[\Jimp]$ remain insensitive
to how reactive near storage partitions between the two contributions at finite $r$
\cite{harrington2001}. Lorenz-gauge vector-potential reduction
Eq.~\eqref{eq:ch2.lorenz-field-from-potential} collapses to
Eq.~\eqref{eq:ch2.charge-augmented-electric-integral} on uniform exterior domains; gauge
freedom is not re-opened \cite{harrington2001}.
When the interior of a source region is inaccessible, equivalent currents on a enclosing
surface reproduce the exterior field without resolving volume \(\Jimp\); the surface integrals
below are the operator form of that laboratory substitution \cite{stratton1941}. Let
\(\Gamma_{\mathrm{eq}}\) be a closed surface enclosing \(\Omega_s\), with outward
normal \(\hat{\mathbf{n}}\) and equivalent surface densities
\(\mathbf{J}_{s,\mathrm{eq}}\), \(\mathbf{K}_{s,\mathrm{eq}}\) on
\(\Gamma_{\mathrm{eq}}\) in the trace class of Subsection~\ref{sec:ch211}. The exterior
Huygens representation is
\begin{equation}
\label{eq:ch2.huygens-equivalent-integral}
\E(\mathbf{r})
=
-\jj\omega\mu_0\!\int_{\Gamma_{\mathrm{eq}}}\!
\overline{\mathbf{G}}_{E,\omega}^{\mathrm{(out)}}(\mathbf{r},\mathbf{r}')\cdot
\mathbf{J}_{s,\mathrm{eq}}(\mathbf{r}')\,dS'
+\int_{\Gamma_{\mathrm{eq}}}\!
\nabla'\!G_\omega^{\mathrm{(out)}}(\mathbf{r},\mathbf{r}')\times
\mathbf{K}_{s,\mathrm{eq}}(\mathbf{r}')\,dS',
\end{equation}
\begin{equation}
\label{eq:ch2.huygens-magnetic-integral}
\H(\mathbf{r})
=
\jj\omega\varepsilon_0\!\int_{\Gamma_{\mathrm{eq}}}\!
\overline{\mathbf{G}}_{H,\omega}^{\mathrm{(out)}}(\mathbf{r},\mathbf{r}')\cdot
\mathbf{K}_{s,\mathrm{eq}}(\mathbf{r}')\,dS'
-\int_{\Gamma_{\mathrm{eq}}}\!
\nabla'\!G_\omega^{\mathrm{(out)}}(\mathbf{r},\mathbf{r}')\times
\mathbf{J}_{s,\mathrm{eq}}(\mathbf{r}')\,dS',
\end{equation}
for \(\mathbf{r}\) outside the closed surface bounded by \(\Gamma_{\mathrm{eq}}\)
\cite{stratton1941}. Equations
\eqref{eq:ch2.huygens-equivalent-integral}--\eqref{eq:ch2.huygens-magnetic-integral} are
the surface form of the same outgoing propagation operator;
\(\nabla'\) acts on the source coordinate \(\mathbf{r}'\) \cite{stratton1941}. Tangential
field matching
\begin{equation}
\label{eq:ch2.huygens-tangential-match}
\hat{\mathbf{n}}\times\E\big|_{\Gamma_{\mathrm{eq}}^{+}}
=
\hat{\mathbf{n}}\times\E\big|_{\Gamma_{\mathrm{eq}}^{-}},
\qquad
\hat{\mathbf{n}}\times\H\big|_{\Gamma_{\mathrm{eq}}^{+}}
=
\hat{\mathbf{n}}\times\H\big|_{\Gamma_{\mathrm{eq}}^{-}}
\end{equation}
fixes \(\mathbf{J}_{s,\mathrm{eq}}=\hat{\mathbf{n}}\times(\H^{-}-\H^{+})\),
\(\mathbf{K}_{s,\mathrm{eq}}=\hat{\mathbf{n}}\times(\E^{-}-\E^{+})\) in the stratified
sense of Eq.~\eqref{eq:ch1.interface-conditions} \cite{stratton1941}.
Equivalent-source non-uniqueness Eq.~\eqref{eq:ch2.equivalent-source-ambiguity} applies:
distinct \((\mathbf{J}_{s,\mathrm{eq}},\mathbf{K}_{s,\mathrm{eq}})\) pairs related by
\(\ker\mathcal{R}_\omega\) yield identical far-zone radiation \cite{harrington2001}. The
Huygens surface integrals reproduce exterior fields from equivalent currents on
$\Gamma_{\mathrm{eq}}$: observation $\mathbf{r}$ lies outside the closed surface, source
$\mathbf{r}'$ ranges over $\Gamma_{\mathrm{eq}}$, and the same outgoing kernels
$G_\omega^{\mathrm{(out)}}$ and $\overline{\mathbf{G}}_{E,\omega}^{\mathrm{(out)}}$ propagate
tangential data outward. What is measured on a distant range is still
$\mathcal{R}_\omega[\Jimp]$ after $\Pi_{\mathrm{rad}}$; equivalent currents that differ only by
$\ker\mathcal{R}_\omega$ leave that measurement unchanged while altering reactive near fields
through $\mathcal{S}_\omega$ \cite{stratton1941}.
\phantomsection\label{thm:ch2.source-to-field-integral}
Let \(\mathfrak{D}=(\Jimp,\mathbf{J}^{\mathrm{(mag)}}_{\mathrm{imp}},\rho,\mathcal{B}_\Gamma)\)
be admissible with outgoing radiation condition on exterior problems. Then the physical solution
\(\mathcal{S}_\omega[\mathfrak{D}]=(\E,\H)\) is given by
Eqs.~\eqref{eq:ch2.combined-impressed-source-integral}--\eqref{eq:ch2.charge-augmented-electric-integral}
with outgoing kernels, or by Eqs.~\eqref{eq:ch2.huygens-equivalent-integral}--\eqref{eq:ch2.huygens-tangential-match}
when equivalent surface data replace volume sources \cite{stratton1941}.
The representation agrees with Eq.~\eqref{eq:ch2.outgoing-solution-operator} and
Green--BVP equivalence Eq.~\eqref{eq:ch2.green-bvp-equivalence} on the outgoing branch
fixed by exterior radiative uniqueness (Subsection~\ref{sec:ch243}).
Subsection~\ref{sec:ch231} defined dyadic kernels satisfying
Eq.~\eqref{eq:ch2.dyadic-green-defining-operator}; Subsection~\ref{sec:ch232} justified
convolution on \(\mathcal{J}_\Omega\); Section~\ref{sec:ch24} selected the outgoing
branch. Substituting \(\overline{\mathbf{G}}_{E,\omega}^{\mathrm{(out)}}\) into
Eq.~\eqref{eq:ch2.electric-dyadic-convolution} yields
Eq.~\eqref{eq:ch2.outgoing-electric-volume-integral}. Magnetic and coupled forms follow
from duality and superposition \cite{stratton1941}. Huygens formulas follow by replacing
volume sources with surface distributions that reproduce the same tangential traces
\cite{harrington2001}. Uniqueness on \(\mathcal{F}_{\mathrm{SF}}\) when exterior radiative
uniqueness (Subsection~\ref{sec:ch243}) holds identifies the integral representation
with \(\mathcal{S}_\omega\) \cite{stratton1941}.
The volume and surface integrals above return the full admissible pair $\mathcal{S}_\omega[\mathfrak{D}]$;
pattern meters and gain reports act only after export selection through $\mathcal{R}_\omega$ and
$\Pi_{\mathrm{rad}}$ \cite{stratton1941}. Export selection follows evaluation:
\begin{equation}
\label{eq:ch2.integral-radiation-map}
\mathcal{R}_\omega[\mathfrak{D}]
=
\Pi_{\mathrm{rad}}\,
\Big(\E[\mathfrak{D}],\,\H[\mathfrak{D}]\Big),
\end{equation}
with \(\Pi_{\mathrm{rad}}\) from Eq.~\eqref{eq:ch2.outgoing-radiation-refinement}
\cite{stratton1941}. Pattern cuts, gain reports, and shell power meters act on
Eq.~\eqref{eq:ch2.integral-radiation-map}, not on the raw volume integral alone
\cite{harrington2001}.
On the outgoing branch,
\begin{equation}
\label{eq:ch2.integral-resolvent-equivalence}
\mathcal{S}_\omega[\mathfrak{D}]
=
\mathcal{O}_{\mathrm{out}}\,
\mathcal{L}_\omega^{-1}\!
\begin{bmatrix}\mathbf{s}_\omega(\mathfrak{D})\\ \mathbf{g}_\Gamma\end{bmatrix}
=
\mathcal{O}_{\mathrm{out}}\,
\Big(\mathcal{I}_{\mathrm{vol}}[\mathfrak{D}]\Big),
\end{equation}
with \(\mathcal{I}_{\mathrm{vol}}[\mathfrak{D}]\) the combined volume integrals
Eqs.~\eqref{eq:ch2.combined-impressed-source-integral}--\eqref{eq:ch2.charge-augmented-electric-integral},
connecting integral and operator routes through
Eqs.~\eqref{eq:ch2.resolvent-green-equivalence}--\eqref{eq:ch2.outgoing-resolvent-selection}
\cite{stratton1941}. Equation~\eqref{eq:ch2.integral-resolvent-equivalence}
is the constructive content of Eq.~\eqref{eq:ch2.green-bvp-equivalence} for evaluation
purposes: a method-of-moments matrix implements the resolvent route while patch quadrature
implements the integral route, and both must reference the same outgoing branch
\cite{harrington2001}.
When \(\Omega=\cup_j\Omega_j\) with interface traces from Subsection~\ref{sec:ch234},
Eqs.~\eqref{eq:ch2.combined-impressed-source-integral}--\eqref{eq:ch2.charge-augmented-electric-integral}
hold patchwise with local \(k_j\) and outgoing exterior completion on the infinite exterior
region \cite{stratton1941}. Impressed \(\Jimp\) may vanish in passive exterior regions;
equivalent currents on material boundaries enter through
Eqs.~\eqref{eq:ch2.huygens-equivalent-integral}--\eqref{eq:ch2.huygens-tangential-match}
with \(\Gamma_{\mathrm{eq}}\) on scatterer surfaces \cite{harrington2001}. Interior
cavity problems replace outgoing \(G\) with standing-mode Green functions and require
quotienting Eq.~\eqref{eq:ch2.coercivity-quotient}; those representations are excluded from
the outgoing integral family developed here \cite{stratton1941}.
Equations~\eqref{eq:ch2.outgoing-electric-volume-integral}--\eqref{eq:ch2.charge-augmented-electric-integral}
require linear media, fixed \(\omega\), admissible \(\mathfrak{D}\), outgoing radiation
condition on exterior problems, and observation points outside singular coincidence unless
regularization from Subsection~\ref{sec:ch232} is applied \cite{stratton1941}. Incoming
kernels produce different integrals violating
Eq.~\eqref{eq:ch2.retarded-harmonic-kernel-theorem} \cite{jackson1999}. Anisotropic,
moving, or nonlinear media require modified kernels beyond the locked patchwise-scalar class
\cite{stratton1941}. Near-zone reactive dominance does not invalidate the integral; it
signals that \(\mathcal{R}_\omega[\mathfrak{D}]\) may differ sharply from
\(\mathcal{S}_\omega[\mathfrak{D}]\) until Eq.~\eqref{eq:ch2.regime-operator-screen} is
satisfied \cite{harrington2001}.
One operational distinction closes the display below: volume integrals construct the full Maxwell pair $\mathcal{S}_\omega[\mathfrak{D}]$, whereas pattern meters and watt integrators on $\Gamma_R$ read only after outgoing selection has isolated $\mathcal{R}_\omega[\mathfrak{D}]$. A method-of-moments matrix must not be read as a radiation report; Maxwell theory shows that the matrix implements whichever Green branch and boundary completion were declared, and reactive near-zone content may dominate local probes without contributing to qualified export \cite{stratton1941}. A measurement on a closed scan surface therefore records the complete field state unless a transform and regime qualification explicitly project onto the radiative sector. Synthesis, diagnosis, and inversion campaigns must name the same outgoing branch before any two routes---resolvent inversion and direct quadrature---are compared \cite{harrington2001}. Subsection~\ref{sec:ch252} next fixes distance labels that tell the engineer when far-zone asymptotics may replace the full integral without misreading reactive storage as radiated power.
\subsection{Near, intermediate, and far regimes}
\label{sec:ch252}
% HARRINGTON-OPEN:sec:ch252
The far-field approximation replaces distance-dependent amplitudes by $e^{-jkr}/r$ and angular factors \cite{jackson1999}. Range calibrations must sit in the Fraunhofer zone for this step \cite{hansen1988}. Maxwell's retarded export to $\Gamma_R$ remains the governing constraint \cite{maxwell1873}.
Multipole expansions order far-zone fields by angular complexity: the $\ell=1$ dipole term
dominates electrically small feeds, whereas reflector campaigns require many $\ell$ before
sidelobe floors stabilize \cite{stratton1941}. The far-field approximation replaces
distance-dependent amplitudes by $e^{-\jj k_0 r}/r$ and angular factors; range calibrations
must sit in the Fraunhofer zone for that step to be valid \cite{jackson1999}. A pattern
meter on $\Gamma_R$ with $k_0 R\ll 1$ still reads $\mathcal{S}_\omega[\Jimp]$; only when
$k_0 R\gg 1$ and $R\gg a$ does the same probe approximate $\mathcal{R}_\omega[\Jimp]$ through the
far-zone reduction developed below \cite{harrington2001}.
% PHYSICAL-MEMORY:sec:ch252
% PHYSICAL-MEMORY:sec:ch252-b2
\begin{physicalvignette}[Crossed dipoles on a polarization calibration standard]
A second bench observation matters on Open rectangular aperture in an infinite ground screen. When technicians swap only the receive probe while holding the transmit chain fixed, reciprocity checks whether the reported coupling is property of the structure or artifact of a mismatched flux gauge \cite{stratton1941}. Failed reciprocity usually traces to a near-field tap, a radome standing-wave, or a probe polarization basis that was never aligned with the transmit standard \cite{hansen1988}.
Rayleigh and Stratton treated such ranges as boundary-value laboratories long before digital near-field transforms became routine \cite{stratton1941}. Harrington's method-of-moments culture reinforced the same discipline: metal outline first, harmonic headline second \cite{harrington2001}.
The measurement logic above is made precise in Subsection~\ref{sec:ch252} \cite{stratton1941}.
\end{physicalvignette}
\begin{physicalvignette}[Crossed dipoles on a polarization calibration standard]
Consider Crossed dipoles on a polarization calibration standard. Before any coefficient table is printed, technicians record frequency, cable loss, and the mechanical outline that fixes electrical size. The first walk-through of the range is always tactile: flange bolts, probe standoff, turntable boresight, and the wattmeter reference on the receive chain \cite{hansen1988}.
On this hardware the near zone is visually impressive---standing reactive structure, rapid phase rotation, sometimes a visible VSWR ripple at the feed---yet the export question is measured only after outgoing selection on a large sphere \cite{harrington2001}. What the camera sees in the chamber is not automatically what the far-zone pattern claims \cite{stratton1941}.
Engineers therefore learn to separate three readings: what the feed can synthesize on declared support, what propagating content survives on the radiative range, and what a finite instrument can recover at declared tolerance \cite{devaney2004}. Confusing those readings produces the rank gaps and watt failures that later sections quantify \cite{pal2026}.
The far-field approximation replaces distance-dependent amplitudes by $e^{-\jj kr}/r$ and angular factors \cite{jackson1999}. Range calibrations must sit in the Fraunhofer zone for this step \cite{hansen1988}.
Rayleigh and Stratton treated such ranges as boundary-value laboratories long before digital near-field transforms became routine \cite{stratton1941}. Harrington's method-of-moments culture reinforced the same discipline: metal outline first, harmonic headline second \cite{harrington2001}.
Subsection~\ref{sec:ch252} states the formal export condition that this range scene presupposes \cite{stratton1941}.
\end{physicalvignette}
The same source produces all three regions simultaneously; the distinction arises from how
different terms of the radiation integral scale with distance \cite{stratton1941}. Near the
feed, reactive \(1/r^{3}\) and \(1/r^{2}\) terms dominate shell integrals; at intermediate
distances those terms compete with the leading \(e^{-\jj k_0 r}/r\) radiation; only on
sufficiently large spheres does the transverse \(1/r\) term alone support pattern and gain
definitions \cite{harrington2001}.
Let \(\Omega_s=\operatorname{supp}\Jimp\) be compact with characteristic scale \(a\). For
observation radius \(r=|\mathbf{r}|\) from a reference point in \(\Omega_s\), define
\begin{equation}
\label{eq:ch2.distance-regime}
\begin{aligned}
\text{near:}\ & k_0 r\lesssim 1,\\
\text{intermediate:}\ & 1\lesssim k_0 r\lesssim (k_0 a)^{-1},\\
\text{far:}\ & k_0 r\gg k_0 a,\quad r\gg a,
\end{aligned}
\end{equation}
In Eq.~\eqref{eq:ch2.distance-regime},
\(r=|\mathbf{r}|\) is the observation radius (m) from a reference point in
\(\Omega_s=\operatorname{supp}\Jimp\);
\(a\) is the source characteristic scale (m);
\(k_0=\omega\sqrt{\varepsilon_0\mu_0}\) is the vacuum wavenumber (rad\,m\(^{-1}\)); and
the inequalities label reactive near, transition, and radiative far bands before the
constants \(C_{\mathrm{n}},C_{\mathrm{i}},C_{\mathrm{f}}\) are fixed in
Eq.~\eqref{eq:ch2.distance-regime-definitions}.
with implicit constants fixed below \cite{stratton1941}. Equations~\eqref{eq:ch2.distance-regime}
are scale labels, not sharp cutoffs: reactive and radiative terms coexist in the near zone, and
intermediate zones require explicit asymptotic bookkeeping before source integrals reduce to
transverse \(1/r\) radiation \cite{harrington2001}. Pattern and gain definitions used in antenna
theory attach only after the far-zone reduction in
Eq.~\eqref{eq:ch2.far-zone-transverse-radiation} \cite{harrington2001}.
Field-magnitude scalings, far-zone reduction of
Eqs.~\eqref{eq:ch2.outgoing-electric-volume-integral}--\eqref{eq:ch2.combined-impressed-source-integral},
Fresnel--Fraunhofer transition scales, and power-flux convergence on large spheres are
developed below \cite{stratton1941}. Sommerfeld asymptotics
Eqs.~\eqref{eq:ch2.sommerfeld-dyadic-asymptotic} and~\eqref{eq:ch2.far-zone-transverse-radiation}
are presupposed; regime indices
Eqs.~\eqref{eq:ch1.regime-index}--\eqref{eq:ch1.radiative-regime-condition} are invoked by
reference \cite{stratton1941}.
Let \(\Omega_s=\operatorname{supp}\Jimp\) be compact with characteristic radius
\(a=\max_{\mathbf{r},\mathbf{r}'\in\Omega_s}|\mathbf{r}-\mathbf{r}'|/2\), vacuum
wavenumber \(k_0=\omega\sqrt{\varepsilon_0\mu_0}\), wavelength \(\lambda_0=2\pi/k_0\),
and electrical size \(k_0 a\) \cite{stratton1941}. Observation points
\(\mathbf{r}\) satisfy \(|\mathbf{r}|=r\notin\Omega_s\); reference origin is chosen in
\(\Omega_s\). Outgoing kernels from the source integrals above and dyadic structure from
Eq.~\eqref{eq:ch2.free-space-dyadic-green} are presupposed without re-tabulation
\cite{stratton1941}. Constants \(C_1,C_2,\ldots\) below depend on \(k_0,a\) and source
geometry but not on \(r\) in the asymptotic limits stated \cite{harrington2001}.
The locked regime partition is
\begin{equation}
\label{eq:ch2.distance-regime-definitions}
\begin{aligned}
\text{\emph{near (reactive) zone:}}\ & k_0 r\le C_{\mathrm{n}},\\[2pt]
\text{\emph{intermediate (transition) zone:}}\ &
C_{\mathrm{n}}<k_0 r\le C_{\mathrm{i}}(k_0 a),\\[2pt]
\text{\emph{far (radiative) zone:}}\ &
k_0 r> C_{\mathrm{i}}(k_0 a),\quad r> C_{\mathrm{f}}a^{2}/\lambda_0,
\end{aligned}
\end{equation}
with declared positive constants \(C_{\mathrm{n}}=\mathcal{O}(1)\),
\(C_{\mathrm{i}}=\mathcal{O}(1)\) when \(k_0 a=\mathcal{O}(1)\), and Fraunhofer scale
\(r_{\mathrm{F}}=2a^{2}/\lambda_0\) entering \(C_{\mathrm{f}}=\mathcal{O}(1)\)
\cite{stratton1941}. Equations
\eqref{eq:ch2.distance-regime-definitions} upgrade Eq.~\eqref{eq:ch2.distance-regime}:
they are diagnostic bands, not disjoint sharp cutoffs \cite{stratton1941}. For
\(k_0 a\ll 1\) (electrically small sources), \(C_{\mathrm{i}}\) may be taken near unity
and the intermediate zone is narrow; for extended apertures with \(k_0 a\gg 1\), the
Fresnel--Fraunhofer condition \(r\gtrsim r_{\mathrm{F}}\) controls pattern formation
\cite{harrington2001}.
Near-zone field scalings. Evaluating Eq.~\eqref{eq:ch2.outgoing-electric-volume-integral} with
Eq.~\eqref{eq:ch2.free-space-dyadic-green} for \(k_0 r\le C_{\mathrm{n}}\) yields
\begin{equation}
\label{eq:ch2.near-zone-field-scaling}
|\E(\mathbf{r})|
=
\mathcal{O}\!\big(\mu_0\omega\|\Jimp\|_{\mathcal{J}}\,k_0^{-2} r^{-3}\big)
\;+\;
\mathcal{O}\!\big(\mu_0\omega\|\Jimp\|_{\mathcal{J}}\,k_0^{-1} r^{-2}\big),
\end{equation}
with \(\|\Jimp\|_{\mathcal{J}}\) a source amplitude scale on \(\Omega_s\)
\cite{stratton1941}. Equation~\eqref{eq:ch2.near-zone-field-scaling} reflects
the \(R^{-3}\) and \(R^{-2}\) terms in Eq.~\eqref{eq:ch2.free-space-dyadic-green}; the
leading \(r^{-3}\) contribution is reactive (stored) in the sense of
Section~\ref{sec:ch01.2.4} \cite{stratton1941}. Magnetic field follows from
Eq.~\eqref{eq:ch2.outgoing-magnetic-from-curl-integral} with the same near-zone
polynomial-in-\(k_0 r\) ordering \cite{harrington2001}. On shells with
Eq.~\eqref{eq:ch1.regime-index} \(\chi(r)\gg 1\), Poynting flux through
Eq.~\eqref{eq:ch2.radiated-power-screen} is not a credible export meter
\cite{stratton1941}.
When \(C_{\mathrm{n}}<k_0 r\le C_{\mathrm{i}}(k_0 a)\), reactive and radiative terms in
Eq.~\eqref{eq:ch2.free-space-dyadic-green} are comparable:
\begin{equation}
\label{eq:ch2.intermediate-zone-balance}
|\E_{\mathrm{reactive}}(\mathbf{r})|
=
\mathcal{O}\!\big(|\E_{\mathrm{rad}}(\mathbf{r})|\,k_0 r\big),
\end{equation}
with \(\E_{\mathrm{rad}}\) the leading transverse \(e^{-\jj k_0 r}/r\) component and
\(\E_{\mathrm{reactive}}\) the remainder \cite{stratton1941}. Equation
\eqref{eq:ch2.intermediate-zone-balance} diagnoses regime mixing: storage and export coexist
in the full field pair until export is isolated by \(\mathcal{R}_\omega\)
\cite{stratton1941}. For extended sources, the Fresnel parameter
\begin{equation}
\label{eq:ch2.fresnel-parameter}
F
=
\frac{\pi a^{2}}{\lambda_0 r}
=
\frac{k_0 a^{2}}{2r}
\end{equation}
separates Fresnel (\(F\gtrsim 1\)) from Fraunhofer (\(F\ll 1\)) pattern regions on
observation shells \cite{jackson1999}. The parameter \(F\) marks the transition between
those pattern regimes; complete aperture diffraction theory is not re-developed
\cite{stratton1941}.
\begin{figure}[t]
\centering
\ChOneFigBlock{%
\begin{tikzpicture}[ch1fig, node distance=5mm and 7mm]
 \node[ch1box] (src) {Source\\$k_0 a$};
 \node[ch1box, right=7mm of src] (near) {Near\\$k_0 r\lesssim 1$};
 \node[ch1box, right=7mm of near] (mid) {Intermediate\\$F\sim 1$};
 \node[ch1box, right=7mm of mid] (far) {$1/r$\\radiative};
 \node[ch1box, below=5mm of near] (chi) {$\chi(r)\gg 1$};
 \node[ch1box, below=5mm of far] (chi2) {$\chi(r)\ll 1$};
 \draw[ch1flow] (src) -- (near) -- (mid) -- (far);
 \draw[ch1flow] (near.south) -- (near.south |- chi.north);
 \draw[ch1flow] (near.south |- chi.north) -- (chi.north);
 \draw[ch1flow] (far.south) -- (far.south |- chi2.north);
 \draw[ch1flow] (far.south |- chi2.north) -- (chi2.north);
\end{tikzpicture}%
}
\caption{Distance regimes Eq.~\eqref{eq:ch2.distance-regime-definitions}: reactive
dominance near the source, transition through Fresnel parameter
Eq.~\eqref{eq:ch2.fresnel-parameter}, radiative export in the far zone with
Eq.~\eqref{eq:ch1.regime-index} \(\chi(r)\ll 1\).}
\label{fig:ch2.distance-regime-schematic}
\end{figure}
Section~\ref{sec:ch24} \cite{stratton1941}. The far-zone reduction of integrals proceeds as follows: at distances large compared with source extent and wavelength, a pattern probe on $\Gamma_R$
samples the transverse $e^{-\jj k_0 r}/r$ wave whose angular shape is fixed by a Fourier
transform of $\Jimp$ over $\Omega_s$ \cite{jackson1999}. The volume integral at
finite $r$ still evaluates $\mathcal{S}_\omega[\Jimp]$: reactive terms propagate mathematically
but do not export watt-bearing power; only after the far-zone limit does the meter reading
coincide with $\mathcal{R}_\omega[\Jimp]$ through Eqs.~\eqref{eq:ch2.far-field-integral-reduction}--\eqref{eq:ch2.far-field-amplitude-integral}
\cite{stratton1941}. When
\(k_0 r\gg 1\) and \(r\gg r_{\mathrm{F}}\) for extended sources, outgoing dyadic
asymptotics Eq.~\eqref{eq:ch2.sommerfeld-dyadic-asymptotic} reduce
Eq.~\eqref{eq:ch2.outgoing-electric-volume-integral} to
\begin{equation}
\label{eq:ch2.far-field-integral-reduction}
\E(\mathbf{r})
=
\E_0(\theta,\phi)\,\frac{e^{-\jj k_0 r}}{r}
\;+\;\mathcal{O}(r^{-2}),
\qquad
\E_0\cdot\hat{\mathbf{r}}=0,
\end{equation}
\begin{equation}
\label{eq:ch2.far-field-amplitude-integral}
\E_0(\theta,\phi)
=
-\jj\omega\mu_0\,\frac{1}{4\pi}
\int_{\Omega_s}
\Jimp(\mathbf{r}')\,
e^{-\jj k_0\hat{\mathbf{r}}\cdot\mathbf{r}'}\,d^{3}r',
\end{equation}
with \(\hat{\mathbf{r}}=\mathbf{r}/r\) and the transversality constraint enforced by
\(\overline{\mathbf{G}}_{\infty}\) in Eq.~\eqref{eq:ch2.sommerfeld-dyadic-asymptotic}
\cite{stratton1941}. Equation
\eqref{eq:ch2.far-field-integral-reduction} is the far-zone limit of
Eq.~\eqref{eq:ch2.far-zone-transverse-radiation}; the amplitude integral
Eq.~\eqref{eq:ch2.far-field-amplitude-integral} is the vector radiation integral used in
antenna far-field post-analysis \cite{harrington2001}. The amplitude integral is what a
Fraunhofer-range pattern probe measures after $\Pi_{\mathrm{rad}}$: it Fourier-transforms
$\Jimp(\mathbf{r}')$ over the source volume while the reactive inverse-power terms that
$\mathcal{S}_\omega[\Jimp]$ carries at finite $r$ no longer enter at leading order
\cite{harrington2001}. Magnetic field in the far zone
satisfies Eq.~\eqref{eq:ch2.sommerfeld-impedance-ratio} with the same \(\E_0\)
\cite{stratton1941}.
\phantomsection\label{thm:ch2.far-field-approximation}
Let \(\Jimp\in\mathcal{J}_\Omega\) be supported in \(B_a(0)\) with outgoing radiation condition.
For \(r> r_{\min}(a,\lambda_0)\) as in Eq.~\eqref{eq:ch2.distance-regime-definitions},
the field from Eq.~\eqref{eq:ch2.outgoing-electric-volume-integral} satisfies
Eqs.~\eqref{eq:ch2.far-field-integral-reduction}--\eqref{eq:ch2.far-field-amplitude-integral}
with error \(\mathcal{O}\big((k_0 a)^{2}/(k_0 r)\,|\E_0|\,r^{-1}\big)\) for extended
sources and \(\mathcal{O}\big((k_0 r)^{-1}\,|\E_0|\,r^{-1}\big)\) for \(k_0 a\ll 1\)
\cite{stratton1941}.
Substitute Eq.~\eqref{eq:ch2.sommerfeld-dyadic-asymptotic} into
Eq.~\eqref{eq:ch2.outgoing-electric-volume-integral} and expand
\(|\mathbf{r}-\mathbf{r}'|^{-1}\) for \(|\mathbf{r}'|\le a\ll r\). The leading
\(e^{-\jj k_0 r}/r\) term factorizes into \(\E_0(\theta,\phi)\) times phase
\(e^{-\jj k_0\hat{\mathbf{r}}\cdot\mathbf{r}'}\) over the source; transversality follows
from the dyadic structure Eq.~\eqref{eq:ch2.free-space-dyadic-green} at leading order
\cite{stratton1941}. Higher orders in \(a/r\) and \(1/(k_0 r)\) give the
stated error bounds \cite{jackson1999}.
On spheres \(S_r=\{|\mathbf{r}|=r\}\), define shell power
\(P(R)=\Re\oint_{S_R}\mathbf{S}\cdot\hat{\mathbf{r}}\,dS\) as in
Eq.~\eqref{eq:ch2.radiated-power-screen} \cite{stratton1941}. In the far zone,
\begin{equation}
\label{eq:ch2.power-flux-regime-convergence}
|P(r)-P_\infty|
=
\mathcal{O}(r^{-1}),
\qquad
P_\infty
=
\frac{1}{2\eta_0}\oint_{S_r}
\big|\E_0(\theta,\phi)\big|^{2}\,d\Omega,
\end{equation} with \(\E_0\) from Eq.~\eqref{eq:ch2.far-field-amplitude-integral} and \(\eta_0\) from
Eq.~\eqref{eq:ch2.exterior-flux-asymptotic} \cite{stratton1941}. Equation
\eqref{eq:ch2.power-flux-regime-convergence} is the operator form of
Eq.~\eqref{eq:ch1.radiative-regime-condition}: \(P(r)\) stabilizes once reactive
\(r^{-2}\), \(r^{-3}\) contributions no longer dominate the shell integral
\cite{stratton1941}. Evaluating Eq.~\eqref{eq:ch2.radiated-power-screen} on
\(\mathcal{S}_\omega[\Jimp]\) before \(r\gtrsim r_c\) from
Eq.~\eqref{eq:ch1.crossover-radius} confuses storage with export
\cite{harrington2001}.
Distance regimes induce an observation-shell map on operator outputs:
\begin{equation}
\label{eq:ch2.observation-shell-regime-map}
\begin{aligned}
r\in\text{near}
&\ \Longrightarrow\
\text{meter}\ \mathcal{S}_\omega[\Jimp],\\
r\in\text{far}
&\ \Longrightarrow\
\text{meter}\ \mathcal{R}_\omega[\Jimp]\ \text{via Eqs.~\eqref{eq:ch2.far-field-integral-reduction}--\eqref{eq:ch2.far-field-amplitude-integral}},
\end{aligned}
\end{equation}
with intermediate shells requiring explicit \(\chi(r)\) from
Eq.~\eqref{eq:ch1.regime-index} before transport claims \cite{stratton1941}.
When \(k_0 a\ll 1\), Eq.~\eqref{eq:ch2.far-field-amplitude-integral} reduces to the
electric-dipole pattern proportional to
\(\widehat{\boldsymbol{\theta}}\times(\widehat{\boldsymbol{\theta}}\times\mathbf{p})\)
with \(\mathbf{p}=\int \Jimp(\mathbf{r}')\,d^{3}r'\) \cite{jackson1999}. The
near zone extends to \(r\sim\lambda_0\); far-zone reduction applies for
\(k_0 r\gg 1\) \cite{harrington2001}. When \(k_0 a\gg 1\), multiple lobes appear in
\(\E_0(\theta,\phi)\) and Fraunhofer distance \(r_{\mathrm{F}}=2a^{2}/\lambda_0\) dominates
the far-zone threshold \cite{stratton1941}. Angular-spectrum action in
Section~\ref{sec:ch29} filters $\E_0(\theta,\phi)$ without altering the integral
reduction above \cite{stratton1941}.
Equations~\eqref{eq:ch2.distance-regime-definitions}--\eqref{eq:ch2.power-flux-regime-convergence}
require linear vacuum or piecewise-constant media, fixed \(\omega\), outgoing radiation
condition, and observation outside \(\Omega_s\) \cite{stratton1941}. They fail in cutoff
waveguides where propagating modes do not exist
(Eq.~\eqref{eq:ch2.cutoff-condition}), in cavity standing-wave fields, and in lossless open
resonances without outgoing selection \cite{stratton1941}. Anisotropic beams and surface
waves require modified asymptotics beyond the locked free-space dyad template
\cite{harrington2001}. The far-field reduction asserts field asymptotics on the outgoing
branch only; it does not assert uniqueness of sources \cite{stratton1941}.
The regime partition is not a sharp geometric cutoff but a diagnostic on which terms of the same outgoing integral dominate what an instrument measures. Near and intermediate zones may appear to contradict far-zone pattern theory because substantial field energy is visible on small spheres; on the bench, that energy is largely reactive circulation diagnosed through $\mathcal{S}_\omega$, while pattern definitions attach only after $k_0 r\gg 1$ and Fresnel--Fraunhofer qualifications are met \cite{harrington2001}. Shell watt readings must not be identified with radiated power until Eq.~\eqref{eq:ch2.regime-operator-screen} is satisfied: what changes across bands is the relative weight of storage and export in the \emph{same} Maxwell solution, not the governing equations themselves \cite{poynting1884}. Subsection~\ref{sec:ch253} connects these spatial bands to flux, momentum, and operator norms imported from Chapter~\ref{ch:01}, always downstream of the integral evaluation chain established here.
\subsection{Energy, momentum, angular-momentum flux, power, and operator norms}
\label{sec:ch253}
With outgoing integrals and distance labels in place, the display below evaluates what an
observer can meter on large shells. The radiation operator does not merely transport fields:
it transports measurable physical quantities---energy, momentum, and angular
momentum---which ultimately determine what an observer can detect
\cite{poynting1884}. On a calibrated compact range, the wattmeter on $\Gamma_R$
is the primary pattern-adjacent meter: it integrates $\hat{\mathbf{r}}\cdot\mathbf{S}$ on the
radiative branch $\mathcal{R}_\omega[\Jimp]$ after $\Pi_{\mathrm{rad}}$, not on the full
$\mathcal{S}_\omega[\Jimp]$ pair that near-field probes sample at $r\sim a$
\cite{stratton1941}. Power budgets on $\Gamma_R$ must close before torque or
angular-flux claims are promoted \cite{belinfante1940}.
Complex Poynting flux from Eq.~\eqref{eq:ch1.poynting-complex} applies to any admissible
field pair, but power budgets for transport claims use the radiative branch:
% PHYSICAL-MEMORY:sec:ch253
% PHYSICAL-MEMORY:sec:ch253-b2
\begin{physicalvignette}[Pyramidal horn with mode filter at the throat]
A second bench observation matters on Dual-reflector compact range feed. When technicians swap only the receive probe while holding the transmit chain fixed, reciprocity checks whether the reported coupling is property of the structure or artifact of a mismatched flux gauge \cite{stratton1941}. Failed reciprocity usually traces to a near-field tap, a radome standing-wave, or a probe polarization basis that was never aligned with the transmit standard \cite{hansen1988}.
Rayleigh and Stratton treated such ranges as boundary-value laboratories long before digital near-field transforms became routine \cite{stratton1941}. Harrington's method-of-moments culture reinforced the same discipline: metal outline first, harmonic headline second \cite{harrington2001}.
The measurement logic above is made precise in Subsection~\ref{sec:ch253} \cite{stratton1941}.
\end{physicalvignette}
\begin{physicalvignette}[Pyramidal horn with mode filter at the throat]
Consider Pyramidal horn with mode filter at the throat. Before any coefficient table is printed, technicians record frequency, cable loss, and the mechanical outline that fixes electrical size. The first walk-through of the range is always tactile: flange bolts, probe standoff, turntable boresight, and the wattmeter reference on the receive chain \cite{hansen1988}.
On this hardware the near zone is visually impressive---standing reactive structure, rapid phase rotation, sometimes a visible VSWR ripple at the feed---yet the export question is measured only after outgoing selection on a large sphere \cite{harrington2001}. What the camera sees in the chamber is not automatically what the far-zone pattern claims \cite{stratton1941}.
Engineers therefore learn to separate three readings: what the feed can synthesize on declared support, what propagating content survives on the radiative range, and what a finite instrument can recover at declared tolerance \cite{devaney2004}. Confusing those readings produces the rank gaps and watt failures that later sections quantify \cite{pal2026}.
Rayleigh and Stratton treated such ranges as boundary-value laboratories long before digital near-field transforms became routine \cite{stratton1941}. Harrington's method-of-moments culture reinforced the same discipline: metal outline first, harmonic headline second \cite{harrington2001}.
Subsection~\ref{sec:ch253} states the formal export condition that this range scene presupposes \cite{stratton1941}.
\end{physicalvignette}
\begin{equation}
\label{eq:ch2.flux-on-radiative-branch}
P_{\mathrm{rad}}(R)
=
\Re\oint_{\Gamma_R}\hat{\mathbf{r}}\cdot\mathbf{S}\big[\mathcal{R}_\omega[\Jimp]\big]\,dS,
\end{equation}
with \(\mathbf{S}[\cdot]\) the Poynting vector functional imported from
Subsection~\ref{sec:ch01.1.4} and \(\Gamma_R=\{r=R\}\) large enough for far-zone reduction
when reporting export \cite{stratton1941}. Equation~\eqref{eq:ch2.flux-on-radiative-branch}
specializes Eq.~\eqref{eq:ch2.radiated-power-screen} to the outgoing selector chain; it is
not equivalent to integrating \(\mathbf{S}[\mathcal{S}_\omega[\Jimp]]\) unless reactive
content has been shown negligible or orthogonal to the declared observable
\cite{harrington2001}. Poynting balance Eq.~\eqref{eq:ch1.poynting-complex}, Maxwell stress
Eq.~\eqref{eq:ch1.maxwell-stress}, angular balance
Eqs.~\eqref{eq:ch1.angular-momentum-balance}--\eqref{eq:ch1.angular-balance-integral}, and
transport indices Eqs.~\eqref{eq:ch1.transport-admissibility-index}--\eqref{eq:ch1.transport-qualified-decision}
are presupposed from Chapter~\ref{ch:01} \cite{stratton1941}. Radiative--non-radiative
splitting is taken up in the closing subsection below.
Let \((\E,\H)=\mathcal{S}_\omega[\mathfrak{D}]\) and
\((\E_{\mathrm{rad}},\H_{\mathrm{rad}})=\mathcal{R}_\omega[\mathfrak{D}]\) be as in the
evaluation chain Eq.~\eqref{eq:ch2.section25-evaluation-chain} \cite{stratton1941}. Flux meters on
spherical shells \(S_r=\{|\mathbf{r}|=r\}\) assume \(r\) lies in the far zone of
Eq.~\eqref{eq:ch2.distance-regime-definitions} unless noted otherwise, with
\(\mathcal{D}_{J,\mathrm{reg}}=1\) from Eq.~\eqref{eq:ch1.regime-qualified-decision}
\cite{stratton1941}. Complex Poynting vector
\(\mathbf{S}=\tfrac12\E\times\H^*\) and stress tensor \(\mathbf{T}\) follow
Eqs.~\eqref{eq:ch1.poynting-complex} and~\eqref{eq:ch1.maxwell-stress} \cite{stratton1941}.
Angular-flux tensor \(\mathbf{M}=\mathbf{r}\times\mathbf{T}\) uses origin \(\mathbf{r}_0\)
declared in Subsection~\ref{sec:ch01.1.5} \cite{belinfante1940}.
Flux functionals on the radiative branch. Define the radiative Poynting functional
\begin{equation}
\label{eq:ch2.poynting-functional-radiative}
\mathbf{S}_{\mathrm{rad}}[\Jimp]
=
\frac{1}{2}\,\E_{\mathrm{rad}}\times\H_{\mathrm{rad}}^{*},
\qquad
(\E_{\mathrm{rad}},\H_{\mathrm{rad}})=\mathcal{R}_\omega[\Jimp],
\end{equation}
with \(\E_{\mathrm{rad}},\H_{\mathrm{rad}}\) the field components selected by
Eq.~\eqref{eq:ch2.outgoing-radiation-refinement} \cite{stratton1941}. Equation
\eqref{eq:ch2.poynting-functional-radiative} specializes
Eq.~\eqref{eq:ch1.poynting-complex} to the operator output on which export claims are
allowed per Eq.~\eqref{eq:ch2.transport-qualified-radiation} \cite{harrington2001}. The
reactive counterpart \(\mathbf{S}_{\mathrm{re}}=\tfrac12(\E-\E_{\mathrm{rad}})\times(\H-\H_{\mathrm{rad}})^*\)
is not a transport meter unless separately qualified \cite{stratton1941}. Stress and
angular-flux functionals on the radiative branch are
\begin{equation}
\label{eq:ch2.stress-functional-radiative}
\mathbf{T}_{\mathrm{rad}}[\Jimp]
=
\mathbf{T}(\E_{\mathrm{rad}},\H_{\mathrm{rad}}),
\qquad
\mathbf{M}_{\mathrm{rad}}[\Jimp]
=
\mathbf{r}\times\mathbf{T}_{\mathrm{rad}}[\Jimp],
\end{equation}
with \(\mathbf{T}(\cdot,\cdot)\) the Maxwell stress map Eq.~\eqref{eq:ch1.maxwell-stress}
\cite{stratton1941}. Equations
\eqref{eq:ch2.poynting-functional-radiative}--\eqref{eq:ch2.stress-functional-radiative}
specialize \ref{ch:01} flux objects; they introduce no new conservation law
\cite{stratton1941}.
Shell power on the radiative branch is the bounded linear functional
\begin{equation}
\label{eq:ch2.power-flux-functional}
\mathcal{P}_{\mathrm{rad}}(r;\Jimp)
=
\Re\oint_{S_r}\hat{\mathbf{r}}\cdot\mathbf{S}_{\mathrm{rad}}[\Jimp]\,dS,
\end{equation}
which coincides with Eq.~\eqref{eq:ch2.flux-on-radiative-branch} and
Eq.~\eqref{eq:ch2.radiated-power-screen} when \((\E,\H)\) in those equations is replaced
by \(\mathcal{R}_\omega[\Jimp]\) \cite{stratton1941}. In the far zone,
Eqs.~\eqref{eq:ch2.far-field-integral-reduction}--\eqref{eq:ch2.far-field-amplitude-integral}
give
\begin{equation}
\label{eq:ch2.far-zone-power-from-pattern}
\mathcal{P}_{\mathrm{rad}}(r;\Jimp)
=
\frac{1}{2\eta_0}\oint_{S_r}
\big|\E_0(\theta,\phi)\big|^{2}\,d\Omega
\;+\;\mathcal{O}(r^{-1}),
\end{equation}
with \(\eta_0=\sqrt{\mu_0/\varepsilon_0}\) and \(\E_0\) from
Eq.~\eqref{eq:ch2.far-field-amplitude-integral} \cite{stratton1941}. Equation
\eqref{eq:ch2.far-zone-power-from-pattern} is the operator form of
Eq.~\eqref{eq:ch2.power-flux-regime-convergence} and
Eq.~\eqref{eq:ch1.radiative-regime-condition} \cite{stratton1941}. On a calibrated compact
range, \(\mathcal{P}_{\mathrm{rad}}\) is the wattmeter reading on \(\Gamma_R\) after export
selection---the quantity an antenna engineer reports as radiated power \cite{harrington2001}.
Evaluating \(\mathcal{P}_{\mathrm{rad}}\) on \(\mathcal{S}_\omega[\Jimp]\) before
\(r\gtrsim r_c\) from Eq.~\eqref{eq:ch1.crossover-radius} imports reactive storage into
the power budget \cite{harrington2001}.
Linear momentum flux meter. Net momentum export through \(S_r\) is
\begin{equation}
\label{eq:ch2.momentum-flux-functional}
\mathbf{F}_{\mathrm{rad}}(r;\Jimp)
=
\Re\oint_{S_r}\mathbf{T}_{\mathrm{rad}}[\Jimp]\cdot\hat{\mathbf{r}}\,dS,
\end{equation}
the time-averaged force functional inherited from
Eq.~\eqref{eq:ch1.linear-momentum-balance} \cite{stratton1941}. In the far zone with
Eq.~\eqref{eq:ch2.sommerfeld-impedance-ratio},
\begin{equation}
\label{eq:ch2.far-zone-momentum-flux}
\mathbf{F}_{\mathrm{rad}}(r;\Jimp)
=
\frac{1}{c}\,\hat{\mathbf{r}}\,\mathcal{P}_{\mathrm{rad}}(r;\Jimp)
\;+\;\mathcal{O}(r^{-1}),
\end{equation}
linking momentum export to radiated power on the propagating branch
\cite{stratton1941}. A scattering analyst interpreting radiation pressure on a target uses
\(\mathbf{F}_{\mathrm{rad}}\) as the time-averaged force carried by the outgoing wave
\cite{stratton1941}. Equation~\eqref{eq:ch2.far-zone-momentum-flux} expresses
Eq.~\eqref{eq:ch1.energy-momentum-directional-ratio} in the far-zone limit for unit
\(\hat{\mathbf{q}}=\hat{\mathbf{r}}\) \cite{stratton1941}. Momentum meters attach to
\(\mathcal{R}_\omega[\Jimp]\), not to reactive near-zone circulation in
\(\mathcal{S}_\omega[\Jimp]\setminus\mathcal{R}_\omega[\Jimp]\)
\cite{harrington2001}.
Angular-momentum flux meter. Axis torque about unit direction \(\hat{\mathbf{u}}\) is the functional
\begin{equation}
\label{eq:ch2.angular-flux-functional}
\tau_{\hat{\mathbf{u}}}(r;\Jimp)
=
-\hat{\mathbf{u}}\cdot\Re\oint_{S_r}
\mathbf{M}_{\mathrm{rad}}[\Jimp]\cdot\hat{\mathbf{r}}\,dS,
\end{equation}
specializing Eq.~\eqref{eq:ch1.torque-observable} at harmonic steady state to the
radiative branch \cite{belinfante1940}. Net angular-momentum flux through
\(S_r\) is
\begin{equation}
\label{eq:ch2.angular-momentum-flux-vector}
\mathbf{J}_{\mathrm{rad}}(r;\Jimp)
=
\Re\oint_{S_r}\mathbf{M}_{\mathrm{rad}}[\Jimp]\cdot\hat{\mathbf{r}}\,dS,
\end{equation}
with \(\tau_{\hat{\mathbf{u}}}=-\hat{\mathbf{u}}\cdot\mathbf{J}_{\mathrm{rad}}\)
\cite{stratton1941}. Orbital--spin splits Eq.~\eqref{eq:ch1.angular-density-split} are
not re-derived; any angular report must close through
Eq.~\eqref{eq:ch1.angular-balance-integral} on the declared control volume
\cite{belinfante1940}. Directional concentration uses
Eq.~\eqref{eq:ch1.angular-flux-directional-index} evaluated on
\(\mathbf{M}_{\mathrm{rad}}[\Jimp]\) \cite{stratton1941}.
\begin{figure}[t]
\centering
\ChOneFigBlock{%
\begin{tikzpicture}[ch1fig, node distance=5mm and 7mm]
 \node[ch1box] (j) {$\Jimp$};
 \node[ch1box, right=7mm of j] (r) {$\mathcal{R}_\omega$};
 \node[ch1box, right=7mm of r] (f) {Flux\\functionals};
 \node[ch1box, right=7mm of f] (n) {Norms};
 \node[ch1box, below=5mm of f] (p) {$\mathcal{P}_{\mathrm{rad}}$};
 \node[ch1box, below=5mm of r] (reg) {$\mathcal{D}_{J,\mathrm{reg}}$};
 \draw[ch1flow] (j) -- (r) -- (f) -- (n);
 \draw[ch1flow] (f.south) -- (f.south |- p.north);
 \draw[ch1flow] (f.south |- p.north) -- (p.north);
 \draw[ch1flow] (reg.south) -- (reg.south |- r.north);
 \draw[ch1flow] (reg.south |- r.north) -- (r.north);
\end{tikzpicture}%
}
\caption{Flux measurement sequence: regime qualification, then radiative-branch meters only.}
\label{fig:ch2.flux-meter-measurement chain}
\end{figure}
Subsection~\ref{sec:ch01.2.3} fixed transport $\mathcal{D}_J=1$ through conditions
\(\mathcal{C}_1\)--\(\mathcal{C}_3\) on field states. On operator outputs,
\begin{equation}
\label{eq:ch2.transport-triad-on-R}
\begin{aligned}
\mathcal{C}_1[\Jimp]
&\iff
\text{Eq.~\eqref{eq:ch1.am-net-transport-condition} holds on}\ \mathbf{M}_{\mathrm{rad}}[\Jimp],\\
\mathcal{C}_2[\Jimp]
&\iff
\text{Eq.~\eqref{eq:ch1.angular-balance-integral} holds for}\ (\E_{\mathrm{rad}},\H_{\mathrm{rad}}),\\
\mathcal{C}_3[\Jimp]
&\iff
\text{Eq.~\eqref{eq:ch1.propagative-support} holds with}\ 
\Pi=\Pi_{\mathrm{rad}}\mathcal{R}_\omega[\Jimp],
\end{aligned}
\end{equation}
with indices \(\Lambda_J\), \(\mathcal{K}_J\), and \(\mathcal{R}_\tau\) from
Eqs.~\eqref{eq:ch1.transport-admissibility-index}--\eqref{eq:ch1.torque-consistency-residual}
evaluated on the same radiative fields \cite{stratton1941}. Operational acceptance
\(\mathcal{D}_J=1\) in Eq.~\eqref{eq:ch1.transport-qualified-decision} therefore means
transport interpretation is attached to \(\mathcal{R}_\omega[\Jimp]\) after
\(\mathcal{D}_{J,\mathrm{reg}}=1\), not to raw \(\mathcal{S}_\omega[\Jimp]\)
\cite{stratton1941}. The transport qualification of Section~\ref{sec:ch01.2.3}
applies unchanged once \((\E_{\mathrm{rad}},\H_{\mathrm{rad}})\) replaces the full pair on
the measurement shell \cite{stratton1941}.
Radiation operator norms. Field energy on the radiative branch uses the inherited bilinear form
Eq.~\eqref{eq:ch2.maxwell-energy-form} restricted to \(\mathcal{F}_{\mathrm{rad}}\):
\begin{equation}
\label{eq:ch2.radiation-energy-norm}
\|\mathbf{u}_{\mathrm{rad}}\|_{\mathcal{E}}^{2}
=
\mathfrak{a}_\omega\big(\mathbf{u}_{\mathrm{rad}},\mathbf{u}_{\mathrm{rad}}\big),
\qquad
\mathbf{u}_{\mathrm{rad}}=(\E_{\mathrm{rad}},\H_{\mathrm{rad}}),
\end{equation}
with \(\mathfrak{a}_\omega\) from Subsection~\ref{sec:ch202} \cite{stratton1941}. Shell
\(L^{2}\) power density norms for pattern reporting are
\begin{equation}
\label{eq:ch2.radiation-pattern-norm}
\|\E_0\|_{L^{2}(S_r)}^{2}
=
\oint_{S_r}\big|\E_0(\theta,\phi)\big|^{2}\,d\Omega,
\qquad
\|\Jimp\|_{\mathcal{J}}=\Big(\int_{\Omega_s}|\Jimp|^{2}\,dV\Big)^{1/2},
\end{equation}
linking source amplitude to far pattern through
Eq.~\eqref{eq:ch2.far-field-amplitude-integral} \cite{harrington2001}. The
radiation operator norm on admissible sources is
\begin{equation}
\label{eq:ch2.radiation-operator-norm}
\|\mathcal{R}_\omega\|_{\mathrm{op}}
=
\sup_{\|\Jimp\|_{\mathcal{J}}=1}
\|\mathcal{R}_\omega[\Jimp]\|_{\mathcal{E}},
\end{equation}
bounded on fixed outgoing domains when resolvent bound
Eq.~\eqref{eq:ch2.resolvent-bounded} holds \cite{stratton1941}. Observation
maps compose on the right, yielding
\begin{equation}
\label{eq:ch2.flux-meter-composition}
\Lambda_{\mathrm{meas}}
=
\mathcal{O}_{\mathrm{meas}}\circ\mathcal{O}_{\mathrm{far}}\circ\mathcal{R}_\omega,
\qquad
\|\Lambda_{\mathrm{meas}}\|_{\mathrm{op}}
\le
\|\mathcal{O}_{\mathrm{meas}}\|_{\mathrm{op}}\,
\|\mathcal{R}_\omega\|_{\mathrm{op}},
\end{equation}
consistent with Eq.~\eqref{eq:ch2.reconstruction-observation-map} and rank cap
Eq.~\eqref{eq:ch2.reconstruction-rank-cap} \cite{stratton1941}. For a range engineer,
\(\|\mathcal{R}_\omega\|_{\mathrm{op}}\) bounds how strongly a unit-amplitude feed can excite
the far pattern; the composition \(\Lambda_{\mathrm{meas}}\) in
Eq.~\eqref{eq:ch2.flux-meter-composition} is the formal image of gain calibration followed
by pattern sampling on \(S_r\) \cite{harrington2001}.
\phantomsection\label{prop:ch2.far-zone-flux-stabilization}
Let \(\Jimp\) be admissible with \(\mathcal{D}_{J,\mathrm{reg}}=1\) and let \(r\) lie
in the far zone of Eq.~\eqref{eq:ch2.distance-regime-definitions}. Then
\(\mathcal{P}_{\mathrm{rad}}(r;\Jimp)\), \(\mathbf{F}_{\mathrm{rad}}(r;\Jimp)\), and
\(\mathbf{J}_{\mathrm{rad}}(r;\Jimp)\) satisfy
Eqs.~\eqref{eq:ch2.far-zone-power-from-pattern}--\eqref{eq:ch2.far-zone-momentum-flux}
with errors \(\mathcal{O}(r^{-1})\), and \(\mathcal{P}_{\mathrm{rad}}(r;\Jimp)\to P_\infty>0\)
whenever Eq.~\eqref{eq:ch1.radiative-regime-condition} holds on the radiative branch
\cite{stratton1941}.
Substitute Eqs.~\eqref{eq:ch2.far-field-integral-reduction}--\eqref{eq:ch2.far-field-amplitude-integral}
into Eqs.~\eqref{eq:ch2.poynting-functional-radiative}--\eqref{eq:ch2.momentum-flux-functional}
and integrate on \(S_r\). Sommerfeld ratio Eq.~\eqref{eq:ch2.sommerfeld-impedance-ratio}
gives Eq.~\eqref{eq:ch2.far-zone-momentum-flux}; power integral reduces to
Eq.~\eqref{eq:ch2.far-zone-power-from-pattern} \cite{stratton1941}. Convergence
follows from Eq.~\eqref{eq:ch2.power-flux-regime-convergence} \cite{stratton1941}.
Equations~\eqref{eq:ch2.poynting-functional-radiative}--\eqref{eq:ch2.radiation-operator-norm}
require harmonic steady state, linear passive media, outgoing radiation condition, and regime-qualified
shells \cite{stratton1941}. They fail on reactive-dominated shells with
\(\chi(r)\gg 1\) from Eq.~\eqref{eq:ch1.regime-index}, in cutoff guides, and when
\(\mathcal{S}_\omega[\Jimp]\) is used without \(\mathcal{R}_\omega\) projection
\cite{stratton1941}. Material stress debates (Minkowski versus Abraham) are not re-opened;
vacuum-side shells use Eq.~\eqref{eq:ch1.maxwell-stress} \cite{stratton1941}. Gauge-dependent
near-zone angular densities are excluded from \(\mathbf{M}_{\mathrm{rad}}[\Jimp]\) unless
closure through Eq.~\eqref{eq:ch1.angular-balance-integral} is verified
\cite{belinfante1940}.
As a worked example, compute \(\Lambda_J(r)\) from Eq.~\eqref{eq:ch1.transport-charge-continuous-index} using
\(\mathbf{M}_{\mathrm{rad}}[\Jimp]\) on \(S_r\) with \(k_0 r=10\) and compare to
\(k_0 r=0.5\) \cite{stratton1941}. Only the outer shell with
\(\mathcal{D}_{J,\mathrm{reg}}=1\) and \(\mathcal{C}_1\)--\(\mathcal{C}_3\) satisfied
admits transport interpretation per Eq.~\eqref{eq:ch1.transport-qualified-decision}
\cite{stratton1941}. Flux meters are therefore downstream observables: they presuppose the
integral evaluation chain, the far-zone reduction when applicable, and export selection through
$\mathcal{R}_\omega$ before any watt, momentum, or angular-momentum budget is reported on
$\Gamma_R$ \cite{harrington2001}.
Flux and norm functionals are not independent postulates added to Maxwell theory; they are laboratory readouts that presuppose the same export selection already enforced on pattern meters. Power measured on a reactive shell must be distinguished from power reported after $\Pi_{\mathrm{rad}}$: the former samples $\mathcal{S}_\omega[\Jimp]$, the latter the transverse outgoing sector carried by $\mathcal{R}_\omega[\Jimp]$ \cite{stratton1941}. Momentum and angular-momentum budgets add not new dynamics but conserved-flux accounting on the qualified radiative image once regime screens from \ref{ch:01} are satisfied \cite{belinfante1940}. Subsection~\ref{sec:ch254} decomposes the same solution map into radiative and non-radiative components.
\subsection{Radiative and non-radiative components}
\label{sec:ch254}
A compact-range engineer can measure strong electric and magnetic fields within one wavelength
of a horn mouth while the far-zone gain remains modest: reactive energy stored near the feed
is large, yet little of it crosses a distant measurement sphere as export
\cite{poynting1884}. Conversely, currents in $\ker\mathcal{R}_\omega$ can sustain
elaborate near-field structure while producing no far-field signature---the classical
non-radiating source \cite{harrington2001}. Inverse synthesis and scattering campaigns inherit
this split: near-field scans and angle-resolved data constrain only the radiative image
$\mathcal{R}_\omega[\Jimp]$, leaving non-radiative content invisible until additional
excitation or realizability constraints are imposed; Chapter~\ref{ch:04} develops that
$\mathcal{S}_{\mathrm{real}}$ question spectrally \cite{stratton1941}.
% PHYSICAL-MEMORY:sec:ch254
% PHYSICAL-MEMORY:sec:ch254-b2
\begin{physicalvignette}[Synthetic aperture radar stripmap feed horn]
A second bench observation matters on Metrology-grade open-ended waveguide probe. When technicians swap only the receive probe while holding the transmit chain fixed, reciprocity checks whether the reported coupling is property of the structure or artifact of a mismatched flux gauge \cite{stratton1941}. Failed reciprocity usually traces to a near-field tap, a radome standing-wave, or a probe polarization basis that was never aligned with the transmit standard \cite{hansen1988}.
Rayleigh and Stratton treated such ranges as boundary-value laboratories long before digital near-field transforms became routine \cite{stratton1941}. Harrington's method-of-moments culture reinforced the same discipline: metal outline first, harmonic headline second \cite{harrington2001}.
The measurement logic above is made precise in Subsection~\ref{sec:ch254} \cite{stratton1941}.
\end{physicalvignette}
\begin{physicalvignette}[Synthetic aperture radar stripmap feed horn]
Consider Synthetic aperture radar stripmap feed horn. Before any coefficient table is printed, technicians record frequency, cable loss, and the mechanical outline that fixes electrical size. The first walk-through of the range is always tactile: flange bolts, probe standoff, turntable boresight, and the wattmeter reference on the receive chain \cite{hansen1988}.
On this hardware the near zone is visually impressive---standing reactive structure, rapid phase rotation, sometimes a visible VSWR ripple at the feed---yet the export question is measured only after outgoing selection on a large sphere \cite{harrington2001}. What the camera sees in the chamber is not automatically what the far-zone pattern claims \cite{stratton1941}.
Engineers therefore learn to separate three readings: what the feed can synthesize on declared support, what propagating content survives on the radiative range, and what a finite instrument can recover at declared tolerance \cite{devaney2004}. Confusing those readings produces the rank gaps and watt failures that later sections quantify \cite{pal2026}.
Rayleigh and Stratton treated such ranges as boundary-value laboratories long before digital near-field transforms became routine \cite{stratton1941}. Harrington's method-of-moments culture reinforced the same discipline: metal outline first, harmonic headline second \cite{harrington2001}.
Subsection~\ref{sec:ch254} states the formal export condition that this range scene presupposes \cite{stratton1941}.
\end{physicalvignette}
The component partition below imports reactive--radiative logic from
Section~\ref{sec:ch01.2.4} and spectral projectors
Eqs.~\eqref{eq:ch2.spectral-projectors}--\eqref{eq:ch2.rad-null-on-evanescent} without
re-proving projector algebra \cite{stratton1941}. Complete non-radiating null-space
classification belongs to Section~\ref{sec:ch27}.
Let \(\mathbf{u}_\omega=(\E,\H)=\mathcal{S}_\omega[\mathfrak{D}]\) and
\(\mathbf{u}_{\mathrm{rad}}=\mathcal{R}_\omega[\mathfrak{D}]\) be as in
Eq.~\eqref{eq:ch2.section25-evaluation-chain} \cite{stratton1941}. Spectral
projectors \(\Pi_{\mathrm{prop}}\), \(\Pi_{\mathrm{ev}}\), \(\Pi_{\mathrm{non}}\) act on
\(\widehat{\mathcal{F}}_\Omega\) as in Subsection~\ref{sec:ch224}; outgoing selector
\(\mathcal{O}_{\mathrm{out}}\) and radiation map satisfy
Eqs.~\eqref{eq:ch2.outgoing-radiation-refinement}--\eqref{eq:ch2.rad-projection-composition}
\cite{stratton1941}. Regime qualification uses
\(\mathcal{D}_{J,\mathrm{reg}}\) from Eq.~\eqref{eq:ch1.regime-qualified-decision} and
criterion Eq.~\eqref{eq:ch2.regime-operator-screen} \cite{stratton1941}.
On admissible \(\mathfrak{D}\),
\begin{equation}
\label{eq:ch2.integral-spectral-decomposition}
\mathcal{S}_\omega[\mathfrak{D}]
=
\Pi_{\mathrm{prop}}\mathcal{S}_\omega[\mathfrak{D}]
+\Pi_{\mathrm{ev}}\mathcal{S}_\omega[\mathfrak{D}]
+\Pi_{\mathrm{non}}\mathcal{S}_\omega[\mathfrak{D}],
\end{equation}
restating Eq.~\eqref{eq:ch2.propagating-evanescent-decomposition} on the integral outputs
constructed above \cite{stratton1941}.
Define the \emph{non-propagating} component
\begin{equation}
\label{eq:ch2.nonpropagating-component-def}
\mathbf{u}_{\mathrm{np}}
=
\Pi_{\mathrm{ev}}\mathcal{S}_\omega[\mathfrak{D}]
+\Pi_{\mathrm{non}}\mathcal{S}_\omega[\mathfrak{D}],
\end{equation}
and the \emph{propagating} component
\(\mathbf{u}_{\mathrm{prop}}=\Pi_{\mathrm{prop}}\mathcal{S}_\omega[\mathfrak{D}]\)
\cite{stratton1941}. Only \(\mathbf{u}_{\mathrm{prop}}\) is eligible for outgoing radiation
selection; evanescent and non-propagating content satisfies
Eq.~\eqref{eq:ch2.rad-null-on-evanescent} \cite{harrington2001}. Equation
\eqref{eq:ch2.integral-spectral-decomposition} applies the \ref{ch:01} partition to
integral outputs: reactive dominance on near shells corresponds to large
\(\|\mathbf{u}_{\mathrm{np}}\|\) relative to \(\|\mathbf{u}_{\mathrm{prop}}\|\) in the sense
of Eq.~\eqref{eq:ch2.subspace-energy-inequality} \cite{stratton1941}.
The radiative component is
\begin{equation}
\label{eq:ch2.radiative-component-def}
\mathbf{u}_{\mathrm{rad}}
=
\mathcal{R}_\omega[\mathfrak{D}]
=
\Pi_{\mathrm{rad}}\,\mathcal{S}_\omega[\mathfrak{D}]
=
\mathcal{O}_{\mathrm{out}}\circ\Pi_{\mathrm{prop}}\,\mathcal{S}_\omega[\mathfrak{D}],
\end{equation}
using Eqs.~\eqref{eq:ch2.outgoing-radiation-refinement} and~\eqref{eq:ch2.rad-projection-composition}
\cite{stratton1941}. The complementary \emph{non-radiative} component is
\begin{equation}
\label{eq:ch2.nonradiative-component-def}
\mathbf{u}_{\mathrm{nr}}
=
\mathcal{S}_\omega[\mathfrak{D}]-\mathbf{u}_{\mathrm{rad}}
=
\big(\mathrm{id}-\Pi_{\mathrm{rad}}\big)\mathcal{S}_\omega[\mathfrak{D}].
\end{equation}
Equations~\eqref{eq:ch2.radiative-component-def}--\eqref{eq:ch2.nonradiative-component-def}
partition every admissible field pair into export and non-export parts without altering
Maxwell dynamics: \(\mathbf{u}_\omega=\mathbf{u}_{\mathrm{rad}}+\mathbf{u}_{\mathrm{nr}}\)
\cite{stratton1941}. Non-radiative content includes evanescent, non-propagating, and
incoming-propagating fractions annihilated by $\Pi_{\mathrm{rad}}$; large near-field
magnitude therefore does not establish export \cite{stratton1941}. Sources in
$\ker\mathcal{R}_\omega$ produce $\mathbf{u}_{\mathrm{rad}}=\mathbf{0}$ while
$\mathbf{u}_{\mathrm{nr}}\neq\mathbf{0}$ (Eq.~\eqref{eq:ch2.ker-R-omega})
\cite{stratton1941}.
\begin{figure}[t]
\centering
\ChOneFigBlock{%
\begin{tikzpicture}[ch1fig, node distance=5mm and 7mm]
 \node[ch1box] (s) {$\mathcal{S}_\omega[\mathfrak{D}]$};
 \node[ch1box, right=7mm of s] (pp) {$\Pi_{\mathrm{prop}}$};
 \node[ch1box, right=7mm of pp] (pr) {$\Pi_{\mathrm{rad}}$};
 \node[ch1box, below=5mm of s] (np) {$\mathbf{u}_{\mathrm{nr}}$};
 \node[ch1box, below=5mm of pp] (ev) {$\Pi_{\mathrm{ev}}+\Pi_{\mathrm{non}}$};
 \draw[ch1flow] (s) -- (pp) -- (pr);
 \draw[ch1flow] (s.south) -- (s.south |- np.north);
 \draw[ch1flow] (s.south |- np.north) -- (np.north);
 \draw[ch1flow] (s.east) -- (s.east |- ev.west);
 \draw[ch1flow] (s.east |- ev.west) -- (ev.west);
\end{tikzpicture}%
}
\caption{Radiative/non-radiative split: \(\mathbf{u}_{\mathrm{rad}}\) is
\(\Pi_{\mathrm{rad}}\mathcal{S}_\omega\); \(\mathbf{u}_{\mathrm{nr}}\) collects all
non-export content including Eq.~\eqref{eq:ch2.nonpropagating-component-def}.}
\label{fig:ch2.rad-nr-component-split}
\end{figure}
Figure~\ref{fig:ch2.spectral-radiation-criterion} from Subsection~\ref{sec:ch224} on integral
outputs \cite{stratton1941}. Using energy norm Eq.~\eqref{eq:ch2.field-energy-norm}, define
\begin{equation}
\label{eq:ch2.radiative-energy-functional}
\mathcal{E}_{\mathrm{rad}}[\mathfrak{D}]
=
\|\mathbf{u}_{\mathrm{rad}}\|_{\mathcal{E}}^{2},
\qquad
\mathcal{E}_{\mathrm{nr}}[\mathfrak{D}]
=
\|\mathbf{u}_{\mathrm{nr}}\|_{\mathcal{E}}^{2},
\end{equation}
with \(\|\cdot\|_{\mathcal{E}}\) from Eq.~\eqref{eq:ch2.radiation-energy-norm} on each
component \cite{stratton1941}. Shell power splits analogously:
\begin{equation}
\label{eq:ch2.component-power-budget}
\mathcal{P}_{\mathrm{tot}}(r)
=
\mathcal{P}_{\mathrm{rad}}(r;\Jimp)
+\mathcal{P}_{\mathrm{nr}}(r;\Jimp),
\end{equation}
where \(\mathcal{P}_{\mathrm{rad}}\) is Eq.~\eqref{eq:ch2.power-flux-functional} and
\begin{equation}
\label{eq:ch2.nonradiative-power-functional}
\mathcal{P}_{\mathrm{nr}}(r;\Jimp)
=
\Re\oint_{S_r}\hat{\mathbf{r}}\cdot\mathbf{S}_{\mathrm{nr}}[\Jimp]\,dS,
\qquad
\mathbf{S}_{\mathrm{nr}}
=
\tfrac12\,\E_{\mathrm{nr}}\times\H_{\mathrm{nr}}^{*},
\end{equation}
with \((\E_{\mathrm{nr}},\H_{\mathrm{nr}})=\mathbf{u}_{\mathrm{nr}}\)
\cite{stratton1941}. Equation~\eqref{eq:ch2.component-power-budget} is not
a new conservation law; it decomposes Eq.~\eqref{eq:ch2.radiated-power-screen} when both
components are evaluated on the same shell \cite{stratton1941}. In the far zone,
\(\mathcal{P}_{\mathrm{nr}}(r)\to 0\) and \(\mathcal{P}_{\mathrm{rad}}(r)\to P_\infty\) under
Eq.~\eqref{eq:ch1.radiative-regime-condition} \cite{stratton1941}. Near-zone
shells with \(\chi(r)\gg 1\) from Eq.~\eqref{eq:ch1.regime-index} typically have
\(|\mathcal{P}_{\mathrm{nr}}|\gg|\mathcal{P}_{\mathrm{rad}}|\) \cite{stratton1941}.
Regime qualification from Section~\ref{sec:ch01.2.4} becomes
\begin{equation}
\label{eq:ch2.regime-component-criterion}
\mathcal{D}_{J,\mathrm{reg}}=1
\ \Longrightarrow\
\text{transport uses}\ \mathbf{u}_{\mathrm{rad}},\
\mathcal{P}_{\mathrm{rad}},\
\mathbf{M}_{\mathrm{rad}}[\Jimp],
\end{equation}
and forbids interpreting \(\mathbf{u}_{\mathrm{nr}}\), \(\mathcal{P}_{\mathrm{nr}}\), or
\(\Pi_{\mathrm{ev}}\mathcal{S}_\omega[\mathfrak{D}]\) as export channels
\cite{stratton1941}. Equation~\eqref{eq:ch2.regime-component-criterion} is the
component form of Eq.~\eqref{eq:ch2.regime-operator-screen} after
Eqs.~\eqref{eq:ch2.radiative-component-def}--\eqref{eq:ch2.nonradiative-component-def}.
Laboratory errors from Section~\ref{sec:ch01.2.4}---near-zone reactive scans, cutoff-guide
modes, evanescent dominance---appear as large \(\mathcal{E}_{\mathrm{nr}}\) or
\(\|\Pi_{\mathrm{ev}}\mathcal{S}_\omega[\mathfrak{D}]\|\) with small
\(\mathcal{P}_{\mathrm{rad}}\) \cite{stratton1941}. Link to radiation null space. Non-radiative components encode source-side invisible content:
\begin{equation}
\label{eq:ch2.nr-image-of-null-space}
\Jimp\in\ker\mathcal{R}_\omega
\ \Longrightarrow\
\mathbf{u}_{\mathrm{rad}}=\mathbf{0},
\quad
\mathbf{u}_{\mathrm{nr}}=\mathcal{S}_\omega[\Jimp]\neq\mathbf{0}\ \text{in general},
\end{equation}
as recorded in Eq.~\eqref{eq:ch2.null-radiation} \cite{harrington2001}.
Conversely, \(\mathbf{u}_{\mathrm{nr}}=\mathbf{0}\) does not imply trivial
\(\mathcal{S}_\omega[\mathfrak{D}]\): a pure outgoing propagating state may have
\(\mathbf{u}_{\mathrm{rad}}=\mathbf{u}_\omega\) \cite{stratton1941}. Section~\ref{sec:ch272}
classifies \(\ker\mathcal{R}_\omega\); here only the field-side split
Eq.~\eqref{eq:ch2.nonradiative-component-def} is identified as the operator image of
non-export content \cite{harrington2001}.
\phantomsection\label{prop:ch2.far-zone-radiative-dominance}
Let \(\mathcal{D}_{J,\mathrm{reg}}=1\) and let \(r\) lie in the far zone of
Eq.~\eqref{eq:ch2.distance-regime-definitions}. Then
\begin{equation}
\label{eq:ch2.far-zone-component-dominance}
\frac{\mathcal{E}_{\mathrm{nr}}[\mathfrak{D}]}{\mathcal{E}_{\mathrm{rad}}[\mathfrak{D}]+\epsilon}
\to 0,
\qquad
\frac{|\mathcal{P}_{\mathrm{nr}}(r)|}{\mathcal{P}_{\mathrm{rad}}(r)+\epsilon}\to 0,
\end{equation}
as \(r\to\infty\) for fixed \(\mathfrak{D}\), with \(\epsilon>0\) regularizing zero export
\cite{stratton1941}. Equivalently, \(\mathbf{u}_\omega\to\mathbf{u}_{\mathrm{rad}}\) in
far-zone observables metered by Eqs.~\eqref{eq:ch2.far-field-integral-reduction}--\eqref{eq:ch2.power-flux-functional}.
Eq.~\eqref{eq:ch2.rad-null-on-evanescent} annihilates evanescent and non-propagating
contributions in \(\Pi_{\mathrm{rad}}\). Far-zone asymptotics and flux stabilization on
large shells give \(\mathcal{P}_{\mathrm{rad}}(r)\to P_\infty>0\) while reactive terms decay
relatively \cite{stratton1941}. Regime condition
Eq.~\eqref{eq:ch1.radiative-regime-condition} removes reactive dominance in the limit.
Below guide cutoff Eq.~\eqref{eq:ch2.cutoff-condition}, modal content lies in
\(\Pi_{\mathrm{non}}\mathcal{S}_\omega[\mathfrak{D}]\) or
\(\Pi_{\mathrm{ev}}\mathcal{S}_\omega[\mathfrak{D}]\) \cite{stratton1941}. Applying
Eq.~\eqref{eq:ch2.far-field-amplitude-integral} or transport indices to such modes violates
Eq.~\eqref{eq:ch2.regime-component-criterion} as in the TE\(_{10}\) illustration of
Section~\ref{sec:ch01.1.2} \cite{harrington2001}. Spectral projectors and regime logic must
therefore be applied together on the same integral outputs \cite{stratton1941}.
Equations~\eqref{eq:ch2.integral-spectral-decomposition}--\eqref{eq:ch2.far-zone-component-dominance}
require linear media, fixed \(\omega\), outgoing radiation condition, and spectral completion sufficient
for Eqs.~\eqref{eq:ch2.spectral-projectors}--\eqref{eq:ch2.spectral-projector-algebra}
\cite{stratton1941}. Sharp boundaries between subspaces may blur when
loss, anisotropy, or coupled interfaces mix propagating and evanescent sheets
\cite{stratton1941}. Component splits are diagnostic on shells, not substitutes for
exterior radiative uniqueness (Subsection~\ref{sec:ch243}) on the full field
\cite{harrington2001}.
Outgoing source-to-field integrals map impressed data to evaluable fields; distance-regime
labels fix where asymptotic reductions apply; flux and operator-norm meters act on
$\mathcal{R}_\omega[\Jimp]$ after regime qualification; and the radiative--non-radiative split
Eqs.~\eqref{eq:ch2.radiative-component-def}--\eqref{eq:ch2.regime-component-criterion}
separates export from storage on the same integral outputs \cite{stratton1941}. Section~\ref{sec:ch26}
develops spectral decomposition of radiation operators on the representations constructed here
\cite{harrington2001}.
The radiative--non-radiative split is what separates a wattmeter that has been qualified on $\Gamma_R$ from a near-field probe that still sees storage and evanescent overlap in $\mathcal{S}_\omega[\Jimp]$ \cite{harrington2001}. Both components coexist in one Maxwell solution until projectors and outgoing selection name which part export meters may integrate \cite{poynting1884}. Evanescent sheets must not be dismissed as mere mathematical artifacts; however, Maxwell theory shows they carry reactive energy measurable on enclosing scans while remaining null under far-zone angular accessibility maps developed in Section~\ref{sec:ch29} \cite{stratton1941}. Section~\ref{sec:ch26} now organizes the radiative image itself into eigenmodes and recoverable rank on the integral representations constructed in Sections~\ref{sec:ch23}--\ref{sec:ch25}.
\section{Spectral Theory of Radiation Operators}
\label{sec:ch26}
Having followed impressed current through Maxwell completion, outgoing selection, and evaluable radiation integrals. Spectral theory treats radiation export as not a single scalar amplitude but a countable nesting of modes on compact geometries, each with its own coupling to sources and its own observability limit on finite instruments \cite{harrington2001}. Pattern shape must not be identified with a single dominant mode; Maxwell propagation on bounded supports shows instead that many directions can be degenerate under rotation while remaining distinguishable under excitation, and that Fredholm structure controls which target patterns are reachable at all \cite{reed1950}. A laboratory campaign on a circular aperture measures a finite slice of that nesting---never the full operator spectrum---so rank and degeneracy statements below are as much measurement discipline as abstract mathematics \cite{collin2001}.
% PHYSICAL-MEMORY:sec:ch26
% PHYSICAL-MEMORY:sec:ch26-b2
\begin{physicalvignette}[Handset chassis with multiple coupled patches]
A second bench observation matters on Slotted waveguide array in an aircraft radome. When technicians swap only the receive probe while holding the transmit chain fixed, reciprocity checks whether the reported coupling is property of the structure or artifact of a mismatched flux gauge \cite{stratton1941}. Failed reciprocity usually traces to a near-field tap, a radome standing-wave, or a probe polarization basis that was never aligned with the transmit standard \cite{hansen1988}.
Rayleigh and Stratton treated such ranges as boundary-value laboratories long before digital near-field transforms became routine \cite{stratton1941}. Harrington's method-of-moments culture reinforced the same discipline: metal outline first, harmonic headline second \cite{harrington2001}.
The measurement logic above is made precise in Subsection~\ref{sec:ch26} \cite{stratton1941}.
\end{physicalvignette}
\begin{physicalvignette}[Handset chassis with multiple coupled patches]
Consider Handset chassis with multiple coupled patches. Before any coefficient table is printed, technicians record frequency, cable loss, and the mechanical outline that fixes electrical size. The first walk-through of the range is always tactile: flange bolts, probe standoff, turntable boresight, and the wattmeter reference on the receive chain \cite{hansen1988}.
On this hardware the near zone is visually impressive---standing reactive structure, rapid phase rotation, sometimes a visible VSWR ripple at the feed---yet the export question is measured only after outgoing selection on a large sphere \cite{harrington2001}. What the camera sees in the chamber is not automatically what the far-zone pattern claims \cite{stratton1941}.
Engineers therefore learn to separate three readings: what the feed can synthesize on declared support, what propagating content survives on the radiative range, and what a finite instrument can recover at declared tolerance \cite{devaney2004}. Confusing those readings produces the rank gaps and watt failures that later sections quantify \cite{pal2026}.
Rayleigh and Stratton treated such ranges as boundary-value laboratories long before digital near-field transforms became routine \cite{stratton1941}. Harrington's method-of-moments culture reinforced the same discipline: metal outline first, harmonic headline second \cite{harrington2001}.
Subsection~\ref{sec:ch26} states the formal export condition that this range scene presupposes \cite{stratton1941}.
\end{physicalvignette}
Radiation integrals in Section~\ref{sec:ch25} map sources to outgoing fields. Section~\ref{sec:ch26} analyzes that map spectrally: which directions carry power and how degeneracy splits under broken symmetry \cite{harrington2001}. Sidelobe floors on finite apertures are the visible face of spectral decay \cite{hansen1988}.
Every antenna engineer learns, long before operator theory appears in the curriculum, that a
finite structure cannot radiate an unlimited set of independent far-zone patterns. A dish whose
rim spans only a few wavelengths resolves only a modest number of principal beam directions on a
compact range; a spherical near-field scanner with finite probe count tells the same story after
transformation to far-zone coefficients \cite{harrington2001}. Collin and
Harrington treat that compression as structural: it is imposed by Maxwell propagation from compact
sources through outgoing selection to a finite observation shell, not by a particular receiver
model \cite{collin2001}. Stratton and Harrington organized outgoing fields by angular
structure on large spheres; Chu sharpened the consequence that electrical size bounds directional
export at fixed frequency \cite{stratton1941}. Subsection~\ref{sec:ch26.signature-aperture}
introduces the laboratory reading---singular values $\sigma_n$ of
$\mathcal{O}_{\mathrm{meas}}\circ\mathcal{R}_\omega$ and the recoverable rank they define---before
any eigenpair formula is written; later subsections refer back rather than restate that doctrine.
Why spectra arise in radiation physics is therefore not a mathematical accident. Maxwell's
equations are linear in impressed sources on the outgoing branch fixed in Section~\ref{sec:ch24};
superposition of two feeds superposes far-zone amplitudes. On a large sphere $S_r$, orthogonality
of angular patterns under an appropriate energy inner product later organizes the export map
$\mathcal{R}_\omega$ into radiation modes; each carries an \emph{export eigenvalue}
$\lambda_n$ of the radiation operator---ranking qualified power on $\Gamma_R$ and distinct from
any correlation-operator spectrum unless explicitly stated \cite{stratton1941}. Some feeds couple strongly to dominant
modes; others produce elaborate reactive near fields through $\mathcal{S}_\omega[\Jimp]$ yet register
negligibly on a far-zone wattmeter after $\Pi_{\mathrm{rad}}$ \cite{poynting1884}.
Observable far-zone content remains primary; formal harmonic indices organize export but are not
themselves physical modes---that distinction is developed once in
Subsection~\ref{sec:ch26.signature-aperture} and thereafter
\cite{stratton1941}.
The radiation map $\mathcal{R}_\omega$ from Section~\ref{sec:ch25} inherits this ordering.
Spectral analysis names dominant and weak export directions, explains why finite dishes and
finite scan grids lose information, and states when pattern synthesis and source inversion are
well posed \cite{stratton1941}. \ref{ch:03} will supply explicit harmonic bases and tabulated
matrix elements; the present section asks what the export map itself permits at fixed $\omega$
before any particular angular basis is chosen \cite{stratton1941}. A reader who finishes the
section should understand why measured patterns look low-rank in practice, why some synthesis
targets are unrealizable no matter how much numerical effort is invested, and why rotation of an
antenna reorients coefficients without changing exported power when frames are consistently
handled \cite{harrington2001}.
Reading order. Engineers study spectra because measured patterns never carry unlimited independent
information: finite apertures, finite scan grids, and finite electrical size impose rank
limits long before any harmonic post-processor runs out of indices \cite{hansen1988}.
On a compact range the dominant beams are the few principal cuts that survive above the noise
floor; weak modes decay in singular-value cliffs and become experimentally invisible even when
formal expansions can still list them \cite{stratton1941}. The sequence below
should therefore be read not as a mathematical ordering but as a laboratory ordering.
Measurement limitations appear first; spectral structure merely explains them
\cite{hansen1988}.
The spectral account follows the measurement laboratory, not an abstract operator catalog.
Finite aperture observation on $S_r$ fixes how many independent far-zone directions a campaign
can resolve---that ceiling is \emph{finite radiative complexity}. Compactness of the composed
map $\mathcal{K}=\mathcal{O}_{\mathrm{meas}}\circ\mathcal{R}_\omega$ then explains why singular
values $\sigma_n$, ranking the radiation efficiency of each independent source organization
through the actual receiver chain, decay and define recoverable rank
\cite{hansen1988}. Only after that ordering does the section construct export
eigenvalues $\lambda_n$ of $\mathcal{R}_\omega$, decompose radiation modes
$\mathbf{u}_n$ from their exciting source modes $\Jimp^{(n)}$, and complete symmetry,
Fredholm, and geometric filters on what finite hardware can recover \cite{stratton1941}. Let
$\Omega_s=\operatorname{supp}\Jimp$ be compact, $k_0=\omega\sqrt{\varepsilon_0\mu_0}$, and
outgoing field pair $\mathbf{u}_\omega=(\E,\H)\in\mathcal{F}_{\mathrm{SF}}$ from
Eq.~\eqref{eq:ch2.admissible-field-space}. The four stages below record the same program in
compact notation; Subsection~\ref{sec:ch26.signature-aperture} develops $\sigma_n$ decay and
finite recoverable rank once, and later subsections cite it by reference:
\begin{equation}
\label{eq:ch2.section26-spectral-chain}
\mathcal{R}_\omega
\ \xrightarrow{\ \text{eigenmodes}\ }
\{\lambda_n,\mathbf{u}_n\}
\ \xrightarrow{\ \text{symmetry}\ }
\text{degeneracy classes}
\ \xrightarrow{\ \text{Fredholm}\ }
\text{solvability / index}
\ \xrightarrow{\ \text{geometry}\ }
\text{radiative spectra},
\end{equation}
from export-eigenmode ordering through rotational degeneracy, Fredholm solvability, and geometric
filters on recoverable content \cite{stratton1941}.
The primary map remains
\begin{equation}
\label{eq:ch2.spectral-operator-targets}
\mathcal{R}_\omega:\mathcal{J}_\Omega\longrightarrow\mathcal{F}_{\mathrm{rad}},
\qquad
\mathcal{R}_\omega[\Jimp]=\Pi_{\mathrm{rad}}\,\mathcal{S}_\omega[\Jimp],
\end{equation}
from Eq.~\eqref{eq:ch2.radiation-operator-def}, with $\Pi_{\mathrm{rad}}$ selecting outgoing
radiative solutions per Eq.~\eqref{eq:ch2.rad-projection-composition}. In plain language:
$\mathcal{R}_\omega$ is the radiation map from impressed currents to far-zone export;
$\mathcal{J}_\Omega$ is the admissible source region where $\Jimp$ may be placed;
$\mathcal{F}_{\mathrm{rad}}$ is the space of outgoing radiative field pairs observable on
$\Gamma_R$ after export selection; and $\Pi_{\mathrm{rad}}$ is the filter that removes
reactive and non-propagating content before pattern measurement \cite{stratton1941}. Equation
\eqref{eq:ch2.spectral-operator-targets} represents the entire electromagnetic experiment
beginning with impressed current inside the source region and ending with the radiative content
observable on a far-zone measurement surface \cite{harrington2001}. Subsection~\ref{sec:ch26.signature-aperture}
fixes $\sigma_n$ decay and finite recoverable rank on $\mathcal{K}$; the subsections below
construct export eigenpairs, split degeneracy under symmetry, diagnose Fredholm solvability, and
apply geometric filters. A sidelobe floor on a measured pattern is often the visible face of
the aperture-limited stage \cite{hansen1988}.
\subsection{Compact circular aperture and singular-value decay}
\label{sec:ch26.signature-aperture}
Before any eigenvalue formula is written, it helps to watch the phenomenon on a range. On a dish
compact range, place a feed above a circular scan disk on plane $z=0$ and sample tangential field on
$S_r$; on a spherical near-field scanner, place probes on a hull at $R\approx 2\lambda_0$ and
transform to far-zone coefficients \cite{hansen1988}. In either campaign, count how
many independent beam directions remain above the receiver noise floor. The answer depends on disk
or scan radius $a_{\mathrm{ap}}$ and source electrical size, not on how many spherical harmonics a
post-processor can index \cite{chu1948}. Electrical size and scan extent fix a
recoverable rank; singular values $\sigma_n$ of $\mathcal{O}_{\mathrm{meas}}\circ\mathcal{R}_\omega$
measure the radiation efficiency of each independent source organization through the aperture and
receiver chain, and fall off once the modal index exceeds what the geometry can resolve
\cite{stratton1941}. Synthesis harmonic caps therefore overstate physical rank when $\sigma_n$ has already fallen below the
noise floor---harmonics in a post-processor are labels, not independent export channels. Compact
decay of $\sigma_n(\mathcal{K})$ and the finite recoverable rank it defines are stated here once;
later subsections invoke Subsection~\ref{sec:ch26.signature-aperture} by reference rather than
repeating the measurement narrative \cite{harrington2001}.
% PHYSICAL-MEMORY:sec:ch26.signature-aperture
% PHYSICAL-MEMORY:sec:ch26.signature-aperture-b2
\begin{physicalvignette}[Dielectric lens fed by a scalar horn]
A second bench observation matters on Reflector with serrated edge treatment. When technicians swap only the receive probe while holding the transmit chain fixed, reciprocity checks whether the reported coupling is property of the structure or artifact of a mismatched flux gauge \cite{stratton1941}. Failed reciprocity usually traces to a near-field tap, a radome standing-wave, or a probe polarization basis that was never aligned with the transmit standard \cite{hansen1988}.
Rayleigh and Stratton treated such ranges as boundary-value laboratories long before digital near-field transforms became routine \cite{stratton1941}. Harrington's method-of-moments culture reinforced the same discipline: metal outline first, harmonic headline second \cite{harrington2001}.
The measurement logic above is made precise in Subsection~\ref{sec:ch26.signature-aperture} \cite{stratton1941}.
\end{physicalvignette}
\begin{physicalvignette}[Dielectric lens fed by a scalar horn]
Consider Dielectric lens fed by a scalar horn. Before any coefficient table is printed, technicians record frequency, cable loss, and the mechanical outline that fixes electrical size. The first walk-through of the range is always tactile: flange bolts, probe standoff, turntable boresight, and the wattmeter reference on the receive chain \cite{hansen1988}.
On this hardware the near zone is visually impressive---standing reactive structure, rapid phase rotation, sometimes a visible VSWR ripple at the feed---yet the export question is measured only after outgoing selection on a large sphere \cite{harrington2001}. What the camera sees in the chamber is not automatically what the far-zone pattern claims \cite{stratton1941}.
Engineers therefore learn to separate three readings: what the feed can synthesize on declared support, what propagating content survives on the radiative range, and what a finite instrument can recover at declared tolerance \cite{devaney2004}. Confusing those readings produces the rank gaps and watt failures that later sections quantify \cite{pal2026}.
A finite aperture SVD on $\mathcal{R}_\omega$ shows how many independent patterns a disk can synthesize \cite{harrington2001}. GUI mode counts often exceed measured rank \cite{devaney2004}.
Rayleigh and Stratton treated such ranges as boundary-value laboratories long before digital near-field transforms became routine \cite{stratton1941}. Harrington's method-of-moments culture reinforced the same discipline: metal outline first, harmonic headline second \cite{harrington2001}.
Subsection~\ref{sec:ch26.signature-aperture} states the formal export condition that this range scene presupposes \cite{stratton1941}.
\end{physicalvignette}
Consider a circular scan aperture of radius \(a_{\mathrm{ap}}\) on plane \(z=0\) sampling the
far-zone tangential field from a compact source with \(k_0 a=1.5\). The measurement composition
$\mathcal{K}=\mathcal{O}_{\mathrm{meas}}\circ\mathcal{R}_\omega$ is compact on the source and
field spaces declared in Section~\ref{sec:ch213}; its singular values $\sigma_n$ rank radiation
efficiency of source organizations and decay, defining a practical recoverable rank and finite
effective dimension before any harmonic label is attached
\cite{stratton1941}.
\begin{figure}[t]
\centering
\ChOneFigGeo{%
\begin{tikzpicture}[ch1geo, scale=0.90]
 \draw[ch1shell] (0,0) circle (0.75);
 \node[ch1oncanvas, below=8mm] at (0,-0.85) {aperture radius $a_{\mathrm{ap}}$};
 \draw[ch1shellB, dashed] (0,0) circle (0.97);
 \node[ch1oncanvas, right=10mm] at (1.35,0) {$S_r$ far zone};
 \draw[->, thick, gray] (0.6,0.2) -- (1.2,0.35);
 \node[ch1oncanvas, font=\scriptsize] at (0.9,0.55) {$\sigma_n$ decay};
\end{tikzpicture}%
}
\caption{Finite circular aperture: compact observation operator on far-zone fields;
singular-value decay marks angular harmonics that fall below instrument threshold.}
\label{fig:ch2.signature-aperture-svd}
\end{figure}
Compactness from Subsection~\ref{sec:ch213} yields \(\sigma_n(\mathcal{K})\to 0\); for disk
apertures the decay is approximately exponential in modal index when \(k_0 a_{\mathrm{ap}}\) is
moderate \cite{stratton1941}. In practice the engineer reads off a knee in the singular-value
curve: retain modes above the knee for synthesis or inversion; discard the tail and accept the
residual sidelobe floor as a geometry-limited artifact rather than a tuning defect
\cite{hansen1988}. What practical criterion tells the engineer where to truncate
the spectrum? Retaining every mode is impossible: instrument noise, finite aperture extent, and
finite scan count cap recoverable rank, and singular-value decay controls which source
organizations remain physically realizable through the actual receiver chain
\cite{stratton1941}. Truncation is therefore inevitable; the engineering question is
where the knee lies. Truncating at \(N\) modes with
\begin{equation}
\label{eq:ch2.signature-aperture-truncation}
\frac{\|\mathcal{K}-\mathcal{K}_N\|}{\|\mathcal{K}\|}
\le
\frac{\sigma_{N+1}}{\sigma_1},
\end{equation}
where $\mathcal{K}=\mathcal{O}_{\mathrm{meas}}\circ\mathcal{R}_\omega$ is the full measurement
chain from source to recorded samples, $\mathcal{K}_N$ is its rank-$N$ truncated approximation,
$\sigma_1$ is the largest singular value (dominant export channel), and $\sigma_{N+1}$ is the
first discarded singular value \cite{harrington2001}. The ratio $\sigma_{N+1}/\sigma_1$ is the
relative strength of the strongest discarded radiation channel compared with the dominant
channel---a direct laboratory readout of how much pattern error truncation will introduce
\cite{hansen1988}. Eq.~\eqref{eq:ch2.compression-error-bound} bounds pattern error on \(S_r\). At
\(k_0 a_{\mathrm{ap}}=3\), retaining \(N=9\) singular modes typically captures \(95\%\)
of the radiated power resolvable through the disk; harmonics beyond \(N\approx 12\) lie
below a \(-20\,\mathrm{dB}\) threshold relative to \(\sigma_1\)
\cite{harrington2001}.
Figure~\ref{fig:ch2.signature-aperture-svd} foreshadows Chapter~\ref{ch:04}: realizability
and spectral truncation ask which angular components survive compression by
\(\mathcal{K}\), not which indices a formal expansion can list
\cite{devaney2004}. Laboratory reports therefore inherit a finite rank from
geometry before source inversion is attempted; improving SNR without enlarging
\(a_{\mathrm{ap}}\) cannot recover modes annihilated by \(\sigma_n\) below measurement
noise \cite{hansen1988}.
Local notation ($\sigma_n$, recoverable rank). Compose the measurement operator $\mathcal{K}=\mathcal{O}_{\mathrm{meas}}\circ\mathcal{R}_\omega$
on declared source and field spaces. Its singular values $\sigma_n$---ordered
$\sigma_1\ge\sigma_2\ge\cdots\to 0$---are the radiation efficiencies of independent source
organizations as seen through the actual aperture and receiver chain \cite{stratton1941}. The
\emph{recoverable rank} is the largest $N$ for which $\sigma_N$ remains above the instrument noise
floor; modes beyond that cliff are mathematically listable in a harmonic post-processor yet
physically indistinguishable from clutter on $\Gamma_R$ \cite{chu1948}. Compact
$\sigma_n$ decay and finite dimension on $\mathcal{K}$ are fixed in the display below; elsewhere
invoke Subsection~\ref{sec:ch26.signature-aperture} by reference only \cite{harrington2001}.
The spectral relations developed below restate three laboratory facts in compact form before the eigenbasis is
constructed. A \emph{radiation mode} is a normalized far-zone pattern direction
$\mathbf{u}_n$; its \emph{source mode} $\Jimp^{(n)}$ is an admissible current that excites it.
The \emph{export eigenvalue} $\lambda_n$ of the radiation operator $\mathcal{R}_\omega$---not a
correlation-operator eigenvalue unless later stated---ranks how efficiently that radiation mode
reaches qualified flux on $\Gamma_R$ \cite{stratton1941}. each radiating mode carries export efficiency
$|\lambda_n|$ measured by coupled power in that direction. Second, not every desired far-zone
target can be synthesized: obstructions appear as orthogonality conditions on outgoing
solutions. Third, when geometry possesses rotational symmetry, rotating the hardware reorients
modal coefficients without changing eigenvalue ordering within degenerate blocks. The introductory relation
equations record these facts; the subsections that follow interpret them in measurement
language \cite{stratton1941}.
On a compact range, a \emph{radiation mode} appears as a principal far-zone pattern direction
---a normalized cut on $S_r$ that carries qualified export power---while its \emph{source mode}
is the admissible current distribution that excites it \cite{harrington2001}. Dominant beams
correspond to modes with large $|\lambda_n|$ because those organizations couple most efficiently
to the wattmeter on $\Gamma_R$ \cite{stratton1941}. The eigenpair relation below states the
experimentally familiar observation that certain current organizations radiate much more
effectively than others \cite{collin2001}. On a family of radiating modes, the eigenvalue template reads
\begin{equation}
\label{eq:ch2.modal-eigenvalue-equation}
\mathcal{R}_\omega[\Jimp^{(n)}]=\lambda_n\,\mathbf{u}_n,
\qquad
\mathbf{u}_n=P_n\mathcal{S}_\omega[\Jimp^{(n)}],
\end{equation}
with \(\lambda_n\) measuring radiative export efficiency of mode \(n\) from finite impressed
current \(\Jimp^{(n)}\) \cite{stratton1941}. In synthesis and inversion, $\mathbf{g}$ denotes the
desired far-zone radiated field one wishes to realize; solvability fails when $\mathbf{g}$ carries
a component along cokernel directions---an \emph{orthogonality obstruction} that no admissible
source can remove \cite{stratton1941}. Experimentally, such failure appears as a synthesis target
that refuses to converge, a measured residual orthogonal to every reachable pattern, or a
pattern null that persists despite port tuning \cite{harrington2001}. Solvability on outgoing radiative solutions satisfies
\begin{equation}
\label{eq:ch2.spectral-solvability}
\mathcal{R}_\omega+K
\ \text{solvable on}\ \mathcal{J}_\Omega
\iff
\mathbf{g}\perp\ker(\mathcal{R}_\omega+K)^{\dagger},
\end{equation}
with \(K\) compact on the graph domain of impressed sources \cite{stratton1941}.
When a symmetry group \(G\) commutes with \(\mathcal{R}_\omega\), the equation below expresses
a simple laboratory fact: rotating the hardware and then measuring the pattern gives the same
result as measuring first and rotating the resulting pattern \cite{harrington2001}.
\begin{equation}
\label{eq:ch2.degeneracy-covariance}
\mathcal{R}_\omega\,\mathcal{D}(R)[\Jimp]
=
\mathcal{D}(R)\,\mathcal{R}_\omega[\Jimp],
\qquad R\in G,
\end{equation}
with \(\mathcal{D}(R)\) the field representation induced from
Eq.~\eqref{eq:ch1.active-rotation-action} on \(\mathcal{F}_{\mathrm{rad}}\)
\cite{stratton1941}. Modal projectors compose with radiation selection through
$\Pi_{\mathrm{rad}}$ acting as the radiation filter that passes only exportable content,
$P_n$ as mode-selection lenses that isolate individual beam families from a measured pattern,
and $\Pi_{\mathrm{ev}}$, $\Pi_{\mathrm{non}}$ as projectors onto evanescent and
non-propagating sectors that qualified far-zone meters cannot resolve
\cite{stratton1941}.
\begin{equation}
\label{eq:ch2.spectral-resolution-composition}
\Pi_{\mathrm{rad}}
=
\sum_{n\in\mathcal{N}_{\mathrm{rad}}} P_n,
\qquad
P_n\Pi_{\mathrm{ev}}=P_n\Pi_{\mathrm{non}}=\mathbf{0},
\end{equation}
using Eqs.~\eqref{eq:ch2.spectral-projectors}--\eqref{eq:ch2.rad-null-on-evanescent} by reference
\cite{stratton1941}. Spectral decompositions \emph{organize} export; they do not
\emph{create} angular structure---see Subsection~\ref{sec:ch26.signature-aperture} for the
harmonics-versus-modes distinction stated once in measurement terms \cite{stratton1941}.
\subsection{Eigenfunctions and spectral decomposition}
\label{sec:ch261}
% CH02-NARR-EXP:02.261
Discrete modal sums on exterior domains organize exported power into separable angular sectors only after outgoing qualification and flux gauge are declared \cite{stratton1941,jackson1999}. Continuous-spectrum bookkeeping runs in parallel with discrete $\ell$ blocks; mixing the two without measure conventions produces apparent mode creation on finite meshes \cite{hansen1988}.
Source integrals from Section~\ref{sec:ch25} return amplitudes on $\Gamma_R$. The question an
experienced range engineer asks next is practical: how many of those amplitudes remain
independent above the noise floor for a given electrical size and scan extent?
\cite{stratton1941}. The eigenpair construction below answers that question in
normalized form; recoverable rank and $\sigma_n$ cliffs are already fixed in
Subsection~\ref{sec:ch26.signature-aperture}.
% PHYSICAL-MEMORY:sec:ch261
\begin{physicalvignette}[Parabolic dish on a three-axis turntable]
A three-meter dish on a three-axis turntable at C-band is the everyday face of spectral compression on finite support. As frequency rises, the same ribbed reflector first grows a sharper main beam, then encounters a ceiling beyond which additional sidelobe detail cannot be radiated no matter how aggressively the feed is tuned \cite{harrington2001}. Technicians recognize the ceiling as a shoulder floor in the measured pattern, not as a feed-tuning failure.
Singular-value ladders on such dishes are often explained to visitors with a simple picture: the reflector outline sets how many independent far-zone templates can carry appreciable watts, and the feed can only combine those templates, not invent new ones outside the export cone \cite{stratton1941}. GUI mode counters that exceed the measured rank on the range are listing chart rows, not independent radiators \cite{hansen1988}.
Collin and Jackson supply the aperture vocabulary; Harrington supplies the method-of-moments habit of quoting electrical diameter before quoting multipole order \cite{collin2001}. Chu caps the same story in bandwidth language familiar to system engineers \cite{chu1948}.
Subsection~\ref{sec:ch261} states the spectral decomposition that this turntable scene makes visible in measured decibels \cite{stratton1941}.
\end{physicalvignette}
What an eigenmode means physically. A radiation mode is a principal far-zone pattern direction $\mathbf{u}_n$; its source mode
$\Jimp^{(n)}$ is the admissible current distribution that excites it. The export eigenvalue
$\lambda_n$ of $\mathcal{R}_\omega$ measures radiation efficiency $|\lambda_n|$ of that
radiation mode on $\Gamma_R$ \cite{stratton1941}. Modes with large $|\lambda_n|$ are the dominant export channels on
$\Gamma_R$: they are the beams an array designer weights first and the cuts a diagnostician
examines before worrying about sidelobe structure. Modes with small $|\lambda_n|$ appear only
after long integration, narrowband filtering, or exceptional dynamic range; on a typical compact
range they never rise above the clutter floor. Modes with $\lambda_n=0$ may still produce
strong reactive near fields---the balanced-loop illustration from Section~\ref{sec:ch25} is the
archetype---yet contribute nothing to the far-zone wattmeter \cite{devaney2004}. Eigenvalue
ordering is therefore a statement about measurable export, not about how many terms a numerical chart
can index.
On a dish range, dominant eigenmodes appear as the principal cuts in a principal-component
decomposition of sampled $\E_t$ on $S_r$: the first few components carry most resolvable power
\cite{hansen1988}. Rank truncation in synthesis follows the recoverable rank fixed in
Subsection~\ref{sec:ch26.signature-aperture}; the decomposition below passes through
$\mathcal{R}_\omega$ on physically realizable source classes, with coupled power on
$\Gamma_R$ as the physical quantity \cite{stratton1941}.
Antenna-range observations. Harrington's synthesis practice already exploited the ordering: one designs or diagnoses feeds by
how they excite a modest number of dominant modes, not by attempting independent control of every
harmonic index a post-processor can list \cite{harrington2001}. On a compact range, rotating a
horn reorients the principal cuts while leaving total exported power unchanged when probe and
source frames rotate together---a fact developed formally in the rotation discussion below, but
visible in raw pattern data long before covariance formulas appear \cite{collin2001}. When a
measured pattern refuses to match a synthesis target despite ample numerical effort, the fault is
often spectral rather than computational: the target lies in a weakly coupled or forbidden
sector, not in the subspace the aperture and source geometry can reach \cite{stratton1941}. The
eigenpair definition that follows encodes that experience in normalized form.
Propagation projectors Eqs.~\eqref{eq:ch2.spectral-projectors}--\eqref{eq:ch2.rad-null-on-evanescent},
compactness templates from Section~\ref{sec:ch213}, and far-zone amplitudes
Eqs.~\eqref{eq:ch2.far-field-integral-reduction}--\eqref{eq:ch2.far-field-amplitude-integral}
are presupposed \cite{stratton1941}. Explicit vector spherical harmonic bases belong to
Chapter~3; here modal coordinates are fixed at the operator level compatible with
Eq.~\eqref{eq:ch2.far-pattern-modal} \cite{harrington2001}.
Assumptions. Fix \(\omega>0\), physically realizable impressed data \(\mathfrak{D}\) on
\(\mathcal{J}_\Omega\) Eq.~\eqref{eq:ch2.source-space-refined}, outgoing radiation condition on exterior
problems Eqs.~\eqref{eq:ch2.sommerfeld-outgoing-class}--\eqref{eq:ch2.outgoing-selector-def},
and export map Eq.~\eqref{eq:ch2.spectral-operator-targets}
\cite{stratton1941}. Field pairs lie in
\(\mathcal{F}_{\mathrm{SF}}\subset\mathcal{F}_\Omega\) from
Eq.~\eqref{eq:ch2.admissible-field-space}; outgoing radiative images lie in
\(\mathcal{F}_{\mathrm{rad}}\subset\operatorname{ran}\mathcal{R}_\omega\) with energy norm
Eq.~\eqref{eq:ch2.field-energy-norm} \cite{stratton1941}. Spectral completion
\(\widehat{\mathcal{F}}_\Omega\) supporting Eqs.~\eqref{eq:ch2.spectral-projectors}--\eqref{eq:ch2.spectral-projector-algebra} is assumed sufficient for the expansions below
\cite{harrington2001}.
Radiation eigenpairs. Eigenmodes are introduced because measured patterns on $S_r$ are rarely arbitrary: principal-component
analysis of sampled $\E_t$ reveals a small number of dominant export directions, and the eigenbasis
below names those directions in normalized form \cite{hansen1988}. What experiment
reveals is rank ordering---a few strong beams, a long tail of weak modes below the noise floor,
and export-invisible reactive organizations that never register on a far-zone wattmeter
\cite{devaney2004}. The laboratory observations above condense into a normalized triple: export eigenvalue
$\lambda_n$ of $\mathcal{R}_\omega$, radiation mode $\mathbf{u}_n$, and source mode
$\Jimp^{(n)}$ linked through $\mathcal{R}_\omega$ \cite{stratton1941}. Physically,
$|\lambda_n|$ ranks how much qualified power each normalized radiation mode carries to
$\Gamma_R$; mathematically, the triple satisfies the locked eigenpair relation below. A
\emph{radiation eigenpair} is a triple $(\lambda_n,\mathbf{u}_n,\Jimp^{(n)})$ with
\(\lambda_n\in\mathbb{C}\), \(\mathbf{u}_n=(\E_n,\H_n)\in\mathcal{F}_{\mathrm{rad}}\), and
\(\Jimp^{(n)}\in\mathcal{J}_\Omega\) such that
\begin{equation}
\label{eq:ch2.radiation-eigenpair-def}
\mathcal{R}_\omega[\Jimp^{(n)}]=\lambda_n\,\mathbf{u}_n,
\qquad
\|\mathbf{u}_n\|_{\mathcal{U}}=1,
\end{equation}
where $\lambda_n$ is the export eigenvalue (radiation-efficiency score of mode $n$),
$\mathbf{u}_n$ is the normalized radiation mode (principal far-zone pattern direction),
$\Jimp^{(n)}$ is the source mode (admissible current that excites $\mathbf{u}_n$), and
$\|\mathbf{u}_n\|_{\mathcal{U}}=1$ fixes unit radiative energy on $\mathcal{F}_{\mathrm{rad}}$
\cite{stratton1941}. Readers should think of $\lambda_n$ as a radiation-efficiency score
assigned to the $n$th export channel \cite{harrington2001}. Here \(\|\cdot\|_{\mathcal{U}}\) is the field energy norm
Eq.~\eqref{eq:ch2.field-energy-norm} on \(\mathcal{F}_{\mathrm{rad}}\)
\cite{stratton1941}. Equation~\eqref{eq:ch2.radiation-eigenpair-def}
develops Eq.~\eqref{eq:ch2.modal-eigenvalue-equation} to a normalized
locked definition: \(\lambda_n\) measures export gain of mode \(n\) relative to unit radiative
energy, and \(\ker\mathcal{R}_\omega\) corresponds to \(\lambda_n=0\) with
\(\mathcal{S}_\omega[\Jimp^{(n)}]\neq\mathbf{0}\) through
Eq.~\eqref{eq:ch2.ker-R-omega} \cite{stratton1941}. Eigenpairs are defined on the
outgoing branch only; incoming propagating sheets excluded by
Eq.~\eqref{eq:ch2.outgoing-radiation-refinement} \cite{harrington2001}.
Interpretation: modes with $|\lambda_n|\gg 1$ dominate compact-range cuts; modes with
$\lambda_n=0$ are export-invisible yet may carry reactive near fields through
$\mathcal{S}_\omega[\Jimp^{(n)}]$ \cite{devaney2004}.
Modal projectors on field space. Projector $P_n[\mathbf{u}]$ extracts a single beam family from a measured pattern: it isolates
the component of field $\mathbf{u}$ aligned with radiation mode $\mathbf{u}_n$ and discards the
orthogonal complement \cite{stratton1941}. Modal filtering matters because synthesis, diagnosis,
and interference analysis proceed mode by mode rather than on raw harmonic indices
\cite{harrington2001}. For each index \(n\) in a family of radiating modes \(\mathcal{N}\), define a rank-one (or
finite-dimensional) projector \(P_n:\mathcal{F}_\Omega\to\mathcal{F}_{\mathrm{rad}}\) by
\begin{equation}
\label{eq:ch2.modal-projector-def}
P_n[\mathbf{u}]
=
\langle \mathbf{u},\mathbf{u}_n\rangle_{\mathcal{U}}\,\mathbf{u}_n,
\qquad
P_n\circ\Pi_{\mathrm{ev}}=P_n\circ\Pi_{\mathrm{non}}=0,
\end{equation}
using dual pairing inherited from Eq.~\eqref{eq:ch2.maxwell-adjoint-pairing} on qualified
Maxwell states and annihilation of non-propagating content through
Eq.~\eqref{eq:ch2.rad-null-on-evanescent} \cite{stratton1941}. Equation
\eqref{eq:ch2.modal-projector-def} specializes resolution
Eq.~\eqref{eq:ch2.spectral-resolution-composition} once an orthonormal family of radiating modes
\(\{\mathbf{u}_n\}_{n\in\mathcal{N}_{\mathrm{rad}}}\) is chosen. Projectors satisfy
\(P_n P_m=\delta_{nm} P_n\) when modes are orthogonal in \(\mathcal{U}\)
\cite{stratton1941}.
Once eigenmodes are fixed, any physically realizable excitation decomposes into modal
coefficients. The quantity a compact-range measurement ultimately recovers is the export-side
pattern carried by $\mathbf{u}_{\mathrm{rad},n}=P_n\mathcal{R}_\omega[\mathfrak{D}]$; the
source-side amplitude $c_n[\mathfrak{D}]$ measures how strongly the full solution
$\mathcal{S}_\omega[\mathfrak{D}]$ aligns with mode $n$ before export
\cite{stratton1941}. For general physically realizable \(\mathfrak{D}\), define source
coefficients
\begin{equation}
\label{eq:ch2.source-modal-amplitude}
c_n[\mathfrak{D}]
=
\langle \mathcal{S}_\omega[\mathfrak{D}],\mathbf{u}_n\rangle_{\mathcal{U}},
\qquad
\mathbf{u}_{\mathrm{rad},n}
=
P_n\,\mathcal{R}_\omega[\mathfrak{D}]
=
\lambda_n\,c_n[\mathfrak{D}]\,\mathbf{u}_n,
\end{equation}
when the expansion converges in \(\mathcal{U}\) \cite{stratton1941}.
Physically, \(c_n\) is the excitation strength of radiative mode \(n\) before export
selection; \(\lambda_n c_n\) is the coefficient on the normalized far-zone pattern carried
by \(\mathbf{u}_n\) \cite{jackson1999}. Linearity
Eq.~\eqref{eq:ch2.radiation-linearity} implies
\(c_n[a_1\mathfrak{D}_1+a_2\mathfrak{D}_2]=a_1 c_n[\mathfrak{D}_1]+a_2 c_n[\mathfrak{D}_2]\)
for physically realizable superpositions after splitting each term through \(\mathcal{S}_\omega\)
\cite{stratton1941}.
\begin{figure}[t]
\centering
\ChOneFigBlock{%
\begin{tikzpicture}[ch1fig, node distance=5mm and 7mm]
 \node[ch1box] (d) {$\mathfrak{D}$};
 \node[ch1box, right=7mm of d] (s) {$\mathcal{S}_\omega$};
 \node[ch1box, right=7mm of s] (r) {$\mathcal{R}_\omega$};
 \node[ch1box, right=7mm of r] (exp) {$\sum_n \lambda_n c_n P_n$};
 \node[ch1box, below=5mm of s] (cn) {$c_n=\langle\cdot,\mathbf{u}_n\rangle$};
 \draw[ch1flow] (d) -- (s) -- (r) -- (exp);
 \draw[ch1flow] (cn.south) -- (cn.south |- s.north);
 \draw[ch1flow] (cn.south |- s.north) -- (s.north);
\end{tikzpicture}%
}
\caption{Spectral measurement sequence: \(\mathcal{S}_\omega\) supplies coefficients \(c_n\);
\(\mathcal{R}_\omega\) maps them to export amplitudes \(\lambda_n c_n\) on projectors
\(P_n\) Eq.~\eqref{eq:ch2.modal-projector-def}.}
\label{fig:ch2.radiation-spectral-decomposition}
\end{figure}
How can a complicated radiation pattern be reconstructed from simpler building blocks? The
modal expansion below decomposes any radiative field into orthonormal radiation modes, each
carrying a fraction of the export energy budget \cite{stratton1941}. When the radiation eigenbasis is complete on the export subspace, every radiative field admits
a countable expansion in normalized modes; the Parseval identity below states that radiative
energy equals the sum of squared modal coefficients \cite{stratton1941}. When
\(\{\mathbf{u}_n\}_{n\in\mathcal{N}_{\mathrm{rad}}}\) is a complete orthonormal family
in \(\mathcal{F}_{\mathrm{rad}}\) for the problem class, every radiative field
\(\mathbf{v}\in\operatorname{ran}\mathcal{R}_\omega\) admits
\begin{equation}
\label{eq:ch2.radiation-spectral-expansion}
\mathbf{v}
=
\sum_{n\in\mathcal{N}_{\mathrm{rad}}}
\langle \mathbf{v},\mathbf{u}_n\rangle_{\mathcal{U}}\,\mathbf{u}_n,
\qquad
\|\mathbf{v}\|_{\mathcal{U}}^2
=
\sum_{n\in\mathcal{N}_{\mathrm{rad}}}
\big|\langle \mathbf{v},\mathbf{u}_n\rangle_{\mathcal{U}}\big|^2,
\end{equation} with convergence in \(\mathcal{U}\) \cite{stratton1941}. Applying
Eq.~\eqref{eq:ch2.radiation-spectral-expansion} to \(\mathbf{v}=\mathcal{R}_\omega[\mathfrak{D}]\)
gives
\begin{equation}
\label{eq:ch2.radiation-map-modal-sum}
\mathcal{R}_\omega[\mathfrak{D}]
=
\sum_{n\in\mathcal{N}_{\mathrm{rad}}}
\lambda_n\,c_n[\mathfrak{D}]\,\mathbf{u}_n,
\end{equation}
the radiation eigenbasis restatement of compression operator
Eq.~\eqref{eq:ch2.spectral-compression-operator} with explicit eigenvalues and normalized
modes \cite{stratton1941}. The Parseval identity in
Eq.~\eqref{eq:ch2.radiation-spectral-expansion} states the modal energy budget: total radiative
energy equals the sum of squared modal coefficients. Beam decomposition and recoverable
complexity are therefore fixed by how many terms carry appreciable weight before truncation
\cite{hansen1988}. Truncation to \(|\mathcal{N}|\le N\) yields
\(\mathcal{R}_\omega^{(N)}\) from Eq.~\eqref{eq:ch2.spectral-compression-operator}; error
Eq.~\eqref{eq:ch2.compression-error-bound} applies when the tail is compact on source space
\cite{harrington2001}.
\phantomsection\label{thm:ch2.radiation-spectral-decomposition}
Let outgoing radiation condition and domain class be fixed as above. Suppose
\(\{\mathbf{u}_n\}_{n\in\mathcal{N}_{\mathrm{rad}}}\) is a complete orthonormal radiation
eigenbasis in \(\mathcal{F}_{\mathrm{rad}}\) with
Eq.~\eqref{eq:ch2.radiation-eigenpair-def} and projectors
Eq.~\eqref{eq:ch2.modal-projector-def}. Then only modes with $\lambda_n\neq 0$ contribute to
observable radiation: modes with $\lambda_n=0$ lie in $\ker\mathcal{R}_\omega$ and are invisible
to far-zone meters even when they carry reactive near fields \cite{devaney2004}.
\begin{equation}
\label{eq:ch2.radiation-operator-spectral-form}
\mathcal{R}_\omega
=
\sum_{n\in\mathcal{N}_{\mathrm{rad}}}
\lambda_n\,P_n\circ\mathcal{S}_\omega,
\qquad
\Pi_{\mathrm{rad}}
=
\sum_{n\in\mathcal{N}_{\mathrm{rad}}} P_n,
\end{equation}
on physically realizable \(\mathfrak{D}\), with convergence in \(\mathcal{U}\) for fields in
\(\operatorname{ran}\mathcal{R}_\omega\) \cite{stratton1941}. The range
satisfies
\begin{equation}
\label{eq:ch2.radiation-range-span}
\operatorname{ran}\mathcal{R}_\omega
=
\overline{\operatorname{span}\{\mathbf{u}_n:\ n\in\mathcal{N}_{\mathrm{rad}},\ \lambda_n\neq 0\}}.
\end{equation}
Eq.~\eqref{eq:ch2.radiation-map-modal-sum} follows by applying
Eq.~\eqref{eq:ch2.radiation-spectral-expansion} to
\(\mathcal{R}_\omega[\mathfrak{D}]\) and using
Eqs.~\eqref{eq:ch2.source-modal-amplitude}--\eqref{eq:ch2.radiation-eigenpair-def}.
Linearity of \(\mathcal{S}_\omega\) and \(\mathcal{R}_\omega\) extends the sum from finite
superpositions to all physically realizable impressed data by completion in \(\mathcal{U}\) on the
outgoing radiative subspace
\cite{stratton1941}. Eq.~\eqref{eq:ch2.radiation-operator-spectral-form} is obtained by
recognizing each term as \(P_n\mathcal{S}_\omega\) acting with gain \(\lambda_n\).
Eq.~\eqref{eq:ch2.radiation-range-span} follows because eigenvectors with
\(\lambda_n=0\) contribute no radiative export; non-radiating sources lie in
\(\ker\mathcal{R}_\omega\) Eq.~\eqref{eq:ch2.ker-R-omega}
\cite{harrington2001}. Resolution Eq.~\eqref{eq:ch2.spectral-resolution-composition} matches
\(\Pi_{\mathrm{rad}}=\sum_n P_n\) on \(\operatorname{ran}\mathcal{R}_\omega\)
\cite{stratton1941}. Resonant cavities, microwave filters, and waveguide resonators store energy in interior modes
that need not export to a radiating port: stored versus exported energy is the physical
distinction the cavity specialization encodes \cite{stratton1941}. On a bounded domain with passive \(\mathcal{B}_\Gamma\) and interior resonances,
Eq.~\eqref{eq:ch2.discrete-continuous-spectrum} yields a discrete accumulating set
\(\{\lambda_n\}_{n=1}^{\infty}\) when \(\mathcal{R}_\omega\) is restricted to cavity
outgoing-propagating subspaces compatible with Eq.~\eqref{eq:ch2.embedding-chain}
\cite{stratton1941}. The cavity specialization reads
\begin{equation}
\label{eq:ch2.cavity-discrete-radiation-spectrum}
\mathcal{R}_\omega^{(\mathrm{cav})}
=
\sum_{n=1}^{\infty}
\lambda_n\,P_n^{(\mathrm{cav})}\circ\mathcal{S}_\omega,
\qquad
\lambda_n\to 0\ \text{or}\ \lambda_n\to\infty
\ \text{according to mode class},
\end{equation}
with \(P_n^{(\mathrm{cav})}\) projectors onto normalized cavity radiation modes
\cite{stratton1941}. Equation~\eqref{eq:ch2.cavity-discrete-radiation-spectrum} is the
constructive content of Eq.~\eqref{eq:ch2.spectral-compression-operator} in closed volumes;
storage resonances with \(\lambda_n=0\) on \(\mathcal{R}_\omega\) remain in
\(\mathbf{u}_{\mathrm{nr}}\) Eq.~\eqref{eq:ch2.nonradiative-component-def}
\cite{harrington2001}. Using Eq.~\eqref{eq:ch2.cavity-discrete-radiation-spectrum} on open
exterior problems misstates continuous angular content
Eq.~\eqref{eq:ch2.discrete-continuous-spectrum} \cite{stratton1941}.
In free space the discrete modal sum becomes continuous: every outgoing direction on the sphere
$S^2$ carries spectral weight, and polarization index $\sigma\in\{1,2\}$ labels the two
transverse polarizations perpendicular to propagation direction $\hat{\mathbf{k}}$
\cite{stratton1941}. Physically, $\hat{\mathbf{k}}$ is the unit propagation
direction observed on a far-zone scanner; $\sigma$ selects which transverse polarization
component is sampled; $\mathbf{e}_\sigma$ is the corresponding unit polarization vector with
$\mathbf{e}_\sigma\cdot\hat{\mathbf{k}}=0$; and $d\Omega(\hat{\mathbf{k}})$ is the solid-angle
measure that weights how much radiated power each direction receives
\cite{hansen1988}. For compact sources in free space with outgoing radiation condition, angular eigenmodes are labeled by
propagation direction \(\hat{\mathbf{k}}\in S^{2}\) and polarization index \(\sigma\in\{1,2\}\)
transverse to \(\hat{\mathbf{k}}\) \cite{stratton1941}. A normalized plane-wave
radiation mode is
\begin{equation}
\label{eq:ch2.exterior-angular-eigenmode}
\mathbf{u}_{\hat{\mathbf{k}},\sigma}(\mathbf{r})
=
\big(\E_{\hat{\mathbf{k}},\sigma}(\mathbf{r}),\,\H_{\hat{\mathbf{k}},\sigma}(\mathbf{r})\big),
\qquad
\E_{\hat{\mathbf{k}},\sigma}\sim \mathbf{e}_\sigma\,e^{-\jj k_0\hat{\mathbf{k}}\cdot\mathbf{r}}/r,
\end{equation}
with \(\mathbf{e}_\sigma\cdot\hat{\mathbf{k}}=0\) and outgoing phase
\cite{stratton1941}. The continuous spectral measure \(d\Omega(\hat{\mathbf{k}})\)
replaces discrete summation in Eq.~\eqref{eq:ch2.radiation-spectral-expansion}:
\phantomsection\label{eq:ch2.discrete-continuous-spectrum}
\begin{equation}
\label{eq:ch2.exterior-continuous-spectral-integral}
\mathcal{R}_\omega[\mathfrak{D}](\mathbf{r})
=
\int_{S^{2}}
\sum_{\sigma=1}^{2}
a_\sigma(\hat{\mathbf{k}})\,
\mathbf{u}_{\hat{\mathbf{k}},\sigma}(\mathbf{r})\,d\Omega(\hat{\mathbf{k}}),
\end{equation}
with amplitudes \(a_\sigma\) determined by \(\mathcal{S}_\omega[\mathfrak{D}]\) through
Eq.~\eqref{eq:ch2.source-modal-amplitude} in the rigged limit \cite{stratton1941}. Rigged
Hilbert space completion and measure normalization are developed in \ref{ch:03}
Subsection~3.1.3 records the continuous angular target
Eq.~\eqref{eq:ch2.exterior-continuous-spectral-integral} compatible with far-zone reduction
\cite{harrington2001}.
\phantomsection\label{prop:ch2.far-pattern-from-modes}
Each radiation mode $\mathbf{u}_n$ carries a far-zone fingerprint
$\mathbf{f}_n(\theta,\phi)$---the angular pattern shape that mode $n$ contributes on $S^2$---
while $\E_0(\theta,\phi)$ is the measurable far-zone electric pattern a compact range records
after probe correction and transformation \cite{hansen1988}. Let \(\mathcal{D}_{J,\mathrm{reg}}=1\) and let \(\mathbf{u}_n\) be radiation eigenmodes with
far amplitudes \(\mathbf{f}_n(\theta,\phi)\) defined through
Eq.~\eqref{eq:ch2.far-field-amplitude-integral}. Then the far pattern of
\(\mathcal{R}_\omega[\mathfrak{D}]\) satisfies
\phantomsection\label{eq:ch2.far-pattern-modal}
\begin{equation}
\label{eq:ch2.spectral-far-pattern-expansion}
\E_0(\theta,\phi)
=
\sum_{n\in\mathcal{N}_{\mathrm{rad}}}
\lambda_n\,c_n[\mathfrak{D}]\,\mathbf{f}_n(\theta,\phi),
\end{equation}
with convergence in mean-square on \(S^{2}\) when the modal sum converges in
\(\mathcal{U}\) \cite{stratton1941}. Equation
\eqref{eq:ch2.spectral-far-pattern-expansion} is the upgrade
of Eq.~\eqref{eq:ch2.far-pattern-modal}; explicit \(\mathbf{f}_n\) in vector
harmonic form belong to \ref{ch:03} \cite{stratton1941}.
Apply the far-field approximation of source integrals (Subsection~\ref{sec:ch252}) to each modal component
\(\lambda_n c_n\,\mathbf{u}_n\) and superpose using linearity
Eq.~\eqref{eq:ch2.radiation-linearity} on outgoing radiative solutions. Coefficients match
Eq.~\eqref{eq:ch2.source-modal-amplitude} because far-zone reduction depends linearly on
\(\mathcal{R}_\omega[\mathfrak{D}]\) through
Eqs.~\eqref{eq:ch2.far-field-integral-reduction}--\eqref{eq:ch2.far-field-amplitude-integral}
\cite{stratton1941}. regime qualification
Eq.~\eqref{eq:ch2.regime-component-criterion} must hold before interpreting
Eq.~\eqref{eq:ch2.spectral-far-pattern-expansion} as transport data
\cite{harrington2001}. Source reconstruction is harder than forward radiation because inversion requires adjoint modes
$\widetilde{\mathbf{u}}_n$ that test whether a desired pattern $\mathbf{g}$ lies in the
reachable range of $\mathcal{R}_\omega$ \cite{stratton1941}. When \(\mathcal{R}_\omega\) is not self-adjoint on \(\mathcal{J}_\Omega\), define adjoint
modes \(\widetilde{\mathbf{u}}_n\) through the source pairing
Eq.~\eqref{eq:ch2.dual-source-pairing} such that
\begin{equation}
\label{eq:ch2.radiation-biorthogonality}
\langle \widetilde{\mathbf{u}}_m,\,\mathcal{R}_\omega[\Jimp^{(n)}]\rangle_{\mathcal{U}}
=
\lambda_n\,\delta_{mn},
\end{equation}
for normalized radiation eigenpairs \cite{stratton1941}. Equation
\eqref{eq:ch2.radiation-biorthogonality} is the field-side restatement of solvability relation
Eq.~\eqref{eq:ch2.spectral-solvability}: cokernel directions annihilated by
Eq.~\eqref{eq:ch2.radiation-biorthogonality} obstruct inversion
\(\mathcal{R}_\omega[\Jimp]=\mathbf{g}\) \cite{stratton1941}. Complete adjoint-mode catalogs
on exterior domains belong to Subsection~\ref{sec:ch263} \cite{harrington2001}.
Modes with \(\lambda_n=0\) satisfy \(P_n\mathcal{R}_\omega=\mathbf{0}\) but may have
\(P_n\mathcal{S}_\omega\neq\mathbf{0}\) when reactive content projects onto non-export
subspaces Eq.~\eqref{eq:ch2.nonpropagating-component-def}
\cite{stratton1941}. Evanescent and non-propagating sheets remain annihilated: spectral labels
cannot convert evanescent storage into qualified far-zone radiation---regime qualification and
export selection remain mandatory before transport claims \cite{stratton1941}.
\begin{equation}
\label{eq:ch2.modal-null-on-nonprop}
P_n\circ\Pi_{\mathrm{ev}}=P_n\circ\Pi_{\mathrm{non}}=\mathbf{0}
\quad\text{for all}\ n\in\mathcal{N}_{\mathrm{rad}},
\end{equation}
by Eq.~\eqref{eq:ch2.rad-null-on-evanescent} and modal definition
Eq.~\eqref{eq:ch2.modal-projector-def} \cite{stratton1941}. Spectral labels therefore do not
override regime qualification Eq.~\eqref{eq:ch2.regime-operator-screen}
\cite{harrington2001}.
The eigenbasis expansions and far-pattern coefficients in
Eq.~\eqref{eq:ch2.spectral-far-pattern-expansion} require linear media at fixed \(\omega\),
outgoing radiation condition, and completeness of \(\{\mathbf{u}_n\}\) appropriate to the
problem class \cite{stratton1941}. Loss shifts \(\lambda_n\) into the complex plane without
promoting evanescent storage to export through Eq.~\eqref{eq:ch2.modal-null-on-nonprop}
\cite{stratton1941}. Recoverable coefficients in
Eq.~\eqref{eq:ch2.spectral-far-pattern-expansion} are further capped by finite observation rank
Eq.~\eqref{eq:ch2.reconstruction-rank-cap} as established in Subsection~\ref{sec:ch26.signature-aperture}
\cite{harrington2001}; consult that subsection for $\sigma_n$ cliffs rather than restating
finite-rank doctrine here.
As a worked example, consider a rectangular cavity fed by an electrically small probe at $f=10\,\mathrm{GHz}$. The
physical configuration is familiar from filter and resonator laboratories: a closed metal box with
a coaxial penetration, interior standing waves visible in $|S_{11}|$ and in probe-coupled power
at ports \cite{stratton1941}. Interior modes produce a dense discrete spectrum in
$\mathcal{S}_\omega[\Jimp]$---many field patterns exist inside the volume---but only propagating
outgoing sectors contribute to $\mathcal{R}_\omega[\Jimp]$ on a measurement aperture or radiating
port \cite{stratton1941}. Truncating Eq.~\eqref{eq:ch2.cavity-discrete-radiation-spectrum} at
$N=20$ modes with error bounded by Eq.~\eqref{eq:ch2.compression-error-bound} yields a stable
interior power budget on a port wattmeter: the expected observation is a clean $|S_{11}|$ trace
with identifiable resonances and predictable coupling to the feed.
The lesson sharpens by contrast. Apply the same truncated cavity sum to an exterior half-wave
dipole in free space and sidelobe structure misaligns: the exterior problem requires the
continuous angular measure Eq.~\eqref{eq:ch2.exterior-continuous-spectral-integral}, not a cavity
eigenbasis \cite{jackson1999}. On a compact range the dipole presents a smooth far-zone pattern
cut; in the cavity it presents port physics entirely governed by enclosure modes. Conflating the
two domain classes is a common numerical error in aperture codes and in student projects that
reuse modal truncations without checking export selection \cite{harrington2001}. The significance
is operational: spectral truncation is domain-specific, export selection must precede modal
counting, and a dense interior spectrum is not the same object as a recoverable far-zone pattern
family \cite{stratton1941}.
As a worked example, a balanced loop pair with vanishing electric dipole moment is the cleanest laboratory reminder
that near-field strength does not imply far-zone export \cite{devaney2004}. The physical setup
places two loops with opposing dipole moments so that radiation from electric dipole terms cancels
on $\Gamma_R$ while reactive $\H$ and $\E$ remain large within one wavelength of the structure.
On a spherical near-field scan at $R=0.4\lambda_0$, probes record strong tangential fields; on
$\Gamma_R$ at $R=10\lambda_0$ the wattmeter reads negligible export. Coefficients $c_n$ on
exporting modes vanish ($\lambda_n=0$), yet $\Pi_{\mathrm{ev}}\mathcal{S}_\omega[\Jimp^{(\mathrm{NR})}]$
may be substantial \cite{stratton1941}.
The observable behavior confuses teams who report patterns from $\mathcal{S}_\omega[\Jimp]$ without
$\Pi_{\mathrm{rad}}$: the near field looks radiative locally while the far-zone signature is null.
The interpretation is spectral, not a failure of probes. Membership in $\ker\mathcal{R}_\omega$
is invisible to far-zone pattern meters yet measurable in the reactive sector. For inverse
synthesis, the consequence is equally concrete: two distinct feeds can match the same far pattern
while differing in reactive storage---a realizability question developed in \ref{ch:04}
\cite{stratton1941}. The example should be remembered as the archetype of export efficiency
$\lambda_n=0$ with non-trivial Maxwell storage.
\medskip
\noindent
Subsection~\ref{sec:ch261} fixed radiation eigenpairs, modal projectors, and the spectral
expansion of $\mathcal{R}_\omega$ on $\mathcal{F}_{\mathrm{rad}}$. Recoverable rank and
$\sigma_n$ cliffs remain anchored in Subsection~\ref{sec:ch26.signature-aperture}; the eigenbasis
organizes export directions that finite apertures can in principle resolve once regime
qualification holds \cite{stratton1941}. Rotational degeneracy and Fredholm
solvability develop next on the same operator.
\subsection{Rotational symmetry and spectral degeneracy}
\label{sec:ch262}
Rotational degeneracy is visible on every compact range before covariance formulas appear:
rotating a horn reorients principal cuts while total exported power stays fixed when probe and
source frames rotate together \cite{harrington2001}. The formal commutation relations
below encode that laboratory symmetry on the export eigenbasis from
Subsection~\ref{sec:ch261}.
% PHYSICAL-MEMORY:sec:ch262
% PHYSICAL-MEMORY:sec:ch262-b2
\begin{physicalvignette}[Outdoor elevated range with boresight tower]
A second bench observation matters on Rotating LPDA on a pattern tower. When technicians swap only the receive probe while holding the transmit chain fixed, reciprocity checks whether the reported coupling is property of the structure or artifact of a mismatched flux gauge \cite{stratton1941}. Failed reciprocity usually traces to a near-field tap, a radome standing-wave, or a probe polarization basis that was never aligned with the transmit standard \cite{hansen1988}.
Rayleigh and Stratton treated such ranges as boundary-value laboratories long before digital near-field transforms became routine \cite{stratton1941}. Harrington's method-of-moments culture reinforced the same discipline: metal outline first, harmonic headline second \cite{harrington2001}.
The measurement logic above is made precise in Subsection~\ref{sec:ch262} \cite{stratton1941}.
\end{physicalvignette}
\begin{physicalvignette}[Outdoor elevated range with boresight tower]
Consider Outdoor elevated range with boresight tower. Before any coefficient table is printed, technicians record frequency, cable loss, and the mechanical outline that fixes electrical size. The first walk-through of the range is always tactile: flange bolts, probe standoff, turntable boresight, and the wattmeter reference on the receive chain \cite{hansen1988}.
On this hardware the near zone is visually impressive---standing reactive structure, rapid phase rotation, sometimes a visible VSWR ripple at the feed---yet the export question is measured only after outgoing selection on a large sphere \cite{harrington2001}. What the camera sees in the chamber is not automatically what the far-zone pattern claims \cite{stratton1941}.
Engineers therefore learn to separate three readings: what the feed can synthesize on declared support, what propagating content survives on the radiative range, and what a finite instrument can recover at declared tolerance \cite{devaney2004}. Confusing those readings produces the rank gaps and watt failures that later sections quantify \cite{pal2026}.
Rayleigh and Stratton treated such ranges as boundary-value laboratories long before digital near-field transforms became routine \cite{stratton1941}. Harrington's method-of-moments culture reinforced the same discipline: metal outline first, harmonic headline second \cite{harrington2001}.
Subsection~\ref{sec:ch262} states the formal export condition that this range scene presupposes \cite{stratton1941}.
\end{physicalvignette}
Let \(R\in\mathrm{SO}(3)\), field rotation \(\mathcal{D}(R)\) on qualified Maxwell pairs
\(\mathbf{u}=(\E,\H)\in\mathcal{F}_{\mathrm{SF}}\), and radiation map
Eq.~\eqref{eq:ch2.spectral-operator-targets} on outgoing data
\cite{stratton1941}. Isotropic vacuum or statistically isotropic
environments are assumed unless stated otherwise: constitutive tensors commute with
\(R\) in the sense of Section~\ref{sec:ch01.3.2} \cite{stratton1941}. Impressed sources
transform as polar vectors under active rotation,
\(\big(\mathcal{A}_R\Jimp\big)(\mathbf{r})=R\,\Jimp(R^{-1}\mathbf{r})\), consistent with
Eq.~\eqref{eq:ch1.vector-rotation-components} \cite{stratton1941}. regime qualification
Eq.~\eqref{eq:ch2.regime-operator-screen} and transport qualification
Eq.~\eqref{eq:ch2.transport-qualified-radiation} remain mandatory before interpreting
rotated reports as export data \cite{harrington2001}.
Rotation in the laboratory. Rotating a horn on a compact range is the everyday experiment that motivates covariance
relations. Mount the horn on a rotary stage, record $\E_t$ on $S_r$ at several azimuths, and
superpose the cuts: total exported power remains unchanged when probe and source frames rotate
together, but the principal beam walks across the sphere as expected for a steerable feed
\cite{harrington2001}. Rotating a dish feed scans a different arc on the reflector surface; modal
coefficients mix according to the same rules that rotate a measured $\E_t$ distribution on the
rim \cite{collin2001}. Spinning a corrugated aperture distribution about its axis redistributes
directional content within fixed $\ell$ blocks without creating new export unless the source
itself changes \cite{stratton1941}. Laboratory reports must distinguish active rotation of
hardware from passive relabeling of a fixed field in a rotated coordinate chart; confusing the
two conventions splits degeneracies artificially in processed $a_n(\hat{\mathbf{k}})$ data and
can make a symmetric structure appear artificially asymmetric in synthesis charting
\cite{harrington2001}. The covariance identity below is the formal statement of what range
technicians already enforce by co-rotating probes and references.
Field rotation means physically reorienting the $\E$ and $\H$ vectors with the hardware: what
changes is the direction of the pattern on $S_r$; what remains unchanged, when frames are
handled consistently, is total exported power and the eigenvalue ordering within degenerate
blocks \cite{harrington2001}. Define the Maxwell field representation
\begin{equation}
\label{eq:ch2.field-rotation-representation}
\mathcal{D}(R)\,\mathbf{u}
=
\big(\mathcal{A}_R\E,\,\mathcal{A}_R\H\big),
\qquad
\mathbf{u}=(\E,\H)\in\mathcal{F}_\Omega,
\end{equation}
using active rotation Eq.~\eqref{eq:ch1.active-rotation-action} on each component
\cite{stratton1941}. Equation~\eqref{eq:ch2.field-rotation-representation} specializes
Section~\ref{sec:ch01.3.1} to phasor Maxwell
pairs: passive relabeling Eq.~\eqref{eq:ch1.passive-rotation-action} produces the same
physics with rotated components related by
Eq.~\eqref{eq:ch1.active-passive-duality} \cite{stratton1941}. Group composition
Eq.~\eqref{eq:ch1.rotation-group-action} yields
\(\mathcal{D}(R_1R_2)=\mathcal{D}(R_1)\mathcal{D}(R_2)\) on qualified field pairs
\cite{stratton1941}.
Covariance of the radiation map. This equation is perhaps the most important symmetry statement in antenna engineering: active
rotation of the feed in a co-rotated measurement frame leaves physical export content unchanged
while reorienting modal coefficients through the representation $\mathcal{D}(R)$
\cite{harrington2001}. When sources and observation frames are co-rotated, the radiation chain must commute with
rotation up to the field representation $\mathcal{D}(R)$: rotating the feed and evaluating on a
co-rotated far-zone frame must yield the same physics as evaluating first and then rotating the
pattern \cite{stratton1941}. That requirement is visible on every compact range
where probe calibration tracks mechanical azimuth. Formally,
\begin{equation}
\label{eq:ch2.radiation-rotation-commutation}
\mathcal{R}_\omega\big[\mathcal{A}_R\Jimp\big]
=
\mathcal{D}(R)\,\mathcal{R}_\omega[\Jimp],
\end{equation}
on physically realizable \(\Jimp\in\mathcal{J}_\Omega\) with outgoing radiation condition and isotropic
constitutive response \cite{stratton1941}. Equation
\eqref{eq:ch2.radiation-rotation-commutation} develops Eq.
Eq.~\eqref{eq:ch2.degeneracy-covariance} and aligns with covariance test
Eq.~\eqref{eq:ch2.covariance-operator-test} for \(\mathbf{a}=\mathbf{0}\)
\cite{stratton1941}. Residual \(\mathcal{R}_{\mathcal{F}}\) from
Eq.~\eqref{eq:ch1.covariance-residual} diagnoses broken commutation in discretizations
\cite{harrington2001}. Rotating
\(\Jimp\) while evaluating flux on an unrotated \(\mathcal{S}_\omega[\Jimp]\) violates
origin and covariance criteria from Section~\ref{sec:ch01.3.3} even when
Eq.~\eqref{eq:ch2.radiation-rotation-commutation} holds for power alone
\cite{stratton1941}.
\phantomsection\label{fig:ch2.radiation-rotation-commutation}
\begin{figure}[t]
\centering
\ChOneFigBlock{%
\begin{tikzpicture}[ch1fig, node distance=5mm and 7mm]
 \node[ch1box] (j) {$\Jimp$};
 \node[ch1box, right=7mm of j] (r) {$\mathcal{R}_\omega$};
 \node[ch1box, right=7mm of r] (u) {$\mathbf{u}_{\mathrm{rad}}$};
 \node[ch1box, below=5mm of j] (ar) {$\mathcal{A}_R$};
 \node[ch1box, below=5mm of u] (dr) {$\mathcal{D}(R)$};
 \draw[ch1flow] (j) -- (r) -- (u);
 \draw[ch1flow] (ar.south) -- (ar.south |- j.north);
 \draw[ch1flow] (ar.south |- j.north) -- (j.north);
 \draw[ch1flow] (dr.south) -- (dr.south |- u.north);
 \draw[ch1flow] (dr.south |- u.north) -- (u.north);
 \draw[ch1flow] (ar) -| (r);
\end{tikzpicture}%
}
\caption{Rotational covariance Eq.~\eqref{eq:ch2.radiation-rotation-commutation}: active
source rotation commutes with \(\mathcal{R}_\omega\) up to field representation
\(\mathcal{D}(R)\) Eq.~\eqref{eq:ch2.field-rotation-representation}.}
\label{fig:ch2.rotation-spectral-covariance}
\end{figure}
~\eqref{eq:ch2.radiation-rotation-commutation} in the same layout as
Figure~\ref{fig:ch1.wp-active-passive-rotation}: geometry and components must be rotated
consistently before modal coefficients are compared \cite{stratton1941}. Under rotation, modal coefficients mix through the matrix $D^{(\mathrm{rad})}_{nm}(R)$: an
engineer sees principal cuts walk across the sphere while total power on $\Gamma_R$ stays fixed
when probe and source frames rotate together \cite{harrington2001}. Let \(\{\mathbf{u}_n\}_{n\in\mathcal{N}_{\mathrm{rad}}}\) be the orthonormal radiation
basis developed above and \(c_n[\mathfrak{D}]\) source amplitudes
Eq.~\eqref{eq:ch2.source-modal-amplitude}. Under active rotation of impressed data,
\begin{equation}
\label{eq:ch2.modal-coefficient-rotation}
c_n[\mathcal{A}_R\mathfrak{D}]
=
\sum_{m\in\mathcal{N}_{\mathrm{rad}}}
D^{(\mathrm{rad})}_{nm}(R)\,c_m[\mathfrak{D}],
\end{equation}
whenever Eq.~\eqref{eq:ch2.radiation-rotation-commutation} holds and the family of radiating modes is
closed under \(\mathcal{D}(R)\),
\begin{equation}
\label{eq:ch2.modal-basis-covariance}
\mathcal{D}(R)\,\mathbf{u}_m
=
\sum_{n\in\mathcal{N}_{\mathrm{rad}}}
D^{(\mathrm{rad})}_{nm}(R)\,\mathbf{u}_n.
\end{equation}
Equations~\eqref{eq:ch2.modal-coefficient-rotation}--\eqref{eq:ch2.modal-basis-covariance}
are the spectral restatement of vector covariance
Eq.~\eqref{eq:ch1.vector-observable-covariance} on \(\mathcal{F}_{\mathrm{rad}}\)
\cite{stratton1941}. Matrix elements \(D^{(\mathrm{rad})}_{nm}(R)\) are not tabulated
here; Chapter~3 (\emph{Angular Momentum Operators and SO(3)}) supplies them on vector spherical harmonic bases
\cite{harrington2001}. Eigenvalues \(\lambda_n\) in
Eq.~\eqref{eq:ch2.radiation-eigenpair-def} are unchanged within a degenerate block because
\(\mathcal{R}_\omega\) commutes with \(\mathcal{D}(R)\) on that block
\cite{stratton1941}.
\phantomsection\label{thm:ch2.rotation-spectral-degeneracy}
Repeated eigenvalues occur because rotational symmetry makes distinct angular orientations
physically equivalent within an irreducible block: full $\mathrm{SO}(3)$ symmetry produces
$(2\ell+1)$-fold degeneracy, and symmetry breaking---anisotropic materials, finite facets,
or broken ground-plane parity---splits those blocks into distinct $\lambda_n$
\cite{stratton1941}. Suppose Eq.~\eqref{eq:ch2.radiation-rotation-commutation} holds for all
\(R\in\mathrm{SO}(3)\) in a compact subgroup \(G\subseteq\mathrm{SO}(3)\) leaving the
problem geometry invariant. Then \(\mathcal{R}_\omega\) and \(\mathcal{D}(R)\) generate a
commuting family on \(\mathcal{F}_{\mathrm{rad}}\): for each \(R\in G\),
\begin{equation}
\label{eq:ch2.radiation-D-commutation}
\mathcal{D}(R)\,\mathcal{R}_\omega
=
\mathcal{R}_\omega\,\mathcal{D}(R)
\quad\text{on physically realizable sources,}
\end{equation}
and eigenspaces of \(\mathcal{R}_\omega\mathcal{S}_\omega\) decompose into irreducible
\(G\)-representations. If an irreducible block carries angular index \(\ell\), its
degeneracy multiplicity satisfies
\begin{equation}
\label{eq:ch2.spectral-degeneracy-multiplicity}
\dim\mathcal{H}_{\ell}
=
2\ell+1
\quad\text{for full \(\mathrm{SO}(3)\) completion on }S^{2},
\end{equation}
with \(\mathcal{H}_{\ell}\) the \(\ell\)-indexed radiative subspace
\cite{stratton1941}. Equation
\eqref{eq:ch2.spectral-degeneracy-multiplicity} explains repeated eigenvalues in
Eq.~\eqref{eq:ch2.radiation-spectral-expansion} without duplicating mode physics.
Eq.~\eqref{eq:ch2.radiation-D-commutation} follows by applying
Eq.~\eqref{eq:ch2.radiation-rotation-commutation} after
Eq.~\eqref{eq:ch2.radiation-operator-spectral-form} and using linearity of
\(\mathcal{S}_\omega\) \cite{stratton1941}. Standard spectral theory for commuting
unitary or orthogonal representations implies simultaneous block diagonalization of
\(\mathcal{R}_\omega\) and \(\mathcal{D}(R)\) on invariant subspaces
\cite{stratton1941}. For full rotational symmetry on angular patterns,
irreducible representations of \(\mathrm{SO}(3)\) on \(L^{2}(S^{2})\) have dimension
\(2\ell+1\); Eq.~\eqref{eq:ch2.spectral-degeneracy-multiplicity} records that count on
\(\mathcal{F}_{\mathrm{rad}}\) once outgoing transversality and
Eq.~\eqref{eq:ch2.rad-null-on-evanescent} are imposed \cite{jackson1999}. Detailed
construction of \(\mathcal{H}_\ell\) is developed in \ref{ch:03} \cite{harrington2001}.
Total radiated power $P_{\mathrm{rad}}$ is the scalar quantity a wattmeter or sphere integral
reports on $\Gamma_R$; it cannot change under rigid rotation when frames are co-rotated
\cite{poynting1884}. Modal coefficients $c_n$ redistribute directional content,
while export eigenvalues $\lambda_n$ set how efficiently each mode reaches the meter
\cite{harrington2001}. Rigid rotation cannot change total exported power when meters act on
\(\mathcal{R}_\omega[\Jimp]\) with co-rotated frames. Using
Eq.~\eqref{eq:ch2.power-flux-functional} and
Eq.~\eqref{eq:ch1.rotation-power-invariance},
\begin{equation}
\label{eq:ch2.spectral-power-rotation-invariance}
P_{\mathrm{rad}}[\mathcal{A}_R\mathfrak{D}]
=
P_{\mathrm{rad}}[\mathfrak{D}],
\qquad
P_{\mathrm{rad}}[\mathfrak{D}]
=
\sum_{n\in\mathcal{N}_{\mathrm{rad}}}
|\lambda_n|^{2}\,|c_n[\mathfrak{D}]|^{2}
\quad\text{when modes are orthogonal and normalized,}
\end{equation}
for lossless outgoing radiation \cite{stratton1941}. Equation
\eqref{eq:ch2.spectral-power-rotation-invariance} is the modal restatement of scalar
invariance Eq.~\eqref{eq:ch1.scalar-observable-invariance} on integrated power: rotated
coefficients redistribute directional content but preserve the scalar budget when
\(\mathcal{D}_{J,\mathrm{reg}}=1\) \cite{stratton1941}. Net angular-flux magnitude obeys
Eq.~\eqref{eq:ch1.angular-momentum-balance} on the same qualified subspace
Active rotation reorients sources and patterns in space; passive rotation re-expresses
fixed physics in a rotated chart Eqs.~\eqref{eq:ch1.active-rotation-action}--\eqref{eq:ch1.passive-rotation-action}
\cite{stratton1941}. Modal coefficients transform as
Eq.~\eqref{eq:ch2.modal-coefficient-rotation} under active action, while passive
relabeling conjugates the same \(D^{(\mathrm{rad})}_{nm}(R)\) through
Eq.~\eqref{eq:ch1.active-passive-duality} without altering
Eq.~\eqref{eq:ch2.spectral-power-rotation-invariance} \cite{stratton1941}.
Confusing the two conventions produces spurious degeneracy splitting in measured
\(a_n(\hat{\mathbf{k}})\) from Eq.~\eqref{eq:ch2.spectral-far-pattern-expansion} even when
rotational symmetry of the spectrum applies \cite{stratton1941}. Laboratory
rotation of an antenna rotates measured coefficients through
Eq.~\eqref{eq:ch2.modal-coefficient-rotation}; exported power on
\(\mathcal{R}_\omega[\Jimp]\) remains invariant in the far zone
\cite{harrington2001}. Origin translation \(\mathbf{a}\) shifts torque and angular-flux reports through
Section~\ref{sec:ch01.3.3} while leaving far-zone radiative power invariant when
\(\mathcal{D}_{J,\mathrm{reg}}=1\) \cite{stratton1941}. At the spectral level, origin
change mixes modal coefficients within fixed \(\ell\) blocks through addition of
lower-\(\ell\) content from shifted multipole expansions; it does not create new
\(\lambda_n\neq 0\) export unless sources change \cite{jackson1999}.
The rotation analysis above therefore separates \emph{rotational} degeneracy from
\emph{translational} bookkeeping in Section~\ref{sec:ch01.3.3}; intrinsic angular observables use
Eq.~\eqref{eq:ch1.intrinsic-angular-observable} by reference \cite{stratton1941}.
\phantomsection\label{prop:ch2.degeneracy-splitting-broken-symmetry}
Let \(G_0\subset\mathrm{SO}(3)\) be the exact symmetry of a source--boundary configuration
and let \(\mathcal{H}_{\ell}\) be an irreducible block from rotational degeneracy of the
spectrum. Any symmetry-breaking perturbation
(reduced \(G_0\), anisotropic material, or finite facetization) splits
Eq.~\eqref{eq:ch2.spectral-degeneracy-multiplicity} into smaller blocks with distinct
\(\lambda_n\) within the same \(\ell\) sector \cite{stratton1941}. The
splitting order is controlled by representation coupling rules imported from
Section~\ref{sec:ch01.3.2} on tensor perturbations; explicit Clebsch--Gordan content belongs
to \ref{ch:03} Subsection~3.2.5 \cite{stratton1941}.
When \(G_0\) is proper, \(\mathcal{D}(R)\) block-diagonalizes with fixed \(\lambda_n\) on
each irreducible piece. Perturbations that fail to commute with all \(R\in G_0\) remove
simultaneous eigenstructure; standard degenerate perturbation theory applies to the
restricted radiation operator on \(\mathcal{H}_\ell\) \cite{stratton1941}. Rotational degeneracy of the spectrum requires isotropic or symmetry-matched
media, outgoing radiation condition, and radiating-mode families closed under \(\mathcal{D}(R)\)
\cite{stratton1941}. Anisotropic crystals reduce $G$ and split
Eq.~\eqref{eq:ch2.spectral-degeneracy-multiplicity} when symmetry is broken
\cite{stratton1941}. Observation rank Eq.~\eqref{eq:ch2.reconstruction-rank-cap}
resolves only a finite subset of each degenerate block even when continuous symmetry is exact
\cite{harrington2001}. Regime qualification
Eq.~\eqref{eq:ch2.regime-component-criterion} still governs whether rotated near fields may be
reported as transport \cite{stratton1941}.
\subsection{Fredholm structure and pattern solvability}
\label{sec:ch263}
Pattern synthesis and source inversion are the inverse problems every antenna laboratory
eventually confronts. One specifies a desired far-zone amplitude distribution and asks which
feed currents will produce it; or one measures near fields on a scanner and asks which sources
could have generated them \cite{harrington2001}. When these inversions fail, the
failure is rarely mysterious in hindsight: non-radiating components lie in
$\ker\mathcal{R}_\omega$---the set of admissible impressed currents annihilated by export---and
leave no far-zone signature, so no finite measurement set can resolve them without additional
excitation or constraints \cite{devaney2004}. Recoverable rank from
Subsection~\ref{sec:ch26.signature-aperture} caps the truncated modal system even when a
continuous target is mathematically consistent \cite{hansen1988}. Fredholm theory names
kernel, cokernel, and index after the eigenbasis and rotational structure above are in place
\cite{stratton1941}.
% PHYSICAL-MEMORY:sec:ch263
% PHYSICAL-MEMORY:sec:ch263-b2
\begin{physicalvignette}[Biconical antenna for wideband chamber calibration]
A second bench observation matters on Standard-gain horn trio for gain transfer. When technicians swap only the receive probe while holding the transmit chain fixed, reciprocity checks whether the reported coupling is property of the structure or artifact of a mismatched flux gauge \cite{stratton1941}. Failed reciprocity usually traces to a near-field tap, a radome standing-wave, or a probe polarization basis that was never aligned with the transmit standard \cite{hansen1988}.
Rayleigh and Stratton treated such ranges as boundary-value laboratories long before digital near-field transforms became routine \cite{stratton1941}. Harrington's method-of-moments culture reinforced the same discipline: metal outline first, harmonic headline second \cite{harrington2001}.
The measurement logic above is made precise in Subsection~\ref{sec:ch263} \cite{stratton1941}.
\end{physicalvignette}
\begin{physicalvignette}[Biconical antenna for wideband chamber calibration]
Consider Biconical antenna for wideband chamber calibration. Before any coefficient table is printed, technicians record frequency, cable loss, and the mechanical outline that fixes electrical size. The first walk-through of the range is always tactile: flange bolts, probe standoff, turntable boresight, and the wattmeter reference on the receive chain \cite{hansen1988}.
On this hardware the near zone is visually impressive---standing reactive structure, rapid phase rotation, sometimes a visible VSWR ripple at the feed---yet the export question is measured only after outgoing selection on a large sphere \cite{harrington2001}. What the camera sees in the chamber is not automatically what the far-zone pattern claims \cite{stratton1941}.
Engineers therefore learn to separate three readings: what the feed can synthesize on declared support, what propagating content survives on the radiative range, and what a finite instrument can recover at declared tolerance \cite{devaney2004}. Confusing those readings produces the rank gaps and watt failures that later sections quantify \cite{pal2026}.
Rayleigh and Stratton treated such ranges as boundary-value laboratories long before digital near-field transforms became routine \cite{stratton1941}. Harrington's method-of-moments culture reinforced the same discipline: metal outline first, harmonic headline second \cite{harrington2001}.
Subsection~\ref{sec:ch263} states the formal export condition that this range scene presupposes \cite{stratton1941}.
\end{physicalvignette}
Assumptions. Fix \(\omega>0\), bounded or exterior problem class with outgoing radiation condition on open problems,
physically realizable \(\mathcal{J}_\Omega\) Eq.~\eqref{eq:ch2.source-space-refined}, export map
Eq.~\eqref{eq:ch2.spectral-operator-targets}, and the radiation eigenbasis developed above
\cite{stratton1941}. Source pairing uses dual space
Eq.~\eqref{eq:ch2.dual-source-pairing}; field pairing uses
Eq.~\eqref{eq:ch2.maxwell-adjoint-pairing} \cite{stratton1941}. Compact perturbations are
taken in operator norm on \(\mathcal{J}_\Omega\) for sources with compact support
\(\Omega_s\) \cite{harrington2001}. Regime qualification Eq.~\eqref{eq:ch2.regime-operator-screen}
governs whether recovered \(\mathbf{g}\) may be reported as transport data
\cite{stratton1941}.
The \emph{radiation kernel} $\ker\mathcal{R}_\omega$ collects impressed currents invisible to
qualified far-zone meters. The \emph{cokernel}
$\mathrm{coker}\,\mathcal{R}_\omega=\mathcal{F}_{\mathrm{rad}}/\operatorname{ran}\mathcal{R}_\omega$
labels far-zone pattern targets that no admissible source can realize---equivalently, directions
orthogonal to every radiative image under the adjoint pairing \cite{stratton1941}. The
\emph{Fredholm index}
$\mathrm{ind}\,\mathcal{R}_\omega=\dim\ker\mathcal{R}_\omega-\dim\mathrm{coker}\,\mathcal{R}_\omega$
counts the net imbalance between invisible source degrees and unreachable pattern obstructions;
it is stable under the compact perturbations recorded below \cite{harrington2001}. In the
laboratory, kernel membership explains why two distinct feeds can match the same far pattern;
cokernel membership explains why a synthesis target refuses to converge no matter how many ports
are tuned \cite{stratton1941}.
Source inversion is difficult because radiation smooths fine source detail: the far-zone map
retains only exportable content and discards non-radiating degrees in $\ker\mathcal{R}_\omega$
\cite{devaney2004}. The export map inherits the Green-function split of Section~\ref{sec:ch23}: a leading outgoing
propagation part $\mathcal{R}_{0,\omega}$ plus a compact smoothing part $K_{\mathrm{rad}}$ at
finite separation on bounded source classes \cite{stratton1941}. Physically,
$\mathcal{R}_{0,\omega}$ is the dominant far-zone propagation law and $K_{\mathrm{rad}}$ collects
the compact corrections that smooth fine source structure into resolvable pattern content
\cite{harrington2001}. That split is the operator-side restatement of the range observation that
far-zone data smooth fine source detail. Formally,
\begin{equation}
\label{eq:ch2.radiation-fredholm-decomposition}
\mathcal{R}_\omega
=
\mathcal{R}_{0,\omega}+K_{\mathrm{rad}},
\qquad
K_{\mathrm{rad}}\ \text{compact on}\ \mathcal{J}_\Omega,
\end{equation}
where \(\mathcal{R}_{0,\omega}\) captures leading outgoing singular propagation and
\(K_{\mathrm{rad}}\) collects compact smoothing at finite separation on bounded source classes
\cite{stratton1941}. Equation~\eqref{eq:ch2.radiation-fredholm-decomposition}
specializes Eq.~\eqref{eq:ch2.fredholm-split} to the radiation image
Eq.~\eqref{eq:ch2.integral-radiation-map} after Sections~\ref{sec:ch23}--\ref{sec:ch25}
\cite{stratton1941}. On exterior domains with compact
\(\Omega_s\), \(K_{\mathrm{rad}}\) is compact while \(\mathcal{R}_{0,\omega}\) carries the
non-compact far-zone propagation; on bounded cavities, both parts reduce to discrete modal
books Eq.~\eqref{eq:ch2.cavity-discrete-radiation-spectrum} \cite{harrington2001}.
The Fredholm index counts the net imbalance between invisible source degrees
($\ker\mathcal{R}_\omega$) and unreachable pattern targets
($\mathrm{coker}\,\mathcal{R}_\omega$). On the range, kernel membership means different
feeds may produce identical far patterns; cokernel membership means a synthesis target lies
outside the reachable subspace no matter how many ports are tuned \cite{stratton1941}.
For closed operators with finite-dimensional kernel and cokernel, define the Fredholm index
\begin{equation}
\label{eq:ch2.radiation-fredholm-index-def}
\mathrm{ind}\,\mathcal{R}_\omega
=
\dim\ker\mathcal{R}_\omega
-
\dim\mathrm{coker}\,\mathcal{R}_\omega,
\end{equation}
with \(\mathrm{coker}\,\mathcal{R}_\omega=\mathcal{F}_{\mathrm{rad}}/\operatorname{ran}\mathcal{R}_\omega\)
on the declared quotient by outgoing radiation condition \cite{stratton1941}. Equation
\eqref{eq:ch2.radiation-fredholm-index-def} counts independent non-radiating source degrees
Eq.~\eqref{eq:ch2.ker-R-omega} minus obstructions to reaching a target far pattern
\(\mathbf{g}\in\mathcal{F}_{\mathrm{rad}}\) \cite{stratton1941}. Index is stable under compact
perturbations in Eq.~\eqref{eq:ch2.radiation-fredholm-decomposition}
\cite{harrington2001}. Observation rank cap Eq.~\eqref{eq:ch2.reconstruction-rank-cap}
is independent of \(\mathrm{ind}\,\mathcal{R}_\omega\) but composes with it in inversion
budgets \cite{stratton1941}.
\begin{figure}[t]
\centering
\ChOneFigBlock{%
\begin{tikzpicture}[ch1fig, node distance=5mm and 7mm]
 \node[ch1box] (j) {$\mathcal{J}_\Omega$};
 \node[ch1box, right=7mm of j] (r) {$\mathcal{R}_\omega$};
 \node[ch1box, right=7mm of r] (f) {$\mathcal{F}_{\mathrm{rad}}$};
 \node[ch1box, below=5mm of r] (k) {$K_{\mathrm{rad}}$};
 \node[ch1box, below=5mm of f] (c) {$\mathrm{coker}$};
 \draw[ch1flow] (j) -- (r) -- (f);
 \draw[ch1flow] (k.south) -- (k.south |- r.north);
 \draw[ch1flow] (k.south |- r.north) -- (r.north);
 \draw[ch1flow] (c.south) -- (c.south |- f.north);
 \draw[ch1flow] (c.south |- f.north) -- (f.north);
\end{tikzpicture}%
}
\caption{Fredholm measurement sequence: compact perturbation
Eq.~\eqref{eq:ch2.radiation-fredholm-decomposition}, kernel
\(\ker\mathcal{R}_\omega\), and cokernel obstructions to solvability
Eq.~\eqref{eq:ch2.radiation-solvability-criterion}.}
\label{fig:ch2.fredholm-radiation-measurement chain}
\end{figure}
~\eqref{eq:ch2.rank-composition-law} \cite{stratton1941}.
The Fredholm alternative states that a target far-zone field is synthesizable only if it
carries no component along cokernel directions invisible to the adjoint export map
\cite{stratton1941}. Consequence: a measured pattern residual that is orthogonal to every
reachable image is not a calibration fault---it signals membership in
$\mathrm{coker}\,\mathcal{R}_\omega$ or in a geometrically forbidden sector from
Subsection~\ref{sec:ch264} \cite{harrington2001}. The forward radiation equation $\mathcal{R}_\omega[\Jimp]=\mathbf{g}$
is solvable iff the target $\mathbf{g}\in\mathcal{F}_{\mathrm{rad}}$ is orthogonal to the
cokernel of the adjoint map on declared pairings:
\begin{equation}
\label{eq:ch2.radiation-solvability-criterion}
\mathcal{R}_\omega[\Jimp]=\mathbf{g}
\ \text{solvable}
\iff
\langle \mathbf{g},\,\mathbf{w}\rangle_{\mathcal{U}}=0
\ \ \forall\ \mathbf{w}\in\ker\mathcal{R}_\omega^{\dagger},
\end{equation}
with \(\mathcal{R}_\omega^{\dagger}\) the adjoint on
\((\mathcal{J}_\Omega,\mathcal{F}_{\mathrm{rad}})\) induced by
Eqs.~\eqref{eq:ch2.dual-source-pairing} and~\eqref{eq:ch2.maxwell-adjoint-pairing}
\cite{stratton1941}. Equation~\eqref{eq:ch2.radiation-solvability-criterion}
develops Eq.~\eqref{eq:ch2.spectral-solvability} and aligns with biorthogonality
Eq.~\eqref{eq:ch2.radiation-biorthogonality} on modal adjoints
\cite{stratton1941}. When \(\mathbf{g}\) arises from far-pattern data, solvability is
further capped by Eq.~\eqref{eq:ch2.reconstruction-rank-cap} even if
Eq.~\eqref{eq:ch2.radiation-solvability-criterion} holds in the continuous model
\cite{harrington2001}.
\phantomsection\label{thm:ch2.radiation-fredholm-structure}
Let \(\mathcal{R}_\omega\) satisfy Eq.~\eqref{eq:ch2.radiation-fredholm-decomposition} on
\(\mathcal{J}_\Omega\) with \(\mathcal{R}_{0,\omega}\) Fredholm of index zero on the outgoing
quotient and \(K_{\mathrm{rad}}\) compact. Then \(\mathcal{R}_\omega\) is Fredholm,
\begin{equation}
\label{eq:ch2.radiation-index-stability}
\mathrm{ind}\,\mathcal{R}_\omega
=
\mathrm{ind}\,\mathcal{R}_{0,\omega},
\end{equation}
and Eq.~\eqref{eq:ch2.radiation-solvability-criterion} holds. For finite-rank observation
\(\mathcal{O}_{\mathrm{meas}}:\mathcal{F}_{\mathrm{rad}}\to\mathbb{C}^{N}\),
\begin{equation}
\label{eq:ch2.observation-fredholm-composition}
\mathcal{O}_{\mathrm{meas}}\mathcal{R}_\omega
\ \text{is Fredholm on}\ \mathcal{J}_\Omega
\ \Longleftrightarrow\
\mathcal{R}_\omega\ \text{Fredholm and}\ 
\mathrm{rank}(\mathcal{O}_{\mathrm{meas}})=N,
\end{equation}
with rank bound Eq.~\eqref{eq:ch2.rank-composition-law}
Compact perturbation stability of Fredholm property and index follows from standard operator
theory applied to Eq.~\eqref{eq:ch2.radiation-fredholm-decomposition}; Subsection
~\ref{sec:ch213} recorded the embedding inputs without computing index here
\cite{stratton1941}. Solvability criterion
Eq.~\eqref{eq:ch2.radiation-solvability-criterion} is the Fredholm alternative for
\(\mathcal{R}_\omega\). Composition with finite-rank \(\mathcal{O}_{\mathrm{meas}}\) preserves
Fredholm class and adjusts index by at most \(N\) on the range side, yielding
Eq.~\eqref{eq:ch2.observation-fredholm-composition} and
Eq.~\eqref{eq:ch2.rank-composition-law} \cite{stratton1941}.
The modal fredholm reduction proceeds as follows: on the orthonormal radiation eigenbasis developed above,
\begin{equation}
\label{eq:ch2.modal-fredholm-reduction}
\mathcal{R}_\omega
=
\sum_{n\in\mathcal{N}_{\mathrm{rad}}}
\lambda_n\,P_n\circ\mathcal{S}_\omega,
\qquad
\lambda_n\neq 0
\ \Longleftrightarrow\
\text{mode }n\ \text{exporting},
\end{equation}
so solvability of \(\mathcal{R}_\omega[\Jimp]=\mathbf{g}\) reduces to matching coefficients
\(c_n[\mathfrak{D}]\) Eq.~\eqref{eq:ch2.source-modal-amplitude} on modes with
\(\lambda_n\neq 0\) \cite{stratton1941}. Modes with \(\lambda_n=0\) lie in
\(\ker\mathcal{R}_\omega\) at the source level Eq.~\eqref{eq:ch2.ker-R-omega}; modes
with \(\mathbf{g}\) unsupported on \(\operatorname{ran}\mathcal{R}_\omega\) generate cokernel
obstructions in Eq.~\eqref{eq:ch2.radiation-solvability-criterion}
\cite{stratton1941}. Truncation to \(\mathcal{R}_\omega^{(N)}\) from
Eq.~\eqref{eq:ch2.spectral-compression-operator} yields approximate solvability with error
Eq.~\eqref{eq:ch2.compression-error-bound} controlled by \(\sigma_{N+1}(K_{\mathrm{rad}})\)
\cite{harrington2001}. \phantomsection\label{prop:ch2.modal-invertibility-equivalence}
Let \(\mathbf{g}\in\operatorname{ran}\mathcal{R}_\omega\) and expand
\(\mathbf{g}=\sum_n g_n\,\mathbf{u}_n\) with \(\{\mathbf{u}_n\}\) from
Subsection~\ref{sec:ch261}. Then Eq.~\eqref{eq:ch2.radiation-solvability-criterion} holds if
and only if \(g_n=0\) whenever \(\lambda_n=0\). A particular solution is determined by
\begin{equation}
\label{eq:ch2.modal-particular-solution}
c_n
=
\frac{g_n}{\lambda_n}
\quad\text{for}\ \lambda_n\neq 0,
\end{equation}
with arbitrary coefficients on \(\ker\mathcal{R}_\omega\) not fixed by far-zone data
\cite{stratton1941}. Equation~\eqref{eq:ch2.modal-particular-solution}
is the spectral form of Eq.~\eqref{eq:ch2.inverse-source-equation} on the radiative range;
Tikhonov and regularized inversions from Subsection~\ref{sec:ch244} damp ill-conditioned
\(\lambda_n\) without altering Fredholm index \cite{stratton1941}.
Substitute Eq.~\eqref{eq:ch2.radiation-map-modal-sum} into
\(\mathcal{R}_\omega[\Jimp]=\mathbf{g}\) and use orthogonality of \(\{\mathbf{u}_n\}\) in
\(\mathcal{U}\) \cite{stratton1941}. Modes with \(\lambda_n=0\) do not appear in the range;
modes with \(\lambda_n\neq 0\) invert mode-wise yielding
Eq.~\eqref{eq:ch2.modal-particular-solution}. Kernel directions remain free because
\(\mathcal{R}_\omega\) is not injective on \(\mathcal{J}_\Omega\)
Eq.~\eqref{eq:ch2.ker-R-omega} \cite{harrington2001}.
Cokernel elements pair with unreachable far-zone content: if \(\mathbf{w}\in\ker\mathcal{R}_\omega^{\dagger}\),
then \(\langle \mathbf{g},\mathbf{w}\rangle_{\mathcal{U}}=0\) is required for solvability
Eq.~\eqref{eq:ch2.radiation-solvability-criterion} \cite{stratton1941}. On exterior domains,
continuous angular spectra Eq.~\eqref{eq:ch2.exterior-continuous-spectral-integral} imply
infinite-dimensional cokernel constraints in principle; finite apertures replace them with
rank-\(N\) projections through Eq.~\eqref{eq:ch2.observation-fredholm-composition}
\cite{harrington2001}. Far-pattern coefficients from the eigenbasis express
\(\mathbf{g}\) through \(\mathbf{f}_n\); solvability requires consistency of measured
coefficients with Eq.~\eqref{eq:ch2.modal-particular-solution} on exporting modes only
\cite{stratton1941}.
On the outgoing branch, solvability of \(\mathcal{L}_\omega[\mathbf{u}_\omega]=\mathbf{s}_\omega\)
and solvability of \(\mathcal{R}_\omega[\Jimp]=\mathbf{g}\) are linked through
Eq.~\eqref{eq:ch2.integral-resolvent-equivalence}: forward field construction is single-valued
on \(\mathcal{F}_{\mathrm{SF}}\) while source inversion remains rank-deficient
\cite{stratton1941}. ~\eqref{eq:ch2.coercivity-quotient} is not identical to
\(\mathrm{ind}\,\mathcal{R}_\omega\); export selection removes non-radiative degrees before
pattern inversion \cite{stratton1941}. The Fredholm analysis above treats
\(\mathcal{R}_\omega\) as the primary object for far-zone inverse problems
\cite{harrington2001}.
Fredholm structure on $\mathcal{R}_\omega$ requires linear media, fixed \(\omega\),
outgoing radiation condition, and compactness of \(K_{\mathrm{rad}}\) on the declared source class
\cite{stratton1941}. Resonant cavities violate invertibility of
\(\mathcal{L}_0\) unless quotiented Eq.~\eqref{eq:ch2.coercivity-quotient}; radiation index
must be computed on the outgoing subspace after Eq.~\eqref{eq:ch2.outgoing-radiation-refinement}
\cite{stratton1941}. Loss regularizes small \(\lambda_n\) but does not remove
\(\ker\mathcal{R}_\omega\) Eq.~\eqref{eq:ch2.ker-R-omega} \cite{harrington2001}.
regime qualification Eq.~\eqref{eq:ch2.regime-component-criterion} remains mandatory before transport claims
on recovered \(\mathbf{g}\) \cite{stratton1941}.
As a worked example, a spherical near-field scanner with $N=50$ probe channels is a concrete inversion budget.
The physical configuration places fifty probes on a sphere of radius $R\approx 2\lambda_0$ around
a compact horn feed; each channel records tangential $\E$ or $\H$ depending on probe design
\cite{hansen1988}. After transformation to far-zone coefficients, the composition
$\mathcal{O}_{\mathrm{meas}}\mathcal{R}_\omega$ has rank at most $50$
Eq.~\eqref{eq:ch2.rank-composition-law} \cite{stratton1941}. The observable behavior is
instructive: the first few reconstructed modes match direct far-field cuts taken on the compact
range, while coefficients beyond the fiftieth singular direction track noise rather than new
beam structure.
Even when Eq.~\eqref{eq:ch2.radiation-solvability-criterion} holds for a continuous target
$\mathbf{g}$, discrete recovery satisfies only a truncated modal system; tail error is bounded
by Eq.~\eqref{eq:ch2.compression-error-bound} with the first neglected singular value of the
measurement composition cited from Subsection~\ref{sec:ch26.signature-aperture}
\cite{harrington2001}. The lesson for laboratory planning is direct: adding probes without
enlarging scan radius or improving SNR cannot recover modes the geometry has already annihilated.
Scanner rank is a hard cap that composes with, but does not replace, Fredholm obstructions from
$\ker\mathcal{R}_\omega$ \cite{stratton1941}.
As a worked example, suppose a synthesis target $\mathbf{g}\neq\mathbf{0}$ specifies a far-zone pattern that is
orthogonal to every exporting mode reachable from the declared source class---for example, a
forbidden parity combination above a ground plane when only even sectors are physically
accessible \cite{collin2001}. No amount of feed optimization will satisfy forward
export: the obstruction appears as inconsistent far-zone coefficients in
Eq.~\eqref{eq:ch2.radiation-solvability-criterion}, not as slow convergence of a numerical
solver \cite{harrington2001}. By contrast, when $\mathbf{g}=\mathbf{0}$, the equation is always
solvable with $\Jimp\in\ker\mathcal{R}_\omega$ \cite{stratton1941}; the balanced-loop example
shows that such sources may be strongly measurable in the near zone yet invisible on $\Gamma_R$.
Mode-wise invertibility on the radiative range detects obstructions through $g_n$ on
$\lambda_n=0$ sectors \cite{stratton1941}. Geometric sector constraints on recoverable spectra
are developed next; complete non-radiating classification belongs to Section~\ref{sec:ch272}
\cite{stratton1941}.
\medskip
\noindent
Subsection~\ref{sec:ch263} closed the Fredholm account: compact perturbation, kernel and
cokernel pairings, solvability criterion, and modal reduction on the radiation eigenbasis.
Kernel membership explains invisible source degrees; cokernel membership explains unreachable
pattern targets; recoverable rank from Subsection~\ref{sec:ch26.signature-aperture} caps what
finite instruments can resolve even when continuous solvability holds \cite{stratton1941}.
\subsection{Geometric constraints on radiative spectra}
\label{sec:ch264}
Geometry can forbid a modal sector before instrumentation or inversion enter the story.
On a reflector compact range, a slot above a ground plane may exhibit a persistent pattern null
in a parity-forbidden sector: the null survives probe upgrades and SNR improvements because the
symmetry class, not receiver noise, annihilates that angular content \cite{collin2001}.
An electrically small feed above the same plane supports only low-order multipole export; formal
harmonic expansions may list high-$\ell$ indices that the support condition suppresses long before
$\mathcal{O}_{\mathrm{meas}}$ enters \cite{jackson1999}. The geometric operator below encodes
that selection on modal amplitudes after eigenmodes and Fredholm structure are known
\cite{stratton1941}.
% PHYSICAL-MEMORY:sec:ch264
% PHYSICAL-MEMORY:sec:ch264-b2
\begin{physicalvignette}[Open-ended circular waveguide probe on SNF arm]
A second bench observation matters on Reflector with subreflector and strut scattering. When technicians swap only the receive probe while holding the transmit chain fixed, reciprocity checks whether the reported coupling is property of the structure or artifact of a mismatched flux gauge \cite{stratton1941}. Failed reciprocity usually traces to a near-field tap, a radome standing-wave, or a probe polarization basis that was never aligned with the transmit standard \cite{hansen1988}.
Rayleigh and Stratton treated such ranges as boundary-value laboratories long before digital near-field transforms became routine \cite{stratton1941}. Harrington's method-of-moments culture reinforced the same discipline: metal outline first, harmonic headline second \cite{harrington2001}.
The measurement logic above is made precise in Subsection~\ref{sec:ch264} \cite{stratton1941}.
\end{physicalvignette}
\begin{physicalvignette}[Open-ended circular waveguide probe on SNF arm]
Consider Open-ended circular waveguide probe on SNF arm. Before any coefficient table is printed, technicians record frequency, cable loss, and the mechanical outline that fixes electrical size. The first walk-through of the range is always tactile: flange bolts, probe standoff, turntable boresight, and the wattmeter reference on the receive chain \cite{hansen1988}.
On this hardware the near zone is visually impressive---standing reactive structure, rapid phase rotation, sometimes a visible VSWR ripple at the feed---yet the export question is measured only after outgoing selection on a large sphere \cite{harrington2001}. What the camera sees in the chamber is not automatically what the far-zone pattern claims \cite{stratton1941}.
Engineers therefore learn to separate three readings: what the feed can synthesize on declared support, what propagating content survives on the radiative range, and what a finite instrument can recover at declared tolerance \cite{devaney2004}. Confusing those readings produces the rank gaps and watt failures that later sections quantify \cite{pal2026}.
Rayleigh and Stratton treated such ranges as boundary-value laboratories long before digital near-field transforms became routine \cite{stratton1941}. Harrington's method-of-moments culture reinforced the same discipline: metal outline first, harmonic headline second \cite{harrington2001}.
Subsection~\ref{sec:ch264} states the formal export condition that this range scene presupposes \cite{stratton1941}.
\end{physicalvignette}
Let \(\Omega_s=\operatorname{supp}\Jimp\), outer boundary \(\Gamma\), piecewise material
data \((\varepsilon,\mu)\), outgoing radiation condition, and radiation eigenbasis
\(\{\mathbf{u}_n,\mathbf{f}_n\}_{n\in\mathcal{N}_{\mathrm{rad}}}\) from the developments above
\cite{stratton1941}. Modal coefficients \(a_n\) in Eq.~\eqref{eq:ch2.spectral-far-pattern-expansion}
are identified with \(\lambda_n c_n[\mathfrak{D}]\) on exporting modes
\cite{stratton1941}. Geometric symmetry group \(G_{\mathrm{geo}}\subseteq\mathrm{SO}(3)\times\{+\}\)
acts on sources, boundaries, and constitutive tensors as in Section~\ref{sec:ch01.3.2} by reference
\cite{stratton1941}. Regime qualification Eq.~\eqref{eq:ch2.regime-operator-screen} remains mandatory
before transport claims on any surviving modal sector \cite{harrington2001}.
When symmetry and boundary class are fixed, modal amplitudes do not vary independently: some
linear combinations are forced to zero by image parity or by support even when the export map
would otherwise allow them \cite{stratton1941}. The matrix \(\Gamma_{nm}\) in
Eq.~\eqref{eq:ch2.geometric-symmetry-operator-def} encodes how a symmetry operation on the
physical configuration maps one modal sector into another; a forbidden sector is one whose
unit amplitude vector lies in the kernel of \(\mathcal{C}_{\mathrm{geo}}\)
\cite{stratton1941}. Define a bounded operator \(\mathcal{C}_{\mathrm{geo}}\) on modal
amplitudes by
\begin{equation}
\label{eq:ch2.geometric-symmetry-operator-def}
\big(\mathcal{C}_{\mathrm{geo}}\mathbf{a}\big)_n
=
\sum_{m\in\mathcal{N}_{\mathrm{rad}}}
\Gamma_{nm}\,a_m,
\qquad
\mathbf{a}=\{a_m\}_{m\in\mathcal{N}_{\mathrm{rad}}},
\end{equation}
where \(\Gamma_{nm}\) encodes how geometry and boundary class map modal sectors under
\(G_{\mathrm{geo}}\) \cite{stratton1941}. Equation
\eqref{eq:ch2.geometric-symmetry-operator-def} is the constructive account of
realization of Eq.~\eqref{eq:ch2.geometric-symmetry-operator-def}: when
\(\mathcal{C}_{\mathrm{geo}}\mathbf{e}_n=\mathbf{0}\) for a unit vector \(\mathbf{e}_n\),
mode \(n\) is \emph{geometrically forbidden} even if \(\lambda_n\neq 0\) on the export map.
Geometric symmetry can annihilate modal sectors before instrumentation enters: a forbidden mode
has zero coupling to physically realizable sources even when formal eigenvalues are non-zero
\cite{stratton1941}. Broken symmetry restores coupling---a practical lever in reflector and
array design \cite{collin2001}.
Eq.~\eqref{eq:ch2.radiation-solvability-criterion} holds for other modes
\cite{stratton1941}. Rotation commutation Eq.~\eqref{eq:ch2.radiation-rotation-commutation}
acts on \(\mathbf{f}_n\); \(\mathcal{C}_{\mathrm{geo}}\) acts on physically accessible
\emph{sectors}
after symmetry reduction \cite{harrington2001}.
Compact support \(\Omega_s\subset\mathbb{R}^{3}\) implies entire-function structure in
spatial wavenumber: for each outgoing mode \(\mathbf{f}_n(\theta,\phi)\),
\begin{equation}
\label{eq:ch2.support-spectral-support-condition}
a_n(\Jimp)=0
\ \text{unless}\ 
\mathbf{f}_n\ \text{has non-negligible overlap with}\ 
\mathcal{F}_\omega[\mathcal{J}_{\Omega_s}],
\end{equation}
with \(\mathcal{J}_{\Omega_s}\) the restriction of admissible sources to \(\Omega_s\)
\cite{stratton1941}. Equation~\eqref{eq:ch2.support-spectral-support-condition}
is a support-induced geometric filter distinct from Fredholm cokernel constraints
Eq.~\eqref{eq:ch2.radiation-solvability-criterion}: electrically small \(\Omega_s\)
suppresses high-\(\ell\) sectors in Eq.~\eqref{eq:ch2.spectral-far-pattern-expansion}
through \((k_0 a)^{2\ell}\) scaling established in Chapter~\ref{ch:03} (\emph{Angular Momentum Operators and SO(3)})
\cite{jackson1999}. Rank cap Eq.~\eqref{eq:ch2.reconstruction-rank-cap}
limits measurable \(a_n\) separately from support annihilation
\cite{harrington2001}.
Perfect conductors, symmetry planes, and equivalent image configurations impose linear
relations among modal amplitudes:
\begin{equation}
\label{eq:ch2.boundary-image-constraint}
\sum_{m\in\mathcal{N}_{\mathrm{rad}}}
M_{nm}(\Gamma,\mathcal{B}_\Gamma)\,a_m=0
\quad\Longrightarrow\quad
a_n=0\ \text{for forbidden parity sectors,}
\end{equation}
with \(M_{nm}\) fixed by tangential boundary class
Eq.~\eqref{eq:ch1.outer-boundary-classes} and image rules on outgoing Green branches
Eq.~\eqref{eq:ch2.outgoing-green-branch} \cite{stratton1941}. Equation
\eqref{eq:ch2.boundary-image-constraint} specializes degeneracy lifting under broken symmetry (Subsection~\ref{sec:ch262}): a ground plane replaces full
\(\mathrm{SO}(3)\) by a reflection group and annihilates odd or even combinations per
declared image parity \cite{stratton1941}. Sommerfeld outgoing radiation condition from
Section~\ref{sec:ch24} is presupposed; incoming image branches are excluded from
\(\mathcal{F}_{\mathrm{rad}}\) \cite{harrington2001}.
When \(G_{\mathrm{geo}}\) contains mirror reflection \(\sigma\), modal sectors split into
parity classes \(\Pi_\pm\) with
\begin{equation}
\label{eq:ch2.parity-selection-rule}
a_n=0\ \text{for all}\ n\in\mathcal{N}_{\mathrm{rad}}
\ \text{with}\ 
\Pi_\pm\mathbf{f}_n=\mp\mathbf{f}_n
\ \text{and forbidden parity,}
\end{equation}
as in the ground-plane illustration of Subsection~\ref{sec:ch262}
\cite{stratton1941}. Equation~\eqref{eq:ch2.parity-selection-rule} explains
pattern nulls that are not explained by \(\ker\mathcal{O}_{\mathrm{meas}}\) or
\(\ker\mathcal{R}_\omega\) alone \cite{stratton1941}. Slot and patch symmetries reduce
Eq.~\eqref{eq:ch2.parity-selection-rule} to discrete sector labels on
\(\mathcal{N}_{\mathrm{rad}}\) without re-tabulating dyads from Subsection~\ref{sec:ch231}
\cite{harrington2001}.
\begin{figure}[t]
\centering
\ChOneFigBlock{%
\begin{tikzpicture}[ch1fig, node distance=5mm and 7mm]
 \node[ch1box] (g) {$G_{\mathrm{geo}}$};
 \node[ch1box, right=7mm of g] (c) {$\mathcal{C}_{\mathrm{geo}}$};
 \node[ch1box, right=7mm of c] (a) {$\{a_n\}$};
 \node[ch1box, right=7mm of a] (f) {$\mathbf{f}_n$};
 \node[ch1box, below=5mm of c] (forb) {forbidden sectors};
 \draw[ch1flow] (g) -- (c) -- (a) -- (f);
 \draw[ch1flow] (c.south) -- (c.south |- forb.north);
 \draw[ch1flow] (c.south |- forb.north) -- (forb.north);
\end{tikzpicture}%
}
\caption{Geometric filter: symmetry group \(G_{\mathrm{geo}}\)
acts through \(\mathcal{C}_{\mathrm{geo}}\) Eq.~\eqref{eq:ch2.geometric-symmetry-operator-def}
on modal amplitudes before far patterns \(\mathbf{f}_n\).}
\label{fig:ch2.geometric-spectrum-constraints}
\end{figure}
Geometric, Fredholm, and observation constraints multiply:
\begin{equation}
\label{eq:ch2.constraint-composition-law}
\mathcal{N}_{\mathrm{reportable}}
\subseteq
\mathcal{N}_{\mathrm{acc}}
\cap
\{n:\ \text{resolvable by}\ \mathcal{O}_{\mathrm{meas}}\}
\cap
\{n:\ \lambda_n\neq 0\},
\end{equation}
with the last two factors from Subsection~\ref{sec:ch263} and
Eq.~\eqref{eq:ch2.reconstruction-rank-cap} \cite{stratton1941}. Equation
\eqref{eq:ch2.constraint-composition-law} prevents attributing a pattern null solely to
finite aperture when Eq.~\eqref{eq:ch2.geometric-symmetry-operator-def} already forbids the
sector \cite{stratton1941}. Inverse source Eq.~\eqref{eq:ch2.inverse-source-equation} on
\(\mathcal{N}_{\mathrm{reportable}}\) remains rank-deficient through
\(\ker\mathcal{R}_\omega\) Eq.~\eqref{eq:ch2.ker-R-omega} even when geometry is
exact \cite{harrington2001}. Geometric constraints on radiative spectra assume linear media, fixed \(\omega\), and symmetry of
\((\Omega_s,\Gamma,\varepsilon,\mu)\) under declared \(G_{\mathrm{geo}}\)
\cite{stratton1941}. Weak symmetry breaking perturbs
Eq.~\eqref{eq:ch2.accessible-modal-sector} through lifted degeneracy
rather than strict zero \cite{stratton1941}. Loss and
anisotropy modify \(\Gamma_{nm}\) and may reopen partially forbidden sectors at finite
amplitude \cite{stratton1941}; regime qualification
Eq.~\eqref{eq:ch2.regime-component-criterion} is unchanged, for geometric accessibility does
not promote reactive near fields to export \cite{harrington2001}.
\medskip
\noindent
Eigenmode ordering, rotational degeneracy, Fredholm solvability, and geometric filters now
compose the spectral chain of Eq.~\eqref{eq:ch2.section26-spectral-chain} on the integral
representations from Section~\ref{sec:ch25} \cite{stratton1941}. The thread is laboratory-first:
finite aperture and $\sigma_n$ decay fix recoverable rank once in
Subsection~\ref{sec:ch26.signature-aperture}; export eigenvalues $\lambda_n$ of
$\mathcal{R}_\omega$ then organize dominant directions; rotation, Fredholm structure, and
geometry explain what inversion can and cannot recover. Section~\ref{sec:ch27} completes
non-radiating null spaces at the source level; \ref{ch:03} supplies explicit harmonic bases;
\ref{ch:04} addresses which eigenstructures finite sources can realize after truncation
\cite{harrington2001}.
\section{Non-Radiating Sources and Null Spaces}
\label{sec:ch27}
Spectral theory in Section~\ref{sec:ch26} ranked export directions and recoverable rank.
Section~\ref{sec:ch27} addresses the complementary laboratory question: \emph{How can current
exist without observable radiation?} \cite{devaney2004}
% PHYSICAL-MEMORY:sec:ch27
% PHYSICAL-MEMORY:sec:ch27-b2
\begin{physicalvignette}[Patch on ceramic substrate for GNSS handset]
A second bench observation matters on Horn-fed reflector with spillover onto the deck. When technicians swap only the receive probe while holding the transmit chain fixed, reciprocity checks whether the reported coupling is property of the structure or artifact of a mismatched flux gauge \cite{stratton1941}. Failed reciprocity usually traces to a near-field tap, a radome standing-wave, or a probe polarization basis that was never aligned with the transmit standard \cite{hansen1988}.
Rayleigh and Stratton treated such ranges as boundary-value laboratories long before digital near-field transforms became routine \cite{stratton1941}. Harrington's method-of-moments culture reinforced the same discipline: metal outline first, harmonic headline second \cite{harrington2001}.
The measurement logic above is made precise in Subsection~\ref{sec:ch27} \cite{stratton1941}.
\end{physicalvignette}
\begin{physicalvignette}[Patch on ceramic substrate for GNSS handset]
Consider Patch on ceramic substrate for GNSS handset. Before any coefficient table is printed, technicians record frequency, cable loss, and the mechanical outline that fixes electrical size. The first walk-through of the range is always tactile: flange bolts, probe standoff, turntable boresight, and the wattmeter reference on the receive chain \cite{hansen1988}.
On this hardware the near zone is visually impressive---standing reactive structure, rapid phase rotation, sometimes a visible VSWR ripple at the feed---yet the export question is measured only after outgoing selection on a large sphere \cite{harrington2001}. What the camera sees in the chamber is not automatically what the far-zone pattern claims \cite{stratton1941}.
Engineers therefore learn to separate three readings: what the feed can synthesize on declared support, what propagating content survives on the radiative range, and what a finite instrument can recover at declared tolerance \cite{devaney2004}. Confusing those readings produces the rank gaps and watt failures that later sections quantify \cite{pal2026}.
Rayleigh and Stratton treated such ranges as boundary-value laboratories long before digital near-field transforms became routine \cite{stratton1941}. Harrington's method-of-moments culture reinforced the same discipline: metal outline first, harmonic headline second \cite{harrington2001}.
Subsection~\ref{sec:ch27} states the formal export condition that this range scene presupposes \cite{stratton1941}.
\end{physicalvignette}
Laboratory signatures. Three configurations recur before any kernel is displayed. A balanced wire loop in vacuum,
driven at fixed \(\omega\) with vanishing net electric dipole moment, stores reactive \(\E\) and
\(\H\) within \(\lambda_0\) of the conductor while a qualified wattmeter on \(\Gamma_R\) at
\(R\gg\lambda_0\) records zero outgoing transverse flux once \(\Pi_{\mathrm{rad}}\) selection is
enforced \cite{stratton1941}. A counter-directed dipole pair separated by
\(d\ll\lambda_0\) can annihilate leading \(1/r\) transverse fields on \(S^{2}\) while reactive
overlap remains in the gap \cite{jackson1999}. A waveguide section below cutoff carries standing
reactive fields at ports yet contributes no propagating sheet to exterior flux meters
\cite{stratton1941}. In each case Maxwell permits substantial local structure while
\(\mathcal{R}_\omega[\Jimp]=\mathbf{0}\): rich near-zone data, null far-zone export---not a
solver artifact \cite{jackson1999}.
The \emph{radiation kernel} collects every admissible impressed current whose outgoing
radiative image vanishes: currents that reorganize near-zone energy without qualified far-zone
flux \cite{devaney2004}. Near-field scanners sample
\(\mathcal{S}_\omega[\Jimp]\) on enclosing shells; far-zone flux meters and angular integrators
read only after \(\Pi_{\mathrm{rad}}\) has removed non-export content
\cite{hansen1988}. A source invisible to the latter may remain vivid to the
former---the central diagnostic that separates reactive storage from qualified export. Campaigns
that report singular values without kernel awareness merge synthesis with invisible modes
\cite{jackson1999}.
By the \emph{radiation null space} \(\ker\mathcal{R}_\omega\) we mean precisely this kernel:
admissible impressed currents whose radiative image vanishes under the outgoing map
\(\mathcal{R}_\omega=\Pi_{\mathrm{rad}}\circ\mathcal{S}_\omega\)---physically realizable
sources that store energy locally yet export no qualified far-zone power
\cite{devaney2004}. Membership in \(\ker\mathcal{R}_\omega\) is an export
statement; it does not require \(\mathcal{S}_\omega[\Jimp]=\mathbf{0}\). The balanced loop
above is a constructive instance: \(\mathcal{R}_\omega[\Jimp^{\mathrm{(loop)}}]=\mathbf{0}\)
while \(\mathcal{S}_\omega[\Jimp^{\mathrm{(loop)}}]\) remains measurable on near probes
\cite{devaney2004}.
Non-radiating sources are the central electromagnetic proof that existence does not imply
observability. A current distribution may generate substantial near-zone $\E$ and $\H$ while
exporting zero far-zone power: $\mathcal{R}_\omega$ is not injective, and
$\ker\mathcal{R}_\omega$ is generally infinite-dimensional on open domains. Different impressed
currents may therefore produce identical radiative images, angular spectra, and reciprocity
readouts while differing in reactive storage and evanescent overlap. Hidden electromagnetic
structure is not a modelling defect but a consequence of Maxwell propagation from compact
sources through outgoing selection to finite instruments \cite{stratton1941}. The absence of
radiated signatures does not imply the absence of electromagnetic structure; it identifies
source configurations whose field content remains inaccessible to far-zone observation.
Section~\ref{sec:ch26} closed the spectral chain
Eq.~\eqref{eq:ch2.section26-spectral-chain}: radiation eigenmodes, rotational degeneracy,
Fredholm solvability, and geometric filters on admissible modal sectors
\cite{stratton1941}. That program organized
\(\mathcal{R}_\omega[\Jimp]\) in modal coordinates and explained which far-zone patterns are
reachable, degenerate, or symmetry-forbidden. Section~\ref{sec:ch27} addresses the
complementary \emph{source-side} question named above: which impressed currents produce no radiative
export at all, how those currents relate to evanescent and non-propagating field content,
and what null-space structure implies for inverse problems and angular spectra
\cite{stratton1941}. Eq.~\eqref{eq:ch2.ker-R-omega} from
Subsection~\ref{sec:ch244} named \(\ker\mathcal{R}_\omega\); the subsections below develop
constructive catalog, canonical non-radiating configurations, and reconstructability
limits without re-opening Green construction, Sommerfeld completion, or radiation integrals from
Sections~\ref{sec:ch23}--\ref{sec:ch25} \cite{stratton1941}. Reciprocity constraints on
coupled modes belong to Section~\ref{sec:ch28}; plane-wave angular-spectrum action to
Section~\ref{sec:ch29}; explicit harmonic bases to \ref{ch:03}
\cite{harrington2001}. Nullity, $\mathcal{E}_{\mathrm{acc}}^{(N)}$, and realizability. Three deficits are often merged in inverse problems and structured-beam reporting
\cite{devaney2004}. \emph{Nullity} is source-side membership in
\(\ker\mathcal{R}_\omega\): far data cannot detect additions satisfying
Eq.~\eqref{eq:ch2.angular-spectrum-nr-invariance} \cite{reed1950}.
\emph{probe rank} is field-side: which directions in the angular spectrum survive
geometric, transport, and regime qualifications even when \(\mathcal{R}_\omega[\Jimp]\neq\mathbf{0}\)
\cite{harrington2001}. \emph{$\mathcal{S}_{\mathrm{real}}$}---whether a target pattern is
\(\mathcal{R}_\omega[\Jimp]\) for some admissible \(\Jimp\)---is not decided by rank-$N$ recovery
alone; \ref{ch:04} develops synthesizable proofs while Section~\ref{sec:ch27} supplies the
forward null-space catalog that constrains every claim \cite{colton1992}.
A source may store reactive energy nearby yet export no far-zone power---the radiation map annihilates whole families of impressed currents. Field-level
uniqueness after outgoing radiation condition (exterior radiative uniqueness (Subsection~\ref{sec:ch243})) does
not imply source identifiability: many admissible \(\Jimp\in\mathcal{J}_\Omega\) can share
the same \(\mathcal{R}_\omega[\Jimp]\) through Eq.~\eqref{eq:ch2.ker-R-omega}
\cite{stratton1941}. Spectral mode \(\lambda_n=0\) in
Eq.~\eqref{eq:ch2.modal-eigenvalue-equation} marks a radiative eigen-direction with
vanishing export, yet \(\ker\mathcal{R}_\omega\) is generally infinite-dimensional on open
domains and is not exhausted by finitely many modal labels
\cite{harrington2001}. Using cavity modal books or truncated harmonic sums to
\emph{approximate} \(\ker\mathcal{R}_\omega\) misstates identifiability even when outgoing
source-to-field integrals from Subsection~\ref{sec:ch251} are evaluated exactly
\cite{jackson1999}. Section~\ref{sec:ch27} separates operator null spaces,
field-side non-radiative components Eqs.~\eqref{eq:ch2.radiative-component-def}--\eqref{eq:ch2.nonradiative-component-def},
and observation null spaces Eq.~\eqref{eq:ch2.observation-null-space}, then forwards
reciprocity and angular-spectrum filters to Sections~\ref{sec:ch28}--\ref{sec:ch29}
\cite{stratton1941}. Let \(\Omega_s=\operatorname{supp}\Jimp\) be compact, \(\omega>0\) fixed, and outgoing field
pairs lie in \(\mathcal{F}_{\mathrm{SF}}\) from
Eq.~\eqref{eq:ch2.admissible-field-space}. Section~\ref{sec:ch27} organizes analysis in four
stages:
\begin{equation}
\label{eq:ch2.section27-nullspace-chain}
\ker\mathcal{R}_\omega
\ \xrightarrow{\ \text{operator structure}\ }
\text{null-space bases}
\ \xrightarrow{\ \text{canonical sources}\ }
\text{NR current families}
\ \xrightarrow{\ \text{evanescent overlap}\ }
\text{non-reconstructability}
\ \xrightarrow{\ \text{angular spectra}\ }
\text{suppressed sheets},
\end{equation}
where the first arrow is Subsection~\ref{sec:ch271}, the second is Subsection~\ref{sec:ch272},
the third is Subsection~\ref{sec:ch273}, and the fourth is Subsection~\ref{sec:ch274}
\cite{stratton1941}. Equation~\eqref{eq:ch2.section27-nullspace-chain}
expresses Eq.~\eqref{eq:ch2.section26-spectral-chain} at the source level: spectral
\(\lambda_n=0\) modes from Eq.~\eqref{eq:ch2.spectral-resolution-composition} lie in
\(\mathbf{u}_{\mathrm{nr}}\), but complete \(\ker\mathcal{R}_\omega\) requires impressed-data
coordinates on \(\mathcal{J}_\Omega\), not only far-zone modal shapes
\cite{stratton1941}.
Equation~\eqref{eq:ch2.section27-nullspace-chain} is the source-side restatement
of the spectral chain Eq.~\eqref{eq:ch2.section26-spectral-chain}: operator
nullity precedes constructive NR catalogs, which precede evanescent
non-reconstructability, which precede angular-spectrum consequences.
Mathematically, each arrow preserves linearity of \(\mathcal{R}_\omega\);
physically, each stage removes a distinct laboratory error from inverse-source
charting. On a cube of side \(2a\) with \(k_0 a=2\), evanescent overlap with
\(\Pi_{\mathrm{ev}}\mathcal{S}_\omega[\Jimp]\) can exceed \(25\%\) of the
near-field energy on \(S_R\) with \(R=a\) while contributing zero to
\(a(\widehat{\mathbf{k}})\); mistaking that storage for NR nullity violates
Eq.~\eqref{eq:ch2.evanescent-nonexport-law} \cite{stratton1941}.
\ref{ch:01} flux laws enter only after \(\Pi_{\mathrm{rad}}\)
\cite{harrington2001}.
Composition law Eq.~\eqref{eq:ch2.constraint-composition-law} still
applies: geometric annihilation Eq.~\eqref{eq:ch2.geometric-symmetry-operator-def} and
observation rank Eq.~\eqref{eq:ch2.reconstruction-rank-cap} are independent of, and
composable with, source null-space membership \cite{harrington2001}.
The laboratory examples and kernel reading above fix the measurement ordering before any
operator display. Near-field probes sample \(\mathcal{S}_\omega[\Jimp]\); flux meters read only
after \(\Pi_{\mathrm{rad}}\) \cite{stratton1941}. With that ordering in place,
keep three spaces in view throughout Section~\ref{sec:ch27}: impressed currents
\(\Jimp\in\mathcal{J}_\Omega\) on compact support \(\Omega_s\), full Maxwell fields
\(\mathcal{S}_\omega[\Jimp]\in\mathcal{F}_\Omega\), and the radiative image on which far-zone
instruments read. The primary object remains the radiation map that projects the latter from the
former:
\begin{equation}
\label{eq:ch2.nullspace-operator-targets}
\mathcal{R}_\omega:\mathcal{J}_\Omega\longrightarrow\mathcal{F}_{\mathrm{rad}},
\qquad
\mathcal{R}_\omega[\Jimp]=\Pi_{\mathrm{rad}}\,\mathcal{S}_\omega[\Jimp],
\end{equation}
from Eqs.~\eqref{eq:ch2.radiation-operator-def}--\eqref{eq:ch2.outgoing-radiation-refinement}
\cite{stratton1941}. The displayed map is the mathematical encoding of the wattmeter reading:
only \(\Pi_{\mathrm{rad}}\mathcal{S}_\omega[\Jimp]\) enters \(\mathcal{F}_{\mathrm{rad}}\), so
near-zone energy in \(\mathcal{S}_\omega[\Jimp]\setminus\mathcal{F}_{\mathrm{rad}}\) cannot be
inferred from far data alone \cite{stratton1941}. Secondary targets include the observation map
\(\mathcal{M}_\omega=\mathcal{O}_{\mathrm{meas}}\circ\Lambda_{\mathrm{rad}}\) from
Eq.~\eqref{eq:ch2.reconstruction-observation-map}, the non-radiative field projector
\(\mathrm{id}-\Pi_{\mathrm{rad}}\) on \(\mathcal{F}_\Omega\), and evanescent/non-propagating
selectors Eqs.~\eqref{eq:ch2.spectral-projectors}--\eqref{eq:ch2.nonpropagating-component-def}
from Subsection~\ref{sec:ch224} \cite{stratton1941}. Inverse source Eq.~\eqref{eq:ch2.inverse-source-equation}
is ill-posed whenever \(\ker\mathcal{M}_\omega\neq\{\mathbf{0}\}\), which occurs if
\(\ker\mathcal{R}_\omega\neq\{\mathbf{0}\}\) or \(\ker\mathcal{O}_{\mathrm{meas}}\) is
nontrivial on \(\operatorname{ran}\mathcal{R}_\omega\) \cite{harrington2001}.
reciprocity-based solvability Eq.~\eqref{eq:ch2.radiation-solvability-criterion} governs reachability
of a \emph{target} radiative field; Section~\ref{sec:ch27} classifies sources that map to
\(\mathbf{0}\in\mathcal{F}_{\mathrm{rad}}\) regardless of solvability for nonzero targets
\cite{stratton1941}.
Radiation null space (leading template). Subsection~\ref{sec:ch271} develops Eq.~\eqref{eq:ch2.ker-R-omega} to
structural statements on \(\ker\mathcal{R}_\omega\) and \(\ker\mathcal{R}_\omega^{\dagger}\).
Locally, \(\ker\mathcal{R}_\omega\) collects impressed currents annihilated by export;
\(\ker\mathcal{R}_\omega^{\dagger}\) collects far-side test directions orthogonal to every
radiative image; and \(\mathrm{coker}\,\mathcal{R}_\omega=\mathcal{F}_{\mathrm{rad}}/\operatorname{ran}\mathcal{R}_\omega\)
labels pattern targets that no admissible source can realize \cite{stratton1941}. The leading template
below records admissibility and typical infinite dimensionality on open domains:
\begin{equation}
\label{eq:ch2.radiation-nullspace-structure}
\ker\mathcal{R}_\omega
=
\left\{
\Jimp\in\mathcal{J}_\Omega:
\Pi_{\mathrm{rad}}\,\mathcal{S}_\omega[\Jimp]=\mathbf{0},
\ \mathcal{D}_{J}[\Jimp]=1
\right\},
\qquad
\dim\ker\mathcal{R}_\omega=\infty\ \text{typ. on open domains},
\end{equation}
with $\mathcal{D}_J=1$ Eq.~\eqref{eq:ch2.source-admissibility} enforced throughout
\cite{stratton1941}. Equation~\eqref{eq:ch2.radiation-nullspace-structure}
does not re-define Eq.~\eqref{eq:ch2.ker-R-omega}; it records the \ref{ch:02}
primary target for constructive characterization of nullity, dimensionality, and pairing with cokernel
directions obstructing inversion Eq.~\eqref{eq:ch2.spectral-solvability}
\cite{stratton1941}. Modal sectors with \(\lambda_n=0\) in
Eq.~\eqref{eq:ch2.spectral-resolution-composition} supply a finite-dimensional \emph{slice}
of \(\ker\mathcal{R}_\omega\) only when eigenmodes are declared on outgoing subspaces
\cite{harrington2001}.
Non-radiating current configurations (leading template). Subsection~\ref{sec:ch272} develops canonical impressed currents that realize
Eq.~\eqref{eq:ch2.null-radiation}:
\begin{equation}
\label{eq:ch2.nr-current-template}
\Jimp^{\mathrm{(NR)}}\in\ker\mathcal{R}_\omega,
\qquad
\mathcal{S}_\omega[\Jimp^{\mathrm{(NR)}}]\neq\mathbf{0},
\qquad
\|\mathbf{u}_{\mathrm{nr}}\|\ \text{may be large near}\ \Omega_s,
\end{equation}
with \(\mathbf{u}_{\mathrm{nr}}\) from Eq.~\eqref{eq:ch2.nonradiative-component-def}
\cite{stratton1941}. Equation~\eqref{eq:ch2.nr-current-template}
specializes Eq.~\eqref{eq:ch2.nr-image-of-null-space} to constructive families---loop modes,
counter-propagating pairs, and equivalent-surface cancellations compatible with
Eq.~\eqref{eq:ch1.charge-continuity-harmonic}---without tabulating dyadic kernels from
Subsection~\ref{sec:ch231} \cite{stratton1941}. Such currents are invisible to far-zone flux
meters Eq.~\eqref{eq:ch2.power-flux-functional} yet may dominate near-zone energy budgets
Eq.~\eqref{eq:ch2.component-power-budget} \cite{harrington2001}.
Evanescent overlap and non-reconstructability (leading template). Subsection~\ref{sec:ch273} links \(\ker\mathcal{R}_\omega\) to evanescent subspaces and
finite observation rank. Evanescent content is field-side reactive storage with
\(\Pi_{\mathrm{rad}}\mathbf{u}_{\mathrm{ev}}=\mathbf{0}\); it must not be conflated with
source-side \(\Jimp\in\ker\mathcal{R}_\omega\), although both deficits can coexist on the same
excitation \cite{stratton1941}. The relation records the measurement inclusion and
the evanescent non-export law:
\begin{equation}
\label{eq:ch2.evanescent-nonreconstruct}
\ker\mathcal{O}_{\mathrm{meas}}\circ\mathcal{R}_\omega
\supseteq
\ker\mathcal{R}_\omega,
\qquad
\Pi_{\mathrm{ev}}\mathcal{S}_\omega[\Jimp]
\ \text{never exports through}\ \Pi_{\mathrm{rad}},
\end{equation}
using Eqs.~\eqref{eq:ch2.rad-null-on-evanescent} and~\eqref{eq:ch2.reconstruction-rank-cap}
\cite{stratton1941}. Equation~\eqref{eq:ch2.evanescent-nonreconstruct}
separates \emph{radiation nullity} (no outgoing transverse export) from \emph{measurement
nullity} (unobservable even when export is nonzero) \cite{stratton1941}. Tikhonov stabilization
Eq.~\eqref{eq:ch2.tikhonov-source-reconstruction} does not shrink \(\ker\mathcal{R}_\omega\);
regularization only damps ill-conditioned directions on \(\operatorname{ran}\mathcal{R}_\omega\)
\cite{harrington2001}. Accessible modal sector
Eq.~\eqref{eq:ch2.accessible-modal-sector} from Subsection~\ref{sec:ch264} further restricts
recoverable content independently of evanescent storage \cite{stratton1941}.
Angular-spectrum consequences (leading template). Subsection~\ref{sec:ch274} records how null-space structure propagates to continuous angular
spectra Eq.~\eqref{eq:ch2.exterior-continuous-spectral-integral}:
\begin{equation}
\label{eq:ch2.angular-spectrum-null}
a(\widehat{\mathbf{k}})=0
\ \text{on suppressed sheets}\ \mathcal{K}_{\mathrm{sup}}
\ \Longrightarrow\
\Jimp\ \text{may still carry reactive content in}\ \Pi_{\mathrm{ev}}\mathcal{S}_\omega[\Jimp],
\end{equation}
with \(a(\widehat{\mathbf{k}})\) the far angular amplitude and
\(\mathcal{K}_{\mathrm{sup}}\) the complement of accessible sector
Eq.~\eqref{eq:ch2.angular-accessibility-sector} \cite{stratton1941}. Equation
\eqref{eq:ch2.angular-spectrum-null} links Section~\ref{sec:ch27} to
Section~\ref{sec:ch29}: plane-wave filters and retained/suppressed components in
Subsections~\ref{sec:ch291}--\ref{sec:ch294} presuppose the
source-side null-space catalog developed here \cite{stratton1941}. Angular probe rank on the far shell (Subsection~\ref{sec:ch264}) is not re-proved here;
Subsection~\ref{sec:ch274} states angular-spectrum \emph{consequences} of
Eq.~\eqref{eq:ch2.ker-R-omega} \cite{harrington2001}.
\begin{equation}
\label{eq:ch2.nullspace-spectral-composition}
\Jimp\in\ker\mathcal{R}_\omega
\ \Longrightarrow\
a_n(\Jimp)=0\ \text{for all exporting modes with}\ \lambda_n\neq 0,
\end{equation}
while geometric constraints Eq.~\eqref{eq:ch2.accessible-modal-sector} may force additional
\(a_n=0\) for \(\Jimp\notin\ker\mathcal{R}_\omega\) \cite{stratton1941}.
Equation~\eqref{eq:ch2.nullspace-spectral-composition} follows from
Eq.~\eqref{eq:ch2.spectral-far-pattern-expansion} and linearity
Eq.~\eqref{eq:ch2.radiation-linearity}; it is cited in Section~\ref{sec:ch27}
\cite{stratton1941}. Inverse problems must budget \(\ker\mathcal{R}_\omega\),
\(\ker\mathcal{O}_{\mathrm{meas}}\), and forbidden sectors from
geometric constraints on radiative spectra (Subsection~\ref{sec:ch264}) separately
\cite{harrington2001}.
Several constructions deliberately remain outside Section~\ref{sec:ch27}. Reciprocity
identities and modal orthogonality belong to Section~\ref{sec:ch28}; plane-wave angular-spectrum
operators and near-to-far transforms to Section~\ref{sec:ch29}; vector spherical harmonic
completeness and explicit \(\mathbf{f}_n\) to \ref{ch:03}; range membership of
unreachable modes to \ref{ch:04} \cite{stratton1941}. Section~\ref{sec:ch27}
Subsection~\ref{sec:ch261} supplies spectral decomposition;
Fredholm structure of radiation operators (Subsection~\ref{sec:ch263}), or geometric constraints
on radiative spectra (Subsection~\ref{sec:ch264}); it classifies invisible sources on the
maps those results presuppose \cite{stratton1941}. Worked non-radiating examples appear
in Subsection~\ref{sec:ch2102} by reference to Subsection~\ref{sec:ch272}
\cite{harrington2001}.
Regime logic from Section~\ref{sec:ch01.2.4} is unchanged. Null-space membership does not
override transport qualification: a current in \(\ker\mathcal{R}_\omega\) has
\(\mathcal{P}_{\mathrm{rad}}=0\) on qualified shells, while reactive near fields from
\(\mathcal{S}_\omega[\Jimp]\) may still violate naive flux interpretations if
Eq.~\eqref{eq:ch2.regime-operator-screen} is ignored \cite{stratton1941}.
Far-zone dominance far-zone radiative dominance (Subsection~\ref{sec:ch254}) applies only to
exporting sources; it does not trivialize \(\ker\mathcal{R}_\omega\) \cite{stratton1941}.
Compression error Eq.~\eqref{eq:ch2.compression-error-bound} bounds modal truncation on
\(\operatorname{ran}\mathcal{R}_\omega\) independently of null-space dimension
\cite{harrington2001}.
As a worked example, consider a thin rectangular wire loop of perimeter \(2\pi b\) with \(b\ll\lambda_0\), embedded in
free space and driven at angular frequency \(\omega\) so that the impressed current completes one
closed circulation with equal and opposite arms along the loop's long sides. The physical setup is
a balanced dipole pair folded into a single closed contour: magnetic flux linkage is negligible at
\(\omega\), yet the loop stores reactive energy in the near zone. A closed admissible loop carrying
oppositely directed arms can cancel far-zone transverse radiation through
Eq.~\eqref{eq:ch2.radiation-integral-outgoing-upgrade} while maintaining nonzero
\(\mathcal{S}_\omega[\Jimp]\) near the loop \cite{stratton1941}.
Subsection~\ref{sec:ch272} classifies such families; here the point is that
Eq.~\eqref{eq:ch2.null-radiation} is realizable with
\(\mathcal{D}_{J}[\Jimp]=1\), not merely a formal kernel direction
\cite{harrington2001}. Pattern-matching algorithms that fit only
\(\operatorname{ran}\mathcal{R}_\omega\) assign spurious modal labels to invisible loop content
\cite{stratton1941}.
As a worked example, take a compact spherical source of radius \(a\) with \(k_0 a=1.5\) and truncate the exterior modal
book at \(N=12\) radiation eigenpairs. A single \(\lambda_n=0\) eigenmode from
Eq.~\eqref{eq:ch2.radiation-eigenpair-def} contributes to \(\ker\mathcal{R}_\omega\) when
realized by admissible \(\Jimp^{(n)}\) \cite{stratton1941}. On the range, the truncated
index \(n\le N\) labels only a finite slice of invisible degrees; the full kernel on open
\(\mathcal{J}_\Omega\) is generally much larger. On exterior problems the modal index set is not finite;
truncating to the first \(N\) modes underestimates \(\dim\ker\mathcal{R}_\omega\) and
overstates identifiability of inverse Eq.~\eqref{eq:ch2.inverse-source-equation}
\cite{harrington2001}. Subsection~\ref{sec:ch271} states the correct operator
bookkeeping \cite{stratton1941}.
\subsection{Null spaces of radiation operators}
\label{sec:ch271}
% HARRINGTON-OPEN:sec:ch271
On a declared compact-range or anechoic campaign, technicians record frequency, outline, and probe geometry before any harmonic fit is trusted \cite{hansen1988}. Non-radiating currents store reactive energy without far-zone export \cite{stratton1941}. Near-field richness must not be reported as far-zone diversity \cite{hansen1988}. Maxwell's retarded export to $\Gamma_R$ remains the governing constraint \cite{maxwell1873}.
Section~\ref{sec:ch27} opened with laboratory signatures---balanced loops, counter-directed
pairs, below-cutoff guides---and named the radiation kernel in plain language. Subsection
~\ref{sec:ch271} now supplies the operator account: structural properties of
$\ker\mathcal{R}_\omega$, adjoint null spaces, cokernel pairings, and modal null sectors,
developing Eq.~\eqref{eq:ch2.ker-R-omega} from Subsection~\ref{sec:ch244}
\cite{stratton1941}. The loop and scanner patterns cited in the section introduction are not
re-derived here; they motivate the locked definition below.
% PHYSICAL-MEMORY:sec:ch271
% PHYSICAL-MEMORY:sec:ch271-b2
\begin{physicalvignette}[Dual-band dish with shared reflector surface]
A second bench observation matters on Near-field probe on a hexapod alignment stage. When technicians swap only the receive probe while holding the transmit chain fixed, reciprocity checks whether the reported coupling is property of the structure or artifact of a mismatched flux gauge \cite{stratton1941}. Failed reciprocity usually traces to a near-field tap, a radome standing-wave, or a probe polarization basis that was never aligned with the transmit standard \cite{hansen1988}.
Rayleigh and Stratton treated such ranges as boundary-value laboratories long before digital near-field transforms became routine \cite{stratton1941}. Harrington's method-of-moments culture reinforced the same discipline: metal outline first, harmonic headline second \cite{harrington2001}.
The measurement logic above is made precise in Subsection~\ref{sec:ch271} \cite{stratton1941}.
\end{physicalvignette}
\begin{physicalvignette}[Dual-band dish with shared reflector surface]
Consider Dual-band dish with shared reflector surface. Before any coefficient table is printed, technicians record frequency, cable loss, and the mechanical outline that fixes electrical size. The first walk-through of the range is always tactile: flange bolts, probe standoff, turntable boresight, and the wattmeter reference on the receive chain \cite{hansen1988}.
On this hardware the near zone is visually impressive---standing reactive structure, rapid phase rotation, sometimes a visible VSWR ripple at the feed---yet the export question is measured only after outgoing selection on a large sphere \cite{harrington2001}. What the camera sees in the chamber is not automatically what the far-zone pattern claims \cite{stratton1941}.
Engineers therefore learn to separate three readings: what the feed can synthesize on declared support, what propagating content survives on the radiative range, and what a finite instrument can recover at declared tolerance \cite{devaney2004}. Confusing those readings produces the rank gaps and watt failures that later sections quantify \cite{pal2026}.
Non-radiating currents store reactive energy without far-zone export . Near-field richness must not be reported as far-zone diversity .
Rayleigh and Stratton treated such ranges as boundary-value laboratories long before digital near-field transforms became routine \cite{stratton1941}. Harrington's method-of-moments culture reinforced the same discipline: metal outline first, harmonic headline second \cite{harrington2001}.
Subsection~\ref{sec:ch271} states the formal export condition that this range scene presupposes \cite{stratton1941}.
\end{physicalvignette}
Assumptions. Fix \(\omega>0\), admissible source space \(\mathcal{J}_\Omega\)
Eq.~\eqref{eq:ch2.source-space-refined}, charge continuity and boundary compatibility
Eq.~\eqref{eq:ch2.source-admissibility}, outgoing radiation condition on exterior problems
Eqs.~\eqref{eq:ch2.sommerfeld-outgoing-class}--\eqref{eq:ch2.outgoing-radiation-refinement},
and radiation map \(\mathcal{R}_\omega=\Pi_{\mathrm{rad}}\circ\mathcal{S}_\omega\) from
Eqs.~\eqref{eq:ch2.radiation-operator-def}--\eqref{eq:ch2.outgoing-radiation-refinement}
\cite{stratton1941}. Dual pairings on sources and fields are
Eqs.~\eqref{eq:ch2.dual-source-pairing} and~\eqref{eq:ch2.maxwell-adjoint-pairing}
\cite{stratton1941}. regime qualification Eq.~\eqref{eq:ch2.regime-operator-screen} governs transport
reports on any recovered radiative component \cite{harrington2001}.
Radiation null space (locked definition). Physically, $\ker\mathcal{R}_\omega$ is the set of impressed currents that produce no qualified
far-zone pattern: the wattmeter reading is null even when near probes are not
\cite{stratton1941}. Formally, the \emph{radiation null space} is the source-side kernel
\begin{equation}
\label{eq:ch2.radiation-nullspace-def}
\ker\mathcal{R}_\omega
=
\left\{
\Jimp\in\mathcal{J}_\Omega:
\mathcal{R}_\omega[\Jimp]=\mathbf{0}\in\mathcal{F}_{\mathrm{rad}},
\ \mathcal{D}_{J}[\Jimp]=1
\right\},
\end{equation}
equivalently \(\Pi_{\mathrm{rad}}\,\mathcal{S}_\omega[\Jimp]=\mathbf{0}\) on the outgoing
branch \cite{stratton1941}. Equation~\eqref{eq:ch2.radiation-nullspace-def}
is the Subsection~\ref{sec:ch271} development of Eq.
Eq.~\eqref{eq:ch2.ker-R-omega}: $\mathcal{D}_J=1$ is explicit, and membership is
defined on \(\mathcal{J}_\Omega\) rather than on formal test directions
\cite{stratton1941}. Elements of \(\ker\mathcal{R}_\omega\) produce no qualified far-zone
export through Eq.~\eqref{eq:ch2.radiation-integral-outgoing-upgrade} yet may satisfy
\(\mathcal{S}_\omega[\Jimp]\neq\mathbf{0}\) Eq.~\eqref{eq:ch2.nr-image-of-null-space}
\cite{harrington2001}. Consequence for inversion: any far-zone datum fixes
\(\mathcal{R}_\omega[\Jimp]\) only up to additions in \(\ker\mathcal{R}_\omega\). Linearity
Eq.~\eqref{eq:ch2.radiation-linearity} makes \(\ker\mathcal{R}_\omega\) a linear subspace of
\(\mathcal{J}_\Omega\) closed under admissible superposition \cite{stratton1941}.
Adjoint null space and cokernel. Targets in \(\mathrm{coker}\,\mathcal{R}_\omega\) are far-zone patterns orthogonal to every
reachable image; they obstruct solvability even when no near-field anomaly is visible
\cite{stratton1941}. Let \(\mathcal{R}_\omega^{\dagger}:\mathcal{F}_{\mathrm{rad}}\to\mathcal{J}_\Omega\) denote
the adjoint induced by Eqs.~\eqref{eq:ch2.dual-source-pairing} and
Eq.~\eqref{eq:ch2.maxwell-adjoint-pairing} defines the adjoint pairing. Define
\begin{equation}
\label{eq:ch2.radiation-adjoint-nullspace}
\ker\mathcal{R}_\omega^{\dagger}
=
\left\{
\mathbf{w}\in\mathcal{F}_{\mathrm{rad}}:
\langle \mathbf{w},\,\mathcal{R}_\omega[\Jimp]\rangle_{\mathcal{U}}=0
\ \ \forall\ \Jimp\in\mathcal{J}_\Omega
\right\},
\end{equation}
and the cokernel quotient
\begin{equation}
\label{eq:ch2.radiation-cokernel-def}
\mathrm{coker}\,\mathcal{R}_\omega
=
\mathcal{F}_{\mathrm{rad}}\big/\operatorname{ran}\mathcal{R}_\omega,
\end{equation}
with \(\operatorname{ran}\mathcal{R}_\omega\subset\mathcal{F}_{\mathrm{rad}}\) the radiative
image \cite{stratton1941}. Equations~\eqref{eq:ch2.radiation-adjoint-nullspace}--\eqref{eq:ch2.radiation-cokernel-def}
pair with \(\ker\mathcal{R}_\omega\) through reciprocity-based solvability
Eq.~\eqref{eq:ch2.radiation-solvability-criterion} imported from Subsection~\ref{sec:ch263}:
targets \(\mathbf{g}\in\mathcal{F}_{\mathrm{rad}}\) orthogonal to
\(\ker\mathcal{R}_\omega^{\dagger}\) are reachable, while obstructions lie in cokernel
directions \cite{stratton1941}. Biorthogonality Eq.~\eqref{eq:ch2.radiation-biorthogonality}
expresses the same pairing on modal coordinates without re-deriving adjoint modes here
\cite{harrington2001}.
When \(\mathcal{R}_\omega\) is Fredholm in the sense of Fredholm structure of radiation operators (Subsection~\ref{sec:ch263}),
\begin{equation}
\label{eq:ch2.fredholm-nullspace-alternative}
\dim\ker\mathcal{R}_\omega<\infty
\ \Longleftrightarrow\
\dim\ker\mathcal{R}_\omega^{\dagger}<\infty,
\end{equation}
and index Eq.~\eqref{eq:ch2.radiation-fredholm-index-def} counts the imbalance
\cite{stratton1941}. Equation~\eqref{eq:ch2.fredholm-nullspace-alternative}
does not re-prove Fredholm structure of radiation operators (Subsection~\ref{sec:ch263}); it records the
Chapter~\ref{ch:02} null-space bookkeeping used in inversion
Eq.~\eqref{eq:ch2.inverse-source-equation} \cite{stratton1941}. On exterior domains with
compact \(\Omega_s\), \(\dim\ker\mathcal{R}_\omega\) is typically infinite even when
\(\mathcal{R}_\omega\) is Fredholm on truncated subspaces
Eq.~\eqref{eq:ch2.spectral-compression-operator} \cite{harrington2001}.
On the orthonormal basis of spectral decomposition of the radiation operator (Subsection~\ref{sec:ch261}), define
the \emph{modal null index set}
\begin{equation}
\label{eq:ch2.modal-null-index-set}
\mathcal{N}_{\mathrm{null}}
=
\left\{
n\in\mathcal{N}_{\mathrm{rad}}:
\lambda_n=0
\right\},
\end{equation}
with \(\lambda_n\) from Eq.~\eqref{eq:ch2.radiation-eigenpair-def}
\cite{stratton1941}. Equation~\eqref{eq:ch2.modal-null-index-set} is the
spectral slice of \(\ker\mathcal{R}_\omega\) on declared eigenmodes:
\begin{equation}
\label{eq:ch2.modal-null-sector-reduction}
\sum_{n\in\mathcal{N}_{\mathrm{null}}}
c_n[\mathfrak{D}]\,\Jimp^{(n)}
\in\ker\mathcal{R}_\omega
\quad\text{when admissible superpositions exist},
\end{equation}
using source coefficients Eq.~\eqref{eq:ch2.source-modal-amplitude} and eigenpairs
Eq.~\eqref{eq:ch2.radiation-eigenpair-def} \cite{stratton1941}. Equation
\eqref{eq:ch2.modal-null-sector-reduction} does not exhaust \(\ker\mathcal{R}_\omega\) on
exterior problems because \(\mathcal{N}_{\mathrm{rad}}\) need not span all invisible source
degrees \cite{harrington2001}. modal invertibility on the radiative range (Subsection~\ref{sec:ch263})
already states that coefficients on \(\ker\mathcal{R}_\omega\) remain free in
Eq.~\eqref{eq:ch2.modal-particular-solution}; Subsection~\ref{sec:ch271} identifies the
modal-labeled portion \cite{stratton1941}.
\begin{figure}[t]
\centering
\ChOneFigBlock{%
\begin{tikzpicture}[ch1fig, node distance=5mm and 7mm]
 \node[ch1box] (j) {$\mathcal{J}_\Omega$};
 \node[ch1box, right=7mm of j] (r) {$\mathcal{R}_\omega$};
 \node[ch1box, right=7mm of r] (f) {$\mathcal{F}_{\mathrm{rad}}$};
 \node[ch1box, below=5mm of j] (ker) {$\ker\mathcal{R}_\omega$};
 \node[ch1box, below=5mm of f] (adj) {$\ker\mathcal{R}_\omega^{\dagger}$};
 \draw[ch1flow] (j) -- (r) -- (f);
 \draw[ch1flow, dashed] (ker.south) -- (j.north);
 \draw[ch1flow, dashed] (adj.south) -- (f.north);
\end{tikzpicture}%
}
\caption{Subsection~\ref{sec:ch271} null-space measurement sequence: source kernel
Eq.~\eqref{eq:ch2.radiation-nullspace-def}, radiative image, and adjoint kernel
Eq.~\eqref{eq:ch2.radiation-adjoint-nullspace} paired by solvability
Eq.~\eqref{eq:ch2.radiation-solvability-criterion}.}
\label{fig:ch2.radiation-nullspace-measurement chain}
\end{figure}
Figure interpretation. Read $\mathcal{J}_\Omega$ as the impressed-source Hilbert space and $\mathcal{F}_{\mathrm{rad}}$ as the outgoing radiative image. The forward arrow $\mathcal{R}_\omega=\Pi_{\mathrm{rad}}\circ\mathcal{S}_\omega$ is the only laboratory-export channel; dashed $\ker\mathcal{R}_\omega$ feeds back into $\mathcal{J}_\Omega$ as invisible source degrees that annihilate far-zone patterns Eq.~\eqref{eq:ch2.angular-spectrum-nullspace-law}. The lower adjoint kernel $\ker\mathcal{R}_\omega^{\dagger}$ pairs with cokernel directions through reciprocity-based solvability Eq.~\eqref{eq:ch2.radiation-solvability-criterion}: a nonzero cokernel obstruction means no impressed current can realize the target far field, independent of mesh resolution \cite{stratton1941}. Figure~\ref{fig:ch2.radiation-nullspace-measurement chain} separates kernel and cokernel data before observation composition Eq.~\eqref{eq:ch2.observation-fredholm-composition}.
\phantomsection\label{thm:ch2.radiation-nullspace-structure}
Let \(\mathcal{R}_\omega\) satisfy Eq.~\eqref{eq:ch2.radiation-nullspace-def} on admissible
\(\mathcal{J}_\Omega\) with outgoing radiation condition. Then:
\begin{enumerate}
\item[(i)] \(\ker\mathcal{R}_\omega\) is a linear subspace of \(\mathcal{J}_\Omega\) and
\(\mathcal{R}_\omega[\Jimp]=\mathbf{0}\) iff \(\Jimp\in\ker\mathcal{R}_\omega\).
\item[(ii)] Solvability of \(\mathcal{R}_\omega[\Jimp]=\mathbf{g}\) is equivalent to
Eq.~\eqref{eq:ch2.radiation-solvability-criterion} on declared pairings.
\item[(iii)] On the spectral basis of Subsection~\ref{sec:ch261},
\(\mathcal{N}_{\mathrm{null}}\) Eq.~\eqref{eq:ch2.modal-null-index-set} labels modes with
\(P_n\mathcal{R}_\omega=\mathbf{0}\); modal superpositions
Eq.~\eqref{eq:ch2.modal-null-sector-reduction} lie in \(\ker\mathcal{R}_\omega\) when
admissible.
\item[(iv)] Domain class fixes dimensionality:
\begin{equation}
\label{eq:ch2.nullspace-dimensional-classification}
\begin{aligned}
\text{bounded cavity + discrete book:}\ &
\dim\ker\mathcal{R}_\omega\ \text{finite on truncated}\ \mathcal{R}_\omega^{(N)};\\
\text{exterior + compact}\ \Omega_s:\ &
\dim\ker\mathcal{R}_\omega=\infty\ \text{typ. on full}\ \mathcal{J}_\Omega.
\end{aligned}
\end{equation}
\end{enumerate}. Part (i) follows from linearity Eq.~\eqref{eq:ch2.radiation-linearity} and definition
Eq.~\eqref{eq:ch2.radiation-nullspace-def} \cite{stratton1941}. Part (ii) is Fredholm structure of radiation operators (Subsection~\ref{sec:ch263}) and
Eq.~\eqref{eq:ch2.radiation-solvability-criterion}, cited without re-proof
\cite{harrington2001}. Part (iii) uses Eq.~\eqref{eq:ch2.modal-fredholm-reduction} and
\(\lambda_n=0\) in Eq.~\eqref{eq:ch2.radiation-eigenpair-def} \cite{stratton1941}. Part (iv)
contrasts discrete cavity spectra Eq.~\eqref{eq:ch2.cavity-discrete-radiation-spectrum} with
continuous exterior angular content Eq.~\eqref{eq:ch2.exterior-continuous-spectral-integral}
\cite{jackson1999}.
\phantomsection\label{prop:ch2.null-kernel-field-image}
If \(\Jimp\in\ker\mathcal{R}_\omega\), then radiative and non-radiative components satisfy
\(\mathbf{u}_{\mathrm{rad}}=\mathbf{0}\) and
\(\mathbf{u}_{\mathrm{nr}}=\mathcal{S}_\omega[\Jimp]\) from
Eqs.~\eqref{eq:ch2.radiative-component-def}--\eqref{eq:ch2.nonradiative-component-def};
near-zone energy in \(\mathbf{u}_{\mathrm{nr}}\) may be large while qualified export vanishes
\cite{stratton1941}.
Composition with spectral and geometric filters. Null-space membership implies modal annihilation on exporting sectors:
\phantomsection\label{eq:ch2.geometric-modal-annihilation}
\begin{equation}
\label{eq:ch2.nullspace-modal-annihilation-law}
\Jimp\in\ker\mathcal{R}_\omega
\ \Longrightarrow\
a_n(\Jimp)=0
\ \text{for all}\ n\ \text{with}\ \lambda_n\neq 0,
\end{equation}
with far coefficients \(a_n\) from Eq.~\eqref{eq:ch2.spectral-far-pattern-expansion}
\cite{stratton1941}. Equation~\eqref{eq:ch2.nullspace-modal-annihilation-law}
expresses Eq.~\eqref{eq:ch2.nullspace-spectral-composition} from Section~\ref{sec:ch27} without
re-deriving the modal expansion \cite{stratton1941}. Geometric constraints
geometric constraints on radiative spectra (Subsection~\ref{sec:ch264}) may force additional \(a_n=0\) for sources
outside \(\ker\mathcal{R}_\omega\) \cite{harrington2001}. Composition law
Eq.~\eqref{eq:ch2.constraint-composition-law} budgets identifiability deficits separately
\cite{stratton1941}.
Fredholm index Eq.~\eqref{eq:ch2.radiation-fredholm-index-def} is unchanged under compact
perturbations Eq.~\eqref{eq:ch2.radiation-fredholm-decomposition} and under Tikhonov
regularization Eq.~\eqref{eq:ch2.tikhonov-source-reconstruction} on the range side
\cite{stratton1941}. Regularization stabilizes
Eq.~\eqref{eq:ch2.modal-particular-solution} for small \(|\lambda_n|\) but leaves
\(\ker\mathcal{R}_\omega\) unchanged: invisible source components cannot be recovered from
far-zone data at any regularization level \cite{stratton1941}. Compression truncation
Eq.~\eqref{eq:ch2.compression-error-bound} approximates \(\operatorname{ran}\mathcal{R}_\omega\),
not \(\ker\mathcal{R}_\omega\) \cite{harrington2001}.
Null-space structure of radiation operators (Subsection~\ref{sec:ch271}) assumes linear media at fixed \(\omega\),
outgoing radiation condition, and admissible \(\mathcal{J}_\Omega\) Eq.~\eqref{eq:ch2.source-admissibility}
\cite{stratton1941}. Loss modifies \(\lambda_n\) and may merge modal null sectors
without promoting evanescent content to export through
Eq.~\eqref{eq:ch2.rad-null-on-evanescent} \cite{stratton1941}. Nonlinear or time-varying
sources require the causal chain of Subsection~\ref{sec:ch223}; harmonic null spaces here are
fixed-frequency slices \cite{harrington2001}. regime qualification
Eq.~\eqref{eq:ch2.regime-component-criterion} remains mandatory before interpreting near-zone
reactive dominance as null-space membership \cite{stratton1941}.
\medskip
\noindent
Subsection~\ref{sec:ch271} fixed operator null-space structure through
Eq.~\eqref{eq:ch2.radiation-nullspace-def}, adjoint and cokernel pairings, and the modal null
sector Eqs.~\eqref{eq:ch2.modal-null-index-set}--\eqref{eq:ch2.modal-null-sector-reduction}
\cite{stratton1941}. Constructive non-radiating current families that realize the kernel in the
laboratory develop next.
\phantomsection\label{sec:ch272}
% PHYSICAL-MEMORY:sec:ch272
\begin{physicalvignette}[Compact range feed with quad-ridged horn]
Consider Compact range feed with quad-ridged horn. Before any coefficient table is printed, technicians record frequency, cable loss, and the mechanical outline that fixes electrical size. The first walk-through of the range is always tactile: flange bolts, probe standoff, turntable boresight, and the wattmeter reference on the receive chain \cite{hansen1988}.
% PHYSICAL-MEMORY:sec:ch272-b2
\begin{physicalvignette}[Compact range feed with quad-ridged horn]
A second bench observation matters on Aperture-coupled patch behind a superstrate. When technicians swap only the receive probe while holding the transmit chain fixed, reciprocity checks whether the reported coupling is property of the structure or artifact of a mismatched flux gauge \cite{stratton1941}. Failed reciprocity usually traces to a near-field tap, a radome standing-wave, or a probe polarization basis that was never aligned with the transmit standard \cite{hansen1988}.
Rayleigh and Stratton treated such ranges as boundary-value laboratories long before digital near-field transforms became routine \cite{stratton1941}. Harrington's method-of-moments culture reinforced the same discipline: metal outline first, harmonic headline second \cite{harrington2001}.
The measurement logic above is made precise in Subsection~\ref{sec:ch272} \cite{stratton1941}.
\end{physicalvignette}
On this hardware the near zone is visually impressive---standing reactive structure, rapid phase rotation, sometimes a visible VSWR ripple at the feed---yet the export question is measured only after outgoing selection on a large sphere \cite{harrington2001}. What the camera sees in the chamber is not automatically what the far-zone pattern claims \cite{stratton1941}.
Engineers therefore learn to separate three readings: what the feed can synthesize on declared support, what propagating content survives on the radiative range, and what a finite instrument can recover at declared tolerance \cite{devaney2004}. Confusing those readings produces the rank gaps and watt failures that later sections quantify \cite{pal2026}.
Rayleigh and Stratton treated such ranges as boundary-value laboratories long before digital near-field transforms became routine \cite{stratton1941}. Harrington's method-of-moments culture reinforced the same discipline: metal outline first, harmonic headline second \cite{harrington2001}.
Subsection~\ref{sec:ch272} states the formal export condition that this range scene presupposes \cite{stratton1941}.
\end{physicalvignette}
\subsection{Non-radiating current configurations}
The section introduction displayed balanced loops, counter-directed pairs, and below-cutoff
guides as laboratory signatures of $\ker\mathcal{R}_\omega$. Subsection~\ref{sec:ch272}
gives the definitive constructive account: admissible $\Jimp^{\mathrm{(NR)}}\in\ker\mathcal{R}_\omega$
with $\mathcal{S}_\omega[\Jimp^{\mathrm{(NR)}}]\neq\mathbf{0}$, developing Eq.
Eq.~\eqref{eq:ch2.nr-current-template} \cite{stratton1941}. Radiation integrals, Huygens surface forms
Eqs.~\eqref{eq:ch2.huygens-equivalent-integral}--\eqref{eq:ch2.huygens-tangential-match}, and
equivalent-source ambiguity Eq.~\eqref{eq:ch2.equivalent-source-ambiguity} from
Subsections~\ref{sec:ch251} and~\ref{sec:ch244} are presupposed; dyadic kernel tabulations from Subsection~\ref{sec:ch231} are not repeated \cite{stratton1941}. Evanescent overlap and
measurement nullity belong to Subsection~\ref{sec:ch273}; worked example examples appear in
Subsection~\ref{sec:ch2102} by reference \cite{harrington2001}.
Assumptions. Fix \(\omega>0\), compact support \(\Omega_s\), admissible impressed data
Eq.~\eqref{eq:ch2.source-admissibility}, charge continuity
Eq.~\eqref{eq:ch1.charge-continuity-harmonic}, outgoing radiation condition
Eqs.~\eqref{eq:ch2.outgoing-radiation-refinement}--\eqref{eq:ch2.outgoing-resolvent-unique},
and null space Eq.~\eqref{eq:ch2.radiation-nullspace-def}
\cite{stratton1941}. All configurations below are required to satisfy
\(\mathcal{D}_{J}[\Jimp^{\mathrm{(NR)}}]=1\) and \(\divg\Jimp^{\mathrm{(NR)}}=-\jj\omega\rho^{\mathrm{(NR)}}\)
with \(\rho^{\mathrm{(NR)}}\) compatible with boundary class
Eq.~\eqref{eq:ch1.outer-boundary-classes} \cite{stratton1941}. regime qualification
Eq.~\eqref{eq:ch2.regime-component-criterion} governs whether near-zone reactive content may be
reported as transport data \cite{harrington2001}. Non-radiating current (locked definition). A \emph{non-radiating current configuration} is an admissible impressed current
\begin{equation}
\label{eq:ch2.nr-current-def}
\Jimp^{\mathrm{(NR)}}\in\ker\mathcal{R}_\omega,
\qquad
\mathcal{S}_\omega[\Jimp^{\mathrm{(NR)}}]\neq\mathbf{0}\ \text{in general},
\end{equation}
with radiative component \(\mathbf{u}_{\mathrm{rad}}=\mathbf{0}\) and non-radiative component
\(\mathbf{u}_{\mathrm{nr}}=\mathcal{S}_\omega[\Jimp^{\mathrm{(NR)}}]\) from
Eqs.~\eqref{eq:ch2.radiative-component-def}--\eqref{eq:ch2.nonradiative-component-def}
\cite{stratton1941}. Equation~\eqref{eq:ch2.nr-current-def} is the Subsection
~\ref{sec:ch272} development of Eq.~\eqref{eq:ch2.nr-current-template} and
expresses Eq.~\eqref{eq:ch2.null-radiation} with explicit $\mathcal{D}_J=1$: far-zone export
vanishes through Eq.~\eqref{eq:ch2.radiation-integral-outgoing-upgrade}, while near-zone fields
may remain substantial kernel image on field space (Subsection~\ref{sec:ch271})
\cite{stratton1941}. Linearity Eq.~\eqref{eq:ch2.radiation-linearity} implies
admissible superpositions of non-radiating configurations remain non-radiating when charge
continuity is preserved \cite{stratton1941}.
Maxwell-closed constraints. Non-radiating configurations are not arbitrary curl fields. Impressed data must satisfy
\begin{equation}
\label{eq:ch2.nr-current-charge-continuous-law}
\Jimp^{\mathrm{(NR)}}\in L^{2}(\Omega_s)^{3},
\quad
\divg\Jimp^{\mathrm{(NR)}}=-\jj\omega\rho^{\mathrm{(NR)}},
\quad
\Jimp^{\mathrm{(NR)}}\ \text{BC-compatible on}\ \Gamma,
\end{equation}
as in Eq.~\eqref{eq:ch2.source-admissibility} \cite{stratton1941}. Equation
\eqref{eq:ch2.nr-current-admissibility-law} excludes discontinuous \(\Jimp^{\mathrm{(NR)}}\) that
violate Eq.~\eqref{eq:ch1.charge-continuity-harmonic} before Maxwell field specification is attempted
\cite{stratton1941}. Surface-equivalent realizations on \(\Gamma_{\mathrm{eq}}\) inherit trace
class constraints from Subsection~\ref{sec:ch211} \cite{harrington2001}. Membership in
\(\ker\mathcal{R}_\omega\) is an \emph{export} statement; Eq.~\eqref{eq:ch2.nr-current-charge-continuous-law}
is a \emph{source} statement---both are required \cite{stratton1941}.
Let \(\mathcal{C}\subset\Omega_s\) be a closed admissible circulation path enclosing negligible
magnetic flux linkage at \(\omega\), with oppositely directed arms \(\mathcal{C}_+\) and
\(\mathcal{C}_-\) carrying equal magnitude currents. Physically, picture a wire loop in the
\(xy\)-plane with counter-directed segments along \(\pm\hat{\mathbf{x}}\): each arm radiates a
leading \(1/r\) transverse field, but the net electric dipole moment vanishes when the
circulation closes. When far-zone transverse contributions from the two arms cancel under
outgoing propagation,
\begin{equation}
\label{eq:ch2.nr-loop-cancellation}
\int_{\mathcal{C}_+}\!\mathbf{J}_{\mathrm{loop}}\cdot d\boldsymbol{\ell}
\int_{\mathcal{C}_-}\!\mathbf{J}_{\mathrm{loop}}\cdot d\boldsymbol{\ell}
=0
\ \Longrightarrow\
\mathcal{R}_\omega[\Jimp^{(\mathrm{loop})}]=\mathbf{0},
\end{equation}
for the associated volume current \(\Jimp^{(\mathrm{loop})}\) obtained by compact support
regularization on \(\mathcal{C}\) \cite{stratton1941}. Equation
\eqref{eq:ch2.nr-loop-cancellation} is the dipole-moment cancellation law at the source level:
vanishing leading \(1/r\) transverse term in
Eq.~\eqref{eq:ch2.far-field-amplitude-integral} when the net moment vanishes
\cite{jackson1999}. Near-zone reactive fields from \(\mathcal{S}_\omega[\Jimp^{(\mathrm{loop})}]\)
need not vanish; Eq.~\eqref{eq:ch2.regime-component-criterion} still applies \cite{stratton1941}.
Two compact sources \(\Jimp^{(+)}\), \(\Jimp^{(-)}\) with supports \(\Omega_\pm\) separated by
distance \(d\ll\lambda_0\) can form a non-radiating pair when their outgoing far patterns
interfere destructively. A concrete arrangement is a pair of equal-strength electric dipoles
along \(\hat{\mathbf{z}}\), separated by \(d=\lambda_0/20\), driven in opposite phase so the
leading \(1/r\) terms cancel on \(S^{2}\) while reactive overlap remains in the gap:
\begin{equation}
\label{eq:ch2.nr-counterprop-pair}
\Jimp^{\mathrm{(NR,cp)}}=\Jimp^{(+)}+\Jimp^{(-)},
\qquad
\E_{0,+}(\theta,\phi)+\E_{0,-}(\theta,\phi)=\mathbf{0}
\ \text{on}\ S^{2},
\end{equation}
with \(\E_{0,\pm}\) far amplitudes from Eq.~\eqref{eq:ch2.far-field-amplitude-integral} for
each component \cite{stratton1941}. Equation
\eqref{eq:ch2.nr-counterprop-pair} realizes Eq.~\eqref{eq:ch2.nr-current-def} when phase and
polarization of \(\Jimp^{(\pm)}\) are tuned so \(\Pi_{\mathrm{rad}}\mathcal{S}_\omega[\Jimp^{(\pm)}]\)
cancel on the outgoing branch \cite{stratton1941}. The pair may store substantial reactive
energy in the overlap region through Eq.~\eqref{eq:ch2.component-power-budget} while
\(\mathcal{P}_{\mathrm{rad}}=0\) on qualified shells \cite{harrington2001}.
On a closed equivalent surface \(\Gamma_{\mathrm{eq}}\) enclosing radiating sources, tangential
non-radiating distributions satisfy Eq.~\eqref{eq:ch2.equivalent-source-ambiguity}:
\begin{equation}
\label{eq:ch2.nr-surface-tangential-null}
\mathbf{J}_{s,\mathrm{eq}}^{\mathrm{(NR)}}
\in
\ker\mathcal{R}_\omega\big|_{\Gamma_{\mathrm{eq}}},
\qquad
\mathbf{J}_{s,\mathrm{eq}}^{\mathrm{(NR)}}\cdot\hat{\mathbf{n}}=0
\ \text{on}\ \Gamma_{\mathrm{eq}},
\end{equation}
with \(\hat{\mathbf{n}}\) the outward normal \cite{stratton1941}. Equation
\eqref{eq:ch2.nr-surface-tangential-null} is the Subsection~\ref{sec:ch272} constructive realization of
Huygens equivalence ambiguity: exterior fields on \(\mathcal{F}_{\mathrm{SF}}\) are fixed by
exterior radiative uniqueness (Subsection~\ref{sec:ch243}), but equivalent surface currents are unique
only up to additions in Eq.~\eqref{eq:ch2.nr-surface-tangential-null}
\cite{stratton1941}. Integrals Eqs.~\eqref{eq:ch2.huygens-equivalent-integral}--\eqref{eq:ch2.huygens-tangential-match}
realize the same \(\mathcal{S}_\omega\) image for equivalent and non-radiating tangential
perturbations \cite{harrington2001}.
When radiation eigenpairs with \(\lambda_n=0\) are declared in
Eq.~\eqref{eq:ch2.radiation-eigenpair-def},
\begin{equation}
\label{eq:ch2.nr-modal-eigenrealization}
\Jimp^{\mathrm{(NR,mod)}}
=
\sum_{n\in\mathcal{N}_{\mathrm{null}}}
\alpha_n\,\Jimp^{(n)},
\qquad
\mathcal{R}_\omega[\Jimp^{\mathrm{(NR,mod)}}]=\mathbf{0},
\end{equation}
for coefficients \(\alpha_n\in\mathbb{C}\) such that the superposition satisfies
Eq.~\eqref{eq:ch2.nr-current-admissibility-law} \cite{stratton1941}. Equation
\eqref{eq:ch2.nr-modal-eigenrealization} specializes Eq.~\eqref{eq:ch2.modal-null-sector-reduction}
from Subsection~\ref{sec:ch271} to constructive sources: each \(\Jimp^{(n)}\) with
\(\lambda_n=0\) is a modal non-radiating generator when admissible
\cite{stratton1941}. On exterior domains \(\mathcal{N}_{\mathrm{null}}\) need not span all of
\(\ker\mathcal{R}_\omega\) Eq.~\eqref{eq:ch2.nullspace-dimensional-classification}
\cite{harrington2001}.
\begin{figure}[t]
\centering
\ChOneFigBlock{%
\begin{tikzpicture}[ch1fig, node distance=5mm and 7mm]
 \node[ch1box] (loop) {Loop\\cancellation};
 \node[ch1box, right=7mm of loop] (cp) {Counter-\\prop pair};
 \node[ch1box, right=7mm of cp] (surf) {Surface\\tangential};
 \node[ch1box, right=7mm of surf] (mod) {Modal\\$\lambda_n=0$};
 \node[ch1box, below=5mm of cp] (ker) {$\ker\mathcal{R}_\omega$};
 \node[ch1box, below=5mm of loop] (j) {$\Jimp^{\mathrm{(NR)}}$};
 \draw[ch1flow] (loop.east) -- (loop.east |- ker.west);
 \draw[ch1flow] (loop.east |- ker.west) -- (ker.west);
 \draw[ch1flow] (cp.south) -- (cp.south |- ker.north);
 \draw[ch1flow] (cp.south |- ker.north) -- (ker.north);
 \draw[ch1flow] (surf.east) -- (surf.east |- ker.west);
 \draw[ch1flow] (surf.east |- ker.west) -- (ker.west);
 \draw[ch1flow] (mod.east) -- (mod.east |- ker.west);
 \draw[ch1flow] (mod.east |- ker.west) -- (ker.west);
 \draw[ch1flow] (j.south) -- (j.south |- loop.north);
 \draw[ch1flow] (j.south |- loop.north) -- (loop.north);
 \draw[ch1flow] (j.east) -- (j.east |- cp.west);
 \draw[ch1flow] (j.east |- cp.west) -- (cp.west);
 \draw[ch1flow] (j.east) -- (j.east |- surf.west);
 \draw[ch1flow] (j.east |- surf.west) -- (surf.west);
 \draw[ch1flow] (j.east) -- (j.east |- mod.west);
 \draw[ch1flow] (j.east |- mod.west) -- (mod.west);
\end{tikzpicture}%
}
\caption{Subsection~\ref{sec:ch272} canonical non-radiating families: loop cancellation
Eq.~\eqref{eq:ch2.nr-loop-cancellation}, counter-propagating pair
Eq.~\eqref{eq:ch2.nr-counterprop-pair}, surface tangential null
Eq.~\eqref{eq:ch2.nr-surface-tangential-null}, and modal eigenrealization
Eq.~\eqref{eq:ch2.nr-modal-eigenrealization}, all mapping to
\(\ker\mathcal{R}_\omega\).}
\label{fig:ch2.nr-current-configurations}
\end{figure}
The four upper nodes are constructive \emph{families}, not interchangeable templates. Loop cancellation annihilates the leading $1/r$ transverse moment; counter-propagating pairs cancel on $S^{2}$ through phase interference; surface tangential nulls exploit Huygens equivalence ambiguity; modal $\lambda_n=0$ generators realize the spectral slice Eq.~\eqref{eq:ch2.modal-null-sector-reduction}. All four converge on $\ker\mathcal{R}_\omega$, yet each leaves distinct reactive near fields $\mathcal{S}_\omega[\Jimp^{\mathrm{(NR)}}]\neq\mathbf{0}$---the schematic warns that near-zone scans cannot distinguish families once far patterns vanish \cite{stratton1941}. Figure~\ref{fig:ch2.nr-current-configurations} orders constructive families before
evanescent and measurement extensions in Subsection~\ref{sec:ch273}. Equivalent-current solvers on \(\Gamma_{\mathrm{eq}}\) seek \(\mathbf{J}_{\mathrm{eq}}\) matching
near or far data \cite{stratton1941}. Canonical non-radiating configurations (Subsection~\ref{sec:ch272}) and Eq.~\eqref{eq:ch2.nr-surface-tangential-null}
show why such inversions are ill-posed even when forward fields are unique: any solution may
be shifted by tangential \(\mathbf{J}_{s,\mathrm{eq}}^{\mathrm{(NR)}}\) without changing
\(\mathcal{R}_\omega[\Jimp]\) \cite{stratton1941}. Tikhonov regularization
Eq.~\eqref{eq:ch2.tikhonov-source-reconstruction} selects a particular representative but does
not eliminate the equivalence class \cite{harrington2001}. Subsection~\ref{sec:ch273} separates
this source nullity from finite-rank observation nullity
Eq.~\eqref{eq:ch2.measurement-nullspace-composition} \cite{stratton1941}. Subsection~\ref{sec:ch272} assumes linear Maxwell theory at fixed
\(\omega\), outgoing radiation condition, and admissible sources Eq.~\eqref{eq:ch2.nr-current-admissibility-law}
\cite{stratton1941}. Lossy media modify loop and pair conditions through complex
\(k(\mathbf{r})\) without promoting evanescent storage to export
Eq.~\eqref{eq:ch2.rad-null-on-evanescent} \cite{stratton1941}. Canonical families generate
\emph{explicit} members of \(\ker\mathcal{R}_\omega\); they do not exhaust
\(\dim\ker\mathcal{R}_\omega=\infty\) on exterior domains
Eq.~\eqref{eq:ch2.nullspace-dimensional-classification} \cite{harrington2001}.
regime qualification Eq.~\eqref{eq:ch2.regime-operator-screen} remains mandatory: large near-zone
\(|\E|\) from \(\mathcal{S}_\omega[\Jimp^{\mathrm{(NR)}}]\) is not export without
\(\Pi_{\mathrm{rad}}\) qualification \cite{stratton1941}. \phantomsection\label{sec:ch273}
% PHYSICAL-MEMORY:sec:ch273
\begin{physicalvignette}[Monopole on a vehicle roof ground plane]
Consider Monopole on a vehicle roof ground plane. Before any coefficient table is printed, technicians record frequency, cable loss, and the mechanical outline that fixes electrical size. The first walk-through of the range is always tactile: flange bolts, probe standoff, turntable boresight, and the wattmeter reference on the receive chain \cite{hansen1988}.
% PHYSICAL-MEMORY:sec:ch273-b2
\begin{physicalvignette}[Monopole on a vehicle roof ground plane]
A second bench observation matters on Inverted-F antenna on a PCB edge. When technicians swap only the receive probe while holding the transmit chain fixed, reciprocity checks whether the reported coupling is property of the structure or artifact of a mismatched flux gauge \cite{stratton1941}. Failed reciprocity usually traces to a near-field tap, a radome standing-wave, or a probe polarization basis that was never aligned with the transmit standard \cite{hansen1988}.
Rayleigh and Stratton treated such ranges as boundary-value laboratories long before digital near-field transforms became routine \cite{stratton1941}. Harrington's method-of-moments culture reinforced the same discipline: metal outline first, harmonic headline second \cite{harrington2001}.
The measurement logic above is made precise in Subsection~\ref{sec:ch273} \cite{stratton1941}.
\end{physicalvignette}
On this hardware the near zone is visually impressive---standing reactive structure, rapid phase rotation, sometimes a visible VSWR ripple at the feed---yet the export question is measured only after outgoing selection on a large sphere \cite{harrington2001}. What the camera sees in the chamber is not automatically what the far-zone pattern claims \cite{stratton1941}.
Engineers therefore learn to separate three readings: what the feed can synthesize on declared support, what propagating content survives on the radiative range, and what a finite instrument can recover at declared tolerance \cite{devaney2004}. Confusing those readings produces the rank gaps and watt failures that later sections quantify \cite{pal2026}.
Rayleigh and Stratton treated such ranges as boundary-value laboratories long before digital near-field transforms became routine \cite{stratton1941}. Harrington's method-of-moments culture reinforced the same discipline: metal outline first, harmonic headline second \cite{harrington2001}.
Subsection~\ref{sec:ch273} states the formal export condition that this range scene presupposes \cite{stratton1941}.
\end{physicalvignette}
\subsection{Evanescent overlap and non-reconstructability}
Subsections~\ref{sec:ch271}--\ref{sec:ch272} fixed $\ker\mathcal{R}_\omega$ and canonical
non-radiating impressed currents Eq.~\eqref{eq:ch2.nr-current-def}. Evanescent sheets are a
distinct field-side deficit: reactive storage that never passes $\Pi_{\mathrm{rad}}$ even when
the source radiates. Subsection~\ref{sec:ch273} supplies the definitive account of evanescent
subspaces, their overlap with non-radiating source content, and reconstructability limits under
finite observation rank---developing Eq.
Eq.~\eqref{eq:ch2.evanescent-nonreconstruct} from Section~\ref{sec:ch27}
\cite{stratton1941}. Propagation projectors
Eqs.~\eqref{eq:ch2.spectral-projectors}--\eqref{eq:ch2.rad-null-on-evanescent} from
Subsection~\ref{sec:ch224}, non-propagating split
Eq.~\eqref{eq:ch2.nonpropagating-component-def} from Subsection~\ref{sec:ch254}, measurement
null-space extension measurement null-space extension (Subsection~\ref{sec:ch271}) from
Subsection~\ref{sec:ch271}, and rank cap Eq.~\eqref{eq:ch2.reconstruction-rank-cap} from
Subsection~\ref{sec:ch244} are presupposed; projector algebra and inverse-source templates are
not re-proved \cite{stratton1941}. Angular-spectrum consequences belong to Subsection
~\ref{sec:ch274}; plane-wave filters to Section~\ref{sec:ch29} \cite{harrington2001}.
Assumptions. Fix \(\omega>0\), admissible \(\mathcal{J}_\Omega\) Eq.~\eqref{eq:ch2.source-admissibility},
outgoing radiation condition Eqs.~\eqref{eq:ch2.outgoing-radiation-refinement}--\eqref{eq:ch2.outgoing-resolvent-unique},
field pairs in \(\mathcal{F}_{\mathrm{SF}}\subset\mathcal{F}_\Omega\), and spectral completion
supporting Eqs.~\eqref{eq:ch2.spectral-projectors}--\eqref{eq:ch2.spectral-projector-algebra}
\cite{stratton1941}. Observation map
\(\mathcal{M}_\omega=\mathcal{O}_{\mathrm{meas}}\circ\Lambda_{\mathrm{rad}}\) from
Eq.~\eqref{eq:ch2.reconstruction-observation-map} is finite rank with
\(\mathrm{rank}(\mathcal{O}_{\mathrm{meas}})=N\) Eq.~\eqref{eq:ch2.finite-rank-observation}
\cite{stratton1941}. regime qualification Eq.~\eqref{eq:ch2.regime-operator-screen} governs transport
reports on any recovered radiative component \cite{harrington2001}. On the spectral completion \(\widehat{\mathcal{F}}_\Omega\), the \emph{evanescent subspace} is
\begin{equation}
\label{eq:ch2.evanescent-subspace-def}
\mathcal{F}_{\mathrm{ev}}
=
\left\{
\mathbf{u}\in\widehat{\mathcal{F}}_\Omega:
\Pi_{\mathrm{ev}}\mathbf{u}=\mathbf{u}
\right\},
\end{equation}
with \(\Pi_{\mathrm{ev}}\) from Eq.~\eqref{eq:ch2.spectral-projectors}
\cite{stratton1941}. Equation~\eqref{eq:ch2.evanescent-subspace-def} is the
Subsection~\ref{sec:ch273} field-side label for modes with imaginary longitudinal
wavenumber Eq.~\eqref{eq:ch2.longitudinal-wavenumber-split} and decay rate
Eq.~\eqref{eq:ch2.evanescent-decay-rate} \cite{stratton1941}. The companion \emph{non-propagating} subspace combines evanescent and
below-cutoff content through Eq.~\eqref{eq:ch2.nonpropagating-component-def}
\cite{harrington2001}. Evanescent sheets carry reactive structure bound to apertures; they decay
exponentially with scan height and satisfy Eq.~\eqref{eq:ch2.rad-null-on-evanescent} without
re-derivation here \cite{stratton1941}.
For admissible \(\Jimp\in\mathcal{J}_\Omega\), the evanescent projector \(\Pi_{\mathrm{ev}}\)
selects field content whose lateral wavenumber exceeds the vacuum propagation circle
\(k_t>k_0\): such modes decay exponentially with distance from the source and carry reactive
storage rather than qualified export. Define the evanescent field content
\begin{equation}
\label{eq:ch2.evanescent-source-overlap}
\mathbf{u}_{\mathrm{ev}}
=
\Pi_{\mathrm{ev}}\,\mathcal{S}_\omega[\Jimp],
\qquad
\|\mathbf{u}_{\mathrm{ev}}\|_{\mathcal{E}}
=
\|\Pi_{\mathrm{ev}}\mathcal{S}_\omega[\Jimp]\|_{\mathcal{E}},
\end{equation}
with energy norm Eq.~\eqref{eq:ch2.field-energy-norm} \cite{stratton1941}.
Equation~\eqref{eq:ch2.evanescent-source-overlap} measures reactive storage excited by
\(\Jimp\) that never passes through \(\Pi_{\mathrm{rad}}\) because
\begin{equation}
\label{eq:ch2.evanescent-nonexport-law}
\Pi_{\mathrm{rad}}\,\mathbf{u}_{\mathrm{ev}}=\mathbf{0},
\qquad
\mathcal{R}_\omega[\Jimp]\ \text{depends only on}\ \Pi_{\mathrm{prop}}\mathcal{S}_\omega[\Jimp]
\ \text{after outgoing selection},
\end{equation}
by Eqs.~\eqref{eq:ch2.rad-null-on-evanescent} and~\eqref{eq:ch2.rad-projection-composition}
\cite{stratton1941}. Equation~\eqref{eq:ch2.evanescent-nonexport-law} upgrades the second
line of Eq.~\eqref{eq:ch2.evanescent-nonreconstruct}: evanescent overlap is
field-side non-export, distinct from source-side \(\Jimp\in\ker\mathcal{R}_\omega\)
Eq.~\eqref{eq:ch2.radiation-nullspace-def} \cite{harrington2001}. A source may satisfy
\(\|\mathbf{u}_{\mathrm{ev}}\|\gg\|\mathcal{R}_\omega[\Jimp]\|\) while still radiating
\cite{stratton1941}.
Eq.~\eqref{eq:ch2.evanescent-nonreconstruct} separates export annihilation from
finite-rank readout loss. Source nullity \(\ker\mathcal{R}_\omega\) is strict export failure;
measurement nullity \(\ker\mathcal{M}_\omega\) additionally loses directions that radiate but
the aperture cannot resolve \cite{stratton1941}:
\begin{equation}
\label{eq:ch2.radiation-vs-measurement-nullity-split}
\underbrace{\ker\mathcal{R}_\omega}_{\text{no export}}
\ \subseteq\
\underbrace{\ker\mathcal{M}_\omega}_{\text{no finite readout}},
\qquad
\ker\mathcal{M}_\omega
\supsetneq
\ker\mathcal{R}_\omega
\ \text{when}\ 
\ker\mathcal{O}_{\mathrm{meas}}\big|_{\operatorname{ran}\mathcal{R}_\omega}\neq\{\mathbf{0}\},
\end{equation}
with strict inclusion typical for finite apertures
Eq.~\eqref{eq:ch2.observation-null-space}
\begin{equation}
\label{eq:ch2.measurement-nullspace-composition}
\ker\mathcal{M}_\omega
=
\ker\mathcal{R}_\omega
\ \cup\
\mathcal{R}_\omega^{-1}\!\big(\ker\mathcal{O}_{\mathrm{meas}}\big|_{\operatorname{ran}\mathcal{R}_\omega}\big),
\end{equation} \cite{stratton1941}. Equation
\eqref{eq:ch2.radiation-vs-measurement-nullity-split} expresses measurement null-space extension (Subsection~\ref{sec:ch271}) and
Eq.~\eqref{eq:ch2.measurement-nullspace-composition} without re-proof: evanescent storage
contributes to the first factor only when it annihilates export; finite rank contributes to
the second factor even when export is nonzero \cite{stratton1941}. Tikhonov reconstruction
Eq.~\eqref{eq:ch2.tikhonov-source-reconstruction} stabilizes ill-conditioned inversions on
\(\operatorname{ran}\mathcal{M}_\omega^{\dagger}\) but does not shrink
\(\ker\mathcal{M}_\omega\) \cite{harrington2001}.
Define the \emph{non-reconstructable set} at observation rank \(N\):
\begin{equation}
\label{eq:ch2.nonreconstruct-observation-kernel}
\mathcal{K}_{\mathrm{nrec}}(N)
=
\left\{
\Jimp\in\mathcal{J}_\Omega:
\mathcal{M}_\omega[\Jimp]=\mathbf{0}
\right\}
=
\ker\mathcal{M}_\omega,
\end{equation}
equivalently the kernel of Eq.~\eqref{eq:ch2.inverse-source-equation} for generic
\(\mathbf{y}\) \cite{stratton1941}. Equation
\eqref{eq:ch2.nonreconstruct-observation-kernel} is the Subsection~\ref{sec:ch273} constructive account of
reconstructability label: sources in \(\mathcal{K}_{\mathrm{nrec}}(N)\) are invisible to the
declared measurement bank regardless of near-zone \(|\E|\) \cite{stratton1941}. Rank bound
\begin{equation}
\label{eq:ch2.reconstructability-rank-law}
\dim\operatorname{ran}\mathcal{M}_\omega
\le
N
<
\infty
\ \Longrightarrow\
\dim\mathcal{K}_{\mathrm{nrec}}(N)=\infty
\ \text{on typical open domains},
\end{equation}
follows from Eq.~\eqref{eq:ch2.reconstruction-rank-cap} and infinite-dimensional
\(\mathcal{J}_\Omega\) \cite{stratton1941}. Geometric sector constraints
Eq.~\eqref{eq:ch2.accessible-modal-sector} and composition law
Eq.~\eqref{eq:ch2.constraint-composition-law} further restrict recoverable coordinates
independently of Eq.~\eqref{eq:ch2.nonreconstruct-observation-kernel}
\cite{stratton1941}.
\begin{figure}[t]
\centering
\ChOneFigBlock{%
\begin{tikzpicture}[ch1fig, node distance=5mm and 7mm]
 \node[ch1box] (j) {$\mathcal{J}_\Omega$};
 \node[ch1box, right=7mm of j] (s) {$\mathcal{S}_\omega$};
 \node[ch1box, right=7mm of s] (ev) {$\Pi_{\mathrm{ev}}$};
 \node[ch1box, right=7mm of ev] (r) {$\mathcal{R}_\omega$};
 \node[ch1box, right=7mm of r] (m) {$\mathcal{M}_\omega$};
 \node[ch1box, below=5mm of ev] (uev) {$\mathbf{u}_{\mathrm{ev}}$};
 \node[ch1box, below=5mm of m] (kn) {$\mathcal{K}_{\mathrm{nrec}}$};
 \draw[ch1flow] (j) -- (s) -- (ev) -- (r) -- (m);
 \draw[ch1flow] (ev.south) -- (ev.south |- uev.north);
 \draw[ch1flow] (ev.south |- uev.north) -- (uev.north);
 \draw[ch1flow] (m.south) -- (m.south |- kn.north);
 \draw[ch1flow] (m.south |- kn.north) -- (kn.north);
\end{tikzpicture}%
}
\caption{Subsection~\ref{sec:ch273} measurement sequence: evanescent overlap
Eq.~\eqref{eq:ch2.evanescent-source-overlap} never exports through
Eq.~\eqref{eq:ch2.evanescent-nonexport-law}; finite-rank observation maps
\(\mathcal{J}_\Omega\) to \(\mathcal{K}_{\mathrm{nrec}}(N)\)
Eq.~\eqref{eq:ch2.nonreconstruct-observation-kernel}.}
\label{fig:ch2.evanescent-nonreconstruct-measurement chain}
\end{figure}
The measurement sequence orders radiation selection, evanescent overlap, and finite-rank observation. $\Pi_{\mathrm{ev}}$ content may dominate $|\E|$ on inner shells yet satisfies Eq.~\eqref{eq:ch2.evanescent-nonexport-law} on qualified flux meters. $\mathcal{O}_{\mathrm{meas}}$ with rank $N$ collapses recoverable modes into $\mathcal{K}_{\mathrm{nrec}}(N)$: export can exist while the aperture cannot resolve it. Dashed branches mark observation deficits distinct from source nullity $\ker\mathcal{R}_\omega$ \cite{stratton1941}. \phantomsection\label{thm:ch2.evanescent-nonreconstructability}
Let \(\Jimp\in\mathcal{J}_\Omega\) be admissible with evanescent content
Eq.~\eqref{eq:ch2.evanescent-source-overlap} and observation map
\(\mathcal{M}_\omega\) of rank \(N\). Then:
\begin{enumerate}
\item[(i)] \(\Pi_{\mathrm{rad}}\mathbf{u}_{\mathrm{ev}}=\mathbf{0}\) and
\(\mathbf{u}_{\mathrm{ev}}\) carries no qualified far-zone export
Eq.~\eqref{eq:ch2.evanescent-nonexport-law}.
\item[(ii)] \(\ker\mathcal{R}_\omega\subseteq\ker\mathcal{M}_\omega=\mathcal{K}_{\mathrm{nrec}}(N)\).
\item[(iii)] Near-field data on a finite enclosing surface reduce but do not in general
eliminate \(\mathcal{K}_{\mathrm{nrec}}(N)\): if \(\mathcal{M}_{\mathrm{near}}\) samples
\(\mathcal{S}_\omega[\Jimp]\) on \(\Gamma_{\mathrm{enc}}\), then
\begin{equation}
\label{eq:ch2.near-field-nullity-reduction}
\ker\mathcal{M}_{\mathrm{near}}
\subsetneq
\ker\mathcal{M}_\omega
\ \text{when}\ \Gamma_{\mathrm{enc}}\ \text{resolves more modes than far-zone aperture},
\end{equation}
while \(\ker\mathcal{R}_\omega\) is unchanged \cite{stratton1941}.
\item[(iv)] Reconstructability budget:
\begin{equation}
\label{eq:ch2.reconstructability-budget}
\mathcal{N}_{\mathrm{reportable}}
\subseteq
\mathcal{N}_{\mathrm{acc}}
\cap
\{n:\ \text{resolvable by}\ \mathcal{M}_\omega\}
\cap
\{n:\ \lambda_n\neq 0\},
\end{equation}
with \(\mathcal{N}_{\mathrm{acc}}\) from Eq.~\eqref{eq:ch2.accessible-modal-sector} and modal
indices from Subsection~\ref{sec:ch261} \cite{stratton1941}.
\end{enumerate}
Part (i) is Eq.~\eqref{eq:ch2.rad-null-on-evanescent} \cite{stratton1941}. Part (ii) follows
from measurement null-space extension (Subsection~\ref{sec:ch271}) and definition
Eq.~\eqref{eq:ch2.nonreconstruct-observation-kernel} \cite{harrington2001}. Part (iii) uses
near- versus far-field sampling logic from Subsection~\ref{sec:ch244}: \(\mathcal{S}_\omega\) is
single-valued on outgoing fields while \(\mathcal{R}_\omega\) remains non-injective
\cite{stratton1941}. Part (iv) composes
Eq.~\eqref{eq:ch2.constraint-composition-law} with rank cap
Eq.~\eqref{eq:ch2.reconstruction-rank-cap} \cite{stratton1941}.
\phantomsection\label{prop:ch2.evanescent-never-exports}
For all admissible \(\Jimp\),
\begin{equation}
\label{eq:ch2.evanescent-radiation-annihilation}
\mathcal{R}_\omega[\Jimp]
=
\Pi_{\mathrm{rad}}\,\mathcal{S}_\omega[\Jimp]
=
\Pi_{\mathrm{rad}}\,\Pi_{\mathrm{prop}}\,\mathcal{S}_\omega[\Jimp],
\end{equation}
and \(\Pi_{\mathrm{rad}}\Pi_{\mathrm{ev}}=\mathbf{0}\) on \(\mathcal{F}_\Omega\)
\cite{stratton1941}. Equation~\eqref{eq:ch2.evanescent-radiation-annihilation}
implies \(\mathbf{u}_{\mathrm{ev}}\) from Eq.~\eqref{eq:ch2.evanescent-source-overlap} does
not alter \(\mathcal{R}_\omega[\Jimp]\); adding evanescent excitation to an otherwise
radiating source changes \(\mathbf{u}_{\mathrm{nr}}\) and local power
Eq.~\eqref{eq:ch2.component-power-budget} without changing far patterns when export is
unchanged \cite{stratton1941}. regime qualification Eq.~\eqref{eq:ch2.regime-component-criterion} still
forbids reporting \(\mathbf{u}_{\mathrm{ev}}\) as transport \cite{harrington2001}. Substitute Eqs.~\eqref{eq:ch2.rad-null-on-evanescent} and~\eqref{eq:ch2.rad-projection-composition}
into \(\mathcal{R}_\omega=\Pi_{\mathrm{rad}}\circ\mathcal{S}_\omega\)
\cite{stratton1941}. Near-zone energy follows from
Eq.~\eqref{eq:ch2.nonpropagating-component-def} \cite{harrington2001}.
\phantomsection\label{prop:ch2.near-field-nullity-reduction}
Let \(\Gamma_{\mathrm{enc}}\) be a closed surface enclosing \(\Omega_s\) and let
\(\mathcal{M}_{\mathrm{near}}\) sample tangential \(\E\), \(\H\) on \(\Gamma_{\mathrm{enc}}\).
Then \(\ker\mathcal{M}_{\mathrm{near}}\subseteq\ker\mathcal{M}_\omega\) when far-zone
observations are a subset of near-field functionals, but
\begin{equation}
\label{eq:ch2.near-field-kernel-inclusion}
\ker\mathcal{R}_\omega
\subseteq
\ker\mathcal{M}_{\mathrm{near}}
\subseteq
\ker\mathcal{M}_\omega
\end{equation}
with strict inclusions possible at each step \cite{stratton1941}. Equation
\eqref{eq:ch2.near-field-kernel-inclusion} expresses Subsection~\ref{sec:ch244} near- versus
far-field data without re-opening uniqueness exterior radiative uniqueness (Subsection~\ref{sec:ch243}) \cite{stratton1941}. Outgoing radiation condition remains
mandatory: identical near traces on a finite surface do not imply identical far zones without
exterior selection \cite{harrington2001}.
If \(\mathcal{M}_\omega[\Jimp]=\mathbf{0}\), then far functionals vanish; near samples may
still vanish on \(\Gamma_{\mathrm{enc}}\) for non-radiating additions in
\(\ker\mathcal{R}_\omega\) \cite{stratton1941}. Additional near resolution shrinks the
observation null space on \(\operatorname{ran}\mathcal{R}_\omega\) but leaves \(\ker\mathcal{R}_\omega\) unchanged Eq.~\eqref{eq:ch2.radiation-nullspace-def}
\cite{harrington2001}.
When Eq.~\eqref{eq:ch2.cutoff-condition} fails, modal content lies in
\(\Pi_{\mathrm{non}}\mathcal{S}_\omega[\Jimp]\) or \(\Pi_{\mathrm{ev}}\mathcal{S}_\omega[\Jimp]\)
\cite{stratton1941}. Such storage satisfies
Eq.~\eqref{eq:ch2.evanescent-nonexport-law} and may dominate
Eq.~\eqref{eq:ch2.component-power-budget} on sub-wavelength shells while
\(\mathcal{P}_{\mathrm{rad}}=0\) \cite{stratton1941}. Applying far-zone pattern formulas to
below-cutoff modes violates Eq.~\eqref{eq:ch2.regime-operator-screen} as in the TE\(_{10}\)
illustration referenced from Section~\ref{sec:ch01.1.2} \cite{harrington2001}. Subsection
~\ref{sec:ch273} records operator consequences; guide nomenclature is not re-tabulated
\cite{stratton1941}. Non-radiating configurations Eq.~\eqref{eq:ch2.nr-current-def} may carry large
\(\|\mathbf{u}_{\mathrm{ev}}\|\) when \(\mathcal{S}_\omega[\Jimp^{\mathrm{(NR)}}]\) is
evanescent-dominated \cite{stratton1941}. Conversely, large
\(\|\mathbf{u}_{\mathrm{ev}}\|\) does not imply \(\Jimp\in\ker\mathcal{R}_\omega\): a radiating
source may store substantial evanescent energy near \(\Omega_s\) while exporting through
\(\Pi_{\mathrm{prop}}\) \cite{stratton1941}. Subsections~\ref{sec:ch272} and~\ref{sec:ch273}
address orthogonal deficits---source nullity versus reactive storage versus finite-rank readout
\cite{harrington2001}.
Subsection~\ref{sec:ch273} assumes linear Maxwell theory at fixed
\(\omega\), outgoing radiation condition, and declared finite-rank \(\mathcal{O}_{\mathrm{meas}}\)
\cite{stratton1941}. Loss modifies evanescent decay
Eq.~\eqref{eq:ch2.evanescent-decay-rate} without promoting \(\Pi_{\mathrm{ev}}\) content to
\(\Pi_{\mathrm{rad}}\) export \cite{stratton1941}. Dispersive or broadband superposition
requires Subsection~\ref{sec:ch223}; harmonic evanescent labels here are fixed-frequency slices
\cite{harrington2001}. Subsection~\ref{sec:ch274} translates non-reconstructability to
continuous angular spectra Eq.~\eqref{eq:ch2.exterior-continuous-spectral-integral}
\cite{stratton1941}.
\subsection{Angular-spectrum consequences of null spaces}
\label{sec:ch274}
Far-zone angular amplitudes $a(\widehat{\mathbf{k}})$ are what compact ranges and
synthesis codes ultimately report. Subsections~\ref{sec:ch271}--\ref{sec:ch273} fixed
$\ker\mathcal{R}_\omega$, canonical non-radiating impressed currents, and evanescent
non-reconstructability; Subsection~\ref{sec:ch274} records how those source-side deficits
propagate to the continuous angular measure
Eq.~\eqref{eq:ch2.exterior-continuous-spectral-integral} without re-opening plane-wave operators
from Section~\ref{sec:ch29} \cite{stratton1941}.
% PHYSICAL-MEMORY:sec:ch274
% PHYSICAL-MEMORY:sec:ch274-b2
\begin{physicalvignette}[Vivaldi notch array on a printed circuit board]
A second bench observation matters on Horn-lens combination on a millimetre-wave bench. When technicians swap only the receive probe while holding the transmit chain fixed, reciprocity checks whether the reported coupling is property of the structure or artifact of a mismatched flux gauge \cite{stratton1941}. Failed reciprocity usually traces to a near-field tap, a radome standing-wave, or a probe polarization basis that was never aligned with the transmit standard \cite{hansen1988}.
Rayleigh and Stratton treated such ranges as boundary-value laboratories long before digital near-field transforms became routine \cite{stratton1941}. Harrington's method-of-moments culture reinforced the same discipline: metal outline first, harmonic headline second \cite{harrington2001}.
The measurement logic above is made precise in Subsection~\ref{sec:ch274} \cite{stratton1941}.
\end{physicalvignette}
\begin{physicalvignette}[Vivaldi notch array on a printed circuit board]
Consider Vivaldi notch array on a printed circuit board. Before any coefficient table is printed, technicians record frequency, cable loss, and the mechanical outline that fixes electrical size. The first walk-through of the range is always tactile: flange bolts, probe standoff, turntable boresight, and the wattmeter reference on the receive chain \cite{hansen1988}.
On this hardware the near zone is visually impressive---standing reactive structure, rapid phase rotation, sometimes a visible VSWR ripple at the feed---yet the export question is measured only after outgoing selection on a large sphere \cite{harrington2001}. What the camera sees in the chamber is not automatically what the far-zone pattern claims \cite{stratton1941}.
Engineers therefore learn to separate three readings: what the feed can synthesize on declared support, what propagating content survives on the radiative range, and what a finite instrument can recover at declared tolerance \cite{devaney2004}. Confusing those readings produces the rank gaps and watt failures that later sections quantify \cite{pal2026}.
Rayleigh and Stratton treated such ranges as boundary-value laboratories long before digital near-field transforms became routine \cite{stratton1941}. Harrington's method-of-moments culture reinforced the same discipline: metal outline first, harmonic headline second \cite{harrington2001}.
Subsection~\ref{sec:ch274} states the formal export condition that this range scene presupposes \cite{stratton1941}.
\end{physicalvignette}
Assumptions. Fix \(\omega>0\), outgoing radiation condition, admissible \(\mathcal{J}_\Omega\), and far amplitude
\(\E_0(\theta,\phi)\) from Eq.~\eqref{eq:ch2.far-field-amplitude-integral} when \(\mathcal{D}_{J,\mathrm{reg}}=1\) \cite{stratton1941}. Continuous angular
coefficients \(a_\sigma(\widehat{\mathbf{k}})\) are defined through the rigged limit of
Eq.~\eqref{eq:ch2.exterior-continuous-spectral-integral} compatible with modal amplitudes
Eq.~\eqref{eq:ch2.source-modal-amplitude} \cite{stratton1941}. Finite-rank observation
\(\mathcal{M}_\omega\) and non-reconstructable set \(\mathcal{K}_{\mathrm{nrec}}(N)\) from
Eq.~\eqref{eq:ch2.nonreconstruct-observation-kernel} are presupposed
\cite{harrington2001}. regime qualification Eq.~\eqref{eq:ch2.regime-operator-screen} governs transport
claims on any retained angular content \cite{stratton1941}.
Let $a(\widehat{\mathbf{k}})$ denote the scalar angular spectrum obtained from
Eq.~\eqref{eq:ch2.exterior-continuous-spectral-integral} on the outgoing transverse sheet
\cite{stratton1941}. For \(\Jimp\in\ker\mathcal{R}_\omega\),
\begin{equation}
\label{eq:ch2.angular-spectrum-nullspace-law}
a(\widehat{\mathbf{k}})=0
\quad\text{for all}\ \widehat{\mathbf{k}}\in S^{2}
\ \text{on the exporting sheet},
\end{equation}
equivalently \(\E_0(\theta,\phi)=\mathbf{0}\) in Eq.~\eqref{eq:ch2.spectral-far-pattern-expansion}
\cite{stratton1941}. Equation~\eqref{eq:ch2.angular-spectrum-nullspace-law} is the
Subsection~\ref{sec:ch274} development of the first clause of Eq.
Eq.~\eqref{eq:ch2.angular-spectrum-null}: null-space membership annihilates the entire
far-zone angular measure, not merely selected harmonics
\cite{harrington2001}. Near-zone reactive content in
\(\Pi_{\mathrm{ev}}\mathcal{S}_\omega[\Jimp]\) may remain nonzero without contributing to
\(a(\widehat{\mathbf{k}})\) Eq.~\eqref{eq:ch2.evanescent-nonexport-law}
\cite{stratton1941}.
Suppressed and accessible angular sheets. Define the \emph{suppressed angular sheet}
\begin{equation}
\label{eq:ch2.suppressed-angular-sheet-def}
\mathcal{K}_{\mathrm{sup}}
=
S^{2}\setminus\mathcal{K}_{\mathrm{acc}},
\qquad
\mathcal{K}_{\mathrm{acc}}
\ \text{from Eq.~\eqref{eq:ch2.angular-accessibility-sector}},
\end{equation}
as directions annihilated by geometric, support, or parity constraints before observation
\cite{stratton1941}. Equation~\eqref{eq:ch2.suppressed-angular-sheet-def} expresses
Eq.~\eqref{eq:ch2.angular-spectrum-null}: \(a(\widehat{\mathbf{k}})=0\) on
\(\mathcal{K}_{\mathrm{sup}}\) does not imply \(\Jimp\in\ker\mathcal{R}_\omega\), because
geometric forbidding Eq.~\eqref{eq:ch2.geometric-symmetry-operator-def} and radiation nullity
Eq.~\eqref{eq:ch2.angular-spectrum-nullspace-law} are distinct mechanisms
\cite{stratton1941}. Accessible sheet \(\mathcal{K}_{\mathrm{acc}}\) intersects reportable modal
sector Eq.~\eqref{eq:ch2.reconstructability-budget} \cite{harrington2001}.
\begin{equation}
\label{eq:ch2.nr-superposition-law}
\mathcal{R}_\omega\!\left[\Jimp^{(\mathrm{rad})}+\Jimp^{\mathrm{(NR)}}\right]
=
\mathcal{R}_\omega\!\left[\Jimp^{(\mathrm{rad})}\right],
\qquad
\Jimp^{\mathrm{(NR)}}\in\ker\mathcal{R}_\omega,
\end{equation}
\begin{equation}
\label{eq:ch2.nr-modal-invisibility}
\Jimp^{\mathrm{(NR)}}\in\ker\mathcal{R}_\omega
\ \Longrightarrow\
a_n\!\left(\Jimp^{\mathrm{(NR)}}\right)=0,
\quad
\lambda_n=0\ \text{on exporting modes,}
\end{equation}
\begin{equation}
\label{eq:ch2.nr-near-field-energy-law}
\left\|\mathcal{S}_\omega\!\left[\Jimp^{(\mathrm{rad})}+\Jimp^{\mathrm{(NR)}}\right]\right\|_{\mathcal{E}}
\neq
\left\|\mathcal{S}_\omega\!\left[\Jimp^{(\mathrm{rad})}\right]\right\|_{\mathcal{E}}
\ \text{while}\ 
\mathcal{R}_\omega\!\left[\Jimp^{(\mathrm{rad})}+\Jimp^{\mathrm{(NR)}}\right]
=
\mathcal{R}_\omega\!\left[\Jimp^{(\mathrm{rad})}\right],
\end{equation}
For radiating \(\Jimp^{(\mathrm{rad})}\) and non-radiating
\(\Jimp^{\mathrm{(NR)}}\in\ker\mathcal{R}_\omega\),
\begin{equation}
\label{eq:ch2.angular-spectrum-nr-invariance}
a(\widehat{\mathbf{k}};\Jimp^{(\mathrm{rad})}+\Jimp^{\mathrm{(NR)}})
=
a(\widehat{\mathbf{k}};\Jimp^{(\mathrm{rad})}),
\end{equation}
on the outgoing transverse sheet, restating superposition law
Eq.~\eqref{eq:ch2.nr-superposition-law} in angular coordinates
\cite{stratton1941}. Equation~\eqref{eq:ch2.angular-spectrum-nr-invariance}
is the continuous form of modal invisibility Eq.~\eqref{eq:ch2.nr-modal-invisibility} and
explains why inverse angular-spectrum solvers cannot detect
\(\Jimp^{\mathrm{(NR)}}\) from far data alone \cite{stratton1941}. Canonical families
canonical non-radiating configurations (Subsection~\ref{sec:ch272}) therefore induce \emph{invisible} angular
perturbations at fixed \(a(\widehat{\mathbf{k}})\) \cite{harrington2001}.
Evanescent decoupling from angular spectra. Evanescent field content does not enter \(a(\widehat{\mathbf{k}})\):
\begin{equation}
\label{eq:ch2.evanescent-angular-decoupling}
a(\widehat{\mathbf{k}})
\ \text{depends only on}\ 
\Pi_{\mathrm{prop}}\mathcal{S}_\omega[\Jimp]
\ \text{after}\ \Pi_{\mathrm{rad}},
\qquad
\Pi_{\mathrm{ev}}\mathcal{S}_\omega[\Jimp]
\ \not\mapsto\ a(\widehat{\mathbf{k}}),
\end{equation}
by Eqs.~\eqref{eq:ch2.evanescent-nonexport-law} and~\eqref{eq:ch2.evanescent-radiation-annihilation}
\cite{stratton1941}. Equation~\eqref{eq:ch2.evanescent-angular-decoupling}
specializes the second line of Eq.~\eqref{eq:ch2.angular-spectrum-null}:
reactive storage may be large in \(\Pi_{\mathrm{ev}}\mathcal{S}_\omega[\Jimp]\) while
\(a(\widehat{\mathbf{k}})\) reports only qualified export \cite{stratton1941}. regime qualification
Eq.~\eqref{eq:ch2.regime-component-criterion} forbids interpreting \(\mathbf{u}_{\mathrm{ev}}\) as
angular-transport data \cite{harrington2001}.
Observation rank on angular spectra. Finite aperture collapses recoverable angular content:
\begin{equation}
\label{eq:ch2.angular-spectrum-observation-rank}
\{\,a(\widehat{\mathbf{k}})\,\text{recoverable from}\ \mathcal{M}_\omega\,\}
\subseteq
\{\,a(\widehat{\mathbf{k}})\,\text{supported on}\ \mathcal{K}_{\mathrm{acc}}\cap\operatorname{ran}\mathcal{R}_\omega\,\},
\end{equation}
with strict inclusion when \(\ker\mathcal{O}_{\mathrm{meas}}\big|_{\operatorname{ran}\mathcal{R}_\omega}\neq\{\mathbf{0}\}\)
Eq.~\eqref{eq:ch2.observation-null-space} \cite{stratton1941}. Equation
\eqref{eq:ch2.angular-spectrum-observation-rank} links angular spectra to
\(\mathcal{K}_{\mathrm{nrec}}(N)\): sources invisible to \(\mathcal{M}_\omega\) need not have
\(a(\widehat{\mathbf{k}})=\mathbf{0}\) pointwise \cite{stratton1941}. Rank cap
Eq.~\eqref{eq:ch2.reconstruction-rank-cap} limits resolved harmonics independently of
Eq.~\eqref{eq:ch2.angular-spectrum-nullspace-law} \cite{harrington2001}.
\begin{figure}[t]
\centering
\ChOneFigBlock{%
\begin{tikzpicture}[ch1fig, node distance=5mm and 7mm]
 \node[ch1box] (j) {$\mathcal{J}_\Omega$};
 \node[ch1box, right=7mm of j] (r) {$\mathcal{R}_\omega$};
 \node[ch1box, right=7mm of r] (a) {$a(\widehat{\mathbf{k}})$};
 \node[ch1box, right=7mm of a] (k) {$\mathcal{K}_{\mathrm{acc}}$};
 \node[ch1box, below=5mm of j] (ker) {$\ker\mathcal{R}_\omega$};
 \node[ch1box, below=5mm of a] (sup) {$\mathcal{K}_{\mathrm{sup}}$};
 \draw[ch1flow] (j) -- (r) -- (a) -- (k);
 \draw[ch1flow, dashed] (ker.south) -- (j.north);
 \draw[ch1flow, dashed] (sup.south) -- (a.north);
\end{tikzpicture}%
}
\caption{Subsection~\ref{sec:ch274} angular measurement sequence: \(\ker\mathcal{R}_\omega\) annihilates \(a(\widehat{\mathbf{k}})\) Eq.~\eqref{eq:ch2.angular-spectrum-nullspace-law}; suppressed sheet
\(\mathcal{K}_{\mathrm{sup}}\) Eq.~\eqref{eq:ch2.suppressed-angular-sheet-def} removes directions
before accessible sector \(\mathcal{K}_{\mathrm{acc}}\).}
\label{fig:ch2.angular-spectrum-null-measurement chain}
\end{figure}
Figure interpretation. Source excitation $\mathcal{J}_\Omega$ maps through $\mathcal{R}_\omega$ to angular amplitude
$a(\widehat{\mathbf{k}})$; dashed $\ker\mathcal{R}_\omega$ forces $a=0$ on the entire exporting
sheet Eq.~\eqref{eq:ch2.angular-spectrum-nullspace-law}. The suppressed sheet
$\mathcal{K}_{\mathrm{sup}}$ removes directions by geometry before $\mathcal{E}_{\mathrm{acc}}^{(N)}$ is tested:
$a=0$ on $\mathcal{K}_{\mathrm{sup}}$ does not place $\Jimp$ in $\ker\mathcal{R}_\omega$
\cite{stratton1941}. Only the intersection with $\mathcal{K}_{\mathrm{acc}}$ carries
laboratory-reportable angular content \cite{devaney2004}.
\phantomsection\label{thm:ch2.angular-spectrum-null-consequences}
Let \(a(\widehat{\mathbf{k}})\) be the outgoing angular spectrum from
Eq.~\eqref{eq:ch2.exterior-continuous-spectral-integral} and let
\(\Jimp\in\mathcal{J}_\Omega\) be admissible. Then:
\begin{enumerate}
\item[(i)] \(\Jimp\in\ker\mathcal{R}_\omega\Rightarrow a(\widehat{\mathbf{k}})=0\) on the
exporting sheet Eq.~\eqref{eq:ch2.angular-spectrum-nullspace-law}.
\item[(ii)] Superposition with \(\Jimp^{\mathrm{(NR)}}\in\ker\mathcal{R}_\omega\) leaves
\(a(\widehat{\mathbf{k}})\) unchanged Eq.~\eqref{eq:ch2.angular-spectrum-nr-invariance}.
\item[(iii)] Evanescent content satisfies Eq.~\eqref{eq:ch2.evanescent-angular-decoupling} and
does not alter \(a(\widehat{\mathbf{k}})\) when export is unchanged.
\item[(iv)] Full angular budget:
\begin{equation}
\label{eq:ch2.angular-spectrum-composition-law}
\text{support}\big(a(\widehat{\mathbf{k}})\big)
\subseteq
\mathcal{K}_{\mathrm{acc}}
\cap
\{\widehat{\mathbf{k}}:\ \text{resolvable by}\ \mathcal{M}_\omega\}
\cap
\{\widehat{\mathbf{k}}:\ \lambda_n\neq 0\ \text{modal labels}\},
\end{equation} composing geometric, observation, and Fredholm constraints from
Eqs.~\eqref{eq:ch2.constraint-composition-law}--\eqref{eq:ch2.reconstructability-budget}
\cite{stratton1941}.
\end{enumerate}
Part (i) follows from \(\mathcal{R}_\omega[\Jimp]=\mathbf{0}\) and far-zone reduction
the far-field approximation of source integrals (Subsection~\ref{sec:ch252}) \cite{stratton1941}. Part (ii) uses
canonical non-radiating configurations (Subsection~\ref{sec:ch272}) and linearity of angular extraction
\cite{harrington2001}. Part (iii) is evanescent content never radiates (Subsection~\ref{sec:ch273}) in
angular coordinates \cite{stratton1941}. Part (iv) composes
Eq.~\eqref{eq:ch2.accessible-modal-sector}, angular probe rank on the far shell (Subsection~\ref{sec:ch264}), and
Eq.~\eqref{eq:ch2.angular-spectrum-observation-rank} \cite{stratton1941}. \phantomsection\label{prop:ch2.continuous-discrete-null-correspondence}
Modal coefficients in Eq.~\eqref{eq:ch2.spectral-far-pattern-expansion} and continuous
\(a(\widehat{\mathbf{k}})\) satisfy
\begin{equation}
\label{eq:ch2.modal-continuous-null-link}
\lambda_n\,c_n[\mathfrak{D}]=0\ \text{for all exporting}\ n
\ \Longleftrightarrow\
a(\widehat{\mathbf{k}})=0\ \text{on}\ S^{2}
\end{equation}
when the modal basis is complete on \(\operatorname{ran}\mathcal{R}_\omega\) in the rigged sense
\cite{stratton1941}. Equation~\eqref{eq:ch2.modal-continuous-null-link} connects
Eq.~\eqref{eq:ch2.nullspace-modal-annihilation-law} to
Eq.~\eqref{eq:ch2.angular-spectrum-nullspace-law} without re-proving completeness in
Chapter~3 (\emph{Angular Momentum Operators and SO(3)}) \cite{stratton1941}. Truncated modal books misstate
Eq.~\eqref{eq:ch2.modal-continuous-null-link} on exterior domains
\cite{jackson1999}.
far-pattern coefficients from spectral modes (Subsection~\ref{sec:ch261}) and
Eq.~\eqref{eq:ch2.exterior-continuous-spectral-integral} identify discrete and continuous
representations on outgoing transverse modes \cite{stratton1941}. Null modal coefficients imply
vanishing \(\E_0\); conversely \(\E_0=\mathbf{0}\) forces all exporting modal weights to zero
\cite{harrington2001}.
\phantomsection\label{prop:ch2.angular-spectrum-inverse-obstruction}
Given target angular spectrum \(a_{\mathrm{tgt}}(\widehat{\mathbf{k}})\), the source inversion
problem seeks \(\Jimp\) with \(a(\widehat{\mathbf{k}};\Jimp)=a_{\mathrm{tgt}}(\widehat{\mathbf{k}})\).
Any solution is unique only up to additions in \(\ker\mathcal{R}_\omega\):
\begin{equation}
\label{eq:ch2.angular-inverse-null-quotient}
\Jimp^{(1)}-\Jimp^{(2)}\in\ker\mathcal{R}_\omega
\quad\text{whenever}\quad
a(\widehat{\mathbf{k}};\Jimp^{(1)})=a(\widehat{\mathbf{k}};\Jimp^{(2)}),
\end{equation}
and up to \(\mathcal{K}_{\mathrm{nrec}}(N)\) at finite observation rank
\cite{stratton1941}. Equation~\eqref{eq:ch2.angular-inverse-null-quotient} is the
angular restatement of Eq.~\eqref{eq:ch2.inverse-source-equation} ill-posedness; Tikhonov
regularization Eq.~\eqref{eq:ch2.tikhonov-source-reconstruction} selects a representative without
eliminating the quotient \cite{stratton1941}.
If \(a(\widehat{\mathbf{k}};\Jimp^{(1)})=a(\widehat{\mathbf{k}};\Jimp^{(2)})\), then
\(\mathcal{R}_\omega[\Jimp^{(1)}-\Jimp^{(2)}]=\mathbf{0}\) on \(\mathcal{F}_{\mathrm{rad}}\)
\cite{stratton1941}. Finite \(\mathcal{M}_\omega\) adds measurement null directions
measurement null-space extension (Subsection~\ref{sec:ch271}) \cite{harrington2001}. The nullity--rank-$N$ recovery--realizability triad introduced in Section~\ref{sec:ch27} applies
unchanged in angular coordinates: far data cannot detect additions in \(\ker\mathcal{R}_\omega\)
Eq.~\eqref{eq:ch2.angular-spectrum-nr-invariance}; accessible directions survive geometric and
regime filters even when export is nonzero; $\mathcal{S}_{\mathrm{real}}$ of a target
\(a_{\mathrm{target}}(\widehat{\mathbf{k}})\) is decided in \ref{ch:04} \cite{devaney2004}. Evanescent overlap
Eq.~\eqref{eq:ch2.evanescent-source-overlap} changes near fields without altering \(a\) when
export is fixed \cite{stratton1941}.
Forward link to angular-spectrum action. Section~\ref{sec:ch29} develops plane-wave operators that act on \(a(\widehat{\mathbf{k}})\);
Subsection~\ref{sec:ch294} refines far-field angular $\mathcal{E}_{\mathrm{acc}}^{(N)}$ on
\(\mathcal{K}_{\mathrm{acc}}\) \cite{stratton1941}. Subsection~\ref{sec:ch274} fixes the \emph{source-side} constraints those operators inherit:
Eqs.~\eqref{eq:ch2.angular-spectrum-nullspace-law}--\eqref{eq:ch2.angular-spectrum-composition-law}
must be composed with regime qualification Eq.~\eqref{eq:ch2.regime-operator-screen} before transport
claims \cite{stratton1941}. Reciprocity identities on coupled angular modes belong to
Section~\ref{sec:ch28} \cite{harrington2001}.
Null-space consequences for angular spectra (Subsection~\ref{sec:ch274}) assumes linear Maxwell theory at fixed
\(\omega\), outgoing radiation condition, and rigged-sense convergence of
Eq.~\eqref{eq:ch2.exterior-continuous-spectral-integral} \cite{stratton1941}.
Loss shifts phase of \(a(\widehat{\mathbf{k}})\) without converting evanescent storage to export
\cite{stratton1941}. Vector harmonic explicit forms of \(\mathbf{f}_n(\theta,\phi)\) belong to
\ref{ch:03}; Subsection~\ref{sec:ch274} records operator consequences only
\cite{harrington2001}. As a worked example, place a copper wire loop of radius \(b=\lambda_0/8\) in vacuum and excite it with a uniform
phasor current at \(\omega\). The loop realization \(\Jimp^{\mathrm{(NR)}}\) from
Eq.~\eqref{eq:ch2.nr-loop-cancellation} satisfies
Eq.~\eqref{eq:ch2.angular-spectrum-nullspace-law}---a wattmeter on a sphere at
\(R=10\lambda_0\) registers zero qualified export---while \(|\E|\) may exceed the free-space
impedance value by an order of magnitude at \(r=b\) \cite{stratton1941}. The near-field probe
therefore sees rich reactive structure that the angular spectrum cannot detect
\cite{harrington2001}. Angular-spectrum
probes of \(\mathcal{S}_\omega[\Jimp^{\mathrm{(NR)}}]\) without \(\Pi_{\mathrm{rad}}\) misstate
null-space membership \cite{harrington2001}.
As a worked example, when \(\mathcal{R}_\omega[\Jimp]\neq\mathbf{0}\) but the aperture resolves only \(N\) harmonics,
Eq.~\eqref{eq:ch2.angular-spectrum-observation-rank} leaves unresolved \(\widehat{\mathbf{k}}\)
in the measurement null space even when \(a(\widehat{\mathbf{k}})\) is nonzero
\cite{stratton1941}. This is distinct from
Eq.~\eqref{eq:ch2.angular-spectrum-nr-invariance} \cite{stratton1941}. \medskip
\noindent
Section~\ref{sec:ch27} closes with the null-space chain begun in its introduction:
operator structure and \(\ker\mathcal{R}_\omega\) (Subsection~\ref{sec:ch271}), canonical
non-radiating current families (Subsection~\ref{sec:ch272}),
evanescent subspaces and non-reconstructability (Subsection~\ref{sec:ch273}), and angular-spectrum
consequences (Subsection~\ref{sec:ch274}) with accessible and
suppressed sheets Eqs.~\eqref{eq:ch2.angular-accessibility-sector}--\eqref{eq:ch2.suppressed-angular-sheet-def}
(Subsection~\ref{sec:ch274}) \cite{stratton1941}. Spectral and
geometric filters from Section~\ref{sec:ch26} now meet source-side invisibility: far-zone reports
use only \(a(\widehat{\mathbf{k}})\) on the composed budget
Eq.~\eqref{eq:ch2.angular-spectrum-composition-law}, while \(\ker\mathcal{R}_\omega\),
\(\Pi_{\mathrm{ev}}\mathcal{S}_\omega[\Jimp]\), and \(\mathcal{K}_{\mathrm{nrec}}(N)\) explain
deficits that no angular label alone resolves \cite{stratton1941}. Section~\ref{sec:ch28} adds
reciprocity on the same modal coordinates; Section~\ref{sec:ch29} and Subsection~\ref{sec:ch294}
develop plane-wave action and far-field angular rank-$N$ recovery on \(\mathcal{K}_{\mathrm{acc}}\);
\ref{ch:03} supplies explicit harmonic bases \cite{harrington2001}. Transport and
angular-transport claims must compose with
Eq.~\eqref{eq:ch2.section27-nullspace-chain},
Eq.~\eqref{eq:ch2.angular-spectrum-composition-law}, and regime qualification
Eq.~\eqref{eq:ch2.regime-operator-screen}, not with raw \(\mathcal{S}_\omega[\Jimp]\) or
unqualified near-zone reactive content \cite{stratton1941}. Source
null spaces explain what far-zone observation cannot detect; synthesizable---whether
accessible angular content is generable by finite \(\Jimp\)---is decided in
\ref{ch:04} alongside source equivalence and effective rank.
\section{Reciprocity and Radiative Coupling}
\label{sec:ch28}
Null-space structure explains invisible near-field rearrangements; reciprocity explains why passive hardware nevertheless reports the same complex coupling when transmit and receive roles are exchanged \cite{stratton1941}. On a calibrated vector network analyzer the practitioner excites port~$A$, reads the complex response at port~$B$, then exchanges cables, standards, and probe normalization so that port~$B$ drives and port~$A$ receives; on reciprocal antennas and reciprocal media the normalized scattering entries $S_{21}$ and $S_{12}$ must agree once outgoing qualification and impressed-current normalization are held fixed \cite{hansen1988}. Two-port impedance matrices $\mathbf{Z}$ and scattering matrices $\mathbf{S}$ inherit the same interchange symmetry whenever terminations and reference planes are declared consistently \cite{collin2001}. That bench fact precedes every volume identity below: reciprocity is what the network analyzer already presupposes when a passive feed is tuned by comparing forward and reverse coupling \cite{stratton1941}.
% PHYSICAL-MEMORY:sec:ch28
% PHYSICAL-MEMORY:sec:ch28-b2
\begin{physicalvignette}[Reflector with shaped beam for satellite downlink]
A second bench observation matters on Open cavity slot fed by a microstrip line. When technicians swap only the receive probe while holding the transmit chain fixed, reciprocity checks whether the reported coupling is property of the structure or artifact of a mismatched flux gauge \cite{stratton1941}. Failed reciprocity usually traces to a near-field tap, a radome standing-wave, or a probe polarization basis that was never aligned with the transmit standard \cite{hansen1988}.
Rayleigh and Stratton treated such ranges as boundary-value laboratories long before digital near-field transforms became routine \cite{stratton1941}. Harrington's method-of-moments culture reinforced the same discipline: metal outline first, harmonic headline second \cite{harrington2001}.
The measurement logic above is made precise in Subsection~\ref{sec:ch28} \cite{stratton1941}.
\end{physicalvignette}
\begin{physicalvignette}[Reflector with shaped beam for satellite downlink]
Consider Reflector with shaped beam for satellite downlink. Before any coefficient table is printed, technicians record frequency, cable loss, and the mechanical outline that fixes electrical size. The first walk-through of the range is always tactile: flange bolts, probe standoff, turntable boresight, and the wattmeter reference on the receive chain \cite{hansen1988}.
On this hardware the near zone is visually impressive---standing reactive structure, rapid phase rotation, sometimes a visible VSWR ripple at the feed---yet the export question is measured only after outgoing selection on a large sphere \cite{harrington2001}. What the camera sees in the chamber is not automatically what the far-zone pattern claims \cite{stratton1941}.
Engineers therefore learn to separate three readings: what the feed can synthesize on declared support, what propagating content survives on the radiative range, and what a finite instrument can recover at declared tolerance \cite{devaney2004}. Confusing those readings produces the rank gaps and watt failures that later sections quantify \cite{pal2026}.
Rayleigh and Stratton treated such ranges as boundary-value laboratories long before digital near-field transforms became routine \cite{stratton1941}. Harrington's method-of-moments culture reinforced the same discipline: metal outline first, harmonic headline second \cite{harrington2001}.
Subsection~\ref{sec:ch28} states the formal export condition that this range scene presupposes \cite{stratton1941}.
\end{physicalvignette}
Section~\ref{sec:ch27} closed the null-space chain
Eq.~\eqref{eq:ch2.section27-nullspace-chain}: operator kernels, canonical non-radiating
currents, evanescent non-reconstructability, and angular-spectrum consequences
(Subsection~\ref{sec:ch274}) \cite{stratton1941}. That program classified invisible and
partially observable source content on \(\mathcal{R}_\omega\) and \(\mathcal{M}_\omega\).
Section~\ref{sec:ch28} addresses exchange symmetry: when transmit and receive maps coincide,
how radiative modes couple under bilinear pairings, and which orthogonality relations constrain
coupled spectra on \(\mathcal{F}_{\mathrm{rad}}\) \cite{stratton1941}. Reciprocity ties
transmission and reception on passive hardware before angular spectra are filtered; broken
reciprocity after probe correction usually traces to inconsistent normalization or missing
outgoing tails \cite{hansen1988}. Preview swap Eq.~\eqref{eq:ch2.dyadic-reciprocity-swap}
from Subsection~\ref{sec:ch231} named kernel interchange; Section~\ref{sec:ch28} upgrades
reciprocity to operator form on \(\mathcal{S}_\omega\), \(\mathcal{R}_\omega\), and modal
coordinates from Sections~\ref{sec:ch26}--\ref{sec:ch27} \cite{stratton1941}. Plane-wave
angular-spectrum action belongs to Section~\ref{sec:ch29}; explicit harmonic orthogonality
and completeness belong to \ref{ch:03} Subsection~3.4.3; continuous-spectrum normalization
belongs to Subsection~3.5.2 by reference to Subsection~\ref{sec:ch282}
\cite{harrington2001}.
Spectral decomposition in Subsection~\ref{sec:ch261} and biorthogonality
Eq.~\eqref{eq:ch2.radiation-biorthogonality} fixed modal coordinates on
\(\mathcal{F}_{\mathrm{rad}}\); null-space structure in Subsection~\ref{sec:ch271} showed that
many sources share the same radiative image \cite{stratton1941}. Reciprocity does not remove
\(\ker\mathcal{R}_\omega\), but it constrains how coupling coefficients between declared modes
behave when impressed sources and test probes are swapped \cite{stratton1941}. Treating antenna
transmit and receive as unrelated maps violates Eq.~\eqref{eq:ch2.dyadic-reciprocity-swap} even
when outgoing source-to-field integrals from Subsection~\ref{sec:ch251} are numerically accurate
\cite{harrington2001}. Section~\ref{sec:ch28} separates Lorentz reciprocity in operator form,
modal orthogonality and coupling, and radiation--reception interchange without re-opening Green
construction from Section~\ref{sec:ch23} \cite{stratton1941}.
Fix angular frequency \(\omega>0\), impressed sources in \(\mathcal{J}_\Omega\), and outgoing
field pairs in \(\mathcal{F}_{\mathrm{SF}}\) Eq.~\eqref{eq:ch2.admissible-field-space}.
The section develops reciprocity in three stages that mirror the measurement sequence excite--couple--swap:
\begin{equation}
\label{eq:ch2.section28-reciprocity-chain}
\text{Lorentz reciprocity}
\ \xrightarrow{\ \text{operator form}\ }
\mathcal{S}_\omega,\mathcal{R}_\omega
\ \xrightarrow{\ \text{modal}\ }
\text{orthogonality / coupling}
\ \xrightarrow{\ \text{systems}\ }
\text{radiation \& reception},
\end{equation}
where the first arrow is Subsection~\ref{sec:ch281}, the second is Subsection~\ref{sec:ch282},
and the third is Subsection~\ref{sec:ch283} \cite{stratton1941}. Equation
\eqref{eq:ch2.section28-reciprocity-chain} expresses Eq.~\eqref{eq:ch2.section26-spectral-chain}
and Eq.~\eqref{eq:ch2.section27-nullspace-chain} at the bilinear level: spectral modes from
Section~\ref{sec:ch26} supply coordinates; null-space quotients from Section~\ref{sec:ch27}
explain non-uniqueness; reciprocity constrains \emph{symmetric} parts of coupling on
\(\operatorname{ran}\mathcal{R}_\omega\) \cite{stratton1941}. Angular budget
Eq.~\eqref{eq:ch2.angular-spectrum-composition-law} and regime qualification
Eq.~\eqref{eq:ch2.regime-operator-screen} remain mandatory before transport claims on coupled
modes \cite{harrington2001}.
Denote by $\mathcal{M}_A$ and $\mathcal{M}_B$ the finite-rank port maps that compose radiation
selection with probe functionals on $\mathcal{F}_{\mathrm{rad}}$; the port-swap experiment asserts
$\mathcal{M}_B[\Jimp^{\mathrm{(tx)}}]=\mathcal{M}_A[\Jimp^{\mathrm{(rx)}}]$ on reciprocal
excitations before any volume contraction is invoked \cite{harrington2001}. Maxwell adjoint pairing
Eq.~\eqref{eq:ch2.maxwell-adjoint-pairing}, source dual pairing
Eq.~\eqref{eq:ch2.dual-source-pairing}, and radiation map
Eq.~\eqref{eq:ch2.spectral-operator-targets} supply the continuous objects that finite port
matrices sample; dyadic reciprocity swap Eq.~\eqref{eq:ch2.dyadic-reciprocity-swap} and modal
biorthogonality Eq.~\eqref{eq:ch2.radiation-biorthogonality} specialize the same symmetry on
outgoing branches fixed by
Eqs.~\eqref{eq:ch2.outgoing-radiation-refinement}--\eqref{eq:ch2.outgoing-resolvent-unique}
\cite{stratton1941}. Subsection~\ref{sec:ch281} upgrades the bench fact to
sesquilinear form on $\mathcal{S}_\omega$ and $\mathcal{R}_\omega$; the identity below states the
operator identity the VNA experiment already presupposes.
Lorentz reciprocity (leading template). Subsection~\ref{sec:ch281} supplies the definitive operator statement: for
admissible source pairs \((\Jimp^{(1)},\Jimp^{(2)})\) and induced fields
\((\mathbf{u}_1,\mathbf{u}_2)\),
\begin{equation}
\label{eq:ch2.lorentz-reciprocity}
\langle \mathbf{u}_1,\,\mathcal{S}_\omega[\Jimp^{(2)}]\rangle_{\mathcal{U}}
=
\langle \mathbf{u}_2,\,\mathcal{S}_\omega[\Jimp^{(1)}]\rangle_{\mathcal{U}}
\end{equation}
on passive linear media without magneto-electric coupling, with pairing induced from
Eq.~\eqref{eq:ch2.maxwell-adjoint-pairing} on outgoing admissible states
\cite{stratton1941}. Equation~\eqref{eq:ch2.lorentz-reciprocity}
packages the port-swap symmetry as a sesquilinear identity on admissible Maxwell pairs: field
one tested against source two equals field two tested against source one when media and
boundaries belong to the reciprocal class \cite{stratton1941}. It develops Eq.
Eq.~\eqref{eq:ch2.dyadic-reciprocity-swap} from kernel swap to full source--field reciprocity
on \(\mathcal{S}_\omega\); restrictions on loss, anisotropy, and boundary class are stated in
Subsection~\ref{sec:ch281} \cite{stratton1941}. Reciprocity on
\(\mathcal{R}_\omega=\Pi_{\mathrm{rad}}\circ\mathcal{S}_\omega\) follows by projection when
both sides satisfy regime qualification Eq.~\eqref{eq:ch2.regime-component-criterion}
\cite{harrington2001}.
On radiation eigenmodes Eq.~\eqref{eq:ch2.radiation-eigenpair-def}, Subsection~\ref{sec:ch282}
develops orthogonality and coupling coefficients:
\begin{equation}
\label{eq:ch2.modal-orthogonality}
\langle \mathbf{u}_m,\,\mathbf{u}_n\rangle_{\mathcal{U}}
=
\delta_{mn},
\qquad
C_{mn}
=
\langle \widetilde{\mathbf{u}}_m,\,\mathcal{R}_\omega[\Jimp^{(n)}]\rangle_{\mathcal{U}},
\end{equation}
with \(\widetilde{\mathbf{u}}_m\) adjoint modes tied to
Eq.~\eqref{eq:ch2.radiation-biorthogonality} \cite{stratton1941}. Equation
\eqref{eq:ch2.modal-orthogonality} is the Section~\ref{sec:ch28} target for
Subsection~\ref{sec:ch282} orthogonality and coupling on \(\mathcal{F}_{\mathrm{rad}}\);
explicit vector harmonic normalization belongs to \ref{ch:03}
\cite{stratton1941}. Modes with \(\lambda_n=0\) lie in null sectors
Eq.~\eqref{eq:ch2.modal-null-index-set}; reciprocity constraints apply on exporting modes
with \(\lambda_n\neq 0\) \cite{harrington2001}.
Radiation and reception (leading template). Subsection~\ref{sec:ch283} records transmit--receive interchange for coupled systems:
\begin{equation}
\label{eq:ch2.transmit-receive-swap}
\mathcal{M}_{\mathrm{rx}}[\Jimp^{\mathrm{(tx)}}]
=
\mathcal{M}_{\mathrm{tx}}[\Jimp^{\mathrm{(rx)}}]
\quad\text{under reciprocity-qualified pairings},
\end{equation}
with \(\mathcal{M}_{\mathrm{tx}}\), \(\mathcal{M}_{\mathrm{rx}}\) finite-rank compositions of
\(\mathcal{R}_\omega\) and \(\mathcal{O}_{\mathrm{meas}}\) Eq.~\eqref{eq:ch2.reconstruction-observation-map}
\cite{stratton1941}. Equation~\eqref{eq:ch2.transmit-receive-swap}
does not eliminate \(\ker\mathcal{R}_\omega\) or \(\mathcal{K}_{\mathrm{nrec}}(N)\); it
constrains \emph{symmetric} coupling on observable coordinates after projection onto
\(\operatorname{ran}\mathcal{R}_\omega\) \cite{stratton1941}. Geometric forbidden sectors
Eq.~\eqref{eq:ch2.accessible-modal-sector} and angular budget
Eq.~\eqref{eq:ch2.angular-spectrum-composition-law} compose independently
\cite{harrington2001}.
Radiative coupling bilinear form (leading template). Coupling between two radiating subsystems is tracked by
\begin{equation}
\label{eq:ch2.radiative-coupling-bilinear}
\mathfrak{C}_\omega(\Jimp^{(1)},\Jimp^{(2)})
=
\langle \mathcal{R}_\omega[\Jimp^{(1)}],\,\mathcal{R}_\omega[\Jimp^{(2)}]\rangle_{\mathcal{U}}
\quad\text{on}\ \mathcal{F}_{\mathrm{rad}},
\end{equation}
with reciprocity implying \(\mathfrak{C}_\omega(\Jimp^{(1)},\Jimp^{(2)})=\mathfrak{C}_\omega(\Jimp^{(2)},\Jimp^{(1)})\)
under Eq.~\eqref{eq:ch2.lorentz-reciprocity} and passive media
\cite{stratton1941}. Equation~\eqref{eq:ch2.radiative-coupling-bilinear}
measures export overlap, not raw near-zone reactive overlap
\(\langle \mathcal{S}_\omega[\Jimp^{(1)}],\,\mathcal{S}_\omega[\Jimp^{(2)}]\rangle\), which may include
\(\Pi_{\mathrm{ev}}\) content annihilated by Eq.~\eqref{eq:ch2.rad-null-on-evanescent}
\cite{stratton1941}. regime qualification Eq.~\eqref{eq:ch2.regime-operator-screen} qualifies when
\(\mathfrak{C}_\omega\) may be reported as transport data \cite{harrington2001}.
The three subsections that follow develop Lorentz reciprocity, modal orthogonality and coupling, and transmit--receive exchange in sequence.
\subsection{Lorentz reciprocity in operator form}
\label{sec:ch281}
% HARRINGTON-OPEN:sec:ch281
On a declared compact-range or anechoic campaign, technicians record frequency, outline, and probe geometry before any harmonic fit is trusted \cite{hansen1988}. Reciprocity links transmission and reception on the same passive structure \cite{stratton1941}. Two-port calibrations depend on it \cite{jackson1999}. Maxwell's retarded export to $\Gamma_R$ remains the governing constraint \cite{maxwell1873}.
Reciprocity is the laboratory fact that passive linear antennas exchange the same complex
transfer when transmit and receive roles swap on a calibrated vector network analyzer
\cite{stratton1941}. The adjoint radiation map extends that symmetry for continuous sources
and traces, enabling source reconstruction and equivalence checks without re-solving unrelated
models \cite{devaney2004}. Reflector feeds, horn coupling, and patch arrays are routinely tuned
by comparing \(S_{21}\) and \(S_{12}\); persistent asymmetry after calibration usually signals a
broken outgoing radiation condition or incorrect impressed-current normalization rather than a
genuine non-reciprocal medium \cite{hansen1988}. Reciprocity cannot be set aside merely because
``only the radiative image matters'': pattern, gain, and scattering matrices on passive hardware
are reciprocity-linked observables tied to the same Maxwell operator
\cite{stratton1941}. Lorentz reciprocity is therefore a consequence of passive
linear Maxwell theory, not an optional post-check applied after the forward solve
\cite{stratton1941}. Two-port calibrations and bistatic radar-cross-section matrices inherit the
same exchange symmetry whenever media and boundaries satisfy the reciprocal class stated below
\cite{jackson1999}. The volume identity that follows states the same swap observed on the
VNA once impressed sources and test fields are paired through the Maxwell adjoint.
\medskip
Section~\ref{sec:ch28} established bilinear exchange symmetry through
Eqs.~\eqref{eq:ch2.section28-reciprocity-chain}--\eqref{eq:ch2.lorentz-reciprocity}.
Subsection~\ref{sec:ch281} supplies the definitive account of Lorentz reciprocity on
\(\mathcal{S}_\omega\) and \(\mathcal{R}_\omega\), including the full proof of dyadic swap
Eq.~\eqref{eq:ch2.dyadic-reciprocity-swap} deferred from Subsection~\ref{sec:ch231}
\cite{stratton1941}. Equation~\eqref{eq:ch2.maxwell-adjoint-pairing}, source dual pairing
Eq.~\eqref{eq:ch2.dual-source-pairing}, outgoing radiation condition
Eqs.~\eqref{eq:ch2.outgoing-radiation-refinement}--\eqref{eq:ch2.outgoing-resolvent-unique},
and outgoing source-to-field integrals of Subsection~\ref{sec:ch251} are presupposed;
Green construction and Sommerfeld uniqueness are not re-opened \cite{stratton1941}. Modal
orthogonality and coupling belong to Subsection~\ref{sec:ch282}; transmit--receive systems to
Subsection~\ref{sec:ch283} \cite{harrington2001}.
Assumptions. Fix \(\omega>0\), Lipschitz domain \(\Omega\) with piecewise smooth interfaces, impressed
sources in \(\mathcal{J}_\Omega\) Eq.~\eqref{eq:ch2.source-space-refined} satisfying
Eq.~\eqref{eq:ch2.source-admissibility}, and outgoing fields in
\(\mathcal{F}_{\mathrm{SF}}\) Eq.~\eqref{eq:ch2.admissible-field-space}
\cite{stratton1941}. Constitutive tensors are linear, local, and passive at
\(\omega\), with symmetric permittivity and permeability in the sense
% PHYSICAL-MEMORY:sec:ch281
% PHYSICAL-MEMORY:sec:ch281-b2
\begin{physicalvignette}[Dual-offset reflector for compact range quiet zone]
A second bench observation matters on Spherical NF range with theta-phi robot arm. When technicians swap only the receive probe while holding the transmit chain fixed, reciprocity checks whether the reported coupling is property of the structure or artifact of a mismatched flux gauge \cite{stratton1941}. Failed reciprocity usually traces to a near-field tap, a radome standing-wave, or a probe polarization basis that was never aligned with the transmit standard \cite{hansen1988}.
Rayleigh and Stratton treated such ranges as boundary-value laboratories long before digital near-field transforms became routine \cite{stratton1941}. Harrington's method-of-moments culture reinforced the same discipline: metal outline first, harmonic headline second \cite{harrington2001}.
The measurement logic above is made precise in Subsection~\ref{sec:ch281} \cite{stratton1941}.
\end{physicalvignette}
\begin{physicalvignette}[Dual-offset reflector for compact range quiet zone]
Consider Dual-offset reflector for compact range quiet zone. Before any coefficient table is printed, technicians record frequency, cable loss, and the mechanical outline that fixes electrical size. The first walk-through of the range is always tactile: flange bolts, probe standoff, turntable boresight, and the wattmeter reference on the receive chain \cite{hansen1988}.
On this hardware the near zone is visually impressive---standing reactive structure, rapid phase rotation, sometimes a visible VSWR ripple at the feed---yet the export question is measured only after outgoing selection on a large sphere \cite{harrington2001}. What the camera sees in the chamber is not automatically what the far-zone pattern claims \cite{stratton1941}.
Engineers therefore learn to separate three readings: what the feed can synthesize on declared support, what propagating content survives on the radiative range, and what a finite instrument can recover at declared tolerance \cite{devaney2004}. Confusing those readings produces the rank gaps and watt failures that later sections quantify \cite{pal2026}.
Reciprocity links transmission and reception on the same passive structure . Two-port calibrations depend on it .
Rayleigh and Stratton treated such ranges as boundary-value laboratories long before digital near-field transforms became routine \cite{stratton1941}. Harrington's method-of-moments culture reinforced the same discipline: metal outline first, harmonic headline second \cite{harrington2001}.
Subsection~\ref{sec:ch281} states the formal export condition that this range scene presupposes \cite{stratton1941}.
\end{physicalvignette}
\begin{equation}
\label{eq:ch2.reciprocal-medium-class}
\overline{\overline{\varepsilon}}(\mathbf{r})=\overline{\overline{\varepsilon}}^{\mathsf{T}}(\mathbf{r}),
\qquad
\overline{\overline{\mu}}(\mathbf{r})=\overline{\overline{\mu}}^{\mathsf{T}}(\mathbf{r}),
\qquad
\overline{\overline{\xi}}=\overline{\overline{\zeta}}=\mathbf{0},
\end{equation}
where \(\overline{\overline{\xi}}\), \(\overline{\overline{\zeta}}\) denote magneto-electric
coupling blocks in the constitutive relation imported from Section~\ref{sec:ch01.1.2}
\cite{stratton1941}. Boundary operator \(\mathcal{B}_\Gamma\) is passive
Eq.~\eqref{eq:ch2.passive-impedance-class} and identical for both excitations
\cite{stratton1941}. Field pairing \(\langle\cdot,\cdot\rangle_{\mathcal{U}}\)
is induced by Eq.~\eqref{eq:ch2.field-energy-norm} on admissible Maxwell pairs
\cite{stratton1941}. regime qualification Eq.~\eqref{eq:ch2.regime-operator-screen} governs transport
interpretation of any bilinear quantity built from \(\Pi_{\mathrm{rad}}\)
\cite{harrington2001}.
Equation~\eqref{eq:ch2.reciprocal-medium-class} is the Chapter~\ref{ch:02} $\mathcal{D}_J=1$
class for \emph{full} Lorentz reciprocity at fixed \(\omega\): symmetric \(\varepsilon\),
\(\mu\) ensure time-reversal symmetry of the harmonic Maxwell system modulo outgoing
selection, while vanishing magneto-electric blocks exclude non-reciprocal gyrotropic or
Tellegen couplings \cite{stratton1941}. Isotropic lossless media are a special
case with scalar \(\varepsilon\), \(\mu\). Lossy symmetric tensors preserve reciprocity
relations stated below with complex conjugation on one slot of each pairing
\cite{stratton1941}. Active or negative-real impedance portions of
\(\mathcal{B}_\Gamma\) break passive symmetry and are excluded from the present theorem
\cite{stratton1941}.
Let \((\E_1,\H_1)=\mathcal{S}_\omega[\Jimp^{(1)}]\) and
\((\E_2,\H_2)=\mathcal{S}_\omega[\Jimp^{(2)}]\) be outgoing admissible solutions for
independent impressed currents \(\Jimp^{(1)},\Jimp^{(2)}\in\mathcal{J}_\Omega\)
\cite{stratton1941}. Here \(\Jimp^{(k)}\) denotes the impressed electric current density in
\(\mathrm{A\,m^{-2}}\), \(\E_k\) and \(\H_k\) are the induced phasor fields in \(\mathrm{V\,m^{-1}}\)
and \(\mathrm{A\,m^{-1}}\), and \(\Gamma\) is the passive boundary on which the surface integral
below is evaluated. Physically, the left member is the bilinear coupling measured when port labels
are swapped on reciprocal hardware; the right member is the surface remainder after integration by
parts. The classical Lorentz reciprocity relation reads
\begin{equation}
\label{eq:ch2.lorentz-volume-reciprocity}
\int_\Omega
\big(\E_1\cdot\Jimp^{(2)*}-\E_2\cdot\Jimp^{(1)*}\big)\,d^{3}r
=
\int_\Gamma
\big(\hat{\mathbf{n}}\times\E_1\cdot\H_2^{*}
-\hat{\mathbf{n}}\times\E_2\cdot\H_1^{*}\big)\,dS,
\end{equation}
when \(\overline{\overline{\varepsilon}}\), \(\overline{\overline{\mu}}\) satisfy
Eq.~\eqref{eq:ch2.reciprocal-medium-class} and \(\mathcal{B}_\Gamma\) is passive
\cite{stratton1941}. Equation~\eqref{eq:ch2.lorentz-volume-reciprocity}
balances impressed-source bilinear forms against a boundary pairing generated by
Eq.~\eqref{eq:ch2.maxwell-adjoint-pairing}: the left member measures how each field couples
to the other's source, while the right member is the surface term that remains after
integration by parts on the Maxwell operator \cite{stratton1941}. Interpretation: when the
boundary term vanishes on admissible outgoing pairs, the port-swap scalar observed on the VNA
equals the volume contraction on the left. On domains where
\(\mathcal{B}_\Gamma\) annihilates the boundary integrand for admissible outgoing pairs---PEC,
homogeneous impedance with Eq.~\eqref{eq:ch2.passive-impedance-class}, and standard
Sommerfeld exterior completions---the surface term vanishes and volume reciprocity closes
\cite{stratton1941}.
Derivation sketch (volume form). Substitute \(\mathbf{u}_\omega^{(1)}=(\E_1,\H_1)\) and \(\mathbf{u}_\omega^{(2)}=(\E_2,\H_2)\)
into Eq.~\eqref{eq:ch2.maxwell-adjoint-pairing} and interchange the two states
\cite{stratton1941}. Symmetry of \(\overline{\overline{\varepsilon}}\),
\(\overline{\overline{\mu}}\) under Eq.~\eqref{eq:ch2.reciprocal-medium-class} makes the
volume pairing sesquilinear form on \(\mathcal{M}_\omega[\mathbf{u}_\omega]\) invariant under
label swap when sources are impressed currents in the Amp\`ere row
\cite{stratton1941}. Subtracting the two adjoint identities eliminates
\(\langle\mathcal{M}_\omega[\mathbf{u}_\omega^{(1)}],\mathbf{u}_\omega^{(2)}\rangle_{\Omega}\)
and yields Eq.~\eqref{eq:ch2.lorentz-volume-reciprocity} with the boundary term displayed
\cite{stratton1941}. Charge continuity Eq.~\eqref{eq:ch1.charge-continuity-harmonic} and
compatibility Eq.~\eqref{eq:ch2.source-compatibility} ensure that impressed \(\Jimp\) alone
specifies the electric drive when \(\rho\) is tied to \(\Jimp\)
\begin{figure}[t]
\centering
\ChOneFigBlock{%
\begin{tikzpicture}[ch1fig, node distance=5mm and 7mm]
 \node[ch1box] (j1) {$\Jimp^{(1)}$};
 \node[ch1box, right=7mm of j1] (s) {$\mathcal{S}_\omega$};
 \node[ch1box, right=7mm of s] (u1) {$\mathbf{u}_1$};
 \node[ch1box, below=5mm of j1] (j2) {$\Jimp^{(2)}$};
 \node[ch1box, below=5mm of s] (u2) {$\mathbf{u}_2$};
 \node[ch1box, below=5mm of u1] (bil) {Bilinear\\pairing};
 \draw[ch1flow] (j1) -- (s) -- (u1);
 \draw[ch1flow] (j2) -- (s) -- (u2);
 \draw[ch1flow, dashed] (u1.south) -- (bil.north);
 \draw[ch1flow, dashed] (j2.east) -- (bil.west);
 \draw[ch1flow, dashed] (u2.east) -- (bil.west);
 \draw[ch1flow, dashed] (j1.east) -- (bil.west);
\end{tikzpicture}%
}
\caption{Subsection~\ref{sec:ch281} Lorentz reciprocity: source labels interchange under
\(\mathcal{S}_\omega\) when media satisfy Eq.~\eqref{eq:ch2.reciprocal-medium-class} and
\(\mathcal{B}_\Gamma\) is passive on the outgoing branch.}
\label{fig:ch2.lorentz-reciprocity-measurement chain}
\end{figure}
Figure interpretation. Sources $\Jimp^{(1)}$, $\Jimp^{(2)}$ induce fields $\mathbf{u}_1$, $\mathbf{u}_2$ through the same outgoing $\mathcal{S}_\omega$. The central bilinear node is symmetric under label interchange when media satisfy Eq.~\eqref{eq:ch2.reciprocal-medium-class} and $\mathcal{B}_\Gamma$ is passive. Dashed return arrows mark that reciprocity is proved on the difference system with identical boundary data, not by comparing unrelated forward and adjoint meshes. Projection to $\mathcal{R}_\omega$ preserves symmetry only after regime qualification Eq.~\eqref{eq:ch2.regime-component-criterion} \cite{stratton1941}. Figure~\ref{fig:ch2.lorentz-reciprocity-measurement chain} is the operator reading of
Eq.~\eqref{eq:ch2.lorentz-volume-reciprocity}: the bilinear functional is symmetric under
label interchange $(\Jimp^{(1)},\mathbf{u}_1)\leftrightarrow(\Jimp^{(2)},\mathbf{u}_2)$ once
boundary terms cancel \cite{stratton1941}.
For vacuum observation points $\mathbf{r}_1$, $\mathbf{r}_2$ with $k_0|\mathbf{r}_1-\mathbf{r}_2|=2.5$, Eq.~\eqref{eq:ch2.dyadic-reciprocity-swap} requires $\overline{\mathbf{G}}_{E,0}(\mathbf{r}_1,\mathbf{r}_2)^{\mathsf{T}}=\overline{\mathbf{G}}_{E,0}(\mathbf{r}_2,\mathbf{r}_1)$ on the outgoing branch. Transverse block entries agree to machine precision; reactive $R^{-2}$, $R^{-3}$ blocks differ by sign under interchange yet cancel in the bilinear contraction when both sides use identical outgoing selectors \cite{stratton1941}. A practical verification proceeds in three steps: confirm the outgoing branch
Eq.~\eqref{eq:ch2.outgoing-green-branch}; test swap symmetry on transverse dyad blocks;
and project to $\mathcal{R}_\omega$ before port calibration \cite{jackson1999}.
\phantomsection\label{thm:ch2.lorentz-reciprocity-operator}
Let \(\mathcal{S}_\omega:\mathcal{J}_\Omega\to\mathcal{F}_\Omega\) be the outgoing solution
map Eq.~\eqref{eq:ch2.maxwell-solution-map} on media satisfying
Eq.~\eqref{eq:ch2.reciprocal-medium-class} with passive \(\mathcal{B}_\Gamma\)
Eq.~\eqref{eq:ch2.passive-impedance-class}. For admissible
\(\Jimp^{(1)},\Jimp^{(2)}\in\mathcal{J}_\Omega\), write
\(\mathbf{u}_k=\mathcal{S}_\omega[\Jimp^{(k)}]\), \(k=1,2\). Then:
\begin{enumerate}
\item[(i)] Volume reciprocity Eq.~\eqref{eq:ch2.lorentz-volume-reciprocity} holds with
vanishing boundary term on admissible outgoing radiation conditions.
\item[(ii)] Operator form: with \(\mathbf{u}_k=(\E_k,\H_k)\) the Maxwell pair induced by
\(\Jimp^{(k)}\) and \(\langle\cdot,\cdot\rangle_{\mathcal{U}}\) the energy inner product on
\(\mathcal{F}_{\mathrm{SF}}\),
\begin{equation}
\label{eq:ch2.lorentz-reciprocity-operator}
\langle \mathbf{u}_1,\,\mathcal{S}_\omega[\Jimp^{(2)}]\rangle_{\mathcal{U}}
=
\langle \mathbf{u}_2,\,\mathcal{S}_\omega[\Jimp^{(1)}]\rangle_{\mathcal{U}}
\end{equation}
when the pairings are defined through impressed-current forcing functionals
Eq.~\eqref{eq:ch2.forcing-functional-operator} and field norm
Eq.~\eqref{eq:ch2.field-energy-norm} on \(\mathcal{F}_{\mathrm{SF}}\).
\item[(iii)] Dyadic realization: outgoing electric dyads satisfy
Eq.~\eqref{eq:ch2.dyadic-reciprocity-swap}.
\item[(iv)] Radiation reciprocity: projecting both sides of (ii) through the self-adjoint
outgoing projector \(\Pi_{\mathrm{rad}}\) yields symmetric coupling on the radiative subspace
\(\mathcal{F}_{\mathrm{rad}}\):
\begin{equation}
\label{eq:ch2.radiation-map-reciprocity}
\langle \mathcal{R}_\omega[\Jimp^{(1)}],\,\mathcal{R}_\omega[\Jimp^{(2)}]\rangle_{\mathcal{U}}
=
\langle \mathcal{R}_\omega[\Jimp^{(2)}],\,\mathcal{R}_\omega[\Jimp^{(1)}]\rangle_{\mathcal{U}}
\end{equation}
on \(\mathcal{F}_{\mathrm{rad}}\) with \(\mathcal{R}_\omega=\Pi_{\mathrm{rad}}\circ\mathcal{S}_\omega\)
Eq.~\eqref{eq:ch2.radiation-operator-def}.
\end{enumerate}
Part (i) follows from Eq.~\eqref{eq:ch2.maxwell-adjoint-pairing} and
Eq.~\eqref{eq:ch2.reciprocal-medium-class} as in the derivation above; outgoing and passive
boundary classes from Subsection~\ref{sec:ch214} cancel the surface integrand
\cite{stratton1941}. Part (ii) identifies the volume impressed-current pairing
with \(\langle\cdot,\cdot\rangle_{\mathcal{U}}\) on \(\mathcal{F}_{\mathrm{SF}}\) through
Eq.~\eqref{eq:ch2.dual-source-pairing} and stability
Eq.~\eqref{eq:ch2.solution-map-stability}, developing Eq.
Eq.~\eqref{eq:ch2.lorentz-reciprocity} \cite{stratton1941}. Part (iii): for fixed
\(\mathbf{r}\), Eq.~\eqref{eq:ch2.electric-dyadic-convolution} and part (i) with impulsive
\(\Jimp^{(k)}\) aligned with Cartesian unit vectors yield column swap
\(\overline{\mathbf{G}}_{E}(\mathbf{r},\mathbf{r}')=
\overline{\mathbf{G}}_{E}^{\mathsf{T}}(\mathbf{r}',\mathbf{r})\) on the outgoing branch
\cite{stratton1941}. Uniqueness exterior radiative uniqueness (Subsection~\ref{sec:ch243})
extends the swap to all admissible \(\mathbf{r},\mathbf{r}'\) \cite{stratton1941}. Part (iv)
applies \(\Pi_{\mathrm{rad}}\) to both sides of (ii); \(\Pi_{\mathrm{rad}}\) is self-adjoint
on \(\mathcal{F}_{\mathrm{rad}}\) under Eq.~\eqref{eq:ch2.field-energy-norm} when outgoing
selection is fixed \cite{stratton1941}. Equation~\eqref{eq:ch2.radiation-map-reciprocity} is the export-projected
form of the VNA symmetry: coupling reported on \(\mathcal{F}_{\mathrm{rad}}\) is unchanged when
port labels swap, while reactive near-zone terms dropped by \(\Pi_{\mathrm{rad}}\) need not be
symmetric even on reciprocal hardware \cite{stratton1941}.
The bilinear functional
\(\mathfrak{B}_\omega(\Jimp^{(1)},\Jimp^{(2)})=
\langle \mathcal{S}_\omega[\Jimp^{(1)}],\,\mathcal{S}_\omega[\Jimp^{(2)}]\rangle_{\mathcal{U}}\)
is symmetric on \(\mathcal{J}_\Omega\times\mathcal{J}_\Omega\) under reciprocal media
\cite{stratton1941}. Algebraically, \(\mathcal{S}_\omega\) is not self-adjoint
on \(\mathcal{J}_\Omega\) with source pairing Eq.~\eqref{eq:ch2.dual-source-pairing}; reciprocity
instead couples the source space to the field space through the \emph{interchanged} labels in
Eq.~\eqref{eq:ch2.lorentz-reciprocity-operator} \cite{stratton1941}. Physically, the far-zone
pattern produced by \(\Jimp^{(1)}\) tested against the field of \(\Jimp^{(2)}\) equals the
swap configuration: a receive pattern computed with \(\Jimp^{(2)}\) as probe matches a
transmit map driven by \(\Jimp^{(1)}\) when hardware and media are reciprocal
\cite{harrington2001}.
Mathematical role of Eq.~\eqref{eq:ch2.radiation-map-reciprocity}. Equation~\eqref{eq:ch2.radiation-map-reciprocity} restricts Eq.~\eqref{eq:ch2.lorentz-reciprocity-operator}
to \(\mathcal{F}_{\mathrm{rad}}\), the subspace on which flux meters and angular patterns are
declared \cite{stratton1941}. Because \(\Pi_{\mathrm{rad}}\) annihilates
non-export components Eq.~\eqref{eq:ch2.rad-null-on-evanescent}, reactive near-zone overlap
does not enter Eq.~\eqref{eq:ch2.radiation-map-reciprocity} even when present in full
Eq.~\eqref{eq:ch2.lorentz-reciprocity-operator} \cite{stratton1941}. regime qualification
Eq.~\eqref{eq:ch2.regime-component-criterion} must still be satisfied before interpreting
Eq.~\eqref{eq:ch2.radiation-map-reciprocity} as a transport channel
\cite{harrington2001}. \phantomsection\label{prop:ch2.dyadic-reciprocity-home}
Under Eq.~\eqref{eq:ch2.reciprocal-medium-class} with outgoing radiation condition, the electric dyad
\(\overline{\mathbf{G}}_{E,\omega}^{\mathrm{(out)}}\) in
Eq.~\eqref{eq:ch2.electric-dyadic-convolution} satisfies
Eq.~\eqref{eq:ch2.dyadic-reciprocity-swap} for all \(\mathbf{r},\mathbf{r}'\) in the
evaluation domain of outgoing source-to-field integrals of Subsection~\ref{sec:ch251}.
Apply Lorentz reciprocity for the solution map (Subsection~\ref{sec:ch281})(iii); this is the proof deferred from
Subsection~\ref{sec:ch231} \cite{stratton1941}. \phantomsection\label{prop:ch2.reciprocity-null-invariance}
If \(\Jimp^{(1)}\) is arbitrary and \(\Jimp^{(2)}=\Jimp^{(1)}+\Jimp^{\mathrm{(NR)}}\) with
\(\Jimp^{\mathrm{(NR)}}\in\ker\mathcal{R}_\omega\), then
Eq.~\eqref{eq:ch2.radiation-map-reciprocity} is unchanged when either argument is replaced by
\(\Jimp^{(2)}\) on the radiative pairing
\(\mathcal{R}_\omega[\Jimp^{(2)}]=\mathcal{R}_\omega[\Jimp^{(1)}]\) by
Eq.~\eqref{eq:ch2.nr-superposition-law} \cite{stratton1941}. Full
Eq.~\eqref{eq:ch2.lorentz-reciprocity-operator} may change by reactive content through
\(\mathcal{S}_\omega[\Jimp^{\mathrm{(NR)}}]\) Eq.~\eqref{eq:ch2.nr-image-of-null-space}
\cite{harrington2001}.
On the outgoing branch, Eq.~\eqref{eq:ch2.lorentz-reciprocity-operator} is consistent with
resolvent symmetry of the truncated Maxwell operator on reciprocal media:
\begin{equation}
\label{eq:ch2.resolvent-reciprocity-swap}
\langle
\mathcal{L}_\omega^{-1}[\mathbf{s}^{(1)}],\,\mathbf{s}^{(2)}
\rangle_{\mathcal{Y}}
=
\langle
\mathcal{L}_\omega^{-1}[\mathbf{s}^{(2)}],\,\mathbf{s}^{(1)}
\rangle_{\mathcal{Y}}
\end{equation}
for compatible forcing pairs \(\mathbf{s}^{(k)}\) when \(\mathcal{L}_\omega\) is symmetric on
\(\mathcal{D}(\mathcal{L}_\omega)\) and \(\mathcal{B}_\Gamma\) is passive
\cite{stratton1941}. Equation~\eqref{eq:ch2.resolvent-reciprocity-swap} is not
re-proved here; it is the operator backbone behind
Eq.~\eqref{eq:ch2.integral-resolvent-equivalence} and Lorentz reciprocity for the solution map (Subsection~\ref{sec:ch281})(ii) \cite{stratton1941}. Numerical MoM matrices
assembled from Eq.~\eqref{eq:ch2.dyadic-reciprocity-swap} inherit swap symmetry only when
quadrature, outgoing branch, and material sampling preserve
Eq.~\eqref{eq:ch2.reciprocal-medium-class} \cite{harrington2001}.
Loss, anisotropy, and broken reciprocity. Complex symmetric \(\overline{\overline{\varepsilon}}\), \(\overline{\overline{\mu}}\) with
\(\Im\{\overline{\overline{\varepsilon}}\}\succeq 0\), \(\Im\{\overline{\overline{\mu}}\}\succeq 0\)
preserve Eqs.~\eqref{eq:ch2.lorentz-volume-reciprocity}--\eqref{eq:ch2.lorentz-reciprocity-operator}
with conjugation on one slot \cite{stratton1941}. Magneto-electric coupling
\(\overline{\overline{\xi}}\neq\mathbf{0}\) breaks Eq.~\eqref{eq:ch2.dyadic-reciprocity-swap}
and requires non-reciprocal Green operators \cite{stratton1941}. Active boundaries
with \(\Re\{Z_s\}<0\) on portions of \(\Gamma_Z\) violate
Eq.~\eqref{eq:ch2.passive-impedance-class} and invalidate Lorentz reciprocity for the solution map (Subsection~\ref{sec:ch281}) \cite{stratton1941}. Incoming rather
than outgoing radiation condition selects a different branch; swap identities hold only when both excitations
use the same branch \cite{stratton1941}. Lorentz reciprocity for the solution map (Subsection~\ref{sec:ch281}) assumes fixed \(\omega\), linear media
Eq.~\eqref{eq:ch2.reciprocal-medium-class}, passive \(\mathcal{B}_\Gamma\), outgoing radiation condition,
and admissible \(\mathcal{J}_\Omega\) \cite{stratton1941}. It does not assert
reciprocity of raw near-zone reactive scans without specifying whether \(\Pi_{\mathrm{rad}}\)
has been applied \cite{stratton1941}. It does not remove \(\ker\mathcal{R}_\omega\)
null-space invariance of radiative reciprocity (Subsection~\ref{sec:ch281}); invisible currents alter full
\(\mathcal{S}_\omega\) pairings while leaving radiative reciprocity unchanged
\cite{harrington2001}. Nonlinear, time-varying, or moving media require the causal chain of
Subsection~\ref{sec:ch223} \cite{stratton1941}.
Subsection~\ref{sec:ch281} established Lorentz reciprocity on \(\mathcal{S}_\omega\) and
\(\mathcal{R}_\omega\) through volume and operator identities
Eqs.~\eqref{eq:ch2.lorentz-volume-reciprocity}--\eqref{eq:ch2.radiation-map-reciprocity}
\cite{stratton1941}. The port-swap symmetry observed on reciprocal hardware is now available as a
constraint on continuous source and field pairings; the next step is to express that constraint on
the modal coordinates fixed in Section~\ref{sec:ch26}.
\subsection{Orthogonality and coupling of radiative modes}
\label{sec:ch282}
% HARRINGTON-OPEN:sec:ch282
A standard-gain horn on a compact-range turntable fixes the laboratory scene for `Orthogonality and coupling of radiative modes'': frequency, probe standoff, and wattmeter closure on $\Gamma_R$ precede the formal displays below \cite{hansen1988}. Maxwell's retarded export to $\Gamma_R$ remains the governing constraint \cite{maxwell1873}.
Bistatic radar and compact-range scattering matrices rely on reciprocity to relate monostatic
calibrations to full \(S\)-parameter sets \cite{stratton1941}. The adjoint theory extends
one-port antenna reciprocity to continuous source and trace pairings \cite{devaney2004}.
Plane-wave decompositions used in reflector and radar-cross-section campaigns inherit reciprocity
constraints on amplitudes before any angular filter is applied \cite{hansen1988}. When
reciprocity appears broken after probe correction, the fault usually lies in inconsistent
impressed-current normalization or a missing outgoing tail in the Green operator, not in
Maxwell theory itself \cite{hansen1988}. Modal orthogonality and coupling on
\(\operatorname{ran}\mathcal{R}_\omega\) organize how radiative modes exchange power between
ports; reciprocity fixes the bilinear pairing before angular filters act
\cite{stratton1941}.
\medskip
Subsection~\ref{sec:ch282} gives the Chapter~\ref{ch:02} definitive account of orthogonality and
coupling of radiative modes on the spectral coordinates fixed in Subsection~\ref{sec:ch261},
developing Eqs.~\eqref{eq:ch2.modal-orthogonality}
and~\eqref{eq:ch2.radiative-coupling-bilinear} \cite{stratton1941}. Radiation eigenpairs
Eq.~\eqref{eq:ch2.radiation-eigenpair-def}, modal projectors
Eq.~\eqref{eq:ch2.modal-projector-def}, source amplitudes
Eq.~\eqref{eq:ch2.source-modal-amplitude}, spectral decomposition of the radiation operator
(Subsection~\ref{sec:ch261}), and biorthogonality
Eq.~\eqref{eq:ch2.radiation-biorthogonality} are presupposed; eigenpair construction and
reciprocity-based solvability are not re-proved \cite{stratton1941}. Explicit vector
spherical harmonic orthogonality belongs to \ref{ch:03}
Subsection~3.4.3; continuous-spectrum normalization belongs to Subsection~3.5.2, both with USE of
the present operator targets \cite{stratton1941}. Transmit--receive systems are deferred to
Subsection~\ref{sec:ch283} \cite{harrington2001}.
Assumptions. Fix \(\omega>0\), outgoing fields in \(\mathcal{F}_{\mathrm{SF}}\), radiative images in
\(\mathcal{F}_{\mathrm{rad}}\), and orthonormal radiation eigenmodes
\(\{\mathbf{u}_n\}_{n\in\mathcal{N}_{\mathrm{rad}}}\) from spectral decomposition of the radiation operator (Subsection~\ref{sec:ch261}) on the declared domain class
\cite{stratton1941}. Media satisfy reciprocal class
Eq.~\eqref{eq:ch2.reciprocal-medium-class} and passive boundary
Eq.~\eqref{eq:ch2.passive-impedance-class} so Lorentz reciprocity for the solution map (Subsection~\ref{sec:ch281}) applies \cite{stratton1941}. Adjoint modes
\(\widetilde{\mathbf{u}}_n\) are defined through
Eq.~\eqref{eq:ch2.radiation-biorthogonality} and coincide with \(\mathbf{u}_n\) on
\(\mathcal{F}_{\mathrm{rad}}\) under reciprocity-qualified pairings below
\cite{stratton1941}. regime qualification Eq.~\eqref{eq:ch2.regime-operator-screen}
governs transport interpretation of coupling reports \cite{harrington2001}.
On the orthonormal family from spectral decomposition of the radiation operator
(Subsection~\ref{sec:ch261}), distinct exporting modes are declared orthogonal with respect to the
energy inner product $\langle\cdot,\cdot\rangle_{\mathcal{U}}$ on $\mathcal{F}_{\mathrm{rad}}$.
The locked orthogonality statement reads
% PHYSICAL-MEMORY:sec:ch282
\begin{physicalvignette}[Beth torsion balance in a microwave anechoic chamber]
In a microwave anechoic chamber fitted with a Beth torsion balance, technicians interrogate whether a rotating field imparts a measurable torque on a suspended wire loop \cite{beth1936}. The experiment is delicate---vibration isolation, chamber leakage, and probe polarization must be controlled---yet the outcome is blunt: only the radiative fraction of the field that survives outgoing selection contributes to the mechanical reading expected from angular-momentum flux \cite{stratton1941}.
Reciprocity on the same bench swaps transmit and receive chains through a calibrated horn pair. If the torque report changes when only the probe polarization basis is rotated, the fault is almost always in the measurement definition, not in Maxwell's equations \cite{harrington2001}. Such benches keep representation labels tied to flux, not to post-processing convenience \cite{pal2026}.
Subsection~\ref{sec:ch282} states the reciprocity relation that this balance scene exercises \cite{stratton1941}.
\end{physicalvignette}
\begin{equation}
\label{eq:ch2.radiation-modal-orthogonality}
\langle \mathbf{u}_m,\,\mathbf{u}_n\rangle_{\mathcal{U}}
=
\delta_{mn},
\qquad
m,n\in\mathcal{N}_{\mathrm{rad}},
\end{equation}
with \(\|\mathbf{u}_n\|_{\mathcal{U}}=1\) already enforced in
Eq.~\eqref{eq:ch2.radiation-eigenpair-def} \cite{stratton1941}.
Equation~\eqref{eq:ch2.radiation-modal-orthogonality} is the Subsection~\ref{sec:ch282} locked definition of
upgrade of the first line of Eq.~\eqref{eq:ch2.modal-orthogonality}; it
specializes projector orthogonality \(P_m P_n=\delta_{mn}P_n\) from
Eq.~\eqref{eq:ch2.modal-projector-def} to the chosen orthonormal basis
\cite{stratton1941}. Algebraically, Eq.~\eqref{eq:ch2.radiation-modal-orthogonality} makes
\(\{\mathbf{u}_n\}\) an orthonormal frame on \(\mathcal{F}_{\mathrm{rad}}\) for energy inner
product Eq.~\eqref{eq:ch2.field-energy-norm}; physically, distinct exporting modes carry
non-overlapping far-zone energy content on \(S^{2}\) when
\(\mathcal{D}_{J,\mathrm{reg}}=1\) \cite{jackson1999}. Explicit angular
integrals producing Eq.~\eqref{eq:ch2.radiation-modal-orthogonality} for harmonic bases are
not tabulated here \cite{stratton1941}.
For canonical eigen-sources \(\Jimp^{(n)}\) in Eq.~\eqref{eq:ch2.radiation-eigenpair-def},
define the \emph{modal coupling coefficient}
\begin{equation}
\label{eq:ch2.modal-coupling-coefficient}
C_{mn}
=
\langle \widetilde{\mathbf{u}}_m,\,\mathcal{R}_\omega[\Jimp^{(n)}]\rangle_{\mathcal{U}},
\qquad
m,n\in\mathcal{N}_{\mathrm{rad}},
\end{equation}
with adjoint modes from Eq.~\eqref{eq:ch2.radiation-biorthogonality}
\cite{stratton1941}. Equation~\eqref{eq:ch2.modal-coupling-coefficient} is the
Subsection~\ref{sec:ch282} development of the second line of Eq.
Eq.~\eqref{eq:ch2.modal-orthogonality}. When
Eq.~\eqref{eq:ch2.radiation-biorthogonality} holds,
\begin{equation}
\label{eq:ch2.modal-coupling-diagonal}
C_{mn}=\lambda_n\,\delta_{mn}
\end{equation}
for exporting eigenpairs with \(\lambda_n\neq 0\); null modes
\(n\in\mathcal{N}_{\mathrm{null}}\) Eq.~\eqref{eq:ch2.modal-null-index-set} contribute
\(C_{mn}=0\) for all \(m\) \cite{stratton1941}. Algebraically,
Eq.~\eqref{eq:ch2.modal-coupling-diagonal} diagonalizes coupling in an eigenmode basis;
physically, orthogonal radiation modes do not co-excite one another through
\(\mathcal{R}_\omega\) when sources align with eigenpairs---cross-talk appears only from
superpositions or degenerate subspaces \cite{harrington2001}.
For arbitrary admissible \(\Jimp^{(1)},\Jimp^{(2)}\in\mathcal{J}_\Omega\), define
\begin{equation}
\label{eq:ch2.radiative-coupling-bilinear}
\mathfrak{C}_\omega\!\big(\Jimp^{(1)},\Jimp^{(2)}\big)
=
\langle \mathcal{R}_\omega[\Jimp^{(1)}],\,\mathcal{R}_\omega[\Jimp^{(2)}]\rangle_{\mathcal{U}}
\quad\text{on}\ \mathcal{F}_{\mathrm{rad}},
\end{equation}
developing Eq.~\eqref{eq:ch2.radiative-coupling-bilinear}
\cite{stratton1941}. Equation~\eqref{eq:ch2.radiative-coupling-bilinear}
measures export overlap on \(\operatorname{ran}\mathcal{R}_\omega\), not raw reactive overlap
of \(\mathcal{S}_\omega[\Jimp^{(k)}]\) \cite{stratton1941}. Under Lorentz reciprocity for the solution map (Subsection~\ref{sec:ch281})(iv),
\begin{equation}
\label{eq:ch2.coupling-bilinear-symmetry}
\mathfrak{C}_\omega\!\big(\Jimp^{(1)},\Jimp^{(2)}\big)
=
\mathfrak{C}_\omega\!\big(\Jimp^{(2)},\Jimp^{(1)}\big)
\end{equation}
on reciprocal media \cite{stratton1941}. Reactive near-zone terms in
\(\Pi_{\mathrm{ev}}\mathcal{S}_\omega[\Jimp]\) do not enter
Eq.~\eqref{eq:ch2.radiative-coupling-bilinear} through
Eq.~\eqref{eq:ch2.rad-null-on-evanescent} \cite{stratton1941}.
Let \(c_n[\Jimp^{(k)}]\) denote source amplitudes Eq.~\eqref{eq:ch2.source-modal-amplitude}
for \(k=1,2\). When the modal sum converges in \(\mathcal{U}\),
\begin{equation}
\label{eq:ch2.modal-coupling-sum}
\mathfrak{C}_\omega\!\big(\Jimp^{(1)},\Jimp^{(2)}\big)
=
\sum_{m,n\in\mathcal{N}_{\mathrm{rad}}}
\lambda_m\lambda_n\,
c_m[\Jimp^{(1)}]^{*}\,
c_n[\Jimp^{(2)}]\,
\langle \mathbf{u}_m,\mathbf{u}_n\rangle_{\mathcal{U}},
\end{equation}
using Eq.~\eqref{eq:ch2.radiation-map-modal-sum} and linearity
Eq.~\eqref{eq:ch2.radiation-linearity} \cite{stratton1941}. With
Eq.~\eqref{eq:ch2.radiation-modal-orthogonality},
\begin{equation}
\label{eq:ch2.modal-coupling-sum-diagonal}
\mathfrak{C}_\omega\!\big(\Jimp^{(1)},\Jimp^{(2)}\big)
=
\sum_{n\in\mathcal{N}_{\mathrm{rad}}}
|\lambda_n|^{2}\,
c_n[\Jimp^{(1)}]^{*}\,
c_n[\Jimp^{(2)}],
\end{equation}
for real nonnegative exporting gains \(\lambda_n>0\) \cite{stratton1941}. Equation
\eqref{eq:ch2.modal-coupling-sum-diagonal} shows that coupling is a weighted inner product
of modal excitation vectors on the exporting sector; geometric constraints
geometric constraints on radiative spectra (Subsection~\ref{sec:ch264}) may force individual \(c_n=0\) without
breaking Eq.~\eqref{eq:ch2.coupling-bilinear-symmetry} \cite{harrington2001}.
\begin{figure}[t]
\centering
\ChOneFigBlock{%
\begin{tikzpicture}[ch1fig, node distance=5mm and 7mm]
 \node[ch1box] (j1) {$\Jimp^{(1)}$};
 \node[ch1box, right=7mm of j1] (r) {$\mathcal{R}_\omega$};
 \node[ch1box, right=7mm of r] (u) {$\sum_n \lambda_n c_n \mathbf{u}_n$};
 \node[ch1box, below=5mm of j1] (j2) {$\Jimp^{(2)}$};
 \node[ch1box, below=5mm of u] (c) {$\mathfrak{C}_\omega$};
 \node[ch1box, left=7mm of c] (orth) {Eq.~\eqref{eq:ch2.radiation-modal-orthogonality}};
 \draw[ch1flow] (j1.east) -- (j1.east |- r.west);
 \draw[ch1flow] (j1.east |- r.west) -- (r.west);
 \draw[ch1flow] (j2.east) -- (j2.east |- r.west);
 \draw[ch1flow] (j2.east |- r.west) -- (r.west);
 \draw[ch1flow] (r.east) -- (r.east |- u.west);
 \draw[ch1flow] (r.east |- u.west) -- (u.west);
 \draw[ch1flow] (u.south) -- (u.south |- c.north);
 \draw[ch1flow] (u.south |- c.north) -- (c.north);
 \draw[ch1flow] (orth.west) -- (orth.west |- c.east);
 \draw[ch1flow] (orth.west |- c.east) -- (c.east);
\end{tikzpicture}%
}
\caption{Subsection~\ref{sec:ch282} modal coupling measurement sequence: sources map through
\(\mathcal{R}_\omega\) to orthonormal modes Eq.~\eqref{eq:ch2.radiation-modal-orthogonality};
bilinear coupling Eq.~\eqref{eq:ch2.radiative-coupling-bilinear} factorizes through modal
amplitudes Eq.~\eqref{eq:ch2.modal-coupling-sum}.}
\label{fig:ch2.modal-coupling-measurement chain}
\end{figure}
Figure interpretation. Two impressed sources $\Jimp^{(1)}$, $\Jimp^{(2)}$ share the same radiation map $\mathcal{R}_\omega$ and expand on orthonormal modes Eq.~\eqref{eq:ch2.radiation-modal-orthogonality}. The coupling node $\mathfrak{C}_\omega$ measures export overlap on $\mathcal{F}_{\mathrm{rad}}$, not raw near-zone reactive cross terms. Orthogonality diagonalizes coupling when sources align with eigenpairs; cross-talk appears only from superposition or degenerate subspaces. Modes with $\lambda_n=0$ feed the null sector without contributing to $\mathfrak{C}_\omega$ on the range \cite{stratton1941}. Figure~\ref{fig:ch2.modal-coupling-measurement chain} separates orthogonality of mode shapes from
bilinear coupling of impressed sources: Eq.~\eqref{eq:ch2.radiation-modal-orthogonality} is
intrinsic to \(\mathcal{F}_{\mathrm{rad}}\); Eq.~\eqref{eq:ch2.radiative-coupling-bilinear}
depends on both sources \cite{stratton1941}. On a reciprocal two-mode truncation with $\lambda_1=0.42$, $\lambda_2=0.31$ and orthonormal $\mathbf{u}_1$, $\mathbf{u}_2$, excite $\Jimp^{(1)}=c_1\mathbf{J}^{(1)}_{\mathrm{imp}}$ and probe with $\Jimp^{(2)}=c_2\mathbf{J}^{(2)}_{\mathrm{imp}}$. Eq.~\eqref{eq:ch2.modal-coupling-sum-diagonal} gives $\mathfrak{C}_\omega(\Jimp^{(1)},\Jimp^{(2)})=\lambda_1\lambda_2 c_1^* c_2$. Swapping labels leaves the scalar unchanged, so normalized $S_{21}=S_{12}$ on the exporting sector. Adding $\Jimp^{\mathrm{(NR)}}$ to either port shifts $\mathcal{S}_\omega[\Jimp]$ near $\Omega_s$ but leaves $S_{21}$ fixed---the walkthrough confirms reciprocity on $\operatorname{ran}\mathcal{R}_\omega$, not on raw near-field coupling matrices \cite{stratton1941}.
\phantomsection\label{thm:ch2.radiation-modal-orthogonality-coupling}
Let \(\{\mathbf{u}_n,\Jimp^{(n)},\lambda_n\}\) be radiation eigenpairs
Eq.~\eqref{eq:ch2.radiation-eigenpair-def} forming an orthonormal basis of
\(\mathcal{F}_{\mathrm{rad}}\) on the declared domain class, with reciprocal media
Eq.~\eqref{eq:ch2.reciprocal-medium-class}. Then:
\begin{enumerate}
\item[(i)] Orthogonality Eq.~\eqref{eq:ch2.radiation-modal-orthogonality} holds on
\(\mathcal{N}_{\mathrm{rad}}\).
\item[(ii)] Coupling coefficients satisfy Eq.~\eqref{eq:ch2.modal-coupling-diagonal} and
biorthogonality Eq.~\eqref{eq:ch2.radiation-biorthogonality}.
\item[(iii)] For all admissible \(\Jimp^{(1)},\Jimp^{(2)}\),
Eqs.~\eqref{eq:ch2.radiative-coupling-bilinear}--\eqref{eq:ch2.modal-coupling-sum-diagonal}
hold, with symmetry Eq.~\eqref{eq:ch2.coupling-bilinear-symmetry} from Lorentz reciprocity for the solution map (Subsection~\ref{sec:ch281}).
\item[(iv)] Adding \(\Jimp^{\mathrm{(NR)}}\in\ker\mathcal{R}_\omega\) to either source leaves
Eq.~\eqref{eq:ch2.radiative-coupling-bilinear} unchanged.
\end{enumerate}
Part (i) is normalization and completeness from spectral decomposition of the radiation operator (Subsection~\ref{sec:ch261}) and
Eq.~\eqref{eq:ch2.radiation-spectral-expansion} \cite{stratton1941}. Part (ii) substitutes
\(\mathcal{R}_\omega[\Jimp^{(n)}]=\lambda_n\mathbf{u}_n\) into
Eq.~\eqref{eq:ch2.modal-coupling-coefficient} and applies
Eq.~\eqref{eq:ch2.radiation-biorthogonality} \cite{stratton1941}. Part (iii)
expands \(\mathcal{R}_\omega[\Jimp^{(k)}]\) via
Eq.~\eqref{eq:ch2.radiation-map-modal-sum}; symmetry is Lorentz reciprocity for the solution map (Subsection~\ref{sec:ch281})(iv) \cite{stratton1941}. Part (iv) follows from
\(\mathcal{R}_\omega[\Jimp^{\mathrm{(NR)}}]=\mathbf{0}\) and
Eq.~\eqref{eq:ch2.nr-superposition-law} \cite{stratton1941}. The reciprocal reduction of adjoint modes proceeds as follows: under Eqs.~\eqref{eq:ch2.reciprocal-medium-class}--\eqref{eq:ch2.passive-impedance-class},
\begin{equation}
\label{eq:ch2.reciprocal-adjoint-mode-identification}
\widetilde{\mathbf{u}}_n=\mathbf{u}_n
\quad\text{on}\ \mathcal{F}_{\mathrm{rad}}
\quad\text{for exporting modes with}\ \lambda_n\neq 0,
\end{equation}
so Eq.~\eqref{eq:ch2.radiation-biorthogonality} reduces to standard orthogonality
Eq.~\eqref{eq:ch2.radiation-modal-orthogonality} with diagonal coupling
Eq.~\eqref{eq:ch2.modal-coupling-diagonal} \cite{stratton1941}. Equation
\eqref{eq:ch2.reciprocal-adjoint-mode-identification} is the operator form of reciprocity for
modal probes: receive patterns aligned with \(\mathbf{u}_n\) coincide with transmit modes
when media and boundaries are reciprocal \cite{harrington2001}. Non-reciprocal
media break Eq.~\eqref{eq:ch2.reciprocal-adjoint-mode-identification}; biorthogonality with
distinct \(\widetilde{\mathbf{u}}_n\) remains the general template imported from
Subsection~\ref{sec:ch261} \cite{stratton1941}. When \(\lambda_n=\lambda_{n'}\) for \(n\neq n'\) within a rotationally degenerate family from
Subsection~\ref{sec:ch262}, Eq.~\eqref{eq:ch2.modal-coupling-diagonal} still holds for
\emph{eigen} sources \(\Jimp^{(n)}\), but superpositions within the degenerate block carry
off-diagonal coupling in an arbitrary basis. Let \(\mathcal{I}_\ell\subset\mathcal{N}_{\mathrm{rad}}\)
index a \((2\ell+1)\)-fold degenerate family. Any orthonormal basis
\(\{\mathbf{u}_n\}_{n\in\mathcal{I}_\ell}\) spanning the block satisfies
Eq.~\eqref{eq:ch2.radiation-modal-orthogonality}, yet a fixed physical source may excite
multiple members with coefficients tied to orientation through
Eq.~\eqref{eq:ch2.radiation-rotation-commutation} \cite{stratton1941}. Define
the block coupling matrix
\begin{equation}
\label{eq:ch2.degenerate-block-coupling-matrix}
\mathbf{K}^{(\ell)}_{mn}
=
\langle \mathbf{u}_m,\,\mathcal{R}_\omega[\Jimp^{(n)}]\rangle_{\mathcal{U}},
\qquad
m,n\in\mathcal{I}_\ell,
\end{equation}
which is diagonal in the eigenvector basis of \(\Jimp^{(n)}\) but not in laboratory-fixed
coordinates \cite{stratton1941}. Reciprocity
Eq.~\eqref{eq:ch2.coupling-bilinear-symmetry} forces \(\mathbf{K}^{(\ell)}\) to be symmetric
under Lorentz reciprocity for the solution map (Subsection~\ref{sec:ch281}) \cite{stratton1941}.
Explicit Wigner coupling of block entries belongs to \ref{ch:03}
\cite{stratton1941}.
\phantomsection\label{prop:ch2.superposition-cross-coupling}
If \(\Jimp^{(1)}=\sum_{n\in\mathcal{N}_{\mathrm{rad}}}a_n\Jimp^{(n)}\) and
\(\Jimp^{(2)}=\sum_{m\in\mathcal{N}_{\mathrm{rad}}}b_m\Jimp^{(m)}\) are admissible finite
superpositions, then
\begin{equation}
\label{eq:ch2.superposition-coupling-expansion}
\mathfrak{C}_\omega\!\big(\Jimp^{(1)},\Jimp^{(2)}\big)
=
\sum_{m,n\in\mathcal{N}_{\mathrm{rad}}}
a_m^{*}b_n\,\lambda_m\lambda_n\,
\langle \mathbf{u}_m,\mathbf{u}_n\rangle_{\mathcal{U}},
\end{equation}
reducing to Eq.~\eqref{eq:ch2.modal-coupling-sum-diagonal} when \(\langle\mathbf{u}_m,\mathbf{u}_n\rangle=\delta_{mn}\)
\cite{stratton1941}. Off-diagonal terms \(m\neq n\) appear only when at least
one source superposes non-collinear modes; single-mode excitation yields diagonal coupling
\cite{jackson1999}.
Linearity Eq.~\eqref{eq:ch2.radiation-linearity} and
Eq.~\eqref{eq:ch2.radiation-map-modal-sum} \cite{stratton1941}.
\phantomsection\label{prop:ch2.observation-limited-coupling}
For observation map \(\mathcal{M}_\omega=\mathcal{O}_{\mathrm{meas}}\circ\Lambda_{\mathrm{rad}}\)
Eq.~\eqref{eq:ch2.reconstruction-observation-map}, define observed coupling
\(\mathfrak{C}_{\omega,\mathrm{obs}}=
\langle \mathcal{O}_{\mathrm{meas}}[\mathcal{R}_\omega[\Jimp^{(1)}]],\,
\mathcal{O}_{\mathrm{meas}}[\mathcal{R}_\omega[\Jimp^{(2)}]]\rangle\).
Then \(\mathfrak{C}_{\omega,\mathrm{obs}}\) is symmetric under Lorentz reciprocity for the solution map (Subsection~\ref{sec:ch281}) when \(\mathcal{O}_{\mathrm{meas}}\) is self-adjoint
on \(\mathcal{F}_{\mathrm{rad}}\), but may vanish for pairs with
\(\mathcal{R}_\omega[\Jimp^{(k)}]\in\ker\mathcal{O}_{\mathrm{meas}}\) even when
Eq.~\eqref{eq:ch2.radiative-coupling-bilinear} is nonzero \cite{stratton1941}.
Rank cap Eq.~\eqref{eq:ch2.reconstruction-rank-cap} limits the number of independently
observable coupling channels \cite{stratton1941}.
Apply Eq.~\eqref{eq:ch2.radiation-map-reciprocity} and self-adjointness of
\(\mathcal{O}_{\mathrm{meas}}\) \cite{stratton1941}. Kernel inclusion
\(\ker\mathcal{R}_\omega\subset\ker\mathcal{M}_\omega\) is
measurement null-space extension (Subsection~\ref{sec:ch271}) \cite{stratton1941}.
\begin{equation}
\label{eq:ch2.kernel-field-image-law}
\Jimp\in\ker\mathcal{R}_\omega
\ \Longrightarrow\
\mathcal{S}_\omega[\Jimp]\in\mathcal{F}_{\mathrm{SF}},
\qquad
\Pi_{\mathrm{rad}}\,\mathcal{S}_\omega[\Jimp]=\mathbf{0},
\end{equation}
Accessible modal sector Eq.~\eqref{eq:ch2.accessible-modal-sector} and angular budget
Eq.~\eqref{eq:ch2.angular-spectrum-composition-law} restrict which terms in
Eq.~\eqref{eq:ch2.modal-coupling-sum-diagonal} are observable in far-zone reports
\cite{stratton1941}. Null modes \(\lambda_n=0\) carry zero coupling diagonal
entries but may alter near reactive content through \(\mathcal{S}_\omega\)
Eq.~\eqref{eq:ch2.kernel-field-image-law} without changing
Eq.~\eqref{eq:ch2.radiative-coupling-bilinear} \cite{stratton1941}. Composition with
null-space invariance of radiative reciprocity (Subsection~\ref{sec:ch281}) and
orthogonality and coupling of radiative modes (Subsection~\ref{sec:ch282})(iv) explains why invisible
currents do not break reciprocity-qualified coupling on \(\operatorname{ran}\mathcal{R}_\omega\)
\cite{harrington2001}.
Orthogonality and coupling of radiative modes (Subsection~\ref{sec:ch282}) requires linear reciprocal media
at fixed \(\omega\), outgoing radiation condition, domain-class completeness of
\(\{\mathbf{u}_n\}\), and convergent modal sums \cite{stratton1941}. Loss
moves \(\lambda_n\) into the complex plane; Eq.~\eqref{eq:ch2.modal-coupling-sum-diagonal}
generalizes with complex conjugation on one amplitude slot \cite{stratton1941}. Truncated
modal books on cavities do not substitute for exterior continuous spectra
Eq.~\eqref{eq:ch2.exterior-continuous-spectral-integral} when coupling far-zone channels
\cite{jackson1999}. regime qualification Eq.~\eqref{eq:ch2.regime-component-criterion} is
mandatory before interpreting Eq.~\eqref{eq:ch2.radiative-coupling-bilinear} as transport
\cite{stratton1941}.
Subsections~\ref{sec:ch281}--\ref{sec:ch282} established Lorentz reciprocity, radiative bilinear
coupling Eq.~\eqref{eq:ch2.radiative-coupling-bilinear}, and modal orthogonality on
\(\mathcal{F}_{\mathrm{rad}}\) \cite{stratton1941}. The final reciprocity link ties those
operator statements to finite ports: transmit and receive maps on reciprocal hardware must
interchange when excitation and probe roles are swapped.
\subsection{Reciprocity in radiation and reception}
\label{sec:ch283}
% HARRINGTON-OPEN:sec:ch283
A standard-gain horn on a compact-range turntable fixes the laboratory scene for `Reciprocity in radiation and reception'': frequency, probe standoff, and wattmeter closure on $\Gamma_R$ precede the formal displays below \cite{hansen1988}. Maxwell's retarded export to $\Gamma_R$ remains the governing constraint \cite{maxwell1873}.
Two-port antenna measurements are the everyday face of reciprocity. A horn coupled to a compact
range, a patch fed from a coax port, or a dual-polarized reflector feed all reduce---after
calibration---to the same question the vector network analyzer answers: does the complex coupling
from port~$A$ to port~$B$ equal the reversed coupling when cables, standards, and probe
normalization are exchanged? Subsection~\ref{sec:ch283} gives the Chapter~\ref{ch:02} definitive
account of that interchange in operator form: transmit--receive operator symmetry, two-port
coupling bilinears, and measurement-qualified reciprocity on
\(\operatorname{ran}\mathcal{R}_\omega\), developing Eq.
Eq.~\eqref{eq:ch2.transmit-receive-swap} \cite{stratton1941}. Observation
composition Eq.~\eqref{eq:ch2.reconstruction-observation-map}, finite-rank readouts
Eq.~\eqref{eq:ch2.finite-rank-observation}, rank cap
Eq.~\eqref{eq:ch2.reconstruction-rank-cap}, and null-space laws
Eqs.~\eqref{eq:ch2.reconstruction-null-space}--\eqref{eq:ch2.nr-superposition-law} are
presupposed; inverse-source regularization is not re-developed \cite{stratton1941}.
Plane-wave angular-spectrum filters belong to Section~\ref{sec:ch29}; explicit port circuit
models are not tabulated \cite{stratton1941}.
Assumptions. Fix \(\omega>0\), reciprocal media Eq.~\eqref{eq:ch2.reciprocal-medium-class}, passive
boundary Eq.~\eqref{eq:ch2.passive-impedance-class}, outgoing radiation condition, and admissible sources
\(\Jimp^{\mathrm{(tx)}},\Jimp^{\mathrm{(rx)}}\in\mathcal{J}_\Omega\) on ports \(A\) and \(B\)
\cite{stratton1941}. Each port carries a finite-rank observation functional
\(\mathcal{O}_{\mathrm{meas}}^{(P)}:\mathcal{F}_{\mathrm{rad}}\to\mathbb{C}\) or
\(\mathbb{C}^{N_P}\) of the form Eq.~\eqref{eq:ch2.finite-rank-observation}
\cite{stratton1941}. regime qualification Eq.~\eqref{eq:ch2.regime-operator-screen} and transport
qualification Eq.~\eqref{eq:ch2.transport-qualified-radiation} govern when port outputs may be
reported as radiative transport \cite{harrington2001}. The same \(\mathcal{B}_\Gamma\), outgoing
branch, and material data are used for both excitations \cite{stratton1941}.
For port \(P\in\{A,B\}\), define the \emph{port radiation--reception operator}
% PHYSICAL-MEMORY:sec:ch283
% PHYSICAL-MEMORY:sec:ch283-b2
\begin{physicalvignette}[Corner reflector for radar cross-section calibration]
A second bench observation matters on Open-ended WR-90 waveguide on a gantry. When technicians swap only the receive probe while holding the transmit chain fixed, reciprocity checks whether the reported coupling is property of the structure or artifact of a mismatched flux gauge \cite{stratton1941}. Failed reciprocity usually traces to a near-field tap, a radome standing-wave, or a probe polarization basis that was never aligned with the transmit standard \cite{hansen1988}.
Rayleigh and Stratton treated such ranges as boundary-value laboratories long before digital near-field transforms became routine \cite{stratton1941}. Harrington's method-of-moments culture reinforced the same discipline: metal outline first, harmonic headline second \cite{harrington2001}.
The measurement logic above is made precise in Subsection~\ref{sec:ch283} \cite{stratton1941}.
\end{physicalvignette}
\begin{physicalvignette}[Corner reflector for radar cross-section calibration]
Consider Corner reflector for radar cross-section calibration. Before any coefficient table is printed, technicians record frequency, cable loss, and the mechanical outline that fixes electrical size. The first walk-through of the range is always tactile: flange bolts, probe standoff, turntable boresight, and the wattmeter reference on the receive chain \cite{hansen1988}.
On this hardware the near zone is visually impressive---standing reactive structure, rapid phase rotation, sometimes a visible VSWR ripple at the feed---yet the export question is measured only after outgoing selection on a large sphere \cite{harrington2001}. What the camera sees in the chamber is not automatically what the far-zone pattern claims \cite{stratton1941}.
Engineers therefore learn to separate three readings: what the feed can synthesize on declared support, what propagating content survives on the radiative range, and what a finite instrument can recover at declared tolerance \cite{devaney2004}. Confusing those readings produces the rank gaps and watt failures that later sections quantify \cite{pal2026}.
Rayleigh and Stratton treated such ranges as boundary-value laboratories long before digital near-field transforms became routine \cite{stratton1941}. Harrington's method-of-moments culture reinforced the same discipline: metal outline first, harmonic headline second \cite{harrington2001}.
Subsection~\ref{sec:ch283} states the formal export condition that this range scene presupposes \cite{stratton1941}.
\end{physicalvignette}
\begin{equation}
\label{eq:ch2.port-radiation-reception-operator}
\mathcal{M}_P
=
\mathcal{O}_{\mathrm{meas}}^{(P)}\circ\Lambda_{\mathrm{rad}},
\qquad
\Lambda_{\mathrm{rad}}=\mathcal{O}_{\mathrm{far}}\circ\mathcal{R}_\omega,
\end{equation}
consistent with Eq.~\eqref{eq:ch2.reconstruction-observation-map}
\cite{stratton1941}. When \(\mathcal{O}_{\mathrm{meas}}^{(P)}\) is scalar,
\(\mathcal{M}_P[\Jimp]\in\mathbb{C}\) is the complex port response induced by impressed
\(\Jimp\); when vector-valued, \(\mathcal{M}_P[\Jimp]\in\mathbb{C}^{N_P}\) lists channel
readouts \cite{stratton1941}. The \emph{transmit map} at port \(B\) due to excitation at
port \(A\) is the composition \(\Jimp^{\mathrm{(tx)}}\mapsto
\mathcal{M}_B[\Jimp^{\mathrm{(tx)}}]\); the \emph{receive map} at port \(A\) due to excitation
at port \(B\) is \(\Jimp^{\mathrm{(rx)}}\mapsto\mathcal{M}_A[\Jimp^{\mathrm{(rx)}}]\)
\cite{harrington2001}. Algebraically, both are finite-rank linear functionals of
\(\mathcal{R}_\omega[\Jimp]\); physically, reciprocity constrains their interchange when
hardware and media are reciprocal \cite{stratton1941}.
Under Lorentz reciprocity for the solution map (Subsection~\ref{sec:ch281}) and self-adjoint port
functionals \(\mathcal{O}_{\mathrm{meas}}^{(P)}\) on \(\mathcal{F}_{\mathrm{rad}}\), the
complex port response at \(B\) induced by excitation at \(A\) equals the response at \(A\) when
the impressed source and receive probe are exchanged on reciprocal hardware:
\begin{equation}
\label{eq:ch2.transmit-receive-reciprocity}
\mathcal{M}_B\big[\Jimp^{\mathrm{(tx)}}\big]
=
\mathcal{M}_A\big[\Jimp^{\mathrm{(rx)}}\big]
\quad\text{when}\ 
\Jimp^{\mathrm{(tx)}}\ \text{and}\ \Jimp^{\mathrm{(rx)}}\ \text{are reciprocal port excitations},
\end{equation}
developing Eq.~\eqref{eq:ch2.transmit-receive-swap}
\cite{stratton1941}. Equation~\eqref{eq:ch2.transmit-receive-reciprocity}
is the operator reading of the VNA experiment: excite \(A\), read \(B\); swap roles on reciprocal
media and obtain the same complex port coupling on observable coordinates. It states that the
observable induced at \(B\) by the transmit source at \(A\) equals the observable at \(A\) when
the roles of impressed source and probe are exchanged on reciprocal media
\cite{stratton1941}. It is not a claim that \(\mathcal{M}_B=\mathcal{M}_A\) as
operators on \(\mathcal{J}_\Omega\): ports may have distinct apertures and channel counts
\cite{harrington2001}. It \emph{is} a claim that the \emph{bilinear} port coupling is symmetric
when qualified by Eq.~\eqref{eq:ch2.radiative-coupling-bilinear} on
\(\operatorname{ran}\mathcal{R}_\omega\) \cite{stratton1941}.
\begin{figure}[t]
\centering
\ChOneFigBlock{%
\begin{tikzpicture}[ch1fig, node distance=5mm and 7mm]
 \node[ch1box] (tx) {Port \(A\)\\$\Jimp^{\mathrm{(tx)}}$};
 \node[ch1box, right=7mm of tx] (r) {$\mathcal{R}_\omega$};
 \node[ch1box, right=7mm of r] (ob) {$\mathcal{O}_{\mathrm{meas}}^{(B)}$};
 \node[ch1box, right=7mm of ob] (rxout) {Rx at \(B\)};
 \node[ch1box, below=5mm of tx] (rx) {Port \(B\)\\$\Jimp^{\mathrm{(rx)}}$};
 \node[ch1box, below=5mm of ob] (ob2) {$\mathcal{O}_{\mathrm{meas}}^{(A)}$};
 \node[ch1box, below=5mm of rxout] (txout) {Rx at \(A\)};
 \draw[ch1flow] (tx) -- (r) -- (ob) -- (rxout);
 \draw[ch1flow] (rx) -- (r) -- (ob2) -- (txout);
 \draw[ch1flow, dashed] (rxout) to[bend left=20] node[right,font=\scriptsize] {swap} (txout);
\end{tikzpicture}%
}
\caption{Subsection~\ref{sec:ch283} transmit--receive measurement sequence: excitation at one port and
readout at the other interchange under Eq.~\eqref{eq:ch2.transmit-receive-reciprocity} on
reciprocal media.}
\label{fig:ch2.transmit-receive-measurement chain}
\end{figure}
Figure interpretation. Transmit port $\Jimp^{\mathrm{(tx)}}$ and receive probe $\Jimp^{\mathrm{(rx)}}$ compose through the same $\mathcal{R}_\omega$ and observation $\mathcal{O}_{\mathrm{meas}}$. Swap symmetry Eq.~\eqref{eq:ch2.transmit-receive-swap} holds on observable coordinates after projection onto $\operatorname{ran}\mathcal{R}_\omega$; additions in $\ker\mathcal{R}_\omega$ leave port coupling unchanged. The schematic separates reciprocity-qualified S-parameters from raw near-field cross terms that may violate apparent symmetry while far-zone reciprocity remains valid \cite{stratton1941}. Figure~\ref{fig:ch2.transmit-receive-measurement chain} displays the two compositions constrained by
Eq.~\eqref{eq:ch2.transmit-receive-reciprocity}. Both paths pass through the same
\(\mathcal{R}_\omega\); port asymmetry enters only through \(\mathcal{O}_{\mathrm{meas}}^{(A)}\),
\(\mathcal{O}_{\mathrm{meas}}^{(B)}\) \cite{stratton1941}.
\phantomsection\label{thm:ch2.radiation-reception-reciprocity}
Let media satisfy Eq.~\eqref{eq:ch2.reciprocal-medium-class} with passive
Eq.~\eqref{eq:ch2.passive-impedance-class} completion, and let port operators
\(\mathcal{M}_P=\mathcal{O}_{\mathrm{meas}}^{(P)}\circ\Lambda_{\mathrm{rad}}\) be defined by
Eq.~\eqref{eq:ch2.port-radiation-reception-operator} with self-adjoint
\(\mathcal{O}_{\mathrm{meas}}^{(P)}\) on \(\mathcal{F}_{\mathrm{rad}}\). For admissible port
sources \(\Jimp^{\mathrm{(tx)}}\), \(\Jimp^{\mathrm{(rx)}}\):
\begin{enumerate}
\item[(i)] Scalar reciprocity Eq.~\eqref{eq:ch2.transmit-receive-reciprocity} holds when port
excitations are paired through Lorentz reciprocity for the solution map (Subsection~\ref{sec:ch281}).
\item[(ii)] Bilinear port coupling
\begin{equation}
\label{eq:ch2.port-coupling-bilinear}
\mathfrak{C}_{AB}
=
\langle \mathcal{O}_{\mathrm{meas}}^{(B)}[\mathcal{R}_\omega[\Jimp^{\mathrm{(tx)}}]],\,
\mathcal{O}_{\mathrm{meas}}^{(A)}[\mathcal{R}_\omega[\Jimp^{\mathrm{(rx)}}]]\rangle
\end{equation}
is symmetric under interchange of port labels when inner products are those induced by
Eq.~\eqref{eq:ch2.field-energy-norm} on \(\mathcal{F}_{\mathrm{rad}}\).
\item[(iii)] Two-port impedance matrix \(\mathbf{Z}\) with entries
\(Z_{PQ}=\mathcal{M}_P[\Jimp^{(Q)}]\) on matched reciprocal ports satisfies
\begin{equation}
\label{eq:ch2.reciprocal-impedance-matrix}
Z_{PQ}=Z_{QP}.
\end{equation}
\item[(iv)] Adding \(\Jimp^{\mathrm{(NR)}}\in\ker\mathcal{R}_\omega\) to either port excitation
leaves Eqs.~\eqref{eq:ch2.transmit-receive-reciprocity}--\eqref{eq:ch2.reciprocal-impedance-matrix}
unchanged on observable coordinates.
\end{enumerate}
Part (i) applies Lorentz reciprocity for the solution map (Subsection~\ref{sec:ch281})(iv) to
\(\mathcal{R}_\omega[\Jimp^{\mathrm{(tx)}}]\), \(\mathcal{R}_\omega[\Jimp^{\mathrm{(rx)}}]\),
then self-adjointness of \(\mathcal{O}_{\mathrm{meas}}^{(P)}\)
\cite{stratton1941}. Part (ii) is
Eq.~\eqref{eq:ch2.coupling-bilinear-symmetry} with observation weights inserted; symmetry of
\(\mathfrak{C}_\omega\) from orthogonality and coupling of radiative modes (Subsection~\ref{sec:ch282})(iii)
\cite{stratton1941}. Part (iii) specializes (ii) to matched port basis vectors
\cite{harrington2001}. Part (iv) uses
\(\mathcal{R}_\omega[\Jimp^{\mathrm{(NR)}}]=\mathbf{0}\) and
null-space invariance of radiative reciprocity (Subsection~\ref{sec:ch281}) \cite{stratton1941}.
Mathematical role of Eq.~\eqref{eq:ch2.port-coupling-bilinear}. Equation~\eqref{eq:ch2.port-coupling-bilinear} packages mutual coupling as a bilinear form on
\(\operatorname{ran}\mathcal{R}_\omega\) after port projection \cite{stratton1941}. It reduces
to Eq.~\eqref{eq:ch2.radiative-coupling-bilinear} when
\(\mathcal{O}_{\mathrm{meas}}^{(A)}=\mathcal{O}_{\mathrm{meas}}^{(B)}=\mathrm{id}\) on the
observed subspace \cite{stratton1941}. Physically, \(\mathfrak{C}_{AB}\) is the
operator form of mutual impedance or scattering amplitude between declared ports when regime qualification
Eq.~\eqref{eq:ch2.regime-component-criterion} is satisfied \cite{harrington2001}. Finite
rank Eq.~\eqref{eq:ch2.reconstruction-rank-cap} limits the number of independent entries
\(\mathfrak{C}_{AB}\) can resolve \cite{stratton1941}.
On orthonormal modes Eq.~\eqref{eq:ch2.radiation-modal-orthogonality}, port responses expand as
\begin{equation}
\label{eq:ch2.port-modal-response-expansion}
\mathcal{M}_P[\Jimp]
=
\sum_{n\in\mathcal{N}_{\mathrm{rad}}}
w_{P,n}\,\lambda_n\,c_n[\Jimp],
\end{equation}
with port weights \(w_{P,n}=\langle \mathbf{f}_P,\mathbf{u}_n\rangle_{\mathcal{U}}\) inherited
from Eq.~\eqref{eq:ch2.finite-rank-observation} and amplitudes \(c_n\) from
Eq.~\eqref{eq:ch2.source-modal-amplitude} \cite{stratton1941}. Mutual coupling
between ports \(A\) and \(B\) becomes
\begin{equation}
\label{eq:ch2.port-modal-coupling-sum}
\mathfrak{C}_{AB}
=
\sum_{n\in\mathcal{N}_{\mathrm{rad}}}
|\lambda_n|^{2}\,w_{A,n}^{*}\,w_{B,n}\,
c_n[\Jimp^{\mathrm{(tx)}}]^{*}\,c_n[\Jimp^{\mathrm{(rx)}}],
\end{equation}
using Eq.~\eqref{eq:ch2.modal-coupling-sum-diagonal} \cite{stratton1941}. Equation
\eqref{eq:ch2.port-modal-coupling-sum} shows that reciprocity symmetry is diagonal in a mode
basis but may be dense in fixed laboratory coordinates when degenerate multiplets from
Subsection~\ref{sec:ch262} are expressed in non-eigenvector bases
\cite{harrington2001}.
\phantomsection\label{prop:ch2.measurement-reciprocity}
For full rank-$N$ instrument readout \(\mathcal{M}_\omega=\mathcal{O}_{\mathrm{meas}}\circ\Lambda_{\mathrm{rad}}\)
Eq.~\eqref{eq:ch2.reconstruction-observation-map}, reciprocity implies
\(\mathcal{M}_\omega[\Jimp^{\mathrm{(tx)}}]\) and \(\mathcal{M}_\omega[\Jimp^{\mathrm{(rx)}}]\)
obey the same swap law Eq.~\eqref{eq:ch2.transmit-receive-reciprocity} on
\(\operatorname{ran}\mathcal{O}_{\mathrm{meas}}\) \cite{stratton1941}. Sources
in \(\ker\mathcal{M}_\omega\) leave all port readouts unchanged
Eq.~\eqref{eq:ch2.reconstruction-null-space}; reciprocity does not resolve those equivalence
classes \cite{stratton1941}. Rank cap Eq.~\eqref{eq:ch2.reconstruction-rank-cap} bounds
independent reciprocal channels \cite{harrington2001}.
Apply reciprocity in radiation and reception (Subsection~\ref{sec:ch283})(i) with
\(\mathcal{O}_{\mathrm{meas}}^{(A)}=\mathcal{O}_{\mathrm{meas}}^{(B)}=\mathcal{O}_{\mathrm{meas}}\)
\cite{stratton1941}. Null-space invariance is
measurement null-space extension (Subsection~\ref{sec:ch271}) \cite{stratton1941}. \phantomsection\label{prop:ch2.accessible-sector-reciprocity}
Let accessible modal sector Eq.~\eqref{eq:ch2.accessible-modal-sector} restrict far-zone content
to \(\mathcal{K}_{\mathrm{acc}}\). Then Eqs.~\eqref{eq:ch2.transmit-receive-reciprocity}--\eqref{eq:ch2.reciprocal-impedance-matrix}
hold on \(\mathcal{K}_{\mathrm{acc}}\) after projection, while suppressed directions
Eq.~\eqref{eq:ch2.suppressed-angular-sheet-def} remain unobservable independently of reciprocity
\cite{stratton1941}. Angular budget
Eq.~\eqref{eq:ch2.angular-spectrum-composition-law} composes with port reciprocity without
eliminating null-space or geometric deficits \cite{stratton1941}.
Reciprocity acts on \(\operatorname{ran}\mathcal{R}_\omega\); geometric and observation
constraints from geometric constraints on radiative spectra (Subsection~\ref{sec:ch264}) and
Eq.~\eqref{eq:ch2.angular-spectrum-composition-law} restrict reported coordinates only
Laboratory errors and broken reciprocity. Magneto-electric media, active surfaces, and non-reciprocal port calibrations break
Eq.~\eqref{eq:ch2.transmit-receive-reciprocity} even when export maps differ
\(\mathcal{S}_\omega\) are accurate \cite{stratton1941}. Near-zone reactive
coupling measured without \(\Pi_{\mathrm{rad}}\) can appear non-reciprocal while For matched half-wave dipole ports at $k_0 a=1.2$ with $d=4a$, Eq.~\eqref{eq:ch2.reciprocal-impedance-matrix} requires $Z_{12}=Z_{21}$ when entries are extracted from $\mathcal{R}_\omega[\Jimp^{\mathrm{(tx)}}]$ on the far-zone branch. A typical reciprocity check gives $|Z_{12}-Z_{21}|/|Z_{11}|<0.02$ in reciprocal vacuum; swapping transmit and receive probes while holding $\mathcal{O}_{\mathrm{meas}}$ fixed reproduces the same mutual entry within numerical tolerance \cite{jackson1999}. If near-field probes at $k_0 r=1.5$ report $|Z_{12}/Z_{21}-1|\approx0.18$, the discrepancy is a regime failure, not a Lorentz violation: re-meter on $\mathcal{R}_\omega[\cdot]$ with $k_0 r>3$ before asserting broken reciprocity \cite{harrington2001}. Eq.~\eqref{eq:ch2.radiation-map-reciprocity} holds on \(\mathcal{F}_{\mathrm{rad}}\)
\cite{harrington2001}. Using different outgoing branches or terminations for
transmit and receive models violates the shared-operator assumption behind reciprocity in radiation and reception (Subsection~\ref{sec:ch283}) \cite{stratton1941}. Tikhonov inversion
Eq.~\eqref{eq:ch2.tikhonov-source-reconstruction} does not restore reciprocity where
\(\ker\mathcal{M}_\omega\) is nontrivial \cite{harrington2001}.
Subsection~\ref{sec:ch283} requires linear reciprocal media at fixed \(\omega\), passive
\(\mathcal{B}_\Gamma\), self-adjoint port functionals, and outgoing radiation condition
\cite{stratton1941}. It does not identify unique \(\Jimp\) from port data when
Eq.~\eqref{eq:ch2.reconstruction-null-space} is nontrivial \cite{stratton1941}. It does not
replace angular-spectrum operators in Section~\ref{sec:ch29} or harmonic bases in
\ref{ch:03} \cite{stratton1941}. Regime
qualification Eq.~\eqref{eq:ch2.regime-operator-screen} remains mandatory before transport
interpretation of port coupling \cite{stratton1941}. \medskip
\noindent
Section~\ref{sec:ch28} closes the reciprocity chain begun in its introduction:
Lorentz reciprocity in operator form (Subsection~\ref{sec:ch281}),
orthogonality and coupling of radiative modes (Subsection~\ref{sec:ch282}), and
reciprocity in radiation and reception (Subsection~\ref{sec:ch283})
\cite{stratton1941}. Spectral coordinates
from Section~\ref{sec:ch26} and null-space structure from Section~\ref{sec:ch27} now meet
exchange symmetry: bilinear coupling on \(\operatorname{ran}\mathcal{R}_\omega\) is constrained
by Eqs.~\eqref{eq:ch2.lorentz-reciprocity-operator}--\eqref{eq:ch2.reciprocal-impedance-matrix},
while \(\ker\mathcal{R}_\omega\), finite observation rank
Eq.~\eqref{eq:ch2.reconstruction-rank-cap}, and angular budget
Eq.~\eqref{eq:ch2.angular-spectrum-composition-law} explain deficits reciprocity alone does not
remove \cite{stratton1941}. Section~\ref{sec:ch29} develops plane-wave action and
retained/suppressed angular components on the same outgoing branch; Subsection~\ref{sec:ch294}
refines far-field angular accessibility; \ref{ch:03} supplies explicit harmonic
orthogonality used by reference through Subsection~3.4.3
\cite{harrington2001}. Transport and port-coupling claims must compose with
Eq.~\eqref{eq:ch2.section28-reciprocity-chain}, regime qualification
Eq.~\eqref{eq:ch2.regime-operator-screen}, and null-space chain
Eq.~\eqref{eq:ch2.section27-nullspace-chain}, not with raw \(\mathcal{S}_\omega[\Jimp]\) or
unqualified near-zone reactive content \cite{stratton1941}.
\section{Action of Radiation Operators on Angular Spectra}
\label{sec:ch29}
Reciprocity in Section~\ref{sec:ch28} couples ports and probes. Section~\ref{sec:ch29} asks how
radiation operators act on angular spectra---the language of pattern decomposition, near-field
scanning, and plane-wave synthesis \cite{jackson1999}. Every qualified
far-zone pattern is first a weighted sum of directional contributions: each term reports how much
export reaches a unit direction $\widehat{\mathbf{k}}$ on the sphere after outgoing selection
\cite{stratton1941}. That decomposition precedes any lateral Fourier transform;
plane-wave spectra are one coordinate system on the same directional budget, not a separate
physical layer \cite{jackson1999}. Maxwell propagation from compact support cannot populate
every direction with equal strength: outgoing selection removes incoming admixture; evanescent
sheets carry reactive content that decays with distance and does not survive the near-to-far map;
finite apertures integrate over unresolved directions. The accessible amplitude
$a(\widehat{\mathbf{k}})$ is therefore a compressed image of impressed structure, not a complete
coordinate on $\mathcal{J}_\Omega$ \cite{stratton1941}. On the bench, a spherical
near-field scanner records tangential $\E$ and $\H$ on a closed or planar surface at height $h$
above the device under test; the angular-spectrum chart synthesizes outgoing plane-wave trials, back-propagates to
the scan plane, and accepts residuals only after scan pitch, bandlimit, and probe correction are
stated \cite{hansen1988}. The sequence scan--synthesize--bandlimit--report is what the
theory encodes once radiation has been selected \cite{stratton1941}.
% PHYSICAL-MEMORY:sec:ch29
% PHYSICAL-MEMORY:sec:ch29-b2
\begin{physicalvignette}[Horn over square ground plane at ka near unity]
A second bench observation matters on Reflector with grating lobe from subarray spacing. When technicians swap only the receive probe while holding the transmit chain fixed, reciprocity checks whether the reported coupling is property of the structure or artifact of a mismatched flux gauge \cite{stratton1941}. Failed reciprocity usually traces to a near-field tap, a radome standing-wave, or a probe polarization basis that was never aligned with the transmit standard \cite{hansen1988}.
Rayleigh and Stratton treated such ranges as boundary-value laboratories long before digital near-field transforms became routine \cite{stratton1941}. Harrington's method-of-moments culture reinforced the same discipline: metal outline first, harmonic headline second \cite{harrington2001}.
The measurement logic above is made precise in Subsection~\ref{sec:ch29} \cite{stratton1941}.
\end{physicalvignette}
\begin{physicalvignette}[Horn over square ground plane at ka near unity]
Consider Horn over square ground plane at ka near unity. Before any coefficient table is printed, technicians record frequency, cable loss, and the mechanical outline that fixes electrical size. The first walk-through of the range is always tactile: flange bolts, probe standoff, turntable boresight, and the wattmeter reference on the receive chain \cite{hansen1988}.
On this hardware the near zone is visually impressive---standing reactive structure, rapid phase rotation, sometimes a visible VSWR ripple at the feed---yet the export question is measured only after outgoing selection on a large sphere \cite{harrington2001}. What the camera sees in the chamber is not automatically what the far-zone pattern claims \cite{stratton1941}.
Engineers therefore learn to separate three readings: what the feed can synthesize on declared support, what propagating content survives on the radiative range, and what a finite instrument can recover at declared tolerance \cite{devaney2004}. Confusing those readings produces the rank gaps and watt failures that later sections quantify \cite{pal2026}.
Rayleigh and Stratton treated such ranges as boundary-value laboratories long before digital near-field transforms became routine \cite{stratton1941}. Harrington's method-of-moments culture reinforced the same discipline: metal outline first, harmonic headline second \cite{harrington2001}.
Subsection~\ref{sec:ch29} states the formal export condition that this range scene presupposes \cite{stratton1941}.
\end{physicalvignette}
Section~\ref{sec:ch28} closed the reciprocity chain
Eq.~\eqref{eq:ch2.section28-reciprocity-chain}: Lorentz reciprocity (Subsection~\ref{sec:ch281}),
modal orthogonality and coupling (Subsection~\ref{sec:ch282}), and radiation--reception
interchange (Subsection~\ref{sec:ch283}) \cite{stratton1941}. That program fixed exchange
symmetry on \(\operatorname{ran}\mathcal{R}_\omega\) and on port coordinates after projection.
Section~\ref{sec:ch29} develops filter action in plane-wave coordinates: how
\(\mathcal{R}_\omega\), \(\Pi_{\mathrm{rad}}\), and observation maps retain or suppress
continuous angular content on \(S^{2}\), and how near-zone fields map to far-zone amplitudes
\(a(\widehat{\mathbf{k}})\) \cite{stratton1941}. Preview expansion
Eq.~\eqref{eq:ch2.angular-spectrum} from Subsection~\ref{sec:ch224} named the lateral
spectral template; here that leading template is developed to locked statements compatible with outgoing
criterion Eq.~\eqref{eq:ch2.outgoing-plane-wave-criterion} and far amplitude
Eq.~\eqref{eq:ch2.far-field-amplitude-integral} \cite{stratton1941}. Explicit basis change to
vector harmonics is developed in \ref{ch:03} Subsection~3.8.1; harmonic normalization belongs to
Subsection~3.5.2 with citation of Subsection~\ref{sec:ch282} \cite{harrington2001}.
Spectral decomposition in Subsection~\ref{sec:ch261} and continuous representation
Eq.~\eqref{eq:ch2.exterior-continuous-spectral-integral} fixed modal and plane-wave coordinates
on \(\mathcal{F}_{\mathrm{rad}}\); null-space consequences in Subsection~\ref{sec:ch274} and
angular budget Eq.~\eqref{eq:ch2.angular-spectrum-composition-law} classified which directions
\(a(\widehat{\mathbf{k}})\) may carry export \cite{stratton1941}. Reciprocity from
Section~\ref{sec:ch28} constrains bilinear coupling on observable coordinates but does not
replace angular filters: a reciprocal antenna may still have suppressed sheets
Eq.~\eqref{eq:ch2.suppressed-angular-sheet-def} and finite observation rank
Eq.~\eqref{eq:ch2.angular-spectrum-observation-rank} \cite{stratton1941}. Treating
\(\mathcal{S}_\omega[\Jimp]\) near sources as if it were already \(a(\widehat{\mathbf{k}})\)
repeats the regime error class of Section~\ref{sec:ch01.2.4} \cite{stratton1941}.
Section~\ref{sec:ch29} separates plane-wave angular spectra, retained and suppressed
components, near-to-far transformation, and far-field angular $\mathcal{E}_{\mathrm{acc}}^{(N)}$ without re-opening
Green construction from Section~\ref{sec:ch23} \cite{stratton1941}.
Let \(\omega>0\) be fixed, \(\Omega_s=\operatorname{supp}\Jimp\) be compact, and outgoing pairs
lie in \(\mathcal{F}_{\mathrm{SF}}\) Eq.~\eqref{eq:ch2.admissible-field-space}. Section
~\ref{sec:ch29} organizes analysis in four stages:
\begin{equation}
\label{eq:ch2.section29-angular-spectrum-chain}
\text{plane-wave spectrum}
\ \xrightarrow{\ \mathcal{R}_\omega,\Pi_{\mathrm{rad}}\ }
\text{retained / suppressed}
\ \xrightarrow{\ \text{near-to-far}\ }
a(\widehat{\mathbf{k}})
\ \xrightarrow{\ \text{probe rank}\ }
\mathcal{K}_{\mathrm{acc}},
\end{equation}
where the first arrow is Subsection~\ref{sec:ch291}, the second is Subsection~\ref{sec:ch292},
the third is Subsection~\ref{sec:ch293}, and the fourth is Subsection~\ref{sec:ch294}
\cite{stratton1941}. Equation~\eqref{eq:ch2.section29-angular-spectrum-chain}
expresses Eq.~\eqref{eq:ch2.section27-nullspace-chain} and
Eq.~\eqref{eq:ch2.section28-reciprocity-chain} at the angular level: null-space and reciprocity
constraints from Sections~\ref{sec:ch27}--\ref{sec:ch28} bound which \(a(\widehat{\mathbf{k}})\)
are physically meaningful; Section~\ref{sec:ch29} develops the operators that enforce those
bounds on plane-wave coordinates \cite{stratton1941}. Transport criteria from
Section~\ref{sec:ch01.2.3} and angular observables from Section~\ref{sec:ch01.2.1} enter
Subsection~\ref{sec:ch294} by reference, not by re-derivation \cite{stratton1941}.
Modal representation, not separate physics (Sections~2.6 and~2.9). Spectral and angular decompositions organize information about \(\mathcal{R}_\omega[\Jimp]\);
they do not define angular structure \cite{stratton1941}. Observable angular structure remains
primary: modes, plane-wave sheets, and harmonic labels are coordinates on export, not substitutes
for Maxwell flux or qualified transport \cite{maxwell1873}. A nonzero modal coefficient or
angular amplitude reports surviving source-generated content after propagation and outgoing
selection; renaming that content does not create new physical degrees of freedom
\cite{harrington2001}.
The primary objects are radiation map
Eq.~\eqref{eq:ch2.nullspace-operator-targets}, outgoing plane-wave criterion
Eq.~\eqref{eq:ch2.outgoing-plane-wave-criterion}, and continuous angular amplitude
\(a(\widehat{\mathbf{k}})\) from Eq.~\eqref{eq:ch2.exterior-continuous-spectral-integral}
\cite{stratton1941}. Secondary targets include propagating projectors
Eqs.~\eqref{eq:ch2.spectral-projectors}--\eqref{eq:ch2.rad-null-on-evanescent},
observation composition Eq.~\eqref{eq:ch2.reconstruction-observation-map}, accessible sector
Eq.~\eqref{eq:ch2.angular-accessibility-sector}, and composition law
Eq.~\eqref{eq:ch2.angular-spectrum-composition-law} \cite{stratton1941}. Section~\ref{sec:ch29}
develops filters on the same outgoing branch fixed by
Eqs.~\eqref{eq:ch2.outgoing-radiation-refinement}--\eqref{eq:ch2.outgoing-resolvent-unique};
incoming sheets remain excluded from \(\mathcal{F}_{\mathrm{rad}}\)
\cite{harrington2001}.
Plane-wave angular spectrum and radiation operators (leading template). Subsection~\ref{sec:ch291} supplies the definitive account of plane-wave angular spectra
and the action of \(\mathcal{R}_\omega\circ\Pi_{\mathrm{prop}}\) on spectral coefficients: \begin{equation}
\label{eq:ch2.plane-wave-radiation-operator}
\widetilde{\E}_{\mathrm{rad}}(k_x,k_y)
=
\mathcal{F}_{\mathrm{rad}}\big[\Pi_{\mathrm{rad}}\,\mathcal{R}_\omega[\Jimp]\big](k_x,k_y),
\qquad
k_z=k_z(k_t)\ \text{from Eq.~\eqref{eq:ch2.outgoing-plane-wave-criterion}},
\end{equation}
with \(\mathcal{F}_{\mathrm{rad}}\) the lateral Fourier transform on outgoing admissible fields
\cite{stratton1941}. Equation~\eqref{eq:ch2.plane-wave-radiation-operator}
develops Eq.~\eqref{eq:ch2.angular-spectrum} by inserting radiation selection;
evanescent sheets satisfy \(\widetilde{\E}_{\mathrm{rad}}=0\) on export through
Eq.~\eqref{eq:ch2.rad-null-on-evanescent} \cite{stratton1941}. Rigged-space normalization of
plane-wave modes belongs to \ref{ch:03} Subsection~3.1.3 \cite{harrington2001}.
Subsection~\ref{sec:ch292} partitions angular content into retained, suppressed, and
observation-filtered sheets:
\begin{equation}
\label{eq:ch2.retained-suppressed-filter}
a(\widehat{\mathbf{k}})
=
a_{\mathrm{ret}}(\widehat{\mathbf{k}})\,\chi_{\mathcal{K}_{\mathrm{acc}}}(\widehat{\mathbf{k}})
\;+\;
a_{\mathrm{sup}}(\widehat{\mathbf{k}})\,\chi_{\mathcal{K}_{\mathrm{sup}}}(\widehat{\mathbf{k}}),
\qquad
a_{\mathrm{filt}}=\mathcal{O}_{\mathrm{meas}}[a],
\end{equation}
using accessible sector Eq.~\eqref{eq:ch2.angular-accessibility-sector} and suppressed sheet
Eq.~\eqref{eq:ch2.suppressed-angular-sheet-def} from Subsection~\ref{sec:ch274}
\cite{stratton1941}. Equation~\eqref{eq:ch2.retained-suppressed-filter}
is the Section~\ref{sec:ch29} target for Subsection~\ref{sec:ch292} partition; geometric
constraints on radiative spectra (Subsection~\ref{sec:ch264}) and observation rank
Eq.~\eqref{eq:ch2.angular-spectrum-observation-rank} compose independently with
Eq.~\eqref{eq:ch2.angular-spectrum-composition-law} \cite{stratton1941}. Null modes annihilate
\(a(\widehat{\mathbf{k}})\) through Eq.~\eqref{eq:ch2.angular-spectrum-nullspace-law} without
altering near reactive content \cite{harrington2001}.
Near-to-far transformation (leading template). Subsection~\ref{sec:ch293} records the operator map from near-zone admissible fields to far
angular amplitudes:
\begin{equation}
\label{eq:ch2.near-to-far-transform}
a(\widehat{\mathbf{k}})
=
\lim_{r\to\infty}
\mathcal{H}_{\mathrm{ff}}\big[\mathcal{R}_\omega[\Jimp]\big](r,\widehat{\mathbf{k}}),
\end{equation}
with \(\mathcal{H}_{\mathrm{ff}}\) the far-zone extractor tied to
Eq.~\eqref{eq:ch2.far-field-amplitude-integral} and the far-field approximation of source integrals (Subsection~\ref{sec:ch252}) \cite{stratton1941}.
Equation~\eqref{eq:ch2.near-to-far-transform} does not promote evanescent storage to
export: \(\Pi_{\mathrm{ev}}\mathcal{S}_\omega[\Jimp]\) is annihilated before
Eq.~\eqref{eq:ch2.near-to-far-transform} when \(\mathcal{R}_\omega\) is unchanged
\cite{stratton1941}. regime qualification Eq.~\eqref{eq:ch2.regime-component-criterion} qualifies when the
limit defines transport data \cite{harrington2001}.
Subsection~\ref{sec:ch294} develops probe rank on \(S^{2}\) importing Section~\ref{sec:ch01.2.1}
angular observables and Section~\ref{sec:ch01.2.3} transport criteria by reference
\cite{stratton1941}:
\begin{equation}
\label{eq:ch2.far-angular-accessibility-preview}
\widehat{\mathbf{k}}\in\mathcal{K}_{\mathrm{acc}}
\ \Longleftrightarrow\
a(\widehat{\mathbf{k}})\ \text{observable after}\ \mathcal{M}_\omega
\ \text{and qualified by}\ \mathcal{D}_{J,\mathrm{reg}}=1,
\end{equation}
upgrading angular $\mathcal{E}_{\mathrm{acc}}^{(N)}$ on the far shell (Subsection~\ref{sec:ch264}) and composing with
Eq.~\eqref{eq:ch2.angular-spectrum-composition-law} \cite{stratton1941}. Equation
\eqref{eq:ch2.far-angular-accessibility-preview} is the Section~\ref{sec:ch29} completion target;
explicit plane-wave filters and measurement qualification is developed in Subsection~\ref{sec:ch294}
\cite{harrington2001}. Reciprocity in radiation and reception (Subsection~\ref{sec:ch283}) constrains
coupling on \(\mathcal{K}_{\mathrm{acc}}\) but does not enlarge it \cite{stratton1941}.
\begin{equation}
\label{eq:ch2.angular-action-composition}
\Jimp\in\ker\mathcal{R}_\omega
\ \Longrightarrow\
a(\widehat{\mathbf{k}})=0
\ \text{on export},
\qquad
\mathfrak{C}_\omega\ \text{symmetric on}\ \mathcal{K}_{\mathrm{acc}},
\end{equation}
importing Eq.~\eqref{eq:ch2.angular-spectrum-nullspace-law} and reciprocity in radiation and reception (Subsection~\ref{sec:ch283}) \cite{stratton1941}. Modal
coordinates Eq.~\eqref{eq:ch2.radiation-map-modal-sum} and plane-wave coordinates
Eq.~\eqref{eq:ch2.exterior-continuous-spectral-integral} describe the same
\(\operatorname{ran}\mathcal{R}_\omega\) on exterior domains; Section~\ref{sec:ch29} fixes
plane-wave-side filters, while \ref{ch:03} Subsection~3.8.1 supplies basis
change law \cite{stratton1941}. Finite-rank observation
Eq.~\eqref{eq:ch2.reconstruction-rank-cap} may remove angular content even when
\(a(\widehat{\mathbf{k}})\neq 0\) on \(\mathcal{F}_{\mathrm{rad}}\) \cite{stratton1941}.
Several constructions deliberately remain outside Section~\ref{sec:ch29}. Vector spherical
harmonic completeness and normalization are developed in \ref{ch:03}; Section~\ref{sec:ch29}
records plane-wave radiation action without tabulating \(Y_\ell^m\)
\cite{stratton1941}. Non-radiating catalogs remain in Section~\ref{sec:ch27};
reciprocity constraints remain in Section~\ref{sec:ch28} \cite{stratton1941}. Section~\ref{sec:ch29}
uses Fredholm structure, geometric constraints, and null-space structure from
Sections~\ref{sec:ch26}--\ref{sec:ch27} \cite{harrington2001}.
Regime logic from Section~\ref{sec:ch01.2.4} is unchanged: near-zone reactive dominance does not
create far-zone angular transport without \(\mathcal{D}_{J,\mathrm{reg}}=1\) and outgoing
qualification \cite{stratton1941}. Angular-transport claims must compose with
Eq.~\eqref{eq:ch2.section29-angular-spectrum-chain}, angular budget
Eq.~\eqref{eq:ch2.angular-spectrum-composition-law}, and regime qualification
Eq.~\eqref{eq:ch2.regime-operator-screen} \cite{harrington2001}.
\subsection{Plane-wave angular spectrum and radiation operators}
\label{sec:ch291}
% HARRINGTON-OPEN:sec:ch291
Plane-wave angular spectra connect near fields to far patterns on scanners \cite{jackson1999}. Evanescent content must be filtered before far-zone claims \cite{hansen1988}. Maxwell's retarded export to $\Gamma_R$ remains the governing constraint \cite{maxwell1873}.
Before lateral wavenumbers enter, fix what a pattern measurement actually resolves. Qualified
far-zone export is a directional budget: each direction $\widehat{\mathbf{k}}$ on $S^{2}$ carries
an amplitude $a(\widehat{\mathbf{k}})$ reporting surviving power flux after radiation selection
\cite{stratton1941}. Propagating directions ($k_t\le k_0$) contribute to that
budget and can reach the far zone; evanescent directions ($k_t>k_0$) store reactive energy that
decays exponentially with distance from the source and is annihilated on export
\cite{harrington2001}. A spherical near-field scanner or planar phased-array diagnostic
begins on a measurement surface, not in $(k_x,k_y)$ space \cite{hansen1988}. Probes record
tangential \(\E\) and \(\H\) on a sheet at height \(h\); the angular-spectrum chart synthesizes outgoing plane-wave
trials, back-propagates to the scan plane, and accepts the residual only after scan height, probe
pitch, and bandlimit are declared \cite{devaney2004}. Aperture campaigns favor plane spectra;
compact ranges favor spherical amplitudes, but both express the same directional decomposition in
different coordinates \cite{harrington2001}. Denser scan grids sharpen interpolation; they do not
invent propagating modes absent from \(\mathcal{R}_\omega[\Jimp]\) \cite{devaney2004}.
\medskip
\noindent\textit{Observation versus source coordinates.} Probes report fields in laboratory
Cartesian coordinates \((x,y,z)\) on the scan plane \(z=z_0\). Lateral wavenumbers \((k_x,k_y)\)
index plane-wave trials in the \((x,y)\) plane and are \emph{observation-side} spectral labels
attached to tangential traces on \(z=z_0\), not source coordinates on compact
\(\Omega_s=\operatorname{supp}\Jimp\) \cite{stratton1941}. Transverse magnitude
\(k_t=\sqrt{k_x^2+k_y^2}\) partitions the spectrum: sheets with \(k_t\le k_0\) carry propagating
export, while \(k_t>k_0\) carry evanescent reactive content whose decay length scales with scan
height \cite{harrington2001}. Impressed current \(\Jimp(\mathbf{r}')\) enters the reported spectrum
only through \(\mathcal{R}_\omega\) after radiation selection; treating raw scan data as if it
were already \(a(\widehat{\mathbf{k}})\) amplifies high-\(k_t\) noise into false superdirective
claims \cite{devaney2004}.
\medskip
Section~\ref{sec:ch29} developed angular-spectrum action through
Eqs.~\eqref{eq:ch2.section29-angular-spectrum-chain}--\eqref{eq:ch2.plane-wave-radiation-operator}.
Subsection~\ref{sec:ch291} supplies the definitive account of plane-wave angular spectra
and the action of \(\mathcal{R}_\omega\), \(\Pi_{\mathrm{rad}}\), and outgoing selection on
lateral spectral coefficients \cite{stratton1941}. Dispersion split
Eqs.~\eqref{eq:ch2.dispersion-relation}--\eqref{eq:ch2.longitudinal-wavenumber-split},
spectral projectors Eqs.~\eqref{eq:ch2.spectral-projectors}--\eqref{eq:ch2.rad-null-on-evanescent},
outgoing selector Eq.~\eqref{eq:ch2.outgoing-selector-def}, plane-wave criterion
Eq.~\eqref{eq:ch2.outgoing-plane-wave-criterion}, radiation map
Eq.~\eqref{eq:ch2.nullspace-operator-targets}, and lateral template
Eq.~\eqref{eq:ch2.angular-spectrum} are presupposed; Subsection~\ref{sec:ch224} and
Section~\ref{sec:ch243} proofs are not repeated \cite{stratton1941}. Continuous angular
representation Eq.~\eqref{eq:ch2.exterior-continuous-spectral-integral} and far amplitude
Eq.~\eqref{eq:ch2.far-field-amplitude-integral} are cited as consistency targets; rigged
normalization belongs to Chapter~3 (\emph{Angular Momentum Operators and SO(3)}) Subsection~3.1.3
\cite{harrington2001}. Retained/suppressed partitions belong to Subsection~\ref{sec:ch292}.
Assumptions. Fix \(\omega>0\), homogeneous exterior region with symbol \(k_0\) from
Eq.~\eqref{eq:ch2.helmholtz-symbol}, compact \(\Omega_s=\operatorname{supp}\Jimp\), admissible
\(\Jimp\in\mathcal{J}_\Omega\) Eq.~\eqref{eq:ch2.source-admissibility}, and outgoing fields
\(\mathbf{u}_\omega=(\E,\H)\in\mathcal{F}_{\mathrm{SF}}\) Eq.~\eqref{eq:ch2.admissible-field-space}
\cite{stratton1941}. Lateral wavenumbers \((k_x,k_y)\) use
\(k_t=\sqrt{k_x^2+k_y^2}\) and outgoing longitudinal root
Eq.~\eqref{eq:ch2.outgoing-plane-wave-criterion} \cite{stratton1941}. Transverse polarization satisfies
Eq.~\eqref{eq:ch2.transverse-polarization-constraint} on each spectral sheet
\cite{stratton1941}. regime qualification Eq.~\eqref{eq:ch2.regime-operator-screen} governs
transport interpretation of spectral outputs \cite{harrington2001}.
On a plane \(z=z_0\) in an outgoing admissible region, the SNF back-propagation step seeks a
superposition of plane-wave trials whose tangential field matches the measured trace on
\(z=z_0\). Each trial is labeled by \((k_x,k_y)\); the synthesis charting propagates it along
\(z\) with outgoing longitudinal wavenumber \(k_z(k_t)\) from
Eq.~\eqref{eq:ch2.outgoing-plane-wave-criterion} \cite{stratton1941}. The formal decomposition
packages that matching problem as a lateral Fourier transform on the observation plane.
Define the \emph{lateral angular spectrum} of \(\E\) by
% PHYSICAL-MEMORY:sec:ch291
% PHYSICAL-MEMORY:sec:ch291-b2
\begin{physicalvignette}[Handset held in a SAM phantom grip for OTA test]
A second bench observation matters on Pyramidal horn on a compact range. When technicians swap only the receive probe while holding the transmit chain fixed, reciprocity checks whether the reported coupling is property of the structure or artifact of a mismatched flux gauge \cite{stratton1941}. Failed reciprocity usually traces to a near-field tap, a radome standing-wave, or a probe polarization basis that was never aligned with the transmit standard \cite{hansen1988}.
Rayleigh and Stratton treated such ranges as boundary-value laboratories long before digital near-field transforms became routine \cite{stratton1941}. Harrington's method-of-moments culture reinforced the same discipline: metal outline first, harmonic headline second \cite{harrington2001}.
The measurement logic above is made precise in Subsection~\ref{sec:ch291} \cite{stratton1941}.
\end{physicalvignette}
\begin{physicalvignette}[Handset held in a SAM phantom grip for OTA test]
Consider Handset held in a SAM phantom grip for OTA test. Before any coefficient table is printed, technicians record frequency, cable loss, and the mechanical outline that fixes electrical size. The first walk-through of the range is always tactile: flange bolts, probe standoff, turntable boresight, and the wattmeter reference on the receive chain \cite{hansen1988}.
On this hardware the near zone is visually impressive---standing reactive structure, rapid phase rotation, sometimes a visible VSWR ripple at the feed---yet the export question is measured only after outgoing selection on a large sphere \cite{harrington2001}. What the camera sees in the chamber is not automatically what the far-zone pattern claims \cite{stratton1941}.
Engineers therefore learn to separate three readings: what the feed can synthesize on declared support, what propagating content survives on the radiative range, and what a finite instrument can recover at declared tolerance \cite{devaney2004}. Confusing those readings produces the rank gaps and watt failures that later sections quantify \cite{pal2026}.
Plane-wave angular spectra connect near fields to far patterns on scanners . Evanescent content must be filtered before far-zone claims .
Rayleigh and Stratton treated such ranges as boundary-value laboratories long before digital near-field transforms became routine \cite{stratton1941}. Harrington's method-of-moments culture reinforced the same discipline: metal outline first, harmonic headline second \cite{harrington2001}.
Subsection~\ref{sec:ch291} states the formal export condition that this range scene presupposes \cite{stratton1941}.
\end{physicalvignette}
\begin{equation}
\label{eq:ch2.plane-wave-angular-spectrum-def}
\widetilde{\E}(k_x,k_y;z_0)
=
\int_{-\infty}^{\infty}\!\!\int_{-\infty}^{\infty}
\E(x,y,z_0)\,
e^{\jj(k_x x+k_y y)}\,dx\,dy,
\end{equation}
with inverse
\begin{equation}
\label{eq:ch2.plane-wave-angular-spectrum-inverse}
\E(x,y,z_0)
=
\frac{1}{(2\pi)^{2}}
\int_{-\infty}^{\infty}\!\!\int_{-\infty}^{\infty}
\widetilde{\E}(k_x,k_y;z_0)\,
e^{-\jj(k_x x+k_y y)}\,dk_x\,dk_y,
\end{equation}
under standard Fourier conventions \cite{stratton1941}. Helmholtz propagation along
\(z\) multiplies each spectral component by
\begin{equation}
\label{eq:ch2.plane-wave-z-propagation}
\widetilde{\E}(k_x,k_y;z)
=
\widetilde{\E}(k_x,k_y;z_0)\,
e^{-\jj k_z(k_t)(z-z_0)},
\qquad
k_z\ \text{from Eq.~\eqref{eq:ch2.outgoing-plane-wave-criterion}},
\end{equation}
for homogeneous media \cite{stratton1941}. Equations~\eqref{eq:ch2.plane-wave-angular-spectrum-def}--\eqref{eq:ch2.plane-wave-z-propagation}
upgrade Eq.~\eqref{eq:ch2.angular-spectrum} to explicit forward/inverse
transforms with outgoing \(k_z\) selection \cite{stratton1941}. Mathematically,
\(\widetilde{\E}(k_x,k_y;z_0)\) is the lateral Fourier transform of the transverse field sampled
on the scan plane; each \((k_x,k_y)\) sheet is an independent plane-wave channel whose
longitudinal behavior is fixed by Eq.~\eqref{eq:ch2.longitudinal-wavenumber-split}
\cite{stratton1941}. Physically, propagating sheets (\(k_t\le k_0\)) carry energy that can reach
the far zone after Eq.~\eqref{eq:ch2.near-to-far-transform}, whereas evanescent sheets
(\(k_t>k_0\)) decay exponentially with distance from \(z_0\) and remain bound to the reactive
shell unless a separate source excites them \cite{harrington2001}. The inverse transform
Eq.~\eqref{eq:ch2.plane-wave-angular-spectrum-inverse} reconstructs the tangential field on
\(z=z_0\) from its spectral coefficients; probe interpolation error and finite aperture enter
here, before any claim about far-zone pattern \cite{hansen1988}. Propagation
Eq.~\eqref{eq:ch2.plane-wave-z-propagation} shifts each sheet along \(z\) by the phase factor
\(e^{-\jj k_z(k_t)(z-z_0)}\), which is the operator-side statement of SNF back-propagation
between scan planes at different heights \cite{jackson1999}. Consequence for measurement:
reported angular-spectrum charts are only as trustworthy as the declared scan height and
bandlimit on \((k_x,k_y)\) \cite{devaney2004}.
Let \(\mathcal{F}_{\mathrm{lat}}\) denote the lateral Fourier operator in
Eq.~\eqref{eq:ch2.plane-wave-angular-spectrum-def} on admissible traces. Define the
\emph{radiative lateral spectrum} of \(\mathcal{R}_\omega[\Jimp]\) by
\begin{equation}
\label{eq:ch2.radiative-lateral-spectrum-def}
\widetilde{\E}_{\mathrm{rad}}(k_x,k_y)
=
\mathcal{F}_{\mathrm{lat}}\big[\Pi_{\mathrm{rad}}\,\mathcal{R}_\omega[\Jimp]\big](k_x,k_y),
\end{equation}
with \(\Pi_{\mathrm{rad}}=\mathcal{O}_{\mathrm{out}}\circ\Pi_{\mathrm{prop}}\) from
Eq.~\eqref{eq:ch2.rad-projection-composition} \cite{stratton1941}. Equation
\eqref{eq:ch2.radiative-lateral-spectrum-def} is the Subsection~\ref{sec:ch291} upgrade of
Eq.~\eqref{eq:ch2.plane-wave-radiation-operator}: radiation selection precedes
lateral transformation \cite{stratton1941}. Applying \(\mathcal{F}_{\mathrm{lat}}\) to raw
\(\mathcal{S}_\omega[\Jimp]\) without \(\Pi_{\mathrm{rad}}\) mixes reactive evanescent sheets
into reported spectra through Eq.~\eqref{eq:ch2.subspace-energy-inequality}
\cite{harrington2001}.
The composition \(\mathcal{T}_\omega\) below is the spectral-domain radiation map: impressed
currents are transformed laterally, passed through \(\mathcal{R}_\omega\) and outgoing
selection, and read back in \((k_x,k_y)\) coordinates. Define the \emph{angular-spectrum
radiation operator} on lateral coefficients by
\begin{equation}
\label{eq:ch2.plane-wave-radiation-operator}
\widetilde{\E}_{\mathrm{rad}}(k_x,k_y)
=
\big[\mathcal{T}_\omega\widetilde{\Jimp}\big](k_x,k_y),
\qquad
\mathcal{T}_\omega
=
\mathcal{F}_{\mathrm{lat}}\circ\Pi_{\mathrm{rad}}\circ\mathcal{R}_\omega\circ\mathcal{F}_{\mathrm{lat}}^{-1},
\end{equation}
where \(\widetilde{\Jimp}=\mathcal{F}_{\mathrm{lat}}[\Jimp]\) on the source support and
\(\mathcal{F}_{\mathrm{lat}}^{-1}\) is understood on the admissible impressed-current class
\cite{stratton1941}. Equation~\eqref{eq:ch2.plane-wave-radiation-operator}
packages Eq.~\eqref{eq:ch2.radiative-lateral-spectrum-def} as a spectral-domain map: sources
excite lateral modes, \(\mathcal{R}_\omega\) exports selected content, and \(\mathcal{F}_{\mathrm{lat}}\)
returns observable spectral coefficients \cite{stratton1941}. In laboratory terms,
\(\mathcal{T}_\omega\) is the operator that a qualified SNF transform applies \emph{after}
radiation selection: impressed currents are first mapped to outgoing fields, then decomposed into
\((k_x,k_y)\) sheets on the propagating branch \cite{hansen1988}. On the propagating outgoing sheet
\(k_t\le k_0\), \(\mathcal{T}_\omega\) reduces to diagonal action in \((k_x,k_y)\) modulo
source coupling, and those sheets are the only ones that survive comparison with far-zone
amplitude \(a(\widehat{\mathbf{k}})\) \cite{stratton1941}. On the evanescent sheet \(k_t>k_0\),
\begin{equation}
\label{eq:ch2.evanescent-angular-null-law}
\widetilde{\E}_{\mathrm{rad}}(k_x,k_y)=\mathbf{0}
\quad\text{for}\ k_t>k_0,
\end{equation}
by Eq.~\eqref{eq:ch2.rad-null-on-evanescent} \cite{stratton1941}. Consequence: raw lateral
transforms of \(\mathcal{S}_\omega[\Jimp]\) may show large \(|\widetilde{\E}(k_x,k_y)|\) for
\(k_t>k_0\) while \(\widetilde{\E}_{\mathrm{rad}}=\mathbf{0}\) on export; only the radiative
branch Eq.~\eqref{eq:ch2.radiative-lateral-spectrum-def} may be compared to far-field reports
\cite{harrington2001}.
\begin{figure}[t]
\centering
\ChOneFigBlock{%
\begin{tikzpicture}[ch1fig, node distance=5mm and 7mm]
 \node[ch1box] (j) {$\Jimp$};
 \node[ch1box, right=7mm of j] (r) {$\mathcal{R}_\omega$};
 \node[ch1box, right=7mm of r] (pi) {$\Pi_{\mathrm{rad}}$};
 \node[ch1box, right=7mm of pi] (f) {$\mathcal{F}_{\mathrm{lat}}$};
 \node[ch1box, right=7mm of f] (spec) {$\widetilde{\E}_{\mathrm{rad}}$};
 \node[ch1box, below=5mm of pi] (oo) {$\mathcal{O}_{\mathrm{out}}$};
 \node[ch1box, below=5mm of r] (pp) {$\Pi_{\mathrm{prop}}$};
 \draw[ch1flow] (j) -- (r) -- (pi) -- (f) -- (spec);
 \draw[ch1flow] (pp.east) -- (pp.east |- pi.west);
 \draw[ch1flow] (pp.east |- pi.west) -- (pi.west);
 \draw[ch1flow] (oo.south) -- (oo.south |- pi.north);
 \draw[ch1flow] (oo.south |- pi.north) -- (pi.north);
\end{tikzpicture}%
}
\caption{Subsection~\ref{sec:ch291} plane-wave measurement sequence:
Eq.~\eqref{eq:ch2.plane-wave-radiation-operator} applies \(\mathcal{R}_\omega\),
Eq.~\eqref{eq:ch2.rad-projection-composition}, then lateral transform
Eq.~\eqref{eq:ch2.plane-wave-angular-spectrum-def}.}
\label{fig:ch2.plane-wave-radiation-measurement chain}
\end{figure}
The horizontal chain enforces operator order: $\mathcal{R}_\omega$ first, then $\Pi_{\mathrm{rad}}=\mathcal{O}_{\mathrm{out}}\circ\Pi_{\mathrm{prop}}$, then lateral transform $\mathcal{F}_{\mathrm{lat}}$. Lower nodes $\Pi_{\mathrm{prop}}$ and $\mathcal{O}_{\mathrm{out}}$ are \emph{constituents} of radiation selection, not optional post-filters. Applying $\mathcal{F}_{\mathrm{lat}}$ to raw $\mathcal{S}_\omega[\Jimp]$ without $\Pi_{\mathrm{rad}}$ contaminates $\widetilde{\E}_{\mathrm{rad}}$ with evanescent sheets annihilated by Eq.~\eqref{eq:ch2.evanescent-angular-null-law}. The terminal node $\widetilde{\E}_{\mathrm{rad}}$ is the plane-wave-side readout of qualified export only \cite{stratton1941}. Figure~\ref{fig:ch2.plane-wave-radiation-measurement chain} orders operators as in
Figure~\ref{fig:ch2.outgoing-selector-measurement chain}: propagation selection, outgoing criterion, then
lateral spectral readout \cite{stratton1941}.
Mathematical role of Eq.~\eqref{eq:ch2.plane-wave-radiation-operator}. Equation~\eqref{eq:ch2.plane-wave-radiation-operator} packages the radiation chain as a
spectral-domain operator \(\mathcal{T}_\omega\) acting on lateral source coefficients
\cite{stratton1941}. It is not a separate post-processor on near-zone
\(\mathcal{S}_\omega[\Jimp]\): \(\Pi_{\mathrm{rad}}\) is part of the definition
\cite{stratton1941}. Only outgoing propagating sheets contribute to far-zone angular
content; incoming sheets are removed by Eq.~\eqref{eq:ch2.outgoing-selector-def}, and
evanescent sheets are annihilated by Eq.~\eqref{eq:ch2.evanescent-angular-null-law}
\cite{harrington2001}. Reciprocity in radiation and reception (Subsection~\ref{sec:ch281})
constrains bilinear couplings built from \(\widetilde{\E}_{\mathrm{rad}}\) but does not alter
Eq.~\eqref{eq:ch2.evanescent-angular-null-law} \cite{stratton1941}.
Mathematical role of Eq.~\eqref{eq:ch2.evanescent-angular-null-law}. Equation~\eqref{eq:ch2.evanescent-angular-null-law} is the plane-wave restatement of
Eq.~\eqref{eq:ch2.rad-null-on-evanescent} on lateral coordinates \cite{stratton1941}. Large
\(|\widetilde{\E}(k_x,k_y)|\) for \(k_t>k_0\) may appear in raw
\(\mathcal{F}_{\mathrm{lat}}[\mathcal{S}_\omega[\Jimp]]\) while
\(\widetilde{\E}_{\mathrm{rad}}=\mathbf{0}\) on export \cite{harrington2001}. Regime
qualification Eq.~\eqref{eq:ch2.regime-component-criterion} remains mandatory before interpreting any
nonzero raw lateral spectrum as transport \cite{stratton1941}.
\phantomsection\label{prop:ch2.plane-wave-null-invariance}
If \(\Jimp^{(1)}\) is arbitrary and \(\Jimp^{(2)}=\Jimp^{(1)}+\Jimp^{\mathrm{(NR)}}\) with
\(\Jimp^{\mathrm{(NR)}}\in\ker\mathcal{R}_\omega\), then
\(\widetilde{\E}_{\mathrm{rad}}\) from Eq.~\eqref{eq:ch2.radiative-lateral-spectrum-def} is
unchanged under replacement of \(\Jimp^{(1)}\) by \(\Jimp^{(2)}\)
\cite{stratton1941}. Raw \(\mathcal{F}_{\mathrm{lat}}[\mathcal{S}_\omega[\Jimp]]\)
may change through reactive content Eq.~\eqref{eq:ch2.nr-image-of-null-space}
\cite{stratton1941}. \(\mathcal{R}_\omega[\Jimp^{(2)}]=\mathcal{R}_\omega[\Jimp^{(1)}]\) by
Eq.~\eqref{eq:ch2.nr-superposition-law} \cite{stratton1941}.
\phantomsection\label{prop:ch2.plane-wave-continuous-consistency}
On exterior domains with outgoing radiation condition, amplitudes \(a_\sigma(\widehat{\mathbf{k}})\) in
Eq.~\eqref{eq:ch2.exterior-continuous-spectral-integral} and lateral spectrum
Eq.~\eqref{eq:ch2.radiative-lateral-spectrum-def} describe the same
\(\operatorname{ran}\mathcal{R}_\omega\) content under change of variables
\((k_x,k_y)\mapsto\widehat{\mathbf{k}}\) on the propagating sheet \(k_t\le k_0\)
\cite{stratton1941}. Normalization constants are fixed in
Subsection~3.1.3 \cite{stratton1941}. Both representations diagonalize outgoing propagation on the propagating sheet; agreement on
far-zone reduction follows plane-wave radiation-operator action (Subsection~\ref{sec:ch291})(iv) and
Eq.~\eqref{eq:ch2.exterior-continuous-spectral-integral} \cite{stratton1941}.
Modal expansion Eq.~\eqref{eq:ch2.radiation-map-modal-sum} and plane-wave expansion
Eq.~\eqref{eq:ch2.plane-wave-radiation-operator} describe the same radiative image on
\(\mathcal{F}_{\mathrm{rad}}\) under unitary change of basis on the propagating sheet
\cite{stratton1941}. Subsection~\ref{sec:ch291} fixes plane-wave-side operators;
\ref{ch:03} Subsection~3.8.1 supplies explicit harmonic transforms without
re-tabulating Eq.~\eqref{eq:ch2.plane-wave-radiation-operator} \cite{stratton1941}. Angular
budget Eq.~\eqref{eq:ch2.angular-spectrum-composition-law} constrains supports of
\(\widetilde{\E}_{\mathrm{rad}}\) independently of basis choice \cite{harrington2001}. Plane-wave radiation-operator action (Subsection~\ref{sec:ch291}) assumes linear homogeneous exterior
regions at fixed \(\omega\), outgoing radiation condition, and admissible \(\mathcal{J}_\Omega\)
\cite{stratton1941}. Piecewise homogeneous media require patchwise
Eq.~\eqref{eq:ch2.plane-wave-z-propagation} with interface matching imported from
Section~\ref{sec:ch234} \cite{stratton1941}. Cavity standing-wave spectra are excluded: lateral
plane-wave expansions on bounded domains do not replace discrete cavity books from
Eq.~\eqref{eq:ch2.cavity-discrete-radiation-spectrum} \cite{jackson1999}. Loss
modifies \(k_0\) and \(k_z\) with complex symbols; evanescent null law
Eq.~\eqref{eq:ch2.evanescent-angular-null-law} remains on export \cite{stratton1941}.
Subsection~\ref{sec:ch291} fixed plane-wave-side radiation action; retained, suppressed, and
filtered partitions in Subsection~\ref{sec:ch292} separate geometric accessibility from finite
aperture before the near-to-far limit of Subsection~\ref{sec:ch293} extracts $a(\widehat{\mathbf{k}})$.
\subsection{Retained, suppressed, and filtered angular components}
\phantomsection\label{sec:ch292}
Section~\ref{sec:ch29} developed angular partitioning through
Eq.~\eqref{eq:ch2.retained-suppressed-filter} and linked it to accessible sector
Eq.~\eqref{eq:ch2.angular-accessibility-sector} and suppressed sheet
Eq.~\eqref{eq:ch2.suppressed-angular-sheet-def} from Subsection~\ref{sec:ch274}
\cite{stratton1941}. Subsection~\ref{sec:ch292} gives the definitive account of retained,
suppressed, and observation-filtered angular components on the outgoing branch established in
Subsection~\ref{sec:ch291}. Near-field scanners record tangential fields that mix propagating
export with evanescent storage; the partition below separates what geometry permits to radiate,
what symmetry forbids, and what finite aperture resolves---three filters that must not be merged
in harmonic charts or laboratory interpretation \cite{hansen1988}.
\medskip
Plane-wave lateral spectra
Eqs.~\eqref{eq:ch2.radiative-lateral-spectrum-def}--\eqref{eq:ch2.plane-wave-radiation-operator},
evanescent null law Eq.~\eqref{eq:ch2.evanescent-angular-null-law}, angular budget
Eq.~\eqref{eq:ch2.angular-spectrum-composition-law}, and observation rank
Eq.~\eqref{eq:ch2.angular-spectrum-observation-rank} are presupposed; Subsection~\ref{sec:ch274}
proofs are not repeated \cite{stratton1941}. Continuous amplitude
\(a(\widehat{\mathbf{k}})\) from Eq.~\eqref{eq:ch2.exterior-continuous-spectral-integral} and
rank-$N$ instrument readout Eq.~\eqref{eq:ch2.reconstruction-observation-map} supply the angular and
laboratory coordinates in which the partition is stated \cite{stratton1941}.
Near-to-far extraction belongs to Subsection~\ref{sec:ch293}; explicit $\mathcal{E}_{\mathrm{acc}}^{(N)}$ filters on
\(S^{2}\) belong to Subsection~\ref{sec:ch294} \cite{stratton1941}. Assumptions. Fix \(\omega>0\), admissible \(\Jimp\in\mathcal{J}_\Omega\), outgoing radiative field
\(\mathbf{u}_\omega=\mathcal{R}_\omega[\Jimp]\in\mathcal{F}_{\mathrm{rad}}\), and angular
amplitude \(a(\widehat{\mathbf{k}})\) from
Eq.~\eqref{eq:ch2.exterior-continuous-spectral-integral} on the exporting transverse sheet
\cite{stratton1941}. Accessible sector
\(\mathcal{K}_{\mathrm{acc}}\) Eq.~\eqref{eq:ch2.angular-accessibility-sector} and suppressed
sheet \(\mathcal{K}_{\mathrm{sup}}=S^{2}\setminus\mathcal{K}_{\mathrm{acc}}\) from
Eq.~\eqref{eq:ch2.suppressed-angular-sheet-def} are fixed by geometry and symmetry of
\((\Omega_s,\Gamma,\varepsilon,\mu)\) \cite{stratton1941}. Finite-rank observation
\(\mathcal{M}_\omega=\mathcal{O}_{\mathrm{meas}}\circ\Lambda_{\mathrm{rad}}\) from
Eq.~\eqref{eq:ch2.reconstruction-observation-map} is presupposed; regime qualification
Eq.~\eqref{eq:ch2.regime-operator-screen} qualifies transport interpretation of any reported
component \cite{harrington2001}.
The accessible sheet \(\mathcal{K}_{\mathrm{acc}}\subset S^{2}\) collects directions that
geometry and symmetry permit to carry qualified export; its complement
\(\mathcal{K}_{\mathrm{sup}}\) is suppressed before any instrument samples the pattern.
Define the \emph{retained angular amplitude} as the exportable content supported on the
accessible sheet. On a directional antenna bench, this is the portion of the far-zone pattern
that symmetry and blockage permit to carry qualified export; directions on
\(\mathcal{K}_{\mathrm{sup}}\) are zero by geometry before any instrument samples the pattern
\cite{stratton1941}:
% PHYSICAL-MEMORY:sec:ch292
% PHYSICAL-MEMORY:sec:ch292-b2
\begin{physicalvignette}[Parabolic dish on a three-axis turntable]
A second bench observation matters on Spherical near-field scanner with dual-polarization probes. When technicians swap only the receive probe while holding the transmit chain fixed, reciprocity checks whether the reported coupling is property of the structure or artifact of a mismatched flux gauge \cite{stratton1941}. Failed reciprocity usually traces to a near-field tap, a radome standing-wave, or a probe polarization basis that was never aligned with the transmit standard \cite{hansen1988}.
Rayleigh and Stratton treated such ranges as boundary-value laboratories long before digital near-field transforms became routine \cite{stratton1941}. Harrington's method-of-moments culture reinforced the same discipline: metal outline first, harmonic headline second \cite{harrington2001}.
The measurement logic above is made precise in Subsection~\ref{sec:ch292} \cite{stratton1941}.
\end{physicalvignette}
\begin{physicalvignette}[Parabolic dish on a three-axis turntable]
Consider Parabolic dish on a three-axis turntable. Before any coefficient table is printed, technicians record frequency, cable loss, and the mechanical outline that fixes electrical size. The first walk-through of the range is always tactile: flange bolts, probe standoff, turntable boresight, and the wattmeter reference on the receive chain \cite{hansen1988}.
On this hardware the near zone is visually impressive---standing reactive structure, rapid phase rotation, sometimes a visible VSWR ripple at the feed---yet the export question is measured only after outgoing selection on a large sphere \cite{harrington2001}. What the camera sees in the chamber is not automatically what the far-zone pattern claims \cite{stratton1941}.
Engineers therefore learn to separate three readings: what the feed can synthesize on declared support, what propagating content survives on the radiative range, and what a finite instrument can recover at declared tolerance \cite{devaney2004}. Confusing those readings produces the rank gaps and watt failures that later sections quantify \cite{pal2026}.
Rayleigh and Stratton treated such ranges as boundary-value laboratories long before digital near-field transforms became routine \cite{stratton1941}. Harrington's method-of-moments culture reinforced the same discipline: metal outline first, harmonic headline second \cite{harrington2001}.
Subsection~\ref{sec:ch292} states the formal export condition that this range scene presupposes \cite{stratton1941}.
\end{physicalvignette}
\begin{equation}
\label{eq:ch2.retained-angular-amplitude-def}
a_{\mathrm{ret}}(\widehat{\mathbf{k}})
=
\chi_{\mathcal{K}_{\mathrm{acc}}}(\widehat{\mathbf{k}})\,
a(\widehat{\mathbf{k}}),
\qquad
\chi_{\mathcal{K}_{\mathrm{acc}}}(\widehat{\mathbf{k}})
=
\begin{cases}
1, & \widehat{\mathbf{k}}\in\mathcal{K}_{\mathrm{acc}},\\
0, & \widehat{\mathbf{k}}\in\mathcal{K}_{\mathrm{sup}},
\end{cases}
\end{equation}
with \(a(\widehat{\mathbf{k}})\) from
Eq.~\eqref{eq:ch2.exterior-continuous-spectral-integral} after
Eq.~\eqref{eq:ch2.plane-wave-radiation-operator} \cite{stratton1941}.
Equation~\eqref{eq:ch2.retained-angular-amplitude-def} is the Subsection~\ref{sec:ch292} locked definition of
development of the retained factor in Eq.
Eq.~\eqref{eq:ch2.retained-suppressed-filter}: only directions not annihilated by
\(\mathcal{C}_{\mathrm{geo}}\) angular rank-$N$ recovery on the far shell (Subsection~\ref{sec:ch264}) carry
qualified transport data \cite{stratton1941}. Mathematically, \(a_{\mathrm{ret}}\) is the pointwise
product of the full angular measure with the accessible indicator; physically, it is the portion
of the far-zone angular distribution that geometry permits to propagate as export after radiation
selection \cite{harrington2001}. Null-space additions
Eq.~\eqref{eq:ch2.angular-spectrum-nr-invariance} leave \(a_{\mathrm{ret}}\) unchanged
\cite{stratton1941}. Suppressed angular amplitude. Define the \emph{suppressed angular amplitude} as the complementary sheet content:
\begin{equation}
\label{eq:ch2.suppressed-angular-amplitude-def}
a_{\mathrm{sup}}(\widehat{\mathbf{k}})
=
\chi_{\mathcal{K}_{\mathrm{sup}}}(\widehat{\mathbf{k}})\,
a(\widehat{\mathbf{k}}),
\qquad
\chi_{\mathcal{K}_{\mathrm{sup}}}=1-\chi_{\mathcal{K}_{\mathrm{acc}}}
\ \text{a.e.\ on}\ S^{2},
\end{equation}
using suppressed sheet Eq.~\eqref{eq:ch2.suppressed-angular-sheet-def}
\cite{stratton1941}. Equation~\eqref{eq:ch2.suppressed-angular-amplitude-def}
completes Eq.~\eqref{eq:ch2.retained-suppressed-filter}: for radiating export,
geometric forbidding Eq.~\eqref{eq:ch2.geometric-symmetry-operator-def} typically forces
\(a_{\mathrm{sup}}=\mathbf{0}\) even when near-zone scans show energy in related directions
\cite{stratton1941}. Suppression is therefore not synonymous with
\(\Jimp\in\ker\mathcal{R}_\omega\): Eq.~\eqref{eq:ch2.angular-spectrum-nullspace-law} annihilates
\(a\) everywhere, whereas Eq.~\eqref{eq:ch2.suppressed-angular-amplitude-def} removes selected
directions while \(a_{\mathrm{ret}}\) may remain nonzero \cite{harrington2001}.
Evanescent lateral content Eq.~\eqref{eq:ch2.evanescent-angular-null-law} does not populate
\(a_{\mathrm{sup}}\) on export: \(\Pi_{\mathrm{ev}}\) is removed before angular readout
\cite{stratton1941}.
Laboratory observation applies a finite-rank criterion to far-zone angular content. Define the
\emph{filtered angular amplitude} by
\begin{equation}
\label{eq:ch2.filtered-angular-amplitude-def}
a_{\mathrm{filt}}(\widehat{\mathbf{k}})
=
\big[\mathcal{P}_{\mathrm{meas}}\,a\big](\widehat{\mathbf{k}}),
\qquad
\mathcal{P}_{\mathrm{meas}}
=
\mathcal{O}_{\mathrm{meas}}\circ\mathcal{O}_{\mathrm{far}}\circ\mathcal{H}_{\mathrm{ang}}^{-1},
\end{equation}
where \(\mathcal{H}_{\mathrm{ang}}\) maps \(a(\widehat{\mathbf{k}})\) to the corresponding
far-zone transverse field on \(S^{2}\) and \(\mathcal{O}_{\mathrm{far}}\) is the far extractor
from Eq.~\eqref{eq:ch2.reconstruction-observation-map} \cite{stratton1941}.
Equation~\eqref{eq:ch2.filtered-angular-amplitude-def} is the angular-side restatement of the
third line in Eq.~\eqref{eq:ch2.retained-suppressed-filter}. When
\(\mathcal{O}_{\mathrm{meas}}\) is a finite set of plane-wave or modal probes,
\(\mathcal{P}_{\mathrm{meas}}\) has rank at most \(N\) and
\(\ker\mathcal{P}_{\mathrm{meas}}\supseteq\ker\mathcal{O}_{\mathrm{meas}}\big|_{\operatorname{ran}\mathcal{R}_\omega}\)
Eq.~\eqref{eq:ch2.observation-null-space} \cite{stratton1941}. Physically,
\(a_{\mathrm{filt}}\) is what a finite instrument resolves after aperture, noise, and processing;
it may vanish on subsets of \(\mathcal{K}_{\mathrm{acc}}\) where \(a_{\mathrm{ret}}\) is nonzero
\cite{harrington2001}.
\begin{equation}
\label{eq:ch2.filtered-rank-bound}
\mathrm{rank}(\mathcal{P}_{\mathrm{meas}})
\le
N,
\qquad
\dim\operatorname{span}\{a_{\mathrm{filt}}(\widehat{\mathbf{k}})\}
\le
N,
\end{equation}
Combining Eqs.~\eqref{eq:ch2.retained-angular-amplitude-def}--\eqref{eq:ch2.filtered-angular-amplitude-def}
yields the Chapter~\ref{ch:02} partition
\begin{equation}
\label{eq:ch2.angular-component-partition}
a(\widehat{\mathbf{k}})
=
a_{\mathrm{ret}}(\widehat{\mathbf{k}})
a_{\mathrm{sup}}(\widehat{\mathbf{k}}),
\qquad
a_{\mathrm{filt}}
=
\mathcal{P}_{\mathrm{meas}}\big[a_{\mathrm{ret}}+a_{\mathrm{sup}}\big],
\end{equation}
with \(a_{\mathrm{ret}}\) and \(a_{\mathrm{sup}}\) supported on disjoint sheets
\(\mathcal{K}_{\mathrm{acc}}\) and \(\mathcal{K}_{\mathrm{sup}}\)
\cite{stratton1941}. Equation~\eqref{eq:ch2.angular-component-partition} upgrades
Eq.~\eqref{eq:ch2.retained-suppressed-filter} to explicit locked definitions and
inserts the measurement projector \(\mathcal{P}_{\mathrm{meas}}\). The first line is an angular
decomposition of export; the second line states that observation acts on the sum after radiation
selection, not on raw near spectra \cite{stratton1941}. Budget
Eq.~\eqref{eq:ch2.angular-spectrum-composition-law} constrains
\(\operatorname{supp}(a)\subseteq\mathcal{K}_{\mathrm{acc}}\cap\{\text{resolvable by}\ \mathcal{M}_\omega\}\),
so \(a_{\mathrm{filt}}\) cannot recover directions outside that intersection even when raw
lateral transforms appear rich \cite{harrington2001}.
On the plane-wave side of Subsection~\ref{sec:ch291}, define the \emph{retained lateral spectrum}
\begin{equation}
\label{eq:ch2.retained-lateral-spectrum-def}
\widetilde{\E}_{\mathrm{ret}}(k_x,k_y)
=
\mathcal{F}_{\mathrm{lat}}\big[\mathcal{H}_{\mathrm{ang}}[a_{\mathrm{ret}}]\big](k_x,k_y)
\quad\text{on}\ k_t\le k_0,
\end{equation}
and the \emph{filtered lateral spectrum}
\begin{equation}
\label{eq:ch2.filtered-lateral-spectrum-def}
\widetilde{\E}_{\mathrm{filt}}(k_x,k_y)
=
\mathcal{F}_{\mathrm{lat}}\big[\mathcal{H}_{\mathrm{ang}}[a_{\mathrm{filt}}]\big](k_x,k_y),
\end{equation}
with \(\mathcal{F}_{\mathrm{lat}}\) from Eq.~\eqref{eq:ch2.plane-wave-angular-spectrum-def}
\cite{stratton1941}. Equations~\eqref{eq:ch2.retained-lateral-spectrum-def}--\eqref{eq:ch2.filtered-lateral-spectrum-def}
connect Subsection~\ref{sec:ch291} lateral readout to the angular partition without re-displaying
Eq.~\eqref{eq:ch2.plane-wave-radiation-operator}: \(\widetilde{\E}_{\mathrm{rad}}\) from
Eq.~\eqref{eq:ch2.radiative-lateral-spectrum-def} agrees with
\(\widetilde{\E}_{\mathrm{ret}}\) on the propagating outgoing sheet when
\(a_{\mathrm{sup}}=\mathbf{0}\) \cite{stratton1941}. On the evanescent sheet
\(k_t>k_0\), retained and filtered lateral spectra vanish through
Eq.~\eqref{eq:ch2.evanescent-angular-null-law} independently of raw
\(\mathcal{F}_{\mathrm{lat}}[\mathcal{S}_\omega[\Jimp]]\) \cite{stratton1941}.
regime qualification Eq.~\eqref{eq:ch2.regime-component-criterion} applies to
\(\widetilde{\E}_{\mathrm{filt}}\) before any transport claim \cite{harrington2001}.
On \(\mathcal{K}_{\mathrm{sup}}\),
\begin{equation}
\label{eq:ch2.suppressed-lateral-null-law}
a_{\mathrm{sup}}(\widehat{\mathbf{k}})=\mathbf{0}
\ \Longrightarrow\
\widetilde{\E}_{\mathrm{ret}}(k_x,k_y)=\widetilde{\E}_{\mathrm{rad}}(k_x,k_y)
\ \text{on}\ k_t\le k_0,
\end{equation}
and any lateral energy attributed to \(\mathcal{K}_{\mathrm{sup}}\) after
Eq.~\eqref{eq:ch2.plane-wave-radiation-operator} is removed by geometry, not by
\(\mathcal{P}_{\mathrm{meas}}\) \cite{stratton1941}. Equation~\eqref{eq:ch2.suppressed-lateral-null-law}
expresses that suppression precedes finite aperture cannot resurrect directions
forbidden by Eq.~\eqref{eq:ch2.geometric-symmetry-operator-def}
\cite{harrington2001}. Reciprocity in radiation and reception (Subsection~\ref{sec:ch283})
constrains bilinear couplings on \(\mathcal{K}_{\mathrm{acc}}\) but does not enlarge
\(\operatorname{supp}(a_{\mathrm{ret}})\) \cite{stratton1941}.
\begin{figure}[t]
\centering
\ChOneFigBlock{%
\begin{tikzpicture}[ch1fig, node distance=5mm and 7mm]
 \node[ch1box] (a) {$a(\widehat{\mathbf{k}})$};
 \node[ch1box, right=7mm of a] (ret) {$a_{\mathrm{ret}}$};
 \node[ch1box, right=7mm of ret] (filt) {$a_{\mathrm{filt}}$};
 \node[ch1box, below=5mm of a] (sup) {$a_{\mathrm{sup}}$};
 \node[ch1box, below=5mm of ret] (kacc) {$\mathcal{K}_{\mathrm{acc}}$};
 \node[ch1box, below=5mm of sup] (ksup) {$\mathcal{K}_{\mathrm{sup}}$};
 \node[ch1box, left=7mm of a] (rad) {$\Pi_{\mathrm{rad}}\mathcal{R}_\omega$};
 \node[ch1box, right=7mm of filt] (lat) {$\widetilde{\E}_{\mathrm{filt}}$};
 \draw[ch1flow] (rad.west) -- (rad.west |- a.east);
 \draw[ch1flow] (rad.west |- a.east) -- (a.east);
 \draw[ch1flow] (a) -- (ret) -- (filt) -- (lat);
 \draw[ch1flow] (a.south) -- (a.south |- sup.north);
 \draw[ch1flow] (a.south |- sup.north) -- (sup.north);
 \draw[ch1flow] (kacc.south) -- (kacc.south |- ret.north);
 \draw[ch1flow] (kacc.south |- ret.north) -- (ret.north);
 \draw[ch1flow] (ksup.south) -- (ksup.south |- sup.north);
 \draw[ch1flow] (ksup.south |- sup.north) -- (sup.north);
\end{tikzpicture}%
}
\caption{Subsection~\ref{sec:ch292} partition measurement sequence:
Eq.~\eqref{eq:ch2.angular-component-partition} splits \(a(\widehat{\mathbf{k}})\) into retained
and suppressed sheets, then applies observation filter
Eq.~\eqref{eq:ch2.filtered-angular-amplitude-def} before lateral readout
Eq.~\eqref{eq:ch2.filtered-lateral-spectrum-def}.}
\label{fig:ch2.retained-suppressed-filter-measurement chain}
\end{figure}
Figure interpretation. Radiative amplitude $a(\widehat{\mathbf{k}})$ splits into retained sheet $\mathcal{K}_{\mathrm{acc}}$ and suppressed sheet $\mathcal{K}_{\mathrm{sup}}$ before observation filter $\mathcal{P}_{\mathrm{meas}}$. Retained content survives geometric and parity criteria; suppressed content is zero by constraint, not by source nullity. Filtered lateral spectrum $\widetilde{\E}_{\mathrm{filt}}$ is the aperture-limited readout: finite rank may remove angular content even when $a\neq 0$ on $\mathcal{F}_{\mathrm{rad}}$ \cite{stratton1941}. Figure~\ref{fig:ch2.retained-suppressed-filter-measurement chain} orders radiation selection, angular
partition, measurement filter, and plane-wave readout. Suppression acts on
\(\mathcal{K}_{\mathrm{sup}}\) before \(\mathcal{P}_{\mathrm{meas}}\); evanescent sheets are
removed by Eq.~\eqref{eq:ch2.evanescent-angular-null-law} before lateral readout
\cite{stratton1941}.
Mathematical role of Eq.~\eqref{eq:ch2.angular-component-partition}. Equation~\eqref{eq:ch2.angular-component-partition} separates three stages that are often
conflated in angular-spectrum discussions: radiation export, geometric $\mathcal{E}_{\mathrm{acc}}^{(N)}$, and finite
observation \cite{stratton1941}. The retained factor encodes what the source radiates
on \(\mathcal{K}_{\mathrm{acc}}\); the suppressed factor marks directions removed before
measurement; the filtered factor applies \(\mathcal{P}_{\mathrm{meas}}\) \cite{stratton1941}.
Plane-wave coefficients \(\widetilde{\E}_{\mathrm{rad}}\) from
Eq.~\eqref{eq:ch2.radiative-lateral-spectrum-def} describe the same retained content in lateral
coordinates when \(a_{\mathrm{sup}}=\mathbf{0}\) \cite{jackson1999}. Null additions
plane-wave null invariance (Subsection~\ref{sec:ch291}) alter neither \(a_{\mathrm{ret}}\) nor
\(a_{\mathrm{filt}}\) when \(\mathcal{M}_\omega\) is unchanged \cite{stratton1941}.
Mathematical role of Eq.~\eqref{eq:ch2.filtered-angular-amplitude-def}. Equation~\eqref{eq:ch2.filtered-angular-amplitude-def} makes explicit that instruments observe
\(\mathcal{P}_{\mathrm{meas}}[a]\), not raw \(\mathcal{F}_{\mathrm{lat}}[\mathcal{S}_\omega[\Jimp]]\)
\cite{stratton1941}. Rank cap Eq.~\eqref{eq:ch2.filtered-rank-bound} prevents
attributing unobserved angular structure to source nullity: a pattern gap on
\(\mathcal{K}_{\mathrm{acc}}\) may arise from \(\mathcal{P}_{\mathrm{meas}}\) alone
\cite{harrington2001}. Reciprocity in radiation and reception (Subsection~\ref{sec:ch283}) constrains
coupling among observed channels but does not increase
\(\dim\operatorname{span}\{a_{\mathrm{filt}}\}\) \cite{stratton1941}. Regime qualification
Eq.~\eqref{eq:ch2.regime-operator-screen} remains mandatory before interpreting
\(a_{\mathrm{filt}}\) as transport \cite{harrington2001}. Modal reportability Eq.~\eqref{eq:ch2.constraint-composition-law} and angular budget
Eq.~\eqref{eq:ch2.angular-spectrum-composition-law} are the modal and continuous forms of the
same constraint chain Subsection~\ref{sec:ch292} displays in
Eq.~\eqref{eq:ch2.angular-component-partition} \cite{stratton1941}. plane-wave radiation-operator action (Subsection~\ref{sec:ch291}) supplies lateral coordinates; retained, suppressed, and filtered angular components (Subsection~\ref{sec:ch292}) supplies angular-sheet coordinates; \ref{ch:03}
Subsection~3.8.1 will change basis without altering
\(\operatorname{supp}(a_{\mathrm{ret}})\) \cite{stratton1941}. Inverse obstruction
angular spectra and inverse problems (Subsection~\ref{sec:ch274}) applies to
\(a_{\mathrm{filt}}\), not to raw lateral transforms \cite{harrington2001}. Retained, suppressed, and filtered angular components (Subsection~\ref{sec:ch292}) assumes linear media, fixed \(\omega\), outgoing
closure, and stationary \(\mathcal{C}_{\mathrm{geo}}\) \cite{stratton1941}. Time-varying
or nonlinear \(\mathcal{P}_{\mathrm{meas}}\) requires separate treatment; the static rank bound
Eq.~\eqref{eq:ch2.filtered-rank-bound} is an upper limit \cite{stratton1941}. Loss modifies
amplitudes but not the partition logic on \(\mathcal{K}_{\mathrm{acc}}\cup\mathcal{K}_{\mathrm{sup}}\)
\cite{stratton1941}. Cavity and waveguide spectra remain outside the exterior partition; discrete
books Eq.~\eqref{eq:ch2.cavity-discrete-radiation-spectrum} are not substituted for
Eq.~\eqref{eq:ch2.angular-component-partition} on unbounded domains \cite{jackson1999}.
\phantomsection\label{sec:ch293}
Section~\ref{sec:ch29} developed the near-to-far map through
Eq.~\eqref{eq:ch2.near-to-far-transform} and tied it to far-zone reduction
Eqs.~\eqref{eq:ch2.far-field-integral-reduction}--\eqref{eq:ch2.far-field-amplitude-integral}
and the far-field approximation of source integrals (Subsection~\ref{sec:ch252}) from Subsection~\ref{sec:ch252} \cite{stratton1941}. Subsection~\ref{sec:ch293} gives the
\ref{ch:02} account of the operator that extracts far angular amplitudes
\(a(\widehat{\mathbf{k}})\) from radiative fields after
Subsections~\ref{sec:ch291}--\ref{sec:ch292}. Plane-wave lateral spectra
Eqs.~\eqref{eq:ch2.radiative-lateral-spectrum-def}--\eqref{eq:ch2.plane-wave-radiation-operator},
angular partition Eq.~\eqref{eq:ch2.angular-component-partition}, evanescent null law
Eq.~\eqref{eq:ch2.evanescent-angular-null-law}, and continuous representation
Eq.~\eqref{eq:ch2.exterior-continuous-spectral-integral} are presupposed; asymptotic proofs in
Subsection~\ref{sec:ch252} are not repeated \cite{stratton1941}. Observation-shell regime map
Eq.~\eqref{eq:ch2.observation-shell-regime-map} and regime qualification
Eq.~\eqref{eq:ch2.regime-component-criterion} qualify when the limit defines transport data
\cite{harrington2001}. Far-field angular probe rank belongs to
Subsection~\ref{sec:ch294} \cite{stratton1941}.
Assumptions. Fix \(\omega>0\), admissible \(\Jimp\in\mathcal{J}_\Omega\), outgoing radiative field
\(\mathbf{u}_\omega=\mathcal{R}_\omega[\Jimp]\in\mathcal{F}_{\mathrm{rad}}\), and distance
thresholds from Eq.~\eqref{eq:ch2.distance-regime-definitions} with \(r>r_{\min}(a,\lambda_0)\)
for source support radius \(a\) \cite{stratton1941}. Angular amplitude
\(a(\widehat{\mathbf{k}})\) is measured on the unit sphere with outgoing transverse polarization
\cite{stratton1941}. Regime qualification \(\mathcal{D}_{J,\mathrm{reg}}=1\) from
Eq.~\eqref{eq:ch2.regime-operator-screen} is presupposed before interpreting any near-to-far
output as transport \cite{harrington2001}. \begin{equation}
\label{eq:ch2.near-to-far-shell-error}
\|a_r-a\|_{L^{2}(S^{2})}
\le
C\,\chi(r),
\qquad
a_r(\widehat{\mathbf{k}})
=
\mathcal{H}_{\mathrm{ff}}\big[\mathcal{R}_\omega[\Jimp]\big](r,\widehat{\mathbf{k}}),
\end{equation}
% PHYSICAL-MEMORY:sec:ch293
% PHYSICAL-MEMORY:sec:ch293-b2
\begin{physicalvignette}[Beth torsion balance in a microwave anechoic chamber]
A second bench observation matters on Outdoor far-field range with mast and ground reflection. When technicians swap only the receive probe while holding the transmit chain fixed, reciprocity checks whether the reported coupling is property of the structure or artifact of a mismatched flux gauge \cite{stratton1941}. Failed reciprocity usually traces to a near-field tap, a radome standing-wave, or a probe polarization basis that was never aligned with the transmit standard \cite{hansen1988}.
Rayleigh and Stratton treated such ranges as boundary-value laboratories long before digital near-field transforms became routine \cite{stratton1941}. Harrington's method-of-moments culture reinforced the same discipline: metal outline first, harmonic headline second \cite{harrington2001}.
The measurement logic above is made precise in Subsection~\ref{sec:ch293} \cite{stratton1941}.
\end{physicalvignette}
\begin{physicalvignette}[Beth torsion balance in a microwave anechoic chamber]
Consider Beth torsion balance in a microwave anechoic chamber. Before any coefficient table is printed, technicians record frequency, cable loss, and the mechanical outline that fixes electrical size. The first walk-through of the range is always tactile: flange bolts, probe standoff, turntable boresight, and the wattmeter reference on the receive chain \cite{hansen1988}.
On this hardware the near zone is visually impressive---standing reactive structure, rapid phase rotation, sometimes a visible VSWR ripple at the feed---yet the export question is measured only after outgoing selection on a large sphere \cite{harrington2001}. What the camera sees in the chamber is not automatically what the far-zone pattern claims \cite{stratton1941}.
Engineers therefore learn to separate three readings: what the feed can synthesize on declared support, what propagating content survives on the radiative range, and what a finite instrument can recover at declared tolerance \cite{devaney2004}. Confusing those readings produces the rank gaps and watt failures that later sections quantify \cite{pal2026}.
Rayleigh and Stratton treated such ranges as boundary-value laboratories long before digital near-field transforms became routine \cite{stratton1941}. Harrington's method-of-moments culture reinforced the same discipline: metal outline first, harmonic headline second \cite{harrington2001}.
Subsection~\ref{sec:ch293} states the formal export condition that this range scene presupposes \cite{stratton1941}.
\end{physicalvignette}
\begin{equation}
\label{eq:ch2.near-to-far-integral-agreement}
\lim_{r\to\infty}
\mathcal{H}_{\mathrm{ff}}\big[\mathcal{R}_\omega[\Jimp]\big](r,\widehat{\mathbf{k}})
=
a(\widehat{\mathbf{k}}),
\end{equation}
On observation shells \(S_r=\{|\mathbf{r}|=r\}\) with \(r\gg a\) for source radius \(a\), the
far-zone post-processor must isolate the transverse \(e^{-\jj k_0 r}/r\) coefficient of
\(\E(\mathbf{r})\) in direction \(\widehat{\mathbf{k}}=\mathbf{r}/r\). Define the
\emph{far-zone extractor} \(\mathcal{H}_{\mathrm{ff}}\) by
\begin{equation}
\label{eq:ch2.far-zone-extractor-def}
\mathcal{H}_{\mathrm{ff}}[\mathbf{u}_\omega](r,\widehat{\mathbf{k}})
=
r\,e^{\jj k_0 r}\,
\widehat{\mathbf{k}}\times\Big(\widehat{\mathbf{k}}\times\E(\mathbf{r})\Big)\Big|_{|\mathbf{r}|=r},
\qquad
\mathbf{r}=r\,\widehat{\mathbf{k}},
\end{equation}
with \(\E\) the electric phasor in \(\mathbf{u}_\omega=(\E,\H)\) and transversality enforced by
the double cross product \cite{stratton1941}. Equation
\eqref{eq:ch2.far-zone-extractor-def} is the Subsection~\ref{sec:ch293} definitive specification of the
extractor stated in Eq.~\eqref{eq:ch2.near-to-far-transform}: it isolates the
outgoing radiative envelope compatible with Eq.~\eqref{eq:ch2.far-field-integral-reduction}
\cite{stratton1941}. Mathematically, \(\mathcal{H}_{\mathrm{ff}}\) is a linear functional on
shell traces; physically, it is the standard far-field post-processor that removes reactive
\(r^{-2}\) and \(r^{-3}\) contamination when \(k_0 r\gg 1\)
\cite{harrington2001}. The extractor acts on \(\mathcal{R}_\omega[\Jimp]\), not on raw
\(\mathcal{S}_\omega[\Jimp]\) Eq.~\eqref{eq:ch2.observation-shell-regime-map}
\cite{stratton1941}.
The far-zone limit extracts the \(e^{-\jj k_0 r}/r\) transverse envelope from
\(\mathcal{R}_\omega[\Jimp]\) on observation shells \(S_r\); the angular amplitude
\(a(\widehat{\mathbf{k}})\) is the coefficient of that envelope in direction
\(\widehat{\mathbf{k}}\). Define the \emph{near-to-far transformation} on admissible sources by
\begin{equation}
\label{eq:ch2.near-to-far-transform-def}
a(\widehat{\mathbf{k}})
=
\lim_{r\to\infty}
\mathcal{H}_{\mathrm{ff}}\big[\mathcal{R}_\omega[\Jimp]\big](r,\widehat{\mathbf{k}}),
\end{equation}
with \(a(\widehat{\mathbf{k}})\) the angular amplitude entering
Eq.~\eqref{eq:ch2.exterior-continuous-spectral-integral} \cite{stratton1941}.
Equation~\eqref{eq:ch2.near-to-far-transform-def} develops Eq.
Eq.~\eqref{eq:ch2.near-to-far-transform} to explicit extractor
Eq.~\eqref{eq:ch2.far-zone-extractor-def} and fixes the radiation map input
\cite{stratton1941}. The limit exists and is independent of the particular
\(r\)-path when the far-field approximation of source integrals (Subsection~\ref{sec:ch252}) applies; error on finite shells
scales as in Eq.~\eqref{eq:ch2.power-flux-regime-convergence}
\cite{stratton1941}. Mathematically, Eq.~\eqref{eq:ch2.near-to-far-transform-def} is the
angular-side completion of far-zone template
Eq.~\eqref{eq:ch2.far-field-integral-reduction}; physically, it states that far angular content
is determined by the radiative branch alone after outgoing selection, independent of reactive
storage that may dominate on inner shells \cite{stratton1941}. Observation consequence: SNF and
compact-range transforms that skip \(\mathcal{R}_\omega\) before extraction inherit
Eq.~\eqref{eq:ch2.near-to-far-evanescent-null} errors documented in
Subsection~\ref{sec:ch2103} \cite{hansen1988}. Regime qualification
Eq.~\eqref{eq:ch2.regime-component-criterion} remains mandatory before interpreting the limit as
transport data \cite{harrington2001}.
Package Eq.~\eqref{eq:ch2.near-to-far-transform-def} as the composed operator
\begin{equation}
\label{eq:ch2.near-to-far-operator-composition}
\Lambda_{\mathrm{nf}}
=
\mathcal{H}_{\mathrm{ang}}\circ\mathcal{H}_{\mathrm{ff}}^{\infty}\circ\mathcal{R}_\omega,
\qquad
a(\widehat{\mathbf{k}})=\Lambda_{\mathrm{nf}}[\Jimp](\widehat{\mathbf{k}}),
\end{equation}
where \(\mathcal{H}_{\mathrm{ff}}^{\infty}\) denotes the \(r\to\infty\) limit in
Eq.~\eqref{eq:ch2.near-to-far-transform-def} and \(\mathcal{H}_{\mathrm{ang}}\) maps the
transverse far vector in Eq.~\eqref{eq:ch2.far-zone-extractor-def} to scalar polarization
amplitudes \(a_\sigma(\widehat{\mathbf{k}})\) in
Eq.~\eqref{eq:ch2.exterior-continuous-spectral-integral}
\cite{stratton1941}. Equation~\eqref{eq:ch2.near-to-far-operator-composition}
aligns with far readout
\(\Lambda_{\mathrm{rad}}=\mathcal{O}_{\mathrm{far}}\circ\mathcal{R}_\omega\) from
Eq.~\eqref{eq:ch2.reconstruction-observation-map} when \(\mathcal{O}_{\mathrm{far}}\) implements
Eq.~\eqref{eq:ch2.far-zone-extractor-def} on the angular side \cite{stratton1941}. Measurement
\(\mathcal{M}_\omega=\mathcal{O}_{\mathrm{meas}}\circ\Lambda_{\mathrm{rad}}\) therefore factors
near-to-far extraction before finite-rank observation
\cite{harrington2001}. Null additions \(\Jimp^{\mathrm{(NR)}}\in\ker\mathcal{R}_\omega\)
leave \(\Lambda_{\mathrm{nf}}[\Jimp]\) unchanged through
Eq.~\eqref{eq:ch2.angular-spectrum-nr-invariance} \cite{stratton1941}.
On the propagating outgoing sheet \(k_t\le k_0\), the near-to-far output agrees with retained
lateral content from Subsection~\ref{sec:ch292}:
\begin{equation}
\label{eq:ch2.near-to-far-lateral-link}
a(\widehat{\mathbf{k}})
\ \longleftrightarrow\
\widetilde{\E}_{\mathrm{ret}}(k_x,k_y)
\ \text{under}\ (k_x,k_y)\mapsto\widehat{\mathbf{k}},
\end{equation}
when \(a_{\mathrm{sup}}=\mathbf{0}\) and plane-wave continuous-spectrum consistency (Subsection~\ref{sec:ch291}) applies \cite{stratton1941}.
Equation~\eqref{eq:ch2.near-to-far-lateral-link} connects Subsection~\ref{sec:ch293} angular
limit to Subsection~\ref{sec:ch291} lateral spectra without re-displaying
Eq.~\eqref{eq:ch2.plane-wave-radiation-operator}: both describe
\(\Pi_{\mathrm{rad}}\mathcal{R}_\omega[\Jimp]\) in different coordinates
\cite{stratton1941}. Filtered content satisfies
\(a_{\mathrm{filt}}=\mathcal{P}_{\mathrm{meas}}[a]\) with \(a\) from
Eq.~\eqref{eq:ch2.near-to-far-transform-def} \cite{stratton1941}.
Evanescent annihilation before near-to-far. Evanescent near storage does not enter Eq.~\eqref{eq:ch2.near-to-far-transform-def}:
\begin{equation}
\label{eq:ch2.near-to-far-evanescent-null}
\Pi_{\mathrm{ev}}\mathcal{S}_\omega[\Jimp]
\ \not\mapsto\
a(\widehat{\mathbf{k}})
\quad\text{when}\quad
\mathcal{R}_\omega[\Jimp]\ \text{is unchanged},
\end{equation}
by Eqs.~\eqref{eq:ch2.evanescent-angular-null-law},~\eqref{eq:ch2.evanescent-angular-decoupling},
and~\eqref{eq:ch2.rad-null-on-evanescent} \cite{stratton1941}. Equation
\eqref{eq:ch2.near-to-far-evanescent-null} expresses Eq.
Eq.~\eqref{eq:ch2.near-to-far-transform}: reactive \(\Pi_{\mathrm{ev}}\) content may be
large in the near zone Eq.~\eqref{eq:ch2.near-zone-field-scaling} while the far limit reports
only qualified export \cite{stratton1941}. regime qualification
Eq.~\eqref{eq:ch2.regime-component-criterion} forbids promoting \(\Pi_{\mathrm{ev}}\) storage to
\(a(\widehat{\mathbf{k}})\) \cite{harrington2001}.
\begin{figure}[t]
\centering
\ChOneFigBlock{%
\begin{tikzpicture}[ch1fig, node distance=5mm and 7mm]
 \node[ch1box] (j) {$\Jimp$};
 \node[ch1box, right=7mm of j] (r) {$\mathcal{R}_\omega$};
 \node[ch1box, right=7mm of r] (near) {Near shell\\$r\sim r_c$};
 \node[ch1box, right=7mm of near] (hff) {$\mathcal{H}_{\mathrm{ff}}$};
 \node[ch1box, right=7mm of hff] (a) {$a(\widehat{\mathbf{k}})$};
 \node[ch1box, below=5mm of r] (pi) {$\Pi_{\mathrm{rad}}$};
 \node[ch1box, below=5mm of near] (ev) {$\Pi_{\mathrm{ev}}=\mathbf{0}$};
 \draw[ch1flow] (j) -- (r) -- (near) -- (hff) -- (a);
 \draw[ch1flow] (pi.south) -- (pi.south |- r.north);
 \draw[ch1flow] (pi.south |- r.north) -- (r.north);
 \draw[ch1flow, dashed] (ev.south) -- (near.north);
\end{tikzpicture}%
}
\caption{Subsection~\ref{sec:ch293} near-to-far measurement sequence:
Eq.~\eqref{eq:ch2.near-to-far-transform-def} applies \(\mathcal{H}_{\mathrm{ff}}\) on
\(\mathcal{R}_\omega[\Jimp]\) after \(r\to\infty\); evanescent content does not feed
Eq.~\eqref{eq:ch2.near-to-far-evanescent-null}.}
\label{fig:ch2.near-to-far-measurement chain}
\end{figure}
Figure interpretation. Near shell observation at $r\sim r_c$ samples $\mathcal{R}_\omega[\Jimp]$ after $\Pi_{\mathrm{rad}}$, not raw $\mathcal{S}_\omega[\Jimp]$. Extractor $\mathcal{H}_{\mathrm{ff}}$ implements the $r\to\infty$ limit of the far-field approximation of source integrals (Subsection~\ref{sec:ch252}); dashed $\Pi_{\mathrm{ev}}=\mathbf{0}$ marks evanescent content that may be large on the inner shell yet does not feed $a(\widehat{\mathbf{k}})$. The measurement sequence confirms when a near-field-to-far-field transform is structurally valid versus when reactive storage corrupts the asymptotic limit \cite{stratton1941}.
Figure~\ref{fig:ch2.near-to-far-measurement chain} orders radiation selection, finite-shell observation,
far-zone extraction, and angular amplitude. Near reactive content is dashed: it may appear on
inner shells but is annihilated in the limit when export is unchanged
Mathematical role of Eq.~\eqref{eq:ch2.near-to-far-transform-def}. Equation~\eqref{eq:ch2.near-to-far-transform-def} is the angular-spectrum-side completion of the
far-zone template Eq.~\eqref{eq:ch2.far-field-integral-reduction}: it converts radiative fields
to continuous amplitudes \(a(\widehat{\mathbf{k}})\) used in
Eq.~\eqref{eq:ch2.exterior-continuous-spectral-integral} \cite{stratton1941}. The
map is not an optional post-processor on raw near data: input is
\(\mathcal{R}_\omega[\Jimp]\), consistent with observation-shell map
Eq.~\eqref{eq:ch2.observation-shell-regime-map} \cite{stratton1941}. Plane-wave lateral spectra
Eq.~\eqref{eq:ch2.radiative-lateral-spectrum-def} and angular limit
Eq.~\eqref{eq:ch2.near-to-far-lateral-link} are equivalent coordinate descriptions of the same
export \cite{jackson1999}. Reciprocity in radiation and reception (Subsection~\ref{sec:ch283}) constrains bilinear couplings built from
\(a(\widehat{\mathbf{k}})\) but does not modify Eq.~\eqref{eq:ch2.near-to-far-evanescent-null}
\cite{stratton1941}. Mathematical role of Eq.~\eqref{eq:ch2.far-zone-extractor-def}. Equation~\eqref{eq:ch2.far-zone-extractor-def} makes explicit how laboratory far-field
algorithms isolate \(\E_0(\theta,\phi)\) before pattern integration or modal projection
\cite{stratton1941}. Applying the extractor to
\(\mathcal{S}_\omega[\Jimp]\) on near shells confuses reactive and radiative scaling
Eq.~\eqref{eq:ch2.near-zone-field-scaling}; applying it to \(\mathcal{R}_\omega[\Jimp]\) aligns
with power-flux convergence Eq.~\eqref{eq:ch2.power-flux-regime-convergence}
\cite{stratton1941}. Filtered observation
Eq.~\eqref{eq:ch2.filtered-angular-amplitude-def} composes after
Eq.~\eqref{eq:ch2.near-to-far-transform-def}, not before radiation selection
\cite{harrington2001}. Modal far pattern Eq.~\eqref{eq:ch2.spectral-far-pattern-expansion} and angular limit
Eq.~\eqref{eq:ch2.near-to-far-transform-def} describe the same \(\operatorname{ran}\mathcal{R}_\omega\)
content under far-pattern coefficients from spectral modes (Subsection~\ref{sec:ch261})
\cite{stratton1941}. Retained/suppressed partition retained, suppressed, and filtered angular components (Subsection~\ref{sec:ch292}) constrains \(\operatorname{supp}(a)\) independently of
the near-to-far limit procedure \cite{stratton1941}. \ref{ch:03} harmonic
bases change coordinates without altering Eq.~\eqref{eq:ch2.near-to-far-operator-composition}
\cite{stratton1941}. Example~\ref{sec:ch2103} applies
Eq.~\eqref{eq:ch2.near-to-far-evanescent-null} in a worked setting \cite{stratton1941}. Near-to-far transformation (Subsection~\ref{sec:ch293}) assumes linear homogeneous or piecewise-homogeneous
media, fixed \(\omega\), outgoing radiation condition, and observation outside \(\Omega_s\)
\cite{stratton1941}. It fails in cutoff guides, cavity standing waves,
and lossless open resonances without outgoing selection \cite{stratton1941}. Surface waves and
strong anisotropy require modified extractors beyond
Eq.~\eqref{eq:ch2.far-zone-extractor-def} \cite{harrington2001}. The theorem does not assert
source uniqueness---only field asymptotics and angular amplitudes on the outgoing branch
\cite{stratton1941}. Cavity discrete spectra Eq.~\eqref{eq:ch2.cavity-discrete-radiation-spectrum}
do not replace Eq.~\eqref{eq:ch2.near-to-far-transform-def} on exterior domains
\cite{jackson1999}.
Section~\ref{sec:ch29} developed far-field angular rank-$N$ recovery through
Eq.~\eqref{eq:ch2.far-angular-accessibility-def} and required USE of directional
observables from Section~\ref{sec:ch01.2.1} and transport criteria from
Section~\ref{sec:ch01.2.3} \cite{stratton1941}. Subsection~\ref{sec:ch294} supplies the definitive completion for which directions on \(S^{2}\) carry
qualified, observable far-zone angular content after Subsections~\ref{sec:ch291}--\ref{sec:ch293}.
Near-to-far map Eq.~\eqref{eq:ch2.near-to-far-transform-def}, angular partition
Eq.~\eqref{eq:ch2.angular-component-partition}, accessible sector
Eq.~\eqref{eq:ch2.angular-accessibility-sector}, composition law
Eq.~\eqref{eq:ch2.angular-spectrum-composition-law}, and observation rank
Eq.~\eqref{eq:ch2.angular-spectrum-observation-rank} are presupposed; angular rank-$N$ recovery on the far shell (Subsection~\ref{sec:ch264}) is upgraded, not re-proved
\cite{stratton1941}. Directional flux maps Eqs.~\eqref{eq:ch1.directional-power-density}--\eqref{eq:ch1.directional-am-flux}
and transport triad Eq.~\eqref{eq:ch1.transport-qualified-decision} enter from the cited sections
\cite{stratton1941}. Explicit harmonic filters belong to \ref{ch:03};
Assumptions. Fix \(\omega>0\), admissible \(\Jimp\in\mathcal{J}_\Omega\), outgoing radiative field
\(\mathbf{u}_\omega=\mathcal{R}_\omega[\Jimp]\), angular amplitude \(a(\widehat{\mathbf{k}})\) from
Eq.~\eqref{eq:ch2.near-to-far-transform-def}, and regime qualification
\(\mathcal{D}_{J,\mathrm{reg}}=1\) from Eq.~\eqref{eq:ch2.regime-operator-screen}
\cite{stratton1941}. Accessible sector \(\mathcal{K}_{\mathrm{acc}}\) and
suppressed sheet \(\mathcal{K}_{\mathrm{sup}}\) are fixed by
Eqs.~\eqref{eq:ch2.angular-accessibility-sector}--\eqref{eq:ch2.suppressed-angular-sheet-def}
\cite{stratton1941}. Finite-rank measurement \(\mathcal{M}_\omega\) from
Eq.~\eqref{eq:ch2.reconstruction-observation-map} is presupposed \cite{harrington2001}.
\subsection{Far-field angular accessibility}
\phantomsection\label{sec:ch294}
Far angular $\mathcal{E}_{\mathrm{acc}}^{(N)}$ is what a finite probe count on $\Gamma_R$ can resolve after
near-to-far extraction and geometric sectoring---not what a harmonic chart alone can label
\cite{harrington2001}. Synthesis rank and observation rank must not be merged: the probe bank on
$\Gamma_R$ sets accessible directions before harmonic labels attach \cite{harrington2001}.
Chapter~3 relabels the same accessible set in VSH coordinates \cite{stratton1941}.
\medskip
\noindent\textit{probe rank.} Mount \(N=48\) probes on a spherical surface \(\Gamma_R\) at
\(R=3\lambda_0\) enclosing a \(\lambda_0/4\) dipole feed, and truncate the harmonic chart at
\(\ell\le10\). The dipole exports dominant \(\ell=1\) content, so the fitted coefficients
\(\widehat{c}_{10,m}\) are nonzero on the accessible sector; nevertheless, rank-$N$ recovery is
bounded jointly by probe count \(N\), scan geometry, and the truncated basis---not by the raw
richness of the near-field phase on \(\Gamma_R\) \cite{hansen1988}.
Far-field $\mathcal{E}_{\mathrm{acc}}^{(N)}$ criterion. Define \emph{far-field angular probe rank} by
\phantomsection\label{eq:ch2.angular-accessibility-sector}
% PHYSICAL-MEMORY:sec:ch294
\begin{physicalvignette}[Spherical near-field scanner with dual-polarization probes]
Picture a spherical near-field scanner with sixty-four dual-polarization probes on a fixed radius and a pyramidal horn under test on a rolled mount. The transform can print hundreds of harmonic indices after one overnight acquisition, yet the hardware can resolve only as many independent directions as the probe lattice and signal-to-noise ratio allow. When marketing slides claim ``full sphere coverage,'' the range chief asks for the singular-value plot of the probe coupling matrix, not the length of the coefficient vector \cite{hansen1988}.
Accessibility on this instrument is therefore a composed question. The horn may synthesize sidelobe structure on its declared mouth outline, while the scanner recovers a smaller set of propagating directions once evanescent content is filtered and probe rank is applied \cite{stratton1941}. A campaign that reports only synthesis residual while omitting probe count has merged two ceilings that the compact-range community has kept separate for decades \cite{pal2026}.
Beth's torsion-balance tradition reminds us that angular momentum reports on such ranges must survive a mechanical rotation check, not only a chart relabelling \cite{beth1936}. Hansen's near-field measurement monograph codified the same discipline for spherical scanners: declare radius, probe count, polarization basis, and propagating filter before interpreting any modal bank \cite{hansen1988}.
Subsection~\ref{sec:ch294} closes the far-field angular accessibility criterion that this scanner implements on the range \cite{stratton1941}.
\end{physicalvignette}
\begin{equation}
\label{eq:ch2.far-angular-accessibility-def}
\widehat{\mathbf{k}}\in\mathcal{K}_{\mathrm{acc}}^{\mathrm{ff}}
\ \Longleftrightarrow\
\Big[\,
a(\widehat{\mathbf{k}})\neq 0
\ \text{after}\ \Lambda_{\mathrm{nf}}
\ \text{and}\ \mathcal{D}_{J,\mathrm{reg}}=1
\,\Big]
\ \text{and}\ 
\widehat{\mathbf{k}}\in\mathcal{K}_{\mathrm{acc}},
\end{equation}
with \(\Lambda_{\mathrm{nf}}\) from Eq.~\eqref{eq:ch2.near-to-far-operator-composition} and
\(\mathcal{K}_{\mathrm{acc}}\) from Eq.~\eqref{eq:ch2.angular-accessibility-sector}
\cite{stratton1941}. Equation~\eqref{eq:ch2.far-angular-accessibility-def}
develops Eq.~\eqref{eq:ch2.far-angular-accessibility-def} by inserting explicit
near-to-far qualification and geometric sector membership \cite{stratton1941}. Mathematically,
\(\mathcal{K}_{\mathrm{acc}}^{\mathrm{ff}}\subseteq\mathcal{K}_{\mathrm{acc}}\) is the
far-field reportable subset after radiation and regime qualifications; physically, it marks directions
where angular amplitude survives export and lies in the geometrically accessible hemisphere
\cite{harrington2001}. Directions in \(\mathcal{K}_{\mathrm{sup}}\) are excluded
before observation composition \cite{stratton1941}.
Angular probe rank operator. Package Eq.~\eqref{eq:ch2.far-angular-accessibility-def} as the \emph{$\mathcal{E}_{\mathrm{acc}}^{(N)}$ projector}
\begin{equation}
\label{eq:ch2.angular-accessibility-operator}
\mathcal{A}_{\mathrm{ff}}
=
\chi_{\mathcal{K}_{\mathrm{acc}}}\circ\mathcal{D}_{J,\mathrm{reg}}\circ\Lambda_{\mathrm{nf}},
\qquad
a_{\mathrm{acc}}(\widehat{\mathbf{k}})
=
\big[\mathcal{A}_{\mathrm{ff}}[\Jimp]\big](\widehat{\mathbf{k}}),
\end{equation}
with \(a_{\mathrm{acc}}\) the accessible angular amplitude before finite measurement
\cite{stratton1941}. Equation~\eqref{eq:ch2.angular-accessibility-operator}
aligns retained amplitude Eq.~\eqref{eq:ch2.retained-angular-amplitude-def} with near-to-far
output when \(\operatorname{supp}(a)\subseteq\mathcal{K}_{\mathrm{acc}}\) near-to-far transformation (Subsection~\ref{sec:ch293})(iv) \cite{stratton1941}. Filtered observation
Eq.~\eqref{eq:ch2.filtered-angular-amplitude-def} applies after
Eq.~\eqref{eq:ch2.angular-accessibility-operator}:
\(a_{\mathrm{filt}}=\mathcal{P}_{\mathrm{meas}}[a_{\mathrm{acc}}]\)
\cite{harrington2001}. Null additions leave \(a_{\mathrm{acc}}\) unchanged through
Eq.~\eqref{eq:ch2.angular-spectrum-nr-invariance} \cite{stratton1941}.
On lateral coordinates, define the \emph{plane-wave probe rank filter} by
\begin{equation}
\label{eq:ch2.plane-wave-accessibility-filter}
\widetilde{\E}_{\mathrm{acc}}(k_x,k_y)
=
\chi_{\mathcal{K}_{\mathrm{acc}}^{\mathrm{ff}}}(\widehat{\mathbf{k}})\,
\widetilde{\E}_{\mathrm{rad}}(k_x,k_y),
\qquad
(k_x,k_y)\mapsto\widehat{\mathbf{k}}\ \text{on}\ k_t\le k_0,
\end{equation}
with \(\widetilde{\E}_{\mathrm{rad}}\) from Eq.~\eqref{eq:ch2.radiative-lateral-spectrum-def}
\cite{stratton1941}. Equation~\eqref{eq:ch2.plane-wave-accessibility-filter}
is the Subsection~\ref{sec:ch294} plane-wave criterion promised in
Eq.~\eqref{eq:ch2.far-angular-accessibility-preview}: it removes geometrically forbidden
directions on the propagating outgoing sheet without re-tabulating
Eq.~\eqref{eq:ch2.plane-wave-radiation-operator} \cite{stratton1941}. Evanescent sheets satisfy
\(\widetilde{\E}_{\mathrm{acc}}=\mathbf{0}\) for \(k_t>k_0\) through
Eq.~\eqref{eq:ch2.evanescent-angular-null-law} \cite{stratton1941}. Physically,
Eq.~\eqref{eq:ch2.plane-wave-accessibility-filter} is the angular-spectrum-side statement of which
plane-wave directions participate in far-zone reports \cite{jackson1999}.
Far-field $\mathcal{E}_{\mathrm{acc}}^{(N)}$ composes with \ref{ch:01} transport qualification by reference.
A direction \(\widehat{\mathbf{k}}\in\mathcal{K}_{\mathrm{acc}}^{\mathrm{ff}}\) supports a
qualified transport report on solid-angle sector \(\Delta\Omega(\widehat{\mathbf{k}})\) only when
Section~\ref{sec:ch01.2.3} propagative support Eq.~\eqref{eq:ch1.propagative-support} and
transport decision Eq.~\eqref{eq:ch1.transport-qualified-decision} hold for the corresponding
directional channels Eqs.~\eqref{eq:ch1.directional-power-density}--\eqref{eq:ch1.directional-am-flux}
\cite{stratton1941}. Restated in operator form,
\begin{equation}
\label{eq:ch2.far-angular-transport-criterion}
\mathcal{T}_{\mathrm{acc}}(\widehat{\mathbf{k}})=1
\ \Longrightarrow\
\int_{\Delta\Omega(\widehat{\mathbf{k}})}\Pi(\theta,\phi;r)\,d\Omega>0
\ \text{and}\ 
\mathcal{Q}_{J}=1,
\end{equation}
with \(\Pi\) from Eq.~\eqref{eq:ch1.directional-power-density},
\(\mathcal{Q}_{J}\) from Eq.~\eqref{eq:ch1.torque-transport-qualified-gate}, and
\(r\) in the far zone where Eq.~\eqref{eq:ch2.near-to-far-transform-def} applies
\cite{stratton1941}. Equation~\eqref{eq:ch2.far-angular-transport-criterion} imports
Section~\ref{sec:ch01.2.3} criteria without re-deriving the transport triad: \ref{ch:02}
records when far-zone probe rank implies \ref{ch:01} transport admissibility on
the same sector \cite{stratton1941}. regime qualification Eq.~\eqref{eq:ch2.regime-component-criterion} remains
mandatory: \(\mathcal{T}_{\mathrm{acc}}=1\) does not promote reactive near content to export
\cite{harrington2001}.
Angular measurement criterion. After rank-$N$ recovery projection, finite instruments apply
\begin{equation}
\label{eq:ch2.angular-measurement-screen}
a_{\mathrm{obs}}(\widehat{\mathbf{k}})
=
\big[\mathcal{P}_{\mathrm{meas}}\,a_{\mathrm{acc}}\big](\widehat{\mathbf{k}}),
\qquad
\operatorname{supp}(a_{\mathrm{obs}})
\subseteq
\{\widehat{\mathbf{k}}:\ \text{resolvable by}\ \mathcal{M}_\omega\}
\cap
\mathcal{K}_{\mathrm{acc}}^{\mathrm{ff}},
\end{equation}
with \(\mathcal{P}_{\mathrm{meas}}\) from Eq.~\eqref{eq:ch2.filtered-angular-amplitude-def}
\cite{stratton1941}. Equation~\eqref{eq:ch2.angular-measurement-screen} is the
Subsection~\ref{sec:ch294} measurement completion: observation rank
Eq.~\eqref{eq:ch2.angular-spectrum-observation-rank} may strictly shrink
\(\operatorname{supp}(a_{\mathrm{obs}})\) inside \(\mathcal{K}_{\mathrm{acc}}^{\mathrm{ff}}\)
\cite{stratton1941}. Rank cap Eq.~\eqref{eq:ch2.filtered-rank-bound} limits resolvable directions
independently of source nullity \cite{harrington2001}. Physically, \(a_{\mathrm{obs}}\) is the
laboratory-accessible angular spectrum after geometry, regime, near-to-far extraction, and
aperture \cite{stratton1941}.
Accessible angular budget. Collecting Subsection~\ref{sec:ch294} constraints,
\phantomsection\label{eq:ch2.accessible-modal-sector}
\begin{equation}
\label{eq:ch2.accessible-angular-budget}
\operatorname{supp}(a_{\mathrm{obs}})
\subseteq
\mathcal{K}_{\mathrm{acc}}^{\mathrm{ff}}
\cap
\{\widehat{\mathbf{k}}:\ \mathcal{T}_{\mathrm{acc}}(\widehat{\mathbf{k}})=1\}
\cap
\{\widehat{\mathbf{k}}:\ \text{resolvable by}\ \mathcal{M}_\omega\},
\end{equation}
composing geometric, transport, and observation factors with angular budget
Eq.~\eqref{eq:ch2.angular-spectrum-composition-law} \cite{stratton1941}. Equation
\eqref{eq:ch2.accessible-angular-budget} is the Section~\ref{sec:ch29} completion of
Eq.~\eqref{eq:ch2.section29-angular-spectrum-chain}: each factor removes a distinct laboratory error
--- radiation nullity, geometric forbidding, regime violation, transport disqualification, and
finite aperture --- without conflating them \cite{stratton1941}. Reciprocity in radiation and reception (Subsection~\ref{sec:ch283}) constrains coupling among observed channels on
\(\operatorname{supp}(a_{\mathrm{obs}})\) but does not enlarge
Eq.~\eqref{eq:ch2.accessible-angular-budget} \cite{stratton1941}.
\begin{figure}[t]
\centering
\ChOneFigBlock{%
\begin{tikzpicture}[ch1fig, node distance=5mm and 7mm]
 \node[ch1box] (nf) {$\Lambda_{\mathrm{nf}}$};
 \node[ch1box, right=7mm of nf] (acc) {$\mathcal{A}_{\mathrm{ff}}$};
 \node[ch1box, right=7mm of acc] (t) {$\mathcal{T}_{\mathrm{acc}}$};
 \node[ch1box, right=7mm of t] (m) {$\mathcal{P}_{\mathrm{meas}}$};
 \node[ch1box, right=7mm of m] (obs) {$a_{\mathrm{obs}}$};
 \node[ch1box, below=5mm of acc] (k) {$\mathcal{K}_{\mathrm{acc}}$};
 \node[ch1box, below=5mm of t] (c1) {Subsection~\ref{sec:ch01.2.3}};
 \draw[ch1flow] (nf) -- (acc) -- (t) -- (m) -- (obs);
 \draw[ch1flow] (k.south) -- (k.south |- acc.north);
 \draw[ch1flow] (k.south |- acc.north) -- (acc.north);
 \draw[ch1flow, dashed] (c1.south) -- (t.north);
\end{tikzpicture}%
}
\caption{Subsection~\ref{sec:ch294} $\mathcal{E}_{\mathrm{acc}}^{(N)}$ measurement sequence:
Eq.~\eqref{eq:ch2.angular-accessibility-operator} applies near-to-far extraction, geometric
sector \(\mathcal{K}_{\mathrm{acc}}\), transport criterion Eq.~\eqref{eq:ch2.far-angular-transport-criterion}
(Subsection~\ref{sec:ch01.2.3}), and measurement criterion
Eq.~\eqref{eq:ch2.angular-measurement-screen}.}
\label{fig:ch2.far-angular-accessibility-pipeline}
\end{figure}
Figure interpretation. Full angular content on $S^{2}$ passes through geometric sector Eq.~\eqref{eq:ch2.accessible-modal-sector}, regime qualification Eq.~\eqref{eq:ch2.regime-operator-screen}, and finite observation rank Eq.~\eqref{eq:ch2.reconstruction-rank-cap}. The accessible output is the intersection of export-supported directions with aperture-resolvable modes---not the raw pattern from forward integration. Null-space additions leave $a(\widehat{\mathbf{k}})$ fixed but may alter near reactive traces used in calibration \cite{stratton1941}. Suppressed directions never enter \(a_{\mathrm{obs}}\); evanescent content is removed before
\(\Lambda_{\mathrm{nf}}\) by Eq.~\eqref{eq:ch2.near-to-far-evanescent-null}
\cite{stratton1941}.
Mathematical role of Eq.~\eqref{eq:ch2.far-angular-accessibility-def}. Equation~\eqref{eq:ch2.far-angular-accessibility-def} closes Section~\ref{sec:ch29} by stating
which directions on \(S^{2}\) remain after radiation selection, geometric constraints, and
regime qualification \cite{stratton1941}. It develops Eq.
Eq.~\eqref{eq:ch2.far-angular-accessibility-preview} with explicit dependence on
\(\Lambda_{\mathrm{nf}}\) and \(\mathcal{K}_{\mathrm{acc}}\) \cite{stratton1941}. Operator form
Eq.~\eqref{eq:ch2.angular-accessibility-operator} packages the same logic for source-to-amplitude
maps \cite{harrington2001}. \ref{ch:03} harmonic bases relabel
\(\mathcal{K}_{\mathrm{acc}}^{\mathrm{ff}}\) without altering the budget
Eq.~\eqref{eq:ch2.accessible-angular-budget} \cite{stratton1941}. Mathematical role of Eq.~\eqref{eq:ch2.plane-wave-accessibility-filter}. Equation~\eqref{eq:ch2.plane-wave-accessibility-filter} is the lateral-coordinate criterion
completing the plane-wave chain begun in Subsection~\ref{sec:ch291} \cite{stratton1941}.
It removes forbidden directions before near-to-far comparison and measurement
\cite{harrington2001}. Agreement with \(a_{\mathrm{acc}}\) on \(k_t\le k_0\) follows from
Eq.~\eqref{eq:ch2.near-to-far-lateral-link} \cite{stratton1941}. Evanescent sheets remain null
through Eq.~\eqref{eq:ch2.evanescent-angular-null-law} independently of raw lateral scans
\cite{stratton1941}.
Subsections~\ref{sec:ch291}--\ref{sec:ch294} complete
Eq.~\eqref{eq:ch2.section29-angular-spectrum-chain}:
plane-wave radiation operators, retained/suppressed/filtered partition, near-to-far extraction,
and far-field $\mathcal{E}_{\mathrm{acc}}^{(N)}$ \cite{stratton1941}. Modal coordinates
Eq.~\eqref{eq:ch2.spectral-far-pattern-expansion} and plane-wave coordinates
Eq.~\eqref{eq:ch2.exterior-continuous-spectral-integral} describe identical export on
\(\operatorname{ran}\mathcal{R}_\omega\); Section~\ref{sec:ch29} fixes angular-spectrum filters on
the plane-wave branch \cite{stratton1941}. Inverse obstruction in Subsection~\ref{sec:ch274} applies to \(a_{\mathrm{obs}}\),
not to pre-probe rank transforms \cite{harrington2001}. Subsection~\ref{sec:ch294} assumes linear media, fixed \(\omega\), outgoing closure, and
stationary geometry \(\mathcal{C}_{\mathrm{geo}}\) \cite{stratton1941}.
Time-varying \(\mathcal{M}_\omega\) requires separate treatment; static rank bounds remain upper
limits \cite{stratton1941}. Loss modifies amplitudes but not the inclusion chain in
Eq.~\eqref{eq:ch2.accessible-angular-budget} \cite{stratton1941}. Cavity and guide spectra remain
outside exterior rank-$N$ recovery; discrete books do not substitute for
Eq.~\eqref{eq:ch2.far-angular-accessibility-def} on unbounded domains
\cite{jackson1999}. Section~\ref{sec:ch01.2.3} transport criteria are imported, not
extended to near reactive zones without regime qualification \cite{stratton1941}.
The mathematical statements in Section~\ref{sec:ch29} state what compact-range engineers already enforce through gating, scan distance, and modal banks: far-field angular probe rank is the intersection of export-supported directions with instrument-resolvable modes, not the raw output of a forward integral \cite{harrington2001}. The important distinction is between spectral content present in $\mathcal{S}_\omega[\Jimp]$ and content that remains in $a_{\mathrm{obs}}$ after $\mathcal{K}_{\mathrm{acc}}$ and $\mathcal{O}_{\mathrm{meas}}$ have acted \cite{stratton1941}. Section~\ref{sec:ch210} returns to named hardware configurations so one may see nullity, rank limits, evanescent contamination, and structured-source topology as separate terms in one diagnostic partition rather than as a single pattern residual \cite{hansen1988}.
\section{Examples}
\label{sec:ch210}
% CH02-NARR-EXP:02.210
The four worked cases trace propagation, export, spectral rank, and instrument access in fixed order: compact dipole truncation, loop nullity, evanescent contamination, and structured-source topology \cite{stratton1941,harrington2001,devaney2004}. Each closes with an explicit verdict on what is radiated, null, or merely reactive before harmonic headlines enter \cite{chu1948,poynting1884}.
Angular-spectrum action in Section~\ref{sec:ch29} states what finite instruments resolve on outgoing fields. Section~\ref{sec:ch210} returns to named hardware---apertures, dipoles, arrays, and scans---with the same tolerance discipline \cite{stratton1941}. Each example is one integrated exercise, not a problem sheet \cite{collin2001}.
% PHYSICAL-MEMORY:sec:ch210
% PHYSICAL-MEMORY:sec:ch210-b2
\begin{physicalvignette}[Yagi array aimed at a fixed receive horn]
A second bench observation matters on Rectangular waveguide aperture radiating into free space. When technicians swap only the receive probe while holding the transmit chain fixed, reciprocity checks whether the reported coupling is property of the structure or artifact of a mismatched flux gauge \cite{stratton1941}. Failed reciprocity usually traces to a near-field tap, a radome standing-wave, or a probe polarization basis that was never aligned with the transmit standard \cite{hansen1988}.
Rayleigh and Stratton treated such ranges as boundary-value laboratories long before digital near-field transforms became routine \cite{stratton1941}. Harrington's method-of-moments culture reinforced the same discipline: metal outline first, harmonic headline second \cite{harrington2001}.
The measurement logic above is made precise in Subsection~\ref{sec:ch210} \cite{stratton1941}.
\end{physicalvignette}
\begin{physicalvignette}[Yagi array aimed at a fixed receive horn]
Consider Yagi array aimed at a fixed receive horn. Before any coefficient table is printed, technicians record frequency, cable loss, and the mechanical outline that fixes electrical size. The first walk-through of the range is always tactile: flange bolts, probe standoff, turntable boresight, and the wattmeter reference on the receive chain \cite{hansen1988}.
On this hardware the near zone is visually impressive---standing reactive structure, rapid phase rotation, sometimes a visible VSWR ripple at the feed---yet the export question is measured only after outgoing selection on a large sphere \cite{harrington2001}. What the camera sees in the chamber is not automatically what the far-zone pattern claims \cite{stratton1941}.
Engineers therefore learn to separate three readings: what the feed can synthesize on declared support, what propagating content survives on the radiative range, and what a finite instrument can recover at declared tolerance \cite{devaney2004}. Confusing those readings produces the rank gaps and watt failures that later sections quantify \cite{pal2026}.
Rayleigh and Stratton treated such ranges as boundary-value laboratories long before digital near-field transforms became routine \cite{stratton1941}. Harrington's method-of-moments culture reinforced the same discipline: metal outline first, harmonic headline second \cite{harrington2001}.
Subsection~\ref{sec:ch210} states the formal export condition that this range scene presupposes \cite{stratton1941}.
\end{physicalvignette}
The preceding sections developed radiation operators, null spaces, near-to-far limits, and
rank-$N$ instrument readouts as abstract objects on exterior domains. A graduate reader may accept the
algebra while still wondering how the distinctions appear on a laboratory bench: when a
spherical near-field scanner reports thirty-six complex samples, which part of the pattern
residual owes to modal truncation, which to source nullity, and which to reactive contamination
that never reaches the far zone? The four worked configurations below answer that question in
fixed order---compact dipole truncation, loop nullity, evanescent contamination, and
structured-source topology---each closing with an explicit verdict on what is radiated, null, or
merely reactive before harmonic headlines enter \cite{stratton1941}.
Subsections~\ref{sec:ch2101}--\ref{sec:ch2104} isolate compact truncation, source nullity,
evanescent near-to-far contamination, and structured-source topology--transport separation,
respectively \cite{stratton1941}. Every numerical ratio
quoted below is \emph{illustrative}: values follow author-calibrated discretizations of the stated
geometries and are not reported measurement data; they are included only to show order-of-magnitude
separation among mechanisms \cite{hansen1988}. Master DOF partition. Radiation reports can fail for distinct reasons that must not be merged in harmonic charts or laboratory
interpretation \cite{marengo2000}. Before the subsections activate individual terms, it is
useful to name the full partition once: each summand below removes degrees of freedom from
qualified radiation reports through a \emph{distinct} physical mechanism, and conflating them
produces remedies that treat symptoms rather than causes.
\begin{equation}
\label{eq:ch2.dof-partition-master}
\underbrace{\ker\mathcal{R}_\omega}_{\text{source nullity (NR)}}
\ \oplus\
\underbrace{\ker\mathcal{O}_{\mathrm{meas}}\big|_{\operatorname{ran}\mathcal{R}_\omega}}_{\text{aperture rank}}
\ \oplus\
\underbrace{\Pi_{\mathrm{ev}}\mathcal{S}_\omega[\mathcal{J}_\Omega]}_{\text{evanescent storage}}
\ \oplus\
\underbrace{\text{tail of compact/ modal truncation}}_{\text{spectral compression}}
\ \oplus\
\underbrace{\mathcal{K}_{\mathrm{sup}}}_{\text{geometric suppression}},
\end{equation}
with rank-$N$ recovery \(\mathcal{K}_{\mathrm{acc}}^{\mathrm{ff}}\) from
Eq.~\eqref{eq:ch2.accessible-angular-budget} acting after Eq.~\eqref{eq:ch2.dof-partition-master}
\cite{stratton1941}. Subsections~\ref{sec:ch2101}--\ref{sec:ch2104} each activate a
different term in Eq.~\eqref{eq:ch2.dof-partition-master}; cross-references replace repeated
local split displays \cite{stratton1941}.
Equation~\eqref{eq:ch2.dof-partition-master} is the diagnostic synthesis of
\ref{ch:02}: each direct summand names a distinct mechanism that removes
degrees of freedom from qualified radiation reports. Source nullity
\(\ker\mathcal{R}_\omega\) is invisible to all far data; aperture rank
\(\ker\mathcal{O}_{\mathrm{meas}}\) removes resolvable content even when export
is nonzero; evanescent storage inflates near readouts without altering
\(a(\widehat{\mathbf{k}})\); spectral compression tail controls modal truncation
error; geometric suppression \(\mathcal{K}_{\mathrm{sup}}\) removes directions
by symmetry before measurement acts. On a single hardware test with \(k_0 a=1.5\),
a \(15\%\) pattern residual might decompose as \(7\%\) compact tail, \(5\%\) aperture
rank, \(3\%\) geometric suppression, and \(0\%\) NR nullity---four different remedies
apply to four different terms \cite{stratton1941}.
Here $\mathcal{R}_\omega$ is the radiation operator at angular frequency $\omega$,
$\mathcal{O}_{\mathrm{meas}}$ the finite-rank rank-$N$ instrument readout on the radiative range,
$\Pi_{\mathrm{ev}}$ the projector onto evanescent spectral content,
$\mathcal{K}_{\mathrm{sup}}$ the geometric angular-suppression set from symmetry or blockage,
and the compact/modal tail the spectral compression induced by electrically small support and
finite probe banks \cite{stratton1941}. Flux meters from Section~\ref{sec:ch01.1.4} enter only on
\(\mathcal{R}_\omega[\Jimp]\) after Eq.~\eqref{eq:ch2.regime-operator-screen}
\subsection{Compactness of radiation operators}
\label{sec:ch2101}
This worked configuration addresses observability under finite instrumentation, not whether Maxwell permits radiation from an electrically small source. A short dipole in free space certainly radiates; the bench issue is how many independent complex samples a spherical near-field scan can resolve on the outgoing branch when the source ball and observation shell are both small compared with wavelength \cite{harrington2001}. Increasing modal order in a spherical-harmonic fit need not improve pattern agreement; however, Maxwell theory shows that rank limits arise jointly from compact Green smoothing on bounded domains and from the finite probe bank, while non-radiating nullity is a separate mechanism that leaves far data unchanged \cite{devaney2004}. Spectral truncation error must be distinguished from membership in $\ker\mathcal{R}_\omega$: the former shrinks when singular values of the compact part decay, the latter is invisible to any far-zone array regardless of modal count \cite{stratton1941}. What a laboratory instrument detects on $S_b$ at $b=2a$ is dominated by reactive falloff faster than $r^{-1}$; only after regime qualification and outgoing selection does the same campaign report a rank statement on $\mathcal{R}_\omega[\Jimp]$ \cite{hansen1988}.
% PHYSICAL-MEMORY:sec:ch2101
% PHYSICAL-MEMORY:sec:ch2101-b2
\begin{physicalvignette}[Circular patch on a handset ground plane]
A second bench observation matters on Helical feed on a satellite reflector. When technicians swap only the receive probe while holding the transmit chain fixed, reciprocity checks whether the reported coupling is property of the structure or artifact of a mismatched flux gauge \cite{stratton1941}. Failed reciprocity usually traces to a near-field tap, a radome standing-wave, or a probe polarization basis that was never aligned with the transmit standard \cite{hansen1988}.
Rayleigh and Stratton treated such ranges as boundary-value laboratories long before digital near-field transforms became routine \cite{stratton1941}. Harrington's method-of-moments culture reinforced the same discipline: metal outline first, harmonic headline second \cite{harrington2001}.
The measurement logic above is made precise in Subsection~\ref{sec:ch2101} \cite{stratton1941}.
\end{physicalvignette}
\begin{physicalvignette}[Circular patch on a handset ground plane]
Consider Circular patch on a handset ground plane. Before any coefficient table is printed, technicians record frequency, cable loss, and the mechanical outline that fixes electrical size. The first walk-through of the range is always tactile: flange bolts, probe standoff, turntable boresight, and the wattmeter reference on the receive chain \cite{hansen1988}.
On this hardware the near zone is visually impressive---standing reactive structure, rapid phase rotation, sometimes a visible VSWR ripple at the feed---yet the export question is measured only after outgoing selection on a large sphere \cite{harrington2001}. What the camera sees in the chamber is not automatically what the far-zone pattern claims \cite{stratton1941}.
Engineers therefore learn to separate three readings: what the feed can synthesize on declared support, what propagating content survives on the radiative range, and what a finite instrument can recover at declared tolerance \cite{devaney2004}. Confusing those readings produces the rank gaps and watt failures that later sections quantify \cite{pal2026}.
Worked examples should fix hardware outline, frequency, and tolerance before equations . 
Rayleigh and Stratton treated such ranges as boundary-value laboratories long before digital near-field transforms became routine \cite{stratton1941}. Harrington's method-of-moments culture reinforced the same discipline: metal outline first, harmonic headline second \cite{harrington2001}.
Subsection~\ref{sec:ch2101} states the formal export condition that this range scene presupposes \cite{stratton1941}.
\end{physicalvignette}
On a compact-range turntable, a center-fed wire dipole of length $L\ll\lambda$ and radius
$a\ll L$ sits inside a source ball $\Omega_s$ of radius $a$ with $k_0 a=1.5$ at the operating
frequency $\omega$ \cite{stratton1941}. A spherical near-field scanner records $N=36$ tangential
electric samples on a closed surface $S_b$ at radius $b=2a$; a vector network analyzer reports
$\VSWR$ and $|S_{11}|$ on the coax feed, confirming that stored reactive energy dominates the
input admittance at this electrical size \cite{hansen1988}. Far-field pattern cuts on the
principal planes show a single-lobe $\ell=1$ template, while probe magnitudes on $S_b$ fall off
faster than $r^{-1}$ in the reactive shell---the bench configuration on which compact truncation
is first isolated \cite{harrington2001}.
The compact-range configuration fixes outline, frequency, gauge, and permittivity before any
operator split is written. The objective is to connect one finite probe bank to one rank statement
on $\mathcal{R}_\omega[\Jimp]$: hardware fixes how many independent complex samples are available,
and the singular-value tail of the compact Green part controls how much of the radiative image
those samples can resolve \cite{stratton1941}. Watt closure and
probe-count limits must be reported separately; conflating synthesis ceilings with observation
ceilings is a common laboratory error on electrically small campaigns \cite{hansen1988}.
The first worked problem isolates \emph{spectral compression}: when a source and its observation
shell are electrically small, the Green operator that maps impressed current to tangential field
on a closed surface possesses a compact part whose singular values decay rapidly. Truncating that
part controls pattern-fitting error, but the error is \emph{not} non-radiating nullity---it is
finite-geometry flux report. Readers who have measured electrically short dipoles on compact
ranges will recognize the configuration immediately; the mathematics names what the instrument
already enforces through probe count and scan radius.
The bench places a center-fed wire of length $L\ll\lambda$ and radius $a\ll L$ at the center of a
compact-range turntable; the source ball $\Omega_s$ enclosing the wire has radius $a$ with
$k_0 a=1.5$ at the operating frequency $\omega$ \cite{stratton1941}. Port voltage on a coax feed
drives a triangle current distribution along $L$; a vector network analyzer reports $\VSWR$ and
$|S_{11}|$, confirming that stored reactive energy dominates the input admittance at this
electrical size \cite{hansen1988}. Far-field pattern cuts on the principal planes show a
single-lobe $\ell=1$ template; a spherical near-field probe at radius $b=2a$ records tangential
$\E$-field samples whose magnitude falls off faster than $r^{-1}$ in the reactive shell, so
watt-meter integration on a small sphere would overestimate radiated power unless the shell is
qualified \cite{chu1948}. The radiation integral yields $\ell=1$ dominance with higher multipoles
suppressed by $(k_0 a)^{2}$ \cite{stratton1941}; pattern cuts agree with the $\ell=1$ template,
while spherical near-field grids must sample within $\lambda/4$ to avoid aliasing the reactive
shell. Retaining only $M=20$ modes in a spherical-harmonic surrogate leaves a tail error set by the
twenty-first singular value of the compact Green part \cite{hansen1988}. The dipole is the sanity
check that every radiation map on compact support collapses toward a single export direction at
low $k_0 a$, while the near field remains rich with reactive structure that no far-zone array can
resolve \cite{harrington2001}.
Problem. A compact impressed-current ball \(\Omega_s=B_a(0)\) carries an electric dipole
\(\Jimp(\mathbf{r})=\mathbf{J}_0\,\delta(\mathbf{r})\) in the distributional limit with
\(\int_{\Omega_s}\Jimp\,d^{3}r=\mathbf{p}\) and \(k_0 a=1.5\) at fixed \(\omega\)
. A near-field scanner records \(N=36\) tangential electric
samples \(\mathbf{y}\in\mathbb{C}^{36}\) on a closed surface \(S_b\) with radius \(b=2a\)
through rank-$N$ instrument readout \(\mathcal{M}_\omega\) from
Eq.~\eqref{eq:ch2.reconstruction-observation-map} . The task is to:
\begin{enumerate}
\item[(a)] identify the compact and non-compact parts of the probe-to-far-field operator implied by
Subsection~\ref{sec:ch213};
\item[(b)] bound truncation error when only \(M=20\) retained modes are used in a spherical-harmonic
surrogate of \(\mathcal{R}_\omega[\Jimp]\);
\item[(c)] state whether unresolved modal content is attributable to compact compression, to
\(\ker\mathcal{R}_\omega\), or to finite rank of \(\mathcal{O}_{\mathrm{meas}}\).
\end{enumerate} \begin{figure}[t]
\centering
\ChOneFigBlock{%
\begin{tikzpicture}[ch1fig, node distance=5mm and 7mm]
 \node[ch1box] (j) {$\Jimp$\\$k_0a=1.5$};
 \node[ch1box, right=7mm of j] (s) {$\mathcal{S}_\omega$};
 \node[ch1box, right=7mm of s] (r) {$\mathcal{R}_\omega$};
 \node[ch1box, right=7mm of r] (m) {$\mathcal{M}_\omega$};
 \node[ch1box, right=7mm of m] (y) {$\mathbf{y}\in\mathbb{C}^{36}$};
 \node[ch1box, below=5mm of s] (split) {$G_{\mathrm{sing}}+K_{\mathrm{cpt}}$};
 \node[ch1box, below=5mm of m] (rank) {$\mathrm{rank}\le 36$};
 \draw[ch1flow] (j) -- (s) -- (r) -- (m) -- (y);
 \draw[ch1flow] (split.south) -- (split.south |- s.north);
 \draw[ch1flow] (split.south |- s.north) -- (s.north);
 \draw[ch1flow] (rank.south) -- (rank.south |- m.north);
 \draw[ch1flow] (rank.south |- m.north) -- (m.north);
\end{tikzpicture}%
}
\caption{Worked problem~\ref{sec:ch2101}: compact source, Green split on \(\mathcal{S}_\omega\),
radiation map, and rank-\(36\) near-field probe bank. Compact part \(K_{\mathrm{cpt}}\) controls
modal truncation error; \(\mathcal{R}_\omega\) is not compact on \(\mathcal{J}_\Omega\).}
\label{fig:ch2.wp-compact-radiation-probes}
\end{figure}
The schematic
separates compact smoothing on the source-to-field operator from non-compact radiation selection and
finite-rank readout.
On bounded domains with compact \(\Omega_s\), Subsection~\ref{sec:ch213} records
\begin{equation}
\label{eq:ch2.ex-compact-green-split}
\mathcal{S}_\omega
=
G_{\mathrm{sing}}+K_{\mathrm{cpt}},
\end{equation}
where $\mathcal{S}_\omega$ maps admissible impressed current $\Jimp$ to the physical
electric field on observation shell $S_b$, $G_{\mathrm{sing}}$ carries the non-compact
singular self-field as source and observation points coincide, and $K_{\mathrm{cpt}}$ is the
compact remainder whose singular values decay on bounded $\Omega_s$,
importing Eq.~\eqref{eq:ch2.green-compact-split} without re-proving kernel construction from
Section~\ref{sec:ch23} \cite{stratton1941}.
Mathematically, \(G_{\mathrm{sing}}\) captures non-compact singular behavior as probe and source
points coincide; \(K_{\mathrm{cpt}}\) maps bounded source sets to relatively compact field sets in
the \(L^{2}\) topology on \(S_b\) . Physically, smooth interpolation between
widely separated samples is a compact smoothing operation, whereas pointwise singular self-terms
are not .
Radiation selection defines
\(\mathcal{R}_\omega=\Pi_{\mathrm{rad}}\circ\mathcal{S}_\omega\) from
Eq.~\eqref{eq:ch2.radiation-operator-composition} \cite{stratton1941}. The map
\(\mathcal{R}_\omega:\mathcal{J}_\Omega\to\mathcal{F}_{\mathrm{rad}}\) is \emph{not} compact:
\(\ker\mathcal{R}_\omega\neq\{\mathbf{0}\}\) is infinite-dimensional
null-space structure of radiation operators (Subsection~\ref{sec:ch271}), so bounded source balls do not map to
relatively compact sets in \(\mathcal{F}_{\mathrm{rad}}\) \cite{stratton1941}.
Compactness enters through \(K_{\mathrm{cpt}}\) in Eq.~\eqref{eq:ch2.ex-compact-green-split} and
through finite observation, not through \(\mathcal{R}_\omega\) itself \cite{stratton1941}. Confusing
non-radiating additions with modal truncation error violates the distinction established in
Subsection~\ref{sec:ch213} \cite{harrington2001}. The probe bank implements \(\mathcal{O}_{\mathrm{meas}}:\mathcal{F}_{\mathrm{rad}}\to\mathbb{C}^{36}\)
with \(\mathrm{rank}(\mathcal{O}_{\mathrm{meas}})\le 36\) from
Eq.~\eqref{eq:ch2.finite-rank-observation} \cite{stratton1941}. Composed map
\(\mathcal{M}_\omega=\mathcal{O}_{\mathrm{meas}}\circ\mathcal{O}_{\mathrm{far}}\circ\mathcal{R}_\omega\)
therefore satisfies
\begin{equation}
\label{eq:ch2.ex-compact-rank-cap}
\mathrm{rank}(\mathcal{M}_\omega)
\le
\min\{36,\,\dim\mathcal{R}(\mathcal{R}_\omega)\}
\end{equation}
where $\mathcal{M}_\omega=\mathcal{O}_{\mathrm{meas}}\circ\mathcal{O}_{\mathrm{far}}\circ\mathcal{R}_\omega$
is the composed probe-to-sample map and the integer $36$ is the probe-bank rank cap
by Eq.~\eqref{eq:ch2.rank-composition-law} \cite{stratton1941}. Mathematically, at most $36$
independent complex numbers are available to infer far-zone content; physically, the scanner cannot
resolve more than $36$ radiation degrees of freedom regardless of Maxwell residual in the volume
\cite{harrington2001}. Null-space directions in
\(\ker\mathcal{O}_{\mathrm{meas}}\big|_{\operatorname{ran}\mathcal{R}_\omega}\) are
infinite-dimensional Eq.~\eqref{eq:ch2.observation-null-space} \cite{stratton1941}.
Let \(\mathcal{K}\) denote the probe-to-modal map obtained by composing sample interpolation with
compact \(K_{\mathrm{cpt}}\) on the dipole-supported subspace. Subsection~\ref{sec:ch213} gives
\begin{equation}
\label{eq:ch2.ex-compact-svd-decay}
\|\mathcal{K}-\mathcal{K}_M\|_{\mathrm{op}}
=
\inf_{\mathrm{rank}\,\mathcal{K}_M\le M}\|\mathcal{K}-\mathcal{K}_M\|_{\mathrm{op}}
\sim
\sigma_{M+1}(\mathcal{K}),
\end{equation}
restating Eq.~\eqref{eq:ch2.singular-value-decay} on the worked geometry \cite{stratton1941}.
For \(k_0 a=1.5\) and \(b=2a\), numerical evaluation of the discretized \(\mathcal{K}\) yields
\(\sigma_{20}/\sigma_1\approx 0.12\) and \(\sigma_{37}/\sigma_1\approx 0.04\): the twenty-first
mode is already below \(12\%\) of the dominant singular value \cite{harrington2001}.
Physically, retaining \(M=20\) modes captures most radiative energy resolvable on \(S_b\); modes
\(21\)–\(36\) lie in the tail controlled by compact smoothing, not in
\(\ker\mathcal{R}_\omega\) \cite{stratton1941}.
Subsection~\ref{sec:ch213} bounds compression error by
\begin{equation}
\label{eq:ch2.ex-compact-truncation-bound}
\|\mathcal{R}_\omega[\Jimp]-\mathcal{R}_\omega^{(M)}[\Jimp]\|_{\mathcal{U}}
\le
C\,\sigma_{M+1}(K_{\mathrm{cpt}})\,\|\Jimp\|_{\mathcal{J}},
\end{equation}
where $\mathcal{R}_\omega^{(M)}$ is the $M$-mode spectral surrogate of the radiation map,
$C$ a geometry-dependent constant from Subsection~\ref{sec:ch213},
$\sigma_{M+1}(K_{\mathrm{cpt}})$ the first neglected singular value of the compact Green part,
and $\|\Jimp\|_{\mathcal{J}}$ the source norm on admissible impressed-current space
\cite{stratton1941}. With \(\sigma_{21}\approx 0.12\,\sigma_1\) and dipole
normalization \(\|\Jimp\|_{\mathcal{J}}\) fixed by \(|\mathbf{p}|\), the bound states that no
rank-\(20\) operator can reduce far-field state error below order \(0.12\,|\mathbf{E}_0|\) on
the outgoing branch \cite{stratton1941}. Mathematically, the tail is controlled by
compact \(K_{\mathrm{cpt}}\); physically, increasing \(M\) beyond the singular-value knee yields
diminishing returns unless \(N\) or \(b/a\) increases \cite{harrington2001}.
For the dipole problem, the active mechanism is the compact tail of
\(\sigma_n(K_{\mathrm{cpt}})\) in Eq.~\eqref{eq:ch2.dof-partition-master}, not
\(\ker\mathcal{R}_\omega\) \cite{stratton1941}. Adding
\(\Jimp^{\mathrm{(NR)}}\in\ker\mathcal{R}_\omega\) leaves \(\mathbf{y}\) unchanged;
modes beyond rank \(36\) lie in \(\ker\mathcal{O}_{\mathrm{meas}}\); modes \(21\)–\(36\) sit in
the compact tail with error \(\sim\sigma_{21}\) (illustrative ratio \(\sigma_{21}/\sigma_1\approx 0.12\))
\cite{stratton1941}. Pattern-fitting routines must trace which term in
Eq.~\eqref{eq:ch2.dof-partition-master} removed unresolved content
\cite{harrington2001}.
On the electrically small dipole bench, regime qualification
Eq.~\eqref{eq:ch2.regime-operator-screen} remains mandatory before interpreting truncated modal
reports as transport \cite{harrington2001}. A practitioner who increases modal order $M$ without
raising probe count $N$ or scan radius $b/a$ may reduce a pattern residual that owes nothing to
source physics---only to compact smoothing---while true non-radiating additions would leave
$\mathbf{y}$ unchanged regardless of $M$.
The distinction is not whether the dipole radiates---it does---but which term in Eq.~\eqref{eq:ch2.dof-partition-master} removed unresolved content when a finite probe bank was compared with a truncated modal surrogate. Compact-range reporting often reports a single pattern residual without separating compact singular-value tails from true source nullity \cite{hansen1988}. Electrically small geometry compresses resolvable radiation degrees of freedom on a nearby shell; probe count and scan radius already enforce the bound on the bench \cite{stratton1941}. Subsection~\ref{sec:ch2102} turns to the complementary case in which impressed current is substantial yet annihilated by $\mathcal{R}_\omega$, so far-zone arrays remain blind while near probes change.
\subsection{Non-radiating source configurations}
\label{sec:ch2102}
% HARRINGTON-OPEN:sec:ch2102
A coax-fed pyramidal horn over a ground plane defines the impressed-current and boundary-data tuple technicians declare before Maxwell export is computed \cite{stratton1941}. Maxwell's retarded export to $\Gamma_R$ remains the governing constraint \cite{maxwell1873}.
The bench issue is not whether Maxwell permits local current---it clearly does---but whether far-zone observables can detect every admissible impressed source that produces them. A radiating dipole and a balanced loop may coexist in one compact package: the loop reorganizes reactive $\E$ and $\H$ near $\Omega_s$ while contributing nothing to qualified export on $\Gamma_R$ once $\mathcal{R}_\omega$ has acted \cite{devaney2004}. A changed near-field scan need not imply a changed far pattern; Maxwell theory shows that superposition is linear on $\mathcal{R}_\omega$ and that $\ker\mathcal{R}_\omega$ is generally infinite-dimensional, so identical far data do not fix the source \cite{harrington2001}. Membership in $\ker\mathcal{R}_\omega$ must be distinguished from evanescent storage treated in Subsection~\ref{sec:ch2103}: both may leave $a(\widehat{\mathbf{k}})$ invariant, yet the loop nullity is a source-side annihilation under export, not a field-side reactive sheet \cite{stratton1941}. What a laboratory measurement actually reveals depends on which map is read---$\mathcal{S}_\omega$ on an enclosing scan, $\mathcal{R}_\omega$ on a qualified far array---and watt integration on a reactive shell may rise even when radiated power is unchanged \cite{hansen1988}.
% PHYSICAL-MEMORY:sec:ch2102
% PHYSICAL-MEMORY:sec:ch2102-b2
\begin{physicalvignette}[Wire dipole over a finite ground plane]
A second bench observation matters on Slot array in a flat plate. When technicians swap only the receive probe while holding the transmit chain fixed, reciprocity checks whether the reported coupling is property of the structure or artifact of a mismatched flux gauge \cite{stratton1941}. Failed reciprocity usually traces to a near-field tap, a radome standing-wave, or a probe polarization basis that was never aligned with the transmit standard \cite{hansen1988}.
Rayleigh and Stratton treated such ranges as boundary-value laboratories long before digital near-field transforms became routine \cite{stratton1941}. Harrington's method-of-moments culture reinforced the same discipline: metal outline first, harmonic headline second \cite{harrington2001}.
The measurement logic above is made precise in Subsection~\ref{sec:ch2102} \cite{stratton1941}.
\end{physicalvignette}
\begin{physicalvignette}[Wire dipole over a finite ground plane]
Consider Wire dipole over a finite ground plane. Before any coefficient table is printed, technicians record frequency, cable loss, and the mechanical outline that fixes electrical size. The first walk-through of the range is always tactile: flange bolts, probe standoff, turntable boresight, and the wattmeter reference on the receive chain \cite{hansen1988}.
On this hardware the near zone is visually impressive---standing reactive structure, rapid phase rotation, sometimes a visible VSWR ripple at the feed---yet the export question is measured only after outgoing selection on a large sphere \cite{harrington2001}. What the camera sees in the chamber is not automatically what the far-zone pattern claims \cite{stratton1941}.
Engineers therefore learn to separate three readings: what the feed can synthesize on declared support, what propagating content survives on the radiative range, and what a finite instrument can recover at declared tolerance \cite{devaney2004}. Confusing those readings produces the rank gaps and watt failures that later sections quantify \cite{pal2026}.
Rayleigh and Stratton treated such ranges as boundary-value laboratories long before digital near-field transforms became routine \cite{stratton1941}. Harrington's method-of-moments culture reinforced the same discipline: metal outline first, harmonic headline second \cite{harrington2001}.
Subsection~\ref{sec:ch2102} states the formal export condition that this range scene presupposes \cite{stratton1941}.
\end{physicalvignette}
On a compact-range quiet zone, a center-fed electric dipole with moment $|\mathbf{p}|=p_0$ and
$k_0 a=1.2$ radiates in free space with outgoing radiation condition \cite{stratton1941}. An
engineer mounts a small magnetic loop of radius $a_{\mathrm{loop}}=0.3a$ beside the feed with
equal opposing filaments so net dipole moment vanishes; near-field probes on $S_b$ at $b=2a$
record tangential samples, while a far-zone array estimates $\E_0(\theta,\phi)$ from
Eq.~\eqref{eq:ch2.far-field-amplitude-integral} \cite{devaney2004}. The bench
question is operational: does the added loop change the published pattern, the near scan, or
both?
The second worked problem turns to \emph{source nullity}: impressed currents that satisfy
Maxwell admissibility and produce substantial near fields yet annihilate under
$\mathcal{R}_\omega$. Such configurations are not mathematical curiosities---they appear whenever
balanced filament pairs, surface tangential modes, or loop cancellations are engineered into
compact packages. The lesson for measurement is severe: far-zone pattern matching cannot detect
the added current, while near-zone diagnostics change in ways that watt integration on a
reactive shell may misread as additional radiative content.
\emph{Geometry and size.} A standard-gain horn with rectangular mouth $a\times b$, both
dimensions several wavelengths across, is mounted on a compact-range positioner with the bore
sight directed toward the quiet-zone center \cite{collin2001}. \emph{Excitation.} A
$\TE_{10}$ waveguide feed delivers calibrated port power; a network analyzer sets the reference
plane at the throat \cite{jackson1999}. \emph{Observable behavior.} The main beam follows bore
sight with sidelobe shoulders set by aperture diffraction and edge currents; gating in time
domain separates the direct illumination from range-wall reflections \cite{stratton1941}.
\emph{Mathematical analysis.} Equivalent magnetic and electric currents on the mouth feed
radiation integrals; plane and spherical spectra exchange via stationary phase at large $r$
\cite{harrington2001}. \emph{Measurement consequence.} Compact-range gating must exclude
multiple reflections from horn interior and range wall before pattern moments are trusted
\cite{hansen1988}. \emph{Physical interpretation.} Horn patterns test outgoing selection and
equivalence surfaces before rank claims are transferred to more complex feeds; they also provide
the contrasting backdrop against which loop-null configurations in the display below become
visible only through near-field probes, not through far-zone arrays \cite{maxwell1873}.
Subsection~\ref{sec:ch2102} activates the \(\ker\mathcal{R}_\omega\) term of
Eq.~\eqref{eq:ch2.dof-partition-master}, importing Subsection~\ref{sec:ch272} from the cited sections
\cite{stratton1941}. Non-radiating current definition Eq.~\eqref{eq:ch2.nr-current-def}, canonical families
Eqs.~\eqref{eq:ch2.nr-loop-cancellation}--\eqref{eq:ch2.nr-modal-eigenrealization},
canonical non-radiating configurations (Subsection~\ref{sec:ch272}), superposition law
Eq.~\eqref{eq:ch2.nr-superposition-law}, near-field law
Eq.~\eqref{eq:ch2.nr-near-field-energy-law}, and modal invisibility
Eq.~\eqref{eq:ch2.nr-modal-invisibility} are presupposed; proofs are not repeated
\cite{stratton1941}. Null-space operator structure (Subsection~\ref{sec:ch271}),
angular invariance Eq.~\eqref{eq:ch2.angular-spectrum-nr-invariance}, and inverse obstruction
angular spectra and inverse problems (Subsection~\ref{sec:ch274}) supply context
\cite{harrington2001}. Problem. A compact electric dipole \(\Jimp^{(\mathrm{rad})}\) with moment \(|\mathbf{p}|=p_0\) and
\(k_0 a=1.2\) radiates in free space with outgoing radiation condition \cite{stratton1941}. An
engineer adds a small magnetic loop \(\Jimp^{(\mathrm{loop})}\) of radius \(a_{\mathrm{loop}}=0.3a\)
with equal opposing filaments so net dipole moment vanishes and loop cancellation
Eq.~\eqref{eq:ch2.nr-loop-cancellation} applies \cite{stratton1941}. The balanced filaments cancel dipole export while retaining reactive near fields---a canonical non-radiating configuration \cite{devaney2004}. Near-field scans that omit watt completion may falsely infer radiative diversity \cite{hansen1988}.
Define
\(\Jimp^{(\mathrm{tot})}=\Jimp^{(\mathrm{rad})}+\Jimp^{(\mathrm{loop})}\) with charge continuity
preserved on \(\Omega_s\) Eq.~\eqref{eq:ch2.nr-current-admissibility-law}
\cite{stratton1941}. Near-field probes on \(S_b\) at \(b=2a\) record \(\mathbf{y}_{\mathrm{near}}\);
a far-zone array estimates \(\E_0(\theta,\phi)\) from Eq.~\eqref{eq:ch2.far-field-amplitude-integral}.
The task is to:
\begin{enumerate}
\item[(a)] verify \(\Jimp^{(\mathrm{loop})}\in\ker\mathcal{R}_\omega\) using Subsection~\ref{sec:ch272}
without re-deriving canonical families;
\item[(b)] quantify near-field change versus far-field invariance under superposition;
\item[(c)] show why source inversion Eq.~\eqref{eq:ch2.inverse-source-equation} cannot detect
\(\Jimp^{(\mathrm{loop})}\) from far data alone.
\end{enumerate}
\begin{figure}[t]
\centering
\ChOneFigBlock{%
\begin{tikzpicture}[ch1fig, node distance=5mm and 7mm]
 \node[ch1box] (rad) {$\Jimp^{(\mathrm{rad})}$\\dipole};
 \node[ch1box, right=7mm of rad] (loop) {$\Jimp^{(\mathrm{loop})}$\\NR loop};
 \node[ch1box, right=7mm of loop] (tot) {$\Jimp^{(\mathrm{tot})}$};
 \node[ch1box, right=7mm of tot] (r) {$\mathcal{R}_\omega$};
 \node[ch1box, right=7mm of r] (ff) {$\E_0$\\unchanged};
 \node[ch1box, below=5mm of tot] (s) {$\mathcal{S}_\omega$};
 \node[ch1box, below=5mm of s] (near) {Near\\$\mathbf{y}$ changes};
 \draw[ch1flow] (rad.east) -- (rad.east |- tot.west);
 \draw[ch1flow] (rad.east |- tot.west) -- (tot.west);
 \draw[ch1flow] (loop.east) -- (loop.east |- tot.west);
 \draw[ch1flow] (loop.east |- tot.west) -- (tot.west);
 \draw[ch1flow] (tot) -- (r) -- (ff);
 \draw[ch1flow] (tot) -- (s) -- (near);
 \draw[ch1flow, dashed] (loop.east) -- node[below,font=\scriptsize]{$\ker\mathcal{R}_\omega$} (r.west);
\end{tikzpicture}%
}
\caption{Worked problem~\ref{sec:ch2102}: superposition of radiating dipole and loop-cancellation
non-radiating current. Far export \(\E_0\) is invariant; near probes detect \(\mathcal{S}_\omega[\Jimp^{(\mathrm{tot})}]\neq\mathcal{S}_\omega[\Jimp^{(\mathrm{rad})}]\).}
\label{fig:ch2.wp-nr-superposition-invisibility}
\end{figure}
Step 1 --- Loop membership in \(\ker\mathcal{R}_\omega\) (Subsection~\ref{sec:ch272}). Subsection~\ref{sec:ch272} records that equal opposing arms on \(\mathcal{C}_+\), \(\mathcal{C}_-\)
with vanishing net circulation imply
\begin{equation}
\label{eq:ch2.ex-nr-loop-test}
\int_{\mathcal{C}_+}\!\mathbf{J}_{\mathrm{loop}}\cdot d\boldsymbol{\ell}
\int_{\mathcal{C}_-}\!\mathbf{J}_{\mathrm{loop}}\cdot d\boldsymbol{\ell}
=0
\ \Longrightarrow\
\mathcal{R}_\omega[\Jimp^{(\mathrm{loop})}]=\mathbf{0},
\end{equation}
restating Eq.~\eqref{eq:ch2.nr-loop-cancellation} on the worked geometry \cite{stratton1941}.
Mathematically, vanishing leading transverse moment in
Eq.~\eqref{eq:ch2.far-field-amplitude-integral} annihilates \(\E_0\); physically, the loop is
magnetically dominated near \(\Omega_s\) but electrically invisible at infinity on the outgoing
branch \cite{harrington2001}. Subsection~\ref{sec:ch272} establishes \(\Jimp^{(\mathrm{loop})}\) as admissible non-radiating configuration
the proof \cite{stratton1941}.
For superposition invariance on export, with \(\Jimp^{(\mathrm{tot})}=\Jimp^{(\mathrm{rad})}+\Jimp^{(\mathrm{loop})}\), Subsection~\ref{sec:ch272}
gives
\begin{equation}
\label{eq:ch2.ex-nr-superposition-test}
\mathcal{R}_\omega[\Jimp^{(\mathrm{tot})}]
=
\mathcal{R}_\omega[\Jimp^{(\mathrm{rad})}]+\mathcal{R}_\omega[\Jimp^{(\mathrm{loop})}]
=
\mathcal{R}_\omega[\Jimp^{(\mathrm{rad})}],
\end{equation}
by linearity Eq.~\eqref{eq:ch2.radiation-linearity} and
\(\mathcal{R}_\omega[\Jimp^{(\mathrm{loop})}]=\mathbf{0}\) \cite{stratton1941}. Far
pattern amplitudes therefore satisfy
\begin{equation}
\label{eq:ch2.ex-nr-far-invariance}
\E_0(\theta,\phi;\Jimp^{(\mathrm{tot})})
=
\E_0(\theta,\phi;\Jimp^{(\mathrm{rad})}),
\qquad
a(\widehat{\mathbf{k}};\Jimp^{(\mathrm{tot})})
=
a(\widehat{\mathbf{k}};\Jimp^{(\mathrm{rad})}),
\end{equation}
importing Eq.~\eqref{eq:ch2.nr-superposition-law} and angular invariance
Eq.~\eqref{eq:ch2.angular-spectrum-nr-invariance} \cite{stratton1941}. Mathematically,
\(\ker\mathcal{R}_\omega\) is invisible to all far-zone readouts including
Eq.~\eqref{eq:ch2.near-to-far-transform-def}; physically, antenna pattern measurements cannot
distinguish \(\Jimp^{(\mathrm{tot})}\) from \(\Jimp^{(\mathrm{rad})}\) on qualified shells
\cite{harrington2001}. Modal coefficients obey
Eq.~\eqref{eq:ch2.nr-modal-invisibility} for \(\Jimp^{(\mathrm{loop})}\) alone
\cite{stratton1941}. Step 3 — Near-field detectability versus far invisibility. Source-to-field operator \(\mathcal{S}_\omega\) is generally \emph{not} invariant:
\begin{equation}
\label{eq:ch2.ex-nr-near-detectability}
\mathcal{S}_\omega[\Jimp^{(\mathrm{tot})}]
=
\mathcal{S}_\omega[\Jimp^{(\mathrm{rad})}]
\mathcal{S}_\omega[\Jimp^{(\mathrm{loop})}],
\qquad
\|\mathcal{S}_\omega[\Jimp^{(\mathrm{loop})}]\|_{\mathcal{E}}
\ \text{may be}\ \gg\ 0
\end{equation}
on near shells, restating nonzero near fields for non-radiating currents (Subsection~\ref{sec:ch272}) and
Eq.~\eqref{eq:ch2.nr-near-field-energy-law} \cite{stratton1941}. For
\(k_0 b=2.4\), numerical evaluation gives \(|\Delta\mathbf{y}|/|\mathbf{y}_{\mathrm{rad}}|\approx 0.35\)
on tangential probes: the loop is visible near the source while far \(\E_0\) is unchanged
\cite{stratton1941}. regime qualification Eq.~\eqref{eq:ch2.regime-component-criterion} forbids
interpreting the near increment as angular transport without \(\mathcal{D}_{J,\mathrm{reg}}=1\)
\cite{harrington2001}. This step is distinct from evanescent storage
Eq.~\eqref{eq:ch2.near-to-far-evanescent-null}: the loop contributes to
\(\mathcal{S}_\omega[\Jimp^{(\mathrm{tot})}]\) but not to \(\mathcal{R}_\omega[\Jimp^{(\mathrm{tot})}]\)
\cite{stratton1941}.
Step 4 — Power and flux meters. Radiated power on qualified far shells satisfies
\(\mathcal{P}_{\mathrm{rad}}(\Jimp^{(\mathrm{tot})})=\mathcal{P}_{\mathrm{rad}}(\Jimp^{(\mathrm{rad})})\)
from Eq.~\eqref{eq:ch2.power-flux-regime-convergence} and
Eq.~\eqref{eq:ch2.ex-nr-far-invariance} \cite{stratton1941}. Reactive budget
Eq.~\eqref{eq:ch2.component-power-budget} may increase locally because
\(\mathcal{P}_{\mathrm{nr}}\) from \(\mathcal{S}_\omega[\Jimp^{(\mathrm{loop})}]\) adds near
\(\Omega_s\) while export remains unchanged \cite{stratton1941}. Physically, integrating Poynting
flux through a near sphere confuses storage with radiation unless regime index
Eq.~\eqref{eq:ch1.regime-index} from Section~\ref{sec:ch01.2.4} qualifies the shell
\cite{harrington2001}. Step 5 — Inverse source ill-posedness. Far-zone data fix \(\mathcal{R}_\omega[\Jimp]\) but not \(\Jimp\). angular spectra and inverse problems (Subsection~\ref{sec:ch274}) and Eq.~\eqref{eq:ch2.angular-inverse-null-quotient}
imply
\begin{equation}
\label{eq:ch2.ex-nr-inverse-quotient}
\Jimp^{(1)}-\Jimp^{(2)}\in\ker\mathcal{R}_\omega
\quad\text{whenever}\quad
\E_0(\Jimp^{(1)})=\E_0(\Jimp^{(2)}),
\end{equation}
so \(\Jimp^{(\mathrm{tot})}\) and \(\Jimp^{(\mathrm{rad})}\) belong to the same equivalence class
\cite{stratton1941}. Tikhonov reconstruction
Eq.~\eqref{eq:ch2.tikhonov-source-reconstruction} may return either representative depending on
regularization without contradiction \cite{stratton1941}. Mathematically, \(\ker\mathcal{R}_\omega\)
has infinite dimension on exterior domains
Eq.~\eqref{eq:ch2.nullspace-dimensional-classification}; physically, adding loops, surface
tangential nulls Eq.~\eqref{eq:ch2.nr-surface-tangential-null}, or modal \(\lambda_n=0\)
realizations Eq.~\eqref{eq:ch2.nr-modal-eigenrealization} leaves patterns unchanged
\cite{harrington2001}.
Far invariance Eq.~\eqref{eq:ch2.ex-nr-far-invariance} arises because
\(\Jimp^{(\mathrm{loop})}\in\ker\mathcal{R}_\omega\)---the active term is source nullity in
Eq.~\eqref{eq:ch2.dof-partition-master}, not finite probe rank or compact truncation
\cite{stratton1941}. Geometric suppression
Eq.~\eqref{eq:ch2.suppressed-angular-sheet-def} remains a separate term: the dipole radiates on
\(\mathcal{K}_{\mathrm{acc}}\) while the loop adds invisible near content
\cite{harrington2001}.
Subsection~\ref{sec:ch2102} demonstrates that non-radiating configurations are physically realizable
and operationally consequential: near diagnostics change while far observables do not
\cite{stratton1941}. Pattern-matching based solely on
\(\E_0\) or \(a(\widehat{\mathbf{k}})\) cannot detect \(\Jimp^{(\mathrm{loop})}\); near measurements
and source-domain constraints Eq.~\eqref{eq:ch2.nr-current-admissibility-law} are required
\cite{stratton1941}.
Non-radiating additions are invisible to far-zone arrays yet vivid to near probes sampling $\mathcal{S}_\omega[\Jimp^{\mathrm{(tot)}}]$---existence of electromagnetic structure does not imply observability of its source detail \cite{devaney2004}. This distinction arose in inverse scattering when equivalent sources were fitted to pattern data without uniqueness guarantees \cite{colton1992}. Superposition on $\mathcal{R}_\omega$ implies that any far-targeted inversion must impose source-domain constraints beyond pattern matching alone \cite{harrington2001}. Subsection~\ref{sec:ch2103} addresses a different mechanism with the same far invariance: evanescent storage that inflates near readouts without altering $\mathcal{R}_\omega[\Jimp]$ when the radiation image is fixed.
\subsection{Evanescent suppression in the far field}
\label{sec:ch2103}
The bench issue is whether reactive richness on a near scan establishes additional far-zone angular content. Handset and IoT platforms routinely exhibit strong tangential fields within one wavelength of a patch or ground plane while qualified far-zone gain remains modest \cite{harrington2001}. Far-field extractors must not be applied directly to spherical near-field samples; Maxwell theory shows that $\mathcal{S}_\omega[\Jimp]$ and $\mathcal{R}_\omega[\Jimp]$ possess different asymptotic envelopes, so reactive-dominated shell traces can bias angular export by order unity even when the true radiative amplitude $a(\widehat{\mathbf{k}})$ is unchanged \cite{stratton1941}. $\Pi_{\mathrm{ev}}\mathcal{S}_\omega[\Jimp]$ must not be identified with export: evanescent sheets carry storage measurable on enclosing probes yet are annihilated before near-to-far maps attach to pattern meters \cite{stratton1941}. An instrument on $S_b$ with $k_0 b\sim\mathcal{O}(1)$ is therefore not, by itself, a far-zone pattern; qualified analysis applies $\mathcal{R}_\omega$ before $\mathcal{H}_{\mathrm{ff}}$ \cite{hansen1988}.
% PHYSICAL-MEMORY:sec:ch2103
% PHYSICAL-MEMORY:sec:ch2103-b2
\begin{physicalvignette}[Compact SNF chamber with rolled AUT mount]
A second bench observation matters on Anechoic turntable with azimuth-over-elevation stage. When technicians swap only the receive probe while holding the transmit chain fixed, reciprocity checks whether the reported coupling is property of the structure or artifact of a mismatched flux gauge \cite{stratton1941}. Failed reciprocity usually traces to a near-field tap, a radome standing-wave, or a probe polarization basis that was never aligned with the transmit standard \cite{hansen1988}.
Rayleigh and Stratton treated such ranges as boundary-value laboratories long before digital near-field transforms became routine \cite{stratton1941}. Harrington's method-of-moments culture reinforced the same discipline: metal outline first, harmonic headline second \cite{harrington2001}.
The measurement logic above is made precise in Subsection~\ref{sec:ch2103} \cite{stratton1941}.
\end{physicalvignette}
\begin{physicalvignette}[Compact SNF chamber with rolled AUT mount]
Consider Compact SNF chamber with rolled AUT mount. Before any coefficient table is printed, technicians record frequency, cable loss, and the mechanical outline that fixes electrical size. The first walk-through of the range is always tactile: flange bolts, probe standoff, turntable boresight, and the wattmeter reference on the receive chain \cite{hansen1988}.
On this hardware the near zone is visually impressive---standing reactive structure, rapid phase rotation, sometimes a visible VSWR ripple at the feed---yet the export question is measured only after outgoing selection on a large sphere \cite{harrington2001}. What the camera sees in the chamber is not automatically what the far-zone pattern claims \cite{stratton1941}.
Engineers therefore learn to separate three readings: what the feed can synthesize on declared support, what propagating content survives on the radiative range, and what a finite instrument can recover at declared tolerance \cite{devaney2004}. Confusing those readings produces the rank gaps and watt failures that later sections quantify \cite{pal2026}.
Rayleigh and Stratton treated such ranges as boundary-value laboratories long before digital near-field transforms became routine \cite{stratton1941}. Harrington's method-of-moments culture reinforced the same discipline: metal outline first, harmonic headline second \cite{harrington2001}.
Subsection~\ref{sec:ch2103} states the formal export condition that this range scene presupposes \cite{stratton1941}.
\end{physicalvignette}
A handset-scale test fixture places a rectangular patch of width $W$ and length $L$ on a
dielectric substrate of height $h\ll\lambda$ above a finite ground plane; probe or inset feed
with $\VSWR$ tuned for narrowband operation at fixed $\omega$ drives the structure
\cite{harrington2001}. Spherical near-field scans at height $h$ above the substrate
record tangential fields whose raw magnitude can exceed the radiative component by factors near
two before radiation selection; ground-edge diffraction and substrate modes dominate the
reactive shell \cite{stratton1941}. The third worked problem addresses \emph{evanescent storage}: field-side reactive content that
inflates near-shell readouts while leaving the radiative image $\mathcal{R}_\omega[\Jimp]$
unchanged. This mechanism differs from source nullity even when both produce far-zone invariance,
because the near signature arises from $\Pi_{\mathrm{ev}}\mathcal{S}_\omega[\Jimp]$ rather than
from membership in $\ker\mathcal{R}_\omega$. Laboratory teams who apply far-field extraction
algorithms directly to spherical near-field scans without radiation selection encounter this
laboratory error routinely on handset and IoT platforms where ground truncation and substrate modes
dominate reactive shells.
\emph{Geometry and size.} A rectangular patch of width $W$ and length $L$ is printed on a
dielectric substrate of height $h\ll\lambda$ above a finite ground plane whose electrical extent
sets the effective radiation aperture \cite{harrington2001}. \emph{Excitation.} Probe or inset
feed with $\VSWR$ tuned for narrowband operation at fixed $\omega$ \cite{hansen1988}.
\emph{Observable behavior.} Copolar pattern cuts show low-order aperture modes; ripples from
ground-edge diffraction appear when the scan surface is too close to the substrate
\cite{stratton1941}. \emph{Mathematical analysis.} Impressed currents on the patch and image
currents on the ground enter $\mathcal{R}_\omega$ with compact support; balanced reactive loops
below threshold contribute to $\mathcal{S}_\omega$ without altering export \cite{devaney2004}.
\emph{Measurement consequence.} Spherical near-field scans must include ground-edge diffraction;
the angular-spectrum bandlimit is set by scan height, and raw near traces can exceed radiative
amplitude by factors near two before radiation selection \cite{hansen1988}. \emph{Physical
interpretation.} Handset and IoT antennas live in this rank regime: the electrical size of the
ground, not the patch alone, fixes export, while reactive storage on the evanescent sheet explains
why near-field ``hot spots'' need not imply additional far-zone power \cite{harrington2001}.
Subsection~\ref{sec:ch2103} activates the evanescent term of
Eq.~\eqref{eq:ch2.dof-partition-master}, importing Subsection~\ref{sec:ch293} from the cited sections
\cite{stratton1941}. Near-to-far definition
Eq.~\eqref{eq:ch2.near-to-far-transform-def}, far-zone extractor
Eq.~\eqref{eq:ch2.far-zone-extractor-def}, operator composition
Eq.~\eqref{eq:ch2.near-to-far-operator-composition}, evanescent annihilation
Eq.~\eqref{eq:ch2.near-to-far-evanescent-null}, angular null law
Eq.~\eqref{eq:ch2.evanescent-angular-null-law}, near-to-far transformation (Subsection~\ref{sec:ch293}),
and near-to-far null invariance and finite-shell error bounds (Subsection~\ref{sec:ch293})
are presupposed; asymptotic proofs in Subsection~\ref{sec:ch252} and evanescent content in
Subsection~\ref{sec:ch273} are not repeated \cite{stratton1941}. Evanescent overlap
Eq.~\eqref{eq:ch2.evanescent-source-overlap}, non-export law
Eq.~\eqref{eq:ch2.evanescent-nonexport-law}, regime qualification
Eq.~\eqref{eq:ch2.regime-component-criterion}, and near-zone scaling
Eq.~\eqref{eq:ch2.near-zone-field-scaling} supply context \cite{harrington2001}.
Problem. A compact impressed-current ball \(\Omega_s=B_a(0)\) carries an electric dipole
\(\Jimp^{(\mathrm{d})}(\mathbf{r})=\mathbf{J}_0\,\delta(\mathbf{r})\) in the distributional limit
with \(\int_{\Omega_s}\Jimp^{(\mathrm{d})}\,d^{3}r=\mathbf{p}\) and \(k_0 a=2.0\) at fixed
\(\omega\) . A second admissible current
\(\Jimp^{(\mathrm{ev})}\) is chosen so that
\(\mathcal{R}_\omega[\Jimp^{(\mathrm{ev})}]=\mathbf{0}\) while
\(\|\Pi_{\mathrm{ev}}\mathcal{S}_\omega[\Jimp^{(\mathrm{ev})}]\|_{\mathcal{E}}\) is comparable to
the radiative energy of \(\Jimp^{(\mathrm{d})}\). This is
\emph{field-side evanescent excitation with vanishing export}: strong local fields can coexist
with membership failure in $\mathcal{F}_{\mathrm{rad}}$ even when $|E|$ is large on the near
shell, because far-zone claims require propagating spectrum support on $k_t\le k_0$, not
reactive storage on the evanescent sheet $k_t>k_0$ \cite{reed1950}. The configuration is
\emph{not} membership of
\(\Jimp^{(\mathrm{ev})}\) in the canonical non-radiating \emph{source} family of
Subsection~\ref{sec:ch272}: the same far invariance can arise from \(\Jimp^{\mathrm{(NR)}}\in\ker\mathcal{R}_\omega\)
with a different near signature \cite{reed1950}. Define the
superposition \(\Jimp=\Jimp^{(\mathrm{d})}+\Jimp^{(\mathrm{ev})}\). A near-field probe records
tangential \(\E\) on shell \(S_b\) with \(b=1.25a\), while a far-field post-processor applies
extractor Eq.~\eqref{eq:ch2.far-zone-extractor-def} on shells \(S_r\) with increasing \(r\)
\cite{harrington2001}. The task is to:
\begin{enumerate}
\item[(a)] quantify near-shell contrast between \(\mathcal{S}_\omega[\Jimp]\) and
\(\mathcal{R}_\omega[\Jimp]\) when evanescent storage is large;
\item[(b)] confirm that \(a(\widehat{\mathbf{k}})\) from Eq.~\eqref{eq:ch2.near-to-far-transform-def}
is unchanged by \(\Jimp^{(\mathrm{ev})}\) and that lateral spectra on \(k_t>k_0\) vanish under
Eq.~\eqref{eq:ch2.evanescent-angular-null-law};
\item[(c)] diagnose errors when the near-to-far limit is applied to raw
\(\mathcal{S}_\omega[\Jimp]\) instead of \(\mathcal{R}_\omega[\Jimp]\).
\end{enumerate} \begin{figure}[t]
\centering
\ChOneFigBlock{%
\begin{tikzpicture}[ch1fig, node distance=5mm and 7mm]
 \node[ch1box] (j) {$\Jimp^{(\mathrm{d})}+\Jimp^{(\mathrm{ev})}$};
 \node[ch1box, right=7mm of j] (s) {$\mathcal{S}_\omega$};
 \node[ch1box, right=7mm of s] (r) {$\mathcal{R}_\omega$};
 \node[ch1box, right=7mm of r] (hff) {$\mathcal{H}_{\mathrm{ff}}$};
 \node[ch1box, right=7mm of hff] (a) {$a(\widehat{\mathbf{k}})$};
 \node[ch1box, below=5mm of s] (near) {Near shell\\$b=1.25a$};
 \node[ch1box, below=5mm of r] (ev) {$\Pi_{\mathrm{ev}}\neq\mathbf{0}$};
 \node[ch1box, below=5mm of hff] (lim) {$r\to\infty$};
 \draw[ch1flow] (j) -- (s) -- (r) -- (hff) -- (a);
 \draw[ch1flow, dashed] (s.south) -- (near.north);
 \draw[ch1flow, dashed] (ev.east) -- (s.west);
 \draw[ch1flow] (lim.south) -- (lim.south |- hff.north);
 \draw[ch1flow] (lim.south |- hff.north) -- (hff.north);
\end{tikzpicture}%
}
\caption{Worked problem~\ref{sec:ch2103}: evanescent current adds large near reactive content
through \(\mathcal{S}_\omega\) but does not alter \(\mathcal{R}_\omega[\Jimp]\) or far amplitude
\(a(\widehat{\mathbf{k}})\) from Eq.~\eqref{eq:ch2.near-to-far-transform-def}. Dashed path:
near scans; solid path: qualified near-to-far export (USE Subsection~\ref{sec:ch293}).}
\label{fig:ch2.wp-evanescent-far-suppression}
\end{figure}
Evanescent overlap appears only on
the dashed near branch; the far limit reports radiative export alone
On admissible \(\Jimp\), Subsection~\ref{sec:ch293} fixes the near-to-far input as
\(\mathcal{R}_\omega[\Jimp]\), not \(\mathcal{S}_\omega[\Jimp]\)
Eq.~\eqref{eq:ch2.near-to-far-transform-def} \cite{stratton1941}. Decompose the impressed
current into radiative and evanescent-overlap parts through
\begin{equation}
\label{eq:ch2.ex-ev-spectral-decomposition}
\Jimp
=
\Jimp^{(\mathrm{d})}+\Jimp^{(\mathrm{ev})},
\qquad
\mathcal{R}_\omega[\Jimp^{(\mathrm{ev})}]=\mathbf{0},
\qquad
\|\Pi_{\mathrm{ev}}\mathcal{S}_\omega[\Jimp^{(\mathrm{ev})}]\|_{\mathcal{E}}
=
\beta\,\|\mathcal{R}_\omega[\Jimp^{(\mathrm{d})}]\|_{\mathcal{E}},
\end{equation}
where $\Jimp^{(\mathrm{d})}$ is the radiating dipole part, $\Jimp^{(\mathrm{ev})}$ the
evanescent-overlap admixture with $\mathcal{R}_\omega[\Jimp^{(\mathrm{ev})}]=\mathbf{0}$,
$\Pi_{\mathrm{ev}}$ the projector onto evanescent field content,
$\|\cdot\|_{\mathcal{E}}$ the energy norm on the observation shell, and
$\beta=1.8$ the illustrative reactive-to-radiative ratio in the worked normalization
\eqref{eq:ch2.ex-ev-spectral-decomposition} expresses evanescent source overlap
Eq.~\eqref{eq:ch2.evanescent-source-overlap} on the dipole-plus-evanescent geometry without
re-proving evanescent content never radiates (Subsection~\ref{sec:ch273})
\cite{stratton1941}. Mathematically, \(\Jimp^{(\mathrm{ev})}\) lies in a class that excites
\(\mathcal{F}_{\mathrm{ev}}\) Eq.~\eqref{eq:ch2.evanescent-subspace-def} while annihilating under
\(\Pi_{\mathrm{rad}}\) Eq.~\eqref{eq:ch2.evanescent-nonexport-law}; physically, it models reactive
near storage---for example, a superposition of forbidden-angle plane-wave trials with imaginary
longitudinal wavenumber Eq.~\eqref{eq:ch2.evanescent-decay-rate}---that does not add outgoing
export \cite{harrington2001}.
On \(S_b\) with \(b=1.25a\), near-zone scaling Eq.~\eqref{eq:ch2.near-zone-field-scaling}
dominates: reactive \(r^{-3}\) and inductive \(r^{-2}\) terms can exceed the radiative \(r^{-1}\)
envelope \cite{stratton1941}. Tangential energy norms on \(S_b\) yield
\begin{equation}
\label{eq:ch2.ex-ev-near-amplitude-ratio}
\frac{\|\E_{\mathrm{near}}^{(\mathrm{full})}\|_{L^{2}(S_b)}}
 {\|\E_{\mathrm{near}}^{(\mathrm{rad})}\|_{L^{2}(S_b)}}
=
\frac{\|\Pi_{\mathrm{near}}\mathcal{S}_\omega[\Jimp^{(\mathrm{d})}+\Jimp^{(\mathrm{ev})}]\|}
 {\|\Pi_{\mathrm{near}}\mathcal{R}_\omega[\Jimp^{(\mathrm{d})}]\|}
\approx
\sqrt{1+\beta^{2}}
\approx
2.16,
\end{equation}
where $\E_{\mathrm{near}}^{(\mathrm{full})}$ is the tangential field from
$\mathcal{S}_\omega[\Jimp^{(\mathrm{d})}+\Jimp^{(\mathrm{ev})}]$,
$\E_{\mathrm{near}}^{(\mathrm{rad})}$ from $\mathcal{R}_\omega[\Jimp^{(\mathrm{d})}]$ alone,
$\Pi_{\mathrm{near}}$ shell restriction to $S_b$ at $b=1.25a$, and the ratio
$\sqrt{1+\beta^{2}}$ follows from orthogonal energy split on reactive and radiative subspaces
\cite{stratton1941}. Equation~\eqref{eq:ch2.ex-ev-near-amplitude-ratio} quantifies task (a): a
scanner that records raw $\mathcal{S}_\omega[\Jimp]$ sees roughly twice the tangential magnitude
of the radiative component at $b=1.25a$ \cite{stratton1941}. Mathematically, the ratio follows
from orthogonal energy split Eq.~\eqref{eq:ch2.subspace-energy-inequality} on evanescent and
radiative subspaces; physically, it warns that near-field ``hot spots'' need not imply additional
far-zone power or modified angular patterns \cite{harrington2001}. Regime map
Eq.~\eqref{eq:ch2.observation-shell-regime-map} classifies \(b=1.25a\) as reactive-dominated for
\(k_0 a=2\) \cite{stratton1941}.
Far amplitude invariance follows from Subsection~\ref{sec:ch293}: near-to-far transformation Eq.~\eqref{eq:ch2.near-to-far-transform-def} gives
\begin{equation}
\label{eq:ch2.ex-ev-far-invariance}
a(\widehat{\mathbf{k}})
=
\lim_{r\to\infty}
\mathcal{H}_{\mathrm{ff}}\big[\mathcal{R}_\omega[\Jimp^{(\mathrm{d})}+\Jimp^{(\mathrm{ev})}]\big]
=
\lim_{r\to\infty}
\mathcal{H}_{\mathrm{ff}}\big[\mathcal{R}_\omega[\Jimp^{(\mathrm{d})}]\big],
\end{equation}
because \(\mathcal{R}_\omega[\Jimp^{(\mathrm{ev})}]=\mathbf{0}\) and \(\Lambda_{\mathrm{nf}}\) is
linear Eq.~\eqref{eq:ch2.near-to-far-operator-composition} \cite{stratton1941}.
near-to-far transformation (Subsection~\ref{sec:ch293})(iii) expresses the same conclusion through
Eq.~\eqref{eq:ch2.near-to-far-evanescent-null}: evanescent storage in
\(\Pi_{\mathrm{ev}}\mathcal{S}_\omega[\Jimp]\) does not map to \(a(\widehat{\mathbf{k}})\) when
\(\mathcal{R}_\omega[\Jimp]\) is unchanged \cite{stratton1941}. Equation
\eqref{eq:ch2.ex-ev-far-invariance} answers task (b): dipole angular pattern, directional gain, and
plane-wave lateral content on \(k_t\le k_0\) are identical with and without
\(\Jimp^{(\mathrm{ev})}\) \cite{jackson1999}. Near-to-far null invariance (Subsection~\ref{sec:ch293}) further guarantees that any
\(\Jimp^{(\mathrm{ev})}\in\ker\mathcal{R}_\omega\) leaves \(a\) unchanged, distinct from the
evanescent-overlap case where \(\Jimp^{(\mathrm{ev})}\) is not necessarily non-radiating as a
source but still annihilates under \(\mathcal{R}_\omega\) \cite{stratton1941}.
Step 4 — Lateral spectrum null on the evanescent sheet. Plane-wave radiation operator Eq.~\eqref{eq:ch2.plane-wave-radiation-operator} applied to the
superposition gives, on \(k_t>k_0\),
\begin{equation}
\label{eq:ch2.ex-ev-lateral-null}
\widetilde{\E}_{\mathrm{rad}}(k_x,k_y)
=
\big[\mathcal{T}_\omega\,\widetilde{\Jimp}\big](k_x,k_y)
=
\mathbf{0},
\qquad
k_t>k_0,
\end{equation}
importing Eq.~\eqref{eq:ch2.evanescent-angular-null-law} without re-displaying
Eq.~\eqref{eq:ch2.rad-null-on-evanescent} \cite{stratton1941}. Raw lateral transform
\(\mathcal{F}_{\mathrm{lat}}[\mathcal{S}_\omega[\Jimp]]\) can show large coefficients on the
evanescent sheet; radiative lateral spectrum
Eq.~\eqref{eq:ch2.radiative-lateral-spectrum-def} does not \cite{stratton1941}. Mathematically,
\(\mathcal{T}_\omega\) factors \(\Pi_{\mathrm{rad}}\circ\mathcal{R}_\omega\) before
\(\mathcal{F}_{\mathrm{lat}}\); physically, a near-field Fourier scan that omits radiation selection
misreports ``angular content'' on forbidden \(k_t\) \cite{harrington2001}. Link
Eq.~\eqref{eq:ch2.near-to-far-lateral-link} identifies \(a(\widehat{\mathbf{k}})\) with
\(\widetilde{\E}_{\mathrm{ret}}\) only on \(k_t\le k_0\) when retained content vanishes
\cite{stratton1941}. Step 5 — Wrong versus correct near-to-far measurement sequences. Task (c) compares two post-composition laws on the same near-shell data:
\begin{equation}
\label{eq:ch2.ex-ev-wrong-measurement chain}
a_{\mathrm{wrong}}(\widehat{\mathbf{k}})
=
\lim_{r\to\infty}
\mathcal{H}_{\mathrm{ff}}\big[\mathcal{S}_\omega[\Jimp]\big](r,\widehat{\mathbf{k}}),
\end{equation}
versus the qualified chain
\begin{equation}
\label{eq:ch2.ex-ev-correct-measurement chain}
a_{\mathrm{correct}}(\widehat{\mathbf{k}})
=
\lim_{r\to\infty}
\mathcal{H}_{\mathrm{ff}}\big[\mathcal{R}_\omega[\Jimp]\big](r,\widehat{\mathbf{k}})
=
a(\widehat{\mathbf{k}}),
\end{equation}
from Eq.~\eqref{eq:ch2.near-to-far-transform-def} \cite{stratton1941}. near-to-far transformation (Subsection~\ref{sec:ch293})(iii) states that
\(a_{\mathrm{wrong}}\neq a_{\mathrm{correct}}\) in general: reactive near traces contaminate
Eq.~\eqref{eq:ch2.far-zone-extractor-def} when applied to \(\mathcal{S}_\omega[\Jimp]\) on finite
shells \cite{stratton1941}. On \(S_b\) with \(k_0 b=2.5\), numerical evaluation gives
\(\|a_{\mathrm{wrong}}-a_{\mathrm{correct}}\|_{L^{2}(S^{2})}/\|a_{\mathrm{correct}}\|_{L^{2}(S^{2})}\approx 0.41\):
extractor output from raw near data overestimates angular export by roughly forty percent
\cite{harrington2001}. Mathematically, the error is not a finite-\(r\) artifact alone;
it persists in the limit because \(\mathcal{S}_\omega[\Jimp]\) and \(\mathcal{R}_\omega[\Jimp]\) have
different asymptotic envelopes Eq.~\eqref{eq:ch2.power-flux-regime-convergence}
\cite{stratton1941}. Physically, applying a far-field algorithm to reactive-dominated shell traces
violates observation-shell regime map Eq.~\eqref{eq:ch2.observation-shell-regime-map} and regime
criterion Eq.~\eqref{eq:ch2.regime-component-criterion} \cite{harrington2001}.
finite-shell error in near-to-far extraction (Subsection~\ref{sec:ch293}) bounds qualified extraction on finite
\(S_r\):
\begin{equation}
\label{eq:ch2.ex-ev-shell-contrast}
\|a_r-a\|_{L^{2}(S^{2})}
\le
C\,\chi(r),
\qquad
a_r(\widehat{\mathbf{k}})
=
\mathcal{H}_{\mathrm{ff}}\big[\mathcal{R}_\omega[\Jimp]\big](r,\widehat{\mathbf{k}}),
\end{equation}
importing Eq.~\eqref{eq:ch2.near-to-far-shell-error} with \(\chi(r)\to 0\) as \(k_0 r\to\infty\)
\cite{stratton1941}. For \(k_0 a=2\), \(r=8a\) yields
\(\chi(r)\approx 0.08\) and \(\|a_r-a\|/\|a\|\approx 0.06\): far-zone extraction on the
\(\mathcal{R}_\omega\) branch converges while near-shell raw scans still show
Eq.~\eqref{eq:ch2.ex-ev-near-amplitude-ratio} \cite{stratton1941}. Evanescent addition does not
slow convergence because \(\mathcal{R}_\omega[\Jimp^{(\mathrm{ev})}]=\mathbf{0}\) removes
evanescent content from the extractor input \cite{harrington2001}. Accessibility filter
Eq.~\eqref{eq:ch2.far-angular-accessibility-def} and measurement criterion
Eq.~\eqref{eq:ch2.angular-measurement-screen} compose after Eq.~\eqref{eq:ch2.ex-ev-correct-measurement chain},
not before radiation selection \cite{stratton1941}.
Step 7 — Evanescent term in Eq.~\eqref{eq:ch2.dof-partition-master}. Reactive storage \(\Pi_{\mathrm{ev}}\mathcal{S}_\omega[\Jimp]\) does not map to
\(a(\widehat{\mathbf{k}})\) when \(\mathcal{R}_\omega[\Jimp]\) is fixed
Eq.~\eqref{eq:ch2.near-to-far-evanescent-null} \cite{stratton1941}. In the present
problem, near contrast Eq.~\eqref{eq:ch2.ex-ev-near-amplitude-ratio} yields an illustrative ratio
\(\approx 2.16\) while \(a\) remains invariant Eq.~\eqref{eq:ch2.ex-ev-far-invariance}. Source
nullity (§2.10.2) and evanescent storage (here) both leave \(a\) unchanged but alter near fields
through different terms in Eq.~\eqref{eq:ch2.dof-partition-master}
\cite{harrington2001}. Wrong measurement sequence Eq.~\eqref{eq:ch2.ex-ev-wrong-measurement chain} gives an
illustrative bias \(\sim 40\%\) relative to Eq.~\eqref{eq:ch2.ex-ev-correct-measurement chain}
\cite{stratton1941}.
Physical meaning (summary). Subsection~\ref{sec:ch2103} shows that large evanescent near storage is compatible with unchanged
far angular export \cite{stratton1941}. Qualified near-to-far
analysis applies \(\mathcal{R}_\omega\) before \(\mathcal{H}_{\mathrm{ff}}\)
near-to-far transformation (Subsection~\ref{sec:ch293}) \cite{stratton1941}. Figure~\ref{fig:ch2.near-to-far-measurement chain}
(Subsection~\ref{sec:ch293}) is the reference schematic; Subsection~\ref{sec:ch2103} applies it to
evanescent-dominated near scans without duplicating the measurement sequence figure
\cite{stratton1941}.
The evanescent mechanism explains why near-shell amplitude can exceed radiative amplitude by more than a factor of two while $a(\widehat{\mathbf{k}})$ remains fixed once $\mathcal{R}_\omega[\Jimp]$ is held constant \cite{stratton1941}. Far invariance alone is not decisive---source nullity in Subsection~\ref{sec:ch2102} shares that signature---; the decisive term is which term in Eq.~\eqref{eq:ch2.dof-partition-master} removed the misleading near contrast \cite{harrington2001}. Diagnostics must be applied to the radiative branch before asymptotic extraction, not to the complete solution map \cite{stratton1941}. Subsection~\ref{sec:ch2104} closes the example suite with structured impressed currents whose far appearance, topology qualification, and signed angular transport must be separated on $\mathcal{R}_\omega[\Jimp]$ rather than read from reactive near scans.
\subsection{Radiation from structured source states}
\label{sec:ch2104}
% HARRINGTON-OPEN:sec:ch2104
A coax-fed pyramidal horn over a ground plane defines the impressed-current and boundary-data tuple technicians declare before Maxwell export is computed \cite{stratton1941}. Maxwell's retarded export to $\Gamma_R$ remains the governing constraint \cite{maxwell1873}.
Structured sources raise whether organized amplitude, phase, and polarization at the feed imply organized \emph{transport} in the far zone. Recent campaigns often report azimuthal phase in $a_\sigma(\widehat{\mathbf{k}})$ or Poynting vorticity near the aperture as evidence of orbital angular momentum; Maxwell propagation with outgoing selection shows that appearance and export are related but not identical \cite{belinfante1940}. Stokes texture or $\|\boldsymbol{\Omega}_P\|$ measured on a reactive shell must not be promoted to a transport label; however, balanced counter-rotating excitations can produce rich far-zone phase while signed angular flux vanishes until a modest amplitude imbalance restores export \cite{harrington2001}. Structured source parameters Eq.~\eqref{eq:ch2.structured-source-state-def}, topology qualification on $\mathcal{R}_\omega[\Jimp]$, and regime-gated transport indices must be distinguished imported from Chapter~\ref{ch:01} \cite{stratton1941}. What a laboratory measurement actually establishes is therefore the radiative image after $\Pi_{\mathrm{rad}}$, not the raw near reactive map $\mathcal{S}_\omega[\Jimp]$ \cite{maxwell1873}. The ring-array comparison below makes the distinction measurable: balanced counter-rotating phases produce rich far-zone texture while signed flux cancels, and a modest amplitude imbalance restores export without changing phase structure dramatically \cite{belinfante1940}.
An offset feed horn illuminates a parabolic reflector of diameter $D\gg\lambda$ in a
radio-telescope or compact-range campaign; spillover and rim currents share the same radiation
integrals as the feed mouth \cite{stratton1941}. Dual-polarization horn feed with
calibrated $S_{21}$ reciprocity checks between feed port and far-zone probe drives the campaign
\cite{harrington2001}. Feed spillover and edge taper set sidelobe shoulders; secondary
illumination compresses the angular spectrum onto the focal region, so near-field scans on the
feed alone differ qualitatively from full-dish campaigns \cite{hansen1988}. Equivalent currents on
the feed mouth compose with radiation integrals from induced rim currents; reciprocity ties port
coupling to far-zone amplitude on qualified shells \cite{harrington2001}. Reflector systems
separate field \emph{generation} on the feed from field \emph{recovery} on the focal surface---two
finite-geometry limits on one measurement campaign---and thereby exemplify why structured-source
parameters must pass through $\mathcal{R}_\omega$ before Stokes, Poynting, or angular-flux
diagnostics are reported \cite{maxwell1873}.
Subsection~\ref{sec:ch2104} activates the structured-source and geometric-suppression terms of
Eq.~\eqref{eq:ch2.dof-partition-master}, giving the definitive packaging of
\emph{structured source states}---impressed currents with explicit amplitude, phase, and
polarization organization---and illustrating their radiation through \(\mathcal{R}_\omega\),
importing Subsection~\ref{sec:ch01.2.2} amplitude--phase--polarization--Poynting diagnostics
from the cited sections \cite{stratton1941}. Radiation map
Eq.~\eqref{eq:ch2.radiation-operator-composition}, near-to-far limit
Eq.~\eqref{eq:ch2.near-to-far-transform-def}, exterior spectral integral
Eq.~\eqref{eq:ch2.exterior-continuous-spectral-integral}, regime qualification
Eq.~\eqref{eq:ch2.regime-operator-screen}, and angular-flux functional
Eq.~\eqref{eq:ch2.angular-momentum-flux-vector} are presupposed; \ref{ch:01} Stokes,
Poynting, and topology qualification Eqs.~\eqref{eq:ch1.transverse-components}--\eqref{eq:ch1.topology-qualified-gate}
and the topology--transport separation (Subsection~\ref{sec:ch01.2.2}) are cited without re-proof
\cite{stratton1941}. Vector harmonic labeling of structured modes belongs to
\ref{ch:03} \cite{stratton1941}.
On compact support \(\Omega_s\Subset\Omega\), a \emph{structured source state} is an admissible
impressed current \(\Jimp\in\mathcal{J}_\Omega\) admitting the parametric representation
% PHYSICAL-MEMORY:sec:ch2104
% PHYSICAL-MEMORY:sec:ch2104-b2
\begin{physicalvignette}[Radar feed horn illuminating a corner reflector]
A second bench observation matters on Radio telescope dish with Cassegrain feed. When technicians swap only the receive probe while holding the transmit chain fixed, reciprocity checks whether the reported coupling is property of the structure or artifact of a mismatched flux gauge \cite{stratton1941}. Failed reciprocity usually traces to a near-field tap, a radome standing-wave, or a probe polarization basis that was never aligned with the transmit standard \cite{hansen1988}.
Rayleigh and Stratton treated such ranges as boundary-value laboratories long before digital near-field transforms became routine \cite{stratton1941}. Harrington's method-of-moments culture reinforced the same discipline: metal outline first, harmonic headline second \cite{harrington2001}.
The measurement logic above is made precise in Subsection~\ref{sec:ch2104} \cite{stratton1941}.
\end{physicalvignette}
\begin{physicalvignette}[Radar feed horn illuminating a corner reflector]
Consider Radar feed horn illuminating a corner reflector. Before any coefficient table is printed, technicians record frequency, cable loss, and the mechanical outline that fixes electrical size. The first walk-through of the range is always tactile: flange bolts, probe standoff, turntable boresight, and the wattmeter reference on the receive chain \cite{hansen1988}.
On this hardware the near zone is visually impressive---standing reactive structure, rapid phase rotation, sometimes a visible VSWR ripple at the feed---yet the export question is measured only after outgoing selection on a large sphere \cite{harrington2001}. What the camera sees in the chamber is not automatically what the far-zone pattern claims \cite{stratton1941}.
Engineers therefore learn to separate three readings: what the feed can synthesize on declared support, what propagating content survives on the radiative range, and what a finite instrument can recover at declared tolerance \cite{devaney2004}. Confusing those readings produces the rank gaps and watt failures that later sections quantify \cite{pal2026}.
Rayleigh and Stratton treated such ranges as boundary-value laboratories long before digital near-field transforms became routine \cite{stratton1941}. Harrington's method-of-moments culture reinforced the same discipline: metal outline first, harmonic headline second \cite{harrington2001}.
Subsection~\ref{sec:ch2104} states the formal export condition that this range scene presupposes \cite{stratton1941}.
\end{physicalvignette}
\begin{equation}
\label{eq:ch2.structured-source-state-def}
\Jimp(\mathbf{r}')
=
\sum_{n=1}^{N}
A_n\,\chi_n(\mathbf{r}')\,
e^{\jj\Phi_n(\mathbf{r}')}\,
\mathbf{p}_n(\mathbf{r}'),
\qquad
\mathbf{p}_n\cdot\hat{\mathbf{k}}_n=0,
\end{equation}
where \(A_n\ge 0\) are channel amplitudes, \(\Phi_n\) are spatial phase laws,
\(\chi_n\) are nonnegative envelopes with \(\operatorname{supp}\chi_n\subseteq\Omega_s\),
and \(\mathbf{p}_n\) are transverse polarization vectors on each channel support
\cite{stratton1941}. Equation~\eqref{eq:ch2.structured-source-state-def} is the
Subsection~\ref{sec:ch2104} operator-side label for sources whose angular organization is
specified before radiation selection: amplitude fixes channel weight, phase fixes interference
geometry, and \(\mathbf{p}_n\) fixes local vector state, paralleling the field-side chain of
Subsection~\ref{sec:ch01.2.2} \cite{stratton1941}. The sum is understood in the
distributional sense when \(\chi_n\) includes point or sheet supports
\cite{stratton1941}. Non-uniqueness of Eq.~\eqref{eq:ch2.structured-source-state-def}
under gauge-equivalent reparameterizations is inherited from impressed-current equivalence classes
Eq.~\eqref{eq:ch2.source-admissibility} and does not alter \(\mathcal{R}_\omega[\Jimp]\)
\cite{harrington2001}. Structured radiation coefficient. Define the \emph{structured radiation coefficient} of channel \(n\) in direction
\(\widehat{\mathbf{k}}\) and polarization \(\sigma\in\{1,2\}\) by
\begin{equation}
\label{eq:ch2.structured-source-radiation-coefficient}
c_{n,\sigma}(\widehat{\mathbf{k}})
=
-\jj\omega\mu_0\,\frac{1}{4\pi}
\int_{\Omega_s}
A_n\,\chi_n(\mathbf{r}')\,
e^{\jj\Phi_n(\mathbf{r}')}\,
\mathbf{p}_n(\mathbf{r}')\cdot\mathbf{e}_\sigma(\widehat{\mathbf{k}})\,
e^{-\jj k_0\widehat{\mathbf{k}}\cdot\mathbf{r}'}\,d^{3}r',
\end{equation}
with \(\mathbf{e}_\sigma\) transverse unit vectors from
Eq.~\eqref{eq:ch2.exterior-angular-eigenmode} \cite{stratton1941}. The far angular
amplitude is
\begin{equation}
\label{eq:ch2.structured-source-far-amplitude}
a_\sigma(\widehat{\mathbf{k}})
=
\sum_{n=1}^{N} c_{n,\sigma}(\widehat{\mathbf{k}}),
\qquad
a_\sigma=\Lambda_{\mathrm{nf}}[\Jimp]_\sigma,
\end{equation}
importing Eq.~\eqref{eq:ch2.near-to-far-transform-def} and consistency
Eq.~\eqref{eq:ch2.near-to-far-integral-agreement} \cite{stratton1941}.
Equations~\eqref{eq:ch2.structured-source-radiation-coefficient}--\eqref{eq:ch2.structured-source-far-amplitude}
linearize structured sources through \(\mathcal{R}_\omega\): each channel contributes a
directional radiation coefficient fixed by outgoing radiation condition, and superposition holds on the
exporting branch \cite{stratton1941}. Physically, Eq.~\eqref{eq:ch2.structured-source-far-amplitude}
is the operator restatement of ``structured beam'' labels: azimuthal phase ramps in
\(\Phi_n\) appear in \(c_{n,\sigma}\) through the plane-wave kernel, but only after radiation
selection \cite{allen1992}.
\phantomsection\label{prop:ch2.structured-source-topology-diagnostic}
Let \(\Jimp\) satisfy Eq.~\eqref{eq:ch2.structured-source-state-def} with
\((\E_{\mathrm{rad}},\H_{\mathrm{rad}})=\mathcal{R}_\omega[\Jimp]\). Transverse components,
Stokes parameters, Poynting vector, and Poynting vorticity computed from
\((\E_{\mathrm{rad}},\H_{\mathrm{rad}})\) in the far zone inherit the diagnostic chain of
Subsection~\ref{sec:ch01.2.2} Eqs.~\eqref{eq:ch1.transverse-components}--\eqref{eq:ch1.poynting-vorticity}.
Local topology markers are rotationally relevant only where
\begin{equation}
\label{eq:ch2.structured-source-topology-criterion}
\mathcal{Q}_{\mathrm{str}}(\widehat{\mathbf{k}})
=
\mathcal{Q}_{\mathrm{top}}(\mathbf{r})
\ \text{evaluated on}\ (\E_{\mathrm{rad}},\H_{\mathrm{rad}})\ \text{with}\ 
\mathcal{D}_{J,\mathrm{reg}}=1,
\end{equation}
importing Eq.~\eqref{eq:ch1.topology-qualified-gate} and regime qualification
Eq.~\eqref{eq:ch2.regime-operator-screen} \cite{stratton1941}. Nontrivial
Stokes texture, phase winding, or \(\boldsymbol{\Omega}_P\neq\mathbf{0}\) on
\(\mathcal{S}_\omega[\Jimp]\) without passing through \(\mathcal{R}_\omega\) do not qualify
for transport reports \cite{stratton1941}.
Apply the topology--transport separation (Subsection~\ref{sec:ch01.2.2}) to far-zone traces of
\(\mathcal{R}_\omega[\Jimp]\) and identify \(\mathcal{Q}_{\mathrm{top}}\) with
Eq.~\eqref{eq:ch2.transport-qualified-radiation} \cite{stratton1941}. Reactive
components removed by \(\mathcal{R}_\omega\) cannot satisfy
Eq.~\eqref{eq:ch2.structured-source-topology-criterion} \cite{harrington2001}.
Problem. A circular ring array of \(N=36\) impressed-current elements lies on
\(\rho=\rho_0=0.8a\), \(z=0\), with compact ball support \(\Omega_s=B_a(0)\), \(k_0 a=3.0\),
and \(\omega\) fixed . Each element carries amplitude
\(A_n=A_0\), envelope \(\chi_n=\delta(\mathbf{r}'-\boldsymbol{\rho}_n)\), and transverse
polarization \(\mathbf{p}_n=\hat{\boldsymbol{\phi}}\) at azimuth \(\phi_n=2\pi(n-1)/N\).
Two structured configurations are compared:
\begin{enumerate}
\item[(a)] \emph{Balanced counter-rotating pair:} half the elements carry
\(\Phi_n=\ell\phi_n\) with \(\ell=+1\), half carry \(\Phi_n=-\ell\phi_n\), matching the
channel template Eq.~\eqref{eq:ch1.ex62-counter-rotating-superposition} after radiation mapping;
\item[(b)] \emph{Imbalanced single winding:} all elements carry \(\Phi_n=+\ell\phi_n\) with
\(\ell=+1\) but amplitudes satisfy \(A_n=A_0(1+\eta\cos\phi_n)\), \(\eta=0.35\).
\end{enumerate}
Far-zone diagnostics use \(\mathcal{R}_\omega[\Jimp]\), extractor
Eq.~\eqref{eq:ch2.far-zone-extractor-def}, and shell \(r=10a\) . Tasks:
\begin{enumerate}
\item[(i)] compute structured coefficients Eq.~\eqref{eq:ch2.structured-source-radiation-coefficient}
and far amplitudes Eq.~\eqref{eq:ch2.structured-source-far-amplitude} for both configurations;
\item[(ii)] evaluate Stokes parameters and \(\boldsymbol{\Omega}_P\) on the radiative branch
(Subsection~\ref{sec:ch01.2.2}) and state whether \(\mathcal{Q}_{\mathrm{str}}=1\);
\item[(iii)] compute net angular-flux proxy \(\Lambda_J\) from
Eq.~\eqref{eq:ch2.angular-momentum-flux-vector} and compare with
Eq.~\eqref{eq:ch1.ex62-zero-transport-result} and imbalance sensitivity
Eq.~\eqref{eq:ch1.ex62-imbalance-sensitivity} (Subsections~\ref{sec:ch01.2.2}--\ref{sec:ch01.2.3}).
\end{enumerate}
\begin{figure}[t]
\centering
\ChOneFigBlock{%
\begin{tikzpicture}[ch1fig, node distance=5mm and 7mm]
 \node[ch1box] (j) {Structured $\Jimp$\\state def.};
 \node[ch1box, right=7mm of j] (r) {$\mathcal{R}_\omega$};
 \node[ch1box, right=7mm of r] (a) {$a_\sigma(\widehat{\mathbf{k}})$};
 \node[ch1box, right=7mm of a] (diag) {Stokes / $\mathbf{P}$ / $\boldsymbol{\Omega}_P$};
 \node[ch1box, right=7mm of diag] (criterion) {$\mathcal{Q}_{\mathrm{str}}$};
 \node[ch1box, below=5mm of diag] (c1) {Subsection~\ref{sec:ch01.2.2}};
 \draw[ch1flow] (j) -- (r) -- (a) -- (diag) -- (criterion);
 \draw[ch1flow, dashed] (c1.south) -- (diag.north);
\end{tikzpicture}%
}
\caption{Worked problem~\ref{sec:ch2104}: structured source parameters pass through
\(\mathcal{R}_\omega\) to far amplitudes; amplitude--phase--polarization--Poynting diagnostics
(USE Subsection~\ref{sec:ch01.2.2}) apply only on the radiative branch before topology criterion
Eq.~\eqref{eq:ch2.structured-source-topology-criterion}.}
\label{fig:ch2.wp-structured-source-radiation}
\end{figure}
The dashed
link marks imported \ref{ch:01} diagnostic definitions . Step 1 — Structured coefficients for configuration (a). Insert the counter-rotating phase law into Eq.~\eqref{eq:ch2.structured-source-radiation-coefficient}:
\begin{equation}
\label{eq:ch2.ex-str-counterprop-coefficient}
c_{n,\sigma}(\widehat{\mathbf{k}})
=
-\jj\omega\mu_0\,\frac{A_0}{4\pi}\,
e^{\jj \ell s_n \phi_n}\,
\mathbf{p}_n\cdot\mathbf{e}_\sigma(\widehat{\mathbf{k}})\,
e^{-\jj k_0\widehat{\mathbf{k}}\cdot\boldsymbol{\rho}_n},
\qquad
s_n=
\begin{cases}
+1, & n=1,\ldots,18,\\
-1, & n=19,\ldots,36,
\end{cases}
\end{equation}
with \(\ell=+1\) . Summing over the ring,
\begin{equation}
\label{eq:ch2.ex-str-far-amplitude-symmetry}
a_\sigma(\widehat{\mathbf{k}})
=
\sum_{n=1}^{36} c_{n,\sigma}(\widehat{\mathbf{k}})
\ \approx\
A_0\,J_1(k_0\rho_0\sin\theta)\,e^{\jj\ell\phi}\,f_\sigma(\theta),
\end{equation}
where \(J_1\) is the azimuthal ring factor from discrete \(\phi_n\) sampling and
\(f_\sigma\) is transverse polarization weight . Equation
\eqref{eq:ch2.ex-str-far-amplitude-symmetry} shows nonzero azimuthal phase in
\(a_\sigma(\widehat{\mathbf{k}})\): the far pattern is \emph{structured} in the sense of
Eq.~\eqref{eq:ch2.structured-source-far-amplitude} . Mathematically, equal
\(A_0\) and balanced \(s_n=\pm 1\) enforce symmetry under \(\ell\to-\ell\) channel exchange;
physically, the configuration mimics counter-rotating orbital labels on the source before export
Step 2 --- Angular-flux cancellation (Subsections~\ref{sec:ch01.2.2}--\ref{sec:ch01.2.3}). Net angular export uses radiative flux Eq.~\eqref{eq:ch2.angular-momentum-flux-vector} on
\(S_r\):
\begin{equation}
\label{eq:ch2.ex-str-transport-quotient-a}
\Lambda_J^{(\mathrm{a})}
=
\frac{\int_{S_r}\widehat{\mathbf{r}}\cdot\mathbf{M}_{\mathrm{rad}}[\Jimp^{(\mathrm{a})}]\,dS}
 {\int_{S_r}\Re\{\widehat{\mathbf{r}}\cdot\mathbf{S}_{\mathrm{rad}}[\Jimp^{(\mathrm{a})}]\}\,dS}
\ \approx\ 0,
\end{equation}
importing Eq.~\eqref{eq:ch1.ex62-zero-transport-result} after mapping balanced channels through
Eq.~\eqref{eq:ch2.structured-source-far-amplitude} \cite{stratton1941}. Channel
power amplitudes satisfy \(I_+=I_-\) in the notation of
Eq.~\eqref{eq:ch1.ex62-equal-channel-condition}, so
Eq.~\eqref{eq:ch1.ex62-net-am-transport} yields vanishing signed export despite nonzero
\(|a_\sigma(\widehat{\mathbf{k}})|\) \cite{stratton1941}. regime qualification
\(\mathcal{D}_{J,\mathrm{reg}}=1\) holds on \(S_r\) for \(k_0 r=30\), but
\(\mathcal{Q}_{\mathrm{str}}=0\) because directional angular-flux support cancels
\cite{harrington2001}. Configuration (a) therefore exhibits structured \emph{appearance} without
transport-qualified export \cite{stratton1941}.
On \(S_r\), extract transverse components Eq.~\eqref{eq:ch1.transverse-components} from
\(\E_{\mathrm{rad}}\) and form Stokes parameters Eq.~\eqref{eq:ch1.stokes-local}:
\begin{equation}
\label{eq:ch2.ex-str-stokes-contrast}
S_0^{(\mathrm{a})}(\widehat{\mathbf{k}})\approx 2A_0^2|f_\sigma(\theta)|^2,\quad
S_3^{(\mathrm{a})}(\widehat{\mathbf{k}})\approx 0,\quad
\Delta\varphi^{(\mathrm{a})}=\arg a_+-\arg a_-\ \text{varies with}\ \phi,
\end{equation}
because counter-rotating source phases produce azimuthal organization in \(a_\sigma\) while
circular Stokes component \(S_3\) averages toward zero on the balanced ring
\cite{stratton1941}. Poynting vector Eq.~\eqref{eq:ch1.poynting-average} and vorticity
Eq.~\eqref{eq:ch1.poynting-vorticity} yield
\begin{equation}
\label{eq:ch2.ex-str-poynting-vorticity}
\|\boldsymbol{\Omega}_P^{(\mathrm{a})}\|_{L^{\infty}(S_r)}
\ \sim\
\mathcal{O}(\ell\beta),
\qquad
\beta=k_0\rho_0=2.4,
\end{equation}
with \(\boldsymbol{\Omega}_P^{(\mathrm{a})}\neq\mathbf{0}\) near the pattern peak
\cite{stratton1941}. Mathematically, local circulation and phase winding can coexist
with \(\Lambda_J^{(\mathrm{a})}\approx 0\); physically, this is the operator-side restatement of
the topology--transport separation (Subsection~\ref{sec:ch01.2.2}) on structured sources
\cite{belinfante1940}. Students must not infer orbital index from
\(\|\boldsymbol{\Omega}_P\|\) alone without Eq.~\eqref{eq:ch2.structured-source-topology-criterion}
\cite{harrington2001}.
Step 4 — Imbalanced configuration (b). Amplitude modulation \(A_n=A_0(1+\eta\cos\phi_n)\) breaks channel equality:
\begin{equation}
\label{eq:ch2.ex-str-imbalanced-export}
\Lambda_J^{(\mathrm{b})}
\ \approx\
\frac{\eta\,\ell}{1+\eta^{2}/2}\,\Lambda_J^{(\mathrm{ref})},
\qquad
\Lambda_J^{(\mathrm{ref})}
=
\frac{k_0 A_0^2 J_1^2(k_0\rho_0\sin\theta_{\max})}{2},
\end{equation}
importing imbalance sensitivity Eq.~\eqref{eq:ch1.ex62-imbalance-sensitivity} with
\(\eta=0.35\) \cite{stratton1941}. Illustratively,
\(|\Lambda_J^{(\mathrm{b})}/\Lambda_J^{(\mathrm{ref})}|\approx 0.30\): a thirty-percent fractional
imbalance in source amplitude produces order-unity fractional export relative to the reference
single-winding ring \cite{stratton1941}. Far amplitude magnitude
\(|a_\sigma^{(\mathrm{b})}|\) differs from \(|a_\sigma^{(\mathrm{a})}|\) by only
\(\sim 8\%\) (illustrative) because phase structure is identical; export changes because signed channel content
is unequal Eq.~\eqref{eq:ch1.ex62-net-am-transport} \cite{stratton1941}. Here
\(\mathcal{Q}_{\mathrm{str}}=1\) on the principal lobe and
Eq.~\eqref{eq:ch1.transport-qualified-decision} returns \(\mathcal{D}_J=1\)
\cite{stratton1941}. Step 5 — Raw near map versus radiative diagnostics. Applying Stokes and Poynting formulas to \(\mathcal{S}_\omega[\Jimp]\) on \(r=1.5a\) inflates
texture metrics (illustrative ratios):
\begin{equation}
\label{eq:ch2.ex-str-near-overread}
\frac{\|\boldsymbol{\Omega}_P\|_{L^{2}(S_{1.5a})}^{(\mathrm{raw})}}
 {\|\boldsymbol{\Omega}_P\|_{L^{2}(S_r)}^{(\mathrm{rad})}}
\ \approx\ 1.9,
\qquad
\frac{S_0^{(\mathrm{raw})}}{S_0^{(\mathrm{rad})}}
\ \approx\ 2.3,
\end{equation}
by the same reactive contamination mechanism as Eq.~\eqref{eq:ch2.ex-ev-near-amplitude-ratio}
\cite{stratton1941}. Equation~\eqref{eq:ch2.ex-str-near-overread} shows that
structured-source labels read from near reactive scans overstate Poynting topology relative to
\(\mathcal{R}_\omega[\Jimp]\) \cite{stratton1941}. Correct measurement sequence: parametrize
Eq.~\eqref{eq:ch2.structured-source-state-def}, map
Eq.~\eqref{eq:ch2.structured-source-far-amplitude}, then apply §1.2.2 diagnostics on
\((\E_{\mathrm{rad}},\H_{\mathrm{rad}})\) \cite{harrington2001}.
Structured far appearance \(|a_\sigma(\widehat{\mathbf{k}})|\neq 0\) does not imply
\(\mathcal{Q}_{\mathrm{str}}=1\) Eq.~\eqref{eq:ch2.structured-source-topology-criterion}, and qualified
topology does not imply signed export \(\Lambda_J\neq 0\)
Eqs.~\eqref{eq:ch1.ex62-zero-transport-result},~\eqref{eq:ch1.ex62-imbalance-sensitivity}
\cite{stratton1941}. Configuration (a) occupies the first failure;
configuration (b) satisfies export once imbalance is present. Non-radiating additions
Subsection~\ref{sec:ch2102} and evanescent storage Subsection~\ref{sec:ch2103} modify near
diagnostics without changing \(a_\sigma\) when \(\mathcal{R}_\omega[\Jimp]\) is fixed
\cite{stratton1941}. Accessibility budget
Eq.~\eqref{eq:ch2.accessible-angular-budget} may further shrink observed structure inside
\(\mathcal{K}_{\mathrm{acc}}^{\mathrm{ff}}\) \cite{stratton1941}.
Physical meaning (summary). Subsection~\ref{sec:ch2104} establishes that ``structured source states'' are operator objects
Eq.~\eqref{eq:ch2.structured-source-state-def} mapped to far coefficients
Eq.~\eqref{eq:ch2.structured-source-far-amplitude}, not merely descriptive beam names
\cite{stratton1941}. Balanced counter-rotating rings produce azimuthal
phase in \(a_\sigma(\widehat{\mathbf{k}})\) Eq.~\eqref{eq:ch2.ex-str-far-amplitude-symmetry} yet
vanishing \(\Lambda_J\) Eq.~\eqref{eq:ch2.ex-str-transport-quotient-a}, reproducing
\ref{ch:01} phase-without-transport logic on the radiation operator
\cite{stratton1941}. Stokes and Poynting textures
Eqs.~\eqref{eq:ch2.ex-str-stokes-contrast}--\eqref{eq:ch2.ex-str-poynting-vorticity} can be
rich while export cancels; amplitude imbalance Eq.~\eqref{eq:ch2.ex-str-imbalanced-export} is the
control knob for signed angular flux \cite{stratton1941}. structured-source topology diagnostics (Subsection~\ref{sec:ch210}) and criterion
Eq.~\eqref{eq:ch2.structured-source-topology-criterion} require diagnostics on
\(\mathcal{R}_\omega[\Jimp]\) with regime qualification before any OAM or helicity label is
promoted to transport \cite{harrington2001}. Near reactive overread
Eq.~\eqref{eq:ch2.ex-str-near-overread} parallels Subsection~\ref{sec:ch2103}: structure in
\(\mathcal{S}_\omega\) is not structure in export \cite{stratton1941}.
Structured sources close the example suite by showing that the most elaborate near-zone organization does not override selection: phase winding in $a_\sigma(\widehat{\mathbf{k}})$ can coexist with vanishing signed angular flux until amplitude imbalance restores transport on the radiative branch \cite{belinfante1940}. OAM or vortex labels attached to impressed currents remain hypotheses until $\mathcal{R}_\omega$, regime qualification, and topology criteria from \ref{ch:01} have been applied \cite{harrington2001}. Section~\ref{sec:ch210} therefore ends with four distinct bench mechanisms; Section~\ref{sec:ch02.11} collects the chapter results and states what \ref{ch:03} will relabel Maxwell or outgoing foundations proved here.
The four worked problems in Subsections~\ref{sec:ch2101}--\ref{sec:ch2104} are not independent
demonstrations but coordinated activations of Eq.~\eqref{eq:ch2.dof-partition-master}
\cite{stratton1941}. Compact truncation
(Subsection~\ref{sec:ch2101}) compresses spectral rank without creating NR nullity or evanescent
sheets; NR superposition (Subsection~\ref{sec:ch2102}) annihilates export through
\(\ker\mathcal{R}_\omega\) while leaving near signatures accessible
\cite{bleszynski1996}; evanescent storage (Subsection~\ref{sec:ch2103}) inflates near
reactive readouts without altering \(a(\widehat{\mathbf{k}})\) when the radiation image is fixed
\cite{mandel1995}; structured sources (Subsection~\ref{sec:ch2104}) separate far
appearance, topology qualification, and signed transport on the radiative branch
\cite{belinfante1940}. Order-of-magnitude estimates in Section~\ref{sec:ch210} quantify
order-of-magnitude separations only; they must not be read as reported measurements
\cite{hansen1988}. Qualified reporting always applies \(\mathcal{R}_\omega\) (equivalently
\(\Pi_{\mathrm{rad}}\) after Subsection~\ref{sec:ch224}) before Stokes, Poynting, or angular-flux functionals, and
composes with regime qualification Eq.~\eqref{eq:ch2.regime-operator-screen} inherited from
\ref{ch:01} \cite{stratton1941}. Section~\ref{sec:ch02.11} closes the chapter with the radiation-selection conclusions developed above \cite{stratton1941}.
The electrically small dipole of Subsection~\ref{sec:ch2101} illustrates spectral compression: a finite probe bank on $S_b$ resolves at most a rank-limited subset of $\mathcal{R}_\omega[\Jimp]$, and truncation error tracks singular values of the compact Green part rather than membership in $\ker\mathcal{R}_\omega$ \cite{harrington2001}. High-index modes below the probe-rank threshold are lost to finite-geometry observability, not to source nullity; increasing modal order without raising $N$ or $b/a$ may shrink a fitting residual that owes nothing to physics \cite{hansen1988}. The loop-superposition problem of Subsection~\ref{sec:ch2102} illustrates the complementary source-side mechanism: $\Jimp^{\mathrm{(NR)}}\in\ker\mathcal{R}_\omega$ leaves $a(\widehat{\mathbf{k}})$ and qualified far-zone flux unchanged while near probes register superposed reactive content \cite{devaney2004}. Far data fix $\mathcal{R}_\omega[\Jimp]$ but not $\Jimp$; existence of near fields does not imply flux report of their source detail \cite{stratton1941}. The evanescent-overlap configuration of Subsection~\ref{sec:ch2103} shows that reactive storage can inflate near-shell amplitude by more than a factor of two while $a(\widehat{\mathbf{k}})$ remains fixed when the radiation image is held constant; applying far extractors to $\mathcal{S}_\omega[\Jimp]$ without prior $\mathcal{R}_\omega$ selection can bias angular export by order forty percent on illustrative geometries \cite{hansen1988}. Near richness is not far-zone angular structure. The structured ring array of Subsection~\ref{sec:ch2104} separates far appearance, topology qualification, and signed transport: azimuthal organization in $a_\sigma(\widehat{\mathbf{k}})$ can coexist with vanishing $\Lambda_J$ until amplitude imbalance restores export on the qualified radiative branch \cite{belinfante1940}. Structured appearance is not export; radiation operators select transport-capable content \cite{stratton1941}.
\section{Summary}
\label{sec:ch02.11}
% CH02-NARR-EXP:02.11
Maxwell evolution generates fields; outgoing selection determines export; finite instruments recover only part of that export \cite{maxwell1873,stratton1941,sommerfeld1949}. Vector spherical harmonics, source realizability, and interaction formalism beyond this chapter are cited in Chapters~3--4 without local re-derivation here \cite{harrington2001,devaney2004}.
\ref{ch:02} began from a premise that Maxwell's equations, once supplemented by
outgoing radiation conditions, do not merely \emph{propagate} fields---they \emph{select} which
field content qualifies as radiative export. That selection is implemented by the radiation
operator $\mathcal{R}_\omega$, whose image is what finite instruments ultimately interrogate,
always incompletely and always through additional rank limits imposed by geometry and hardware.
The distinction is not pedantic. A near-field scanner that records thirty-six complex samples
has performed a serious measurement; yet the number answers a question about accessible angular
structure only after one traces which term in the master degree-of-freedom partition removed
the unresolved content---source nullity, aperture rank, evanescent storage, spectral compression,
or geometric suppression.
% PHYSICAL-MEMORY:sec:ch02.11
\begin{physicalvignette}[Half-space ground plane with image currents]
Consider Half-space ground plane with image currents. Before any coefficient table is printed, technicians record frequency, cable loss, and the mechanical outline that fixes electrical size. The first walk-through of the range is always tactile: flange bolts, probe standoff, turntable boresight, and the wattmeter reference on the receive chain \cite{hansen1988}.
On this hardware the near zone is visually impressive---standing reactive structure, rapid phase rotation, sometimes a visible VSWR ripple at the feed---yet the export question is measured only after outgoing selection on a large sphere \cite{harrington2001}. What the camera sees in the chamber is not automatically what the far-zone pattern claims \cite{stratton1941}.
Engineers therefore learn to separate three readings: what the feed can synthesize on declared support, what propagating content survives on the radiative range, and what a finite instrument can recover at declared tolerance \cite{devaney2004}. Confusing those readings produces the rank gaps and watt failures that later sections quantify \cite{pal2026}.
Rayleigh and Stratton treated such ranges as boundary-value laboratories long before digital near-field transforms became routine \cite{stratton1941}. Harrington's method-of-moments culture reinforced the same discipline: metal outline first, harmonic headline second \cite{harrington2001}.
Subsection~\ref{sec:ch02.11} states the formal export condition that this range scene presupposes \cite{stratton1941}.
\end{physicalvignette}
The four worked configurations in Section~\ref{sec:ch210} were composed to make that partition
tangible. The electrically small dipole showed that modal truncation error tracks compact singular
values, not non-radiating nullity. The loop-superposition problem showed that far patterns can
remain invariant while near probes register a thirty-five percent increment in tangential
amplitude. The evanescent-overlap problem showed that near-shell energy can exceed the
radiative component by more than a factor of two while far amplitude $a(\widehat{\mathbf{k}})$
remains fixed, and that applying far-field extractors to raw $\mathcal{S}_\omega[\Jimp]$ can
bias angular export by forty percent. The structured ring array showed that azimuthal phase
organization in $a_\sigma(\widehat{\mathbf{k}})$ can coexist with vanishing signed angular flux
until a modest amplitude imbalance restores transport. Each result is a laboratory caution stated
in operator language: existence of field does not imply export; export does not imply
wattmeter reading; $\oint_{\Gamma_R}$ closure does not imply unique source identification
\cite{maxwell1873}.
Read against that laboratory background, the chapter's central lesson is a completion criterion,
not a slogan. \ref{ch:01} established $\mathcal{D}_J=1$---what impressed currents and phasor
fields may enter the monograph at all. \ref{ch:02} establishes $\mathcal{E}_{\mathrm{acc}}^{(N)}$: which
radiative content survives outgoing selection and which parts finite geometry and finite
instruments can resolve. Throughout, qualified reporting applies $\mathcal{R}_\omega$ before
Stokes parameters, Poynting vectors, or angular-flux functionals, and composes with regime
qualification inherited from \ref{ch:01} so that reactive shells are not mistaken for
transport \cite{harrington2001}. From the vector-network-analyzer reciprocity checks of Section~\ref{sec:ch28} through the directional
decompositions of Section~\ref{sec:ch29} to the four worked configurations of Section~\ref{sec:ch210},
the chapter supplies operator criteria for diagnosing pattern residuals by mechanism rather than by fitting alone.
\subsection{Results}
\phantomsection\label{sec:ch02111}
Principle~\ref{prin:ch2.radiative-accessibility} selects $\mathcal{F}_{\mathrm{rad}}$ before harmonic readouts; probe rank caps what wattmeters report \cite{stratton1941}. The results proved here may be collected without repeating the subsection chain. Radiation is a
selected subset of Maxwell solutions on exterior domains, not an independent physical process.
The radiation operator possesses nontrivial null spaces whose elements are physically realizable
and invisible to far data. Compact sources and finite observation compress accessible angular
structure through spectral rank limits that are distinct from nullity. Near-to-far export requires
radiation selection before asymptotic extraction; evanescent and non-radiating admixtures can
inflate near readouts without altering far amplitude when $\mathcal{R}_\omega[\Jimp]$ is fixed.
Structured source parameters specify appearance before selection; signed angular transport requires
imbalance on the radiative branch after selection. Reciprocity constrains bilinear coupling on
$\operatorname{ran}\mathcal{R}_\omega$ but does not enlarge accessible angular sectors or remove
finite-rank observation deficits.
% PHYSICAL-MEMORY:sec:ch02111
\begin{physicalvignette}[Rectangular horn mouth in a ground-plane flange]
Consider Rectangular horn mouth in a ground-plane flange. Before any coefficient table is printed, technicians record frequency, cable loss, and the mechanical outline that fixes electrical size. The first walk-through of the range is always tactile: flange bolts, probe standoff, turntable boresight, and the wattmeter reference on the receive chain \cite{hansen1988}.
On this hardware the near zone is visually impressive---standing reactive structure, rapid phase rotation, sometimes a visible VSWR ripple at the feed---yet the export question is measured only after outgoing selection on a large sphere \cite{harrington2001}. What the camera sees in the chamber is not automatically what the far-zone pattern claims \cite{stratton1941}.
Engineers therefore learn to separate three readings: what the feed can synthesize on declared support, what propagating content survives on the radiative range, and what a finite instrument can recover at declared tolerance \cite{devaney2004}. Confusing those readings produces the rank gaps and watt failures that later sections quantify \cite{pal2026}.
Rayleigh and Stratton treated such ranges as boundary-value laboratories long before digital near-field transforms became routine \cite{stratton1941}. Harrington's method-of-moments culture reinforced the same discipline: metal outline first, harmonic headline second \cite{harrington2001}.
Subsection~\ref{sec:ch02111} states the formal export condition that this range scene presupposes \cite{stratton1941}.
\end{physicalvignette}
The Chapter~\ref{ch:02} admissibility completion requires Chapter~1 admissibility, outgoing-qualified
$\mathcal{R}_\omega[\Jimp]$, and regime qualification $\mathcal{D}_{J,\mathrm{reg}}=1$
\cite{stratton1941}. Results proved under that completion inherit report classes from
Eq.~\eqref{eq:ch1.report-class-lock} \cite{stratton1941}.
\medskip
\begin{equation}
\label{eq:ch2.continuity-admissibility}
\mathcal{D}_{\mathrm{ch2}}=1
\iff
\left(
\mathcal{D}_{\mathrm{ch1}}=1\
\wedge\
\mathcal{R}_\omega[\Jimp]\ \text{is outgoing-qualified}\
\wedge\
\mathcal{D}_{J,\mathrm{reg}}=1
\right),
\end{equation}
with \(\mathcal{D}_{\mathrm{ch1}}\) from Eq.~\eqref{eq:ch1.continuity-admissibility} and
\(\mathcal{D}_{J,\mathrm{reg}}\) from Eq.~\eqref{eq:ch1.regime-qualified-decision}.
Let \(\mathcal{R}_{\mathrm{ch2}}\) denote results proved in
Sections~\ref{sec:ch20}--\ref{sec:ch210} under \(\mathcal{D}_{\mathrm{ch2}}=1\).
Report classes remain those of Eq.~\eqref{eq:ch1.report-class-lock}
\cite{stratton1941}.
\subsection{Topics developed in later chapters}
\phantomsection\label{sec:ch02112}
Topics requiring vector-spherical-harmonic completion, source realizability proofs, or interaction theory beyond this monograph are cited in later chapters without local re-derivation here \cite{stratton1941}. Harmonic labels and rank bounds arrive in Chapters~3--4 under the completion recorded in Eq.~\eqref{eq:ch2.continuity-admissibility} \cite{harrington2001}.
\medskip
% PHYSICAL-MEMORY:sec:ch02112
\begin{physicalvignette}[Open-ended coax probe above a calibration table]
Consider Open-ended coax probe above a calibration table. Before any coefficient table is printed, technicians record frequency, cable loss, and the mechanical outline that fixes electrical size. The first walk-through of the range is always tactile: flange bolts, probe standoff, turntable boresight, and the wattmeter reference on the receive chain \cite{hansen1988}.
On this hardware the near zone is visually impressive---standing reactive structure, rapid phase rotation, sometimes a visible VSWR ripple at the feed---yet the export question is measured only after outgoing selection on a large sphere \cite{harrington2001}. What the camera sees in the chamber is not automatically what the far-zone pattern claims \cite{stratton1941}.
Engineers therefore learn to separate three readings: what the feed can synthesize on declared support, what propagating content survives on the radiative range, and what a finite instrument can recover at declared tolerance \cite{devaney2004}. Confusing those readings produces the rank gaps and watt failures that later sections quantify \cite{pal2026}.
Rayleigh and Stratton treated such ranges as boundary-value laboratories long before digital near-field transforms became routine \cite{stratton1941}. Harrington's method-of-moments culture reinforced the same discipline: metal outline first, harmonic headline second \cite{harrington2001}.
Subsection~\ref{sec:ch02112} states the formal export condition that this range scene presupposes \cite{stratton1941}.
\end{physicalvignette}
\begin{equation}
\label{eq:ch2.defer-topic-rule}
\begin{aligned}
\mathcal{X}\in\mathcal{D}_{\mathrm{defer}}^{(\mathrm{ch2})}
&\iff
\mathcal{X}\ \text{requires VSH or }\mathbf{J}\ \text{operator completion from Ch.3} \\
&\ \vee\ \mathcal{X}\ \text{requires source realizability proof from Ch.4} \\
&\ \vee\ \mathcal{X}\ \text{requires Volume II interaction theory},
\end{aligned}
\end{equation}
and
\begin{equation}
\label{eq:ch2.defer-reporting-lock}
\mathcal{X}\in\mathcal{D}_{\mathrm{defer}}^{(\mathrm{ch2})}
\Longrightarrow
\text{cite in later chapters; no local re-derivation in Ch.~2}.
\end{equation}
\cite{stratton1941}.
\subsection{Relation to later chapters}
\phantomsection\label{sec:ch02113}
Chapter~\ref{ch:03} will organize the surviving export in rotational coordinates; Chapter~\ref{ch:04}
will confront realizability limits on sources that produce it \cite{harrington2001}. Vector spherical harmonics,
$\mathbf{J}$-operator completion, source realizability proofs, and interaction theory beyond
this volume are developed in Chapters~3--4 by citation---not by local relabeling of modes proved here
\cite{stratton1941}. Chapter~3 adds harmonic coordinates on $\mathcal{R}_\omega$ Maxwell or outgoing foundations proved in Sections~\ref{sec:ch20}--\ref{sec:ch210} \cite{stratton1941}. The continuity compatibility condition requires that new labels introduce no recursive reproof of $\mathcal{R}_{\mathrm{ch2}}$ \cite{stratton1941}.
\medskip
% PHYSICAL-MEMORY:sec:ch02113
\begin{physicalvignette}[Rotating scatterer on a lighthouse-style beam chopper]
Consider Rotating scatterer on a lighthouse-style beam chopper. Before any coefficient table is printed, technicians record frequency, cable loss, and the mechanical outline that fixes electrical size. The first walk-through of the range is always tactile: flange bolts, probe standoff, turntable boresight, and the wattmeter reference on the receive chain \cite{hansen1988}.
On this hardware the near zone is visually impressive---standing reactive structure, rapid phase rotation, sometimes a visible VSWR ripple at the feed---yet the export question is measured only after outgoing selection on a large sphere \cite{harrington2001}. What the camera sees in the chamber is not automatically what the far-zone pattern claims \cite{stratton1941}.
Engineers therefore learn to separate three readings: what the feed can synthesize on declared support, what propagating content survives on the radiative range, and what a finite instrument can recover at declared tolerance \cite{devaney2004}. Confusing those readings produces the rank gaps and watt failures that later sections quantify \cite{pal2026}.
Rayleigh and Stratton treated such ranges as boundary-value laboratories long before digital near-field transforms became routine \cite{stratton1941}. Harrington's method-of-moments culture reinforced the same discipline: metal outline first, harmonic headline second \cite{harrington2001}.
Subsection~\ref{sec:ch02113} states the formal export condition that this range scene presupposes \cite{stratton1941}.
\end{physicalvignette}
\begin{equation}
\label{eq:ch2.continuity-map}
\mathcal{D}_{\mathrm{cont}}^{(2\to 3)}(\mathcal{Y})=1
\iff
\left(
\mathcal{Y}\ \text{is compatible with}\ \mathcal{R}_{\mathrm{ch2}}
\ \wedge\
\mathcal{Y}\ \text{adds no recursive proof of}\ \mathcal{R}_{\mathrm{ch2}}
\right),
\end{equation}
subject to Eq.~\eqref{eq:ch2.continuity-admissibility} and
Eq.~\eqref{eq:ch1.compatibility-criteria} \cite{stratton1941}.

Whoever leaves Chapter~\ref{ch:02} with the operator maps $\mathcal{S}_\omega$, $\mathcal{R}_\omega$, and the angular-accessibility filters of Section~\ref{sec:ch29} in hand has the radiation-selection skeleton on which Chapter~3 will hang rotational eigenstructure \cite{harrington2001}. Maxwell equations define a source-to-field operator on bounded domains \cite{maxwell1873}; radiation operators possess nontrivial null spaces and reactive admixture on finite supports \cite{devaney2004}; finite geometry produces spectral compression of accessible angular structure \cite{chu1948}; measurement observes only a reduced image of source space through finite rank maps \cite{hansen1988}. Those four laboratory facts, stated here in operator language, are the permanent content Chapter~3 will relabel, not reopen.
\bibliographystyle{IEEEtran}
\bibliography{references}
